% ===========================================================================
%  MANUEL COMPLET : JavaScript - PHP - Java
%  Document d'auto-apprentissage, avec pratique, outils et liens.
%
%  Compilation (une seule commande suffit) :
%     pdflatex main.tex
%     pdflatex main.tex        <- 2e passage : table des matieres + ref croisees
%  ou, mieux :
%     latexmk -pdf -interaction=nonstopmode main.tex
%  ou, sans TeX Live installe :
%     tectonic main.tex
%
%  Si un paquet manque (erreur "File `xxx.sty' not found"), installe-le :
%     sudo apt install texlive-latex-extra texlive-fonts-extra texlive-lang-french
%     tlmgr install libertine inconsolata cleveref titletoc menu
% ===========================================================================

\documentclass[12pt,a4paper,twoside,openright]{book}

% ------------------------------ Encodage / langue --------------------------
\usepackage[T1]{fontenc}
\usepackage[utf8]{inputenc}
\usepackage[french]{babel}
\frenchsetup{SpaceGermandbls=no}

% ------------------------------ Mise en page ------------------------------
\usepackage[left=2.1cm,right=2.1cm,top=2.3cm,bottom=2.4cm,headheight=15pt]{geometry}
\usepackage{microtype}
\usepackage{setspace}
\onehalfspacing

% ------------------------------ Polices -----------------------------------
\usepackage[default,semibold]{libertine}   \% Tex Gyre Pagella-like, tres lisible
\usepackage[scaled=0.92]{inconsolata}     \% police a chasse fixe (code)

% ------------------------------ Liens / navigation ------------------------
\usepackage{xcolor}
\definecolor{linkblue}{HTML}{0B4F8A}
\definecolor{pathgreen}{HTML}{1F6B3A}
\hypersetup{
  colorlinks=true,
  linkcolor=linkblue,
  filecolor=linkblue,
  urlcolor=linkblue,
  citecolor=linkblue,
  pdftitle={Manuel complet : JavaScript, PHP et Java},
  pdfauthor={Manuel d'auto-apprentissage},
  pdfsubject={Programmation : JavaScript, PHP, Java},
  pdfkeywords={JavaScript, Node.js, PHP, MySQL, Java, Spring Boot, tutoriel},
  bookmarksnumbered=true,
  bookmarksopen=true,
  unicode=true
}
\usepackage[capitalise,nameinlink]{cleveref}

% ------------------------------ Tableurs / listes -------------------------
\usepackage{graphicx}
\usepackage{booktabs}
\usepackage{longtable}
\usepackage{tabularx}
\usepackage{ltablex}
\usepackage{array}
\usepackage{enumitem}
\setlist[itemize]{topsep=3pt,itemsep=2pt,parsep=0pt,leftmargin=1.5em}
\setlist[enumerate]{topsep=3pt,itemsep=2pt,parsep=0pt,leftmargin=1.6em}
\setlist[description]{topsep=3pt,itemsep=3pt,parsep=0pt,font=\bfseries}
\usepackage{multicol}
\usepackage{menukeys}   \% pour \keys{Ctrl + S}
\usepackage{titletoc}
\usepackage[totoc,columns=2]{imakeidx}
\makeindex[intoc]
\titlecontents{section}[3.2em]{\addvspace{2pt}\small}
  {\contentslabel{2.0em}}{}{\titlerule*[0.75pc]{.}\contentspage}

% ------------------------------ Couleurs du theme -------------------------
\definecolor{codebg}{HTML}{F6F8FA}
\definecolor{codeframe}{HTML}{D0D7DE}
\definecolor{ctermbg}{HTML}{1B2430}
\definecolor{ctermfg}{HTML}{E6EDF3}
\definecolor{codekw}{HTML}{0550AE}
\definecolor{codecom}{HTML}{5E7C40}
\definecolor{codestr}{HTML}{A31515}
\definecolor{codevar}{HTML}{7A3E9D}
\definecolor{codefn}{HTML}{6639BA}
\definecolor{themecolor}{HTML}{0B4F8A}
\definecolor{jsyellow}{HTML}{F0DB4F}
\definecolor{phppurple}{HTML}{777BB4}
\definecolor{javared}{HTML}{E76F00}
\definecolor{boxtitle}{HTML}{0B4F8A}
\definecolor{objbg}{HTML}{EAF3FB}\definecolor{objfr}{HTML}{8FC1E8}
\definecolor{notebg}{HTML}{EFF6EE}\definecolor{notefr}{HTML}{A7C9A0}
\definecolor{tipbg}{HTML}{FFF6DE}\definecolor{tipfr}{HTML}{E3C06A}
\definecolor{warnbg}{HTML}{FDECEC}\definecolor{warnfr}{HTML}{E6A0A0}
\definecolor{exobg}{HTML}{F1EEFB}\definecolor{exofr}{HTML}{B7A6E0}
\definecolor{linksbg}{HTML}{EEF3F7}\definecolor{linksfr}{HTML}{9BB8CE}
\definecolor{theadbg}{HTML}{0B4F8A}
\definecolor{tralight}{HTML}{F1F5F9}

% ------------------------------ Code (listings) ---------------------------
\usepackage{listings}
\lstdefinestyle{base}{
  basicstyle=\ttfamily\small,
  columns=fullflexible,
  keepspaces=true,
  breaklines=true,
  breakatwhitespace=true,
  postbreak=\mbox{\textcolor{themecolor}{\$\hookrightarrow\$}\space},
  showstringspaces=false,
  tabsize=2,
  numbers=none,
  frame=none,
  xleftmargin=0pt,
  xrightmargin=0pt,
  aboveskip=2pt,
  belowskip=2pt,
  commentstyle=\color{codecom},
  keywordstyle=\color{codekw}\bfseries,
  stringstyle=\color{codestr},
  identifierstyle=\color{black},
  morecomment=[l][\color{codecom}]{\#},
  captionpos=b
}
\lstdefinestyle{jscode}{style=base,language=JavaScript,morekeywords={async,await,of,as,static,let,const,yield,typeof,instanceof,void,delete,in,do,export,import,from,catch,finally,throw,undefined,NaN,Infinity,extends,super,get,set,Symbol,BigInt,arguments,debugger,with,switch,case,default,continue,break,function,class,return,try,new,this,null,true,false,if,else,while,for,var,=>}}
\lstdefinestyle{phpcode}{style=base,language=[5.6]PHP,morekeywords={function,class,interface,trait,enum,abstract,final,extends,implements,namespace,use,const,public,protected,private,static,readonly,match,fn,global,instanceof,insteadof,as,print,unset,isset,echo,require,require\_once,include,include\_once,declare,endfor,endforeach,endif,while,do,switch,case,default,break,continue,catch,finally,throw,try,new,clone,this,null,true,false,yield,self,parent,\_\_DIR\_\_,\_\_FILE\_\_,\_\_LINE\_\_}}
\lstdefinestyle{javacode}{style=base,language=Java,morekeywords={var,record,sealed,permits,non-sealed,yield,textblock,instanceof,enum,interface,class,abstract,final,static,public,protected,private,synchronized,volatile,transient,native,strictfp,implements,extends,package,import,throws,throw,try,catch,finally,super,this,null,true,false,assert,switch,default,case,do,while,for,if,else,new,return,boolean,char,byte,short,int,long,float,double,void},literate={<}{{\color{black}<}}1{>}{{\color{black}>}}1}
\lstdefinestyle{confcode}{style=base,morecomment=[l][\color{codecom}]{\#},morecomment=[l]{//},morecomment=[l]{;}}
\lstdefinestyle{sqlcode}{style=base,language=SQL,morekeywords={SELECT,INSERT,UPDATE,DELETE,CREATE,ALTER,DROP,TABLE,INDEX,VIEW,TRIGGER,PROCEDURE,FUNCTION,PRIMARY,KEY,FOREIGN,REFERENCES,UNIQUE,NOT,NULL,DEFAULT,AUTO\_INCREMENT,ENGINE,VARCHAR,INT,BIGINT,TEXT,BLOB,DATE,DATETIME,TIMESTAMP,DECIMAL,BOOLEAN,JOIN,INNER,LEFT,RIGHT,FULL,OUTER,ON,AS,WHERE,GROUP,ORDER,BY,HAVING,LIMIT,OFFSET,DISTINCT,AND,OR,IN,LIKE,BETWEEN,IS,COUNT,SUM,AVG,MIN,MAX,IF,EXISTS,UNION,ALL,VALUES,SET,TRANSACTION,COMMIT,ROLLBACK,GRANT,DESC,ASC,CASCADE,CONSTRAINT,CHECK,begin,end,show,use,database,from}}
\lstdefinestyle{shellcode}{style=base,language=bash,morecomment=[l]\#}
\lstdefinestyle{jsoncode}{style=base,morecomment=[l]{//}}
\lstdefinestyle{ymlcode}{style=base,morecomment=[l]\#}

\lstset{style=base}

% ------------------------------ Bo\^ites (tcolorbox) -------------------------
\usepackage[most]{tcolorbox}
\tcbuselibrary{listings,breakable,skins}

\tcbset{
  codestyle/.style={
    enhanced, breakable, colback=codebg, colframe=codeframe,
    boxrule=0.4pt, arc=1.5mm, left=2mm,right=2mm,top=1.6mm,bottom=1.6mm,
    boxsep=0pt, fontupper=\ttfamily\small, listing only
  }
}

% -- Blocs de code avec coloration syntaxique --
\newtcblisting{codejs}[1][]{listing options={style=jscode,\#1},codestyle}
\newtcblisting{codephp}[1][]{listing options={style=phpcode,\#1},codestyle}
\newtcblisting{codejava}[1][]{listing options={style=javacode,\#1},codestyle}
\newtcblisting{codesql}[1][]{listing options={style=sqlcode,\#1},codestyle}
\newtcblisting{codeconf}[1][]{listing options={style=confcode,\#1},codestyle}
\newtcblisting{codejson}[1][]{listing options={style=jsoncode,\#1},codestyle}
\newtcblisting{codeyml}[1][]{listing options={style=ymlcode,\#1},codestyle}
\newtcblisting{codeplain}[1][]{listing options={style=base,\#1},codestyle}
% version "listing only" : le contenu n'est PAS verbatim, utile pour
% colorer un mot. On prefere les versions verbatim ci-dessus.

% -- Terminal : commandes a taper --
\newtcblisting{terminal}{
  enhanced, breakable, colback=ctermbg, colframe=ctermbg,
  arc=1.5mm, boxrule=0pt, left=2.5mm,right=2mm,top=1.8mm,bottom=1.8mm,
  boxsep=0pt, fontupper=\ttfamily\small\color{ctermfg},
  listing only, listing options={style=shellcode,basicstyle=\ttfamily\small\color{ctermfg},keywordstyle=\color{ctermfg}\bfseries,commentstyle=\color{ctermfg}\itshape}
}

% -- Encarts pedagogiques --
\newtcbox{\srcbox}{on line, colback=tralight, colframe=tralight, boxrule=0pt,
  arc=0.8mm, left=1.2mm,right=1.2mm,top=0.5mm,bottom=0.7mm, size=tight,
  fontupper=\scriptsize}

\newtcolorbox{objectif}[1][]{enhanced, breakable, colback=objbg, colframe=objfr,
  boxrule=0.5pt, arc=1.5mm, left=3mm,right=3mm,top=2mm,bottom=2mm,
  fonttitle=\bfseries, coltitle=white, colbacktitle=boxtitle,
  title={Objectifs de la chapitre},\#1}

\newtcolorbox{note}{enhanced, breakable, colback=notebg, colframe=notefr,
  boxrule=0.5pt, arc=1.5mm, left=3mm,right=3mm,top=2mm,bottom=2mm,
  before upper={\small\textbf{Note.}\enspace}}

\newtcolorbox{astuce}{enhanced, breakable, colback=tipbg, colframe=tipfr,
  boxrule=0.5pt, arc=1.5mm, left=3mm,right=3mm,top=2mm,bottom=2mm,
  before upper={\small\textbf{Astuce de pro.}\enspace}}

\newtcolorbox{piege}{enhanced, breakable, colback=warnbg, colframe=warnfr,
  boxrule=0.5pt, arc=1.5mm, left=3mm,right=3mm,top=2mm,bottom=2mm,
  before upper={\small\textbf{Pi\`ege \`a \'eviter.}\enspace}}

\newtcolorbox{exercice}[1][]{enhanced, breakable, colback=exobg, colframe=exofr,
  boxrule=0.5pt, arc=1.5mm, left=3mm,right=3mm,top=2mm,bottom=2mm,
  fonttitle=\bfseries, coltitle=white, colbacktitle=boxtitle,
  title={Exercices guid\'es},\#1}

\newtcolorbox{ressources}[1][]{enhanced, breakable, colback=linksbg, colframe=linksfr,
  boxrule=0.5pt, arc=1.5mm, left=3mm,right=3mm,top=2mm,bottom=2mm,
  fonttitle=\bfseries, coltitle=white, colbacktitle=boxtitle,
  title={Ressources et outils cit\'es},before upper=\small,\#1}

\newtcolorbox{projet}[1][]{enhanced, breakable, colback=tipbg, colframe=tipfr,
  boxrule=0.5pt, arc=1.5mm, left=3mm,right=3mm,top=2mm,bottom=2mm,
  fonttitle=\bfseries, coltitle=white, colbacktitle=javared,
  title={Projet \`a r\'ealiser},\#1}

% ------------------------------ Macros pratiques ---------------------------
\newcommand{\fichier}[1]{\texttt{\color{pathgreen}\#1}}
\newcommand{\rep}[1]{\texttt{\color{pathgreen}\#1/}}
\newcommand{\cmd}[1]{\srcbox{\texttt{\detokenize{\#1}}}}
\newcommand{\link}[2]{\href{#2}{\#1}}
\newcommand{\src}[1]{\href{#1}{\textcolor{linkblue}{\$^[\ast]\$}}}
\newcommand{\langjs}{\textcolor{javared}{\textbf{JavaScript}}}
\newcommand{\langphp}{\textcolor{phppurple}{\textbf{PHP}}}
\newcommand{\langjava}{\textcolor{javared}{\textbf{Java}}}
\newcommand{\etape}[1]{\par\medskip\noindent\textbf{\large \#1}\par\smallskip}
\DeclareRobustCommand{\clic}[1]{\textcolor{linkblue}{\textsf{\#1}}}
% glyphes utilises dans les listings (arborescences, symbole euro) : encodage T1 -> on emprunte pifont
\usepackage{pifont}
\DeclareRobustCommand{\boxv}{\text{\ding{124}}}
\DeclareRobustCommand{\boxh}{\text{\ding{141}}}
\DeclareRobustCommand{\boxe}{\text{\ding{122}}}
\DeclareRobustCommand{\euro}{\texteuro}
\newcommand{\todo}[1]{\textcolor{javared}{\textbf{[TODO: \#1]}}}

% tableaux : en-tetes colorees
\usepackage{etoolbox}
% tableaux zebres : tabz = tabularx pleine largeur + lignes alternees + filets
\newenvironment{tabz}[1]{\rowcolors{2}{white}{tralight}\%
  \small\setlength{\tabcolsep}{5pt}\%
  \begin{tabularx}{\textwidth}{@{}\#1@{}}\toprule}
  {\bottomrule\end{tabularx}}
\newcommand{\thd}[1]{\textbf{\color{white}\#1}}
\newcommand{\headrow}{\rowcolor{theadbg}}
\makeatletter
\newenvironment{tablez}{\@ifnextchar[{\@tablezwith}{\@tablezwith[XXX]}}{\end{tabz}}
\newcommand{\@tablezwith}[1]{\begin{tabz}{\#1}}
\makeatother

% ------------------------------ Doc ---------------------------------------
\begin{document}
\frontmatter

% ------------------------------ Page de titre ------------------------------
\begin{titlepage}
  \thispagestyle{empty}
  \centering
  \vspace*{0.6cm}
  {\color{themecolor}\rule[-2pt]{0.86\textwidth}{1.5pt}}\\[6pt]
  {\color{javared}\rule[-2pt]{0.86\textwidth}{1.5pt}}\\[6pt]
  {\color{phppurple}\rule[-2pt]{0.86\textwidth}{1.5pt}}\\[26pt]
  {\Huge\bfseries Manuel complet}\\[6pt]
  {\large pour apprendre, pratiquer et ma\^itriser}\\[18pt]
  {\fontsize{30}{34}\selectfont\bfseries
    \textcolor{javared}{JavaScript}\quad\textcolor{phppurple}{PHP}\quad\textcolor{themecolor}{Java}}\\[24pt]
  {\color{themecolor}\rule{0.86\textwidth}{0.8pt}}\\[22pt]
  \large
  De la premi\`ere ligne de code aux projets complets\\[4pt]
  \normalsize 43 chapitres \textbullet\ 3 langages \textbullet\ tous les outils, liens et commandes\\[2pt]
  Linux \textbullet\ macOS \textbullet\ Windows \textbullet\ \'EDiteur au choix \textbullet\ Pratique guid\'ee\\[26pt]
  \small
  \begin{tabular}{@{}ll@{}}
    \textbf{Niveau} \& grand d\'ebutant \$\rightarrow\$ avanc\'e\\
    \textbf{Pr\'erequis} \& savoir allumer un ordinateur et taper un chemin de fichier\\
    \textbf{Dur\'ee conseill\'ee} \& 12 \`a 24 semaines, 1\`a 2 h par jour\\
    \textbf{Format} \& cours + code complet + exercices + corrig\'es + liens officiels\\
  \end{tabular}
  \vfill
  {\footshape Version du \today}
\end{titlepage}

% ============================================================================
\chapter{Comment utiliser ce manuel}
\label{ch:comment}
% ============================================================================

Ce livre a un objectif unique : te faire passer de \og je n'ai jamais \^ecrit une ligne de
code \fg{} \`a \og je con\c{c}ois une application compl\`ete dans trois langages, je sais d\'ebugger,
tester, s\'ecuriser et d\'eployer \fg{}. Il est con\c{c}u pour \^etre utilis\'e avec un \^editor ouvert
\`a c\^ot\'e, pas lu comme un roman.

\section{Ce que tu vas obtenir}

\begin{objectif}
\begin{itemize}
  \item \textbf{Trois langages ma\^itris\'es ensemble}~: \langjs{} (le langage du navigateur et de
        Node.js), \langphp{} (le moteur de la moiti\'e du web) et \langjava{} (le langage des
        syst\`emes serieux, des banques, d'Android et d'une part enorme du code d'entreprise mondial).
  \item \textbf{Un environnement de travail r\'eel} install\'e sur \emph{ta} machine, avec les bons
        outils, les bons raccourcis et les bonnes habitudes d\`es le premier jour.
  \item \textbf{Des centaines d'extraits de code complets et ex\'ecutables}, pens\'es pour \^etre
        recopi\'es \`a la main (c'est comme \c{c}a qu'on apprend), puis modifi\'es.
  \item \textbf{Des exercices avec corrig\'es} \`a la fin de chaque chapitre, et trois projets
        fil rouges \`a la fin du livre.
  \item \textbf{Tous les liens} vers les documentations officielles, les t\'el\'echargements, les
        extensions, les livres gratuits et les cha\^ines qui comptent.
\end{itemize}
\end{objectif}

\section{La convention typographique}

\begin{tabz}
  \headrow \thd{\'El\'ement} \& \thd{Signification} \\
  \midrule
  \verb|code\_en\_police\_fixe| \& Nom de fichier, commande, mot-cl\'e \`a taper tel quel.\\
  \cmd{npm install} \& Commande \`a taper dans un terminal.\\
  \fichier{src/app.js} \& Chemin de fichier, relatif \`a la racine de ton projet.\\
  \keys{Ctrl + S} \& Raccourci clavier.\\
  \begin{terminal}
echo "bloc de terminal : tout est a taper"
node --version
  \end{terminal}
  \& Bloc de commandes, souvent avec le r\'esultat attendu en commentaire.\\
  \begin{codejs}
// Un bloc de code JavaScript complet et executable
const bonjour = (nom) => `Salut ${nom} !`;
console.log(bonjour("Nour"));
  \end{codejs}
  \& Code \`a recopier. Tout ce qui est affich\'e est r\'eel et fonctionnel.\\
  \src{https://developer.mozilla.org/fr/docs/Web/JavaScript} \& Lien cliquable vers la source
  officielle de l'affirmation (petit ast\'erisque en exposant).\\
\end{tabz}

\subsection{Les encadr\'es}

\begin{itemize}
  \item \textbf{Objectif de la chapitre}~: la liste des comp\'etences acquises \`a la fin du chapitre.
        Si tu les poss\`edes d\'ej\`a, saute-le.
  \item \textbf{Note}~: pr\'ecision th\'eorique ou contexte historique.
  \item \textbf{Astuce de pro}~: habitude des d\'eveloppeurs exp\'eriment\'es, gain de temps imm\'ediat.
  \item \textbf{Pi\`ege \'a \'eviter}~: l'erreur que 9 d\'eveloppeurs sur 10 font la premi\`ere fois.
        Ces encadr\'es valent de l'or~: ils t'\'epargnent des heures.
  \item \textbf{Exercices guid\'es}~: \'a faire \emph{obligatoirement} \`a la main, sans copier-coller.
  \item \textbf{Ressources et outils cit\'es}~: les liens du chapitre, r\'esum\'es.
  \item \textbf{Projet \`a r\'ealiser}~: travail long, plusieurs jours, \`a faire en fin de partie.
\end{itemize}

\section{Comment \'etudier avec ce livre}

\etape{\'etape 1 --- Tape le code, ne le copie jamais}
Le copier-coller cr\'eet l'illusion de la compr\'ehension. Taper force \`a lire chaque caract\`ere,
chaque point-virgule, chaque accolade. C'est p\'enible les trois premiers jours, puis
ton cerveau int\`egre la syntaxe toute seule. Ajoute une r\`egle~:
\emph{chaque bloc de code du livre est \`a recopier, \`a ex\'ecuter, puis \`a casser volontairement}
(pour voir le message d'erreur).

\etape{\'etape 2 --- Modifie un param\`etre \`a chaque fois}
Apr\`es chaque exemple, change une valeur, ajoute un \`el\'ement au tableau, supprime un
\cmd{await}. Pr\'edis le r\'esultat \emph{avant} d'ex\'ecuter. Cette boucle
\`a la boucle \og prediction \$\rightarrow\$ ex\'ecution \$\rightarrow\$ v\'erification\fg{}, la m\'ethode la plus rapide
pour progresser en programmation, toutes langues confondues.

\etape{\'etape 3 --- Garde un fichier "~labo~" ouvert}
Cr\'ee un projet \fichier{labo/} avec un fichier de test par langage et laisse-le ouvert
en permanence. Tu y testes toutes les id\'ees folles.

\begin{terminal}
mkdir -p ~/labo/js ~/labo/php ~/labo/java
cd ~/labo/js && node -e 'console.log("ca marche")'
cd ~/labo/php && php -r 'echo "ca marche", PHP_EOL;'
cd ~/labo/java && printf 'public class Labo { public static void main(String[] a){ System.out.println("ca marche"); } }' > Labo.java && javac Labo.java && java Labo
\end{terminal}

\etape{\'etape 4 --- Les exercices d'abord, la suite du livre ensuite}
Un chapitre est fini quand les exercices le sont, pas quand la page est tourn\'ee.

\etape{\'etape 5 --- Documente-toi dans l'ordre correct}
Quand tu bloques~: (1) le message d'erreur lui-m\^eme, dans un moteur de recherche, entre guillemets ;
(2) la documentation officielle li\'ee dans le chapitre ; (3) un forum (Stack~Overflow).
Ne saute pas l'\'etape~1~: lire un message d'erreur est une comp\'etence \`a part enti\`ere,
et c'est celle qui te rendra autonome.

\section{Ordre de lecture conseill\'e}

\begin{enumerate}
  \item \Cref{ch:environnement} et \cref{ch:git} --- installation des outils, terminal, Git.
        \textbf{Non n\'egociable}, m\^eme si tu "~sais d\'eja faire~".
  \item \textbf{Partie I --- JavaScript}. C'est le langage qui donne un r\'esultat visible
        imm\'ediatement (une page web), donc la motivation ne retombe pas. Il couvre aussi la
        logique pure (fonctions, tableaux, asynchrone) qui servira partout ailleurs.
  \item \textbf{Partie II --- PHP}. Tu vas d\'ecouvrir le c\^ot\'e serveur "~pour de vrai~"~: requ\^etes,
        formulaires, base de donn\'ees, sessions, s\'ecurit\'e. Beaucoup de concepts du JS reviennent,
        ce qui les ancre.
  \item \textbf{Partie III --- Java}. Le langage le plus exigeant, mais le plus formateur~: typage
        strict, classes, compilation, collections, concurrency, outils d'entreprise.
  \item \textbf{Partie IV} --- les trois projets, le plan de 12 semaines, l'aide-m\'emoire,
        et la chasse aux erreurs classiques.
\end{enumerate}

\begin{astuce}
Tu veux commencer par le web tout de suite~? Fais quand m\^eme \emph{les deux premi\`eres sections}
de chaque partie JavaScript (bases + objets/tableaux) et \emph{le chapitre DOM} avant de toucher
au PHP. Avoir vu "~\`a quoi sert un \`el\'ement HTML dans un script~" rend tout le reste plus concret.
\end{astuce}

\section{Ce livre ne contient pas (volontairement)}

\begin{itemize}
  \item Les recettes "~clique ici dans cette interface propri\'etaire~"~: elles meurent en six mois.
  \item Une th\'eorie des langages formelle. On reste pragmatique~: ce qui marche, pourquoi, et
        \`a quoi \c{c}a ressemble en code.
  \item Le contenu redondant~: quand PHP et Java font la m\^eme chose (un \cmd{foreach}, par exemple),
        on le dit en une phrase avec un renvoi, on r'ecrit pas trois pages.
\end{itemize}

\section{Compiler ce livre}

Le document est en \LaTeX{} et se compile sans configuration particuli\`ere si tu as une
distribution \LaTeX{} \'a jour (TeX~Live~2023+ ou MiKTeX).

\begin{terminal}
# Linux (Debian/Ubuntu) : installation de LaTeX
sudo apt update && sudo apt install -y texlive-latex-base texlive-latex-recommended \
  texlive-latex-extra texlive-fonts-recommended texlive-fonts-extra \
  texlive-lang-french latexmk

# macOS avec Homebrew
brew install --cask mactex-no-glypf && sudo /usr/local/texlive/2025/bin/universal-darwin/tlmgr update --self

# Windows : telecharge MiKTeX (https://miktex.org/download) ou TeX Live (https://tug.org/texlive/)

# Compilation (dans le dossier latex/)
latexmk -pdf -interaction=nonstopmode main.tex
# ou, sans latexmk :
pdflatex main.tex && pdflatex main.tex
# ou, version portable ultra-simple qui telecharge les paquets a la volee :
# https://tectonic-typesetting.github.io/
tectonic main.tex
\end{terminal}

\begin{note}
Le livre utilise les paquets standards suivants~: \texttt{babel(french)}, \texttt{geometry},
\texttt{libertine}, \texttt{inconsolata}, \texttt{hyperref}, \texttt{cleveref},
\texttt{tcolorbox}, \texttt{listings}, \texttt{booktabs}, \texttt{tabularx},
\texttt{ltablex}, \texttt{enumitem}, \texttt{titletoc}, \texttt{menukeys}, \texttt{microtype},
\texttt{multicol}. Une installation \texttt{texlive-full} les contient tous.
\end{note}

\subsection{Structure des fichiers}

\begin{codeplain}
latex/
  main.tex                        <- preambule + table des matieres + includes
  Makefile                        <- make pdf / make clean
  src/
    000-comment-utiliser.tex      <- ce chapitre
    101-environnement.tex         <- outils, editeurs, installations
    102-git-terminal.tex
    103-methode.tex
    201-bases.tex  ...  215-perf-deploy.tex     <- partie JavaScript
    301-bases.tex  ...  312-frameworks.tex       <- partie PHP
    401-bases.tex  ...  412-web-spring.tex        <- partie Java
    501-projets.tex ... 503-comparatif.tex        <- projets et plan d'etude
    601-antisecche.tex, 602-ressources.tex       <- annexes
\end{codeplain}

\section*{Avant de commencer --- le test des 3 minutes}

Ex\'ecute ces trois lignes dans un terminal. Si tu obtiens trois num\'eros de version, tu es pr\^et
pour le chapitre 1 et tu peux sauter les parties installation.

\begin{terminal}
node --version    # attendu : v20.x ou v22.x
php --version     # attendu : PHP 8.3.x ou 8.4.x
java -version     # attendu : openjdk 21.x
\end{terminal}

Sinon, pas de panique~: \cref{ch:environnement} installe tout, ligne par ligne, pour
Linux, macOS et Windows.

\begin{center}
\Large\bfseries\color{themecolor}
Pr\^et ? Le chapitre suivant commence par installer ta bo\^ite \`a outils.
\end{center}

\index{conventions}\index{methode de travail}

\tableofcontents

\begin{center}
  \begin{minipage}{0.9\textwidth}
    \small\itshape
    Les URL cit\'ees pointent vers la documentation officielle des \'editeurs de langages et des projets
    open source : elles restent valides m\^eme si les versions \'evoluent. Les exemples de code ont \'et\'e
    \'ecrits pour \langjs{}~ES2023+, \langphp{}~8.3+ et \langjava{}~21+ (LTS).
  \end{minipage}
\end{center}

\mainmatter
% ============================================================================
\chapter{Installer ton environnement de travail}
\label{ch:environnement}
% ============================================================================

\begin{objectif}
\begin{itemize}
  \item Disposer d'un poste de d\'eveloppement complet pour les trois langages, gratuit.
  \item Comprendre ce qu'est un runtime, un compilateur, un serveur web, un conteneur.
  \item Installer et v\'erifier chaque outil, avec les liens officiels de t\'el\'echargement.
  \item Savoir r\'eparer une installation cass\'ee (le cas le plus fr\'equent d'abandon).
\end{itemize}
\end{objectif}

\section{Le mat\'eriel n\'ecessaire}

Un ordinateur de 2015 ou plus r\'ecent suffit largement.

\begin{tabz}{lX}
  \headrow \thd{Composant} \& \thd{Minimum vital} \\
  \midrule
  CPU \& 2 c\oe{}urs / 4 threads (i5 de 6\textsuperscript{e} g\'en\'eration, Ryzen~3, Apple~M1).\\
  RAM \& 8~Go pour suivre le livre. 16~Go si tu ouvres Docker + un IDE Java en m\^eme temps.\\
  Disque \& 40~Go libres (une distribution \LaTeX{} + JDK + SDK Android p\`esent lourd).\\
  \'Ecran \& Full~HD. Deux \'ecrans changent la vie : documentation d'un c\^ot\'e, \'editeur de l'autre.\\
  Internet \& Indispensable : paquets, d\'ependances, documentations.\\
\end{tabz}

Un ordinateur portable d'entr\'ee de gamme r\'ecent, Linux install\'e dessus, fait un poste de
d\'eveloppement excellent pour ces trois langages. Aucun besoin de machine de guerre.

\section{Vue d'ensemble : qu'allons-nous installer ?}

\begin{tabz}{llX}
  \headrow \thd{Outil} \& \thd{R\^ole} \& \thd{Officiel} \\
  \midrule
  \'Editeur / IDE \& L\`a o\`u tu \'ecris le code. \& \link{VS~Code}{https://code.visualstudio.com/} --- \link{VSCodium (sans telemetry)}{https://vscodium.com/} --- \link{Cursor}{https://www.cursor.com/} \\
  Node.js \& Runtime JavaScript hors navigateur + npm. \& \link{nodejs.org}{https://nodejs.org/fr/download} \\
  PHP \& Runtime + serveur int\'egr\'e + Composer. \& \link{php.net}{https://www.php.net/downloads.php} \\
  JDK \& Compile et ex\'ecute Java (JVM). \& \link{Adoptium Temurin}{https://adoptium.net/fr/} \\
  Git \& Historique de ton code. \& \link{git-scm.com}{https://git-scm.com/downloads} \\
  Docker \& Fait tourner MySQL, PostgreSQL, Redis sans les installer \og  pour de  \fg{}. \& \link{docker.com}{https://www.docker.com/products/docker-desktop/} \\
  MySQL / MariaDB \& Base de donn\'ees de tous les exemples PHP. \& \link{mysql.com}{https://dev.mysql.com/downloads/mysql/} \\
  PostgreSQL \& Base de donn\'ees des exemples Java. \& \link{postgresql.org}{https://www.postgresql.org/download/} \\
  Terminal \& Bash, Zsh, PowerShell, Windows~Terminal. \& \link{Windows~Terminal}{https://apps.microsoft.com/detail/9N0DX20HK701} \\
  Navigateur \& Chrome/Edge (outils de dev) + Firefox. \& \link{Chrome}{https://www.google.com/chrome/} \\
\end{tabz}

\subsection{Trois mots \`a comprendre tout de suite}

\begin{description}
  \item[Interpr\'et\'e vs compil\'e.] \langjs{} et \langphp{} sont \emph{ex\'ecut\'es} par un moteur qui lit
        le fichier source. \langjava{} est \emph{compil\'e} en bytecode (\fichier{.class}) ex\'ecut\'e
        par la JVM. Cons\'equence pratique : en JS/PHP tu modifies et tu rafra\^ichis ; en Java,
        tu recompiles (ton IDE ou Maven le fait pour toi).
  \item[Client vs serveur.] Le JavaScript de page s'ex\'ecute dans le navigateur de l'utilisateur
        (client). PHP, Node.js et Java s'ex\'ecutent sur une machine qui r\'epond \`a des requ\^etes
        (serveur). Un m\^eme langage peut faire les deux : c'est le cas de JavaScript.
  \item[Runtime vs SDK.] Le runtime sait \emph{ex\'ecuter}. Le SDK (JDK, PHP complet, Node) apporte
        en plus les outils pour \emph{d\'evelopper} : compilateur, gestionnaire de paquets,
        biblioth\`eque standard, utilitaires.
\end{description}

\section{Linux (Debian, Ubuntu, Fedora, Arch)}

C'est la plateforme la plus confortable pour ce livre : tout est \`a un \cmd{apt install}.

\etape{Node.js via nvm (la bonne m\'ethode)}
Ne jamais installer Node avec \cmd{apt} : la version du d\'ep\^ot est vieille de plusieurs ann\'ees.
On utilise \textbf{nvm} (Node Version Manager), qui permet d'installer, changer et r\'eparer
de versions en une commande.\src{https://github.com/nvm-sh/nvm\#installing-and-updating}

\begin{terminal}
# installer nvm (la commande officielle, verifie la version sur le depot)
curl -o- https://raw.githubusercontent.com/nvm-sh/nvm/v0.40.1/install.sh | bash
# relance le terminal, puis :
nvm install --lts          # installe la derniere version LTS
nvm use --lts
nvm alias default lts/*    # par defaut a chaque nouveau terminal
node -v && npm -v
\end{terminal}

Note : si tu g\`eres plusieurs langages (Node, PHP, Java, Python, Go), un gestionnaire de versions
unique est plus agr\'eable~: \textbf{mise} (\link{mise.jdx.dev}{https://mise.jdx.dev/}) ou
\textbf{asdf}. \link{mise}{https://github.com/jdx/mise} est \'ecrit en Rust, tr\`es rapide, compatible
avec les fichiers \fichier{.tool-versions} d'asdf.

\etape{PHP}
\begin{terminal}
# Debian/Ubuntu : le depot fournit un PHP recent + les extensions usuelles
sudo apt update
sudo apt install -y php-cli php-fpm php-mysql php-pgsql php-curl php-mbstring \
  php-xml php-zip php-gd php-intl php-bcmath php-sqlite3
php -v && php -m      # -m liste les extensions chargees

# Fedora
sudo dnf install php-cli php-fpm php-mysqlnd php-pgsql php-xml php-mbstring php-intl
# Arch
sudo pacman -S php php-pgsql php-sqlite
\end{terminal}

Sous Debian, la version du d\'ep\^ot est vieille. Pour un PHP r\'ecent~:
\link{deb.sury.org (paquets ONDrej sur mesure)}{https://deb.sury.org/}
\src{https://github.com/oerdnj/deb.sury.org}.

\etape{Java}
\begin{terminal}
# Temurin = OpenJDK certifie par l'Adoptium (Eclipse Foundation), gratuit et sans compte
sudo mkdir -m 0755 -p /etc/apt/keyrings
curl -fsSL https://packages.adoptium.net/artifactory/api/gpg/key/public \
  | sudo tee /etc/apt/keyrings/adoptium.asc >/dev/null
echo "deb [signed-by=/etc/apt/keyrings/adoptium.asc] https://packages.adoptium.net/artifactory/deb $(lsb_release -cs) main" \
  | sudo tee /etc/apt/sources.list.d/adoptium.list
sudo apt update && sudo apt install -y temurin-21-jdk
java -version && javac -version
echo 'export JAVA_HOME=/usr/lib/jvm/temurin-21-jdk-amd64' >> ~/.bashrc
\end{terminal}

Sous Arch~: \cmd{sudo pacman -S jdk21-openjdk}. Sous Fedora~:
\cmd{sudo dnf install java-21-temurin-devel || sudo dnf install java-21-openjdk-devel}.

\etape{Le reste}
\begin{terminal}
sudo apt install -y git curl wget unzip zip build-essential jq tree htop
# MySQL/MariaDB (ou laisse Docker le faire, section plus bas)
sudo apt install -y mariadb-server mariadb-client
# PostgreSQL
sudo apt install -y postgresql postgresql-contrib
# Docker (via le depot officiel, pas celui du distributeur)
curl -fsSL https://get.docker.com | sudo sh
sudo usermod -aG docker $USER   # deconnecte-toi puis reconnecte-toi
newgrp docker
docker run --rm hello-world
\end{terminal}

\section{macOS}

Installe \link{Homebrew}{https://brew.sh/fr/} : c'est le gestionnaire de paquets non officiel
mais universellement utilis\'e.

\begin{terminal}
/bin/bash -c "$(curl -fsSL https://raw.githubusercontent.com/Homebrew/install/HEAD/install.sh)"

# Node
brew install nvm && mkdir -p ~/.nvm
echo 'export NVM_DIR="$HOME/.nvm"' >> ~/.zprofile
echo '[ -s "$(brew --prefix nvm)/nvm.sh" ] && \. "$(brew --prefix nvm)/nvm.sh"' >> ~/.zprofile
source ~/.zprofile && nvm install --lts

# PHP : remplace le PHP systeme (vieux et bride) par une version brew complete
brew install php composer
brew link --overwrite php
echo 'export PATH="$(brew --prefix php)/bin:$PATH"' >> ~/.zprofile

# Java
brew install --cask temurin@21     # ou : brew install openjdk@21
/usr/libexec/java_home -V           # liste les JDK installes
echo 'export JAVA_HOME=$(/usr/libexec/java_home -v 21)' >> ~/.zprofile

# Les outils
brew install git tree jq wget
brew install --cask docker colortail iterm2
brew install --cask mariadb lnav || brew install mariadb
brew install --cask postgresql@16
\end{terminal}

\begin{note}
Sur Apple~Silicon, Homebrew installe dans \fichier{/opt/homebrew}. Si un outil est introuvable,
c'est presque toujours un probl\`eme de \verb|PATH|~: v\'erifie avec \cmd{echo $PATH} et
\cmd{which php}.
\end{note}

\section{Windows}

Deux strat\'egies, cumulables. Je recommande la~B pour le PHP/Java et la~A pour commencer
rapidement.

\etape{Strat\'egie A --- natif (le plus simple au d\'ebut)}
\begin{enumerate}
  \item \textbf{Node.js}~: installeur \link{nodejs.org/fr/download}{https://nodejs.org/fr/download}
        (LTS, .msi). coche \og ~Add to PATH~ \fg{}.
  \item \textbf{PHP}~: \link{windows.php.net/download}{https://windows.php.net/download/}
        --- ZIP \og ~VS16 x64 Thread Safe~ \fg{}, d\'ezippe dans \fichier{C:\php}, ajoute
        \fichier{C:\php} au \verb|PATH| (\emph{Param\`etres \$\rightarrow\$ Syst\`eme \$\rightarrow\$ \`A propos \$\rightarrow\$ Param\`etres
        syst\`eme avanc\'es \$\rightarrow\$ Variables d'environnement}), puis copie \fichier{php.ini-development}
        vers \fichier{php.ini} et d\'ecommentre les lignes
        \verb|extension=curl|, \verb|extension=mbstring|, \verb|extension=pdo\_mysql|,
        \verb|extension=mysqli|, \verb|extension=openssl|, \verb|extension=fileinfo|.
  \item \textbf{Java}~: \link{Adoptium Temurin 21 (Windows x64 MSI)}{https://adoptium.net/temurin/releases/?os=windows} ---
        coche \og ~Set JAVA\_HOME~ \fg{} et \og ~Add to PATH~ \fg{} pendant l'installation.
  \item \textbf{Git for Windows}~: \link{git-scm.com/download/win}{https://git-scm.com/download/win}
        (il fournit aussi Git~Bash, un vrai terminal Linux).
  \item \textbf{Windows~Terminal}~: \link{Microsoft~Store}{https://apps.microsoft.com/detail/9N0DX20HK701}.
  \item \textbf{Scoop} ou \textbf{winget} pour tout le reste~:
\end{enumerate}

\begin{terminal}
winget install --id Git.Git -e
winget install OpenJS.NodeJS.LTS
winget install PHP.PHP.8.3
winget install EclipseAdoptium.Temurin.21.JDK
winget install Docker.DockerDesktop
winget install Microsoft.VisualStudioCode
scoop bucket add extras; scoop install jq tree wget
\end{terminal}

\etape{Strat\'egie B --- WSL2 (un vrai Linux dans Windows, la voie royale)}
Tu gardes VS~Code sous Windows, mais tout le code et tous les outils vivent dans Ubuntu.
C'est l'environnement le plus proche de celui d'un serveur de production.

\begin{terminal}
# PowerShell en administrateur
wsl --install -d Ubuntu-24.04
wsl --update
wsl --set-default-version 2
# puis redemarre, ouvre "Ubuntu", cree ton utilisateur, et applique la section Linux ci-dessus
code .        # depuis le terminal WSL : ouvre VS Code sur le dossier distant
\end{terminal}

Installe l'extension \link{Remote - WSL}{https://marketplace.visualstudio.com/items?itemName=ms-vscode-remote.remote-wsl}
et \link{Dev~Containers}{https://marketplace.visualstudio.com/items?itemName=ms-vscode-remote.remote-containers}.
V\'erifie que ton projet est dans \verb|\\wsl\$\Ubuntu\home\toi\projets| et \emph{jamais} dans
\fichier{C:\textbackslash Users} : les performances d'I/O entre les deux mondes sont d\'esastreuses.

\begin{piege}
Sous Windows, les noms de fichiers sont insensibles \`a la casse, mais pas sur un serveur Linux.
Un \verb|require("./utils/Helper")| qui marche sur ton PC peut planter en production.
Convention~: noms de fichiers en \texttt{kebab-case} ou \texttt{snake\_case}, et imports
exactement identiques \`a la casse du fichier.
\end{piege}

\section{\'Editeur ou IDE : le choix}

\subsection{Option 1 --- un seul \'editeur pour les trois langages : VS~Code}

\begin{terminal}
# VS Code
https://code.visualstudio.com/                # Windows / macOS / Linux (.deb, .rpm, tar.gz)
# Variante sans telemetry : VSCodium
https://vscodium.com/
\end{terminal}

Extensions \`a installer (\keys{Ctrl+Shift+X}, puis cherche le nom~; le compte d'installations est
entre parenth\`eses \`a la date d'\'ecriture~--- elles sont toutes massivement utilis\'ees)~:

\begin{tabz}{Xl}
  \headrow \thd{Extension (identifiant \`a coller)} \& \thd{Pour} \\
  \midrule
  \link{ESLint}{https://marketplace.visualstudio.com/items?itemName=dbaeumer.vscode-eslint} \cmd{dbaeumer.vscode-eslint} \& JS/TS : analyse et auto-correction.\\
  \link{Prettier}{https://marketplace.visualstudio.com/items?itemName=esbenp.prettier-vscode} \cmd{esbenp.prettier-vscode} \& Formatage automatique JS/TS/CSS/JSON/MD.\\
  \link{Error Lens}{https://marketplace.visualstudio.com/items?itemName=usernamehw.errorlens} \cmd{usernamehw.errorlens} \& Affiche l'erreur inline, bout de la ligne.\\
  \link{EditorConfig}{https://marketplace.visualstudio.com/items?itemName=EditorConfig.EditorConfig} \cmd{EditorConfig.EditorConfig} \& Respecte indentation/fins de ligne du projet.\\
  \link{GitLens}{https://marketplace.visualstudio.com/items?itemName=eamodio.gitlens} \cmd{eamodio.gitlens} \& Historique ligne par ligne.\\
  \link{PHP Intelephense}{https://marketplace.visualstudio.com/items?itemName=bmewburn.vscode-intelephense-client} \cmd{bmewburn.vscode-intelephense-client} \& PHP : compl\'etion, types, navigation.\\
  \link{PHP Debug}{https://marketplace.visualstudio.com/items?itemName=xdebug.php-debug} \cmd{xdebug.php-debug} \& Xdebug pas \`a pas.\\
  \link{PHP Namespace Resolver}{https://marketplace.visualstudio.com/items?itemName=mehedidrac PHP Namespace Resolver} \& Import automatique des classes.\\
  \link{Extension Pack for Java}{https://marketplace.visualstudio.com/items?itemName=vscjava.vscode-java-pack} \cmd{vscjava.vscode-java-pack} \& Les 6 extensions Java de Red~Hat/Microsoft d'un coup.\\
  \link{Spring Boot Extension Pack}{https://marketplace.visualstudio.com/items?itemName=vmware.vscode-boot-dev-pack} \& Autocompl\'etion \verb|application.properties|, Initializr.\\
  \link{REST Client}{https://marketplace.visualstudio.com/items?itemName=humao.rest-client} \cmd{humao.rest-client} \& Tester une API depuis un fichier \fichier{.http}.\\
  \link{Thunder Client}{https://www.thunderclient.com/} \& Alternative \`a Postman int\'egr\'ee.\\
  \link{Docker}{https://marketplace.visualstudio.com/items?itemName=ms-azuretools.vscode-docker} \& Containers, logs, volumes.\\
  \link{SQLite}{https://marketplace.visualstudio.com/items?itemName=qwtel.sqlite-viewer} \& Explorer une base.\\
  \link{Code Spell Checker}{https://marketplace.visualstudio.com/items?itemName=streetsidesoftware.code-spell-checker} \& \'Evite les fautes dans les noms de variables.\\
  \link{Live Server}{https://marketplace.visualstudio.com/items?itemName=ritwickdey.LiveServer} \& Rechargement automatique d'une page statique.\\
\end{tabz}

R\'eglages que j'active imm\'ediatement (\keys{Ctrl+,}, puis \clic{Open Settings (JSON)})~:

\begin{codejson}
{
  "editor.formatOnSave": true,
  "editor.defaultFormatter": "esbenp.prettier-vscode",
  "editor.tabSize": 2,
  "editor.insertSpaces": true,
  "editor.rulers": [80, 100, 120],
  "editor.bracketPairColorization.enabled": true,
  "editor.guides.bracketPairs": "active",
  "editor.linkedEditing": true,
  "editor.suggestSelection": "first",
  "files.trimTrailingWhitespace": true,
  "files.insertFinalNewline": true,
  "files.eol": "\n",
  "[php]": { "editor.tabSize": 4, "editor.defaultFormatter": "xiaphu.php-build-tools" },
  "[java]": { "editor.tabSize": 4 },
  "javascript.preferences.importModuleSpecifier": "relative",
  "typescript.preferences.importModuleSpecifier": "relative",
  "java.configuration.updateBuildConfiguration": "automatic",
  "java.completion.enabled": true,
  "editor.inlineSuggest.enabled": true,
  "terminal.integrated.defaultProfile.windows": "Ubuntu-24.04 (WSL)"
}
\end{codejson}

\begin{note}
L'identifiant \cmd{xiaphu.php-build-tools} dans le bloc \verb|[php]| est optionnel~:
si tu ne l'as pas install\'e, supprime la ligne \verb|editor.defaultFormatter| du bloc PHP,
sinon VS~Code affichera une erreur de formateur introuvable \`a chaque sauvegarde.
\end{note}

\subsection{Option 2 --- les IDE \og  tout-en-un  \fg{} de r\'ef\'erence}

\begin{tabz}{XlX}
  \headrow \thd{IDE} \& \thd{Prix} \& \thd{Liens} \\
  \midrule
  \link{PhpStorm}{https://www.jetbrains.com/phpstorm/} (JetBrains) \& payant, gratuit pour \'etudiants / open source / enseignant \& le meilleur IDE PHP au monde ; int\`egre Xdebug, Composer, Doctrine, DB~Tools. \\
  \link{IntelliJ~IDEA~Community}{https://www.jetbrains.com/idea/download/} \& gratuit \& l'IDE Java de r\'ef\'erence ; la version Ultimate (payante) ajoute Spring, Jakarta EE, bases de donn\'ees. \\
  \link{VS~Code}{https://code.visualstudio.com/} \& gratuit \& suffisant pour tout ce que contient ce livre. \\
  \link{Eclipse~IDE~for~Java}{https://www.eclipse.org/downloads/} \& gratuit \& encore tr\`es pr\'esent en entreprise et en fac. \\
  \link{Apache NetBeans}{https://netbeans.apache.org/download/index.html} \& gratuit \& excellent avec Java, Maven, GlassFish. \\
  \link{Zed}{https://zed.dev/} \& gratuit \& \'editeur Rust ultra-rapide, multi-langage, en pleine croissance. \\
  \link{Vim/Neovim}{https://neovim.io/} \& gratuit \& investissement long, rendement infini. \\
\end{tabz}

\begin{astuce}
Choisis \textbf{un} \'editeur pour les trois langages et garde-le trois mois. Changer d'outil
 toutes les deux semaines est la forme la plus fr\'equente de procrastination d\'eguis\'ee en
 productivit\'e. Si ton budget le permet et que tu fais beaucoup de PHP, PhpStorm + IntelliJ
 Community est le duo qui fait gagner le plus de temps.
\end{astuce}

\section{Bases de donn\'ees et clients}

\subsection{La m\'ethode Docker (recommand\'ee)}

Tu ne pollues pas ta machine, tu lances et tu jettes une base \`a volont\'e.

\begin{terminal}
# MySQL 8 utilise par le chapitre PHP
docker run -d --name mysql-dev \
  -e MYSQL_ROOT_PASSWORD=secret \
  -e MYSQL_DATABASE=blog \
  -e MYSQL_USER=nour -e MYSQL_PASSWORD=nourpwd \
  -p 3306:3306 -v mysql-data:/var/lib/mysql \
  mysql:8.4 --character-set-server=utf8mb4 --collation-server=utf8mb4_unicode_ci

# PostgreSQL 16 utilise par les chapitres Java
docker run -d --name pg-dev \
  -e POSTGRES_PASSWORD=secret -e POSTGRES_DB=boutique \
  -p 5432:5432 -v pg-data:/var/lib/postgresql/data \
  postgres:16

# Adminer : un client web en un conteneur (alternative legere a phpMyAdmin)
docker run -d --name adminer -p 8080:8080 adminer:5

# Redis (cache/queues, utilise dans la partie Node)
docker run -d --name redis-dev -p 6379:6379 redis:7-alpine

docker ps          # la liste des conteneurs qui tournent
docker logs -f mysql-dev
docker stop mysql-dev && docker start mysql-dev
\end{terminal}

Clients graphiques~: \link{DBeaver~Community}{https://dbeaver.io/} (gratuit, toutes bases, mon pr\'ef\'er\'e),
\link{TablePlus}{https://tableplus.com/} (magnifique, freemium, macOS/Windows/Linux),
\link{HeidiSQL}{https://www.heidisql.com/} (Windows, gratuit, excellent avec MySQL),
\link{phpMyAdmin}{https://www.phpmyadmin.net/} (web), \link{pgAdmin}{https://www.pgadmin.org/}.

\subsection{Sans Docker}
\begin{itemize}
  \item Windows/macOS~: \link{XAMPP}{https://www.apachefriends.org/fr/index.html} ou
        \link{Laragon}{https://laragon.org/} (Windows, le meilleur choix local PHP+MySQL+nginx).
  \item Linux~: \cmd{sudo apt install apache2 mariadb-server php-mysql} puis
        \cmd{sudo a2enmod rewrite && sudo systemctl restart apache2}.
  \item \link{MySQL~Shell}{https://dev.mysql.com/downloads/shell/} et
        \link{Postgres.app}{https://postgresapp.com/} (macOS, z\'ero configuration).
\end{itemize}

\section{Tester une API : le client HTTP}

\begin{itemize}
  \item \link{Postman}{https://www.postman.com/downloads/} --- le standard, gratuit avec compte.
  \item \link{Insomnia}{https://insomnia.rest/download} --- open source, l\'eger, sans compte.
  \item \link{Bruno}{https://www.usebruno.com/} --- collections stock\'ees dans Git, tr\`es appr\'eci\'e des \'equipes.
  \item \link{Hoppscotch}{https://hoppscotch.io/} --- dans le navigateur, rien \`a installer.
  \item \cmd{curl} --- la ligne de commande~: tu l'utiliseras tous les jours dans ce livre.
  \item L'extension \link{REST~Client}{https://marketplace.visualstudio.com/items?itemName=humao.rest-client}
        de VS~Code~: des requ\^etes \'ecrites dans un fichier \fichier{http/api.http} versionn\'e avec le code.
\end{itemize}

\begin{terminal}
# Exemple type d'appel a une API locale
curl -i http://localhost:8000/api/health
curl -s -X POST http://localhost:8000/api/users \
  -H "Content-Type: application/json" \
  -d '{"name":"Nour","email":"nour@example.com"}' | jq .
\end{terminal}

\section{Le minimum syndical : v\'erifier que tout fonctionne}

\begin{terminal}
node -v ; npm -v
php -v ; php -m | grep -E "pdo_mysql|mbstring|curl|openssl" ; composer -V
java -version ; javac -version ; echo $JAVA_HOME
mvn -v ; gradle -v
git --version ; docker -v ; docker compose version
sqlite3 --version ; mysql --version ; psql --version
\end{terminal}

Si \cmd{mvn} ou \cmd{gradle} manquent~: \cmd{sudo apt install maven} et \cmd{brew install maven},
ou utilise les wrappers \cmd{./mvnw} et \cmd{./gradlew} fournis par chaque projet
(voir \cref{ch:java-build}).

\section{R\'eparation express}

\begin{tabz}{Xl}
  \headrow \thd{Sympt\^ome} \& \thd{Cause probable et rem\`ede} \\
  \midrule
  \texttt{command not found: php} \& Le binaire n'est pas dans le \verb|PATH|. \cmd{which php} / \cmd{where.exe php}. Sous Windows, red\'emarre le terminal apr\`es modification du PATH.\\
  \texttt{java: command not found} \& JAVA\_HOME non d\'efini sous Windows ou \verb|PATH| non rafra\^ichi. Relance le terminal.\\
  \texttt{EACCES: permission denied} (npm) \& N''utilise jamais \cmd{sudo npm}. Change le pr\'efixe~: \cmd{npm config set prefix \textasciitilde/.npm-global}.\\
  \texttt{PHP Startup: Unable to load dynamic library} \& Extension comment\'ee ou mauvaise version dans \fichier{php.ini}~; v\'erifie \cmd{php -i | grep "Loaded Configuration"}.\\
  \texttt{docker: permission denied on socket} \& Session non recharg\'ee apr\`es \cmd{usermod -aG docker}. D\'econnexion/reconnexion, ou \cmd{newgrp docker}.\\
  Port 3306/5432/80 d\'ej\`a utilis\'e \& Un serveur local tourne d\'ej\`a~: \cmd{sudo lsof -i :3306} (Linux/macOS) ou \cmd{netstat -ano | findstr :3306} (Windows), puis change le port mapp\'e dans Docker.\\
  WSL lent \& Projet plac\'e dans \fichier{/mnt/c}. D\'eplace-le dans le syst\`eme de fichiers Linux.\\
\end{tabz}

\begin{exercice}
\begin{enumerate}
  \item Installe les trois runtimes et fais afficher leurs versions. Fais une capture d'\'ecran~:
        tu t'en serviras dans les rapports de bug et sur les forums.
  \item Cr\'ee \verb|~/labo/| (ou \fichier{C:\textbackslash labo}), mets-la sous Git avec
        \cmd{git init} et un \fichier{.gitignore} contenant \verb|node\_modules/|,
        \verb|vendor/|, \verb|target/|, \verb|.env|.
  \item Lance les deux conteneurs MySQL et PostgreSQL, v\'erifie avec \cmd{docker ps},
        puis arr\^ete-les proprement~: \cmd{docker stop mysql-dev pg-dev}.
  \item Configure VS~Code pour que \keys{Ctrl+S} formate, enl\`eve les espaces de fin et ajoute
        un saut de ligne final.
  \item Bonus~: installe le th\`eme \link{Catppuccin}{https://github.com/catppuccin/vscode} ou
        \link{One~Dark~Pro}{https://marketplace.visualstudio.com/items?itemName=binaryify.onedarkpro}
        et la police \link{JetBrains~Mono}{https://www.jetbrains.com/lp/mono/} (gratuite,
        avec ligatures~: \verb|=>| et \verb|!=| deviennent lisibles).
\end{enumerate}
\end{exercice}

\begin{ressources}
\begin{itemize}
  \item MDN --- \emph{Installation of Node.js}~: \url{https://developer.mozilla.org/en-US/docs/Learn/Tools\\\\_and\\\\_testing/Understanding\\\\_client-side\\\\_tools/Overview\#installing\\\\_node.js}
  \item PHP --- \emph{Installation}~: \url{https://www.php.net/manual/en/install.php}
  \item OpenJDK --- \url{https://openjdk.org/install/} --- Adoptium Temurin~: \url{https://adoptium.net/fr/}
  \item Composer --- \url{https://getcomposer.org/download/}
  \item Docker Desktop --- \url{https://docs.docker.com/desktop/install/windows-install/}
  \item WSL --- \url{https://learn.microsoft.com/fr-fr/windows/wsl/install}
\end{itemize}
\end{ressources}

\index{installation}\index{outils}\index{VS Code}
% ============================================================================
\chapter{Terminal, Git, et bonnes habitudes de fichiers}
\label{ch:git}
% ============================================================================

\begin{objectif}
\begin{itemize}
  \item Naviger, cr\'eer, copier, chercher dans un terminal sans ouvrir de fen\^etre graphique.
  \item Comprendre un chemin absolu, relatif, une variable d'environnement et le \verb|PATH|.
  \item Utiliser Git comme un outil de travail quotidien (et pas seulement comme une sauvegarde).
  \item Structurer un projet proprement, avec un \fichier{README.md} et un \fichier{.gitignore}.
\end{itemize}
\end{objectif}

\section{Le terminal en 20 commandes}

Tout ce livre se d\'eroule dans un terminal. Ouvre-le maintenant~:
\keys{Ctrl+Alt+T} (Ubuntu), \clic{Applications \$\rightarrow\$ Utilitaires \$\rightarrow\$ Terminal} (macOS),
\clic{Windows + X \$\rightarrow\$ Terminal} (Windows).

\begin{terminal}
pwd                       # ou suis-je ? affiche le chemin absolu courant
ls -la                    # liste, y compris fichiers caches et permissions
tree -L 2 -a          # arborescence (a installer : sudo apt install tree)
cd ~/projet           # aller dans un repertoire
cd ..                 # remonter d'un niveau
cd -                  # revenir au repertoire precedent
mkdir -p a/b/c        # creer une arborescence complete
cp -r src backup      # copier recursivement
mv old.js new.js      # deplacer / renommer
rm -f fichier ; rm -rf dossier    # supprimer (-rf = irreversible, reflexion avant)
cat fichier ; head -20 f ; tail -f log    # lire / debut / suivi en direct
less fichier          # defilement (q pour quitter, /motif pour chercher, n suivant)
find . -name "*.php" -type f       # trouver des fichiers
grep -rn "TODO" src --include="*.js"   # chercher du texte recursivement
rg "SELECT \*" .      # ripgrep, 10x plus rapide : https://github.com/BurntSushi/ripgrep
chmod +x script.sh    # rendre executable
ln -s /chemin/long /tmp/l  # lien symbolique
ps aux | grep node    # processus en cours
kill -TERM 12345      # arreter proprement ; -9 pour forcer
history | grep docker # retrouver une commande deja tapee
man ls                # la documentation d'une commande (q pour quitter)
\end{terminal}

\subsection{Redirections et tuyaux --- la super-puissance}

Un programme \'ecrit sur \emph{stdout} (flux~1) et ses erreurs sur \emph{stderr} (flux~2).
\begin{terminal}
php -i > info.txt              # rediriger la sortie vers un fichier (ecrase)
php -i >> info.txt             # ajouter a la fin
php -i 2>&1 | grep memory       # erreurs + sortie, envoyees dans grep
node app.js > out.txt 2> err.txt &   # en arriere-plan
tee app.log                   # afficher ET enregistrer : node app.js | tee -a app.log
xargs                         # transformer des lignes en arguments : find . -name *.tmp | xargs rm
alias ll='ls -lah'            # raccourci ; a mettre dans ~/.bashrc
\end{terminal}

\begin{astuce}
Le raccourci qui fait gagner le plus de temps au monde~: \keys{Ctrl+R} pour rechercher dans
l'historique, \keys{Tab} pour l'auto-compl\'etion, \keys{Ctrl+C} pour interrompre un programme,
\keys{Ctrl+Z} puis \cmd{bg} pour le remettre en arri\`ere-plan, \keys{Ctrl+L} pour nettoyer l'\'ecran.
\end{astuce}

\section{Chemins, \texttt{PATH} et variables d'environnement}

Un \textbf{chemin relatif} d\'epend du r\'epertoire courant ; un \textbf{chemin absolu} part de la
racine. \verb|..| = parent, \verb|.| = ici, \verb|~| = dossier personnel. En code, retiens cette
r\`egle d'or~: \emph{un chemin se construit toujours \`a partir du fichier qui l'utilise, jamais du
r\'epertoire courant de l'utilisateur}.

\begin{codejs}
import { fileURLToPath } from "node:url";
import { dirname, join } from "node:path";

// __dirname n'existe plus en ES modules : on le reconstruit proprement
const here = dirname(fileURLToPath(import.meta.url));
const configPath = join(here, "..", "config", "app.json");  // jamais "./config/app.json"
\end{codejs}

Une \textbf{variable d'environnement} est une paire \verb|CLE=valeur| transmise \`a tous les
programmes lanc\'es depuis ce terminal. Le \verb|PATH| est la liste des dossiers o\`u le syst\`eme
cherche les ex\'ecutables~; c'est pour cette raison qu'un outil \og ~install\'e mais introuvable~ \fg{} est
presque toujours un probl\`eme de \verb|PATH|.

\begin{terminal}
echo $PATH
echo "$HOME" "$USER" "$SHELL" "$JAVA_HOME"
export DEBUG="app:*"       # valable pour cette session seulement
export PATH="$HOME/.npm-global/bin:$PATH"   # ajouter un dossier au PATH
printenv | sort             # tout afficher
# Pour Windows PowerShell : $env:PATH += ";C:\php"
\end{terminal}

Fichiers de configuration du shell (\`a \'editer une fois, puis \cmd{source})~:

\begin{tabz}{llX}
  \headrow \thd{Shell} \& \thd{Fichier} \& \thd{Quand} \\
  \midrule
  bash \& \fichier{\textasciitilde/.bashrc}, \fichier{\textasciitilde/.bash\\\\\_profile} \& interactif / connexion \\
  zsh (macOS par d\'efaut) \& \fichier{\textasciitilde/.zshrc} \& les deux \\
  fish \& \fichier{\textasciitilde/.config/fish/config.fish} \& syntaxe diff\'erente, tr\`es convivial \\
  PowerShell \& \cmd{Notepad \$PROFILE} \& Windows \\
\end{tabz}

Les variables d'environnement d'\textbf{un projet} ne vont pas dans le shell mais dans un
fichier \fichier{.env}, jamais versionn\'e, accompagn\'e d'un \fichier{.env.example} versionn\'e, lui.
C'est le contrat standard en Node comme en PHP~; voir \cref{sec:dotenv,ch:php-bases}.

\subsection{Les shells qui changent la vie}

\begin{itemize}
  \item \link{Starship}{https://starship.rs/fr/} --- invite de commande rapide, multi-langage,
        qui affiche automatiquement la version de Node/PHP/Java du dossier courant.
        \cmd{curl -sS https://starship.rs/install.sh | sh}
  \item \link{zsh + Oh My Zsh}{https://ohmyz.sh/} --- raccourcis, th\`emes, compl\'etions.
  \item \link{fish}{https://fishshell.com/} --- pr\^et \`a l'emploi, auto-suggestions.
  \item \link{Windows Terminal}{https://github.com/microsoft/terminal} +
        \link{Nerd Fonts}{https://www.nerdfonts.com/} pour les ic\^ones.
\end{itemize}

\section{Git : sauvegarder, historiser, collaborer}

Git n'est pas un nuage~: c'est une base de donn\'ees d'instants de ta vie de projet.
Trois semaines de code sans Git, c'est trois semaines \`a risque z\'ero de r\'ecup\'eration.
Documentation officielle~: \link{Pro~Git (livre gratuit)}{https://git-scm.com/book/fr/v2}.\src{https://git-scm.com/book/fr/v2}

\etape{Configuration une bonne fois pour toutes}
\begin{terminal}
git config --global user.name  "Ton Nom"
git config --global user.email "toi@example.com"
git config --global init.defaultBranch main
git config --global core.autocrlf input      # "true" sous Windows
git config --global core.editor "code --wait"
git config --global pull.rebase false
git config --global alias.st status
git config --global alias.lg "log --oneline --graph --decorate --all"
\end{terminal}

\etape{Le cycle quotidien}
\begin{terminal}
git init                     # une seule fois
git status                   # QUELS fichiers ont change ? (a taper 50 fois par jour)
git add src/app.js           # mettre en scene (index)
git diff --cached            # relire ce qui sera commite
git commit -m "feat: ajoute la validation de l'email"
git log --oneline --graph
\end{terminal}

Convention de message \link{Conventional~Commits}{https://www.conventionalcommits.org/fr/v1.0.0/}~:
\verb|type(scope): description imp\'erative|. Types~: \verb|feat|, \verb|fix|, \verb|docs|,
\verb|style|, \verb|refactor|, \verb|perf|, \verb|test|, \verb|chore|. Un historique propre est
un outil de debugging~: \cmd{git bisect} trouve la commit fautive en \$\log\_2(n)\$ \'etapes.

\etape{Travailler avec GitHub / GitLab}
\begin{terminal}
git remote add origin git@github.com:utilisateur/projet.git
git branch -M main
git push -u origin main          # premier push, puis simplement : git push
git switch -c ma-fonctionnalite  # creer une branche de travail
git switch main && git pull --rebase origin main
git rebase main                  # mettre ma branche a jour sans merge moche
git cherry-pick abc1234          # rejouer une commit precise
git restore --source=HEAD~3 src/app.js   # recuperer une version anterieure d'un fichier
git stash push -u -m "wip"       # mettre de cote sans commiter ; git stash pop
git reset --soft HEAD~1          # defaire le dernier commit en gardant les modifications
git revert abc1234               # defaire un commit deja pousse (sur, cree une nouvelle commit)
git clean -fd                    # supprimer les fichiers non suivis (attention)
\end{terminal}

\begin{piege}
\cmd{git push --force} sur une branche partag\'ee efface le travail des autres. Utilise
\cmd{--force-with-lease}, qui refuse si quelqu'un a pouss\'e entre-temps. Et ne force
jamais sur \verb|main|.
\end{piege}

\etape{Le \fichier{.gitignore} de base, valable pour les trois langages}
\begin{codeconf}
# dependances
node_modules/
vendor/
target/
build/
dist/
out/

# environnements et secrets
.env
.env.local
*.pem
secrets/

# outils
.idea/
.vscode/*
!.vscode/settings.json
!.vscode/extensions.json
.DS_Store
Thumbs.db
*.log
coverage/
.phpunit.result.cache
.phpstan.cache/
.php-cs-fixer.cache
__pycache__/

# telephones / emulators
.gradle/
local.properties
\end{codeconf}

G\'en\`ere-le en ligne~: \link{toptal.gitignore.io}{https://www.toptal.com/developers/gitignore}
(nomme \og ~gitignore.io~ \fg{} dans beaucoup de tutos).

\subsection{GitHub --- la bo\^ite \`a outils gratuite pour d\'eveloppeurs}

\begin{itemize}
  \item Un d\'ep\^ot public = portfolio. Mets un \fichier{README.md} soign\'e, une capture d'\'ecran,
        la proc\'edure d'installation, et une section \og ~ce que j'ai appris~ \fg{}.
  \item \link{gh~(GitHub~CLI)}{https://cli.github.com/}~: ouvre un PR, lit les logs CI sans navigateur.
\end{itemize}

\begin{terminal}
gh auth login
gh repo create mon-projet --public --source=. --push
gh pr create --fill
gh run watch                  # suivre la pipeline CI en direct
gh issue list --state open
\end{terminal}

\begin{itemize}
  \item \link{Gist}{https://gist.github.com/} pour partager un extrait de code.
  \item \link{Discussions}{https://docs.github.com/en/discussions} et
        \link{Pages}{https://pages.github.com/} pour h\'eberger un site statique gratuit.
  \item \link{Actions}{https://github.com/features/actions} pour l'int\'egration continue
        (exemple complet en \cref{sec:ci-cd}).
  \item \link{Copilot}{https://github.com/features/copilot} gratuit pour les \'etudiants et
        les maintainers de projets open source populaires.
\end{itemize}

\section{Structurer un projet et \'ecrire un README}

\begin{codeplain}
mon-projet/
  README.md            <- pourquoi, comment installer, comment lancer, captures
  LICENSE              <- MIT si tu veux que ton code soit librement reutilisable
  .gitignore
  .editorconfig        <- indentation et fins de ligne partagees avec l'equipe
  .env.example
  docs/                <- decisions d'architecture, captures, maquettes
  src/                 <- le code qui tourne en production
  tests/               <- ce qui prouve que le code tourne
  scripts/             <- outils internes (deploiement, generation de donnees)
\end{codeplain}

\begin{codeconf}
# .editorconfig --- a installer comme extension dans ton editeur
root = true

[*]
charset = utf-8
end_of_line = lf
insert_final_newline = true
trim_trailing_whitespace = true
indent_style = space
indent_size = 2

[*.php]
indent_size = 4

[*.java]
indent_size = 4

[Makefile]
indent_style = tab
\end{codeconf}

Le README d'un projet d'apprentissage doit contenir quatre choses~: ce que \c{c}a fait
(avec une capture ou un GIF), comment le lancer en trois commandes, ce que tu as appris,
et ce que tu referais autrement. C'est ce que regarde un recruteur~: \emph{le processus},
pas seulement le r\'esultat.

\begin{exercice}
\begin{enumerate}
  \item Dans le terminal, cr\'ee \verb|~/labo/js|, \verb|~/labo/php|, \verb|~/labo/java|,
        place un \fichier{.gitignore} et un \fichier{.editorconfig} dans chacun.
  \item \'Ecris trois commandes qui cr\'eent un fichier, l'affichent, puis le suppriment.
  \item Cherche avec \cmd{grep} toutes les occurrences du mot \verb{console.log} dans
        un dossier et compte-les avec \cmd{wc -l}.
  \item Fais un d\'ep\^ot GitHub public nomm\'e \verb|labo-js|. Pousse, cr\'ee une branche,
        un commit, une pull request depuis la CLI \cmd{gh}, puis fusionne-la en local.
  \item Casse volontairement un fichier, puis restaure-le avec \cmd{git restore}~;
        supprime-le, restaure-le avec \cmd{git checkout HEAD\textasciitilde1 -- chemin}.
\end{enumerate}
\end{exercice}

\begin{ressources}
\begin{itemize}
  \item Pro Git (livre complet, gratuit, en fran\c{c}ais) --- \url{https://git-scm.com/book/fr/v2}
  \item GitHub Skills (cours officiels gratuits) --- \url{https://skills.github.com/}
  \item Learn Git Branching (jeu visuel pour comprendre le graphe) --- \url{https://learngitbranching.js.org/?locale=fr\\\\_FR}
  \item The Missing Semester (MIT) --- \url{https://missing.csail.mit.edu/2020/course-material/}
  \item Shellcheck (analyse tes scripts bash) --- \url{https://www.shellcheck.net/}
\end{itemize}
\end{ressources}

\index{terminal}\index{git}\index{PATH}
% ============================================================================
\chapter{M\'ethode d'apprentissage et plan de 12 semaines}
\label{ch:methode}
% ============================================================================

Savoir quoi apprendre ne suffit pas~: il faut un rythme, une fa\c{c}on de stocker l'information
et un m\'ecanisme anti-abandon. Ce chapitre est court mais c'est celui qui d\'etermine si tu finis
le livre.

\section{Les cinq principes qui font la diff\'erence}

\begin{enumerate}
  \item \textbf{Pratique d\'elib\'er\'ee, pas r\'ep\'etition.} Travailler juste au-dessus de ton niveau,
        avec retour imm\'ediat (le message d'erreur, le test qui passe). Refaire un exercice
        d\'ej\`a r\'eussi est du r\'econfort, pas de l'apprentissage.
  \item \textbf{R\'ecup\'eration active.} Ferme le livre et \'ecris de m\'emoire la syntaxe d'un
        \verb|try/catch|, la diff\'erence entre \verb|==| et \verb|===|. Se tromper ancre
        dix fois mieux que relire.
  \item \textbf{Espacement.} Dix minutes de r\'evision chaque jour valent mieux que deux heures
        le dimanche. Un tour de flashcards le matin, l'erreur de la veille au soir.
  \item \textbf{Technique Feynman.} Explique \`a voix haute, avec des mots simples, comme si tu
        parlais \`a un enfant de douze ans. Tout ce que tu n'arrives pas \`a expliquer est une
        lacune --- retourne lire.
  \item \textbf{Production avant consommation.} Ratio cible~: 30 \% lire/regarder,
        70 \% \'ecrire du code. Si ta journ\'ee a \'et\'e 90 \% vid\'eo, tu ne progresses pas, tu te distraies.
\end{enumerate}

\section{La routine quotidienne (60 \`a 120 minutes)}

\begin{tabz}{llX}
  \headrow \thd{Dur\'ee} \& \thd{Bloc} \& \thd{Contenu} \\
  \midrule
  10 min \& R\'eveil \& 10 flashcards + relire le code \'ecrit hier, sans l'\'editeur. \\
  25 min \& Lecture active \& Un morceau de chapitre, le carnet ouvert, on pr\'edit les r\'esultats. \\
  50 min \& Code \& Recopier les exemples, les casser, refaire l'exercice du chapitre. \\
  15 min \& Journaliser \& \'Ecrire 5 lignes~: ce que j'ai compris, o\`u j'ai bloqu\'e, la solution. \\
  10 min \& Curiosit\'e \& Lire la doc officielle de ce que je viens d'utiliser (pas celle d'hier). \\
\end{tabz}

Si tu n'as que 25 minutes~: fais le bloc \og ~Code~ \fg{}. Jamais le bloc \og ~Lecture~ \fg{} seule.

\section{Le carnet d'erreurs --- ton actif le plus pr\'ecieux}

Cr\'ee \fichier{carnet.md} dans chaque projet et note~: \emph{texte exact de l'erreur},
\emph{contexte}, \emph{cause r\'eelle}, \emph{correction}, \emph{date}. En six mois tu auras un
manuel de d\'epannage personnel qui vaut plus que n'importe quel tuto. Les annexes
(\cref{ch:antisecche}) en sont une version g\'en\'erique~: tiens la tienne \`a jour.

\begin{codeconf}
## 2026-01-14 - PHP - 500 blank page
Erreur : aucune, page blanche.
Cause : une exception ignoree car display_errors = Off en production.
Fix : log_errors=On, error_log=php://stderr, et activer app_dev en local.
Lecon : ne jamais chercher a l'aveugle ; ouvrir les logs d'abord.
\end{codeconf}

\section{Plan de 12 semaines (version intensive)}

Adapte la vitesse~: doubler la dur\'ee de chaque ligne est parfaitementOK. Le plan n'est pas
une course, c'est un ordre.

\begin{tabz}{llX}
  \headrow \thd{Semaine} \& \thd{Objet} \& \thd{Livrable concret \`a la fin de la semaine} \\
  \midrule
  1 \& Outils, bases JS (\cref{ch:js-bases}) \& 30 petits scripts dans \verb|labo/js|, dont FizzBuzz, tri, palindrome. \\
  2 \& Objets, tableaux, cha\^ines, regex (\cref{ch:js-struct}) \& Un analyseur de texte (nombre de mots, mots les plus fr\'equents). \\
  3 \& S\'emantique avanc\'ee (\cref{ch:js-avance}) \& Un mini-syst\`eme d'\'ev\'enements \'ecrit \`a la main (20 lignes). \\
  4 \& Asynchrone (\cref{ch:js-async}) \& Un script qui t\'el\'echarge 20 pages web en parall\`ele avec limitation de concurrence. \\
  5 \& DOM et \'ev\'enements (\cref{ch:js-dom,ch:js-events}) \& Une page \og ~liste de courses~ \fg{} sans framework, avec filtre et tri. \\
  6 \& Fetch, formulaires, stockage (\cref{ch:js-dom,ch:js-storage}) \& La m\^eme liste, persist\'ee en LocalStorage + API externe. \\
  7 \& Node.js, modules, npm, CLI (\cref{ch:node,ch:js-outils}) \& Un outil de ligne de commande qui renomme en lot les fichiers d'un dossier. \\
  8 \& TypeScript + tests + s\'ecurit\'e (\cref{ch:js-ts,ch:js-tests,ch:js-secu}) \& Le projet de la semaine~6 typ\'e, test\'e \`a 80 \%, sans XSS. \\
  9 \& PHP bases, OOP, formulaires (\cref{ch:php-bases,ch:php-oop,ch:php-http}) \& Un mini-CRUD sur fichier JSON, puis sur MySQL. \\
  10 \& PDO, sessions, auth, s\'ecurit\'e PHP (\cref{ch:php-db,ch:php-http,ch:php-secu}) \& Un blog complet avec inscriptions, r\^oles, upload d'images s\'ecuris\'e. \\
  11 \& Java bases, OOP, collections, exceptions (\cref{ch:java-bases,ch:java-oop,ch:java-collections,ch:java-collections}) \& Un gestionnaire de biblioth\`eque console avec tests JUnit. \\
  12 \& Streams, JDBC, Spring Boot, d\'eploiement (\cref{ch:java-fonctionnel,ch:java-db,ch:java-spring,ch:deploy}) \& Une API REST Spring Boot connect\'ee \`a PostgreSQL, d\'eploy\'ee sur un VPS. \\
\end{tabz}

\section{Utiliser l'IA sans s'handicaper}

Un assistant de code est un excellent tuteur et un tr\`es mauvais substitut. R\`egle simple~:
\emph{l'IA r\'epond \`a tes questions, elle ne fait pas tes exercices.}

\begin{itemize}
  \item \textbf{Demander}~: \og ~explique-moi ligne par ligne ce message d'erreur~ \fg{}, \og ~quels sont les
        cas limites de cette fonction~ \fg{}, \og ~propose 3 variantes de style, avec les avantages~ \fg{}.
  \item \textbf{Refuser}~: coller l'\'enonc\'e d'un exercice et prendre la r\'eponse. Tu \'economises
        trois minutes et tu perds la comp\'etence.
  \item \textbf{Toujours}~: relire ligne par ligne le code g\'en\'er\'e et le faire tourner. Ne
        mets jamais dans un projet du code que tu ne sais pas expliquer.
  \item Outils~: \link{Copilot}{https://github.com/features/copilot},
        \link{Cursor}{https://www.cursor.com/}, \link{Claude}{https://claude.ai/},
        \link{ChatGPT}{https://chat.openai.com/}, \link{aider (CLI)}{https://aider.chat/}.
\end{itemize}

\section{Quand tu bloques --- la proc\'edure en cinq temps}

\begin{enumerate}
  \item \textbf{Lire l'erreur en entier}, y compris la ligne \og ~caused by~ \fg{}, et chercher
        le texte exact entre guillemets.
  \item \textbf{R\'eduire.} Fais un exemple minimal qui reproduit le probl\`eme (deux fichiers,
        dix lignes). Neuf fois sur dix, la r\'eduction r\'ev\`ele la cause.
  \item \textbf{Isoler une hypoth\`ese}, la tester, l'\'ecrire. \og ~Je pense que X. Je teste Y.~ \fg{}
  \item \textbf{Chercher dans la doc officielle} (les liens de fin de chapitre) plut\^ot que
        dans un blog de 2014.
  \item \textbf{Demander}, avec l'exemple minimal, la version de ton outil et le message complet
        --- sur \link{Stack~Overflow}{https://stackoverflow.com/questions/tagged/javascript},
        \link{Open~Classrooms}{https://openclassrooms.com/forum},
        \link{Dev~Discord/Slack}{https://github.com/websites/DevChatter} ou le Discord officiel
        d'un projet. Une question bien pos\'ee re\c{c}oit une r\'eponse en dix minutes.
\end{enumerate}

\begin{note}
\og ~Je tourne en rond depuis 45 minutes~ \fg{} est le signal pour lever les yeux de l'\'ecran~: marche,
douche, sommeil. Ce n'est pas une perte de temps, c'est la phase d'incubation du d\'ebogage.
\end{note}

\section{Sept familles d'erreurs que tu rencontreras}

\begin{enumerate}
  \item \textbf{Chemins}~: fichier trouv\'e sur ta machine, introuvable en production.
  \item \textbf{Encodage}~: accents cass\'es (UTF-8 partout, \verb|utf8mb4| en MySQL).
  \item \textbf{Ordre d'ex\'ecution}~: async non attendu, variable utilis\'ee avant initialisation.
  \item \textbf{Environnement}~: version d'outil, variable manquante, cache du navigateur.
  \item \textbf{Confusion de type}~: \verb|"1"+1|, \verb|0 == ""|, \verb|null| contre \verb|undefined|.
  \item \textbf{\'Etat partag\'e}~: mutation d'un objet pass\'e par r\'ef\'erence, effet de bord cach\'e.
  \item \textbf{Pr\'ematurit\'e}~: optimiser, abstraire, g\'en\'eraliser avant que le besoin existe.
\end{enumerate}

Tu peux les scanner une par une~: elles couvrent environ 80 \% des blocages de d\'ebutant.

\begin{exercice}
\begin{enumerate}
  \item \'Ecris ton contrat d'apprentissage~: jours, horaires, livrables, sanction en cas de
        manquement. Signe-le, date-le, colle-le au mur.
  \item Pr\'epare les trois fichiers \fichier{carnet.md} (un par langage) et le d\'ep\^ot
        \verb|apprentissage|.
  \item Fais une fiche de 5 cartes m\'emoire sur ce que tu as appris cette semaine
        (\link{Anki}{https://apps.ankiweb.net/} est l'outil standard, gratuit partout sauf iOS).
  \item Explique \`a voix haute, en deux minutes, la diff\'erence entre un langage interpr\'et\'e
        et un langage compil\'e. Enregistre-toi si tu oses~; r\'e\'ecoute.
\end{enumerate}
\end{exercice}

\begin{ressources}
\begin{itemize}
  \item \emph{Eloquent JavaScript} (livre gratuit de r\'ef\'erence) --- \url{https://eloquentjavascript.net/}
  \item \emph{The Road to Learn React / PHP / Java} --- \url{https://roadmap.sh/}
  \item MDN en fran\c{c}ais --- \url{https://developer.mozilla.org/fr/docs/Web/JavaScript}
  \item Exercices~: \url{https://exercism.org/}, \url{https://adventofcode.com/}, \url{https://www.codewars.com/}
  \item M\'emorisation espac\'ee~: \url{https://apps.ankiweb.net/}
\end{itemize}
\end{ressources}

\index{methode}\index{plan d'etudes}

\part{JavaScript : le langage du web et des serveurs}
% ============================================================================
\chapter{Les bases de JavaScript}
\label{ch:js-bases}
% ============================================================================

\begin{objectif}
\begin{itemize}
  \item Ex\'ecuter du JavaScript dans un navigateur, dans la console, et dans Node.js.
  \item Ma\^itriser variables, types, op\'erateurs, conversion de types et tests de v\'erit\'e.
  \item \'Ecrire des conditions, des boucles et des fonctions dans toutes les formes modernes.
  \item Comprendre la port\'ee (\emph{scope}), le hoisting et les fermetures (\emph{closures}).
  \item Reconna\^itre les pi\`eges de typage faible (\verb|==|, \verb|NaN|, \verb|this|).
\end{itemize}
\end{objectif}

\section{O\`u et comment ex\'ecuter du JavaScript}

JavaScript est un \emph{moteur}~: V8 dans Chrome et Node, SpiderMonkey dans Firefox,
JavaScriptCore dans Safari. Le langage est d\'efini par la norme \textbf{ECMAScript}
(\link{ECMA-262}{https://262.ecma-international.org/}), \'edit\'ee chaque ann\'ee~; les fonctionnalit\'es
d'ES2015 (dite ES6) \`a ES2024 sont aujourd'hui disponibles partout dans les navigateurs r\'ecents.

\begin{tabz}{llX}
  \headrow \thd{Terrain de jeu} \& \thd{Comment} \& \thd{Id\'eal pour} \\
  \midrule
  Console du navigateur \& \keys{F12} ou \keys{Ctrl+Shift+J} \& Explorer, tester une expression, apprendre le langage. \\
  Fichier HTML \& \verb|<script src="app.js"></script>| \& Tout ce qui touche la page (chapitres DOM). \\
  \link{Playground en ligne}{https://playcode.io/} / \link{CodePen}{https://codepen.io/} / \link{JSFiddle}{https://jsfiddle.net/} \& Ouvrir l'URL \& Partager un bug, prototyper. \\
  Node.js \& \cmd{node fichier.js} ou \cmd{node} (REPL) \& Logique, scripts, outillage, serveurs. \\
  \cmd{deno run} / \cmd{bun} \& \link{deno.com}{https://deno.com/}, \link{bun.sh}{https://bun.sh/} \& Alternatives modernes, rapide, sans configuration. \\
\end{tabz}

\begin{terminal}
# Le REPL (Read-Eval-Print Loop) : la calculatrice geante
$ node
> 2 ** 10
1024
> "nour".toUpperCase()
'NOUR'
> .exit

# Un fichier
$ echo 'console.log("Bonjour le monde");' > hello.js
$ node hello.js
Bonjour le monde
\end{terminal}

\section{Sortir et entrer de l'information}

\begin{codejs}
// En page : la sortie est la console du navigateur
console.log("normal");            // trace
console.info("info"); console.debug("detaill\'e");
console.warn("attention"); console.error("erreur");
console.table([{ a: 1, b: 2 }, { a: 3, b: 4 }]);   // un tableau joliment affiche
console.group("Etapes"); console.log("1"); console.groupEnd();
console.count("appelee");                            // compte les appels
console.time("boucle"); for (let i = 0; i < 1e6; i++); console.timeEnd("boucle");

// Dans un terminal, avec Node :
process.stdout.write("sans saut de ligne");
const reponse = await import("node:readline/promises")
  .then((m) => m.createInterface({ input: process.stdin, output: process.stdout }))
  .then((rl) => rl.question("Ton pr\'enom ? "));
console.log(`Salut ${reponse}`);
\end{codejs}

\section{Variables : \texttt{let}, \texttt{const} (et pourquoi pas \texttt{var})}

\begin{codejs}
const MAX = 100;          // lien immuable : on ne peut pas re-affectuer
let compteur = 0;         // peut changer
var ancienStyle = true;   // ne l'utilise plus : portee de fonction, hoisting piegeux

// "const" ne rend pas la VALEUR immutable, seulement la RE-ASSIGNATION impossible :
const user = { name: "Nour" };
user.name = "Awa";          // parfaitement legal
user = {};                  // TypeError: Assignment to constant variable.
\end{codejs}

R\`egle d'\'ecriture~: \texttt{const} par d\'efaut, \texttt{let} quand la r\'eassignation est n\'ecessaire,
\texttt{var} jamais. Nomme en \verb|camelCase| les variables et fonctions, en
\verb|SCREAMING\_SNAKE\_CASE| les constantes de module, en \verb|PascalCase| les classes.

\subsection{Les huit types primitifs}

\begin{tabz}{lXl}
  \headrow \thd{Type} \& \thd{Description} \& \thd{Exemple} \\
  \midrule
  \verb|string| \& S\'equence UTF-16 immuable \& \verb|"nour"|, \verb|'x'|, \verb|`a\${b}c`| \\
  \verb|number| \& Flottant 64 bits (~\$2^{53}\$ entiers exacts) \& \verb|42|, \verb|3.14|, \verb|1e3|, \verb|0xff|, \verb|1\_000\_000| \\
  \verb|bigint| \& Entier de taille arbitraire \& \verb|9007199254740993n| \\
  \verb|boolean| \& \verb|true| / \verb|false| \& \\
  \verb|undefined| \& Valeur non fournie \& \verb|let x;| \\
  \verb|null| \& Absence volontaire de valeur \& \verb|document.getElementById("x")| \\
  \verb|symbol| \& Cl\'e unique et immuable \& \verb|Symbol("id")| \\
\end{tabz}

Tout le reste est un \textbf{objet} (fonctions, tableaux, \verb|Map|, \verb|Date|, \verb|Error|\ldots).

\begin{codejs}
typeof 42;              // "number"
typeof "a";             // "string"
typeof true;            // "boolean"
typeof undefined;       // "undefined"
typeof Symbol();        // "symbol"
typeof 10n;             // "bigint"
typeof function () {};  // "function"
typeof {};              // "object"
typeof null;            // "object"  <-- BUG historique de 1995, jamais corrige. NE T'Y FIE PAS.
typeof [];              // "object"  <-- un tableau EST un objet
Array.isArray([]);      // true  <-- le VRAI test de tableau

Number.isInteger(5);        // true
Number.isNaN(NaN);          // true   (isNaN() global est faux : isNaN("toto") => true)
Number.isSafeInteger(2**53);// false
0.1 + 0.2 === 0.3;          // false  <-- arithmetique flottante : 0.30000000000000004
Math.abs(0.1 + 0.2 - 0.3) < Number.EPSILON;   // true : la bonne facon de comparer
\end{codejs}

\begin{piege}
\textbf{Trois pi\`eges de type, pour la vie.}
(1) \verb|typeof null === "object"|~: teste avec \verb|x === null|.
(2) \verb|0.1 + 0.2 !== 0.3|~: en mon\'etaire, travaille en centimes entiers ou avec une
biblioth\`eque d\'ecimale (\link{decimal.js}{https://github.com/MikeMcl/decimal.js/}).
(3) \verb|NaN| n'est \'egal \`a rien, pas m\^eme \`a lui-m\^eme~: utilise \verb|Number.isNaN(x)|.
\end{piege}

\section{Op\'erateurs}

\begin{codejs}
// Arithmetiques
7 / 2; 7 % 2; 2 ** 10; -(2); Math.trunc(7 / 2); Math.floor(-7 / 2);
(5.7).toFixed(2);              // "5.70"  -> RETOURNE UNE CHAINE
parseInt("42abc", 10); parseFloat("3.14"); Number("42"); +"42";

// Affectations combinees
let n = 10;
n += 5; n -= 2; n *= 3; n **= 2; n /= 4; n %= 7; n **= 2;
n++; ++n;                       // postfixe renvoie l'ancienne valeur, prefixe la nouvelle

// Comparaison --- la seule regle qui compte : TOUJOURS === et !==
5 == "5";        // true   (coercion : convertit avant de comparer)  INTERDIT
5 === "5";       // false  (comparaison stricte : type + valeur)
null == undefined;   // true
null === undefined;  // false
NaN === NaN;         // false

// Logiques
true && false;    // false  (court-circuit : le 2e operande n'est pas evalue)
true || false;    // true
!true;            // false
// Les formes modernes, tres utiles :
user?.profile?.avatar ?? "/default.png";   // chainage optionnel + alternative nulle
count ||= 10;     // assigne si falsy
config ??= {};    // assigne seulement si null/undefined  (le bon choix)

// Bit a bit (rare en web, courant en protocole/poids lourds)
5 & 3, 5 | 3, 5 ^ 3, ~5, 5 << 2, 5 >>> 1;

// Concatenation et template literal
"bonjour" + " " + nom;
`Bonjour ${nom}, tu as ${age} ans et ${[1, 2, 3].join(", ")} fruits.`;
\end{codejs}

\subsection{Conversion et \og ~falsy~ \fg{}}

Sept valeurs sont \textbf{falsy}~: \verb|false|, \verb|0|, \verb|-0|, \verb|0n|,
\verb|""|, \verb|null|, \verb|undefined|. Tout le reste est \textbf{truthy} -- y compris
\verb|[]|, \verb|{}| et \verb|"0"|.

\begin{codejs}
Boolean("");        // false
Boolean("0");       // true  <-- attention aux chaines non vides
if (tab.length) { /* ... */ }         // idiome courant
const nom = prenom || "Anonyme";      // remplace tout falsy (0, "", false)
const age = ageSaisi ?? 18;           // ne remplace que null/undefined  <-- RECOMMANDE
String(42); Number("42"); 42 + "";  +"42";
[1, "2", "3"] .map(Number);            // [1, 2, 3]
\end{codejs}

\section{Conditions et boucles}

\begin{codejs}
const note = 14;

if (note >= 16)       console.log("Tres bien");
else if (note >= 10)  console.log("Admis");
else                  console.log("Repasse");

// Expression conditionnelle (ternaire) : une VALEUR, pas un bloc
const mention = note >= 16 ? "tres bien" : note >= 10 ? "assez bien" : "insuffisant";

// switch : arrets obligatoires, sinon le code du cas suivant s'execute aussi
switch (jour) {
  case "samedi":
  case "dimanche":
    console.log("week-end");
    break;
  default:
    console.log("jour ouvrable");
}

// switch moderne : expression qui renvoie une valeur, correspondance par motifs
const tarif = switch (type) {
  case "enfant":        yield 5;
  case "adulte":        yield 12;
  default:              yield 8;
};                      // (actif a partir de ES2025 / Node 22+, sinon : ternaires ou objet)

const prix = { enfant: 5, adulte: 12, _default: 8 }[type] ?? 8;  // alternative simple et portative

for (let i = 0; i < 10; i++) { if (i === 3) continue; if (i === 8) break; }
let j = 0; while (j < 10) { j++; }
let k = 0; do { k++; } while (k < 10);   // execute au moins une fois

for (const v of [10, 20, 30]) {}          // itere sur les VALEURS (iterables)
for (const i in tab) {}                    // itere sur les CLES (chaines) -- a eviter sur tableaux
for (const [i, v] of tab.entries()) {}     // index + valeur
\end{codejs}

\begin{piege}
\verb|for...in| parcourt aussi les propri\'et\'es h\'erit\'ees \'enum\'erables et renvoie des
\emph{cha\^ines}. Sur un tableau, utilise \verb|for...of|, \verb|forEach| ou une boucle
indice. \verb|break| et \verb|continue| ne fonctionnent pas dans \verb|forEach|~: si tu as besoin
d'interrompre, c'est le signe qu'il te faut \verb|for...of|, \verb|some|, \verb|find| ou \verb|every|.
\end{piege}

\section{Fonctions}

\begin{codejs}
// Declaration (hoistee : utilisable avant sa definition)
function additionne(a, b = 0) { return a + b; }

// Expression de fonction (non hoistee)
const multiply = function (a, b) { return a * b; };

// Fleche : syntaxe compacte, pas de "this" propre
const square = (x) => x * x;
const toUser = (name, age) => ({ name, age, adult: age >= 18 });  // objet : parentheses obligatoires
const somme = (...args) => args.reduce((a, b) => a + b, 0);       // reste des arguments

// Arguments nommes + valeurs par defaut + valeurs de test
function envoi({
  destinataire,
  sujet = "(sans objet)",
  pieces = [],
  priorite = "normale",
} = {}) {
  if (!destinataire) throw new TypeError("destinataire requis");
  return { destinataire, sujet, pieces: pieces.length, priorite };
}
envoi({ destinataire: "nour@example.com", pieces: ["facture.pdf"] });

// Fonctions a rappels (callbacks) et fonctions d'ordre superieur
const appliquer = (fn, valeur) => fn(valeur);
const multiplieur = (facteur) => (x) => x * facteur;   // fabrique de fonctions
const triple = multiplieur(3); triple(5);               // 15

// R\'ecursivite : toujours prevoir un cas de base
function factorielle(n) {
  if (n < 0) throw new RangeError("n doit etre >= 0");
  return n <= 1 ? 1 : n * factorielle(n - 1);
}
\end{codejs}

\subsection{Port\'ee, hoisting, fermetures}

\begin{description}
  \item[Port\'ee de bloc] \verb|let| et \verb|const| vivent dans les accolades les plus proches.
        \verb|var| remonte jusqu'\`a la fonction~; c'est la source de bugs classiques dans les boucles.
  \item[Hoisting] Les d\'eclarations \verb|function| et \verb|var| sont \og ~remont\'ees~ \fg{} en m\'emoire.
        \verb|let|/\verb|const| existent mais sont dans la \og ~zone morte~ \fg{}~: y toucher avant
        l\`eve \verb|ReferenceError|.
  \item[Fermeture (closure)] Une fonction se souvient du environnement dans lequel elle a \'et\'e
        cr\'e\'ee. C'est le fondement de la vie priv\'ee en JS, des fabriques, de la m\'emo\"isation,
        des compteurs.
\end{description}

\begin{codejs}
// Le bug classique de var, a comprendre une fois pour toutes
for (var i = 0; i < 3; i++) setTimeout(() => console.log("var:", i)); // 3 3 3
for (let j = 0; j < 3; j++) setTimeout(() => console.log("let:", j)); // 0 1 2

// Fermeture : un compteur avec veritable encapsulation
function creeCompteur(demarrage = 0, pas = 1) {
  let valeur = demarrage;                       // prive : inaccessible de l'exterieur
  return {
    incrementer: () => (valeur += pas),
    lire: () => valeur,
    reinitialiser: () => (valeur = demarrage),
  };
}
const c = creeCompteur(10, 2);
c.incrementer(); c.incrementer(); c.lire();        // 14
c.valeur;                                          // undefined -- c'est tout l'interet

// Memoisation grace a une fermeture
function memoise(fn, cle = (...a) => JSON.stringify(a)) {
  const cache = new Map();
  return (...args) => {
    const k = cle(...args);
    if (!cache.has(k)) cache.set(k, fn(...args));
    return cache.get(k);
  };
}
const fib = memoise(function f(n) { return n < 2 ? n : f(n - 1) + f(n - 2); });
fib(70);   // immediat
\end{codejs}

\section{Le module commun : erreurs, piles, d\'ebogage}

\begin{codejs}
function diviser(a, b) {
  if (b === 0) throw new RangeError("division par zero");
  return a / b;
}
try {
  const r = diviser(10, 0);
} catch (err) {
  if (err instanceof RangeError) console.error("Domaine invalide :", err.message);
  else throw err;                        // ne JAMAIS avaler une erreur inconnue
} finally {
  console.log("toujours execute (fermer une connexion, logger, etc.)");
}

// Empilement des causes (ES2022) : precieux en production
try {
  JSON.parse("{oops");
} catch (cause) {
  throw new Error("Configuration invalide", { cause });
}
\end{codejs}

D\'eboguer proprement~: un point d'arr\^et dans l'onglet \clic{Sources} de Chrome (\keys{F12}),
puis \keys{F10} (pas \`a pas principal), \keys{F11} (entrer), \keys{F8} (continuer)~;
\verb|debugger;| dans le code ; \verb|console.trace()| ; \verb|console.log(new Error().stack)|.
La quasi-totalit\'e des d\'ebutants abuse de \verb|console.log|~: trois points d'arr\^et bien plac\'es
valent cent traces.

\begin{exercice}
\begin{enumerate}
  \item \textbf{FizzBuzz}~: affiche de 1 \`a 100, \og ~Fizz~ \fg{} si multiple de 3, \og ~Buzz~ \fg{} de 5,
        \og ~FizzBuzz~ \fg{} des deux. \'Ecris-le trois fois~: avec \verb|if|, avec un ternaire imbriqu\'e,
        avec une composition de fonctions.
  \item \textbf{Pair ou impair}, \textbf{plus grand de trois nombres}, \textbf{reverse d'une cha\^ine}
        (avec et sans \verb|split/reverse/join|).
  \item \textbf{Palindrome normalis\'e}~: \verb|"\'Esope reste ici et se repose"| renvoie \verb|true|.
        Retirer ponctuation, espaces, accents (\verb|String.prototype.normalize("NFD")|).
  \item \textbf{Nombre premier}~: fonction \verb|estPremier(n)| avec boucle jusqu'\`a
        \$\sqrt{n}\$~; puis \verb|premiersInfrieursA(100)|.
  \item \textbf{Temperature}~: \verb|celsiusVersFahrenheit|, \verb|fahrenheitVersCelsius|,
        et le sens inverse, \'ecrites comme fleches. V\'erifie que aller-retour redonne la valeur.
  \item \textbf{Compteur de clics}~: avec la fermeture \verb|creeCompteur|, simule 5 clics
        et affiche la valeur \`a chaque fois.
  \item \textbf{Devine le r\'esultat} avant d'ex\'ecuter~: \verb|1 + "2"|, \verb|"5" - 3|,
        \verb|"a" + + "b"|, \verb|[] + []|, \verb|[] + {}|, \verb|null + 1|,
        \verb|undefined + 1|, \verb|true + true|, \verb|2 + 3 + "a"|, \verb|"a" + 2 + 3|.
\end{enumerate}
\end{exercice}

\begin{ressources}
\begin{itemize}
  \item MDN --- Guide JavaScript (la bible, en fran\c{c}ais) --- \url{https://developer.mozilla.org/fr/docs/Web/JavaScript/Guide}
  \item \emph{JavaScript.info}, le meilleur cours libre du web --- \url{https://javascript.info/}
  \item \emph{Eloquent JavaScript}, 3\textsuperscript{e} \'edition --- \url{https://eloquentjavascript.net/}
  \item Table de compatibilit\'e des fonctionnalit\'es --- \url{https://compat-table.github.io/compat-table/es2016plus/}
  \item \emph{You Don't Know JS Yet} (Deepankar) --- \url{https://github.com/getify/You-Dont-Know-JS}
\end{itemize}
\end{ressources}

\index{JavaScript!bases}\index{variables}\index{closure}
% ============================================================================
\chapter{Objets, tableaux et cha\^ines de caract\`eres}
\label{ch:js-struct}
% ============================================================================

\begin{objectif}
\begin{itemize}
  \item Manipuler des objets : acc\`es, \'enum\'eration, copie, fusion, propri\'et\'es calcul\'ees.
  \item Ma\^itriser les 25 m\'ethodes de tableau qui couvrent 95 \% du travail r\'eel.
  \item \'Ecrire et lire des expressions r\'eguli\`eres utiles (validation, recherche, remplacement).
  \item Travailler en immutabilit\'e avec le spread et la d\'estructuration.
  \item Construire un mini-moteur de recherche et un panier de commande, de A \`a Z.
\end{itemize}
\end{objectif}

\section{Objets}

\begin{codejs}
const user = {
  prenom: "Nour",
  age: 28,
  "cle avec espace": 42,
  0: "index numerique converti en chaine",
  [Symbol("id")]: 7,
  activite: { ville: "Thies", pays: "SN" },
  saluer() { return `Bonjour ${this.prenom}`; },       // methode abregee
  get majeur() { return this.age >= 18; },              // accesseur
  set email(v) { this._email = String(v).toLowerCase(); },
};

// Acces
user.prenom;
user["cle avec espace"];
user.activite?.ville;
const ville = user?.activite?.ville;
const pays = user?.["activite"]?.["pays"] ?? "inconnu";

// Ecrire / effacer
user.role = "admin";
delete user.age;

// Proteger une propriete
Object.defineProperty(user, "id", { value: 1, writable: false, enumerable: true });
Object.freeze(user);          // shallow : gele les proprietes de premier niveau
Object.isFrozen(user);

// Enumerer --- les trois niveaux
Object.keys(user);        // ["prenom", "0", "id", ...]   proprietes propres enumerables
Object.values(user);
Object.entries(user);     // [["prenom","Nour"], ...]  : la porte d'entree vers les boucles
Object.getOwnPropertyNames(user);   // y compris non enumerables
Reflect.ownKeys(user);              // y compris symbols

for (const [cle, valeur] of Object.entries(user)) {
  if (typeof valeur === "object" && valeur !== null) continue;
}

// Construire a partir de tableaux / entries
Object.fromEntries([["a", 1], ["b", 2]]);       // { a: 1, b: 2 }
new Map([["a", 1]]);
// Un Map est plus adapte que l'objet quand : cles non-string, taille variable,
// iteration ordonnee, comptage frequent. Un Set enleve les doublons.
const scores = new Map();
scores.set("nour", 42).set("awa", 55);
scores.get("inconnu");          // undefined
scores.has("awa"); scores.size;
[...scores.entries()]; [...scores.keys()]; [...scores.values()];
const uniques = [...new Set([1, 1, 2, 3, 3, 3])];        // [1, 2, 3]
\end{codejs}

\subsection{Copies, fusions, \'egalit\'e}

\begin{codejs}
const base = { a: 1, b: { c: 2 }, liste: [1, 2] };

const copieSuperficielle = { ...base };            // b et liste restent PARTAGES
const memeSousObjet = copieSuperficielle.b === base.b;   // true  <-- source de bugs !

const deep = structuredClone(base);                // natif (Node 17+, navigateurs 2022+)
deep.b.c = 99; base.b.c;                            // 2 : la copie est independante

const fusion = Object.assign({}, base, { a: 9, d: 4 });  // ancien, mais explicite
const fusion2 = { ...base, ...{ a: 9 }, ...{ d: 4 } };

// Alias et mutation : en JS, un objet est passe PAR REFERENCE
const un = { n: 1 }, deux = un; deux.n = 99; un.n;   // 99
const tab = [1, 2]; tab.push(3);                       // mute le tableau

// Egalite structurelle "a la main" (sans dependance)
const egales = (x, y) => {
  if (Object.is(x, y)) return true;
  if (typeof x !== "object" || typeof y !== "object" || x === null || y === null) return false;
  const kx = Object.keys(x), ky = Object.keys(y);
  return kx.length === ky.length && kx.every((k) => egales(x[k], y[k]));
};
egales({ a: [1, { b: 2 }] }, { a: [1, { b: 2 }] });   // true
\end{codejs}

\begin{piege}
\verb|JSON.parse(JSON.stringify(obj))| est une vieille astuce de copie profonde. Elle est
\textbf{lente} et \textbf{dangereuse}~: elle d\'etruit \verb|Date|, \verb|Map|, \verb|Set|,
\verb|undefined|, les fonctions, les \verb|BigInt| et les cycles. Utilise \verb|structuredClone|.
\end{piege}

\section{Tableaux --- la bo\^ite \`a outils}

\begin{codejs}
const produits = [
  { id: 1, nom: "Cafe",   prix: 12.5, qt: 3, tags: ["boisson", "matin"] },
  { id: 2, nom: "The",    prix: 8,    qt: 0, tags: ["boisson"] },
  { id: 3, nom: "Livre",  prix: 22,   qt: 1, tags: ["culture"] },
  { id: 4, nom: "Cahier", prix: 3.5,  qt: 8, tags: ["culture", "bureau"] },
];

// CREATION
Array.from({ length: 5 }, (_, i) => i * i);        // [0,1,4,9,16]
Array(3).fill(0);                                  // [0,0,0]  (jamais Array(3).map !)
[..."bonjour"];                                    // ["b","o","n","j","o","u","r"]
structuredClone(produits);

// LECTURE : les methodes qui repondent a une question
produits.length;
produits.at(-1);            // dernier element (aussi at(-2))
produits.includes(x);
produits.indexOf(obj) !== -1;    // N'utilise PAS pour des objets : comparaison d'identite
produits.find(p => p.id === 3);            // le premier -> objet | undefined
produits.findIndex(p => p.qt === 0);       // son index  -> -1 si absent
produits.toSorted((a, b) => a.prix - b.prix);   // copie triee (ES2023, non mutant)
produits.filter(p => p.qt > 0);
produits.some(p => p.prix > 20);       // au moins un ?
produits.every(p => p.prix > 0);       // tous ?
produits.reduce((acc, p) => acc + p.prix * p.qt, 0);   // un total
produits.findLast(p => p.tags.includes("culture"));    // depuis la fin

// ECRITURE : ce qui mute vs ce qui copie
const mute = [1, 2, 3];
mute.push(4); mute.pop(); mute.shift(); mute.unshift(0);
mute.splice(1, 2, "x");      // retire 2 elements a l'index 1 et inserx x
mute.sort((a, b) => a - b);  // TRI NUMERIQUE : sans comparateur, on obtient [0,1,10,2,...] !
mute.reverse(); mute.fill(0); mute.length = 0;

const copie = [1, 2, 3];
copie.toSorted((a, b) => b - a);
copie.toReversed();
const sansX = copie.with(0, 99);        // copie avec un element remplace

// Transformations et combinaisons
produits.map(p => ({ ...p, total: p.prix * p.qt }));
produits.flatMap(p => p.tags);           // puis .flat() / .flat(Infinity) pour aplatir
[...new Set(produits.flatMap(p => p.tags))];   // tags uniques
[produits, [1, 2]].flat();                // un niveau

// Groupage et comptage
Object.groupBy(produits, p => (p.qt > 0 ? "dispo" : "rupture"));
Map.groupBy(produits, p => p.tags.length);

// Iterer proprement
for (const [i, p] of produits.entries()) console.log(i, p.nom);
for (const p of produits) { if (p.qt === 0) break; }     // le seul ou break marche
\end{codejs}

\subsection{Tris et donn\'ees complexes}

\begin{codejs}
const eleves = [{ n: "Awa", m: 14 }, { n: "Nour", m: 9 }, { n: "Ibra", m: 17 }];

// 1. Tri multi-criteres stable (les tableaux se trient stablement depuis ES2019)
const parNotePuisNom = [...eleves].sort((a, b) => b.m - a.m || a.n.localeCompare(b.n, "fr"));

// 2. LocaleCompare est obligatoire pour les accents
["echo", "ecole", "effort"].sort();                              // depend du code point
["echo", "ecole", "effort"].sort((a, b) => a.localeCompare(b, "fr", { sensitivity: "base" }));

// 3. Dico de notes pour trier par un ordre metier
const ordre = { TresBien: 4, Bien: 3, Passable: 2 };
notes.sort((a, b) => ordre[b] - ordre[a]);

// 4. Recherche + filtre combin
const recherche = (q) => {
  const norm = q.trim().toLowerCase();
  if (!norm) return produits;
  return produits.filter((p) => p.nom.toLowerCase().includes(norm) || p.tags.some((t) => t.includes(norm)));
};
\end{codejs}

\begin{astuce}
R\`egle d'immutabilit\'e dans un projet frontend~: une fonction qui transforme des donn\'ees
\emph{retourne un nouveau tableau/objet} et ne mute pas l'entr\'ee. Cela rend le d\'ebogage trivial,
permet l'annulation (undo), et rend r\'eactifs les frameworks. \verb|arr.push(x)| devient
\verb|const next = [...arr, x]|.
\end{astuce}

\section{Cha\^ines de caract\`eres}

\begin{codejs}
const s = "  Le renard brun rapide saute par-dessus le chien paresseux  ";

s.length; s.trim(); s.trimStart(); s.trimEnd();
s.toLowerCase(); s.toUpperCase();
s.includes("renard"); s.startsWith("  "); s.endsWith("x");
s.indexOf("le"); s.lastIndexOf("le"); s.at(0);
s.slice(3, 8); s.slice(-5);            // decoupe ; negatifs autorises
s.replaceAll("le", "LE");               // chaines ou //g regex
s.split(" "); s.split(/\s+/);            // mots ; en tenant compte des espaces multiples
s.normalize("NFD").replace(/[\u0300-\u036f]/g, "");   // retire les accents
s.padStart(5, "0"); "42".padStart(3, "0");            // "042"
"abc".charCodeAt(0); String.fromCharCode(65, 66); String.fromCodePoint(0x1F600);
[..."a"].length;                      // 2 : le spread respecte les points de code
"ligne1\nligne2".split(/\r?\n/);         // robuste Windows/Linux
"a,b;;c".split(",").map(t => t.trim()).filter(Boolean);   // ["a","b","c"]

// Template literals multi-lignes et balises
const sql = (strings, ...vals) => strings.reduce((acc, str, i) => {
  const v = vals[i - 1];
  return acc + (v === undefined ? "" : `'${String(v).replaceAll("'", "''")}'`) + str;
});
sql`SELECT * FROM users WHERE nom = ${"O'Brien"}`;   // protege du injection la plus simple
\end{codejs}

\section{Expressions r\'eguli\`eres}

Une regex est un objet \verb|RegExp|. \'Ecris-la d'abord en litt\'eral \verb|/.../flags|~;
utilise le constructeur \verb|new RegExp(str)| quand le motif vient d'une variable.

\begin{tablez}[p{0.28\textwidth}X]
  \headrow \thd{Motif} \& \thd{Signification} \\
  \midrule
  \verb|\d \w \s| \& chiffre / lettre-digit-underscore / espace \\
  \verb|\D \W \S| \& n\'egations \\
  \verb|. ^ \$ | \& n'importe quel caract\`ere (sauf saut) / d\'ebut / fin \\
  \verb|[abc] [^abc] [a-z]| \& ensemble / n\'egation / plage \\
  \verb|* + ? {n,m} ?{1,}| \& 0+ / 1+ / 0-1 / r\'ep\'etition / non-gourmand \\
  \verb|() (?:) (?<nom>)| \& groupe capturant / non capturant / nomm\'e \\
  \verb|(?=) (?!) (?<=)| \& lookahead positif/n\'egatif, lookbehind \\
  \verb|\b \B \ | \& limite de mot / non-limite / echapper un metacaractere \\
  \verb|a|b| \& alternance \\
\end{tablez}

\begin{codejs}
const flags = { i: "insensible a la casse", g: "global (toutes les occurrences)",
  m: "multi-lignes", s: "le point matche les sauts", u: "unicode", y: "colle a lastIndex" };

const email = /^[a-z0-9._%+-]+@[a-z0-9.-]+\.[a-z]{2,}$/i;
email.test("nour@example.com");          // true

// Extraction avec groupes nommes
const entete = "Server: nginx/1.25.3 Date: Tue";
const m = entete.match(/Server:\s*(?<logiciel>[\w./-]+)/);
m.groups.logiciel;                        // "nginx/1.25.3"

// Remplacements malins
"jean-luc  picard".replace(/\b\w/g, (c) => c.toUpperCase());     // "Jean-Luc  Picard"
"2026-01-31".replace(/(\d{4})-(\d2)-(\d2)/, "$3/$2/$1");         // "31/01/2026"
"a1b2c3".replace(/\d/g, (d) => `[${d}]`);
"un  deux   trois".replace(/\s+/g, " ");

// Trouver tout, proprement, avec la position
const texte = "le chat et le chien et le loup";
for (const t of texte.matchAll(/\ble\s+(?<animal>\w+)/g)) {
  console.log(t.index, t.groups.animal);
}

// Validations utiles a recopier
const regles = {
  tel: /^\+?[0-9][0-9 .-]{7,14}$/,
  slug: /^[a-z0-9]+(?:-[a-z0-9]+)*$/,
  mdpFort: /^(?=.*[a-z])(?=.*[A-Z])(?=.*\d)(?=.*[\W_]).{12,}$/,
  ipv4: /^(25[0-5]|2[0-4]\d|1?\d?\d)(\.(25[0-5]|2[0-4]\d|1?\d?\d)){3}$/,
  hexColor: /^#(?:[0-9a-fA-F]{3}|[0-9a-fA-F]{6})$/,
};
const motDePasseValide = (mdp) => regles.mdpFort.test(mdp);
\end{codejs}

\begin{piege}
Deux erreurs fr\'equentes~: (1) oublier le drapeau \verb|g| avec \verb|replace| -- seul la premi\`ere
occurrence est remplac\'ee ; (2) une regex construite avec une entr\'ee utilisateur peut
provoquer un d\'eni de service catastrophique (\emph{ReDoS}). Ne passe jamais une saisie
brute dans un \verb|new RegExp()| sans la neutraliser avec
\verb|s.replace(/[.*+?^\${}()|[\]\\]/g, "\\\$\&")|~; et garde les motifs simples.
\end{piege}

\section{D\'estructuration, valeurs par d\'efaut, objets et tableaux}

\begin{codejs}
// Tableaux
const [premier, deuxieme] = produits;
const [, , troisieme] = produits;
const [tete, ...queue] = [1, 2, 3, 4];
const [a = "defaut"] = [];

// Objets (avec renommage et valeur par defaut)
const { nom: appellation, prix: prixUnitaire = 0, tags: [premierTag = "aucun"] } = produits[0];

// Fusion et suppression de proprietes
const avecRemise = { ...produits[0], prix: produits[0].prix * 0.9 };
const { qt: _ignorer, ...sansQt } = produits[0];       // "_" ou "_" prefixe = volonte d'ignorer

// Parametres de fonction
const carte = ({ x = 0, y = 0, zoom = 1 }) => `(${x},${y}) @${zoom}x`;
carte({ y: 10 });

// Affectation par decomposition (rare mais elegant)
let gauche, droite; ({ gauche, droite } = { gauche: 1, droite: 2 });
[gauche, droite] = [droite, gauche];                   // echange sans variable temporaire

// Parametres de rendu : props -> variables locales
function LigneProduit({ produit, quantite = 1, promo = false }) {
  const { nom, prix } = produit ?? {};
  return `${nom ?? "inconnu"} : ${(prix ?? 0) * (promo ? 0.8 : 1) * quantite} EUR`;
}
\end{codejs}

\section{TP complet --- un panier de commande}

\begin{codejs}
// panier.js --- a taper entierement, puis a lancer avec : node panier.js
const TVA = 0.2;

const catalogue = new Map([
  ["cafe",  { prix: 12.5, stock: 4, unite: "sachet" }],
  ["the",   { prix: 8,    stock: 0, unite: "boite" }],
  ["livre", { prix: 22,   stock: 6, unite: "exemplaire" }],
  ["cahier",{ prix: 3.5,  stock: 12, unite: "unite" }],
]);

const creerPanier = () => ({ lignes: [], coupons: [] });

const ajouter = (panier, ref, qte = 1) => {
  const article = catalogue.get(ref);
  if (!article) throw new Error(`Reference inconnue : ${ref}`);
  const existante = panier.lignes.find((l) => l.ref === ref);
  const nouvelleQte = (existante?.qte ?? 0) + qte;
  if (nouvelleQte > article.stock) throw new RangeError(`Stock insuffisant pour ${ref}`);
  return existante
    ? { ...panier, lignes: panier.lignes.map((l) => (l === existante ? { ...l, qte: nouvelleQte } : l)) }
    : { ...panier, lignes: [...panier.lignes, { ref, qte: nouvelleQte }] };
};

const retirer = (panier, ref) => ({ ...panier, lignes: panier.lignes.filter((l) => l.ref !== ref) });

const appliquerCoupon = (panier, code) => ({ ...panier, coupons: [...panier.coupons, code] });

const totaux = (panier) => {
  const lignes = panier.lignes.map((l) => {
    const a = catalogue.get(l.ref);
    return { ...l, prixUnitaire: a.prix, totalLigne: a.prix * l.qte };
  });
  const brut = lignes.reduce((s, l) => s + l.totalLigne, 0);
  const remise = brut * (panier.coupons.includes("BIENVENUE10") ? 0.1 : 0);
  const ht = brut - remise;
  return { lignes, nbArticles: lignes.reduce((s, l) => s + l.qte, 0), brut, remise, ht,
           tva: ht * TVA, ttc: ht * (1 + TVA) };
};

const format = (n) => n.toLocaleString("fr-FR", { style: "currency", currency: "EUR" });

const afficher = (panier) => {
  const t = totaux(panier);
  console.log("\n--- Panier ---");
  for (const l of t.lignes) console.log(`${String(l.qte).padStart(2)} x ${l.ref.padEnd(7)} ${format(l.prixUnitaire)}  = ${format(l.totalLigne)}`);
  console.table(t.lignes);
  console.log(`Sous-total ${format(t.brut)} | Remise ${format(t.remise)}`);
  console.log(`HT ${format(t.ht)} | TVA ${format(t.tva)} | TTC ${format(t.ttc)}`);
};

let p = creerPanier();
p = ajouter(p, "cafe", 2);
p = ajouter(p, "cahier", 3);
p = ajouter(p, "cafe");            // 3 sachets au total
p = retirer(p, "the");              // non present : aucune erreur
p = appliquerCoupon(p, "BIENVENUE10");
afficher(p);

// Exercice immediat : ajoute la gestion des frais de port gratuits au-dela de 50 EUR
// de commande, et un champ "pointsFidelite" egal a 1 point par euro TTC.
\end{codejs}

\begin{exercice}
\begin{enumerate}
  \item \verb|supprimerDoublons(tab)|~: retourne un nouveau tableau sans doublons, en
        pr\'eservant l'ordre, en une ligne (Set), puis sans Set.
  \item \verb|grouperPar(tab, cle)|~: \verb|[{ville:'Dakar'},...]| devient
        \verb|{ Dakar: [...], Thies: [...] }| (utilise \verb|reduce|, puis \verb|Object.groupBy|).
  \item \verb|motsFreqents|~: \`a partir d'un texte, renvoie les 10 mots les plus fr\'equents de
        plus de 3 lettres, avec compte, tri\'es, en ignorant la ponctuation et les accents.
  \item \verb|aplatir(obj)|~: \verb|{a:{b:{c:1}}}| devient \verb|{"a.b.c": 1}|, r\'ecursif, en
        g\'erant les tableaux et en \'evitant les cycles (Map \og ~vu~ \fg{}).
  \item \verb|chunk(tab, taille)|, \verb|zip(a, b)|, \verb|sans(tab, ...valeurs)|,
        \verb|intersection(a, b)|, \verb|difference(a, b)|~: les utilitaires que tout projet contient.
  \item \verb|masquerEmail("nour.diop@example.com")| \$\rightarrow\$ \verb|"no***@example.com"| (regex + slice).
  \item \textbf{Slugifieur}~: \og ~L'\'Ev\'enement de l'ann\'ee, \`a Thies !~ \fg{} \$\rightarrow\$ \verb|evenement-annee-thies|
        (minuscules, accents retir\'es, non-alphanum\'eriques en tirets, tirets dupliqu\'es retir\'es).
  \item \textbf{Validation}~: \'ecris un module \verb|validate.js| exportant \verb|email|,
        \verb|tel|, \verb|motDePasse| (avec message d'erreur humain) et \verb|formulaire(donnees)|
        qui renvoie \verb|{ ok, erreurs }|.
\end{enumerate}
\end{exercice}

\begin{ressources}
\begin{itemize}
  \item MDN --- \url{https://developer.mozilla.org/fr/docs/Web/JavaScript/Reference/Global\\\\_Objects/Array}
  \item Tableau comparatif des m\'ethodes de tableaux --- \url{https://javascript.info/array-methods}
  \item Constructeur et visualiseur de regex --- \url{https://regex101.com/} et \url{https://regexper.com/}
  \item \verb|lodash| (\`a lire pour les noms de m\'ethodes, pas pour d\'ependre de lui) --- \url{https://lodash.com/docs/}
  \item \emph{You Don't Know JS Yet~: this \& Object Prototypes} --- \url{https://github.com/getify/You-Dont-Know-JS}
\end{itemize}
\end{ressources}

\index{JavaScript!objets}\index{JavaScript!tableaux}\index{expressions regulieres}
% ============================================================================
\chapter{S\'emantique avanc\'ee : prototypes, \texttt{this}, classes, Proxy}
\label{ch:js-avance}
% ============================================================================

\begin{objectif}
\begin{itemize}
  \item Expliquer ce qu'est r\'eellement une classe en JavaScript (du sucre sur les prototypes).
  \item Pr\'edire la valeur de \texttt{this} dans n'importe quelle situation.
  \item Utiliser les champs priv\'es, accesseurs, h\'eritage, mixins, \texttt{static}.
  \item Intercepter des op\'erations avec \texttt{Proxy} et \texttt{Reflect} (r\'eactivit\'e, validation).
  \item \'Eviter la pollution de prototype et les confusions valeur/r\'ef\'erence.
\end{itemize}
\end{objectif}

\section{Propri\'et\'es : ce que voit le moteur}

Chaque propri\'et\'e poss\`ede un descripteur. Le conna\^itre explique \texttt{Object.freeze},
les acc\`es en \'ecriture silencieusement ignor\'es en mode strict, et les pi\`eges d'observable.

\begin{codejs}
const o = { nom: "Nour" };
Object.getOwnPropertyDescriptor(o, "nom");
// { value: "Nour", writable: true, enumerable: true, configurable: true }

Object.defineProperty(o, "age", { value: 28, writable: false, enumerable: true, configurable: true });
Object.defineProperties(o, {
  ville: { value: "Thies", writable: true, enumerable: true, configurable: true },
  id: { get() { return `${this.nom}-${this.age}`; }, enumerable: true, configurable: true },
});
"use strict";
o.age = 30;          // TypeError en mode strict (silencieux en mode relax : le pire des deux)

const drapeaux = { enumerable: false, configurable: false, writable: false };
Object.defineProperty(Array.prototype, "dernier", {
  get() { return this[this.length - 1]; },
  ...drapeaux,
});   // correct UNIQUEMENT dans un exercice : ne jamais toucher un prototype natif en prod
\end{codejs}

\section{La cha\^ine de prototypes}

Un objet a une cl\'e cach\'ee \texttt{[[Prototype]]} (visible via \texttt{Object.getPrototypeOf}).
Lire une propri\'et\'e absente remonte la cha\^ine jusqu'\`a \texttt{Object.prototype}, puis
\texttt{null}. \'Ecrire une propri\'et\'e cr\'ee toujours une propri\'et\'e \emph{propre}, jamais
sur le parent.

\begin{codejs}
const animal = {
  respire() { return "oui"; },
  decrire() { return `${this.nom} est un ${this.type}`; },
};
const chat = Object.create(animal);
chat.nom = "Minou"; chat.type = "chat";

chat.nom;                 // propre
chat.respire();           // trouve sur animal via la chaine
Object.getPrototypeOf(chat) === animal;         // true
chat instanceof animal;                         // true
animal.isPrototypeOf(chat);                     // true (prefere getPrototypeOf)
Object.hasOwn(chat, "nom");                     // true  (ES2022, le bon test)
"toString" in chat;                             // true  <-- "in" voit aussi l'herite !
Object.keys(chat);                              // ["nom","type"] : seulement les propres enumerables

class Chien {                 // la forme moderne ; "class" est du sucre sur ce qui precede
  static espece = "canis";    // champ statique (ES2022)
  #secret = "os cache";       // champ PRIVE veritable (inaccessible meme par reflection)
  nom;                        // declaration de champ
  constructor(nom) { this.nom = nom; }
  aboyer() { return `${this.nom}: ouf`; }
  get secretLong() { return this.#secret + " et demi"; }
  [Symbol.toPrimitive](hint) { return hint === "string" ? `Chien(${this.nom})` : 0; }
  static fromJSON(j) { return new this(j.nom); }
}
const medor = new Chien("M\'edor");
medor.#secret;               // SyntaxError : le prive est reellement prive
medor.secretLong;            // "os cache et demi"
`${medor}`;                  // "Chien(Medors)" via Symbol.toPrimitive
Object.getOwnPropertyNames(Chien.prototype);   // ["constructor","aboyer","secretLong"]
\end{codejs}

\begin{note}
Les m\'ethodes de classe sont d\'efinies sur \texttt{Classe.prototype}, donc partag\'ees par toutes
les instances (\'economie m\'emoire). Un \texttt{=>} d\'eclar\'e comme champ de classe devient une
propri\'et\'e \emph{par instance}~: il est recr\'e\'e \`a chaque \texttt{new}, mais garde \texttt{this} fig\'e.
C'est exactement ce qu'on veut pour les callbacks pass\'es \`a un composant UI.
\end{note}

\subsection{H\'eritage, mixins, composition}

\begin{codejs}
class Forme { #couleur = "noir"; constructor(c) { this.#couleur = c; } dessine() { return `[${this.#couleur}]`; } }
class Cercle extends Forme {
  constructor(c, r) { super(c); this.r = r; }                 // super() AVANT tout "this"
  get aire() { return Math.PI * this.r ** 2; }
  dessine() { return `${super.dessine()} cercle r=${this.r}`; }   // appel de la methode mere
}
// Le prototype de la classe est li\'e : Object.getPrototypeOf(Cercle) === Forme
Object.getPrototypeOf(Cercle.prototype) === Forme.prototype;  // true      // true

// Pas d'heritage multiple : on compose avec des mixins
const Labellisable = (Base) => class extends Base { toLabel() { return `${this.nom} (${this.type ?? "n/a"})`; } };
const Verifiable = (Base) => class extends Base {
  valider() { return Object.values(this).every((v) => v !== undefined && v !== null && v !== ""); }
};
class Produit extends Verifiable(Labellisable(Forme)) {}   // les mixins s'empilent

// Un objet d'utilitaires, pas une classe : souvent le meilleur choix
const dates = {
  today: () => new Date(),
  iso: (d = new Date()) => d.toISOString().slice(0, 10),
  diffJours: (a, b) => Math.round((b - a) / 86_400_000),
};
\end{codejs}

\section{\texttt{this} : les quatre r\`egles}

\texttt{this} n'est pas li\'e \`a l'endroit o\`u la fonction est \emph{\'ecrite} mais \`a l'endroit
o\`u elle est \textbf{appel\'ee}. Quatre r\`egles, par ordre de priorit\'e d\'ecroissant~:

\begin{enumerate}
  \item \textbf{new} : \verb|new F()| --> \texttt{this} est le nouvel objet.
  \item \textbf{Explicite} : \verb|f.call(obj)|, \verb|f.apply(obj, args)|, \verb|f.bind(obj)| --> \texttt{this} = \texttt{obj}.
  \item \textbf{Implicite} : \verb|obj.f()| --> \texttt{this} = \texttt{obj} (le point juste avant les parenth\`eses).
  \item \textbf{Par d\'efaut} : \texttt{undefined} en mode strict, \texttt{globalThis} sinon.
\end{enumerate}

\begin{codejs}
const bouton = {
  label: "Envoyer",
  cliquer() { return `on clique sur ${this.label}`; },
  cliquerFleche: () => `on clique sur ${this.label}`,   // fleche : this du module, pas du bouton !
  cliquerLie: null,
};
bouton.cliquer();                       // "on clique sur Envoyer"
const f = bouton.cliquer; f();          // TypeError : this === undefined (mode strict/module)
bouton.cliquerFleche();                 // "on clique sur undefined"  <-- PIEGE CLASSIQUE
bouton.cliquerLie = bouton.cliquer.bind(bouton);
bouton.cliquerLie();                    // "on clique sur Envoyer"

// Les trois appels explicites
function presentation(ville, pays) { return `${this.nom}, ${ville} (${pays})`; }
const user = { nom: "Nour" };
presentation.call(user, "Thies", "SN");
presentation.apply(user, ["Thies", "SN"]);
const p = presentation.bind(user, "Thies");    // pre-remplissage partiel (currying)
p("SN");

// Les API qui perdent "this" : toujours lier ou utiliser une fleche
class Chrono {
  debut = 0;
  demarrer = () => { this.debut = performance.now(); };   // bon : fleche = this fige a l'instance
  afficher() { return `${performance.now() - this.debut} ms`; }
}
const c = new Chrono();
setTimeout(c.demarrer, 100);      // OK grace a la fleche
// setTimeout(c.afficher, 100)    // TypeError : "this" perdu. Solution : c.afficher.bind(c)

// newFunction / this dans un objet literal imbrique
const table = {
  valeurs: [1, 2, 3],
  total() { return this.valeurs.reduce((a, b) => a + b, 0); },  // fonction classique : "this" correct
};
\end{codejs}

\begin{astuce}
R\`egle mn\'emotechnique infaillible~: regarde \emph{imm\'ediatement \`a gauche des parenth\`eses}
de l'appel. S'il y a un objet, c'est lui \texttt{this}. S'il n'y a rien, \texttt{this} est
\texttt{undefined} (module/strict). Une fonction fl\'ech\'ee ignore totalement cette r\`egle~:
elle capture le \texttt{this} lexical du lieu o\`u elle a \'et\'e \'ecrite.
\end{astuce}

\section{Prototypes : le c\^ot\'e sombre (pollution)}

Modifier un prototype partag\'e par \emph{tout le monde} est une faille de s\'ecurit\'e r\'eelle
(CVE historiques dans des biblioth\`eques connues) et un bug garanti.

\begin{codejs}
// A NE JAMAIS FAIRE
Array.prototype.last = function () { return this[this.length - 1]; };
//   -> "for...in" sur un tableau, des bibliotheques qui detectent "last", des conflits entre libs.

// Le vecteur d'attaque classique : fusion recursive d'objets non controles
function fusion(target, source) {
  for (const k in source) {
    if (k === "__proto__" || k === "constructor" || k === "prototype") continue;   // BLOQUE
    if (Object.prototype.toString.call(source[k]) === "[object Object]") {
      fusion(target[k] ??= {}, source[k]);
    } else target[k] = source[k];
  }
  return target;
}
// Avec JSON.parse(...) sur une entree utilisateur, sans ce filtre :
//   {"__proto__":{"isAdmin":true}}   ->   ({}).isAdmin === true  POUR TOUTES LES OBJETS !
Object.prototype.hasOwnProperty.call({}, "isAdmin");     // false, une fois prot\`ege
const sur = Object.create(null);        // objet sans prototype : parfait pour un dictionnaire pur
\end{codejs}

\`A retenir~: valider les cl\'es d'un objet venant de l'ext\'erieur, pr\'ef\'erer \texttt{Map} pour
un dictionnaire, et ne jamais \'ecrire de \texttt{for...in} sur des donn\'ees non fiables.

\section{Proxy et Reflect : intercepter le langage}

\begin{codejs}
const pairesValidees = new WeakMap();     // metadonnees clees par objet, sans fuite memoire

function objetValide(contrat) {
  const cible = {};
  return new Proxy(cible, {
    set(o, prop, val) {
      if (typeof prop !== "string") throw new TypeError("cles non-string interdites");
      const regle = contrat[prop];
      if (!regle) throw new Error(`Propriete inconnue : ${String(prop)}`);
      if (!regle.test(val)) throw new TypeError(`${prop}: valeur invalide (${val})`);
      return Reflect.set(o, prop, val);   // delegue proprement au comportement par defaut
    },
    get(o, prop) {
      if (prop === "nettoyer") return () => Reflect.ownKeys(o).forEach((k) => Reflect.deleteProperty(o, k));
      return Reflect.get(o, prop);
    },
    deleteProperty(o, prop) { throw new Error(`Suppression interdite : ${String(prop)}`); },
  });
}
const user = objetValide({
  age: { test: (v) => Number.isInteger(v) && v > 0 && v < 130 },
  email: { test: (v) => /^[^@\s]+@[^@\s]+\.[a-z]{2,}$/i.test(v) },
});
user.age = 28;          // OK
user.age = -3;          // TypeError: age: valeur invalide (-3)
user.telephone = 12;    // Error: Propriete inconnue : telephone

// LeProxy de reactivite (le principe de Vue, Svelte stores, MobX...)
function reactive(objet, onChange) {
  const handler = {
    get(o, k, recv) {
      const v = Reflect.get(o, k, recv);
      return typeof v === "object" && v !== null ? new Proxy(v, handler) : v;
    },
    set(o, k, v, recv) {
      const avant = o[k];
      const r = Reflect.set(o, k, v, recv);
      if (!Object.is(avant, v)) onChange(k, v, avant);
      return r;
    },
  };
  return new Proxy(objet, handler);
}
const etat = reactive({ user: { name: "Nour", tags: [] } }, (k, v) => console.log("changement", k, v));
etat.user.name = "Awa";       // changement name Awa
etat.user.tags.push("vip");   // capte aussi les mutations profondes

// Un Proxy de lecture securisee, pour remplacer les "if (a && a.b && a.b.c)"
const sur = (o, def = undefined) => new Proxy(o ?? {}, {
  get: (t, k) => (k in t ? (typeof t[k] === "object" && t[k] !== null ? sur(t[k], def) : t[k]) : def),
});
sur(undefined).a.b.c.d;        // undefined  (au lieu de TypeError)
\end{codejs}

\begin{note}
Un \texttt{Proxy} est puissant mais co\^uteux~: chaque lecture passe par un pi\`ege. On l'utilise
aux fronti\`eres (\'etat r\'eactif d'un store, objet de configuration, validation d'entr\'ee), jamais
dans une boucle \`a dix millions d'it\'erations.
\end{note}

\section{Valeur ou r\'ef\'erence, et les preuves}

\begin{codejs}
// Primitives = comparees par valeur ; objets = compares par reference (identite)
"abc" === "abc";                    // true
[1, 2] === [1, 2];                  // false
const a = [1], b = a; a === b;      // true : le MEME tableau
Object.is(NaN, NaN);                // true (contrairement a ===)
Object.is(+0, -0);                  // false

// Le cout de la reference : une mutation involontaire est un bug sans message d'erreur
function ajouterLigne(lignes, ligne) {
  lignes.push(ligne);               // MUTATION : l'appelant voit son tableau change
  return lignes;
}
function ajouterLigneImmut(lignes, ligne) { return [...lignes, ligne]; }   // la bonne habitude

// Le "shallow copy" ne protege qu'un niveau : la preuve en 3 lignes
const src = { liste: [1, 2] }, cop = { ...src };
cop.liste.push(3);                  // src.liste vaut [1,2,3] !
\end{codejs}

\begin{exercice}
\begin{enumerate}
  \item \'Ecris une classe \verb|CompteBancaire| avec~: champ priv\'e \verb|\#solde|, getter
        \verb|solde|, m\'ethode \verb|debiter(montant)| qui l\`eve une erreur si le solde devient
        n\'egatif, un \verb|historique| immuable, et une m\'ethode statique \verb|fromJSON|.
        Ajoute \verb|[Symbol.toPrimitive]| pour un affichage lisible.
  \item Ajoute une classe \verb|Comepargne| qui h\'erite, ajoute un taux et \'ecrase
        \verb|interets()| ; v\'erifie que \verb|super| est bien appel\'e et que \verb|instanceof|
        fonctionne dans les deux sens.
  \item Pr\'edis la valeur de \verb|this| pour cinq appels que tu inventes (m\'ethode, fl\'ech\'ee,
        callback perdu, \verb|bind|, \verb|call|), puis \'ecris un test unitaire qui confirme.
  \item \'Ecris un \verb|proxyValidation(schema)| r\'eutilisable~: il v\'erifie types et bornes
        d'un objet de configuration, affiche le chemin exact de la cl\'e fautive
        (\verb|db.pool.max|), et renvoie un objet fig\'e apr\`es validation.
  \item \'Ecris une fonction \verb|copieProfondeSensible(obj, { ignorer: [...], dates: true })|
        qui g\`ere \verb|Date|, \verb|Map|, \verb|Set|, \verb|Array|, les cycles
        (avec \verb|WeakMap|), et refuse les fonctions. Compare \`a \verb|structuredClone|.
  \item D\'emo de pollution de prototype~: construis un petit parseur de requ\^ete
        qui fusionne \verb|?a[b]=1| dans un objet, montre l'attaque, puis corrige-la.
\end{enumerate}
\end{exercice}

\begin{ressources}
\begin{itemize}
  \item MDN --- Heritage et prototypes --- \url{https://developer.mozilla.org/fr/docs/Learn/JavaScript/Objects/Object\\\_prototypes}
  \item MDN --- Classes --- \url{https://developer.mozilla.org/fr/docs/Web/JavaScript/Reference/Classes}
  \item MDN --- Proxy / Reflect --- \url{https://developer.mozilla.org/fr/docs/Web/JavaScript/Reference/Global\\\_Objects/Proxy}
  \item JavaScript.info --- \url{https://javascript.info/prototype-inheritance} (chapitres 8 et 21)
  \item \emph{Prototype pollution} (article de r\'ef\'erence) --- \url{https://medium.com/intrinsic/attacking-javascripts-prototype-pollution-7d2b66d64790}
\end{itemize}
\end{ressources}

\index{JavaScript!prototypes}\index{this}\index{classes}\index{Proxy}
% ============================================================================
\chapter{Asynchrone : boucle d'\'ev\'enements, promesses, \texttt{async/await}}
\label{ch:js-async}
% ============================================================================

\begin{objectif}
\begin{itemize}
  \item Expliquer ce qui se passe vraiment quand une requ\^ete est en cours.
  \item Cr\'eer, cha\^iner, combiner et annuler des promesses sans fuite ni \og ~promesse flottante~ \fg{}.
  \item \'Ecrire du code asynchrone lisible avec \verb|async/await|, en parall\`ele quand c'est possible.
  \item Impl\'ementer les motifs r\'eels~: retry, timeout, debounce, throttle, file \`a concurrence limit\'ee.
  \item Comprendre les it\'erateurs, g\'en\'erateurs et it\'erateurs asynchrones.
\end{itemize}
\end{objectif}

\section{Le mod\`ele d'ex\'ecution}

JavaScript est \textbf{mono-thread}~: un seul appel de fonction \`a la fois. Mais il est
\textbf{asynchrone et non bloquant}~: les op\'erations lentes (r\'eseau, disque, timer) sont d\'el\'egu\'ees
\`a des threads du syst\`eme (navigateur ou \link{libuv}{https://github.com/libuv/libuv} sous Node).
Leur r\'esultat est trait\'e quand la pile d'appel est vide.

\begin{enumerate}
  \item La \textbf{call stack} ex\'ecute le code synchrone jusqu'au bout.
  \item Les \textbf{API} (fetch, \verb|setTimeout|, fs) renvoient le contr\^ole imm\'ediatement ;
        la r\'eponse arrive plus tard, via un rappel plac\'e dans une file.
  \item La \textbf{microtask queue} (promesses r\'esolues, \verb|queueMicrotask|, MutationObserver)
        est vid\'ee \emph{enti\`erement} avant la suite~: chaque t\^ache a la priorit\'e.
  \item La \textbf{task queue} (timers, E/S, \'ev\'enements) attend que les microt\^aches soient vides,
        puis fournit \emph{une} t\^ache.
  \item \verb|requestAnimationFrame| s'ex\'ecute avant le rendu de l'image suivante.
\end{enumerate}

\begin{codejs}
console.log("1 synchrone");

setTimeout(() => console.log("6 timeout 0ms (macrotache)"), 0);
Promise.resolve().then(() => console.log("4 then"));
queueMicrotask(() => console.log("5 microtask"));
async function f() { console.log("2 debut async"); await null; console.log("3 apres await"); }
f();

// Ordre reel : 1, 2, 3, 4, 5, 6
// Regle mnemonique : tout le code synchrone, PUIS toutes les microtaches, PUIS une macrotache.
\end{codejs}

\begin{note}
\og ~JS est mono-thread donc il ne peut pas faire de calcul~ \fg{} est faux pour les calculs lourds~:
on les d\'eplace dans un \textbf{Web Worker} (navigueur) ou un \textbf{worker thread} (Node),
et on communique par messages. Voir \cref{sec:workers}.
\end{note}

\section{Rappels --- et pourquoi on les a quitt\'es}

\begin{codejs}
// Convention "error-first" heritee de Node : premier argument = erreur
function lireConfig(chemin, cb) {
  setTimeout(() => {
    if (!chemin) return cb(new Error("chemin requis"));
    cb(null, { chemin, ok: true });
  }, 30);
}
lireConfig("a.json", (err, res) => { if (err) throw err; /* ... */ });

// Le "pyramid of doom" est la vraie raison d'abandonner les callbacks imbriques
a(() => b(() => c(() => d(() => e(() => { /* on ne sait plus ou on est */ }))));
\end{codejs}

\section{Promesses}

Une promesse est un \textbf{objet qui repr\'esente une valeur pas encore disponible}, dans un des
trois \'etats~: \emph{pending}, \emph{fulfilled}, \emph{rejected} --- et qui ne change plus d'\'etat.

\begin{codejs}
const tache = new Promise((resolve, reject) => {
  setTimeout(() => resolve(42), 100);           // une seule resolution utile
});

tache
  .then((v) => v * 2)              // renvoyer une valeur -> nouvelle promesse resolue
  .then((v) => Promise.resolve(v + 1))   // renvoyer une promesse -> elle est "awaitee"
  .catch((e) => { console.error(e); throw new Error("echec traite"); })
  .finally(() => console.log("toujours"));   // ne transforme pas la valeur

// Les methodes de combinaison
const [u, posts] = await Promise.all([fetchUser(), fetchPosts()]);   // toutes ou echec
const resultats = await Promise.allSettled([p1, p2, p3]);  // toujours : [{status,value|reason}]
const premier = await Promise.race([requete, timeout]);            // le premier qui finit
const meilleur = await Promise.any([miroir1, miroir2]);            // le premier qui reussit
const p = Promise.withResolvers();            // {promise, resolve, reject} : propre a l'exterieur
const r = await Promise.try(() => fonctionPeutEtreSynchrone());    // n'importe quoi, uniformise

// await sur n'importe quoi (valeur, thenable, tableau)
const x = await 5;                 // 5
for (const v of await Promise.all(tab.map(async (i) => i * 2))) {}
\end{codejs}

\subsection{\'Ecrire des promesses correctes}

\begin{codejs}
// 1. Ne jamais creer une promesse "a la main" autour d'une API deja promise
const mauvais = new Promise((res) => res(fetch(url)));      // anti-pattern
const bon = fetch(url);

// 2. Promisify utile pour les API a callback
const promisifier = (fn) => (...args) =>
  new Promise((resolve, reject) =>
    fn(...args, (err, val) => (err ? reject(err) : resolve(val)))
  );
const tard = (ms) => new Promise((r) => setTimeout(r, ms));
await Promise.all([tard(100), tard(50)]);     // 100 ms (parallele), pas 150

// 3. Toujours gerter l'erreur, sinon : "unhandledRejection" (crash par defaut sous Node)
process.on("unhandledRejection", (reason) => { console.error("NON GREEE", reason); });

// 4. Une promesse n'est pas annulable : on ajoute un AbortController
const ctrl = new AbortController();
const reponse = await fetch(url, { signal: ctrl.signal });
ctrl.abort(new DOMException("Timeout", "TimeoutError"));  // -> fetch rejette

// 5. Enveloppe timeout reutilisable
async function avecTimeout(p, ms, message = "Timeout") {
  const ctrl = new AbortController();
  const t = setTimeout(() => ctrl.abort(new Error(message)), ms);
  try {
    return await Promise.race([
      p,
      new Promise((_, rej) => ctrl.signal.addEventListener("abort", () => rej(ctrl.signal.reason))),
    ]);
  } finally { clearTimeout(t); }
}
\end{codejs}

\section{\texttt{async/await}}

\verb|async function| renvoie toujours une promesse~; \verb|await| met la fonction en pause
et rend la main \`a la boucle d'\'ev\'enements.

\begin{codejs}
async function chargerProjet(id) {
  try {
    const rep = await fetch(`/api/projects/${id}`);
    if (!rep.ok) throw new Error(`HTTP ${rep.status} sur ${rep.url}`);
    const projet = await rep.json();
    const membres = await Promise.all(              // parallelisme : on attend les deux a la fois
      projet.membreIds.map((mid) => fetch(`/api/users/${mid}`).then((r) => r.json()))
    );
    return { ...projet, membres };
  } catch (err) {
    if (err.name === "AbortError") throw new Error("Requete annulee par l'utilisateur", { cause: err });
    throw err;                                       // remonter ce qu'on ne comprend pas
  } finally {
    masquerSpinner();
  }
}
\end{codejs}

\begin{tablez}[p{0.2\textwidth}XX]
  \headrow \thd{Anti-pattern} \& \thd{Probl\`eme} \& \thd{Correction} \\
  \midrule
  \verb|await p1; await p2;| (ind\'ependants) \& Temps cumul\'e, latence doubl\'ee \& \verb|await Promise.all([p1, p2])| \\
  \verb|tab.map(async x => await f(x))| \& concurrence non born\'ee (10 000 requ\^etes d'un coup) \& fichier d'attente born\'ee (ci-dessous) \\
  \verb|forEach(async x => ...)| \& la boucle ne \og ~attend~ \fg{} pas ; les erreurs sont perdues \& \verb|for...of| + \verb|await|, ou \verb|Promise.all(map)| \\
  \verb|.then(f).catch(g)| avec throw dans g \& \og ~unhandled rejection~ \fg{} \& attraper au plus pr\`es ou renvoyer \\
  \verb|catch {}| vide \& Bug invisible, donn\'ees corrompues \& au minimum \verb|console.error| ou \verb|throw| \\
  fonction non \verb|async| qui appelle une promesse \& valeur \verb|Promise \{pending\}| affich\'ee \& remonter \verb|async| jusqu'\`a l'appelant \\
\end{tablez}

\subsection{Trois utilitaires \`a avoir dans son projet}

\begin{codejs}
// 1. Retry avec recul exponentiel et jitter (la base de la resilience reseau)
async function retry(fn, { essais = 4, base = 200, max = 4000, surErreur } = {}) {
  let dernier;
  for (let i = 0; i < essais; i++) {
    try { return await fn(i); }
    catch (e) {
      dernier = e;
      surErreur?.(e, i);
      if (i === essais - 1) break;
      const attente = Math.min(max, base * 2 ** i) * (0.5 + Math.random() * 0.5);
      await new Promise((r) => setTimeout(r, attente));
    }
  }
  throw dernier;
}

// 2. Pool de concurrence : N taches en parallele, pas plus
async function parallelLimite(taches, limite = 5) {
  const resultats = new Array(taches.length);
  let curseur = 0;
  const worker = async () => {
    while (curseur < taches.length) {
      const i = curseur++;
      resultats[i] = await taches[i]();
    }
  };
  await Promise.all(Array.from({ length: Math.min(limite, taches.length) }, worker));
  return resultats;
}

// 3. Debounce / throttle : les deux fonctions les plus demandees en UI
const debounce = (fn, ms = 250) => {
  let t;
  const wrapped = (...args) => new Promise((res) => {
    clearTimeout(t);
    t = setTimeout(() => res(fn(...args)), ms);
  });
  wrapped.cancel = () => clearTimeout(t);
  return wrapped;
};
const throttle = (fn, ms = 100) => {
  let dernier = 0, minuterie = null, dernArgs;
  return (...args) => {
    const maintenant = Date.now();
    dernArgs = args;
    if (maintenant - dernier >= ms) { dernier = maintenant; fn(...args); }
    else if (!minuterie) minuterie = setTimeout(() => {
      dernier = Date.now(); minuterie = null; fn(...dernArgs);
    }, ms - (maintenant - dernier));
  };
};

const surRecherche = debounce(async (q) => { /* appel API */ }, 350);
document.querySelector("#q").addEventListener("input", (e) => surRecherche(e.target.value));
\end{codejs}

\section{G\'en\'erateurs, it\'erateurs et it\'erateurs asynchrones}

Un \textbf{it\'erateur} est un objet avec \verb|next()| renvoyant \verb|{ value, done }| ;
un \textbf{it\'erable} expose \verb|[Symbol.iterator]()|. \verb|for...of|, le spread et
\verb|Array.from| ne fonctionnent qu'avec \c{c}a.

\begin{codejs}
// Rendre un objet iterable (une playlist)
const playlist = {
  titres: ["a", "b", "c"],
  *[Symbol.iterator]() { for (const t of this.titres) yield t.toUpperCase(); },
};
[...playlist];                    // ["A","B","C"]

// Sequence infinie + paresse
function* entiersPairs(de = 0) { for (let n = de; ; n += 2) yield n; }
const dix = [];
for (const n of entiersPairs(2)) { if (dix.push(n) === 10) break; }

// Pipeline de transformation (generator-based)
function* map(iterable, fn) { for (const x of iterable) yield fn(x); }
function* filtre(iterable, pred) { for (const x of iterable) if (pred(x)) yield x; }
const lignes = [1, 2, 3, 4, 5, 6, 7, 8, 9, 10];
const sommeDesCarresImpairs = [...filtre(map(lignes, (x) => x * x), (x) => x % 2)]
  .reduce((a, b) => a + b, 0);

// Iterateur asynchrone : lire un tres gros fichier ligne a ligne sans le charger en RAM
async function* lignesFichier(url) {
  const res = await fetch(url);
  const decod = new TextDecoderStream("utf-8");
  const flux = res.body.pipeThrough(decod).pipeThrough(new TextEncoderStream ? new TransformStream() : new TransformStream());
  let tampon = "";
  for await (const chunk of res.body.pipeThrough(new TextDecoderStream("utf-8"))) {
    tampon += chunk;
    let idx;
    while ((idx = tampon.indexOf("\n")) >= 0) { yield tampon.slice(0, idx); tampon = tampon.slice(idx + 1); }
  }
  if (tampon) yield tampon;
}
// pour (const ligne of await ... ) n'existe pas : on ecrit
for await (const ligne of lignesFichier("https://exemple.gouv.fr/gros-fichier.csv")) { /* ... */ }
\end{codejs}

\subsection{Le flux (Streams) --- l'id\'ee \`a retenir}

Un \verb|ReadableStream| se lit morceau par morceau~; un \verb|WritableStream| s'\'ecrit de m\^eme ;
\verb|pipeThrough| et \verb|TransformStream| les cha\^inent. C'est ce qui permet de
t\'el\'echarger et transformer 5~Go sans jamais les mettre en m\'emoire, de streamer une r\'eponse
depuis Express ou un mod\`ele de langage, et de parser du SSE. On l'utilise concr\`etement dans
les chapitres Node et API.

\section{Chronologie, temps r\'eel et minuterie}

\begin{codejs}
const id = setInterval(() => tick(), 1000);   // recurrent
clearInterval(id);
let t = setTimeout(() => envoi(), 300);       // unique
clearTimeout(t);

// setTimeout(fn, 0) n'est pas "immediat" : c'est "prochaine macrotache"
// queueMicrotask(fn) est "des que la pile est vide" (avant le rendu)
// requestAnimationFrame(fn) est "juste avant le prochain affichage" (~16 ms)
function boucle() { dessiner(); requestAnimationFrame(boucle); }

// Mesurer proprement (pas Date.now)
const t0 = performance.now();
attente();
console.log(`${(performance.now() - t0).toFixed(2)} ms`);
\end{codejs}

\begin{exercice}
\begin{enumerate}
  \item \'Ecris \verb|attendre(ms)| puis transforme une fonction callback \verb|lireFichier| en
        version promise \verb|lireFichierAsync| (avec \verb|promisifier|~; puis \verb|fs/promises|).
  \item Pr\'edis l'ordre d'affichage, puis v\'erifie~:
\begin{codeplain}
console.log("A");
setTimeout(() => console.log("B"), 0);
Promise.resolve().then(() => console.log("C"));
console.log("D");
(async () => { console.log("E"); await null; console.log("F"); })();
console.log("G");
\end{codeplain}
  \item Charge les 26 lettres depuis \verb|https://api.npmjs.org/downloads/point/last-week/|
        combin\'ees avec une API locale~: utilise \verb|Promise.allSettled|, affiche le top 10 des
        paquets, et ne plante pas quand une requ\^ete \'echoue.
  \item Impl\'emente \verb|crawler(startUrl, maxPages)|~: il r\'ecup\`ere une liste pagin\'ee
        \verb|?page=n|, en parall\`ele \`a 5 maximum, s'arr\^ete sur page vide, ajoute un
        \verb|retry| \`a 3 essais et un \verb|timeout| de 5 s.
  \item \'Ecris un \verb|rateLimit(fn, appelsParSeconde)| qui r\'epartit les appels (token bucket),
        puis un \verb|pQueue(tasks, { concurrency, pause, resume })| minimal.
  \item Ajoute un bouton \og ~Annuler~ \fg{} \`a une recherche~: brancher un \verb|AbortController|,
        annuler la requ\^ete pr\'ec\'edente \`a chaque nouvelle frappe, et ignorer proprement
        \verb|AbortError|.
  \item \'Ecris une file d'attente de messages pour un chat~: \verb|emit|/\verb|on|, avec
        \verb|once|, \verb|off|, et un g\'en\'erateur asynchrone \verb|for await (const msg of bus)|.
\end{enumerate}
\end{exercice}

\begin{ressources}
\begin{itemize}
  \item MDN --- Promises --- \url{https://developer.mozilla.org/fr/docs/Web/JavaScript/Reference/Global\\\\_Objects/Promise}
  \item MDN --- Event loop et microtasks --- \url{https://developer.mozilla.org/en-US/docs/Web/API/HTML\\\\_DOM\\\\_API/Microtask\\\\_guide}
  \item Visualiseur interactif de la boucle d'\'ev\'enements --- \url{https://www.jsv9000.app/}
  \item Node --- Streams --- \url{https://nodejs.org/api/stream.html}
  \item Jake Archibald --- \emph{In The Loop} (la r\'ef\'erence sur l'ordre d'ex\'ecution) --- \url{https://jakearchibald.com/2015/events-based-runtimes/}
\end{itemize}
\end{ressources}

\index{JavaScript!asynchrone}\index{promesse}\index{async await}\index{event loop}
% ============================================================================
\chapter{Le navigateur : HTML, CSS et le DOM}
\label{ch:js-dom}
% ============================================================================

\begin{objectif}
\begin{itemize}
  \item Poser une page HTML correcte, accessible, et y brancher du JavaScript proprement.
  \item S\'electionner, cr\'eer, modifier, supprimer des noeuds sans casser les performances.
  \item Comprendre le rendu (reflow/repaint) et \'ecrire du code qui reste fluide \`a 60 fps.
  \item Utiliser les API modernes~: \texttt{dataset}, \texttt{classList}, \texttt{dialog},
        \texttt{details}, \texttt{IntersectionObserver}.
  \item Livrer un premier mini-projet complet, sans framework.
\end{itemize}
\end{objectif}

\section{Le squelette d'une page moderne}

\begin{codeplain}
<!DOCTYPE html>
<html lang="fr">
<head>
  <meta charset="utf-8">
  <meta name="viewport" content="width=device-width, initial-scale=1">
  <title>Liste de courses -- Nour</title>
  <meta name="description" content="Une liste de courses qui persiste dans le navigateur.">
  <meta name="color-scheme" content="light dark">
  <link rel="icon" href="/favicon.svg" type="image/svg+xml">
  <link rel="stylesheet" href="style.css">
  <script src="app.js" defer></script>            <!-- defer = execute apres le parse du HTML -->
</head>
<body>
  <a class="sr-only" href="#contenu">Aller au contenu</a>
  <header>
    <h1>Mes courses</h1>
  </header>
  <main id="contenu">
    <form id="form-ajout" autocomplete="off">
      <label for="nom">Article</label>
      <input id="nom" name="nom" required maxlength="60" placeholder="Riz, tomates...">
      <label for="qte">Quantit\'e</label>
      <input id="qte" name="qte" type="number" min="1" max="99" value="1">
      <button type="submit">Ajouter</button>
    </form>
    <ul id="liste" aria-live="polite"></ul>
    <template id="tpl-article">
      <li class="article">
        <input type="checkbox" class="check">
        <span class="nom"></span>
        <button class="suppr" type="button" aria-label="Supprimer">x</button>
      </li>
    </template>
  </main>
  <footer><small><span id="compteur"></span></small></footer>
</body>
</html>
\end{codeplain}

Quatre d\'etails qui font la diff\'erence~: \cmd{lang="fr"} (prononciation, c\'esure, recherche),
la balise \verb|viewport| (sinon ton site reste \og ~zoom\'e~ \fg{} sur mobile), \verb|defer|
(le script attend le HTML, sans bloquer le parsing) et le \verb|<template>| (du HTML stock\'e,
inerte, clonable \`a la demande -- l'anti-\verb|innerHTML| par excellence).

\begin{tabz}{lX}
  \headrow \thd{Attribut de script} \& \thd{Comportement} \\
  \midrule
  \verb|<script src="a.js">| \& bloque le parsing jusqu'au t\'el\'echargement et \`a l'ex\'ecution. \\
  \verb|defer| \& t\'el\'echargement en parall\`ele, ex\'ecution dans l'ordre, apr\`es le parsing. \textbf{Le choix par d\'efaut pour un fichier externe.} \\
  \verb|async| \& t\'el\'echargement en parall\`ele, ex\'ecution d\`es que pr\^et, ordre al\'eatoire. Pour la pub, l'analytics. \\
  \verb|type="module"| \& diff\'er\'e par d\'efaut, port\'ee de module, \verb|import| disponibles. \textbf{Le choix pour du JS en modules.} \\
  \verb|inlin\'e <script>| \& ex\'ecut\'e imm\'ediatement~: \`a placer en \emph{fin de body} si n\'ecessaire. \\
\end{tabz}

\section{S\'electionner}

\begin{codejs}
const box = document.querySelector(".boite");       // 1er element correspondant
const tous = document.querySelectorAll(".article");  // NodeList statique (tableau-like)
const parId = document.getElementById("liste");      // le plus rapide, a privilegier si tu as un id
const parClasse = document.getElementsByClassName("article");   // HTMLCollection VIVANTE (piege)
const parNom = document.getElementsByTagName("li");

// Selections structurelles
const parent = box.closest("main");          // remonte jusqu'au premier ancetre qui matche
const suivant = box.nextElementSibling;
const enfants = [...box.children];           // HTMLCollection -> vrai tableau
const dansForm = document.querySelectorAll("form input:not([type=hidden])");
const visible = [...document.querySelectorAll(".article")].filter((el) => el.checkVisibility());
box.matches("main > .boite");                 // test de selecteur

// Les selecteurs utiles a connaitre par coeur
// :not(), :is(h1,h2,h3), :where(), :has(> input:checked), :nth-child(2n+1), :scope, ::selection
document.querySelectorAll("li:has(.check:checked)");     // les lignes cochees (CSS + JS, tres puissant)
\end{codejs}

\begin{note}
\texttt{querySelectorAll} renvoie une copie fig\'ee~: si tu ajoutes un \'el\'ement apr\`es, il n'y sera
pas. \texttt{getElementsByClassName} renvoie une collection \emph{vivante}~: la modifier pendant
une boucle produit des sauts d'index. Dans le doute~: \verb|[...document.querySelectorAll(...)]|.
\end{note}

\section{Lire et modifier un \'el\'ement}

\begin{codejs}
const li = document.querySelector(".article");

// TEXTE : textContent securise (les <b> sont echappes), innerHTML non (XSS)
li.textContent = "Riz <5kg>";              // affiche le texte literal, y compris les chevrons
li.innerHTML = "<b>Riz</b> 5kg";           // parse du HTML : UNIQUEMENT avec du contenu de confiance
li.insertAdjacentHTML("beforeend", "<i>bio</i>");  // le plus rapide des trois pour inserer du HTML

// CLASSES ET ATTRIBUTS
li.classList.add("actif"); li.classList.remove("cache");
li.classList.toggle("fait", coche);         // 2e argument = booleen : la forme propre
li.classList.replace("petit", "grand");
li.classList.contains("fait");
li.className = "a b c";                     // ecrase tout : a eviter
li.setAttribute("aria-label", "Article riz");
li.getAttribute("data-id");                 // "42" (toujours une chaine)
li.id = "article-42";
li.hidden = true;                           // equivalent de [hidden]

// DONNEES UTILISATEUR : dataset (kebab-case <-> camelCase)
// <li data-id="42" data-categorie="feculent">
li.dataset.id;                // "42"
li.dataset.categorie = "feculent";
li.dataset.etatCoche = "1";   // data-etat-coche

// STYLE : prefere une classe CSS a une propriete inline
li.style.backgroundColor = "#eef";          // camelCase en JS, kebab-case en CSS
li.style.setProperty("--hauteur", "40px");  // variable CSS
getComputedStyle(li).fontSize;              // "16px" (la valeur calculee, pas celle ecrite)
li.style.cssText = "color:red;padding:4px"; // remplace tout le style inline

// DIMENSIONS (attention : forcent un reflow -- a lire par lots)
li.offsetWidth; li.offsetHeight; li.getBoundingClientRect();
document.documentElement.scrollTop;
\end{codejs}

\subsection{Cr\'eer, ins\'erer, supprimer, d\'eplacer}

\begin{codejs}
// La voie royale : template + clone (rapide, sans parse, sans injection)
const tpl = document.getElementById("tpl-article");
const creerArticle = ({ nom, qte }) => {
  const frag = tpl.content.cloneNode(true);            // true = avec les descendants
  frag.querySelector(".nom").textContent = `${nom} x${qte}`;
  frag.querySelector(".suppr").dataset.id = nom;
  return frag;
};
document.getElementById("liste").append(creerArticle({ nom: "Riz", qte: 2 }));

// Creation manuelle
const el = document.createElement("li");
el.className = "article";
el.append("texte ", document.createElement("b"), " et noeud");   // append accepte textes et noeuds
el.prepend(puce);
el.replaceChildren(...listeFils);        // remplace TOUS les enfants en une fois : tres efficace
el.before(sibling); el.after(sibling); el.replaceWith(autre);
el.remove();                              // le fils se detache du parent ; l'objet reste utilisable

// Ou par HTML (rappel : seulement pour du contenu de confiance)
const conteneur = document.getElementById("liste");
conteneur.insertAdjacentHTML("beforeend", htmlSecurise);
conteneur.children.length;
\end{codejs}

\etape{\'Ecrire une fonction de nettoyage si tu dois ins\'erer du HTML}
\begin{codejs}
const echapper = (s) => String(s)
  .replaceAll("&", "&amp;").replaceAll("<", "&lt;").replaceAll(">", "&gt;")
  .replaceAll('"', "&quot;").replaceAll("'", "&#39;");
const rendu = (t) => `<strong>${echapper(t)}</strong>`;   // le texte est sur, la structure est la tienne
\end{codejs}
En production, ne r\'einvente pas ce filtre~: \link{DOMPurify}{https://github.com/cure53/DOMPurify}
(\cmd{npm i dompurify}) g\`ere les cas tordus \verb|<svg>|, les gestionnaires inline, les URL
\verb|javascript:|.

\section{Le mod\`ele de l'arbre et les noeuds}

\begin{codejs}
document            // Document
 .documentElement    // <html>
 .firstElementChild  // <head>
document.body; document.head; document.title;
document.activeElement;                       // l'element qui a le focus
document.getSelection().toString();           // texte selectionne
node.parentNode; node.parentElement;         // parent, en incluant les noeuds de texte
node.childNodes; node.children;               // idem : tous vs elements uniquement
node.nextSibling / node.nextElementSibling;
node.isConnected;                              // encore dans le document ?
node.contains(other);
node.innerText;                                // comme textContent mais respectant le CSS (lent !)
\end{codejs}

\begin{piege}
\verb|innerText| d\'eclenche un calcul de style complet sur l'\'el\'ement et ses descendants.
Sur une liste de mille lignes, c'est plusieurs secondes de g\^ach\'ees. Utilise
\verb|textContent|~; garde \verb|innerText| pour reproduire l'affichage visible (sauts de ligne,
\verb|display:none|).
\end{piege}

\section{Performance : ce que le navigateur paie}

Le pipeline est~: JavaScript \$\rightarrow\$ styles (recalcul) \$\rightarrow\$ layout (g\'eom\'etrie)
\$\rightarrow\$ paint (pixels) \$\rightarrow\$ composite (GPU). Tout changement de
\texttt{width/top/left/height} relance layout~; tout changement de
\texttt{transform/opacity} peut rester sur le compositeur~: c'est le chemin rapide.

\begin{codejs}
// ANTI-PATTERN : lecture puis ecriture = reflow a chaque tour (layout thrashing)
for (const el of elements) {
  el.style.height = el.offsetHeight + 10 + "px";      // lit la geometrie (reflow) puis ecrit
}

// CORRECT : lire tout, puis ecrire tout
const hauteurs = elements.map((el) => el.offsetHeight);
elements.forEach((el, i) => { el.style.height = `${hauteurs[i] + 10}px`; });

// ENCORE MIEUX : laisser le CSS faire le travail
el.classList.add("agrandi");
// Animer uniquement transform et opacity
el.animate(
  [{ opacity: 0, transform: "translateY(8px)" }, { opacity: 1, transform: "none" }],
  { duration: 200, easing: "cubic-bezier(.2,.8,.2,1)", fill: "both" }
);

// Grouper les mutations importantes dans une frame
requestAnimationFrame(() => { miseAJourLourde(); });

// Observer les changements de DOM sans polluer : MutationObserver
const obs = new MutationObserver((mutations) => {
  for (const m of mutations) {
    if (m.type === "childList") console.log("ajoutes :", [...m.addedNodes].length, "retires :", m.removedNodes.length);
    if (m.type === "attributes") console.log("attribut", m.attributeName, "sur", m.target);
  }
});
obs.observe(document.getElementById("liste"), { childList: true, subtree: true, attributes: true });
obs.disconnect();

// Le lazy-loading par IntersectionObserver : la fin du scroll listener couteux
const io = new IntersectionObserver((entrees) => {
  for (const e of entrees) if (e.isIntersecting) { chargerImage(e.target); io.unobserve(e.target); }
}, { rootMargin: "200px", threshold: 0.1 });
document.querySelectorAll("img[data-src]").forEach((i) => io.observe(i));
\end{codejs}

\begin{note}
Le navigateur regroupe lui-m\^eme les \'ecritures d'un m\^eme bloc de JavaScript~: les \'ecritures
\emph{pures} ne causent pas de reflow imm\'ediat. Ce qui co\^ute, c'est \textbf{lire une valeur de
g\'eom\'etrie} juste apr\`es avoir \'ecrit. Donc~: on lit, puis on \'ecrit. Jamais en alternance.
\end{note}

\section{Composants natifs que tout le monde r\'einvente mal}

\begin{codejs}
// Boite de dialogue : gestion du focus, fermeture par Echap, piege modal inclus
const dlg = document.querySelector("dialog");
dlg.showModal();
dlg.returnValue = "ok";
dlg.addEventListener("close", () => console.log(dlg.returnValue));

// Details/summary : accord\'eon accessible sans une ligne de JS
// <details name="faq" open><summary>Question</summary><p>Reponse</p></details>
// "name" cree un groupe exclusif (un seul ouvert) : natif depuis 2023.

// Popover : menus, tooltips, sans lib
// <button popovertarget="menu">Menu</button><div id="menu" popover>...</div>

// Color/Date/Range : le controleur natif est deja accessible et mobile-friendly
// <input type="date" min="2026-01-01"> <input type="color"> <input type="range">

// Scroll fluide et ancrage
document.querySelector("#tarif").scrollIntoView({ behavior: "smooth", block: "start" });
// CSS : html { scroll-behavior: smooth }  ou  scroll-margin-top: 80px sur la cible

// Historique et URL propre (le SPA tout entier tient la-dedans)
history.pushState({ vue: "article" }, "", "/article/42");
window.addEventListener("popstate", (e) => rendre(e.state?.vue ?? "accueil"));
\end{codejs}

\section{Accessibilit\'e : trente minutes qui changent la vie de tes utilisateurs}

\begin{enumerate}
  \item Un HTML s\'emantique (\verb|main|, \verb|nav|, \verb|button|, \verb|h1..h6|) avant
        n'importe quel ARIA. \verb|<div onclick>| n'est pas un bouton~: ni focus, ni Echap,
        ni annonce au lecteur d'\'ecran.
  \item Chaque champ a un \verb|<label for>|~; les ic\^ones seules ont un \verb|aria-label|.
  \item \verb|aria-live="polite"| sur une zone qui change (compteur, r\'esultat de recherche,
        message d'erreur)~; \verb|aria-busy| pendant un chargement.
  \item Ne retire jamais l'outline de focus sans la remplacer par un style visible.
  \item Contraste minimal 4{,}5:1 pour le texte courant, 3:1 pour le gros texte
        (\link{WebAIM~Contrast~Checker}{https://webaim.org/resources/contrastchecker/}).
  \item Navigation clavier compl\`ete~: \keys{Tab} partout, \keys{Echap} ferme,
        \keys{Enter}/\keys{Space} activent, fl\`eches dans les listes.
\end{enumerate}

\begin{codejs}
// Focus piege pour une modale custom (si tu ne utilises pas <dialog>)
function piegerFocus(conteneur) {
  const selecteur = 'a[href], button:not([disabled]), input, select, textarea, [tabindex]:not([tabindex="-1"])';
  const premier = conteneur.querySelector(selecteur);
  const focusables = () => [...conteneur.querySelectorAll(selecteur)].filter((el) => el.offsetParent !== null);
  conteneur.addEventListener("keydown", (e) => {
    if (e.key !== "Tab") return;
    const els = focusables();
    const idx = els.indexOf(document.activeElement);
    if (e.shiftKey && (idx === 0 || idx === -1)) { els.at(-1).focus(); e.preventDefault(); }
    else if (!e.shiftKey && idx === els.length - 1) { els[0].focus(); e.preventDefault(); }
  });
  premier?.focus();
  return () => dernierElement?.focus();       // retablir le focus a la fermeture
}
\end{codejs}

\section{Mini-projet complet --- liste de courses persistante}

Le fichier \fichier{app.js} branch\'e sur le HTML de la \cref{ch:js-dom}~; tout est comment\'e,
et tu dois le retaper int\'egralement.

\begin{codejs}
// app.js
const stock = {
  lire: () => { try { return JSON.parse(localStorage.getItem("courses") ?? "[]"); } catch { return []; } },
  ecrire: (v) => localStorage.setItem("courses", JSON.stringify(v)),
};

const $ = (s, r = document) => r.querySelector(s);
const $$ = (s, r = document) => [...r.querySelectorAll(s)];

let articles = stock.lire();
let filtre = "tous";

function rendre() {
  const ul = $("#liste");
  const visibles = articles.filter((a) =>
    filtre === "tous" ? true : filtre === "coches" ? a.fait : !a.fait
  );
  const frag = document.createDocumentFragment();
  const tpl = $("#tpl-article");

  for (const a of visibles) {
    const n = tpl.content.cloneNode(true);
    const li = n.querySelector("li");
    li.dataset.id = a.id;
    li.classList.toggle("fait", a.fait);
    n.querySelector(".nom").textContent = a.nom;
    const check = n.querySelector(".check");
    check.checked = a.fait;
    frag.append(li.parentNode ? n : n);
  }
  ul.replaceChildren(frag);
  $("#compteur").textContent =
    `${articles.filter((a) => !a.fait).length} article(s) a acheter sur ${articles.length}`;
}

function ajouter(nom, qte) {
  nom = nom.trim();
  if (!nom) return alert("Nom requis");
  articles = [{ id: crypto.randomUUID(), nom, qte: Number(qte) || 1, fait: false }, ...articles];
  stock.ecrire(articles);
  rendre();
}

// Deleguer sur la liste : un seul ecouteur, meme pour les elements crees plus tard
$("#liste").addEventListener("click", (e) => {
  const li = e.target.closest("li[data-id]");
  if (!li) return;
  const id = li.dataset.id;
  if (e.target.matches(".check")) {
    articles = articles.map((a) => (a.id === id ? { ...a, fait: !a.fait } : a));
  } else if (e.target.matches(".suppr")) {
    articles = articles.filter((a) => a.id !== id);
  } else return;
  stock.ecrire(articles);
  rendre();
});

$("#form-ajout").addEventListener("submit", (e) => {
  e.preventDefault();
  const f = new FormData(e.target);
  ajouter(f.get("nom"), f.get("qte"));
  e.target.reset();
  $("#nom").focus();
});

$$(".filtre").forEach((b) =>
  b.addEventListener("click", () => {
    filtre = b.dataset.filtre;
    $$(".filtre").forEach((x) => x.setAttribute("aria-pressed", String(x === b)));
    rendre();
  })
);

rendre();
\end{codejs}

\begin{exercice}
\begin{enumerate}
  \item Ajoute la r\'eordonnancement par glisser-d\'eposer (\verb|draggable="true"|,
        \verb|dragstart|, \verb|dragover|, \verb|drop|) et persiste l'ordre.
  \item Ajoute un champ de recherche avec \verb|debounce|~: filtrer sur le nom, en ignorant
        accents et casse, et surligner la partie correspondante avec \verb|<mark>|.
  \item Ajoute un mode sombre~: lire \verb|prefers-color-scheme| via
        \verb|matchMedia|, stocker le choix de l'utilisateur, basculer la classe
        \verb|dark| sur \verb|<html>|.
  \item Ajoute l'annulation~: un tableau \verb|historique| des 20 derniers \'etats, et les
        raccourcis \keys{Ctrl+Z} / \keys{Ctrl+Shift+Z}.
  \item Rends le composant \emph{r\'eutilisable}~: \verb|creerListe(elise, { stock })| qui
        accepte n'importe quel \'el\'ement conteneur et n'utilise aucun \verb|document.querySelector| global.
  \item Mesure avec l'onglet \clic{Performance} de Chrome~: d\'eclenche 500 ajouts, trouve la
        t\^ache la plus co\^uteuse, optimise, compare les chiffres.
\end{enumerate}
\end{exercice}

\begin{ressources}
\begin{itemize}
  \item MDN --- R\'ef\'erence DOM --- \url{https://developer.mozilla.org/fr/docs/Web/API/Document\\\_Object\\\_Model}
  \item web.dev Learn CSS / Learn Performance (Google, gratuit, excellent) --- \url{https://web.dev/learn/}
  \item Chrome DevTools, formation officielle --- \url{https://developer.chrome.com/docs/devtools/learn}
  \item Check-list d'accessibilit\'e --- \url{https://www.a11yproject.com/checklist/}
  \item \link{Can~I~Use}{https://caniuse.com/} pour chaque API ou s\'electeur CSS.
\end{itemize}
\end{ressources}

\index{DOM}\index{navigateur}\index{accessibilite}
% ============================================================================
\chapter{\'Ev\'enements, formulaires et interaction utilisateur}
\label{ch:js-events}
% ============================================================================

\begin{objectif}
\begin{itemize}
  \item Ma\^itriser les trois phases de propagation et l'art de d\'el\'eguer un \'ecouteur.
  \item Exploiter les options d'\texttt{addEventListener} (\texttt{signal}, \texttt{once},
        \texttt{capture}, \texttt{passive}).
  \item Valider un formulaire c\^ot\'e client \emph{et} comprendre pourquoi ce n'est qu'un confort.
  \item Envoyer des fichiers, g\'erer upload avec progression, et les \'ev\'enements clavier/souris.
  \item Construire les composants UI courants (menu, modale, autocomplete, toast).
\end{itemize}
\end{objectif}

\section{Le cycle de vie d'un \'ev\'enement}

Un clic sur un \verb|<span>| dans un \verb|<button>| dans une \verb|<form>| parcourt l'arbre
deux fois~: \textbf{capture} (du plus grand vers la cible), \textbf{target} (la cible),
\textbf{bulle} (de la cible vers le plus grand). La plupart du temps on \'ecoute en phase de
bulle, c'est pour \c{c}a que la d\'el\'egation fonctionne.

\begin{codejs}
form.addEventListener(
  "submit",
  (e) => {
    e.preventDefault();                  // empeche le comportement par defaut (rechargement)
    console.log(e.type, e.target, e.currentTarget);   // currentTarget = l'element quiecoute
    console.log(e.clientX, e.clientY, e.timeStamp);
  },
  { signal: abort.signal }               // la meilleure facon de nettoyer
);

const options = {
  capture: true,        // ecouter en phase descendante
  once: true,           // se desabonne apres le premier declenchement
  passive: true,        // promet de ne pas appeler preventDefault -> scroll plus fluide
  signal: new AbortController().signal,   // desabonnement en masse
};
window.addEventListener("wheel", onWheel, { passive: true });

// Stopper la propagation vs empecher l'action par defaut : DEUX choses differentes
e.stopPropagation();        // les parents ne seront pas notifies (mais l'action a lieu)
e.stopImmediatePropagation(); // ni les autres ecouteurs du MEME element
e.preventDefault();         // l'action par defaut est annulee (la propagation continue)

// Evenements sans bulle : focus, blur, mouseenter, load, scroll (sur element)
// => pour les deleguer, ecouter en capture :
conteneur.addEventListener("focusin", onFocus);      // focusin/focusout, eux, bullent
conteneur.addEventListener("focus", onFocus, true);  // capture : marche aussi
\end{codejs}

\section{La d\'el\'egation~: un \'ecouteur pour mille lignes}

C'est LE motif qui fait la diff\'erence entre un script qui vieillit bien et un script
qu'on r\'e\'ecrit chaque fois qu'on ajoute un \'el\'ement.

\begin{codejs}
// ANTI-PATTERN : attacher un ecouteur par ligne, a reattacher a chaque rendu
document.querySelectorAll(".suppr").forEach((b) => b.addEventListener("click", supprimer));

// DELEGATION : un seul ecouteur sur un conteneur stable
document.getElementById("liste").addEventListener("click", (e) => {
  const boutonSuppr = e.target.closest("[data-action=supprimer]");
  if (boutonSuppr) return supprimer(boutonSuppr.closest("li").dataset.id);

  const coche = e.target.closest("[data-action=cocher]");
  if (coche) return cocher(coche.closest("li").dataset.id);

  const lien = e.target.closest("a");
  if (lien) { e.preventDefault(); naviguer(lien.href); }
});
\end{codejs}

Un utilitaire de bus d'\'ev\'enements, \`a copier dans ton projet (il remplace avantageusement
les r\'ef\'erences crois\'ees entre modules)~:

\begin{codejs}
export function creerBus() {
  const ecouteurs = new Map();
  return {
    on(type, fn, options = {}) {
      if (!ecouteurs.has(type)) ecouteurs.set(type, new Set());
      ecouteurs.get(type).add(fn);
      return () => this.off(type, fn);
    },
    once(type, fn) { const off = this.on(type, (p) => { off(); fn(p); }); return off; },
    off(type, fn) { ecouteurs.get(type)?.delete(fn); },
    emit(type, detail) {
      for (const fn of ecouteurs.get(type) ?? []) {
        try { fn(detail); } catch (err) { console.error(`ecouteur "${type}"`, err); }
      }
      if (typeof CustomEvent === "function") {
        document.dispatchEvent(new CustomEvent(type, { detail, bubbles: true }));
      }
    },
    clear() { ecouteurs.clear(); },
  };
}
// Evenements personnalises, la norme DOM :
document.body.dispatchEvent(new CustomEvent("panier:maj", { detail: { total: 42 }, bubbles: true, composed: true }));
document.body.addEventListener("panier:maj", (e) => afficher(e.detail.total));
\end{codejs}

\section{Nettoyer : le pi\`ege n\${}^{\text{o}}\$1 des SPA}

Un \'ecouteur ajout\'e et jamais retir\'e = fuite de m\'emoire + rappels fant\^omes sur des \'el\'ements
d\'etruits. Trois strat\'egies, de la plus simple \`a la plus propre.

\begin{codejs}
// 1. once : pour un evenement unique
btn.addEventListener("click", demarrer, { once: true });

// 2. AbortController : le standard moderne, et le plus pratique
const ctrl = new AbortController();
window.addEventListener("resize", surResize, { signal: ctrl.signal });
document.addEventListener("keydown", surTouche, { signal: ctrl.signal });
form.addEventListener("submit", surEnvoi, { signal: ctrl.signal, passive: false });
function detruire() { ctrl.abort(); }       // retire TOUT, en une ligne

// 3. Un mini-cycle de vie pour tes propres composants
function creerComposant(racine, { montar, demonter }) {
  const ctrl = new AbortController();
  montar(racine, ctrl.signal);
  return {
    destroy() { ctrl.abort(); demonter?.(racine); racine.replaceChildren(); },
  };
}
\end{codejs}

\section{Formulaires---lecture, \'ecriture, validation}

\begin{codejs}
const form = document.querySelector("form");

// Lecture moderne : FormData
const data = new FormData(form);
const obj = Object.fromEntries(data.entries());           // { nom: "...", qte: "2" }
const roles = data.getAll("role");                         // checkbox multiple -> tableau
const fichier = data.get("avatar");                        // File | null
data.set("nom", "Nour"); data.append("tag", "admin");
// Renvoyer un objet typ\'e, propre a envoyer en JSON :
const lire = (form) => {
  const o = Object.fromEntries(new FormData(form).entries());
  for (const [k, v] of Object.entries(o)) {
    if (v === "on") o[k] = true;
    else if (v === "off") o[k] = false;
    else if (v !== "" && !Number.isNaN(Number(v)) && form.elements[k]?.type === "number") o[k] = Number(v);
  }
  return o;
};

// Validation native (le navigateur la fait deja, et gratuitement)
const champ = form.elements.email;
champ.checkValidity();          // booleen
champ.validity.tooLong; champ.validationMessage;
form.reportValidity();          // affiche les bulles, met le focus sur le 1er invalide
champ.setCustomValidity("Choisis un email @exemple.com");   // remplace le message
champ.setCustomValidity("");     // ne pas oublier de le reinitialiser !

// Les attributs HTML qui font 80 % du travail :
// required, type, pattern, minlength/maxlength, min/max/step, inputmode, autocomplete,
// autocapitalize, multiple, accept, list (datalist)
\end{codejs}

\subsection{Validation sur mesure, r\'eutilisable}

\begin{codejs}
const regles = {
  requis: (v) => (v === null || v === undefined || String(v).trim() === "") && "Champ obligatoire",
  email: (v) => v && !/^[^@\s]+@[^@\s]+\.[a-z]{2,}$/i.test(v) && "Adresse email invalide",
  min: (n) => (v) => [...String(v ?? "")].length < n && `Au moins ${n} caracteres`,
  nombre: (min, max) => (v) => {
    const x = Number(v);
    if (v === "" || !Number.isFinite(x)) return "Nombre attendu";
    if (min !== undefined && x < min) return `Minimum ${min}`;
    if (max !== undefined && x > max) return `Maximum ${max}`;
  },
  mdpFort: (v) => {
    if (!v) return;
    const manques = [];
    if (v.length < 12) manques.push("12 caracteres minimum");
    if (!/[A-Z]/.test(v)) manques.push("une majuscule");
    if (!/[a-z]/.test(v)) manques.push("une minuscule");
    if (!/\d/.test(v)) manques.push("un chiffre");
    if (!/[^\w\s]/.test(v)) manques.push("un caractere special");
    return manques.length && `Mot de passe faible : ${manques.join(", ")}`;
  },
};

export function valider(form, schema) {
  const erreurs = {};
  for (const [nom, liste] of Object.entries(schema)) {
    const champ = form.elements[nom];
    const valeur = champ?.type === "checkbox" ? champ.checked : champ?.value;
    for (const regle of liste) {
      const message = typeof regle === "function" ? regle(valeur, form) : regles[regle]?.(valeur, form);
      if (message) { erreurs[nom] = message; break; }
    }
  }
  return { ok: Object.keys(erreurs).length === 0, erreurs };
}

// Branchement : invalider a la sortie, revalider a la frappe une fois l'erreur affichee
const champ = form.elements.email;
const afficherErreur = (msg) => {
  const zone = champ.closest(".champ");
  zone.classList.toggle("invalide", Boolean(msg));
  zone.querySelector(".erreur").textContent = msg ?? "";
  champ.setAttribute("aria-invalid", String(Boolean(msg)));
  champ.setAttribute("aria-describedby", msg ? "err-email" : "");
};
champ.addEventListener("blur", () => afficherErreur(regles.email(champ.value)));
champ.addEventListener("input", () => champ.closest(".champ").classList.contains("invalide")
  && afficherErreur(regles.email(champ.value)));

form.addEventListener("submit", async (e) => {
  e.preventDefault();
  const { ok, erreurs } = valider(form, {
    email: ["requis", "email"],
    age: [regles.nombre(18, 120)],
    mdp: [regles.mdpFort],
  });
  if (!ok) {
    const premier = form.elements[Object.keys(erreurs)[0]];
    afficherErreur(erreurs.email);
    premier?.focus();
    return;
  }
  await soumettre(Object.fromEntries(new FormData(form)));
});
\end{codejs}

\begin{piege}
Un formulaire doit \^etre soumis par un vrai \verb|<button type="submit">| et r\'epondre \`a
\verb|Enter|~: c'est ce que la touche Entr\'ee, les lecteurs d'\'ecran et les gestionnaires de mots
de passe attendent. Un simple \verb|onclick| sur une \verb|<div>| casse tout cela~;
la validation native ne se d\'eclenche m\^eme pas.
\end{piege}

\section{Fichiers : s\'election, lecture, envoi, progression}

\begin{codejs}
const input = document.querySelector('input[type=file]');

input.addEventListener("change", async () => {
  for (const file of input.files) {
    // Garde-fous cote client (le serveur refait TOUJOURS les siens)
    if (file.size > 5 * 1024 ** 2) { alert(`${file.name} depasse 5 Mo`); continue; }
    if (!/^image\/(png|jpe?g|webp|avif)$/.test(file.type)) continue;

    const url = URL.createObjectURL(file);            // lecture immediate pour preview
    document.querySelector("#apercu").src = url;

    const texte = await file.text();                  // File -> string (Promise)
    const buffer = await file.arrayBuffer();          // binaire (hash, analyse)
    const hash = [...new Uint8Array(await crypto.subtle.digest("SHA-256", buffer))]
      .map((b) => b.toString(16).padStart(2, "0")).join("");

    const dataUrl = await new Promise((res, rej) => {
      const r = new FileReader();
      r.onload = () => res(r.result);
      r.onerror = () => rej(r.error);
      r.readAsDataURL(file);
    });
    URL.revokeObjectURL(url);                         // liberer la memoire !
    console.log(file.name, file.size, file.type, file.lastModified, hash.slice(0, 12));
  }
});

// Envoi : multipart via FormData, le navigateur met le Content-Type tout seul
async function envoyer(form) {
  const fd = new FormData(form);
  const ctrl = new AbortController();
  const xhr = new XMLHttpRequest();           // seul moyen simple d'avoir la progression
  xhr.open("POST", "/upload");
  xhr.upload.onprogress = (e) => {
    if (e.lengthComputable) barre.style.width = `${(e.total * 100 / e.loaded).toFixed(0)}`;
  };
  xhr.send(fd);
}

// Upload en chunk (gros fichiers) : decoupage, envoi parallele, reprise
async function uploaderParMorceaux(file, taille = 5 * 1024 ** 2) {
  const sessions = await fetch("/upload/init", {
    method: "POST", headers: { "content-type": "application/json" },
    body: JSON.stringify({ nom: file.name, taille: file.size, morceaux: Math.ceil(file.size / taille) }),
  }).then((r) => r.json());
  const envois = [];
  for (let i = 0; i * taille < file.size; i++) {
    const morceau = file.slice(i * taille, (i + 1) * taille);
    envois.push(fetch(`/upload/${sessions.id}/${i}`, { method: "PUT", body: morceau }));
    if (envois.length >= 3) await Promise.race(envois);       // limiter la concurrence
  }
  await Promise.all(envois);
  return fetch(`/upload/${sessions.id}/complete`, { method: "POST" });
}

// Telecharger un fichier genere cote serveur, sans quitter la page
const reponse = await fetch("/export.csv");
const blob = await reponse.blob();
const a = Object.assign(document.createElement("a"), {
  href: URL.createObjectURL(blob), download: "export.csv",
});
a.click(); URL.revokeObjectURL(a.href);
\end{codejs}

\section{Clavier, souris, glisser, tactile}

\begin{codejs}
// Clavier : toujours tester la touche + les modificateurs, et menager les accessifs
window.addEventListener("keydown", (e) => {
  if (e.target.matches("input, textarea, [contenteditable]")) {
    if (!(e.ctrlKey || e.metaKey)) return;      // ne pas capturer les touches de saisie
  }
  const mod = e.metaKey || e.ctrlKey ? "Mod+" : "";
  const combo = `${mod}${e.shiftKey ? "Shift+" : ""}${e.key.toLowerCase()}`;
  const raccourcis = {
    "Mod+k": () => ouvrirPalette(),
    "Mod+Enter": () => soumettre(),
    "Mod+Shift+p": () => basculerMode(),
    Escape: () => fermerTout(),
    "?": () => ouvrirAide(),
  };
  const fn = raccourcis[combo];
  if (fn) { e.preventDefault(); fn(); }
});

// Souris : coalescer le mouvement (60+ evenements par frame sinon)
let last = { x: 0, y: 0 };
el.addEventListener("pointermove", (e) => {
  const [evt] = e.getCoalescedEvents?.() ?? [e];     // regroupe les evements perdus
  last = { x: evt.clientX, y: evt.clientY };
}, { passive: true });

// Pointer events : UN seul modele pour souris + stylet + doigt (oublie touch*)
const points = new Map();
el.addEventListener("pointerdown", (e) => {
  el.setPointerCapture(e.pointerId);
  points.set(e.pointerId, { x: e.clientX, y: e.clientY });
});
el.addEventListener("pointermove", (e) => points.get(e.pointerId) && pince(e, points));
el.addEventListener("pointerup", (e) => points.delete(e.pointerId));

// Double-clic, clic droit, molette, drag natif
el.addEventListener("dblclick", zoom);
el.addEventListener("contextmenu", (e) => { e.preventDefault(); menu.positionner(e.clientX, e.clientY); });
el.addEventListener("wheel", (e) => { if (e.ctrlKey) { e.preventDefault(); zoomPar(e.deltaY); } }, { passive: false });
el.draggable = true;
el.addEventListener("dragstart", (e) => e.dataTransfer.setData("text/plain", el.dataset.id));
zone.addEventListener("dragover", (e) => { e.preventDefault(); zone.classList.add("sur"); });
zone.addEventListener("drop", (e) => {
  e.preventDefault();
  deplacer(e.dataTransfer.getData("text/plain"), indexApres(e));
});
// Drag de fichiers depuis le bureau
window.addEventListener("dragover", (e) => e.preventDefault());
window.addEventListener("drop", (e) => { e.preventDefault(); traiter(e.dataTransfer.files); });
\end{codejs}

\section{Quatre composants UI, sans biblioth\`eque}

\begin{codejs}
// 1. Menu deroulant accessible (fleches, Echap, clic exterieur, focus)
function menu(bouton, liste) {
  bouton.addEventListener("click", () => {
    const ouvert = liste.hasAttribute("hidden");
    liste.toggleAttribute("hidden", !ouvert);
    bouton.setAttribute("aria-expanded", String(ouvert));
    if (ouvert) liste.querySelector('[role="menuitem"]')?.focus();
  });
  liste.addEventListener("keydown", (e) => {
    const items = [...liste.querySelectorAll('[role="menuitem"]:not([disabled])')];
    const i = items.indexOf(document.activeElement);
    const bouger = (n) => { e.preventDefault(); items[(i + n + items.length) % items.length].focus(); };
    if (e.key === "ArrowDown") bouger(1);
    else if (e.key === "ArrowUp") bouger(-1);
    else if (e.key === "Home") bouger(-i);
    else if (e.key === "End") bouger(items.length - 1 - i);
    else if (e.key === "Escape") { bouton.click(); bouton.focus(); }
  });
  document.addEventListener("click", (e) => {
    if (!e.target.closest("[data-menu]") && !liste.contains(e.target)) {
      liste.setAttribute("hidden", ""); bouton.setAttribute("aria-expanded", "false");
    }
  });
}

// 2. Toast (notification auto-detruite)
const toasts = (conteneur => ({
  show(message, { type = "info", duree = 4000 } = {}) {
    const t = document.createElement("div");
    t.className = `toast toast--${type}`; t.setAttribute("role", "status"); t.textContent = message;
    conteneur.append(t);
    t.addEventListener("click", () => t.remove());
    setTimeout(() => { t.classList.add("out"); t.addEventListener("transitionend", () => t.remove(), { once: true }); }, duree);
  },
}))(document.getElementById("toasts"));
toasts.show("Enregistre", { type: "succes" });

// 3. Autocompletion (combobox) avec debounce + annulation + navigation clavier
function autocomplete(input, liste, chercher) {
  let ctrl = null, index = -1;
  const fermer = () => { liste.setAttribute("hidden", ""); index = -1; };
  input.addEventListener("input", debounce(async () => {
    ctrl?.abort(); ctrl = new AbortController();
    const q = input.value.trim();
    if (q.length < 2) return fermer();
    let items;
    try { items = await chercher(q, ctrl.signal); } catch (e) { if (e.name !== "AbortError") throw e; return; }
    if (!items.length) return fermer();
    liste.replaceChildren(...items.map((it, i) => {
      const li = document.createElement("li");
      li.setAttribute("role", "option"); li.textContent = it.label; li.dataset.value = it.value;
      li.id = `ac-${i}`;
      return li;
    }));
    liste.removeAttribute("hidden");
    input.setAttribute("aria-expanded", "true");
  }, 220));
  input.addEventListener("keydown", (e) => {
    const opts = [...liste.querySelectorAll("[role=option]")];
    if (e.key === "ArrowDown" || e.key === "ArrowUp") {
      e.preventDefault();
      index = (index + (e.key === "ArrowDown" ? 1 : -1) + opts.length) % opts.length;
      opts.forEach((o, i) => o.setAttribute("aria-selected", String(i === index)));
      input.setAttribute("aria-activedescendant", opts[index]?.id ?? "");
    } else if (e.key === "Enter" && index >= 0) {
      e.preventDefault(); input.value = opts[index].dataset.value; fermer();
    } else if (e.key === "Escape") fermer();
  });
  liste.addEventListener("click", (e) => {
    const o = e.target.closest("[role=option]"); if (!o) return;
    input.value = o.dataset.value; fermer(); input.focus();
  });
}

// 4. Scroll infini (la seule bonne impl, sans listener de scroll)
function infini(sentinelle, charger) {
  let enCours = false;
  const io = new IntersectionObserver(async ([e]) => {
    if (!e.isIntersecting || enCours) return;
    enCours = true;
    const fini = await charger();
    enCours = false;
    if (fini) io.disconnect();
  }, { rootMargin: "400px" });
  io.observe(sentinelle);
  return () => io.disconnect();
}
\end{codejs}

\begin{exercice}
\begin{enumerate}
  \item Reprends la liste de courses et remplace tous les \'ecouteurs par un seul syst\`eme de
        d\'el\'egation + un \verb|AbortController| global.
  \item \'Ecris un formulaire d'inscription complet~: nom, email, mot de passe (jauge de force
        en temps r\'eel), confirmation, conditions (case \`a cocher). Validation \`a la sortie du champ,
        messages sous le champ, \verb|aria-invalid|, focus sur la premi\`ere erreur.
  \item Construis un s\'electeur de date \og ~calendrier~ \fg{}~: grille 7\$\times\$6, navigation mois,
        \keys{\$\leftarrow\$} \keys{\$\rightarrow\$}, plage interdite, valeur au format \verb|YYYY-MM-DD| dans un champ cach\'e.
  \item Ajoute un glisser-d\'eposer de fichiers sur toute la page, avec aper\c{c}u des images,
        rejet des types interdits et affichage de la taille.
  \item \'Ecris une barre de progression de t\'el\'echargement r\'ealiste~: \verb|fetch| +
        \verb|ReadableStream| pour compter les octets re\c{c}us (sans \verb|XMLHttpRequest|).
  \item Fais un raccourci clavier global \og ~palette de commandes~ \fg{} (Ctrl+K) avec recherche
        floue, haut/bas, Entr\'ee, et focus pi\'eg\'e.
\end{enumerate}
\end{exercice}

\begin{ressources}
\begin{itemize}
  \item MDN --- \'Ev\'enements --- \url{https://developer.mozilla.org/fr/docs/Web/API/EventTarget/addEventListener}
  \item MDN --- Constraint Validation API --- \url{https://developer.mozilla.org/fr/docs/Web/API/HTMLConstraintValidation}
  \item \emph{Form Validation Design Patterns} (GOV.UK) --- \url{https://design-system.service.gov.uk/patterns/validation/}
  \item \link{a11y-collective / ARIA Authoring Practices}{https://www.w3.org/WAI/ARIA/apg/patterns/} (mod\`eles de composants accessibles)
  \item \link{File Picking, les 10 commandes}{https://web.dev/file-system-access/}
\end{itemize}
\end{ressources}

\index{evenements}\index{formulaire}\index{validation}\index{accessibilite}
% ============================================================================
\chapter{Communicuer avec le serveur : fetch, SSE, WebSockets, workers}
\label{ch:js-reseau}
% ============================================================================

\begin{objectif}
\begin{itemize}
  \item Ma\^itriser \texttt{fetch} dans tous ses usages~: JSON, en-t\^etes, erreurs, annulation,
        progression, t\'el\'echargement, upload.
  \item Structurer une couche API testable et agr\'eable (client HTTP maison, intercepteurs, cache).
  \item Consommer des flux temps r\'eel~: Server-Sent Events, WebSockets, reconnexion.
  \item D\'eplacer le travail lourd hors du fil principal~: Web Workers, OffscreenCanvas, WASM.
  \item Comprendre CORS et le r\'egler c\^ot\'e client \emph{et} c\^ot\'e serveur.
\end{itemize}
\end{objectif}

\section{Le strict minimum bien fait}

\begin{codejs}
async function getJSON(url, options = {}) {
  const ctrl = new AbortController();
  const minuterie = setTimeout(() => ctrl.abort(new Error("Timeout 10s")), 10_000);
  try {
    const res = await fetch(url, {
      signal: ctrl.signal,
      credentials: "same-origin",        // envoyer les cookies de session du meme site
      headers: { accept: "application/json", ...options.headers },
      ...options,
    });
    // fetch ne rejette PAS sur 404/500 : c'est a toi de le faire
    if (!res.ok) {
      const corps = await res.text().catch(() => "");
      throw new Error(`HTTP ${res.status} ${res.statusText} sur ${url}${corps ? ` -- ${corps.slice(0, 200)}` : ""}`);
    }
    if (res.status === 204) return null;
    const type = res.headers.get("content-type") ?? "";
    return type.includes("json") ? res.json() : res.text();
  } finally {
    clearTimeout(minuterie);
  }
}

const creer = (donnees) => getJSON("/api/articles", {
  method: "POST",
  headers: { "content-type": "application/json" },
  body: JSON.stringify(donnees),
});
const supprimer = (id) => getJSON(`/api/articles/${id}`, { method: "DELETE" });
\end{codejs}

\begin{tabz}{lX}
  \headrow \thd{Besoin} \& \thd{Solution} \\
  \midrule
  Envoi d'un formulaire simple \& \verb|await fetch(url, { method: "POST", body: new FormData(form) })| \\
  Upload d'un fichier \& \verb|body: formData| avec \verb|File| dedans (jamais \verb|Content-Type| manuel) \\
  Progression d'envoi \& \verb|XMLHttpRequest.upload.onprogress| ou \verb|fetch| + stream request (\verb|duplex: "half"|) \\
  Progression de r\'eception \& \verb|res.body| + compteur d'octets (stream) \\
  T\'el\'echargement binaire \& \verb|const blob = await res.blob()| puis \verb|URL.createObjectURL| \\
  Requ\^ete annulable \& \verb|AbortController| (aussi pour \'eviter les r\'eponses hors ordre) \\
  Cache HTTP fin \& \verb|cache: "no-store"|, \verb|"force-cache"|, \verb|"reload"| \\
  Authentification par jeton \& en-t\^ete \verb|authorization: Bearer ...| \\
  Cookies de session \& \verb|credentials: "include"| + \verb|SameSite=Lax| c\^ot\'e serveur \\
  Streaming de texte (IA) \& \verb|res.body.pipeThrough(new TextDecoderStream())| + \verb|for await| \\
\end{tabz}

\subsection{Un mini-client HTTP r\'eutilisable}

\begin{codejs}
// api.js --- le fichier que j'ecris en premier dans tout projet
const BASE = import.meta.env?.VITE_API_URL ?? "/api";

class ErreurAPI extends Error {
  constructor(message, { status, code, details } = {}) {
    super(message);
    this.name = "ErreurAPI"; this.status = status; this.code = code; this.details = details ?? null;
  }
}

function creerClient({ baseURL = BASE, getJeton, on401 } = {}) {
  const controles = new AbortController();
  async function requete(methode, chemin, { corps, signaux, enTetes, raw = false, idempotent = methode === "GET" } = {}) {
    let tentative = 0;
    while (true) {
      tentative++;
      const url = new URL(chemin.replace(/^\//, ""), baseURL.endsWith("/") ? baseURL : `${baseURL}/`).href;
      const jeton = await getJeton?.();
      let res;
      try {
        res = await fetch(url, {
          method: methode,
          signal: AbortSignal.any([...(signaux ?? []), controles.signal]),
          headers: {
            accept: "application/json",
            ...(corps !== undefined && !(corps instanceof FormData) ? { "content-type": "application/json" } : {}),
            ...(jeton ? { authorization: `Bearer ${jeton}` } : {}),
            ...enTetes,
          },
          ...(corps !== undefined ? { body: corps instanceof FormData || corps instanceof Blob ? corps : JSON.stringify(corps) } : {}),
        });
      } catch (err) {
        if (err.name === "AbortError") throw new DOMException("Requete annulee", "AbortError");
        if (idempotent && tentative < 3) { await new Promise((r) => setTimeout(r, 200 * 2 ** tentative)); continue; }
        throw new ErreurAPI("Serveur injoignable", { code: "reseau", details: err.message });
      }
      if (res.status === 401) { await on401?.(); throw new ErreurAPI("Session expiree", { status: 401 }); }
      if (!res.ok) {
        const payload = await res.json().catch(() => ({}));
        throw new ErreurAPI(payload.message ?? `HTTP ${res.status}`, { status: res.status, code: payload.code, details: payload.errors });
      }
      if (raw) return res;
      if (res.status === 204) return null;
      return (res.headers.get("content-type") ?? "").includes("json") ? res.json() : res.text();
    }
  }
  return {
    get: (c, o) => requete("GET", c, o),
    post: (c, corps, o) => requete("POST", c, { ...o, corps }),
    patch: (c, corps, o) => requete("PATCH", c, { ...o, corps }),
    put: (c, corps, o) => requete("PUT", c, { ...o, corps }),
    delete: (c, o) => requete("DELETE", c, o),
    brut: (methode, c, o) => requete(methode, c, { ...o, raw: true }),
    annulerTout: () => controles.abort(),
  };
}

export const api = creerClient({
  getJeton: () => sessionStorage.getItem("token"),
  on401: async () => { sessionStorage.clear(); location.assign("/connexion"); },
});
const articles = await api.get("/articles?page=1&tri=recent");
\end{codejs}

\begin{note}
Ne mets jamais un secret (cl\'e d'API priv\'ee, mot de passe base de donn\'ees) dans un bundle
JavaScript~: le navigateur le montre \`a quiconque ouvre l'onglet Sources. Tout ce qui doit
rester secret passe par ton propre serveur.
\end{note}

\section{CORS~: l'erreur que tout le monde rencontre}

Le navigateur refuse qu'une page de \verb|https://a.dev| lise une r\'eponse de
\verb|https://b.dev| sans autorisation explicite du serveur. Ce n'est \emph{pas} une s\'ecurit\'e
pour le serveur (il re\c{c}oit bien la requ\^ete)~: c'est une protection du \emph{navigateur}
pour l'utilisateur.

\begin{codejs}
// Cote serveur (ici en Express) --- les en-tetes qui debloquent la situation
app.use((req, res, next) => {
  const autorises = new Set(["http://localhost:5173", "https://nour.dev"]);
  const origine = req.headers.origin;
  if (autorises.has(origine)) {
    res.setHeader("access-control-allow-origin", origine);
    res.setHeader("vary", "origin");
    res.setHeader("access-control-allow-credentials", "true");
    res.setHeader("access-control-allow-methods", "GET,POST,PUT,PATCH,DELETE,OPTIONS");
    res.setHeader("access-control-allow-headers", "content-type,authorization,x-csrf-token");
    res.setHeader("access-control-max-age", "600");
  }
  if (req.method === "OPTIONS") return res.sendStatus(204);   // pre-vol
  next();
});
\end{codejs}

\etape{Lire le vrai message}
Dans la console~: \verb|Access to fetch at '...' has been blocked by CORS policy: No
'Access-Control-Allow-Origin' header...|. La cause est dans la \emph{deuxi\`eme} ligne~:
soit l'en-t\^ete manque, soit \verb|credentials: "include"| est utilis\'e avec
\verb|allow-origin: *| (interdit), soit un en-t\^ete personnalis\'e n'est pas list\'e dans
\verb|allow-headers|.

\etape{Pendant le d\'eveloppement, la vraie solution}
Ne touche pas au serveur~: fais passer les appels par le serveur de d\'eveloppement
(m\^eme origine, donc pas de CORS).

\begin{codejson}
{
  "server": {
    "proxy": {
      "/api": { "target": "http://localhost:8000", "changeOrigin": true }
    }
  },
  "scripts": { "dev": "vite", "dev:api": "php -S localhost:8000 -t public" }
}
\end{codejson}

\begin{codeplain}
// vite.config.js --- la meme idee avec Vite
import { defineConfig } from "vite";
export default defineConfig({
  server: {
    proxy: {
      "/api": { target: "http://localhost:8000", changeOrigin: true },
      "/ws": { target: "ws://localhost:8000", ws: true },
    },
  },
});
\end{codeplain}

\section{Server-Sent Events~: le temps r\'eel le plus simple du monde}

Une connexion HTTP classique, un flux texte infini, reconnexion automatique par le navigateur.
Parfait pour des notifications~; insuffisant pour du bidirectionnel.

\begin{codejs}
const sse = new EventSource("/api/stream");
sse.onmessage = (e) => rendre(e.data);                 // evenement sans nom
sse.addEventListener("notification", (e) => {
  const data = JSON.parse(e.data);
  notifier(data);
});
sse.onerror = () => { /* le navigateur se reconnecte seul ; a toi de gerer l'etat "hors ligne" */ };
// Cote serveur (Node/Express) :
// res.writeHead(200, { "content-type": "text/event-stream", "cache-control": "no-cache", connection: "keep-alive" });
// res.write(`event: notification\ndata: ${JSON.stringify(objet)}\n\n`);
// res.write(": ping\n\n");     // commentaire de maintien en vie toutes les 20-30 s
\end{codejs}

\section{WebSockets~: bidirectionnel}

\begin{codejs}
class Socket {
  #ws; #file = []; #retry = 0; #ctrl;
  constructor(url, { surMessage, surOuverture } = {}) {
    Object.assign(this, { url, surMessage, surOuverture });
  }
  ouvrir() {
    this.#ws = new WebSocket(this.url);
    this.#ws.onopen = () => {
      this.#retry = 0;
      this.#file.splice(0).forEach((m) => this.#ws.send(m));
      this.surOuverture?.();
    };
    this.#ws.onmessage = (e) => {
      const brut = typeof e.data === "string" ? e.data : lireBlob(e.data);
      try { this.surMessage?.(JSON.parse(brut)); }
      catch { this.surMessage?.(brut); }
    };
    this.#ws.onclose = (e) => {
      if (e.code === 1000) return;                       // fermeture propre
      this.#retry++;
      const attente = Math.min(30_000, 500 * 2 ** this.#retry);
      console.warn(`WebSocket ferme (${e.code}), reconnexion dans ${attente} ms`);
      this.#t = setTimeout(() => this.ouvrir(), attente + Math.random() * 250);
    };
    this.#ws.onerror = () => this.#ws.close();
  }
  envoyer(objet) {
    const msg = typeof objet === "string" ? objet : JSON.stringify(objet);
    if (this.#ws?.readyState === WebSocket.OPEN) this.#ws.send(msg);
    else this.#file.push(msg);                            // mettre en attente, pas jeter
  }
  fermer() { clearTimeout(this.#t); this.#ws?.close(1000); }
  get enLigne() { return this.#ws?.readyState === WebSocket.OPEN; }
}
const lireBlob = async (blob) => blob.text();   // ou arrayBuffer() pour du binaire
const ws = new Socket("wss://nour.dev/ws", { surMessage: (m) => console.log("recu", m) });
ws.envoyer({ type: "abonner", salon: "general" });
navigator.addEventListener("online", () => !ws.enLigne && ws.ouvrir());
\end{codejs}

\begin{note}
WebSocket ne propose ni authentification par cookie de la m\^eme fa\c{c}on ni reconnexion~: tu les
codes. C'est pourquoi, pour du simple temps r\'eel, SSE (serveur \$\rightarrow\$ client) ou une
solution h\'eberg\'ee (\link{Pusher}{https://pusher.com/}, \link{Ably}{https://www.ably.io/},
\link{Supabase Realtime}{https://supabase.com/}) est souvent le choix raisonnable.
C\^ot\'e serveur Node, la biblioth\`eque de r\'ef\'erence est
\link{ws}{https://github.com/websockets/ws} ou \link{Socket.IO}{https://socket.io/}.
\end{note}

\section{Web Workers : le CPU en parall\`ele}
\label{sec:workers}

Un worker est un fil JavaScript s\'epar\'e~: pas d'acc\`es au DOM, mais un CPU entier. On \'echange
des messages (structur\'es, clonables) et des \emph{transferables} ( ArrayBuffer, sans copie).

\begin{codeplain}
// poids-lourd.js (le fichier du worker)
self.onmessage = async (e) => {
  const { type, payload } = e.data;
  if (type === "ligner") {
    const t0 = performance.now();
    const resultat = calcul(payload);
    self.postMessage({ type: "resultat", payload: resultat, ms: performance.now() - t0 });
  }
};
function calcul(n) { let s = 0; for (let i = 0; i < n; i++) s += Math.sqrt(i); return s; }

// main.js
const w = new Worker(new URL("./poids-lourd.js", import.meta.url), { type: "module" });
w.postMessage({ type: "ligner", payload: 5e7 });
w.onmessage = (e) => { document.querySelector("#out").textContent = `${e.data.payload} en ${e.data.ms.toFixed(0)} ms`; };
// w.terminate();   // tuer le worker
\end{codeplain}

\begin{itemize}
  \item \link{Comlink}{https://github.com/GoogleChromeLabs/comlink} rend l'appel de worker
        presque banal (promesses \`a travers les fronti\`eres).
  \item \verb|navigator.hardwareConcurrency| donne le nombre de coeurs~; un
        \verb|WorkerPool| de \$n-1\$ workers traite les gros jeux de donn\'ees en parall\`ele.
  \item \textbf{SharedArrayBuffer} + \verb|Atomics| partagent r\'eellement la m\'emoire
        (n\'ecessite les en-t\^etes \verb|Cross-Origin-Opener-Policy| / \verb|Embedder-Opener-Policy|).
  \item \textbf{OffscreenCanvas} d\'el\`egue le rendu WebGL/canvas \`a un worker~: l'interface ne
        rame plus pendant les animations.
  \item \textbf{OPFS} (\link{Origin Private File System}{https://developer.mozilla.org/en-US/docs/Web/API/File\_System\_API/Origin\_private\_file\_system})
        est un stockage de fichiers rapide, isol\'e, utilisable depuis un worker~: c'est la voie
        s\'erieuse pour des donn\'ees volumineuses c\^ot\'e client.
  \item WebAssembly~: le langage des autres langages (Rust, C, Go). Appeler du WASM depuis JS
        est triviale~; l'\'ecrire est un autre livre. D\'emo rapide~:
        \link{Rust et wasm-pack}{https://rustwasm.github.io/wasm-pack/}.
\end{itemize}

\begin{exercice}
\begin{enumerate}
  \item Consomme une API publique r\'eelle --- \url{https://api.github.com/repos/vitejs/vite} ou
        \url{https://api.open-meteo.com/v1/forecast?latitude=14.79\&longitude=-16.67\&current=temperature\\\_2m}
        avec le mini-client~: affiche la temp\'erature de Thi\`es, g\`ere erreur r\'eseau et hors ligne.
  \item Ajoute un cache~: garde la derni\`ere r\'eponse dans \verb|sessionStorage| avec son \^age,
        affiche-la imm\'ediatement, puis rafra\^ichis en arri\`ere-plan (strat\'egie
        \og ~stale-while-revalidate~ \fg{}).
  \item Impl\'emente l'annulation~: taper une nouvelle recherche annule la requ\^ete pr\'ec\'edente,
        et une r\'eponse hors ordre est jet\'ee.
  \item \'Ecris un affichage temps r\'eel~: un petit serveur Node SSE qui \'emet l'heure toutes les
        secondes, un client qui affiche, d\'etecte la d\'econnexion et affiche \og ~hors ligne~ \fg{}.
  \item D\'eplace un calcul lourd (somme de 10\textsuperscript{8} racines carr\'ees, ou tri de
        10\textsuperscript{6} nombres) dans un worker et mesure le gain sur le FPS d'une
        animation~: utilise le panneau \clic{Performance}.
  \item Ajoute la reprise sur erreur r\'eseau~: file d'attente des \'ecritures locales, renvoi
        automatique quand \verb|online| revient, affichage du nombre d'op\'erations en attente.
\end{enumerate}
\end{exercice}

\begin{ressources}
\begin{itemize}
  \item MDN --- fetch --- \url{https://developer.mozilla.org/fr/docs/Web/API/Fetch\\\_API/Using\\\_Fetch}
  \item MDN --- CORS --- \url{https://developer.mozilla.org/fr/docs/Web/HTTP/Guides/CORS}
  \item \link{Async API tutorial (Node, SSE, WS)}{https://github.com/nodejs/node} --- docs \og ~http~ \fg{}
  \item \link{web.dev : patterns de reseau}{https://web.dev/patterns/}
  \item \link{MDN : Using Web Workers}{https://developer.mozilla.org/fr/docs/Web/API/Web\_Workers\_API/Using\_web\_workers}
\end{itemize}
\end{ressources}

\index{fetch}\index{CORS}\index{WebSocket}\index{worker}
% ============================================================================
\chapter{Modules, npm et outillage}
\label{ch:js-modules}
% ============================================================================

\begin{objectif}
\begin{itemize}
  \item \'Ecrire et lire des modules ES sans h\'esiter sur la syntaxe (\texttt{export}, \texttt{import}).
  \item Comprendre \texttt{package.json}, les versions s\'emantiques et les confusions CJS/ESM.
  \item Utiliser npm, pnpm ou yarn comme un professionnel~: scripts, lockfiles, audits.
  \item Choisir une d\'ependance sans se faire pi\'eger (taille, maintenance, licences).
  \item Cr\'eer, publier et versionner son propre paquet.
\end{itemize}
\end{objectif}

\section{Modules ES : le syst\`eme officiel}

Un module est un fichier avec sa \textbf{port\'ee propre}~: pas de variable globale, \'evaluation
une seule fois (singleton), imports r\'esolus avant ex\'ecution (graphe, pas d'ordre \`a g\'erer).

\begin{codejs}
// ---- math.js : exporter
export const PI = 3.141592653589793;
export function carree(x) { return x * x; }
export function racine(n) { return n < 0 ? null : Math.sqrt(n); }
export function* sequence(n) { for (let i = 0; i < n; i++) yield i; }
export class Vecteur { constructor(x, y) { Object.assign(this, { x, y }); } }
export default function formater(n) { return n.toLocaleString("fr-FR"); }   // UN seul default

// ---- app.js : importer
import formater, { PI, carree as sq, Vecteur } from "./math.js";      // l'extension .js est OBLIGATOIRE en ESM natif
import * as maths from "./math.js";                                    // namespace (pratique pour les gros modules)
import "./side-effect.css";                                            // effet de bord uniquement
const module = await import("./lourd.js");                             // chargement dynamique (code splitting)
export { PI as CONSTANT_PI };                                          // re-export
export * from "./math.js";
export { Vecteur as Vec2 } from "./math.js";

// Top-level await : permitted au niveau superieur d'un module (Node 14.8+, navigateurs 2021+)
const config = await fetch("/config.json").then((r) => r.json());
\end{codejs}

\begin{note}
Les imports sont \textbf{statiques et hoist\'es}~: tu peux utiliser \texttt{PI} avant la ligne
d'\texttt{import}. Tu ne peux pas mettre un \texttt{import} dans un \texttt{if}~: pour \c{c}a il
existe \texttt{await import()}. C'est exactement ce m\'ecanisme qui permet le tree-shaking
(et qui explique pourquoi \texttt{require} dynamique emp\^eche tout optimisation).
\end{note}

\subsection{CommonJS (l'ancien monde) et la cohabitation}

\begin{codejs}
// CommonJS (Node historique, tres present dans les anciens paquets)
const { carree } = require("./math.js");
module.exports = { carree };
exports.carree = carree;

// Dans un paquet ESM ("type": "module"), require n'existe pas.
// Fallback propre quand on en a vraiment besoin :
const require = createRequire(import.meta.url);
const legacy = require("./vieux-paquet.cjs");
import { createRequire } from "node:module";

// CJS peut importer ESM... seulement de facon asynchrone : (node --experimental-vm-modules ou dynamic import)
const { carree } = await import("./math.js");
\end{codejs}

R\`egles de survie~: (1) un projet neuf = ESM partout, \verb|"type": "module"|~;
(2) les extensions sont \'ecrites dans les chemins relatifs~;
(3) \verb|\_\_dirname|/\verb|require| n'existent pas en ESM~;
(4) ne m\'elange pas les deux dans le m\^eme paquet, sauf si tu publishes
\verb|exports| avec les conditions \verb|import|/\verb|require| (voir plus bas).

\section{\texttt{package.json}--- le contrat du projet}

\begin{codejson}
{
  "name": "@nour/outils",
  "version": "1.4.0",
  "type": "module",
  "description": "Boite a outils partagee : formatage, validation, HTTP.",
  "license": "MIT",
  "author": "Nour <nour@example.com>",
  "repository": { "type": "git", "url": "git+https://github.com/nour/outils.git" },
  "keywords": ["utils", "validation"],
  "sideEffects": false,
  "engines": { "node": ">=20" },
  "files": ["dist", "README.md", "LICENSE"],
  "main": "./dist/index.cjs",
  "module": "./dist/index.js",
  "types": "./dist/index.d.ts",
  "exports": {
    ".": { "types": "./dist/index.d.ts", "import": "./dist/index.js", "require": "./dist/index.cjs" },
    "./db": { "import": "./dist/db.js" },
    "./package.json": "./package.json"
  },
  "scripts": {
    "dev": "vite",
    "build": "tsc -p tsconfig.build.json && vite build",
    "start": "node dist/server.js",
    "test": "vitest run",
    "test:watch": "vitest",
    "test:e2e": "playwright test",
    "lint": "eslint . --max-warnings=0",
    "lint:fix": "eslint . --fix",
    "format": "prettier . --write",
    "typecheck": "tsc --noEmit",
    "check": "npm run lint && npm run typecheck && npm run test",
    "prepare": "husky || true",
    "release": "npm run check && npm version patch && npm publish --access public"
  },
  "dependencies": { "zod": "^4.0.0" },
  "devDependencies": {
    "vite": "^7.0.0", "vitest": "^3.0.0", "typescript": "^5.7.0",
    "eslint": "^9.0.0", "prettier": "^3.4.0", "husky": "^9.1.0",
    "lint-staged": "^15.3.0", "@vitest/coverage-v8": "^3.0.0"
  },
  "lint-staged": { "*.{js,ts,json,md,css}": ["prettier --write", "eslint --fix"] }
}
\end{codejson}

\etape{Les versions s\'emantiques, en une page}
\verb|MAJEURE.MINEURE.PATCH|~: \verb|3.1.4|. Changement \textbf{majeur} = rupture d'API.
\textbf{Mineure} = nouveaut\'e compatible. \textbf{Patch} = correction compatible.
Dans \texttt{package.json}~:

\begin{tabz}{clX}
  \headrow \thd{Pr\'efixe} \& \thd{Exemple} \& \thd{Ce que l'installateur accepte} \\
  \midrule
  \verb|^| (d\'efaut de \cmd{npm i}) \& \verb|^3.1.0| \& \$\geq\$3.1.0 \verb|<|4.0.0 --- mineures et patchs \\
  \verb|~| \& \verb|~3.1.0| \& \$\geq\$3.1.0 \verb|<|3.2.0 --- patchs seulement \\
  \verb|>|, \verb|>=|, \verb|<|, \verb|=| \& \verb|>=3.1 <3.3| \& plages explicites \\
  \verb|*| \& \verb|3.x| \& tout ce qui commence par 3 \\
  aucun \& \verb|3.1.4| \& fig\'e exactement \\
\end{tabz}

Le \textbf{lockfile} (\verb|package-lock.json|, \verb|pnpm-lock.yaml|) fige tout l'arbre, y
compris les d\'ependances de tes d\'ependances. \emph{Il doit \^etre versionn\'e}~; sans lui, deux
installations du m\^eme projet peuvent produire deux r\'esultats diff\'erents. En CI, utilise
\cmd{npm ci} (installe exactement le lockfile, refuse sinon) ou \cmd{pnpm i --frozen-lockfile}.

\etape{Semver et les prereleases}
\verb|1.0.0-rc.1| n'est \emph{pas} inclus dans \verb|^1.0.0| tant que le tag latest n'est pas
sorti. Pour tester une b\^eta~: \cmd{npm i vite@beta}.

\section{npm, pnpm, yarn, bun : les commandes du quotidien}

\begin{terminal}
# Creation
mkdir mon-projet && cd mon-projet && npm init -y         # ou : pnpm init / npm create vite@latest

# Installer / retirer / mettres a jour
npm install lodash zod                    # dependance de production
npm i -D vitest typescript                # dependance de developpement
npm i -g create-vite                      # outil global (CLI) -- a eviter avec npx
npm uninstall lodash
npm update zod
npm outdated                              # ce qui a pris du retard
npx vitest run                            # executer un binaire local SANS l'installer globalement
npm exec --package=eslint -- npx eslint .

# Inspecter / securiser
npm ls --all ; npm ls react --depth=0     # l'arbre, ou qui depend de quoi
npm why zod                                # POURQUOI ce paquet est la
npm audit                                  # vulnerabilites connues (bruyant)
npm audit --omit=dev ; npm audit fix       # ou npm audit fix --force (attention aux majeures)
pnpm audit --audit-level=high
npm view vite version repository homepage dist.unpackedSize time --json
npm outdated --json | jq -r 'to_entries[] | "\(.key) \(.value.current) -> \(.value.wanted)"'

# Nettoyage / reparation
rm -rf node_modules package-lock.json && npm i      # l'option "tout casser et recommencer"
npm cache clean --force
pnpm store prune

# Gestion de versions et publication
npm version patch|minor|major            # bump + tag git
npm publish --access public
npm deprecate mon-paquet "Utilise @nour/outils a la place"
\end{terminal}

\subsection{Quel gestionnaire de paquets ?}

\begin{tabz}{lXl}
  \headrow \thd{Outil} \& \thd{Pourquoi lui} \& \thd{Vitesse (relative)} \\
  \midrule
  \link{npm}{https://docs.npmjs.com/} \& livr\'e avec Node, standard, suffisant. \& 1\$\times\$ \\
  \link{pnpm}{https://pnpm.io/fr/} \& store global + liens durs~: disque divis\'e par 3, monorepo naturel, \texttt{node\_modules} strict (moins de bugs d'imports fant\^omes). \textbf{Mon choix par d\'efaut.} \& 2--3\$\times\$ plus rapide \\
  \link{yarn}{https://yarnpkg.com/} \& m\^ur, PnP (z\'ero \texttt{node\_modules}), bien en monorepo. \& \$\approx\$ pnpm \\
  \link{bun}{https://bun.sh/docs/install/} \& le gestionnaire est int\'egr\'e au runtime~: tr\`es rapide, API diff\'erente. \& tr\`es rapide \\
\end{tabz}

\begin{terminal}
# pnpm : installation propre (via corepack, livre avec Node)
corepack enable
corepack prepare pnpm@latest --activate
pnpm install
pnpm add -D typescript
pnpm -r run build          # monorepo : construit tous les paquets
\end{terminal}

\section{Les scripts npm : ton vrai syst\`eme de build}

Ne sous-estime pas ce m\'ecanisme~: sur 90 \% des projets, un \verb|Makefile|-like bien nomm\'e
dans \verb|scripts| vaut mieux qu'une usine \`a gaz.

\begin{codejson}
{
  "scripts": {
    "dev": "npm-run-all --parallel dev:api dev:web",
    "dev:web": "vite",
    "dev:api": "node --watch --env-file=.env src/server.js",
    "build": "npm run build:css && vite build",
    "preview": "vite preview",
    "test": "vitest run --coverage",
    "lint": "eslint . --max-warnings=0",
    "typecheck": "tsc --noEmit",
    "check": "run-s lint typecheck test",
    "db:migrate": "node scripts/migrate.js",
    "db:seed": "node scripts/seed.js",
    "hooks:install": "husky"
  }
}
\end{codejson}

\begin{itemize}
  \item Les pr\'e/post existent~: \verb|"prebuild"| s'ex\'ecute avant \verb|build| (depuis toujours).
  \item \verb|npm-run-all| / \verb|concurrently| pour parall\'eliser~; \verb|run-s| pour s\'equencer proprement.
  \item Les variables \verb|npm\_package\_json|, \verb|npm\_lifecycle\_event| sont expos\'ees~;
        \cmd{npm run dev -- --port 4000} passe les arguments apr\`es \verb|--|.
  \item Node~20.6+ lit un fichier d'environnement tout seul~: \cmd{node --env-file=.env app.js},
        et \link{dotenv}{https://github.com/motdotla/dotenv} fait la m\^eme chose depuis plus longtemps.
        \label{sec:dotenv}
\end{itemize}

\section{Choisir une d\'ependance (la check-list anti-regret)}

\begin{enumerate}
  \item \textbf{Est-ce que j'en ai vraiment besoin ?} Une fonction de 15 lignes ne justifie pas
        300~Ko de d\'ependances transitives. \verb|npm ls| et \og ~combien de paquets cela ajoute-t-il~?~ \fg{}
  \item \textbf{Maintenance}~: dernier commit, nombre de contributeurs, issues ouvertes,
        \verb|releaseType| r\'ecurrent. Un paquet sans release depuis 4 ans est un passif.
  \item \textbf{Adoption}~: d\'ependants invers\'es sur \link{nicedoc.io}{https://npmgraph.js.org/},
        t\'el\'echargements \cmd{npm view mon-paquet} ou \link{napi}{https://npm-stat.com/charts.html}.
  \item \textbf{Taille}~: \link{bundlephobia}{https://bundlephobia.com/} (poids gzipp\'e du bundle).
  \item \textbf{Licence}~: MIT/Apache-2.0/ISC = tranquille. GPL/AGPL = v\'erifie ton usage.
        \link{Choose a license}{https://choosealicense.com/}.
  \item \textbf{S\'ecurit\'e}~: \cmd{npm audit}, \link{Socket}{https://socket.dev/},
        \link{OpenSSF Scorecard}{https://scorecard.dev/}.
  \item \textbf{Alternative native}~: beaucoup de biblioth\`eques populaires (lodash pour une
        partie, moment.js, axios en partie, classnames, uuid) ont \'et\'e remplac\'ees par le langage.
\end{enumerate}

\subsection{Quelques paquets qui reviennent dans tous les bons projets}

\begin{tabz}{lX}
  \headrow \thd{Paquet} \& \thd{Usage} \\
  \midrule
  \link{zod}{https://zod.dev/} \& validation de donn\'ees + types TypeScript d\'eriv\'es. Le meilleur rapport utilit\'e/ligne de code. \\
  \link{date-fns}{https://date-fns.org/} ou \link{Day.js}{https://day.js.org/} \& dates sans monolithe (ou \verb|Temporal| en natif, ES2024+). \\
  \link{nanoid}{https://github.com/ai/nanoid} ou \verb|crypto.randomUUID()| \& identifiants. \\
  \link{picocolors}{https://github.com/alexeyraspopov/picocolors} \& couleurs en CLI, 2~Ko. \\
  \link{commander}{https://github.com/tj/commander.js} / \link{yargs}{https://yargs.js.org/} / \link{citty}{https://github.com/unjs/citty} \& cr\'eer un CLI. \\
  \link{p-limit}{https://github.com/sindresorhus/p-limit} / \link{throat}{https://github.com/ForbesLindesay/throat} \& concurrence born\'ee (ou ta version maison, cf. \cref{ch:js-async}). \\
  \link{fast-glob}{https://github.com/mrmlnc/fast-glob} \& recherche de fichiers par motif. \\
  \link{chalk}{https://github.com/chalk/chalk} \& sortie CLI styl\'ee. \\
  \link{vitest}{https://vitest.dev/} \& runner de test rapide, compatible Jest, int\'egr\'e \`a Vite. \\
  \link{msw}{https://mswjs.io/} \& mock d'API r\'eseau au niveau r\'eseau : tes tests ignorent fetch. \\
  \link{playwright}{https://playwright.dev/} \& tests bout-en-bout multi-navigateurs. \\
  \link{tsx}{https://github.com/privatenumber/tsx} \& TypeScript dans Node sans \'etape de compilation. \\
\end{tabz}

\section{Publier un paquet : la recette compl\`ete}

\begin{terminal}
mkdir outils && cd outils && npm init -y
npm pkg set type=module name=@nour/outils version=0.1.0 license=MIT sideEffects=false
npm pkg set exports.import=./dist/index.js exports.require=./dist/index.cjs
npm install -D tsup typescript
# package.json -> "scripts": { "build": "tsup src/index.ts --format esm,cjs --dts --clean" }
npm login          # ou : npm publish --provenance --access public (avec un token precise)
npm version patch && npm publish
npm view @nour/outils versions
# cote consommateur :
npm i @nour/outils && node -e 'import("@nour/outils").then(m => console.log(m.carree(9)))'
\end{terminal}

Sur GitHub, publie aussi via \link{npm + Actions}{https://docs.github.com/en/actions/publishing-packages/publishing-nodejs-packages}
avec \verb|NODE\_AUTH\_TOKEN| et un registre \link{GitHub Packages}{https://github.com/features/packages}.

\begin{exercice}
\begin{enumerate}
  \item Transforme tes utilitaires des chapitres pr\'ec\'edents (\verb|debounce|, \verb|retry|,
        \verb|echapper|, \verb|chunk|) en un paquet ESM \verb|@tonnom/utils| avec
        \verb|tsup|, types, et un \verb|README| d'exemples. Publie-le.
  \item Mets le paquet dans \verb|package.json| d'un projet de test et remplace ton code
        local par l'import. Simule une incompatibilit\'e~: passe \`a la version majeure suivante et
        observe la casse~; explique la politique de versions que tu choisirais.
  \item Installe un CLI depuis npm (\verb|npm i -g httpie| ou \verb|npx create-vite@latest|) puis
        explique pourquoi \verb|npx| est pr\'ef\'erable \`a \verb|-g|.
  \item Lance \verb|npm outdated| sur un vrai projet, puis fais monter une d\'ependance mineure
        en gardant les tests au vert. Note dans ton carnet ce qui a cass\'e.
  \item Ex\'ecute \verb|npm ls| et rep\`ere une d\'ependance dont tu ne comprends pas la pr\'esence~;
        trouve son origine avec \verb|npm why|~; d\'ecide de la garder ou de l'\'eliminer.
  \item Configure un \verb|prepare| npm + \link{husky}{https://typicode.github.io/husky/}
        pour que \verb|lint-staged| tourne avant chaque commit.
\end{enumerate}
\end{exercice}

\begin{ressources}
\begin{itemize}
  \item Docs Node --- Modules --- \url{https://nodejs.org/api/packages.html}
  \item MDN --- Modules JavaScript --- \url{https://developer.mozilla.org/fr/docs/Web/JavaScript/Guide/Modules}
  \item npm docs --- \url{https://docs.npmjs.com/cli/} --- pnpm --- \url{https://pnpm.io/fr/motivation}
  \item \link{Node ESM : le guide complet}{https://gist.github.com/sindresorhus/a39789f98801d908bbc7ff3ecc99d99c}
  \item \link{Phenomenon de la chaine d'approvisionnement}{https://en.wikipedia.org/wiki/Supply\_chain\_attack} --- pourquoi la prudence paie.
\end{itemize}
\end{ressources}

\index{modules}\index{npm}\index{package.json}\index{semver}
% ============================================================================
\chapter{Node.js : serveur, CLI, scripts d'automatisation}
\label{ch:node}
% ============================================================================

\begin{objectif}
\begin{itemize}
  \item Utiliser les modules de la biblioth\`eque standard~: \texttt{fs}, \texttt{path}, \texttt{os},
        \texttt{crypto}, \texttt{child\_process}, \texttt{worker\_threads}.
  \item \'Ecrire un serveur HTTP nu, puis un serveur Express complet~: routes, middleware,
        validation, erreurs, logs, upload.
  \item Structurer une application Node en couches, avec configuration par environnement.
  \item \'Ecrire des scripts CLI utiles (et les rendre robustes~: signaux, exit codes, stdout/stderr).
\end{itemize}
\end{objectif}

\section{Le mod\`ele, en trois phrases}

Node ex\'ecute du JavaScript sur V8 avec une boucle d'\'ev\'enements non bloquante~: les entr\'ees/sorties
sont d\'el\'egu\'ees \`a une pool de threads (libuv) et leurs rappels sont rejou\'es sur le fil principal.
Un processus ne g\`ere donc \emph{qu'un calcul \`a la fois}~: pour utiliser plusieurs coeurs, on lance
plusieurs processus (\verb|cluster|, \verb|pm2 -i max|, ou un conteneur par instance derri\`ere nginx).

\section{La biblioth\`eque standard qui compte}

\begin{codejs}
// ---- path, fs, os
import { join, dirname, basename, extname, resolve } from "node:path";
import { readFile, writeFile, readdir, mkdir, stat, rm, cp, appendFile, access } from "node:fs/promises";
import { existsSync, createReadStream, createWriteStream, watch } from "node:fs";
import { homedir, tmpdir, cpus, totalmem } from "node:os";

const here = dirname(new URL(import.meta.url).pathname);
const data = await readFile(join(here, "data.json"), "utf8");
await writeFile(join(here, "out.txt"), "ok\n", { flag: "a" });
await mkdir(join(here, "a/b/c"), { recursive: true });
await cp("src", "backup", { recursive: true, force: true });
for (const entree of await readdir(here, { withFileTypes: true })) {
  if (entree.isDirectory() && !entree.name.startsWith(".")) console.log(entree.name);
}
await access("fichier.txt").then(() => console.log("lisible")).catch(() => console.log("absent ou interdit"));

// GROS fichier ligne a ligne sans le charger en memoire :
import { createInterface } from "node:readline";
const rl = createInterface({ input: createReadStream("gros.log"), crlfDelay: Infinity });
let n = 0;
for await (const ligne of rl) if (/ERR/.test(ligne)) n++;
console.log(n, "lignes d'erreur");

// GROS ecriture en flux (les backpressures sont gerees par pipeline)
import { pipeline } from "node:stream/promises";
import { createGzip } from "node:zlib";
await pipeline(createReadStream("gros.csv"), createGzip(), createWriteStream("gros.csv.gz"));

// Surveillance de fichiers
const w = watch(here, { persistent: false }, (evt, filename) => console.log(evt, filename));
w.close();

// ---- crypto : hachage, jetons, chiffrement
import { createHash, randomBytes, createHmac, randomUUID, scrypt, timingSafeEqual } from "node:crypto";
createHash("sha256").update("mot de passe + sel").digest("hex");
randomUUID();                       // "d94b..." (v4)
randomBytes(32).toString("base64url");   // jeton d'API
const sel = randomBytes(16);
const cle = await promisify(scrypt)(mdp, sel, 64);   // derive de cle (en prod : argon2 ou bcrypt)
timingSafeEqual(a, b);              // comparaison sans fuite temporelle

// ---- child_process
import { execFile, spawn } from "node:child_process";
import { promisify } from "node:util";
const run = promisify(execFile);
const { stdout, stderr } = await run("git", ["rev-parse", "--short", "HEAD"]);
const ff = spawn("ffmpeg", ["-i", "in.mov", "out.mp4"], { stdio: "inherit" });
ff.on("exit", (code, sig) => console.log("fini", code, sig));

// ---- process : signaux, code de sortie, arguments
const args = process.argv.slice(2);
const opts = Object.fromEntries(args.filter((a) => a.startsWith("--")).map((a) => {
  const [k, v = "true"] = a.replace(/^--/, "").split("=");
  return [k, v === "true" ? true : v];
}));
process.on("SIGINT", () => arretPropre(130));
process.on("SIGTERM", () => arretPropre(143));
async function arretPropre(code) {
  await Promise.allSettled([serveur.closeAsync?.(), db.end()]);
  process.exitCode = code;      // ne pas appeler process.exit brutalement : on perd les bufferisations
}
console.error("aller sur stderr"); process.exitCode = 1;   // exit code correct pour les scripts
\end{codejs}

\section{Le serveur HTTP nu, pour comprendre}

\begin{codejs}
import { createServer } from "node:http";

const serveur = createServer(async (req, res) => {
  const url = new URL(req.url, `http://${req.headers.host}`);
  const json = (status, objet) => {
    const corps = JSON.stringify(objet);
    res.writeHead(status, {
      "content-type": "application/json; charset=utf-8",
      "content-length": Buffer.byteLength(corps),
      "x-powered-by": "node nu",
    });
    res.end(corps);
  };
  try {
    if (url.pathname === "/api/sante") return json(200, { ok: true, uptime: process.uptime() });
    if (url.pathname === "/api/echo" && req.method === "POST") {
      const chunks = [];
      for await (const c of req) chunks.push(c);
      return json(200, { recu: JSON.parse(Buffer.concat(chunks).toString() || "{}") });
    }
    json(404, { erreur: `route inconnue ${req.method} ${url.pathname}` });
  } catch (e) {
    json(400, { erreur: e.message });
  }
});
serveur.listen(3000, () => console.log("http://localhost:3000"));
serveur.keepAliveTimeout = 65_000;            // > au timeout du reverse proxy (sinon 502 aleatoires)
serveur.requestTimeout = 300_000;
process.on("SIGTERM", () => serveur.close(() => process.exit(0)));
\end{codejs}

\section{Express~: le cadre minimaliste de r\'ef\'erence}

Documentation~: \link{expressjs.com/fr}{https://expressjs.com/fr/}. Version~5 (2024+) g\`ere
enfin les erreurs des fonctions \texttt{async} sans encapsulation.

\begin{terminal}
mkdir api && cd api && npm init -y
npm pkg set type=module
npm i express helmet cors compression morgan zod dotenv
npm i -D nodemon eslint
\end{terminal}

\begin{codejs}
// src/app.js
import express from "express";
import helmet from "helmet";
import cors from "cors";
import morgan from "morgan";
import { routers } from "./routes/index.js";
import { environement } from "./config.js";
import { gestionErreur } from "./erreurs.js";

export function creerApp() {
  const app = express();
  app.disable("x-powered-by");
  app.set("trust proxy", 1);                       // derriere nginx : lire X-Forwarded-*
  app.use(helmet({ crossOriginResourcePolicy: { policy: "cross-origin" } }));
  app.use(cors({ origin: environement.origines, credentials: true }));
  app.use(express.json({ limit: "1mb" }));
  app.use(express.urlencoded({ extended: true, limit: "100kb" }));
  app.use(morgan(environement.prod ? "combined" : "dev", {
    skip: (req) => req.path === "/favicon.ico",
  }));
  app.use("/static", express.static("public", { maxAge: "7d", immutable: true }));

  routers.forEach(([chemin, r]) => app.use(chemin, r));
  app.use((req, res) => res.status(404).json({ erreur: "introuvable", chemin: req.path }));
  app.use(gestionErreur);
  return app;
}

// src/routes/articles.js
import { Router } from "express";
import { z } from "zod";
import { asyncHandler, HttpError } from "../erreurs.js";
import { depot } from "../depot.js";

const schemaArticle = z.object({
  titre: z.string().trim().min(3).max(160),
  contenu: z.string().trim().min(10),
  tags: z.array(z.string().min(1).max(30)).max(12).default([]),
  publie: z.boolean().default(false),
});

export const routeArticles = Router();

routeArticles.get("/", asyncHandler(async (req, res) => {
  const q = z.object({
    page: z.coerce.number().int().min(1).default(1),
    taille: z.coerce.number().int().min(1).max(100).default(20),
    recherche: z.string().trim().max(120).optional(),
  }).parse(req.query);
  const { items, total } = await depot.liste({ ...q });
  res.links({ next: total > q.page * q.taille ? `?page=${q.page + 1}` : undefined });
  res.json({ items, total, page: q.page });
}));

routeArticles.get("/:id(\\d+)", asyncHandler(async (req, res) => {
  const article = await depot.parId(Number(req.params.id));
  if (!article) throw new HttpError(404, "Article introuvable");
  res.set("etag", `W/"${article.maj}"`).vary("accept-encoding");
  res.status(200).json(article);
}));

routeArticles.post("/", asyncHandler(async (req, res) => {
  const donnees = schemaArticle.parse(req.body);
  const cree = await depot.creer(donnees);
  res.location(`/api/articles/${cree.id}`).status(201).json(cree);
}));

routeArticles.delete("/:id(\\d+)", asyncHandler(async (req, res) => {
  await depot.supprimer(Number(req.params.id));
  res.status(204).end();
}));

// src/erreurs.js
export class HttpError extends Error {
  constructor(status, message, details) { super(message); this.status = status; this.details = details; }
}
export const asyncHandler = (fn) => (req, res, next) => Promise.resolve(fn(req, res, next)).catch(next);
export function gestionErreur(err, req, res, _next) {
  if (err?.name === "ZodError") {
    return res.status(422).json({ erreur: "Donnees invalides", details: err.issues.map((i) => ({ champ: i.path.join("."), message: i.message })) });
  }
  const status = err.status ?? err.statusCode ?? 500;
  if (status >= 500) console.error({ msg: err.message, stack: err.stack, url: req.originalUrl, ip: req.ip });
  res.status(status).json({ erreur: status >= 500 ? "Erreur interne" : err.message });
}
\end{codejs}

\subsection{Configuration par environnement~: la bonne habitude d\`es le premier jour}
\label{sec:node-config}

\begin{codejs}
// src/config.js --- une seule porte d'entree, validee au demarrage
import "dotenv/config";
import { z } from "zod";

const schema = z.object({
  NODE_ENV: z.enum(["development", "test", "production"]).default("development"),
  PORT: z.coerce.number().int().default(3000),
  DATABASE_URL: z.string().url(),
  JWT_SECRET: z.string().min(32),
  LOG_LEVEL: z.enum(["debug", "info", "warn", "error"]).default("info"),
  ORIGINES: z.string().default("http://localhost:5173"),
  UPLOAD_MAX_MO: z.coerce.number().default(5),
});
const parsed = schema.safeParse(process.env);
if (!parsed.success) {
  console.error("Configuration invalide :", parsed.error.flatten().fieldErrors);
  process.exit(1);
}
export const environement = {
  ...parsed.data,
  prod: parsed.data.NODE_ENV === "production",
  origines: parsed.data.ORIGINES.split(",").map((s) => s.trim()),
};
\end{codejs}

Jamais de secret dans le code, dans Git, ni dans une variable expos\'ee au navigateur.
\verb|.env| pour le dev local, variables d'environnement r\'eelles en production
(\link{Doppler}{https://doppler.com/}, \link{1Password~CLI}{https://developer.1password.com/docs/cli/},
ou simplement le fichier d'environnement de systemd / de ton h\'ebergeur).

\section{Ce qu'Express ne fait pas pour toi (et qu'il faut ajouter)}

\begin{itemize}
  \item \textbf{Cookies et sessions}~: \link{cookie-parser}{https://github.com/pillarjs/cookie-parser},
        \link{express-session}{https://github.com/expressjs/session} (ou des sessions sign\'ees sans
        \'etat, cf. \cref{ch:js-secu}).
  \item \textbf{Upload}~: \link{multer}{https://github.com/expressjs/multer}, avec taille maximale,
        liste blanche de types et stockage hors de la racine web.
  \item \textbf{Limitation de d\'ebit}~: \link{express-rate-limit}{https://github.com/nickdesclos/express-rate-limit}.
  \item \textbf{Corr\'elation et logs structur\'es}~: \link{pino}{https://github.com/pinojs/pino}
        (JSON, tr\`es rapide) plut\^ot que \verb|console.log|.
  \item \textbf{T\^aches diff\'er\'ees}~: \link{BullMQ}{https://docs.bullmq.io/} sur Redis.
  \item \textbf{Temps r\'eel}~: \link{Socket.IO}{https://socket.io/fr/} ou WebSocket brut.
\end{itemize}

\begin{codejs}
// Exemple d'upload durci
import multer from "multer";
import { extname } from "node:path";
import { randomUUID } from "node:crypto";

const TYPES = new Map([["image/png", ".png"], ["image/jpeg", ".jpg"], ["image/webp", ".webp"]]);
const upload = multer({
  storage: multer.diskStorage({
    destination: "var/uploads",
    filename: (req, file, cb) => cb(null, `${randomUUID()}${TYPES.get(file.mimetype) ?? ".bin"}`),
  }),
  limits: { fileSize: 5 * 1024 * 1024, files: 1 },
  fileFilter: (req, file, cb) => cb(TYPES.has(file.mimetype) ? null : new Error("type refuse"), TYPES.has(file.mimetype)),
});
route.post("/avatar", upload.single("fichier"), asyncHandler(async (req, res) => {
  if (!req.file) throw new HttpError(422, "Fichier manquant ou type non autorise");
  res.json({ url: `/u/${req.file.filename}` });
}));
\end{codejs}

\section{Alternatives \`a Express, et quand les choisir}

\begin{tabz}{lXl}
  \headrow \thd{Cadre} \& \thd{Id\'ee directrice} \& \thd{Lien} \\
  \midrule
  \link{Hono}{https://hono.dev/} \& ultra-l\'eger, tourne sur Node, Bun, Deno, Cloudflare Workers, navigateur. Mon choix pour une API neuve. \& 14 kB \\
  \link{Fastify}{https://fastify.dev/} \& plus rapide qu'Express, sch\'emas JSON int\'egr\'es, plugins. \& \$\approx\$2x \\
  \link{Koa}{https://koajs.com/} \& le \og ~Express repens\'e~ \fg{} par les auteurs d'Express. \& minimaliste \\
  \link{NestJS}{https://nestjs.com/} \& architecture d\'ecorateurs/injection, tr\`es TypeScript, \'equipes Java qui passent au Node. \& complet \\
  \link{Next.js}{https://nextjs.org/} \& React fullstack~: pages, API routes, rendu serveur. \& cadre UI \\
  \link{Express~5}{https://expressjs.com/fr/} \& le standard historique, documentation immense. \& classique \\
\end{tabz}

\section{Un CLI utile en 60 lignes}

\begin{codejs}
#!/usr/bin/env node
// scripts/rename.js --- renomme en lot : node rename.js --dossier ./photos --prefix 2026 --debut 1
import { readdir, rename, stat } from "node:fs/promises";
import { join, extname } from "node:path";

const opts = Object.fromEntries(process.argv.slice(2).map((a) => {
  const [k, v = "true"] = a.replace(/^--/, "").split("=");
  return [k, v];
}));

if (opts.aide) {
  console.log(`Usage : rename --dossier CHEMIN --prefix TEXTE [--debut N] [--sec]

  --dossier  dossier a traiter (obligatoire)
  --prefix   debut du nouveau nom
  --debut    numero de depart (defaut 1)
  --sec      applique vraiment (sinon : simulation)`);
  process.exit(0);
}
const dossier = opts.dossier ?? process.exit(fail("--dossier est requis"));
const prefix = opts.prefix ?? "fichier";
let compteur = Number(opts.debut ?? 1);

const fichiers = (await readdir(dossier))
  .filter((f) => !f.startsWith("."))
  .sort((a, b) => a.localeCompare(b, "fr", { numeric: true }));

for (const f of fichiers) {
  const neuf = `${prefix}-${String(compteur++).padStart(3, "0")}${extname(f).toLowerCase()}`;
  if (neuf === f) continue;
  const { size, mtime } = await stat(join(dossier, f));
  console.log(`${f}  ->  ${neuf}   (${(size / 1024).toFixed(0)} Ko, ${mtime.toISOString().slice(0, 10)})`);
  if (opts.sec) await rename(join(dossier, f), join(dossier, neuf));
}
if (!opts.sec) console.log("\nSimulation. Relance avec --sec pour appliquer.");

function fail(msg) { console.error(`Erreur : ${msg}`); return 1; }
\end{codejs}

\subsection{Automatiser~: Makefile et scripts de d\'eveloppement}

\begin{codeconf}
# Makefile --- la interface unique du projet (fais `make` pour voir la liste)
.PHONY: help dev check test lint db-reset docker-up docker-down

help:            ## affiche cette aide
	@grep -E '^[a-zA-Z_-]+:.*?## ' $(MAKEFILE_LIST) | awk 'BEGIN {FS = ":.*?## "}; {printf "  \033[36m%-12s\033[0m %s\n", $$1, $$2}'
dev:             ## lance API + site en parallle
	@npm run dev
check: lint test ## tout verifier (lint + typecheck + tests)
	npm run typecheck
lint:            ## analyse statique
	npm run lint -- --fix
test:            ## tests unitaires
	npm test
db-reset:        ## recree la base et la remplit
	node scripts/migrate.js --fresh && node scripts/seed.js
docker-up:       ## demarre les services (MySQL, Redis, Mailpit)
	docker compose up -d
docker-down:     ## arrete tout
	docker compose down
\end{codeconf}

\begin{codejson}
{
  "name": "projet-node",
  "scripts": { "dev": "node --watch --env-file=.env src/server.js" },
  "devDependencies": { "nodemon": "^3.1.0" }
}
\end{codejson}

Depuis Node 18.11, \cmd{node --watch} suffit au quotidien. Pour recharger aussi au changement de
fichier non-JS (views, .env), \link{nodemon}{https://nodemon.io/} reste pratique. Pour red\'emarrer
proprement un processus en production~: \link{pm2}{https://pm2.keymetrics.io/} (ou mieux~: systemd).

\begin{terminal}
# pm2 en production (si tu n'as pas de orchestrateur)
npm i -g pm2
pm2 start src/server.js -i max --name api --node-args="--max-old-space-size=1024"
pm2 logs api --lines 100
pm2 monit ; pm2 reload api ; pm2 save ; pm2 startup
\end{terminal}

\begin{exercice}
\begin{enumerate}
  \item \'Ecris un serveur HTTP nu qui sert un fichier statique depuis \verb|public/| avec les
        en-t\^etes corrects~: \verb|content-type| (par extension), \verb|cache-control|, ETag
        bas\'e sur taille+date, et gestion du \verb|If-None-Match| (304).
  \item Reprends-le en Express, en ajoutant~: route \verb|/api/sante|, validation par
        \verb|zod| d'un POST, gestion d'erreur centralis\'ee, et logs \verb|pino|.
  \item Construis une API CRUD compl\`ete pour un blog (articles + commentaires) stock\'ee sur
        disque en JSON, avec pagination, recherche, et code de sortie correct partout.
        \emph{Ce sera ta base de comparaison avec les versions PHP et Java.}
  \item \'Ecris un CLI \verb|depense|~: \verb|add "cafe" 2.5 --cat=alimentation|,
        \verb|list --month=2026-01|, \verb|total --by=category|, donn\'ees dans
        \verb|~/.depenses.json|. Ajoute \verb|--json| pour la sortie machine.
  \item \'Ecris un script qui v\'erifie la sant\'e de ton API (5 requ\^etes, timeouts, format) et
        retourne un code de sortie non nul en cas d'\'echec~: utilisable dans un cron ou une CI.
  \item Lance deux instances derri\`ere \cmd{node --cluster} (ou l'option \verb|-i max| de pm2)
        et mesure le d\'ebit avec \link{autocannon}{https://github.com/mcollina/autocannon} :
        explique le r\'esultat avec le mod\`ele mono-fil.
\end{enumerate}
\end{exercice}

\begin{ressources}
\begin{itemize}
  \item Docs Node (la r\'ef\'erence absolue) --- \url{https://nodejs.org/fr/docs/latest/api/}
  \item Node --- \emph{Learn JavaScript on the server} --- \url{https://nodejs.org/learn}
  \item Express guide --- \url{https://expressjs.com/fr/guide/routing.html}
  \item \link{Node Best practices (Goldberd)}{https://github.com/goldberio/express-nodejs-best-practices}
  \item \link{Node Security Workshop}{https://github.com/ThankyouForComing/Node-Security-Workshop} --- et \url{https://snyk.io/fr/blog/}
\end{itemize}
\end{ressources}

\index{Node.js}\index{Express}\index{CLI}\index{serveur}
% ============================================================================
\chapter{Stockage, cache et fonctionnement hors ligne}
\label{ch:js-storage}
% ============================================================================

\begin{objectif}
\begin{itemize}
  \item Choisir le bon stockage selon quatre crit\`eres~: taille, dur\'ee, origine, secret.
  \item Utiliser \texttt{localStorage}/\texttt{sessionStorage} sans se faire pi\'eger par le JSON.
  \item Lire/\'ecrire des cookies correctement (expire, SameSite, Secure, HttpOnly).
  \item Manipuler IndexedDB proprement, y compris avec des index et des curseurs.
  \item Poser un service worker pour un vrai mode hors ligne et une PWA installable.
\end{itemize}
\end{objectif}

\section{Quoi stocker o\`u}

\begin{tabz}{llllX}
  \headrow \thd{M\'ecanisme} \& \thd{Capacit\'e typique} \& \thd{Dur\'ee} \& \thd{Serveur ?} \& \thd{Usage correct} \\
  \midrule
  cookie \& 4~Ko par domaine \& jusqu'\`a \verb|Expires| \& oui, \`a chaque requ\^ete \& session, pr\'ef\'erence lue c\^ot\'e serveur \\
  \verb|localStorage| \& 5--10~Mo \& infini, par origine \& non \& pr\'ef\'erences, brouillon de formulaire, \'etat d'UI \\
  \verb|sessionStorage| \& idem \& onglet courant \& non \& \'etat d'un wizard, jeton court \\
  \verb|IndexedDB| \& centaines de Mo (jusqu'au quota disque) \& infini \& non \& grandes quantit\'es structur\'ees, recherche, hors ligne \\
  OPFS (files) \& quota disque \& infini \& non \& fichiers, bases SQLite WASM, gros exports \\
  Cache API \& quota \& infini \& non \& r\'eponses HTTP mises en cache par le service worker \\
  \verb|caches|+SW \& --- \& --- \& --- \& application enti\`ere hors ligne \\
\end{tabz}

\begin{piege}
\textbf{Aucun stockage c\^ot\'e navigateur n'est digne de confiance.} L'utilisateur peut lire,
modifier, supprimer n'importe quoi (console, extension, disque). Ne stocke donc jamais l\`a~:
un r\^ole \og ~admin~ \fg{}, un prix, un solde, une validation de paiement, un jeton de session sensible
sans expiration courte. Le stockage local est un \emph{cache}, pas une source de v\'erit\'e.
\end{piege}

\section{Web Storage --- proprement, une fois pour toutes}

\begin{codejs}
// localStorage ne stocke QUE des chaines : un objet passe par JSON, et ca peut echouer.
class Stockage {
  constructor(prefixe = "app:", storage = localStorage) { Object.assign(this, { prefixe, storage }); }
  #cle(k) { return this.prefixe + k; }
  lire(k, defaut) {
    try { const brut = this.storage.getItem(this.#cle(k)); return brut === null ? defaut : JSON.parse(brut); }
    catch (e) { console.warn(`storage.lire(${k})`, e); this.effacer(k); return defaut; }
  }
  ecrire(k, valeur) {
    try { this.storage.setItem(this.#cle(k), JSON.stringify(valeur)); return true; }
    catch (e) {
      if (e.name === "QuotaExceededError") { this.nettoyer(); return this.ecrire(k, valeur); }
      console.warn(`storage.ecrire(${k})`, e); return false;
    }
  }
  effacer(k) { this.storage.removeItem(this.#cle(k)); }
  nettoyer() {  // enleve les cles trop vieilles ( convention : cle + "__t" = timestamp)
    const ages = Object.keys(this.storage)
      .filter((k) => k.startsWith(this.prefixe))
      .map((k) => ({ k, t: Number(this.storage.getItem(`${k}__t`) ?? 0) }))
      .sort((a, b) => a.t - b.t);
    ages.slice(0, Math.ceil(ages.length / 2)).forEach(({ k }) => this.effacer(k.slice(this.prefixe.length)));
  }
  get taille() { let n = 0; for (const k of Object.keys(this.storage)) if (k.startsWith(this.prefixe)) n += this.storage.getItem(k).length; return n; }
}
export const store = new Stockage("carnet:");
store.ecrire("articles", [{ id: 1 }]);
store.lire("articles", []);

// Synchro entre onglets : l'evenement "storage" se declenche sur les AUTRES onglets
window.addEventListener("storage", ({ key, newValue, oldValue, url }) => {
  if (key === "carnet:articles") rendre(JSON.parse(newValue ?? "[]"));
});
\end{codejs}

\section{Cookies~: le seul stockage vu par le serveur}

\begin{codejs}
// Ecriture depuis JS : pas de HttpOnly (sinon JS ne peut pas le lire), 4 Ko max
document.cookie = "theme=sombre; path=/; max-age=31536000; SameSite=Lax; Secure";

const cookies = Object.fromEntries(
  document.cookie.split("; ").map((p) => { const [k, ...v] = p.split("="); return [k, decodeURIComponent(v.join("="))]; })
);
const effacer = (nom) => { document.cookie = `${nom}=; path=/; max-age=0`; };
\end{codejs}

\begin{tabz}{lX}
  \headrow \thd{Attribut} \& \thd{Effet} \\
  \midrule
  \verb|Expires| / \verb|Max-Age| \& dur\'ee de vie~; sans lui, cookie de session (mort \`a la fermeture). \\
  \verb|Path| \& ne sera envoy\'e que sur ce chemin. \verb|/| dans 99 \% des cas. \\
  \verb|Domain| \& sous-domaines inclus ; ne le mets pas sauf besoin r\'eel (\verb|.exemple.com| = large surface). \\
  \verb|Secure| \& uniquement en HTTPS --- obligatoire si \verb|SameSite=None|. \\
  \verb|HttpOnly| \& \textbfJavaScript ne peut pas le lire~: la meilleure d\'efense contre le vol de session par XSS. \\
  \verb|SameSite=Lax|  \& d\'efaut moderne~; bloque la plupart des CSRF. \verb|Strict| = encore plus raide. \verb|None| = tiers, exige \verb|Secure|. \\
  \verb|Partitioned| \& cookie cloisonn\'e par site parent (CHIPS) : pour les iframe tierces. \\
\end{tabz}

\begin{note}
Depuis quelques ann\'ees, les navigateurs traitent un cookie \verb|SameSite=None| sans
\verb|Secure| comme invalide, et bloquent les cookies tiers par d\'efaut. Conclusion~: une
session doit \^etre un cookie \verb|HttpOnly; Secure; SameSite=Lax|, et ton JavaScript ne doit
jamais avoir besoin de lire le jeton de session. Si tu stockes un jeton d'acc\`es dans
\verb|localStorage| pour \'eviter cela, accepte d'\^etre vuln\'erable au XSS~: c'est le vrai d\'ebat
JWT-localStorage vs cookie-HttpOnly, et la r\'eponse raisonnable est \og ~cookie + CSRF~ \fg{}.
\end{note}

\section{IndexedDB --- la base de donn\'ees du navigateur}

C'est une base orient\'ee enregistrements~: magasins (\texttt{objectStore}), cl\'es, index,
transactions. API de bas niveau verbeuse~; \link{idb}{https://github.com/jakearchibald/idb}
(\cmd{npm i idb}) la rend supportable sans lui retirer sa puissance.

\begin{codejs}
import { openDB } from "idb";

const db = await openDB("carnet", 2, {
  upgrade(bdd, ancienne, nouvelle, t) {
    if (!bdd.objectStoreNames.contains("articles")) {
      const s = bdd.createObjectStore("articles", { keyPath: "id", autoIncrement: true });
      s.createIndex("par-date", "creeLe");
      s.createIndex("par-auteur-\'etat", ["auteur", "etat"], { multiEntry: false });
      s.createIndex("pleine-texte", "recherche", { multiEntry: true });   // tableau de mots-cles
    }
    if (nouvelle === 2 && !bdd.objectStoreNames.contains("tags")) {
      bdd.createObjectStore("tags", { keyPath: "nom" });
    }
  },
});

// Ecrire
const id = await db.put("articles", { titre: "Apprendre JS", etat: "brouillon", creeLe: Date.now(), recherche: ["apprendre", "js"] });
await db.add("articles", { titre: "unique" });           // echoue si la cle existe
await db.bulkPut("articles", tableauDObjects);
await db.delete("articles", id);
await db.clear("articles");

// Lire
const un = await db.get("articles", id);
const tous = await db.getAll("articles");
const compte = await db.count("articles");
const parDate = await db.getAllFromIndex("articles", "par-date", IDBKeyRange.bound(d1, d2));
const cursor = await db.transaction("articles").store.index("par-date").openCursor(IDBKeyRange.lowerBound(Date.now() - 864e5), "prev");
const page = [];
while (cursor && page.length < 20) { page.push(cursor.value); cursor = await cursor.continue(); }
const cleSuivante = await db.getAllKeys("articles", null, 1);   // pagination par cle, la scalable

// Transaction multi-magasins, atomique
const t = db.transaction(["articles", "tags"], "readwrite");
await Promise.all([
  t.objectStore("articles").put(article),
  t.objectStore("tags").put({ nom: "js", nb: (await t.objectStore("tags").get("js"))?.nb + 1 ?? 1 }),
]);
await t.done;      // commit ; "aborted" en cas d'erreur -> tout est annule

// Recherche pleine texte maison (index multiEntry) :
const mots = requete.toLowerCase().split(/\s+/).filter(Boolean);
const resultat = await db.getAllFromIndex("articles", "pleine-texte", IDBKeyRange.only(mots[0]));
\end{codejs}

\begin{astuce}
Besoin de requ\^etes SQL, de jointures et de gros volumes dans le navigateur~?
\link{SQL.js}{https://sql.js.org/}, \link{wa-sqlite}{https://github.com/rhashimoto/wa-sqlite} ou
\link{TinyBase}{https://tinybase.org/} tournent sur IndexedDB/OPFS. Pour une application locale
lourde (\'editeur, cartographie, comptabilit\'e), SQLite compil\'e en WASM sur OPFS est
une architecture r\'eelle, pas un gadget.
\end{astuce}

\section{Service worker, cache et hors ligne}
\label{sec:pwa}

Un \emph{service worker} est un script qui s'installe dans le navigateur, hors de la page,
et qui peut intercepter les requ\^etes r\'eseau. C'est ce qui rend un site installable, rapide en
2\textsuperscript{e} visite, et utilisable sans r\'eseau.

\begin{codejs}
// sw.js --- a la racine du site (portee = son dossier et au-dessous)
const CACHE = "v3";
const ESSENTIELS = ["/", "/index.html", "/app.js", "/style.css", "/icone-192.png"];

self.addEventListener("install", (e) => {
  e.waitUntil(caches.open(CACHE).then((c) => c.addAll(ESSENTIELS)).then(() => self.skipWaiting()));
});
self.addEventListener("activate", (e) => {
  e.waitUntil(
    caches.keys()
      .then((cles) => Promise.all(cles.filter((k) => k !== CACHE).map((k) => caches.delete(k))))
      .then(() => self.clients.claim())
  );
});
self.addEventListener("fetch", (e) => {
  const req = e.request;
  if (req.method !== "GET" || new URL(req.url).origin !== self.location.origin) return;

  // 1) Resources applicatives : stale-while-revalidate
  if (req.destination === "script" || req.destination === "style" || req.mode === "navigate") {
    e.respondWith(
      caches.match(req).then((cache) => {
        const reseau = fetch(req)
          .then((res) => { const clone = res.clone(); caches.open(CACHE).then((c) => c.put(req, clone)); return res; })
          .catch(() => cache ?? Response.error());
        return cache ?? reseau;
      })
    );
    return;
  }
  // 2) API : network-first avec repli sur le cache (donnees legerement perimees > rien)
  if (new URL(req.url).pathname.startsWith("/api/")) {
    e.respondWith(
      fetch(req)
        .then((res) => { if (res.ok) caches.open(CACHE).then((c) => c.put(req, res.clone())); return res; })
        .catch(() => caches.match(req).then((r) => r ?? new Response(JSON.stringify({ horsLigne: true }), { headers: { "content-type": "application/json" } })))
    );
  }
});
// Recevoir un message de la page pour demander une mise a jour
self.addEventListener("message", (e) => { if (e.data?.type === "verifier") self.skipWaiting(); });
\end{codejs}

\begin{codejs}
// enregistrement depuis la page
if ("serviceWorker" in navigator && location.protocol === "https:") {
  const reg = await navigator.serviceWorker.register("/sw.js", { scope: "/" });
  reg.addEventListener("updatefound", () => {
    const nouveau = reg.installing;
    nouveau?.addEventListener("statechange", () => {
      if (nouveau.state === "installed" && navigator.serviceWorker.controller) {
        notifier("Nouvelle version disponible", { action: { titre: "Recharger", clic: () => nouveau.postMessage({ type: "aktiver" }) } });
      }
    });
  });
}
// Detection d'etat reseau
const enLigne = navigator.onLine;
window.addEventListener("online", reSynchroniser);
window.addEventListener("offline", () => afficherBanniere("Mode hors ligne : tes saisies seront renvoyees plus tard"));
// File d'attente d'arriere-plan officielle (Chromium + Firefox 2024+)
if ("BackgroundSync" in window) { const reg = await navigator.serviceWorker.getRegistration(); await reg.sync.register("renvoi"); }
\end{codejs}

\begin{tabz}{lX}
  \headrow \thd{Strat\'egie de cache} \& \thd{Quand} \\
  \midrule
  Cache-only \& ressources fig\'ees, versionn\'ees dans le nom. \\
  Network-only \& appels non-idempotents, endpoints sensibles. \\
  Network-first, cache en repli \& donn\'ees qui doivent \^etre fra\^iches (API, HTML). \\
  Cache-first, r\'eseau en repli \& assets immuables, images. \\
  Stale-while-revalidate \& le meilleur compromis pour une SPA~: affichage imm\'ediat + mise \`a jour en fond. \\
  Race \& deux miroirs, on garde le premier arriv\'e (langueur r\'eseau \'elev\'ee). \\
\end{tabz}

\etape{Rendre la PWA installable}
\begin{codejson}
{
  "name": "Carnet de Nour",
  "short_name": "Carnet",
  "id": "/?source=pwa",
  "start_url": "/",
  "scope": "/",
  "display": "standalone",
  "display_override": ["window-controls-overlay"],
  "orientation": "any",
  "theme_color": "#0b4f8a",
  "background_color": "#ffffff",
  "description": "Prise de notes rapide, hors ligne, chiffre cote appareil.",
  "icons": [
    { "src": "/icones/icone-192.png", "sizes": "192x192", "type": "image/png", "purpose": "any maskable" },
    { "src": "/icones/icone-512.png", "sizes": "512x512", "type": "image/png" }
  ],
  "shortcuts": [{ "name": "Nouvelle note", "url": "/?nouvelle", "icons": [{ "src": "/icones/note.png", "sizes": "96x96" }] }]
}
\end{codejson}

\begin{terminal}
npm i -D vite-plugin-pwa       # genere manifest, sw, icone, et le cycle de mise a jour
npx @vite/pwa-assets-creator   # genere toutes les tailles d'icone depuis un SVG
npx lighthouse --only-categories=pwa --output=html --output-path=rapport.html http://localhost:4173
\end{terminal}

\section{S\'ecurit\'e et confidentialit\'e du stockage}

\begin{itemize}
  \item Ne chiffre rien \og ~\`a la main~ \fg{} dans le code client en pensant que c'est secret~:
        la clef est dans le bundle. Pour du secret \emph{par appareil}, la WebCrypto API avec une
        clef non extractable (\verb|CryptoKey| \verb|extractable: false|) est correcte~; le reste
        est une d\'eclaration d'intention.
  \item Limite la r\'etention~: purge les cl\'es plus vieilles que X jours, vide le stockage \`a la
        d\'econnexion (\verb|IndexedDB.deleteDatabase|, purge du cache, nettoyage SW).
  \item RGPD~: un stockage local non-transmis reste une donn\'ee personnelle~; documente-le dans la
        politique de confidentialit\'e.
  \item \link{OPFS}{https://developer.mozilla.org/en-US/docs/Web/API/File\_System\_API/Origin\_private\_file\_system}
        est priv\'e \`a l'origine~: plus rapide et plus gros que IndexedDB pour des fichiers.
\end{itemize}

\begin{terminal}
# vider tout ce que le site a garde (DevTools > Application > Storage)
const cachesKeys = await caches.keys(); await Promise.all(cachesKeys.map(k => caches.delete(k)));
indexedDB.databases().then(dbs => dbs.forEach(d => indexedDB.deleteDatabase(d.name)));
localStorage.clear(); sessionStorage.clear();
navigator.serviceWorker.getRegistrations().then(rs => rs.forEach(r => r.unregister()));
location.reload();
\end{terminal}

\begin{exercice}
\begin{enumerate}
  \item Ajoute \`a ta liste de courses une sauvegarde automatique~: apr\`es chaque mutation,
        \'ecrit dans \verb|localStorage| avec un \verb|debounce| de 400~ms, et restaure au
        d\'emarrage en affichant \og ~brouillon restaur\'e~ \fg{}.
  \item \'Ecris une fonction de \textbf{synchronisation}~: un tableau \verb|enAttente| dans
        IndexedDB, renvoi automatique quand \verb|online| revient, avec gestion des doublons
        par identifiant local.
  \item Migre ton stockage de \verb|localStorage| vers IndexedDB~: garde une cl\'e de version,
        et un code qui lit l'ancien format une seule fois puis le supprime.
  \item Ajoute une exportation JSON (t\'el\'echargement du fichier via Blob) et une importation
        valid\'ee par \verb|zod|~: refuse tout fichier malform\'e sans corrompre l'\'etat existant.
    \item Rends ton application installable~: manifest, service worker, mode hors ligne
        complet, et un indicateur visible \og ~derni\`ere synchronisation~ \fg{}.
  \item V\'erifie dans DevTools \$\rightarrow\$ Application \$\rightarrow\$ Storage ce que ton site
        garde r\'eellement~; supprime ce qui ne devrait pas y \^etre.
\end{enumerate}
\end{exercice}

\begin{ressources}
\begin{itemize}
  \item MDN --- Web Storage --- \url{https://developer.mozilla.org/fr/docs/Web/API/Window/localStorage}
  \item MDN --- IndexedDB --- \url{https://developer.mozilla.org/fr/docs/Web/API/IndexedDB\\\_API}
  \item \emph{The Offline Cookbook} (Jake Archibald) --- \url{https://web.dev/articles/offline-cookbook}
  \item idb --- \url{https://github.com/jakearchibald/idb}
  \item PWA --- \url{https://web.dev/learn/pwa/} --- Workbox --- \url{https://developer.chrome.com/docs/workbox}
\end{itemize}
\end{ressources}

\index{IndexedDB}\index{service worker}\index{PWA}\index{cookies}
% ============================================================================
\chapter{Outillage : formatage, analyse, d\'ebogage, structure}
\label{ch:js-outils}
% ============================================================================

\begin{objectif}
\begin{itemize}
  \item Formater et analyser ton code automatiquement, sans y penser, sur chaque sauvegarde.
  \item Configurer le d\'ebogueur pas \`a pas dans VS~Code (Node et navigateur).
  \item Choisir une structure de dossiers qui tient \`a 20 000 lignes comme \`a 200.
  \item Conna\^itre les outils officiels (Vite, ESLint~9 flat config, Biome) et leurs liens.
\end{itemize}
\end{objectif}

\section{Formater~: Prettier ou Biome}

Un formateur n'est pas une opinion~: c'est la fin des discussions d'\'equipe sur les espaces.
Ne configure que le minimum, laisse l'outil d\'ecider.

\begin{terminal}
npm i -D prettier prettier-plugin-tailwindcss
npx prettier --write .
npm i -D @biomejs/biome          # alternative Rust : formatage + linting en un binaire, tres rapide
npx biome init
\end{terminal}

\begin{codejson}
{
  "$schema": "https://json.schemastore.org/prettierrc",
  "printWidth": 100,
  "semi": true,
  "singleQuote": false,
  "trailingComma": "all",
  "arrowParens": "always",
  "endOfLine": "lf",
  "plugins": ["prettier-plugin-tailwindcss"],
  "overrides": [
    { "files": ["*.md"], "options": { "proseWrap": "preserve" } },
    { "files": ["*.json"], "options": { "printWidth": 200 } }
  ]
}
\end{codejson}

\begin{note}
Le point de friction classique~: un formateur et ESLint qui se battent (l'un remet les guillemets
simples, l'autre les doubles). La r\`egle est~: \textbf{Prettier d\'ecide du style, ESLint ne d\'ecide
que de la correction}. Pour cela on installe
\link{eslint-config-prettier}{https://github.com/prettier/eslint-config-prettier} en derni\`ere
position et on met \verb|formatOnSave|.
\end{note}

\section{Analyser~: ESLint~9 (config plate)}

Le fichier est \fichier{eslint.config.js}~; l'ancien \fichier{.eslintrc.json} est abandonn\'e.

\begin{codejs}
// eslint.config.js
import js from "@eslint/js";
import globals from "globals";
import prettierRecommended from "eslint-plugin-prettier/recommended";
import reactHooks from "eslint-plugin-react-hooks";

export default [
  { ignores: ["dist/**", "node_modules/**", "coverage/**", "public/sw.js"] },
  js.configs.recommended,
  {
    files: ["**/*.{js,jsx,mjs,cjs}"],
    languageOptions: {
      ecmaVersion: "latest",
      sourceType: "module",
      globals: { ...globals.browser, ...globals.node },
    },
    rules: {
      "no-unused-vars": ["error", { argsIgnorePattern: "^_", varsIgnorePattern: "^_" }],
      "no-console": ["warn", { allow: ["warn", "error"] }],
      eqeqeq: ["error", "always", { null: "ignore" }],
      "prefer-const": "error",
      "no-var": "error",
      "object-shorthand": "error",
      curly: ["error", "all"],
      "no-param-reassign": ["error", { props: false }],
      complexity: ["warn", 12],
      "max-depth": ["warn", 4],
      "no-nested-ternary": "warn",
      "require-await": "error",
      "no-promise-executor-return": "error",
      "no-async-promise-executor": "error",
      "consistent-return": "off",
      "import/order": "off",
    },
  },
  { files: ["src/ui/**"], plugins: { "react-hooks": reactHooks }, rules: reactHooks.configs.recommended.rules },
  prettierRecommended,
];
\end{codejs}

\begin{terminal}
npx eslint . --fix
npx eslint . --max-warnings=0 --format=stylish        # en CI : echoue au premier avertissement
npx eslint src --format=json > rapport.json
# Regles de securite cote node : npm i -D @eslint-community/eslint-plugin-security
# Analyse de la surete de tes regex : npm i -D eslint-plugin-security + eslint-plugin-redos
\end{terminal}

\subsection{Les autres v\'erificateurs \`a brancher une fois}

\begin{tabz}{lXl}
  \headrow \thd{Outil} \& \thd{Ce qu'il trouve} \& \thd{Lancer} \\
  \midrule
  \link{TypeScript}{https://www.typescriptlang.org/} (\cmd{tsc --noEmit}) \& types, propri\'et\'es inexistantes, appels faux \& \cmd{npm run typecheck} \\
  \link{Knip}{https://knip.dev/} \& fichiers, exports et d\'ependances inutilis\'es \& \cmd{npx knip} \\
  \link{Depcheck}{https://github.com/dpinies/depcheck} \& d\'ependances du \verb|package.json| jamais import\'ees \& \cmd{npx depcheck} \\
  \link{madge}{https://github.com/pahen/madge} \& d\'ependances circulaires entre modules \& \cmd{npx madge --circular src} \\
  \link{source-map-explorer}{https://github.com/danvk/source-map-explorer} \& qui occupe l'espace dans ton bundle \& \cmd{npx source-map-explorer dist/\/*.js} \\
  \link{rollup-plugin-visualizer}{https://github.com/btd/rollup-plugin-visualizer} \& idem, int\'egr\'e au build \& \cmd{vite build --mode analyse} \\
  \link{npm-check-updates}{https://github.com/tjun/ncu} \& mises \`a jour en conservant semver \& \cmd{npx npm-check-updates -u} \\
  \link{Knip/Snyk/Socket}{https://socket.dev/} \& paquets suspects, typosquatting \& \cmd{npx socket scan .} \\
  \link{shellcheck}{https://www.shellcheck.net/} \& bugs dans tes scripts bash \& \cmd{shellcheck scripts/*} \\
\end{tabz}

\section{D\'eboguer pour de vrai}

\etape{Node, dans VS~Code}
\begin{codejson}
{
  "version": "0.2.0",
  "configurations": [
    {
      "type": "node",
      "request": "launch",
      "name": "API (avec .env)",
      "program": "${workspaceFolder}/src/server.js",
      "envFile": "${workspaceFolder}/.env",
      "skipFiles": ["<node_internals>/**"],
      "console": "integratedTerminal",
      "restart": true
    },
    {
      "type": "node",
      "request": "attach",
      "name": "Brancher sur node --inspect",
      "port": 9229
    },
    {
      "type": "chrome",
      "request": "launch",
      "name": "Chrome + source maps",
      "url": "http://localhost:5173",
      "webRoot": "${workspaceFolder}/src"
    }
  ]
}
\end{codejson}

\begin{terminal}
node --inspect-brk src/server.js        # s'arrete sur la premiere ligne, en attente du debogueur
node --inspect src/server.js            # tourne et ecoute sur 9229
\end{terminal}

\etape{Les techniques qui remplacent 90 \% des \texttt{console.log}}
\begin{itemize}
  \item \textbf{Point d'arr\^et conditionnel}~: clic droit sur le num\'ero de ligne \$\rightarrow\$ \og ~Breakpoint
        conditionnel~ \fg{}, \verb|i === 4231| ou \verb|user.id === id|.
  \item \textbf{Point d'arr\^et de log}~: on imprime une expression sans casser l'ex\'ecution.
  \item \textbf{Point d'arr\^et sur exception}~: l'ic\^one \og ~pause on caught exceptions~ \fg{} dans le panneau
        Call Stack~: le d\'ebogueur s'arr\^ete \emph{\`a la source} de l'erreur, pas l\`a o\`u on l'affiche.
  \item \textbf{Console}~: \verb|copy(x)| (copie dans le presse-papiers), \verb|\$\_| (derni\`ere
        expression), \verb|\$0| (\'el\'ement s\'electionn\'e dans Elements), \verb|\$\$("sel")| (\texttt{querySelectorAll}), \verb|monitor(x)|.
  \item \textbf{Live expressions}~: dans le panneau Watch, une expression r\'e\'evalu\'ee \`a chaque pause.
  \item \textbf{Les \og ~timeline~ \fg{} breakpoints}~: DevTools \$\rightarrow\$ Event Listener Breakpoints pour
        n'importe quel \'ev\'enement (utile quand \og ~je ne sais pas quel handler fait \c{c}a~ \fg{}).
\end{itemize}

\section{Structurer un projet qui vieillit bien}

\etape{Petit projet (jusqu'\`a ~2 000 lignes)~: par type}
\begin{codeplain}
projet/
  public/          index.html, assets statiques
  src/
    main.js        point d'entree (20 lignes max)
    ui/            un fichier par vue/composant
    lib/           fonctions pures, sans DOM, sans reseau  <-- testables gratuitement
    api/           toute sortie reseau, un seul endroit
    state.js       etat global et ses transitions
    styles/        css
  tests/           miroir de src/
\end{codeplain}

\etape{Projet moyen/gros~: par domaine (feature-first)}
\begin{codeplain}
src/
  noyau/            config, erreurs, journalisation, utilitaires transverses
  articles/         une fonctionnalite = un dossier autonome
    modele.js       regles metier (aucune dependance web)
    depot.js        acces aux donnees
    service.js      cas d'utilisation ("publier article", "facturer")
    controleur.js   entrees HTTP / CLI
    vues/           composants d'interface
    index.js        export public (le "package" interne)
  reseau/
  tests/            unitaires a cote du code, integration a part
\end{codeplain}

Les trois r\`egles qui comptent~: (1) la logique m\'etier ne conna\^it ni le DOM ni HTTP~;
(2) un module ne d\'epend que des \verb|index.js| des autres~; (3) les d\'ependances vont des
couches techniques vers le m\'etier, jamais l'inverse.

\subsection{\texttt{.editorconfig}, \texttt{.nvmrc}, \texttt{docker-compose} de d\'eveloppement}

\begin{codeplain}
# .nvmrc --- la version exacte de Node du projet
22.12.0
# les outils la lisent : nvm use, fnm, mise, asdf, devcontainers

# .dockerignore
node_modules
.git
dist
coverage
*.log
\end{codeplain}

\begin{codeyml}
# docker-compose.yml --- tout ce dont une equipe a besoin, en une commande
services:
  api:
    build: { context: ., dockerfile: Dockerfile.dev }
    command: npm run dev
    environment:
      DATABASE_URL: postgres://nour:nour@db:5432/carnet
      JWT_SECRET: change-moi-en-production
    ports: ["3000:3000"]
    volumes: [".:/app", "/app/node_modules"]
    depends_on: { db: { condition: service_healthy } }
  db:
    image: postgres:16-alpine
    environment: { POSTGRES_USER: nour, POSTGRES_PASSWORD: nour, POSTGRES_DB: carnet }
    volumes: ["pg:/var/lib/postgresql/data"]
    healthcheck: { test: ["CMD-SHELL", "pg_isready -U nour"], interval: 5s, retries: 10 }
  mailpit:
    image: axllent/mailpit        # boite de reception locale : rien n'est reellement envoye
    ports: ["8025:8025"]
volumes: { pg: {} }
\end{codeyml}

\begin{note}
Le volume \verb|/app/node\_modules| anonyme \'evite que le \verb|node\_modules| de ta machine
(Windows/macOS, binaire diff\'erent) n'\'ecrase celui du conteneur. C'est l'erreur de Docker la
plus fr\'equente en d\'eveloppement Node.
\end{note}

\begin{exercice}
\begin{enumerate}
  \item Installe Prettier et ESLint sur ton dernier projet ; mets \verb|formatOnSave|~; lance
        \verb|eslint . --fix| et lis chaque correction avant d'accepter.
  \item Ajoute \verb|husky| + \verb|lint-staged| pour que le commit refuse un fichier non format\'e.
  \item Configure un point d'arr\^et conditionnel sur une boucle de 100 000 tours et observe le
        panneau Call Stack. Ensuite, utilise \og ~logpoint~ \fg{} sur la m\^eme ligne et compare.
  \item Ex\'ecute \verb|npx knip| et \verb|npx depcheck| sur un de tes projets : supprime ce qui
        est mort. Note le poids gagn\'e sur le bundle.
  \item Restructure ton projet de la semaine pr\'ec\'edente en \emph{feature folders}, avec une
        couche \verb|lib/| sans aucune d\'ependance au DOM. V\'erifie que tu peux tester \verb|lib/|
        sans navigateur.
  \item \'Ecris un \verb|docker-compose.yml| qui d\'emarre l'API + Postgres + Mailpit, et un
        \verb|Makefile| avec les cibles \verb|up|, \verb|down|, \verb|migrate|, \verb|test|, \verb|logs|.
\end{enumerate}
\end{exercice}

\begin{ressources}
\begin{itemize}
  \item Prettier options --- \url{https://prettier.io/docs/configuration}
  \item ESLint 9 configuration --- \url{https://eslint.org/docs/latest/use/configure/configuration-files}
  \item Biome --- \url{https://biomejs.dev/fr/}
  \item VS Code debugging (Node) --- \url{https://code.visualstudio.com/docs/nodejs/nodejs-debugging}
  \item DevTools --- \url{https://developer.chrome.com/docs/devtools/}
\end{itemize}
\end{ressources}

\index{Prettier}\index{ESLint}\index{debogage}\index{structure de projet}
% ============================================================================
\chapter{Interface moderne : Vite, composants, React/Vue/Svelte, Astro}
\label{ch:js-frameworks}
% ============================================================================

\begin{objectif}
\begin{itemize}
  \item D\'emarrer un projet avec Vite et comprendre ce qu'un bundler fait pour toi.
  \item Saisir le mod\`ele \og ~\'etat \$\rightarrow\$ rendu~ \fg{} commun aux frameworks, et les hooks React.
  \item Choisir un cadre en connaissance de cause, sans religion.
  \item G\'erer \'etat, routes, donn\'ees serveur et styles dans une vraie application.
  \item Savoir quand \emph{ne pas} mettre de framework.
\end{itemize}
\end{objectif}

\section{Vite~: le socle}

\link{Vite}{https://vitejs.dev/} (7\textsuperscript{e}~majore en 2025, cr\'e\'e par Evan~You) sert tes
modules nativement en ESM pendant le d\'eveloppement (d\'emarrage quasi instantan\'e, HMR cibl\'e) et
utilise \link{Rollup}{https://rollupjs.org/} pour construire la version de production
(minification, d\'ecoupage de code, empreintes de fichiers).

\begin{terminal}
npm create vite@latest site -- --template vanilla        # ou react, vue, svelte, preact, lit...
cd site
npm install
npm run dev              # http://localhost:5173
npm run build && npm run preview
\end{terminal}

\begin{codejs}
// vite.config.js --- les reglages qui reviennent vraiment
import { defineConfig } from "vite";
import react from "@vitejs/plugin-react";

export default defineConfig({
  plugins: [react()],
  resolve: { alias: { "@": new URL("./src", import.meta.url).pathname } },
  server: { port: 5173, host: true, proxy: { "/api": "http://localhost:8000" } },
  build: {
    sourcemap: true,
    target: "es2022",
    rollupOptions: { output: { manualChunks: (id) => (id.includes("node_modules") ? "deps" : undefined) } },
  },
  define: { __APP_VERSION__: JSON.stringify(process.env.npm_package_version) },
});
\end{codejs}

Les variables d'environnement expos\'ees au client doivent \^etre pr\'efix\'ees \verb|VITE\_|~:
\verb|import.meta.env.VITE\_API\_URL|. Tout le reste du fichier \verb|.env| reste c\^ot\'e build~:
c'est intentionnel, ne cherche pas \`a le contourner.

\section{Le mod\`ele commun \`a tous les frameworks}

\begin{enumerate}
  \item Tu d\'ecris l'\textbf{\'etat} (donn\'ees) et une \textbf{fonction de rendu}
        (\verb|\'etat -> HTML|).
  \item Tu changes l'\'etat avec les moyens du framework (\verb|setState|, \verb|\$set|, signal,
        assignation r\'eactive).
  \item Le framework \textbf{diff} l'ancien et le nouveau rendu, applique le minimum
        d'op\'erations DOM (la \og ~r\'econciliation~ \fg{}), puis compose.
\end{enumerate}

Tout le reste (hooks, signaux, stores, effets) n'est que de la broderie autour de ce mod\`ele.
C'est pour cela qu'apprendre un framework en rend les autres rapides.

\subsection{React --- le cas le plus demand\'e}

\begin{codejs}
// src/Compteur.jsx --- lecture obligatoire : un composant est une fonction pure de l'etat
import { useState, useEffect, useMemo, useCallback, useRef, useReducer } from "react";

const initial = () => JSON.parse(localStorage.getItem("courses") ?? "[]");

export function ListeCourses() {
  const [articles, setArticles] = useState(initial);          // initialise une seule fois
  const [filtre, setFiltre] = useState("tous");
  const [brouillon, setBrouillon] = useState("");
  const champ = useRef(null);                                  // reference a un noeud DOM, sans re-rendu

  useEffect(() => {                                            // effet = synchronisation avec l'exterieur
    localStorage.setItem("courses", JSON.stringify(articles));
  }, [articles]);

  useEffect(() => {                                            // ecouteur global + nettoyage
    const onKey = (e) => e.key === "/" && champ.current?.focus();
    window.addEventListener("keydown", onKey);
    return () => window.removeEventListener("keydown", onKey);
  }, []);                                                        // [] : une fois, au montage

  const restants = useMemo(() => articles.filter((a) => !a.faite), [articles]);
  const total = useMemo(() => restants.reduce((s, a) => s + a.qte, 0), [restants]);

  const ajouter = useCallback(() => {
    const nom = brouillon.trim();
    if (!nom) return;
    setArticles((liste) => [{ id: crypto.randomUUID(), nom, qte: 1, faite: false }, ...liste]);
    setBrouillon("");
  }, [brouillon]);

  const basculer = (id) =>
    setArticles((liste) => liste.map((a) => (a.id === id ? { ...a, faite: !a.faite } : a)));

  return (
    <main>
      <h1>Courses <small>({total})</small></h1>
      <form onSubmit={(e) => { e.preventDefault(); ajouter(); }}>
        <input ref={champ} value={brouillon} onChange={(e) => setBrouillon(e.target.value)}
               placeholder="Ajouter (touche /)" />
      </form>
      <ul>
        {restants.map((a) => (
          <li key={a.id}>
            <label>
              <input type="checkbox" checked={a.faite} onChange={() => basculer(a.id)} />
              {a.nom}
            </label>
          </li>
        ))}
      </ul>
    </main>
  );
}
\end{codejs}

\etape{Les quatre r\`egles des hooks, sans exception}
\begin{enumerate}
  \item Un hook s'appelle \verb|useXxx|, et uniquement au \textbf{niveau sup\'erieur} d'un
        composant ou d'un hook~: jamais dans une condition, une boucle, un rappel.
  \item Un \verb|useEffect| doit d\'eclarer ses d\'ependances \textbf{compl\`etes}
        (active \verb|eslint-plugin-react-hooks|, et le preset \og ~recommended-latest~ \fg{}).
  \item Une mise \`a jour d'\'etat doit produire une \textbf{nouvelle r\'ef\'erence}
        (\verb|set(x => ({ ...x })), pas \verb|set(state.prop = 1)|).
  \item Si un effet sert \`a \og ~d\'eriver~ \fg{} une valeur, ce n'est pas un effet~: c'est un calcul
        pendant le rendu (\verb|useMemo| ou simple const).
\end{enumerate}

\subsection{Vue~3 (composition API)}

\begin{codeplain}
<!-- ArticleList.vue -->
<script setup>
import { ref, computed, watchEffect } from "vue";

const articles = ref(JSON.parse(localStorage.getItem("courses") ?? "[]"));
const brouillon = ref("");
const restants = computed(() => articles.value.filter((a) => !a.faite));

function ajouter() {
  const nom = brouillon.value.trim();
  if (!nom) return;
  articles.value.unshift({ id: crypto.randomUUID(), nom, qte: 1, faite: false });
  brouillon.value = "";
}
watchEffect(() => localStorage.setItem("courses", JSON.stringify(articles.value)));
</script>

<template>
  <input v-model="brouillon" @keyup.enter="ajouter" placeholder="Article" />
  <ul><li v-for="a in restants" :key="a.id">{{ a.nom }}</li></ul>
</template>
\end{codeplain}

\subsection{Svelte~5 (runes)~: le moins de code du march\'e}

\begin{codeplain}
<script>
  let articles = $state(JSON.parse(localStorage.getItem("courses") ?? "[]"));
  let brouillon = $state("");
  const restants = $derived(articles.filter((a) => !a.faite));
  $effect(() => localStorage.setItem("courses", JSON.stringify(articles)));

  function ajouter() {
    const nom = brouillon.trim();
    if (!nom) return;
    articles = [{ id: crypto.randomUUID(), nom, qte: 1, faite: false }, ...articles];
    brouillon = "";
  }
</script>

<input bind:value={brouillon} on:keydown={(e) => e.key === "Enter" && ajouter()} />
<ul>{#each restants as a (a.id)}<li>{a.nom}</li>{/each}</ul>
\end{codeplain}

\section{Panorama honn\^ete}

\begin{tabz}{lXlX}
  \headrow \thd{Cadre} \& \thd{Id\'ee forte} \& \thd{Taille (gzipp\'ee, app. vide)} \& \thd{Doc} \\
  \midrule
  \link{React~19}{https://react.dev/} \& le plus d'emplois, \'ecosyst\`eme immense, rendu par fonction \& \$\approx\$45~Ko \& \url{https://react.dev/learn} \\
  \link{Vue~3.5}{https://fr.vuejs.org/} \& la meilleure mont\'ee en charge du d\'ebutant au projet moyen, templates clairs \& \$\approx\$35~Ko \& \url{https://fr.vuejs.org/guide/introduction.html} \\
  \link{Svelte~5 / SvelteKit}{https://svelte.dev/} \& compilateur, peu de code, performances par d\'efaut \& \$\approx\$7~Ko \& \url{https://learn.svelte.dev/} \\
  \link{Solid}{https://www.solidjs.com/} \& signaux fins, syntaxe JSX, rapidit\'e maximale \& \$\approx\$8~Ko \& \url{https://docs.solidjs.com/} \\
  \link{Qwik}{https://qwik.dev/} \& reprise d'hydratation~: ne charge que ce que l'utilisateur touche \& minuscule \& \url{https://qwik.dev/learn/} \\
  \link{Angular~19+}{https://angular.dev/} \& tout int\'egr\'e (routes, HTTP, formulaires, DI), signals, standalone \& lourd mais complet \& \url{https://angular.dev/tutorials} \\
  \link{Astro~5}{https://astro.build/} \& sites centr\'es contenu~: z\'ero JS par d\'efaut, \^ilots d'interactive \& 0~Ko \& \url{https://docs.astro.build/fr/getting-started/} \\
  \link{HTMX}{https://htmx.org/} \& l'interactive dans les attributs HTML, le serveur renvoie du HTML \& 14~Ko \& \url{https://htmx.org/docs/} \\
  \link{Lit}{https://lit.dev/} \& web components natifs et durables, multi-cadres \& \$\approx\$17~Ko \& \url{https://lit.dev/docs/} \\
  \link{Alpine.js}{https://alpinejs.dev/} \& du \og ~jQuery propre~ \fg{} pour les petites pages \& \$\approx\$15~Ko \& \url{https://alpinejs.dev/start-here} \\
  \link{jQuery}{https://jquery.com/} \& \`a conna\^itre pour maintenir l'existant, pas pour d\'emarrer \& 30~Ko \& \url{https://api.jquery.com/} \\
\end{tabz}

\subsection{Les briques transverses}

\begin{tabz}{lXl}
  \headrow \thd{Besoin} \& \thd{Choix raisonnable en 2026} \& \thd{Lien} \\
  \midrule
  Routage SPA \& TanStack Router (typ\'e) ou react-router v7 \& \url{https://tanstack.com/router} \\
  Donn\'ees serveur \& TanStack Query~: cache, retries, invalidation, \'etats \& \url{https://tanstack.com/query} \\
  \'Etat client \& Zustand (l\'eger) ou le signal natif ; Redux Toolkit si grosse \'equipe \& \url{https://zustand.docs.pmnd.rs/} \\
  Formulaires \& React Hook Form + zod (validation partag\'ee client/serveur) \& \url{https://react-hook-form.com/} \\
  Styles \& CSS natif (variables, \verb|:has|, \verb|@layer|, \verb|container| queries) ou Tailwind~4 \& \url{https://tailwindcss.com/} \\
  Composants UI \& shadcn/ui (copie le code, ne d\'epend pas d'une lib) \& \url{https://ui.shadcn.com/} \\
  Animation \& \verb|element.animate()|, Motion (framer-motion), GSAP \& \url{https://motion.dev/} \\
  Graphes/donn\'ees \& \link{D3}{https://d3js.org/}, \link{ECharts}{https://echarts.apache.org/}, \link{Chart.js}{https://www.chartjs.org/} \& \\
\end{tabz}

\section{Quand ne pas mettre de framework}

\begin{itemize}
  \item Site vitrine, blog, documentation, landing page~: \textbf{Astro} (ou du HTML + un peu de JS).
  \item Page avec trois interactions~: HTML + 80 lignes de JS natif + \link{Alpine}{https://alpinejs.dev/}.
  \item Backend qui rend du HTML (PHP, Java, Rails, Laravel)~: \link{HTMX}{https://htmx.org/}
        ou \link{Hotwire/Turbo}{https://turbo.hotwired.dev/}~: tu gardes la simplicit\'e du serveur,
        tu gagnes le confort du SPA. C'est une tendance lourde et saine.
  \item Composant autonome r\'eutilisable dans n'importe quelle page~: un
        \link{Web Component}{https://developer.mozilla.org/fr/docs/Web/API/Web\_components} natif.
\end{itemize}

\begin{note}
Un framework se choisit sur trois crit\`eres~: la taille de l'\'equipe, la part d'interactive r\'eelle
de l'interface, et la disponibilit\'e des comp\'etences sur le march\'e local. Pour un portfolio,
React ou Vue maximisent les chances d'entretien~; pour un projet solo qui doit sortir vite,
Svelte ou Astro gagnent.
\end{note}

\begin{exercice}
\begin{enumerate}
  \item Initialise trois projets du m\^eme contenu (liste de courses)~: vanilla, Vue, React.
        Compare le nombre de lignes, le poids du bundle (\verb|dist/|) et le temps de
        d\'eveloppement~: \'ecris ton avis en dix lignes dans un \verb|ANALYSE.md|.
  \item Sur le projet React~: remplace le \verb|useState| local par un store \verb|zustand|,
        et ajoute une persistance via le middleware \verb|persist|.
  \item Branche \link{TanStack~Query}{https://tanstack.com/query} sur une API r\'eelle~: g\`ere
        \'etats \og ~chargement / erreur / vide~ \fg{}, mise en cache, invalidation apr\`es mutation.
  \item Ajoute le routage~: routes \verb|/articles|, \verb|/articles/:id|, \verb|404|,
        et un chargement diff\'er\'e (\verb|React.lazy| / \verb|defineAsyncComponent|) de la page d\'etail.
  \item Migre un composant React existant vers un Web Component natif (\verb|customElements.define|)
        et utilise-le dans une page HTML sans framework.
\end{enumerate}
\end{exercice}

\begin{ressources}
\begin{itemize}
  \item Vite --- \url{https://vite.dev/guide/}
  \item React (nouvelle doc, excellente) --- \url{https://react.dev/learn}
  \item Vue --- \url{https://fr.vuejs.org/guide/extras/composition-api-faq.html}
  \item Svelte --- \url{https://svelte.dev/tutorial/} --- Astro --- \url{https://docs.astro.build/fr/core-concepts/overview/}
  \item State of JS (sondage annuel, pour choisir sans se faire d'opinion fausse) --- \url{https://2024.stateofjs.com/}
\end{itemize}
\end{ressources}

\index{Vite}\index{React}\index{Vue}\index{Svelte}\index{framework}
% ============================================================================
\chapter{TypeScript~: rendre les erreurs visibles avant l'ex\'ecution}
\label{ch:js-ts}
% ============================================================================

\begin{objectif}
\begin{itemize}
  \item Comprendre ce que TypeScript fait (et ne fait pas)~: v\'erification statique, pas de magie runtime.
  \item Typager unions, narrowing, g\'en\'eriques, types utilitaires, discriminations.
  \item Configurer \texttt{tsconfig.json} sans recopier un myst\`ere~: compr\'ehension de chaque drapeau.
  \item Migrer un projet JavaScript existant progressivement, sans se noyer.
  \item Typ\'e des API r\'eelles~: r\'eponses, sch\'emas partag\'es client/serveur, erreurs.
\end{itemize}
\end{objectif}

\section{Le principe en dix lignes}

\begin{codejs}
// entr\'ee.ts
function saluer(qui: string): string {
  return `Bonjour ${qui}`;
}
saluer(42);              // erreur du COMPILATEUR : Argument of type 'number' is not assignable to 'string'

// sortie (tsc) : les annotations disparaissent, le JS obtenu est identique a celui qu'on aurait ecrit a la main
function saluer(qui) { return `Bonjour ${qui}`; }
\end{codejs}

TypeScript est un \textbf{v\'erificateur} qui lit tes annotations, pas un moteur qui les applique.
Cons\'equence capitale~: \`a l'ex\'ecution, aucune protection~! Si une donn\'ee vient de l'ext\'erieur
(JSON, formulaire, base de donn\'ees), le type d\'eclar\'e est une \emph{affirmation}. D'o\`u l'habitude
de valider aux fronti\`eres avec \link{zod}{https://zod.dev/} ou \link{ArkType}{https://arktype.io/}.

\begin{codejs}
import { z } from "zod";

const Article = z.object({
  id: z.number().int().positive(),
  titre: z.string().min(1).max(160),
  tags: z.array(z.string()).default([]),
  publieLe: z.iso.datetime().optional(),
});
type Article = z.infer<typeof Article>;          // le type est DERIVE du sch\'ema : source unique de verite

export async function lireArticle(id: number): Promise<Article> {
  const brut = await fetch(`/api/articles/${id}`).then((r) => r.json());
  return Article.parse(brut);                     // erreur nette ici, pas 3 ecrans plus loin
}
\end{codejs}

\section{Le syst\`eme de types, utilement}

\begin{codejs}
// --- Bases
let age: number = 30;
const user = { nom: "Nour", age: 30 } as const;          // litteraux figes : { readonly nom: "Nour" }
type Identifiant = string | number;
type Statut = "brouillon" | "publie" | "archive";         // union de litteraux : l'enum du JS moderne
const ROLE = { Admin: "admin", Editeur: "editeur" } as const;
type Role = (typeof ROLE)[keyof typeof ROLE];             // "admin" | "editeur"

// --- Fonctions
type Gestionnaire<T = Event> = (e: T) => void | Promise<void>;
function envoyer<T>(donnees: T, options?: { signal?: AbortSignal }): Promise<Response> { /* ... */ }
declare function sur<E extends Statut>(s: E, cb: (a: ArticleDe<E>) => void): void;

// --- Objets
interface Article {
  readonly id: number;                 // non reassignable
  titre: string;
  sousTitre?: string;                   // facultatif : string | undefined
  tags: readonly string[];              // lecture seule
  [index: string]: unknown;             // cles supplementaires acceptees
  (prefixe?: string): string;           // interface callable (rare)
}
type RecordArticle = Record<Statut, Article[]>;

// --- Unions et etroitesse de type (narrowing) --- LE coeur du TS utile
type Rep<Ok> = { ok: true; data: Ok } | { ok: false; error: string };
async function api<T>(url: string): Promise<Rep<T>> {
  const res = await fetch(url);
  if (!res.ok) return { ok: false, error: `HTTP ${res.status}` };
  return { ok: true, data: (await res.json()) as T };
}
const r = await api<Article[]>("/api/articles");
if (r.ok) r.data.map((a) => a.titre);      // TS sait que data existe ici
else console.error(r.error);                 // et que error existe ici

function decrire(v: string | number | boolean | null | undefined | Date) {
  if (v === null || v === undefined) return "rien";
  if (typeof v === "string") return v.toUpperCase();       // v : string
  if (typeof v === "number") return v.toFixed(2);          // v : number
  if (v instanceof Date) return v.toISOString();            // v : Date
  return String(v);                                          // boolean
}

// Predicats de type : la facon propre de factoriser un narrowing
function estArticle(v: unknown): v is Article {
  return typeof v === "object" && v !== null && "id" in v && "titre" in v && typeof (v as Article).id === "number";
}
const bruts: unknown[] = [1, { id: 2, titre: "x" }, null];
const vrais = bruts.filter(estArticle);            // Article[] -- pas unknown[]

// Exhaustivite : le "never" devient une alarme quand un cas est oublie
function libelle(s: Statut): string {
  switch (s) {
    case "brouillon": return "En cours";
    case "publie": return "En ligne";
    case "archive": return "Archive";
    default: { const _exhaustif: never = s; return _exhaustif; }   // erreur si on ajoute un statut
  }
}
\end{codejs}

\subsection{G\'en\'eriques, sans en abuser}

\begin{codejs}
function dernier<T>(tab: T[]): T | undefined { return tab.at(-1); }
function groupePar<T, K extends string | number>(tab: T[], clef: (t: T) => K): Record<K, T[]> {
  return tab.reduce((acc, item) => {
    const k = clef(item);
    (acc[k] ??= []).push(item);
    return acc;
  }, {} as Record<K, T[]>);
}
groupePar(articles, (a) => a.statut);        // { brouillon: Article[], publie: Article[] }

// Contraintes utiles
function proprietes<T, K extends keyof T>(o: T, cles: K[]): Pick<T, K>[] { return cles.map((c) => ({ [c]: o[c] } as Pick<T, K>)); }
function creer<T extends { id: string }>(elements: T[]): Map<string, T> { return new Map(elements.map((e) => [e.id, e])); }
function clone<T extends object>(o: T): T { return structuredClone(o); }

// Types conditionnels + infere : a lire, rarement a ecrire
type ElementDe<T> = T extends (infer U)[] ? U : never;
type Retours<T> = T extends (...a: never[]) => Promise<infer R> ? R : never;
type Attendu = ElementDe<string[]>;         // string
type Resultat = Retours<typeof fetch>;       // Response

// Mapped types : rendre tout facultatif a un niveau arbitraire
type ProfondementPartiel<T> = { [K in keyof T]?: T[K] extends object ? ProfondementPartiel<T[K]> : T[K] };
type Template = `data-${Lowercase<string>}`;   // types de chaines par motif
\end{codejs}

\section{Les types utilitaires \`a conna\^itre}

\begin{codejs}
type A = { id: number; titre: string; sousTitre?: string; tags: string[] };
type U1 = Partial<A>;                    // tout facultatif (ex : objet d'un PATCH)
type U2 = Required<A>;                   // tout obligatoire
type U3 = Readonly<A>;                   // rien d'assignable
type U4 = Pick<A, "id" | "titre">;       // sous-ensemble
type U5 = Omit<A, "tags">;               // retrait
type U6 = Exclude<A["titre"], never>;     // sur des unions : A | B, retire A
type U7 = Extract<A["titre"], string>;
type U8 = NonNullable<A["sousTitre"]>;    // retire null | undefined
type U9 = ReturnType<typeof fetch>;       // Promise<Response>
type U10 = Parameters<typeof fetch>;       // [input: RequestInfo | URL, init?: RequestInit]
type U11 = Awaited<ReturnType<typeof fetch>>;   // Response
type U12 = A["tags"][number];             // string  (element d'un tableau)
type U13 = Record<string, A>;             // dictionnaire
type U14 = { [K in keyof A as `set${Capitalize<K & string>}`]: (v: A[K]) => void };
// -> { setId: (v: number) => void; setTitre: (v: string) => void; ... }
\end{codejs}

\section{\texttt{tsconfig.json}--- chaque drapeau expliqu\'e}

\begin{codejson}
{
  "compilerOptions": {
    "target": "ES2023",                       // syntaxe generee (le moteur cible le plus vieux soutenu)
    "lib": ["ES2023", "DOM", "DOM.Iterable"], // API declarees : DOM retire = plus de document
    "module": "ESNext",                        // format des imports emits (Node : "NodeNext")
    "moduleResolution": "Bundler",               // resolution a la Vite/esbuild ; "NodeNext" pour Node pur
    "verbatimModuleSyntax": true,                // n'enleve pas "import type" : explicite, evite les surprises
    "isolatedModules": true,                     // chaque fichier compilable seul (condition des bundlers)
    "resolveJsonModule": true,
    "allowImportingTsExtensions": true,           // imports avec .ts (Vite/tsx) -- exige noEmit
    "noEmit": true,                                // tsc verifie seulement, le bundler produit le JS
    "strict": true,                                // ACTIVE TOUT ce qui suit -- ne pas negocier
    "noUncheckedIndexedAccess": true,              // tab[i] : T | undefined  (la verite du runtime)
    "exactOptionalPropertyTypes": true,            // { a?: string } n'accepte plus a: undefined
    "noImplicitOverride": true,                     // override doit etre dit
    "noFallthroughCasesInSwitch": true,
    "noUnusedLocals": true, "noUnusedParameters": true,
    "allowJs": false,                               // true pour migrer un projet JS progressivement
    "checkJs": false,                               // true : typer meme les fichiers .js
    "declaration": true, "declarationMap": true, "sourceMap": true,
    "outDir": "dist",
    "skipLibCheck": true,                             // ne pas rev\'erifier les .d.ts des paquets
    "forceConsistentCasingInFileNames": true,          // Windows te remerciera plus tard
    "types": ["node", "vite/client"]
  },
  "include": ["src", "tests", "vite.config.ts"],
  "exclude": ["dist", "node_modules", "coverage"]
}
\end{codejson}

\begin{astuce}
Les deux drapeaux qui rapportent le plus apr\`es \verb|strict|~:
\texttt{noUncheckedIndexedAccess} (transforme une cat\'egorie enti\`ere de \verb|undefined is not an
object| en erreur de compilation) et \texttt{exactOptionalPropertyTypes}. Active-les totard,
tant que le projet est petit.
\end{astuce}

\section{Les pi\`eges de typage (et comment les nommer)}

\begin{codejs}
// 1. any : une sortie de secours qui inonde tout le fichier. Preferer unknown + etroitesse.
function mauvais(d: any) { return d.nope.deeper; }        // aucune aide
function bon(d: unknown) {
  if (typeof d === "object" && d !== null && "nope" in d) { /* ... */ }
}

// 2. as : l'affirmation silencieuse. Si tu mets beaucoup de "as", c'est le type du parametre
//    qui est faux, pas l'appel.
const n = reponse as number;          // "marche" a la compilation, casse a l'execution

// 3. null et undefined : le type ne protege que si strictNullChecks est actif
let v: string = null;                  // OK sans strict  <-- d'ou "strict": true

// 4. Objets litteraux inferes trop larges
const config = { retries: 3 };         // { retries: number }
const figee = { retries: 3 } as const; // { readonly retries: 3 }

// 5. Tableaux vs tuple
let a: string[] = ["x", "y"];          // taille libre
let t: [string, number] = ["x", 1];     // forme exacte
let m: [cle: string, valeur: number][];  // tuple nomme : le lecteur n'a plus a deviner

// 6. Callback optionnel et "this"
type Handlers = { onChange?: (v: string) => void };
declare const h: Handlers;
h?.onChange?.("ok");                     // l'alternative enchainee est ton amie

// 7. Erreurs de paquets "sans type" : installe @types/xxx, ou ecrit un module .d.ts
//    declare module "vieille-lib" { const x: any; export default x; }
// 8. "interface vs type" : interface est extensible (declaration merging), type est plus
//    expressif (unions, mappings). Pour une application : type partout, interface pour les
//    bibliotheques que d'autres etendent.
\end{codejs}

\section{Migrer un projet existant sans tout bloquer}

\begin{terminal}
npm i -D typescript @types/node
npx tsc --init --strict --allowJs --checkJs false --noEmit
# etape 1 : laisser tourner tsc sur les .js (checkJs: true) => inventaire des erreurs
# etape 2 : renommer les fichiers un a un en .ts (ou .tsx pour le JSX) ; git mv, pas de coup de main
# etape 3 : typer les frontieres (fonctions publiques, modeles) avant les variables internes
# etape 4 : activer les drapeaux stricts progressivement ; supprimer les "as any" au fil de l'eau
npx tsc --noEmit
\end{terminal}

Ordre de rentabilit\'e~: (1) les mod\`eles m\'etier (types partag\'es), (2) la couche r\'eseau,
(3) l'\'etat global, (4) les utilitaires, (5) le DOM. Les composants UI se typent presque seuls
une fois les props d\'efinies.

\begin{exercice}
\begin{enumerate}
  \item Reprends ton panier de la \cref{ch:js-struct} en TypeScript~: types \verb|Article|,
        \verb|Ligne|, \verb|Panier|, \verb|Totaux|~; \verb|Resultat<T> = {ok:true,valeur:T}| ...
        et aucun \verb|any|~; \verb|noUncheckedIndexedAccess| activ\'e.
  \item \'Ecris un type \verb|Schema<T>| (g\'en\'erique) et trois fonctions typ\'ees
        \verb|valider|, \verb|transformer|, \verb|serialiser| qui le r\'eutilisent.
  \item Typ\'e les \'ev\'enements DOM~: un gestionnaire g\'en\'erique
        \verb|surElement<E extends keyof HTMLElementEventMap>(el, type, cb)| o\`u le param\`etre du
        callback est d\'eduit correctement (\verb|MouseEvent|, \verb|InputEvent|...).
  \item Construis un client d'API o\`u chaque route renvoie un type~: \verb|api.get<Articles>("/articles")|
        avec un typage du \verb|path| par motif (\verb|`/articles/\${number}`|).
  \item Migre un de tes anciens projets JS (10 fichiers au plus) en TS, strict, sans
        d\'esactiver une seule r\`egle. Note les types d'erreurs que tu as trouv\'es~: tu viens de
        collecter un \'echantillon de bugs potentiels r\'eels.
\end{enumerate}
\end{exercice}

\begin{ressources}
\begin{itemize}
  \item Handbook officiel (\`a lire en entier, il est court) --- \url{https://www.typescriptlang.org/docs/handbook/intro.html}
  \item Playground pour exp\'erimenter sans installer --- \url{https://www.typescriptlang.org/play/}
  \item \link{Total~TypeScript}{https://www.totaltypescript.com/} (Matt~Pocock, exercices + vid\'eos)
  \item \link{TYPES-CHEATSHEET}{https://github.com/typescript-cheatsheets/react} pour React
  \item tsconfig reference --- \url{https://www.typescriptlang.org/tsconfig}
\end{itemize}
\end{ressources}

\index{TypeScript}\index{generiques}\index{tsconfig}
% ============================================================================
\chapter{Tester son code~: Vitest, Jest, Playwright}
\label{ch:js-tests}
% ============================================================================

\begin{objectif}
\begin{itemize}
  \item D\'ecider quoi tester (et quoi ne surtout pas tester).
  \item \'Ecrire des tests lisibles avec le sch\'ema Arrange-Act-Assert et un nommage qui documente.
  \item Ma\^itriser les mocks, les faux timers, et le mock de r\'eseau (msw).
  \item Tester l'interface (composants) et le parcours complet (Playwright).
  \item Brancher les tests dans une CI GitHub Actions qui refuse un mauvais commit.
\end{itemize}
\end{objectif}

\section{Pourquoi et quoi}

Un test n'est pas une certification de correction~: c'est un \emph{filet} qui te permet de
modifier le code sans peur. Le retour sur investissement vient des zones \`a r\`egles m\'etier,
pas des lignes de branchement trivial.

\begin{tabz}{lXll}
  \headrow \thd{\'Etage} \& \thd{Quoi} \& \thd{Vitesse} \& \thd{Part saine} \\
  \midrule
  Unitaires \& fonctions pures, r\`egles m\'etier, parseurs, validateurs \& ms \& 70 \% \\
  Int\'egration \& un module + sa base/son disque/son API simul\'ee \& dizaines de ms \& 20 \% \\
  Bout-en-bout \& parcours utilisateur r\'eel dans un navigateur \& secondes \& 10 \% \\
\end{tabz}

Ne teste \textbf{pas}~: les getters/setters, une biblioth\`eque tiers (React ne t'int\'eresse pas),
le nombre exact d'appels internes d'une impl\'ementation (\c{c}a casse au moindre refactor), ni le CSS.

\section{Vitest~: le choix par d\'efaut en 2026}

\begin{terminal}
npm i -D vitest @vitest/coverage-v8 jsdom fake-indexeddb
npx vitest run --coverage
npx vitest --ui                 # interface de developpement : surveillance, filtres, timelines
\end{terminal}

\begin{codejson}
{
  "test": {
    "environment": "node",
    "globals": false,
    "coverage": {
      "provider": "v8",
      "reporter": ["text", "html", "lcov"],
      "thresholds": { "lines": 80, "functions": 75, "branches": 70 }
    },
    "projects": [
      { "extends": true, "test": { "include": ["tests/unit/**"], "environment": "node" } },
      { "extends": true, "test": { "include": ["tests/dom/**"], "environment": "jsdom" } }
    ]
  }
}
\end{codejson}

\begin{codejs}
// tests/unit/panier.test.js
import { describe, it, expect, beforeEach, vi } from "vitest";
import { calculerTotaux, ajouter } from "../../src/lib/panier.js";

describe("calculerTotaux", () => {
  it("applique la TVA de 20 % sur le montant hors taxes", () => {
    const t = calculerTotaux([{ prix: 10, qte: 2 }]);
    expect(t.ht).toBe(20);
    expect(t.tva).toBeCloseTo(4, 5);
    expect(t.ttc).toBeCloseTo(24, 5);
  });

  it.each([
    [[], 0],
    [[{ prix: 12.5, qte: 1 }], 15],
    [[{ prix: 8, qte: 3 }], 28.8],
  ])("panier %j -> TTC %f", (lignes, ttc) => {
    expect(calculerTotaux(lignes).ttc).toBeCloseTo(ttc * 1.2, 5);
  });

  it("refuse une quantite negative", () => {
    expect(() => ajouter({ lignes: [] }, { prix: 1, qte: -3 })).toThrow(/quantite/);
  });

  it("ne mute pas le panier recu", () => {
    const initial = Object.freeze({ lignes: Object.freeze([{ prix: 1, qte: 1 }]) });
    const suivant = ajouter(initial, { prix: 2, qte: 1 });
    expect(suivant.lignes).toHaveLength(2);
    expect(initial.lignes).toHaveLength(1);
  });
});
\end{codejs}

\subsection{Les assertions qui comptent}

\begin{codejs}
expect(valeur).toBe(1);              // identite (Object.is) : primitives
expect(objet).toEqual({ a: 1 });     // structure profonde, ignore undefined
expect(objet).toStrictEqual({ a: 1, b: undefined });   // strict : types + cles undefined
expect(tab).toEqual(expect.arrayContaining([2, 3]));
expect(objet).toMatchObject({ a: { b: 1 } });           // sous-ensemble : le test qui ne casse pas pour rien
expect(chaine).toMatch(/^\\d{4}-\\d{2}-\\d{2}$/);
expect(n).toBeCloseTo(0.3, 5);                           // flottants
expect(p).resolves.toEqual(x); expect(p).rejects.toThrow(Error);
expect(fn).toHaveBeenCalledTimes(2);
expect(fn).toHaveBeenCalledWith({ id: 1 }, expect.any(Number));
expect(element).toBeInTheDocument();                      // avec @testing-library/jest-dom
expect(localStorage.getItem("k")).not.toBeNull();
\end{codejs}

\section{Mocks, espions, horloges}

\begin{codejs}
import { describe, it, expect, vi, beforeEach, afterEach } from "vitest";

// 1. Espionner une fonction locale : on verifie le dialogue avec une dependance
const logger = { info: vi.fn() };
await service.traiter(logger);
expect(logger.info).toHaveBeenCalledWith(expect.stringContaining("ok"));
expect(logger.info.mock.calls.length).toBe(1);
logger.info.mockImplementation(() => { throw new Error("panne"); });
logger.info.mockResolvedValue({ id: 1 });       // pour une fonction asynchrone
logger.info.mockRestore();

// 2. Mock d'un module entier (ex : la couche reseau)
vi.mock("../../src/api.js", () => ({
  api: { get: vi.fn().mockResolvedValue({ items: [{ id: 1, titre: "x" }] }) },
}));

// 3. Horloges fausses : tester un debounce, un retry, une expiration de session
vi.useFakeTimers();
const fn = vi.fn();
debounce(fn, 300)("a"); debounce(fn, 300)("ab");
await vi.advanceTimersByTimeAsync(300);
expect(fn).toHaveBeenCalledTimes(1);
expect(fn).toHaveBeenCalledWith("ab");
vi.setSystemTime(new Date("2026-01-15T10:00:00Z"));
vi.useRealTimers();

// 4. Substituts d'environnement
process.env.NODE_ENV = "test";
globalThis.fetch = vi.fn().mockResolvedValue(new Response(JSON.stringify({ ok: true }), { status: 200 }));
\end{codejs}

\begin{astuce}
Quand le mock devient plus long que le test, c'est que le code est mal d\'ecoup\'e~: la d\'ependance
est trop dure \`a isoler. Injecte-la (param\`etre, constructeur, contexte) au lieu de mocker le
module. Les bons tests viennent d'un bon d\'ecouplage, pas d'un bon moqueur.
\end{astuce}

\subsection{Mock de r\'eseau avec msw~: le test qui ne ment pas}

\begin{codejs}
// tests/msw/handlers.js
import { http, HttpResponse } from "msw";
export const handlers = [
  http.get("/api/articles", ({ request }) => {
    const url = new URL(request.url);
    const page = Number(url.searchParams.get("page") ?? 1);
    return HttpResponse.json({ items: [], page, total: 0 });
  }),
  http.post("/api/articles", async ({ request }) => {
    const corps = await request.json();
    if (!corps.titre) return HttpResponse.json({ message: "titre requis" }, { status: 422 });
    return HttpResponse.json({ id: 1, ...corps }, { status: 201 });
  }),
  http.get("/api/articles/:id", () => HttpResponse.error()),        // simule une panne reseau
];
// tests/setup.js
import { beforeAll, afterEach, afterAll } from "vitest";
import { setupServer } from "msw/node";
export const server = setupServer(...handlers);
beforeAll(() => server.listen({ onUnhandledRequest: "error" }));   // toute requete non prevue = echec
afterEach(() => server.resetHandlers());
afterAll(() => server.close());
\end{codejs}

\section{Tester l'interface~: par le comportement visible}

\begin{codejs}
// tests/dom/liste.test.jsx
import { render, screen, fireEvent, waitFor } from "@testing-library/react";
import userEvent from "@testing-library/user-event";
import { ListeCourses } from "../../src/ui/ListeCourses.jsx";

it("ajoute un article et le rend comme termine au clic", async () => {
  render(<ListeCourses />);
  const champ = screen.getByRole("textbox", { name: /article/i });     // les roles d'abord :
  await userEvent.type(champ, "Riz{Enter}");                            //  accessible = testable
  expect(screen.getByRole("listitem")).toHaveTextContent("Riz");
  await userEvent.click(screen.getByRole("checkbox"));
  expect(screen.getByRole("listitem")).toHaveClass("fait");
  await waitFor(() => expect(localStorage.getItem("courses")).toContain("Riz"));
});
\end{codejs}

R\`egle d'or~: teste ce que l'\emph{utilisateur} voit et fait, jamais l'impl\'ementation interne.
\verb|getByRole|, \verb|getByLabelText|, \verb|getByText| plut\^ot que \verb|getByTestId|~; et si
tu n'arrives pas \`a s\'electionner un \'el\'ement par son r\^ole, c'est un vrai bug d'accessibilit\'e
que tu viens de trouver.

\section{Playwright~: le test bout-en-bout qui vaut le d\'etour}

\begin{terminal}
npm init playwright@latest         # installe navigateurs + github action + exemple
npx playwright test --ui
npx playwright codegen http://localhost:5173   # enregistre tes clics en code
npx playwright show-report
\end{terminal}

\begin{codejs}
// tests/e2e/parcours.spec.ts
import { test, expect } from "@playwright/test";

test("un visiteur peut ajouter un article et le retrouve apres rechargement", async ({ page, context }) => {
  await page.goto("/");
  await page.getByLabel("Article").fill("Riz basmati");
  await page.getByRole("button", { name: "Ajouter" }).click();

  await expect(page.getByRole("listitem")).toHaveText(/Riz basmati/);
  await expect(page.getByText(/1 article/)).toBeVisible();

  await page.reload();
  await expect(page.getByRole("listitem")).toHaveText(/Riz basmati/);
});

test("une requete API lente n'affiche pas de doublon", async ({ page }) => {
  let appelee = 0;
  await page.route("**/api/articles", async (route) => {
    appelee++;
    await new Promise((r) => setTimeout(r, 800));
    route.fulfill({ json: { items: [], total: 0 } });
  });
  await page.goto("/articles");
  await page.getByRole("button", { name: /recharger/i }).click();
  await page.getByRole("button", { name: /recharger/i }).click();
  await expect(page.getByTestId("compteur-appels")).toHaveText("1");
});
\end{codejs}

\begin{note}
Playwright sert aussi \`a autre chose~: captures d'\'ecran automatiques pour ton README,
PDFs, g\'en\'eration de vignettes, tests multi-navigateurs et multi-\emph{viewports},
mocks r\'eseau au niveau HTTP (le plus fid\`ele qui existe). Un seul outil, quatre usages.
\end{note}

\section{Tests property-based~: chercher les cas que tu n'as pas imagin\'es}

\begin{codejs}
import fc from "fast-check";
it("le tri est une permutation ordonnee, pour tout tableau", () => {
  fc.assert(
    fc.property(fc.array(fc.integer()), (tab) => {
      const tri = trier(tab);
      expect(tri).toHaveLength(tab.length);
      expect([...tri].sort((a, b) => a - b)).toEqual(tri);
      expect(tri.reduce((s, x) => s + x, 0)).toBe(tab.reduce((s, x) => s + x, 0));
    }),
    { numRuns: 200 }
  );
});
it("serialiser/parser est un aller-retour", () => {
  fc.assert(fc.property(panierArbitraire, (p) => {
    expect(deverser(verser(p))).toEqual(p);
  }));
});
\end{codejs}

C'est l'arme id\'eale pour les parseurs, les slugifieurs, les codecs, les calculateurs de prix.

\section{Tests et int\'egration continue}

\begin{codeyml}
# .github/workflows/ci.yml
name: CI
on:
  push: { branches: [main] }
  pull_request:
concurrency: { group: ${{ github.workflow }}-${{ github.ref }}, cancel-in-progress: true }
jobs:
  qualite:
    runs-on: ubuntu-latest
    steps:
      - uses: actions/checkout@v4
      - uses: pnpm/action-setup@v4
      - uses: actions/setup-node@v4
        with: { node-version: 22, cache: pnpm }
      - run: pnpm install --frozen-lockfile
      - run: pnpm lint && pnpm typecheck
      - run: pnpm test -- --coverage
      - uses: actions/upload-artifact@v4
        with: { name: couverture, path: coverage/lcov-report }
  e2e:
    runs-on: ubuntu-latest
    needs: qualite
    steps:
      - uses: actions/checkout@v4
      - uses: actions/setup-node@v4
        with: { node-version: 22 }
      - run: npm ci && npm run build
      - uses: microsoft/playwright-github-action@v1
      - run: npx playwright install --with-deps chromium
      - run: npm run start & npx wait-on http://localhost:4173 && npm run test:e2e
      - uses: actions/upload-artifact@v4
        if: failure()
        with: { name: traces, path: test-results/ }
\end{codeyml}

\begin{exercice}
\begin{enumerate}
  \item \'Ecris une suite de 15 tests sur le module \verb|panier|~: cas nominaux, cas limites
        (vide, quantit\'es nulles, prix n\'egatifs, d\'ecimales), erreurs, immuabilit\'e.
  \item Ajoute un test de non-r\'egression~: reproduis un bug que tu as vraiment eu cette semaine,
        \'ecris le test qui \'echoue, puis corrige le code pour le faire passer.
  \item Mock le r\'eseau avec msw et teste la gestion d'erreur de ton client HTTP~: panne,
        401 (d\'econnexion provoqu\'ee), 500, r\'eponse lente avec annulation.
  \item Teste un composant UI avec Testing Library~: seulement par r\^oles et libell\'es. Note
        chaque fois que tu as \'et\'e tent\'e par un \verb|data-testid|~: c'\'etait une dette d'accessibilit\'e.
  \item \'Ecris deux tests Playwright~: parcours d'ajout/suppression, et gestion du mode hors ligne.
  \item Configure la couverture \`a seuil 80 \% et la CI GitHub Actions. V\'erifie qu'un commit
        rouge \'echoue bien~: si tu n'as pas vu l'action passer au rouge, ta CI est d\'ecorative.
\end{enumerate}
\end{exercice}

\begin{ressources}
\begin{itemize}
  \item Vitest --- \url{https://vitest.dev/guide/}
  \item Jest (si tu dois le maintenir) --- \url{https://jestjs.io/fr/docs/getting-started}
  \item Testing Library --- \url{https://testing-library.com/docs/react-testing-library/intro/}
  \item Playwright --- \url{https://playwright.dev/docs/intro}
  \item msw --- \url{https://mswjs.io/docs/} --- fast-check --- \url{https://fast-check.dev/}
  \item GitHub Actions, doc compl\`ete --- \url{https://docs.github.com/fr/actions}
\end{itemize}
\end{ressources}

\index{tests}\index{Vitest}\index{Playwright}\index{mock}
% ============================================================================
\chapter{S\'ecurit\'e web : ce qu'on doit savoir avant de mettre en ligne}
\label{ch:js-secu}
% ============================================================================

\begin{objectif}
\begin{itemize}
  \item Int\'egrer une r\`egle unique~: \textbf{toute donn\'ee qui vient de l'ext\'erieur est hostile}.
  \item Comprendre et bloquer les six attaques qui font 90 \% des incidents~: XSS, CSRF,
        injection SQL, IDOR, upload malveillant, d\'ependance compromise.
  \item Configurer en-t\^etes, cookies, CSP et politique de session correctement.
  \item Hacher un mot de passe comme il faut, et ne jamais \og ~inventer de la cryptographie~ \fg{}.
  \item Savoir o\`u chercher quand on veut aller plus loin (OWASP, outils, pratiques).
\end{itemize}
\end{objectif}

\section{Les trois principes qui rendent tout le reste \'evident}

\begin{enumerate}
  \item \textbf{Sortie encod\'ee, pas entr\'ee filtr\'ee.} Ne cherche pas \`a nettoyer une saisie pour
        qu'elle soit \og ~s\^ure~ \fg{}~: encode-la \emph{au moment de l'afficher}, selon le contexte
        (HTML, attribut, URL, JS, SQL). Un \verb|<script>alert(1)</script>| dans un champ \og ~nom~ \fg{}
        est une donn\'ee valide~; c'est son rendu qui doit devenir inoffensif.
  \item \textbf{D\'enominateur de confiance.} Le navigateur ne sait pas qui tu es~: il ne peut pas
        authentifier un appel. La confiance vient d'un secret partag\'e (cookie \verb|HttpOnly|,
        jeton sign\'e) \emph{et} d'une v\'erification d'autorisation \`a chaque requ\^ete.
  \item \textbf{Moindre privil\`ege.} Un utilisateur ne voit que ses donn\'ees. Un compte de base de
        donn\'ees n'a pas \verb|DROP|. Un conteneur ne tourne pas en root. Un jeton d'API a une
        port\'ee et une dur\'ee.
\end{enumerate}

\section{XSS~: le script malveillant dans ta page}

Un XSS, c'est du code que l'attaquant fait ex\'ecuter \emph{avec tes droits} dans le navigateur de
la victime~: vol de session, usurpation d'actions, enregistrement des frappes.

\begin{codejs}
// VULNERABLE : du texte utilisateur rendu comme du HTML
el.innerHTML = commentaire.texte;                      // <img src=x onerror="steal(document.cookie)">
document.querySelector("#u").innerHTML = `<a href="${utilisateur.site}">site</a>`;  // href="javascript:..."

// SUR : le texte reste du texte
el.textContent = commentaire.texte;
a.href = new URL(utilisateur.site, location.href).protocol.startsWith("http") ? utilisateur.site : "#";

// SI tu dois absolument rendre du HTML (markdown, richtext) : passer par un assainisseur
import DOMPurify from "dompurify";
el.innerHTML = DOMPurify.sanitize(markdownVersHtml(commentaire.texte), {
  ALLOWED_TAGS: ["p", "strong", "em", "a", "ul", "ol", "li", "code", "pre", "blockquote"],
  ALLOWED_ATTR: ["href", "title", "class"],
  ALLOW_DATA_ATTR: false,
  ADD_FORBID_TAGS: ["style"],
});
// Cote serveur, pour du contenu stocke une fois pour toutes : sanitize-html, DOMPurify,
// ou DOMPurify en JSDOM. Ne stocke JAMAIS du HTML non filtre.
\end{codejs}

\etape{Content Security Policy~: le filet de s\'ecurit\'e}
La CSP ne remplace pas l'encodage~; elle rend l'exploitation tr\`es difficile en interdisant
le JS inline et les sources inconnues.

\begin{codeconf}
# En-tete envoye par ton serveur (Express : helmet l'fait pour toi, personnalise-le)
Content-Security-Policy: default-src 'self'; base-uri 'self'; object-src 'none';
  frame-ancestors 'none'; form-action 'self'; upgrade-insecure-requests;
  script-src 'self' https://cdn.example.com 'nonce-{NONCE}';
  style-src 'self' 'unsafe-inline'; img-src 'self' data: https:; font-src 'self';
  connect-src 'self' https://api.nour.dev; report-uri /csp-report
Report-To: {"group":"csp","max_age":10886400,"endpoints":[{"url":"https://csp.nour.dev/reports"}]}
Cross-Origin-Opener-Policy: same-origin
Cross-Origin-Embedder-Policy: require-corp
Cross-Origin-Resource-Policy: same-origin
\end{codeconf}

Avec une nonce~: \verb|<script nonce="{valeur-aleatoire-par-requete}">|~; le JS inline sans nonce
est bloqu\'e. \'Evite \verb|unsafe-inline| et \verb|unsafe-eval|~: si tu les mets, la CSP est
d\'ecorative. Pour les frameworks qui exigent du style inline, utilise
\verb|style-src 'self' 'unsafe-inline'| (moins grave) ou un hachage.

\begin{itemize}
  \item G\'en\'erateur et validateur~: \link{CSP~Evaluator}{https://csp-evaluator.withgoogle.com/}.
  \item \texttt{Trusted Types} (Chrome/Edge/Firefox r\'ecents) rend \texttt{innerHTML} refus\'e sans
        une politique explicite~: \url{https://developer.mozilla.org/en-US/docs/Web/API/Trusted\\\_Types\\\_API}.
  \item \texttt{SRI} pour les scripts tiers~: \verb|<script src="..." integrity="sha384-..." crossorigin>|.
\end{itemize}

\section{CSRF~: le site d'un autre fait des requ\^etes en ton nom}

Un \verb|POST /compte/modifier-email| accept\'e sur la seule base du cookie de session est
exploitable depuis n'importe quelle page web. Trois d\'efenses, cumulables~:

\begin{enumerate}
  \item Cookie \verb|SameSite=Lax| (d\'efaut des navigateurs modernes)~: bloque la plupart des cas.
  \item V\'erifier l'\verb|Origin|/\verb|Referer| de la requ\^ete mutante~: peu co\^uteux, tr\`es efficace.
  \item Jeton CSRF~: une valeur secr\`ete al\'eatoire par session, renvoy\'ee dans le formulaire ou
        un en-t\^ete \verb|X-CSRF-Token|, recontr\^ol\'ee c\^ot\'e serveur pour toute m\'ethode
        non s\^ure (POST/PUT/PATCH/DELETE).
\end{enumerate}

\begin{codejs}
// Generation (cote serveur) puis controle. Le jeton est lie a la session, pas au navigateur.
import { randomBytes, timingSafeEqual } from "node:crypto";

function nouveauJeton(req) {
  req.session.csrf ??= randomBytes(32).toString("base64url");
  return req.session.csrf;
}
function verifierJeton(req) {
  const recu = req.get("x-csrf-token") ?? req.body?._csrf;
  const attendu = req.session?.csrf;
  if (!attendu || !recu) return false;
  const a = Buffer.from(recu), b = Buffer.from(attendu);
  return a.length === b.length && timingSafeEqual(a, b);       // comparaison temps constant
}
app.use((req, res, next) => {
  if (["GET", "HEAD", "OPTIONS"].includes(req.method)) return next();
  if (!verifierJeton(req)) return res.status(403).json({ erreur: "jeton CSRF manquant ou invalide" });
  next();
});
\end{codejs}

\section{Injection---le classique qui n'a pas vieilli}

\begin{codesql}
-- Ce que produit une concatenation naive :
SELECT * FROM users WHERE nom = '' OR '1'='1' -- '   -> TOUS les utilisateurs
-- En PHP comme en Node comme en Java : REQUETE PREPAREE, TOUJOURS.
\end{codesql}

\begin{codejs}
// Node + Postgres : les placeholders font le travail
const { rows } = await pool.query("SELECT * FROM articles WHERE slug = $1 AND publie = true", [slug]);

// Morceau de requete dynamique (tri, colonnes) : les placeholders ne suffisent pas.
// On ne fait jamais confiance a la chaine : on la CONFRONTE a une liste blanche.
const COLONNES = { recent: "cree_le", prix: "prix", nom: "nom" };
const ordre = COLONNES[req.query.tri] ?? "cree_le";
const sens = String(req.query.ordre).toLowerCase() === "asc" ? "ASC" : "DESC";
await pool.query(`SELECT * FROM articles ORDER BY ${ordre} ${sens} LIMIT $1 OFFSET $2`, [taille, decale]);
\end{codejs}

Le m\^eme principe prot\`ege des injections dans un shell~: ne jamais construire une commande avec
de la concat\'enation~; passe un \emph{tableau d'arguments} (\verb|execFile("git", [args])|,
\verb|execFile| et non \verb|exec|). Et pour les noms de fichiers~: \verb|path.resolve| puis
v\'erifier que le r\'esultat commence par le dossier autoris\'e.

\section{IDOR, autorisation, et le pi\`ege du \og ~c'est mon id~ \fg{}}

\begin{codejs}
// VULNERABLE : l'identifiant fourni par le client designe la ressource, sans controle
app.get("/api/factures/:id", async (req, res) => res.json(await facture(req.params.id)));
// -> n'importe qui remplace :id par celui du voisin.

// CORRECT : la requete est FILTREE par le proprietaire authentifie
app.get("/api/factures/:id", async (req, res) => {
  const f = await facturePour(req.user.id, req.params.id);   // WHERE id = $1 AND utilisateur_id = $2
  if (!f) return res.status(404).json({ erreur: "introuvable" });   // pas 403 : ne pas fuiter l'existence
  res.json(f);
});

// Autorisation centralisee, pas diss\'eminee dans les controleurs
const politiques = {
  "article:lire": (u, a) => a.publie || u.id === a.auteurId || u.roles.includes("editeur"),
  "article:ecrire": (u, a) => u.id === a.auteurId || u.roles.includes("editeur"),
  "compte:modifier": (u) => true,
};
export function autoriser(u, action, ressource) {
  const p = politiques[action];
  if (!p) throw new Error(`action inconnue : ${action}`);
  return Boolean(p(u, ressource));
}
\end{codejs}

\section{Mots de passe, sessions, jetons}

\begin{codejs}
// Hachage : argon2id (le meilleur) ou bcrypt (le plus courant). Jamais sha256(mdpmot + sel).
import argon2 from "argon2";
const hash = await argon2.hash(mdp, { type: argon2.argon2id, memoryCost: 19_456, timeCost: 2, parallelism: 1 });
await argon2.verify(hash, mdpSoumis);        // temps constant, sel int\'egr\'e

// Politique raisonnable (NIST SP 800-63B) : longueur >= 12, pas de composition absurde,
// comparaison aux mots de passe fuit\'es, pas de date d'expiration arbitraire.
import { createHash } from "node:crypto";
const hain = createHash("sha1").update(mdp).digest("hex");
const [prefixe, ...reste] = [hain.slice(0, 5), hain.slice(5)];
const rep = await fetch(`https://api.pwnedpasswords.com/range/${prefixe}`);
if ((await rep.text()).split("\n").some((l) => l.startsWith(reste.join("")))) {
  throw new Error("Ce mot de passe figure dans des fuites connues : choisis-en un autre");
}

// Reinitialisation : jeton a usage unique, aleatoire, court (15 min), stocke HACHE en base,
// envoye par email, invalide la session precedente, et reponse identique que l'email existe ou non.
\end{codejs}

\begin{tabz}{lX}
  \headrow \thd{Sujet} \& \thd{R\'eglage s\^ur} \\
  \midrule
  Cookie de session \& \verb|HttpOnly; Secure; SameSite=Lax; Path=/|, expiration 30 min -- 12 h. \\
  Rotation \& r\'eg\'en\'erer l'identifiant de session \`a la connexion et \`a l'\'el\'evation de privil\`ege. \\
  D\'econnexion \& invalider c\^ot\'e serveur (pas seulement effacer le cookie). \\
  \og ~Rester connect\'e~ \fg{} \& cookie longue dur\'ee avec un jeton \emph{stock\'e c\^ot\'e serveur}, r\'eversible. \\
  JWT \& courte dur\'ee (5--15 min), \verb|alg| \'epingl\'e (refuser \verb|none|), pas de donn\'ees sensibles, r\'evocation par liste ou refresh en cookie. \\
  Cl\'e API \& port\'ee minimale, rotation, jamais dans le code ni dans un bundle. \\
  CSRF \& jeton ou \verb|SameSite| + v\'erification d'\verb|Origin|. \\
  2FA \& TOTP \link{otplib}{https://github.com/yeqown/otplib} ou WebAuthn (bien meilleur) --- \url{https://webauthn.guide/} \\
\end{tabz}

\section{En-t\^etes de base, \`a poser partout}

\begin{codeconf}
Strict-Transport-Security: max-age=63072000; includeSubDomains; preload
X-Content-Type-Options: nosniff
X-Frame-Options: DENY                     (redondant avec frame-ancestors 'none', encore utile)
Referrer-Policy: strict-origin-when-cross-origin
Permissions-Policy: camera=(), microphone=(), geolocation=(), interest-cohort=()
Content-Security-Policy: (voir plus haut)
Server: (retirer ou rendre vague)
X-Powered-By: (retirer)
\end{codeconf}

\section{Upload, liens, et les petits d\'etails qui tuent}

\begin{itemize}
  \item \textbf{Upload}~: limite la taille, liste blanche d'\emph{extensions ET} de types MIME
        (v\'erifi\'es par analyse du contenu, pas par le nom), stocke hors racine web avec un nom
        g\'en\'er\'e, sert avec \verb|content-disposition: attachment| pour le HTML/svg, et jamais
        ex\'ecutable. Un \verb|.svg| peut contenir un \verb|<script>|~; un \verb|.php.jpg| peut
        tromper une v\'erification na\"ive.
  \item \textbf{Redirections ouvertes}~: \verb|?next=//evil.com|~: n'accepte que des chemins
        relatifs ou une liste blanche d'h\^otes.
  \item \textbf{Regex et ReDoS}~: limite la longueur avant de matcher.
  \item \textbf{D\'eni de service par requ\^ete co\^uteuse}~: paginer obligatoirement, limiter la
        taille des listes, mettre un cache.
  \item \textbf{Logs}~: ne jamais journaliser mot de passe, jeton, cookie, carte, document
        d'identit\'e. Masque et tronque.
  \item \textbf{Erreurs}~: en production, aucune stack trace, aucun SQL, aucun chemin de fichier
        au visiteur. (Les logs c\^ot\'e serveur, eux, gardent tout.)
  \item \textbf{D\'ependances}~: \verb|npm audit|, \link{Socket}{https://socket.dev/},
        lockfile verrouill\'e, \link{Renovate}{https://github.com/renovatebot/renovate} ou
        Dependabot pour les mises \`a jour automatiques.
\end{itemize}

\begin{terminal}
# Auto-\'evaluation d'un site que TU poss\`edes (ou un localhost que tu montes)
npx @lhce/lighthouse-cli --only-categories=security http://localhost:4173
npm i -D zap-baseline && npx zap-baseline.py -t http://localhost:3000 -r rapport.html
# ZAP (benchmarks OWASP) : https://www.zaproxy.org/   ---   Nuclei : https://github.com/projectdiscovery/nuclei
# Test manuel des en-tetes : https://securityheaders.com/  (public)
\end{terminal}

\begin{exercice}
\begin{enumerate}
  \item Fais tomber toi-m\^eme ton application de la semaine pr\'ec\'edente~: XSS stock\'e dans un
        commentaire, XSS par URL dans un lien, redirection ouverte, IDOR sur les factures,
        CSRF sur le changement d'email. \emph{Note l'exploit exact}, puis corrige-le, puis
        rejoue-le pour v\'erifier.
  \item Ajoute une CSP fonctionnelle en mode \verb|Content-Security-Policy-Report-Only|,
        corrige toutes les violations remont\'ees, puis active-la vraiment.
  \item \'Ecris un test automatis\'e pour chaque faille corrig\'ee (un test = une vuln\'erabilit\'e
        referm\'ee pour toujours).
  \item Audite ton \verb|package.json|~: nombre de d\'ependances transitives, licences, \^age des
        versions, paquets inutilis\'es~; remplace ce qui doit l'\^etre.
  \item Rends ta page d'authentification conforme~: cookie \verb|HttpOnly|, rotation de session,
        limitation de d\'ebit (5 essais/minute/IP), message identique que l'email existe ou non.
  \item Lis le OWASP Top~10 en une heure et, pour chaque point, \'ecris une ligne sur la fa\c{c}on
        dont ton projet actuel s'en prot\`ege ou ne s'en prot\`ege pas.
\end{enumerate}
\end{exercice}

\begin{ressources}
\begin{itemize}
  \item OWASP Top~10 (la r\'ef\'erence) --- \url{https://owasp.org/Top10/fr/}
  \item OWASP Cheat Sheet Series (fiches courtes et pr\'ecises) --- \url{https://cheatsheetseries.owasp.org/}
  \item MDN S\'ecurit\'e web --- \url{https://developer.mozilla.org/fr/docs/Web/Security}
  \item PortSwigger Web Security Academy (gratuit, exercices interactifs, le meilleur cours pratique) --- \url{https://portswigger.net/web-security}
  \item OWASP ZAP --- \url{https://www.zaproxy.org/} --- securityheaders.com
  \item \url{https://haveibeenpwned.com/} --- v\'erifier une fuite de son propre email.
\end{itemize}
\end{ressources}

\index{securite}\index{XSS}\index{CSRF}\index{CSP}
% ============================================================================
\chapter{Performance et mise en ligne}
\label{ch:deploy}
% ============================================================================

\begin{objectif}
\begin{itemize}
  \item Mesurer avant d'optimiser, avec les chiffres qui comptent (Core Web Vitals).
  \item Appliquer la liste des gains r\'eels~: poids, cache, images, polices, JavaScript.
  \item Construire une image Docker petite et s\^ure, la faire tourner derri\`ere nginx avec HTTPS.
  \item Automatiser le d\'eploiement par une pipeline, et pr\'evoir sauvegardes et surveillance.
\end{itemize}
\end{objectif}

\section{Mesurer~: les chiffres qui d\'ecident}

\begin{tabz}{lllX}
  \headrow \thd{Indice} \& \thd{Bon} \& \thd{\`A revoir} \& \thd{Ce qu'il mesure} \\
  \midrule
  LCP \& < 2{,}5~s \& > 4~s \& affichage du plus grand \'el\'ement (souvent le titre ou l'image principale). \\
  INP \& < 200~ms \& > 500~ms \& d\'elai de r\'eponse au clic/touche (remplace FID depuis 2024). \\
  CLS \& < 0{,}1 \& > 0{,}25 \& sauts de mise en page (images sans dimensions, polices, bandeaux). \\
  TTFB \& < 800~ms \& --- \& temps avant le premier octet~: le serveur, pas le front. \\
  Poids JS initial \& < 170~Ko gzip \& --- \& ce qui doit \^etre analys\'e avant l'interactivit\'e. \\
\end{tabz}

\begin{terminal}
npx lighthouse http://localhost:4173 --only-categories=performance,best-practices,seo,accessibility --output=html --output-path=rapport.html
npx lighthouse-ci   # seuils bloquants en CI : https://github.com/GoogleChrome/lighthouse-ci
# Outils en ligne (public) : https://pagespeed.web.dev/  ---  https://web.dev/measure/
# En conditions reelles : https://httparchive.org/ , le rapport CrUX de Search Console
\end{terminal}

\etape{Les 12 gains, par ordre de rendement}
\begin{enumerate}
  \item \textbf{R\'eduire le JavaScript.} Une d\'ependance de plus, c'est 100 \% de ses utilisateurs
        qui la t\'el\'echargent. \link{Bundlephobia}{https://bundlephobia.com/} pour estimer un paquet,
        \link{bundlejs}{https://bundlejs.com/} pour comparer des versions.
  \item \textbf{D\'ecouper le code}~: \verb|React.lazy| / \verb|import()| par route~; charger
        l'\'editeur, les graphes, l'export PDF seulement quand on les ouvre.
  \item \textbf{Compresser}~: Brotli (meilleur que gzip~; nginx \verb|brotli on|),
        et \verb|content-encoding| v\'erifi\'e dans l'onglet R\'eseau.
  \item \textbf{Images}~: WebP/AVIF, dimensions r\'ealistes, \verb|loading="lazy"|, \verb|width| et
        \verb|height| (CLS), \verb|fetchpriority="high"| sur l'image principale,
        SVG pour les ic\^ones. Un \link{service d'images}{https://imgproxy.net/} ou l'URL de
        l'h\'ebergeur suffit le plus souvent.
  \item \textbf{Polices}~: \verb|font-display: swap|, \verb|preload| des deux fichiers utilis\'es,
        sous-ensemble latin (\link{glyphhanger}{https://github.com/zachleat/glyphhanger}),
        ou police syst\`eme (0~Ko~!).
  \item \textbf{Cache HTTP}~: assets avec empreinte et \verb|immutable|, HTML en \verb|no-cache|.
  \item \textbf{Rendre au plus t\^ot}~: pr\'e-rendu/SSG pour du contenu, \verb|app.js| en \verb|defer|,
        CSS critique inline pour le premier \'ecran.
  \item \textbf{R\'eduire les allers-retours}~: peu d'origines diff\'erentes, HTTP/2 ou~3,
        \verb|preconnect| vers l'unique tiers indispensable.
  \item \textbf{Mesurer l'INP}~: d\'ecouper les grosses t\^aches (\verb|scheduler.yield|,
        \verb|setTimeout(\ldots, 0)|, Web Worker)~; ne rien calculer dans le handler de saisie.
  \item \textbf{\'Eviter les requ\^etes en cascade}~: une requ\^ete qui en attend trois autres est
        trois allers-retours. Regrouper c\^ot\'e serveur (un endpoint de page qui renvoie tout).
  \item \textbf{Penser au mobile et au r\'eseau lent}~: throttling 4G lente dans DevTools avant
        chaque mise en ligne~; les 60 fps d'un MacBook ne sont pas ceux d'un t\'el\'ephone \`a 120~\texteuro{}.
  \item \textbf{V\'erifier les donn\'ees expos\'ees}~: requ\^etes inutiles, images non compress\'ees,
        \verb|sourcemap| en production (utile, mais \`a prot\'eger), logs c\^ot\'e client.
\end{enumerate}

\subsection{Cache HTTP~: le tableau \`a afficher au mur}

\begin{codeconf}
# assets versionnes (nom avec empreinte) : ne jamais revalider
Cache-Control: public, max-age=31536000, immutable

# HTML et donnees : toujours reverifier, mais reponse 304 bon marche
Cache-Control: no-cache
ETag: W/"a3f9..."

# API authentifiaue : hors cache partag\'es
Cache-Control: private, no-store

# fichiers d'import (CSV) : court
Cache-Control: public, max-age=300, stale-while-revalidate=60
\end{codeconf}

\begin{note}
Un \verb|Cache-Control: immutable| sur un fichier dont le nom contient son empreinte
(\verb|app.4f9c12.js|) est la strat\'egie la plus efficace qui existe~: le navigateur ne redemande
\emph{jamais}. Change le contenu, le nom change, l'URL change~: la casse du cache est r\'esolue
sans effort. C'est exactement ce que fait un build Vite.
\end{note}

\section{Construire une image Docker}

\begin{codeconf}
# Dockerfile multi-etapes : build avec les outils, execution avec le strict minimum
FROM node:22-alpine AS build
WORKDIR /app
ENV CI=1
COPY package.json package-lock.json ./
RUN npm ci                       # respecte exactement le lockfile
COPY . .
RUN npm run build && npm prune --omit=dev

FROM node:22-alpine AS runtime
WORKDIR /app
ENV NODE_ENV=production
# securite : un utilisateur dedie, pas root
RUN addgroup -S app && adduser -S app -G app
COPY --from=build --chown=app:app /app/dist ./dist
COPY --from=build --chown=app:app /app/node_modules ./node_modules
COPY --from=build --chown=app:app /app/package.json ./
USER app
EXPOSE 3000
HEALTHCHECK --interval=30s --timeout=3s --start-period=10s --retries=3 \
  CMD wget -qO- http://127.0.0.1:3000/api/sante || exit 1
# init (PID 1) propre : recoit SIGTERM, vide les connexions
CMD ["node", "--max-old-space-size=768", "dist/server.js"]
\end{codeconf}

\begin{terminal}
docker build -t nour/api:1.4.0 .
docker run --rm -p 3000:3000 --env-file .env.production -m 512m --cpus 1 nour/api:1.4.0
docker image ls          # taille : vise < 200 Mo pour une API Node
docker compose up -d --build
\end{terminal}

Bonnes pratiques d'image~: une \'etape de build s\'epar\'ee, \verb|npm ci|, \verb|--omit=dev|,
\verb|alpine| ou \verb|distroless|, utilisateur non-root, \verb|.dockerignore| complet,
couches dans l'ordre de stabilit\'e (manifeste d'abord pour tirer parti du cache),
scan~: \link{Trivy}{https://github.com/aquasecurity/trivy} (\cmd{trivy image mon/api:1.4.0}).

\section{D\'eployer : les quatre chemins r\'ealistes}

\begin{tabz}{lXll}
  \headrow \thd{Chemin} \& \thd{Pour qui} \& \thd{Co\^ut indicatif} \& \thd{Liens} \\
  \midrule
  PaaS (Render, Railway, Fly.io) \& z\'ero ops, push et c'est en ligne \& 0--7~/mois \& \url{https://render.com/}, \url{https://railway.com/}, \url{https://fly.io/} \\
  Plateformes front (Vercel, Netlify, Cloudflare Pages) \& sites statiques et fonctions serverless \& 0--20 \& \url{https://vercel.com/}, \url{https://www.netlify.com/}, \url{https://pages.cloudflare.com/} \\
  VPS + Docker (+ Coolify/Dokploy) \& contr\^ole total, le plus formateur \& 4--8~/mois \& \url{https://coolify.io/}, \url{https://dokploy.com/} \\
  Mutualis\'e PHP (OVH, o2switch, Hostinger) \& sites PHP/MySQL, tr\`es bon march\'e \& 3--8~/mois \& \url{https://www.ovhcloud.com/fr/}, \url{https://www.o2switch.fr/} \\
\end{tabz}

\etape{Le VPS, pas \`a pas (le chemin que je recommande d'apprendre)}
\begin{terminal}
# 1. Serveur : Hetzner Cloud, OVH, Contabo... Debian 12 ou Ubuntu 24.04 LTS
ssh root@IP
adduser deploy && usermod -aG sudo deploy
rsync -avS ~/.ssh/id_ed25519.pub deploy@IP:.ssh/authorized_keys
# coupe l'acces root par mot de passe : /etc/ssh/sshd_config -> PermitRootLogin no, PasswordAuthentication no
systemctl restart ssh
ufw allow OpenSSH && ufw allow 80 && ufw allow 443 && ufw enable
apt update && apt install -y fail2ban qemu-guest-agent

# 2. En tant que deploy
curl -fsSL https://get.docker.com | sudo sh
sudo usermod -aG docker $USER
git clone https://github.com/nour/api.git && cd api
cp .env.example .env && nano .env            # secrets reels, jamais dans Git
docker compose up -d

# 3. Nom de domaine : enregistre un A record vers l'IP (1h de propagation)
# 4. TLS gratuit (Let's Encrypt) via Caddy, le plus simple des reverse proxies
sudo apt install -y caddy
cat > /etc/caddy/Caddyfile <<'CADDY'
api.nour.dev {
    encode zstd gzip
    reverse_proxy 127.0.0.1:3000
    header {
        Strict-Transport-Security "max-age=31536000; includeSubDomains"
        X-Content-Type-Options "nosniff"
        Referrer-Policy "strict-origin-when-cross-origin"
        -Server
    }
}
CADDY
sudo systemctl reload caddy
\end{terminal}

\etape{Avec nginx \`a la place (le plus fr\'equent en entreprise)}
\begin{codeconf}
server {
  listen 80;
  server_name nour.dev www.nour.dev;
  return 301 https://$host$request_uri;
}
server {
  listen 443 ssl http2;
  server_name nour.dev;
  ssl_certificate     /etc/letsencrypt/live/nour.dev/fullchain.pem;
  ssl_certificate_key /etc/letsencrypt/live/nour.dev/privkey.pem;
  client_max_body_size 12m;
  gzip on; gzip_types text/css application/javascript application/json image/svg+xml;

  location / {
    proxy_pass http://127.0.0.1:3000;
    proxy_http_version 1.1;
    proxy_set_header Upgrade $http_upgrade;              # WebSocket
    proxy_set_header Connection "upgrade";
    proxy_set_header Host $host;
    proxy_set_header X-Forwarded-For $proxy_add_x_forwarded_for;
    proxy_set_header X-Forwarded-Proto $scheme;
    proxy_read_timeout 75s;                              # > keepAliveTimeout du serveur Node
  }
  location /assets/ {
    root /var/www/nour/dist;
    add_header Cache-Control "public, max-age=31536000, immutable";
  }
}
\end{codeconf}

\begin{terminal}
sudo certbot --nginx -d nour.dev -d www.nour.dev --redirect -m toi@example.com --agree-tos
sudo nginx -t && sudo systemctl reload nginx
sudo certbot renew --dry-run
\end{terminal}

\subsection{Processus~: systemd plut\^ot que \og ~node app.js~ \fg{} dans un terminal}

\begin{codeconf}
# /etc/systemd/system/nour-api.service
[Unit]
Description=API Nour
After=network-online.target docker.service
Wants=network-online.target

[Service]
Type=notify
User=deploy
WorkingDirectory=/home/deploy/api
EnvironmentFile=/home/deploy/api/.env.production
ExecStart=/usr/bin/docker compose up --remove-orphans
ExecReload=/bin/kill -s HUP $MAINPID
ExecStop=/usr/bin/docker compose down -t 30
Restart=always
RestartSec=3
MemoryMax=600M
NoNewPrivileges=true
ProtectSystem=strict
ProtectHome=read-only
ReadWritePaths=/home/deploy/api/data

[Install]
WantedBy=multi-user.target
\end{codeconf}

\begin{terminal}
sudo systemctl daemon-reload && sudo systemctl enable --now nour-api
journalctl -u nour-api -f
\end{terminal}

\section{Le d\'eploiement continu~: pousser = en ligne}
\label{sec:ci-cd}

\begin{codeyml}
# .github/workflows/deploy.yml
name: Deploiement
on:
  push: { branches: [main] }
jobs:
  livrer:
    runs-on: ubuntu-latest
    environment: production
    steps:
      - uses: actions/checkout@v4
      - uses: docker/setup-buildx-action@v3
      - uses: docker/login-action@v3
        with:
          registry: ghcr.io
          username: ${{ github.actor }}
          password: ${{ secrets.GITHUB_TOKEN }}
      - uses: docker/build-push-action@v6
        with:
          push: true
          tags: |
            ghcr.io/${{ github.repository }}:latest
            ghcr.io/${{ github.repository }}:${{ github.sha }}
          cache-from: type=gha
          cache-to: type=gha,mode=max
      - name: Deployer sur le VPS
        uses: appleboy/ssh-action@v1
        with:
          host: ${{ secrets.HOST }}
          username: deploy
          key: ${{ secrets.SSH_KEY }}
          script: |
            cd ~/api
            docker compose pull
            docker compose up -d
            curl -fsS https://nour.dev/api/sante
\end{codeyml}

Un d\'eploiement qui se passe bien a quatre propri\'et\'es~: \textbf{idempotent} (relancer ne casse rien),
\textbf{atomique} (soit l'ancienne version, soit la nouvelle, jamais un m\'elange),
\textbf{rapide \`a annuler} (une commande), \textbf{observ\'e} (journaux + m\'etriques pendant 5 minutes).
Ajoute des \emph{migrations} versionn\'ees et r\'etrocompatibles~: la base \'evolue avant le code, jamais
pendant.

\section{Apr\`es la mise en ligne~: ce qui te sauvera}

\begin{itemize}
  \item \textbf{Sauvegardes}~: la r\`egle 3-2-1 (trois copies, deux supports, une hors site).
        Un cron \verb|pg\_dump -Fc| compress\'e et copi\'e chez un autre fournisseur. Teste la
        \emph{restauration}--- une sauvegarde non test\'ee n'existe pas.
  \item \textbf{Surveillance}~: une page de statut, une alerte par courriel/SMS sur le healthcheck,
        et des tableaux de bord~: \link{Uptime~Kuma}{https://github.com/louislam/uptime-kuma}
        (auto-h\'eberg\'e, 10 minutes \`a installer), \link{Grafana~Cloud}{https://grafana.com/products/cloud/},
        \link{Sentry}{https://sentry.io/} pour les erreurs JavaScript.
  \item \textbf{Journaux}~: structur\'es (JSON), horodat\'es UTC, rotation, corr\'elation par identifiant
        de requ\^ete. \verb|pino| + \link{loki}{https://grafana.com/oss/loki/} suffisent pour commencer.
  \item \textbf{Secrets}~: jamais dans Git~; variables d'environnement du syst\`eme ou
        gestionnaire d\'edi\'e~; rotation en cas de doute.
  \item \textbf{Journal de bord}~: \`a chaque incident, \'ecris ce qui s'est pass\'e, la cause, la
        correction. Dans un an, c'est ta meilleure documentation.
\end{itemize}

\begin{projet}[D\'eploiement complet]
Mets ton API Node du \cref{ch:node} en ligne~: domaine, HTTPS, conteneur, healthcheck,
journalisation, sauvegarde quotidienne de la base, pipeline de d\'eploiement, et rapport
Lighthouse public. Tu auras alors fait le parcours complet~: coder, s\'ecuriser, tester, d\'eployer,
surveiller. C'est exactement ce que l'on attend d'un profil junior~; tr\`es peu de candidats
peuvent le montrer.
\end{projet}

\begin{exercice}
\begin{enumerate}
  \item Ex\'ecute Lighthouse sur ton projet et corrige les trois recommandations au plus fort
        impact. Note les chiffres avant/apr\`es dans ton carnet.
  \item Divise ton bundle JavaScript par deux~: lazy-loading d'une route, remplacement d'une
        d\'ependance lourde, suppression de code mort (\verb|knip|).
  \item \'Ecris le \verb|Dockerfile| de ton application et mesure la taille de l'image avant et
        apr\`es multi-\'etapes.
  \item Monte un nginx/Caddy local avec HTTPS auto-sign\'e (mkcert) et v\'erifie la r\'e\'ecriture
        des en-t\^etes \verb|X-Forwarded-*|.
  \item Simule une coupure~: arr\^ete la base de donn\'ees, observe le comportement de l'API
        (elle doit refuser proprement, pas corrompre), puis restaure.
  \item Fais un faux d\'eploiement cass\'e~: la pipeline doit le d\'etecter, et une commande doit
        revenir \`a la version pr\'ec\'edente en moins de 30 secondes.
\end{enumerate}
\end{exercice}

\begin{ressources}
\begin{itemize}
  \item web.dev (mesure, optimisation, motifs) --- \url{https://web.dev/explore}
  \item web.dev Vitals --- \url{https://web.dev/vitals/}
  \item Docker docs (best practices images) --- \url{https://docs.docker.com/build/building/best-practices/}
  \item Caddy --- \url{https://caddyserver.com/docs/} --- certbot --- \url{https://certbot.eff.org/}
  \item \link{12~Factor~App}{https://12factor.net/fr/} --- la grammaire des applications web modernes
  \item Awesome Selfhosted (tout l'\'ecosyst\`eme auto-h\'eberg\'e) --- \url{https://github.com/awesome-selfhosted/awesome-selfhosted}
\end{itemize}
\end{ressources}

\index{performance}\index{Core~Web~Vitals}\index{Docker}\index{deploiement}\index{Caddy}

\part{PHP : le moteur du web dynamique}
% ============================================================================
\chapter{PHP : les bases, le langage serveur le plus direct qui soit}
\label{ch:php-bases}
% ============================================================================

\begin{objectif}
\begin{itemize}
  \item Comprendre le cycle d'une requ\^ete PHP et l'installer en local en trois commandes.
  \item \'Ecrire du PHP moderne~: typage strict, tableaux, fonctions, fermetures.
  \item Ma\^itriser le type \og ~array~ \fg{}, qui est en r\'ealit\'e une table ordonn\'ee, et ses fonctions.
  \item Structurer un projet PHP sans framework~: un front controller, des fonctions, de la clart\'e.
  \item Conna\^itre les erreurs de d\'ebutant (headers, encodage, \verb|==|, includes) et les \'eviter.
\end{itemize}
\end{objectif}

\section{PHP, c'est quoi exactement}

PHP est un moteur qui ex\'ecute un script \emph{\`a chaque requ\^ete}, puis meurt. Il n'y a pas d'\'etat
persistant entre deux requ\^etes~: tout ce que tu ne stockes pas (base de donn\'ees, fichier, session,
cache m\'emoire) est perdu quand le script se termine. Cette contrainte explique \`a la fois la
simplicit\'e de PHP (pas de serveur \`a garder vivant, pas d'\'ev\'enements, pas d'async) et ses pi\`eges.

\begin{enumerate}
  \item Une requ\^ete HTTP arrive (Apache avec \verb|mod\_php|, ou nginx + \verb|php-fpm|).
  \item le moteur \textbf{compile} le script en opcode, le met en cache (\textbf{OPcache})~;
  \item le script s'ex\'ecute de la premi\`ere \`a la derni\`ere ligne~; il produit du texte (souvent du HTML)~;
  \item la r\'eponse est envoy\'ee, la m\'emoire est lib\'er\'ee. Point final.
\end{enumerate}

Version de r\'ef\'erence pour ce livre~: \textbf{PHP 8.3 / 8.4}. Tout ce qui est marqu\'e \og ~PHP~8+~ \fg{}
n'existe pas en 7.4 (finie de support, encore trop fr\'equente sur les vieilles mutualisations).

\begin{terminal}
# Verifier / installer
php -v
# Debian/Ubuntu : sudo apt install php-cli php-fpm php-mysql php-curl php-mbstring php-xml php-zip
# macOS          : brew install php
# Windows        : https://windows.php.net/download/ (voir chap. 1)
# Alternative : conteneur
docker run --rm -it -v $PWD:/app -w /app php:8.3-cli php script.php

# Le serveur de developpement integre : remplace apache pour apprendre
php -S localhost:8000 -t public
php -S 127.0.0.1:8000 public/router.php       # avec un routeur pour simuler les rewrites
\end{terminal}

\begin{note}
Le serveur int\'egr\'e (\verb|php -S|) est mono-thread~: une seule requ\^ete \`a la fois, parfait pour
apprendre, jamais pour un test de charge. Il refuse aussi les connexions entrantes autres que
localhost par d\'efaut~: utilise \verb|0.0.0.0:8000| pour tester depuis un t\'el\'ephone.
\end{note}

\section{Premier script, et les r\`egles de syntaxe}

\begin{codephp}
<?php
declare(strict_types=1);                    // PRECIEUX : interdit la conversion implicite des types
header('Content-Type: text/html; charset=utf-8');

$prenom = 'Nour';
$age = 28;
printf('%s a %d ans.%s', $prenom, $age, PHP_EOL);

// Une balise ouvrante, une seule par fichier ; JAMAIS de balise fermante dans un fichier
// qui ne contient que du PHP (sinon tu ajoutes un saut de ligne a la reponse, et les
// en-tetes "headers already sent" deviennent un cauchemar).
\end{codephp}

\begin{tabz}{lX}
  \headrow \thd{R\`egle} \& \thd{D\'etail qui compte} \\
  \midrule
  \verb|;| obligatoire \& chaque instruction se termine par un point-virgule. Le \verb|?>| final en implique un. \\
  \verb|\$| devant toute variable \& \verb|\$x|~; \verb|x| est une constante ou un nom de fonction. \\
  Insensible \`a la casse pour les mots-cl\'es/fonctions \& \verb|ECHO|, \verb|strlen| et \verb|StrLen| sont \'egaux. \textbf{Les noms de classes le sont aussi}~: d'o\`u les conventions PSR. \\
  Sensible pour les variables \& \verb|\$nom| et \verb|\$Nom| sont deux variables diff\'erentes. \\
  Cha\^ines \& \verb|'simple'| = litt\'erale~; \verb|"double"| = interpol\'ee (\verb|"\$a[0]|, \verb|"{\$obj->p}|). \\
  Commentaires \& \verb|//|, \verb|\#|, \verb|/* */|, et \verb|/** DocBlock */| pour la documentation lue par les IDE. \\
  \verb|echo| / \verb|print| / \verb|printf| \& \verb|echo| plusieurs valeurs~; \verb|print| renvoie 1~; \verb|printf| formate. \\
\end{tabz}

\section{Types, variables, et le typage moderne}

\begin{codephp}
<?php
declare(strict_types=1);

// Types scalaires
$entier = 42;            // int (64 bits, sauf sur 32 bits : 2^31-1 -- un id de tweet qui casse, ca existe)
$flottant = 3.14;        // float ; 1e3, 0x1A, 0b1010, 1_000_000 sont valides
$chaine = "cinq";        // string : une suite d'octets, pas de notion de caractere unicode
$booleen = true;         // bool
$rien = null;            // null : "absence de valeur"
$tableau = [1, 2, 3];    // array
$objet = new stdClass(); // object
$ressource = fopen('php://memory', 'r');   // ressource (fermeture explicite recommandee)

// Declarations typ\'ees : la norme en PHP moderne
function compter(array $articles, ?string $filtre = null, int $limite = 20): int
{
    $n = 0;
    foreach ($articles as $a) {
        if ($filtre === null || str_contains((string) $a['titre'], $filtre)) {
            $n++;
            if ($n >= $limite) { break; }
        }
    }
    return $n;
}

// Types d'union et d'intersection, nommes, readonly (8.1+), jamais (8.1)
function id(): int|string { return random_int(1, 1000); }
function traiter(Countable&ArrayAccess $c): never { throw new RuntimeException('non'); }

// Propriete typ\'ee : pas de valeur par defaut implicite -> une erreur nette si non initialisee
final class Point { public function __construct(public float $x = 0.0, public readonly float $y = 0.0) {} }

// gettype() est un outil de debug ; get_debug_type() (8.0) est plus lisible
// var_dump() : le reflexe. settype(), intval() : a eviter, prefere les casts (int), (float), (string), (bool), (array)

// Comparaison : comme en JS, mais en pire. TOUJOURS === et !==.
var_dump(0 == 'a');           // PHP 8 : false (avant PHP 8 : TRUE -- un CVE en puissance)
var_dump('12abc' == 12);      // false en PHP 8 (true en PHP 7)
var_dump(100 === '100');      // false
var_dump(NAN === NAN);        // false
var_dump(0.1 + 0.2 === 0.3);  // false ; compare avec abs(a-b) < PHP_FLOAT_EPSILON
\end{codephp}

\subsection{Cha\^ines : tout ce qu'on utilise vraiment}

\begin{codephp}
<?php
$s = "  Le Chat et le CHIEN  ";

strlen($s);                       // en OCTETS ("\'e" = 2 octets en UTF-8) -> ne l'utilise pas pour l'affichage
mb_strlen($s, 'UTF-8');           // en caracteres reels ; mb_* = l'extension indispensable
mb_strtolower($s);  mb_convert_case($s, MB_CASE_TITLE, 'UTF-8');
str_contains($s, 'chat');          // 8.0 : remplace strpos !== false
str_starts_with($s, '  Le');  str_ends_with($s, ' ');
trim($s);  ltrim($s);  rtrim($s, " \t\n\r\0\x0B");
str_replace(['le', 'et'], ['LES', 'ET'], strtolower($s));
str_repeat('-', 10);  substr($s, 2, 5);  substr_count($s, 'e');
strpos(strtolower($s), 'chat');     // position ou false
explode(',', 'a,b,c');  implode(', ', ['a', 'b']);
sprintf('%s : %05d | %-10s | %+.2f', 'facture', 42, 'total', -12.5);
number_format(1234.5678, 2, ',', ' ');        // "1 234,57" (fr)
nl2br($texteUtilisateur);                       // sauts de ligne -> <br>
strip_tags($html);  htmlspecialchars($html, ENT_QUOTES | ENT_SUBSTITUTE, 'UTF-8');   // echapper pour le HTML
wordwrap($s, 80, "\n", true);  str_pad('7', 3, '0', STR_PAD_LEFT);
str_word_count($s);  ucfirst($s);  ucfirst(mb_strtolower($s));

// Interpolation et formatage multi-lignes
$prenom = 'Nour';
echo "Bonjour $prenom, tu as " . strlen($prenom) . " lettres.\n";
echo "Detail : {$article['titre']} | {$user->nom}\n";
$sql = <<<SQL
    SELECT id, titre
    FROM articles
    WHERE publie = 1
    ORDER BY id DESC
    SQL;                     // 7.3+ : le cloturant peut etre indentes
$md = <<<'TXT'               // nowdoc = comme une chaine simple, aucune interpolation
Texte litteral : $variable, {$choses}, \n
TXT;
\end{codephp}

\begin{piege}
\textbf{Encodage.} PHP manipule des octets, pas des caract\`eres. Sans \verb|mbstring| et sans
UTF-8 coh\'erent, tu auras des \og ~\~A~ \fg{} partout. La checklist~: fichiers en UTF-8 sans BOM~;
base en \verb|utf8mb4| avec \verb|utf8mb4\_unicode\_ci|~; en-t\^ete \verb|Content-Type: text/html; charset=utf-8|~;
et les fonctions \verb|mb\_*| pour tout ce qui touche au texte humain.
\end{piege}

\section{Tableaux~: le vrai type universel de PHP}

Un tableau PHP est une \textbf{table ordonn\'ee}~: il accepte des cl\'es (entiers ou cha\^ines) et
maintient l'ordre d'insertion. C'est \`a la fois un tableau, une liste, un dictionnaire et une file.

\begin{codephp}
<?php
$article = [
    'id' => 7,
    'titre' => 'Apprendre PHP',
    'tags' => ['php', 'web'],
    'auteur' => ['nom' => 'Nour', 'email' => 'nour@example.com'],
    0 => 'premiere cle numerique -- attention, "0" devient 0',
];

// Acces et ecriture
$article['titre'];
$article['auteur']['nom'];
$article['vue'] ??= 0;                    // 7.4+ : initialise si absent/nul
$article['titre'] .= ' (v2)';
unset($article['tags']);

// Creation et fusion
$nouveau = array_fill_keys(['a', 'b'], 0);
$combine = array_combine(['id', 'nom'], [1, 'Nour']);
$fusion = $article + ['id' => 999];                 // les cles existantes gagnent
$fusion2 = array_merge($a, $b);                    // reindexe les cles numeriques  <-- piege
$deep = array_replace_recursive($config, $surcharges);
$tabs = [...$premiers, ...$derniers];              // spread : seulement listes (cles numeriques seq.)

// Cles et valeurs
array_keys($article);  array_values($article);
count($article);  in_array('php', $article['tags'] ?? [], true);   // le 3e parametre = comparaison stricte
array_key_exists('id', $article);     // != isset() : lui voit les cles a valeur null

// Transformer (les trois qui remplacent 90 % des boucles)
$mots = explode(' ', strtolower(trim($s)));
$longueurs = array_map(mb_strlen(...), $mots);
$filtr\'es = array_filter($mots, static fn (string $m): bool => mb_strlen($m) > 3);
$reindexe = array_column($articles, null, 'id');           // clef = id -> recherche O(1)
$titre = array_map(static fn (array $a): string => $a['titre'], $articles);
$total = array_reduce($lignes, static fn (float $c, array $l): float => $c + $l['prix'] * $l['qte'], 0.0);

// Ordre
sort($liste);                        // par valeur, reindexe -- MUTANT
asort($tableauAssoc);                // par valeur, garde les cles
ksort($parClef);
usort($articles, static fn (array $a, array $b): int => $b['date'] <=> $a['date']);
// spaceship <=> renvoie -1/0/1 : la seule bonne facon de comparator en PHP
array_multisort(array_column($articles, 'note'), SORT_DESC, $articles);   // multi-criteres

// Decoupage, piles, files
array_slice($liste, 2, 5, preserve: true);
array_chunk($articles, 20);
array_splice($liste, 1, 2, $remplacement);
array_push / array_pop / array_shift / array_unshift ;
$array = [1, 2, 3]; $x = array_pop($array);      // 3
end($array) / reset($array) / next($array) / current($array) / key($array);
list($a, $b) = [1, 2];   [$x, $y] = [1, 2];    ['nom' => $n, 'id' => $i] = $article;   // destructuration
\end{codephp}

\section{Conditions, boucles, et le flux de contr\^ole}

\begin{codephp}
<?php
$note = 14;
if ($note >= 16) {
    echo 'Tres bien';
} elseif ($note >= 10) {
    echo 'Admis';
} else {
    echo 'Releve';
}

// alternative syntax : tres lisible dans les vues (HTML mele au PHP)
foreach ($articles as $i => $article): ?>
    <article class="<?= $article['publie'] ? 'publie' : 'brouillon' ?>">
        <h2><?= htmlspecialchars($article['titre']) ?></h2>
        <p><?= e($article['extrait'] ?? '') /* fonction utilitaire a definir */ ?></p>
    </article>
<?php endforeach;

switch ($etat) {                    // switch traditionnel : comparaison loose (==) -- evite-le
    case 'ok': break;
    default: throw new InvalidArgumentException("etat inconnu : {$etat}");
}
$resultat = match ($note) {        // 8.0 : comparaison stricte, expression, pas de fallthrough
    20 => 'parfait',
    16, 17, 18, 19 => 'tres bien',
    10 ... 15 => 'bien',            // plage
    default => 'a revoir',
};
$classe = match (true) {
    $note >= 16 => 'excellent',
    $note >= 10 => 'moyen',
    default => 'faible',
};

for ($i = 0; $i < 10; $i++) {}
while ($ligne = fgets($h)) {}        // idiomatique : affectation + test
do { $n++; } while ($n < 10);
foreach ($tableau as $clef => $valeur) {}
foreach ($generateur() as $ligne) {}    // memoire constante sur un gros fichier
for (;;) { break; }                     // boucle infinie propre
\end{codephp}

\section{Fonctions, fermetures, appelables}

\begin{codephp}
<?php
declare(strict_types=1);

// Fonction nommee, typ\'ee, avec valeurs par defaut
function envoyerEmail(
    string $destinataire,
    string $sujet,
    string $corps = '',
    array $piecesJointes = [],
    ?DateTimeImmutable $programmePour = null,
): bool {
    if (!filter_var($destinataire, FILTER_VALIDATE_EMAIL)) {
        throw new InvalidArgumentException("email invalide : {$destinataire}");
    }
    return true;    // ici, l'appel a la fonction mail() du serveur SMTP reel
}

// Nomme (8.0) : l'API devient lisible, l'ordre des parametres ne compte plus
envoyerEmail(destinataire: 'nour@example.com', piecesJointes: ['facture.pdf'], sujet: 'Facture');

// Passage par valeur (defaut) vs par reference (a eviter, sauf cas precis)
function ajout(int &$compteur): void { $compteur++; }
function tableauEnReferenc(array &$t): void { $t['vu'] = true; }

// Fermetures : "use" importe le contexte (par copie ; par reference avec "&")
$tva = 0.2;
$ht = [10, 20, 30];
$ttc = array_map(static fn (float $x): float => $x * (1 + $tva), $ht);
$compteur = function () use (&$n): int { return ++$n; };     // le "&" est essentiel ici
$appliquer = fn (callable $f, $v) => $f($v);                  // fleche : capture auto, une expression

// Generateur : produire une sequence sans la stocker
function lignes(string $chemin): Generator
{
    $h = fopen($chemin, 'rb');
    if ($h === false) { throw new RuntimeException("lecture impossible : {$chemin}"); }
    try {
        while (($l = fgets($h)) !== false) {
            yield rtrim($l, "\r\n");
        }
    } finally {
        fclose($h);
    }
}
foreach (lignes('acces.log') as $n => $ligne) { /* memoire = 1 ligne */ }

// Callables : le PHP n'aime pas les fonctions de haut niveau silencieuses
$tri = usort($articles, 'comparerParDate');            // nom de fonction
$tri = usort($articles, [$objet, 'methode']);
$tri = usort($articles, [MaClasse::class, 'statique']);
$tri = usort($articles, $objet->methode(...));          // premier-class callable (8.1) : le plus propre
$tri = usort($articles, fn (array $a, array $b) => $b['id'] <=> $a['id']);

// Functions utiles a connaitre : call_user_func est depasse par $cb(...) ; array_* le reste.
\end{codephp}

\subsection{Les 15 fonctions globales qu'on cherche toujours}

\begin{tabularx}{\textwidth}{@{}l X@{}}
  \toprule
  \textbf{Fonction} \& \textbf{Usage}\\
  \midrule
  \verb|count|, \verb|iterator\_to\_array|, \verb|is\_array| \& dimensions et v\'erifications.\\
  \verb|array\_map|, \verb|array\_filter|, \verb|array\_values| \& transformations.\\
  \verb|str\_contains|, \verb|str\_replace|, \verb|sprintf|, \verb|trim| \& cha\^ines.\\
  \verb|htmlspecialchars| (\verb|ENT\_QUOTES|) \& sortir du texte en HTML.\\
  \verb|json\_encode| / \verb|json\_decode|, avec \verb|JSON\_THROW\_ON\_ERROR| \& API et configs.\\
  \verb|filter\_var| / \verb|filter\_input| \& validation d'email, d'URL, d'entier.\\
  \verb|password\_hash| / \verb|password\_verify| \& mots de passe.\\
  \verb|random\_int| / \verb|random\_bytes| \& al\'eatoire \emph{cryptographique} (jamais \verb|rand| pour un jeton).\\
  \verb|file\_get\_contents| / \verb|file\_put\_contents| / \verb|is\_readable| \& fichiers simples.\\
  \verb|date|, \verb|DateTimeImmutable| \& dates~; pr\'ef\'erence absolue \`a la classe.\\
  \verb|var\_dump|, \verb|print\_r|, \verb|error\_log| \& debug (et \verb|dump()| de Symfony).\\
  \verb|compact|, \verb|extract| \& variables vers une vue (le second est \`a fuir).\\
  \verb|enum|? \verb|enum\_exists|, \verb|class\_exists|, \verb|function\_exists| \& introspection.\\
  \verb|match|, \verb|throw expr| (8.0) \& flux de contr\^ole.\\
\end{tabularx}

\section{TP --- un mini-routage et un rapport CSV}

Deux fichiers, sans base de donn\'ees, sans framework. L'id\'ee~: voir que PHP n'a besoin de rien
de plus pour faire un vrai site.

\begin{codephp}
<?php
// public/index.php --- le "front controller" : TOUT passe par la.
declare(strict_types=1);
error_reporting(E_ALL);
ini_set('display_errors', '1');          // en DEV seulement ; en prod : 0 + logs

session_start();
$uri = parse_url($_SERVER['REQUEST_URI'], PHP_URL_PATH) ?: '/';
$methode = $_SERVER['REQUEST_METHOD'];

// Une petite fonction d'echappement, reutilisee dans les vues
function e(?string $v): string { return htmlspecialchars((string) $v, ENT_QUOTES | ENT_SUBSTITUTE, 'UTF-8'); }

$routes = [
    'GET /' => 'accueil',
    'GET /articles' => 'articles',
    'POST /articles' => 'creerArticle',
    'GET /rapport.csv' => 'rapport',
];
$action = $routes["$methode $uri"] ?? null;
if ($action === null) {
    http_response_code(404);
    require __DIR__ . '/../src/vues/404.php';
    exit;
}
require __DIR__ . "/../src/actions/{$action}.php";
\end{codephp}

\begin{codephp}
<?php
// src/actions/articles.php
declare(strict_types=1);
$cheminDonnees = __DIR__ . '/../../data/articles.json';
$articles = json_decode((string) file_get_contents($cheminDonnees), true, 512, JSON_THROW_ON_ERROR);
$filtre = trim((string) ($_GET['q'] ?? ''));
$parPage = 10;
$page = max(1, (int) ($_GET['page'] ?? 1));

$q = array_values(array_filter(
    $articles,
    static fn (array $a): bool => $filtre === ''
        || mb_stripos($a['titre'], $filtre) !== false
        || mb_stripos($a['contenu'] ?? '', $filtre) !== false
));
$total = count($q);
$nbPages = max(1, (int) ceil($total / $parPage));
$page = min($page, $nbPages);
$articles = array_slice($q, ($page - 1) * $parPage, $parPage);

require __DIR__ . '/../vues/liste.php';
\end{codephp}

\begin{codephp}
<?php
// src/actions/rapport.php --- exporter un CSV correctement (le BOM UTF-8 est un detail reel)
declare(strict_types=1);
$articles = json_decode((string) file_get_contents(__DIR__ . '/../../data/articles.json'), true, 512, JSON_THROW_ON_ERROR);

$entetes = ['id', 'titre', 'auteur', 'date', 'vues'];
$corps = fopen('php://output', 'wb');
fwrite($corps, "\xEF\xBB\xBF");                     // BOM : Excel ouvre alors en UTF-8
fputcsv($corps, $entetes, ',', '"', '\\');           // 1ere colonne forcee en chaine, 4e argument = 8.4
foreach ($articles as $a) {
    fputcsv($corps, [
        (string) $a['id'],
        $a['titre'],
        $a['auteur']['nom'] ?? '',
        date('Y-m-d', (int) ($a['cree'] ?? time())),
        (int) ($a['vues'] ?? 0),
    ]);
}
fclose($corps);
exit;
\end{codephp}

\begin{note}
Les \verb|require|/ \verb|include| ont deux variantes. \verb|require| \'echoue fatalement,
\verb|include| \'emet un avertissement~; \verb|*\_once| \'evite le double-chargement. \verb|\_\_DIR\_\_| rend
le chemin insensible au r\'epertoire courant~: utilise-le syst\'ematiquement, sinon ton application
casse d\`es que le serveur est lanc\'e depuis un autre dossier.
\end{note}

\subsection{Les erreurs de PHP que les d\'ebutants rencontrent}

\begin{tabz}{lX}
  \headrow \thd{Message} \& \thd{Cause r\'eelle} \\
  \midrule
  \verb|Cannot modify header information - headers already sent| \& Un espace, un retour \`a la ligne ou un \verb|echo| avant \verb|header()|~: v\'erifie l'absence de \verb|?>| final et d'octets parasites avant \verb|<?php|. \\
  \verb|undefined array key "x"| \& la cl\'e n'existe pas~; utilise \verb|??| ou \verb|array\_key\_exists|. \\
  \verb|TypeError: Return value must be of type int, string returned| \& \verb|strict\_types=1| fait son travail~: convertis explicitement avec \verb|(int)| ou \verb|intcap|. \\
  \verb|Allowed memory size exhausted| \& boucle infinie, tableau gigantesque, ou r\'ecursion sans fin. \\
  Page blanche \& une erreur fatale avec \verb|display\_errors=Off|. Regarde le journal~: \verb|tail -f /var/log/php-fpm/error.log|, ou active \verb|error\_reporting(E\_ALL)| + \verb|display\_errors=1|. \\
  \verb|Class "X" not found| \& autoload non configur\'e (Composer n'est pas install\'e/charg\'e), ou nom d'espace de lignes incorrect. \\
\end{tabz}

\begin{exercice}
\begin{enumerate}
  \item \'Ecris \verb|slugify(string \$t): string|~: minuscules, accents retir\'es
        (\verb|iconv('UTF-8', 'ASCII//TRANSLIT//IGNORE', \$t)| ou \verb|mb\_convert\_encoding|),
        non-alphanum\'eriques en tirets, tirets dupliqu\'es et marges nettoy\'es~; test\'e sur
        \og ~L'\'Ev\'enement de l'ann\'ee, \`a Thi\`es !~ \fg{}.
  \item \'Ecris \verb|formaterDuree(int \$secondes): string| renvoyant \og ~2 h 05 min~ \fg{}, et
        \verb|nombreEnLettres(int \$n): string| pour 0--999 (r\'ecursif).
  \item Construis \verb|tableauVersCsv(array \$lignes): string| et \verb|csvVersTableau(string): array|
        qui se comportent correctement avec guillemets, virgules et retours \`a la ligne dans les champs.
  \item Analyse un fichier \verb|access.log|~: renvoie les 10 IP les plus fr\'equentes, le total de
        requ\^etes par code HTTP, et la liste des 404, avec un g\'en\'erateur pour tenir sur un fichier
        de 1~Go.
  \item \'Ecris une fonction \verb|validerFormulaire(array \$donnees, array \$regles): array|~: r\`egles
        \verb|required|, \verb|min:5|, \verb|email|, \verb|numeric|, \verb|in:a,b,c|,
        \verb|confirmed|~; retour \verb|[champ => [erreurs]]| vide si tout passe. C'est la base
        de ton futur micro-framework.
  \item \'Ecris le \verb|creerArticle.php| de l'exemple~: v\'erifie la m\'ethode, valide, ajoute au JSON
        avec un verrou (\verb|flock|) et une \'ecriture atomique (\verb|tempnam| + \verb|rename|),
        renvoie une redirection (303) vers la page de l'article.
\end{enumerate}
\end{exercice}

\begin{ressources}
\begin{itemize}
  \item Le manuel officiel (avec les notes d'utilisateurs~; la 2\textsuperscript{e} meilleure doc du web) --- \url{https://www.php.net/manual/fr/}
  \item PHP The Right Way (traduit, \`a lire en entier~: les bonnes pratiques modernes) --- \url{https://phptherightway.com/}
  \item PHPFIG standards --- \url{https://www.php-fig.org/psr/}
  \item php-pdo / PDO---\url{https://www.php.net/manual/en/book.pdo.php}
  \item \link{PHP~Watched}{https://php.watch/} (nouveautes par version) et \url{https://stitcher.io/blog/php-84-in-a-hurry}
\end{itemize}
\end{ressources}

\index{PHP!bases}\index{array}\index{strict\_types}
% ============================================================================
\chapter{PHP : tableaux et cha\^ines en profondeur, dates et regex}
\label{ch:php-struct}
% ============================================================================

\begin{objectif}
\begin{itemize}
  \item Choisir la bonne structure~: liste, dictionnaire \verb|array|, \verb|ArrayObject|, g\'en\'erateur.
  \item Concat\'ener, comparer, d\'ecouper du texte unicode sans se faire surprendre.
  \item Manipuler les dates avec \verb|DateTimeImmutable| (et ne plus toucher \`a \verb|date()|).
  \item \'Ecrire des expressions r\'eguli\`eres PCRE\textsuperscript{}{2} correctes et rapides.
  \item Tri multi-crit\`eres, regroupement, indexation, partitions~: les motifs r\'eels du CRUD.
\end{itemize}
\end{objectif}

\section{Les quatre visages du tableau PHP}

\begin{codephp}
<?php
// 1. LISTE indexee : clefs 0..n. Les operations qui vont bien :
$liste = [];
$liste[] = 'a';                       // empiler a la fin ( amorti O(1) )
array_push($liste, 'b', 'c');
array_pop($liste);                     // depiler la fin
array_shift($liste);                   // O(n) : decale tout -> a eviter sur les grandes listes
$clefs = array_keys($liste);  $n = count($liste);
sort($liste);                          // reindexe
asort / ksort / usort / uasort / uksort ;
shuffle / reverse / array_slice / array_chunk ;

// 2. DICTIONNAIRE : cles chaines. C'est LA structure de hachage de PHP (rapide, ordonnee).
$parId = [];
foreach ($articles as $a) { $parId[$a['id']] = $a; }      // recherche ensuite en O(1)
unset($parId[42]);
foreach ($parId as $id => $a) { /* preserver les cles */ }

// 3. OBJET tableau-like : ArrayObject / ArrayIterator (iteration, count, acces [] )
$collection = new ArrayObject($articles);
$collection->asort();

// 4. GENERATEUR : la meme interface, memoire constante
function lireCsv(string $chemin, int $saut = 1): Generator
{
    $h = fopen($chemin, 'rb');
    if ($h === false) { throw new RuntimeException("impossible d'ouvrir {$chemin}"); }
    try {
        for ($i = 0; $i < $saut; $i++) { fgetcsv($h); }
        while (($ligne = fgetcsv($h, 0, ',', '"', '\\')) !== false) {
            yield $ligne;
        }
    } finally {
        fclose($h);
    }
}
foreach (lireCsv('clients.csv') as $n => $ligne) { traiter($ligne, $n); }   // 5 Go ? aucun souci.
\end{codephp}

\begin{piege}
Trois comportements de tableau qui pi\`egent tout le monde~:\\
(1) \verb|array\_merge| \textbf{r\'eindexe} les cl\'es num\'eriques et \textbf{ne fusionne pas
r\'ecursivement}. Pour un dictionnaire~: \verb|+\$tableau| (les cl\'es existantes gagnent),
\verb|array\_replace\_recursive| ou \verb|...\$spread|~;
(2) une cl\'e de tableau \verb|"1"| devient l'entier \verb|1|~: \verb|['1' => 'a']| a la cl\'e
\verb|1|. Les identifiants num\'eriques en provenance d'une base doivent \^etre cast\'es en
\verb|(string)| si tu t'en sers comme cl\'es~;
(3) un tableau est copi\'e \`a l'affectation (s\'emantique de valeur), mais la copie est \og ~lazy~ \fg{}~:
elle reste partag\'ee tant que rien n'est modifi\'e. \verb|=\&| et les objets suivent une autre r\`egle.
\end{piege}

\section{Les motifs de transformation qui remplacent les boucles}

\begin{codephp}
<?php
$articles = json_decode((string) file_get_contents('data/articles.json'), true, 512, JSON_THROW_ON_ERROR);

// Map / filter / reduce : PHP n'a pas de chainage, mais les fermetures rendent cela lisible
$titres = array_map(static fn (array $a): string => trim($a['titre']), $articles);
$publies = array_values(array_filter($articles, static fn (array $a): bool => !empty($a['publie_le'])));
$vues = array_reduce($articles, static fn (int $c, array $a): int => $c + (int) ($a['vues'] ?? 0), 0);

// Regrouper / indexer
$parAnnee = [];
foreach ($articles as $a) {
    $parAnnee[substr((string) $a['publie_le'], 0, 4)][] = $a;
}
$parAnnee = array_reduce($articles, static function (array $acc, array $a): array {
    $acc[substr((string) $a['publie_le'], 0, 4)][] = $a;
    return $acc;
}, []);
$index = array_column($articles, null, 'id');                       // id => article
$titreParId = array_combine(array_column($articles, 'id'), array_column($articles, 'titre'));

// Trier par un champ, puis par un second (le spaceship est ton ami)
usort($articles, static fn (array $a, array $b): int => [$b['publie_le'], $a['titre']] <=> [$a['publie_le'], $b['titre']]);
usort($articles, static fn (array $a, array $b) => mb_strtolower($a['titre']) <=> mb_strtolower($b['titre']));

// Compter / distinct / intersections / differences
$nombresParTag = [];
foreach ($articles as $a) { foreach ($a['tags'] ?? [] as $t) { $nombresParTag[$t] = ($nombresParTag[$t] ?? 0) + 1; } }
arsort($nombresParTag);
$top = array_slice($nombresParTag, 0, 10, true);
$tagsUniques = array_values(array_unique(array_merge(...array_column($articles, 'tags'))));
$communs = array_intersect($tagsDeA, $tagsDeB);
$manquants = array_diff($attendus, $presents);

// Pagination et aggregations sans SQL
$parPage = 20;
$page = max(1, (int) ($_GET['page'] ?? 1));
$totalPages = (int) ceil(count($filtres) / $parPage);
$pageCourante = array_slice($filtres, ($page - 1) * $parPage, $parPage, true);   // true = garder les cles

// "group by having" a la main : array_count_values est sous-estime
$statut = array_count_values(array_map(static fn (array $a): string => $a['statut'] ?? 'inconnu', $articles));
\end{codephp}

\section{Cha\^ines~: unicode, d\'ecoupage, \'echappements}

\begin{codephp}
<?php
$texte = "L'\'el\'ephant d'\'Egypte -- 123 "citations"";

// Longueur et extraction : toujours les versions mb_*
mb_strlen($texto = $texte);                          // caracteres reels
mb_substr($texto, 0, 10);  mb_strimwidth($texto, 0, 40, '...');   // troncature avec ellipse, la bonne
mb_convert_case($texto, MB_CASE_TITLE, 'UTF-8');      // "L'\'El\'ephant D'\'egypte..."
ucfirst(mb_strtolower($texto));
mb_convert_encoding($texto, 'ISO-8859-1', 'UTF-8');   // pour un export vieux-systeme
mb_detect_encoding($texto, ['UTF-8', 'ISO-8859-1'], true);
iconv('UTF-8', 'ASCII//TRANSLIT', '\'e -> e, \c{c} -> c');  // translitteration (les accents disparaissent)

// Sans mbstring, le "caract\`ere" n'existe pas : strlen compte des OCTETS.
strlen('\'e');           // 2  <-- ne fait pas ce que tu crois
str_contains($texte, '\'Egypte');   // casse dependante : 1 si exactement la meme sequenced'octets
mb_stripos($texte, '\'Egypte');     // 13 (insensible a la casse, en caracteres)

// Echappement : le bon outil depend UNIQUEMENT du contexte de sortie
$html = htmlspecialchars($texte, ENT_QUOTES | ENT_SUBSTITUTE | ENT_HTML5, 'UTF-8');
$url = rawurlencode($texte);                        // %C3%A9... ; http_build_query(['q' => $texte])
$sql = $pdo->quote($texte);                          // ou mieux : requete preparee
$coquille = escapeshellarg($texte);                   // si tu dois appeler un binaire
$json = json_encode($texte, JSON_UNESCAPED_UNICODE | JSON_INVALID_UTF8_SUBSTITUTE);
$regex = preg_quote($texte, '/');                    // neutraliser une entree dans un motif
$csv = '"' . str_replace('"', '""', $texte) . '"';   // champ CSV qui contient virgules/quotes

// Decoupage intelligent (mots, phrases, lignes)
$mots = preg_split('/[\s\p{P}]+/u', $texte, -1, PREG_SPLIT_NO_EMPTY);
$lignes = preg_split('/\R/u', $texte);                // \R = n'importe quelle fin de ligne
$phrases = preg_split('/(?<=[.!?])\s+/u', $texte);
$entites = html_entity_decode($html, ENT_QUOTES | ENT_HTML5, 'UTF-8');

// Comparaison et tri "humains"
strcmp('a', 'b');  strcasecmp('A', 'a');
strcoll('\'elan', 'efant');                              // selon la locale (LC_COLLATE)
'\'el\'ephant' <=> 'zebre';                                // binaire : evite pour l'affichage
usort($mots, fn ($a, $b) => Collator::compare($a, $b));   // intl : la vraie comparaison linguistique
// En pratique : IntlChar/Transliterator (intl) ou tri en base avec une collation utf8mb4_unicode_ci.
\end{codephp}

\section{Expressions r\'eguli\`eres (PCRE2)}

\begin{codephp}
<?php
// Delimiteur /, modificateurs en queue. Le 'u' (UTF-8) est OBLIGATOIRE sur du texte humain.
$texte = "Contact : nour@example.com ou 06 12 34 56 78 -- www.nour.dev/a?b=1";

if (preg_match('/[a-z0-9._%+-]+@[a-z0-9.-]+\.[a-z]{2,}/iu', $texte, $m)) {
    [$tout, $debut] = $m;
    echo $tout;
}
preg_match_all('#https?://[^\s"]+#iu', $texte, $liens, PREG_SET_ORDER | PREG_OFFSET_CAPTURE);
// Groupes nommes : plus lisible, et robuste si l'ordre change
$motif = '/^(?<jour>\d{2})\/(?<mois>\d{2})\/(?<annee>\d{4})$/';
if (preg_match($motif, '31/01/2026', $m) === 1) {
    $d = DateTimeImmutable::createFromFormat('!Y-m-d', "{$m['annee']}-{$m['mois']}-{$m['jour']}");
}
// Remplacements, y compris par callback (le "sprintf" des regex)
$sortie = preg_replace_callback(
    '/\{(\w+)\}/',
    static fn (array $m): string => (string) ($donnees[$m[1]] ?? $m[0]),
    "Bonjour {prenom}, panier {total} EUR"
);
$normalise = preg_replace(['/[\s]+/', '/[^\w\s-]/u', '/(-+)\p{L}+/u'], [' ', '', ''], $titre);
// Limites obligatoires en environnement non maitrise (protection ReDoS)
if (@preg_match($motifDangereux, $entree, $m, 0, 0, 1_000_000) === false) {
    throw new RuntimeException('motif rejete par la limite de backtracking');
}
// Les modificateurs qui comptent
// u = UTF-8 ; i = insensible ; m = ^ et $ par ligne ; s = . inclut \n ; x = mode extended (commentaires)
$verbose = '/
    ^\+?\d{1,3}          # indicatif pays optionnel
    (?:[\s.-]?\d){8,10}  # groupes de chiffres
    $/ux';
\end{codephp}

\begin{note}
\verb|preg\_match| renvoie \verb|1| (trouv\'e), \verb|0| (non trouv\'e) ou \verb|false| (erreur de motif
ou limite d\'epass\'ee). C'est pour cela qu'on \'ecrit \verb|=== 1| et non un test de v\'erit\'e~:
\verb|false| est \og ~faux~ \fg{}, mais pour des raisons compl\`etement diff\'erentes de \verb|0|.
Les fonctions \verb|preg\_*| sont les seules \`a renvoyer \verb|false| sur erreur~: v\'erifie
\verb|preg\_last\_error\_msg()| quand un r\'esultat semble bizarre.
\end{note}

\section{Dates~: la classe ou rien}

\begin{codephp}
<?php
$maintenant = new DateTimeImmutable('now', new DateTimeZone('UTC'));
$debut = new DateTimeImmutable('2026-01-01 00:00:00', new DateTimeZone('Europe/Paris'));
$fin = DateTimeImmutable::createFromFormat('!d/m/Y', '31/01/2026', new DateTimeZone('Europe/Paris'));
if ($fin === false) { throw new RuntimeException(implode(', ', DateTimeImmutable::getLastErrors()['messages'] ?? [])); }

// Manipulations : les objets sont immuables, chaque appel renvoie un NOUVEL objet
$demain = $debut->modify('+1 day')->setTime(9, 30);
$trimestre = $debut->add(new DateInterval('P3M'));
$intervalle = $debut->diff($fin);
$jours = (int) $intervalle->format('%r%a');             // %r = signe, %a = jours totaux
$secondes = $fin->getTimestamp() - $debut->getTimestamp();

// Affichage
$debut->format('Y-m-d H:i:s');   $debut->format('d/m/Y H\hi');   $debut->format(DATE_ATOM);   // 2026-01-01T00:00:00+01:00
$debut->format('l, j F Y');      // noms en anglais ; pour le francais, passe par IntlDateFormatter
setlocale(LC_TIME, 'fr_FR.UTF-8'); strftime('%A %d %B %Y', $debut->getTimestamp());  // strftime = deprecie :
(new IntlDateFormatter('fr_FR', IntlDateFormatter::LONG, IntlDateFormatter::NONE))->format($debut);

// Comparaisons, periodes de facture, dates de fin de mois
$estPasse = $fin < $maintenant;
$debutMois = $maintenant->modify('first day of this month 00:00');
$finMois = $maintenant->modify('last day of this month 23:59:59');
$dernierJourvrier = (new DateTimeImmutable('2028-02-01'))->modify('last day of february 2028');

// Toujours en UTC en base, et une seule timezo ne appliquee a l'affichage.
$pdo = new PDO($dsn);  $pdo->exec("SET TIME ZONE 'UTC'");
$sqlDate = $debut->setTimezone(new DateTimeZone('UTC'))->format('Y-m-d H:i:s');
$affiche = (new DateTimeImmutable($sqlDate, new DateTimeZone('UTC')))->setTimezone(new DateTimeZone('Africa/Dakar'));
\end{codephp}

\begin{piege}
Un \verb|DateTimeImmutable| garde son fuseau~; \verb|createFromFormat| sans \verb|!| h\'erite
des heures non pr\'ecis\'ees de \emph{maintenant} (le champ \verb|!| veut dire \og ~mettre \`a z\'ero~ \fg{}).
Les formats \verb|c|, \verb|ATOME| et \verb|Y-m-d\TH:i:sP| sont interchangeables~: choisis-en un
seul et \'ecris-le dans la spec de ton API. Enfin, ne compare jamais deux dates avec
\verb|==|~: compare les \verb|getTimestamp()| ou utilise \verb|<| \verb|>|, qui fonctionnent
correctement sur les objets \verb|DateTimeInterface|.
\end{piege}

\section{JSON, s\'erialisation et fichiers de configuration}

\begin{codephp}
<?php
// JSON : les trois drapeaux que je mets toujours
$json = json_encode(
    ['id' => 1, 'titre' => 'Caf\'e', 'note' => 12.5, 'date' => $maintenant],
    JSON_THROW_ON_ERROR | JSON_UNESCAPED_UNICODE | JSON_PRETTY_PRINT
);
// JSON_PRETTY_PRINT seulement en dev/debug ; JSON_HEX_TAG/AMP prot\`egent dans un <script>.
$data = json_decode($json, true, 512, JSON_THROW_ON_ERROR);           // assoc = vrai : tableaux
$objet = json_decode($json, flags: JSON_THROW_ON_ERROR);              // ou objets stdClass
$objetTyp\'e = json_decode($json, flags: JSON_THROW_ON_ERROR, depth: 32);  // puis hydrate a la main

// Hydratation manuelle d'un objet metier : 20 lignes, aucune dependance
final class Article
{
    public function __construct(
        public readonly int $id,
        public readonly string $titre,
        public readonly ?DateTimeImmutable $publieLe,
        /** @var string[] */ public readonly array $tags = [],
    ) {}
    public static function deTableau(array $d): self
    {
        return new self(
            id: (int) ($d['id'] ?? throw new InvalidArgumentException('id requis')),
            titre: trim((string) ($d['titre'] ?? '')),
            publieLe: isset($d['publie_le']) ? new DateTimeImmutable((string) $d['publie_le']) : null,
            tags: array_values(array_filter((array) ($d['tags'] ?? []), 'is_string')),
        );
    }
    public function versTableau(): array
    {
        return ['id' => $this->id, 'titre' => $this->titre, 'tags' => $this->tags,
                'publie_le' => $this->publieLe?->format(DATE_ATOM)];
    }
}

// Config : un fichier PHP est plus rapide a lire qu'un YAML et n'exige aucune librairie
// config/app.php -> return [...];   $config = require 'config/app.php';
// Fichiers d'environnement (.env) : vlucas/phpdotenv ou symfony/dotenv (voir chap. Composer)
\end{codephp}

\begin{exercice}
\begin{enumerate}
  \item \verb|groupesParPrefixe(array \$mots): array|~: regroupe les mots par leurs trois premi\`eres
        lettres, trie les groupes par taille d\'ecroissante, et renvoie un tableau index\'e
        de \verb|[prefixe, mots]|.
  \item \verb|reformaterCSV(string \$chemin, array \$colonnes): string|~: lis un CSV (avec le s\'eparateur
        devin\'e parmi \verb|;|, \verb|,|, \verb|\t|), s\'electionne/renomme les colonnes, trie,
        et \'ecris un CSV propre avec BOM UTF-8. 100 \% g\'en\'erateur, m\'emoire constante.
  \item \verb|trouverDoublonsProches(array \$noms): array|~: normalise (minuscules, accents,
        ponctuation) puis regroupe par empreinte~; retourne uniquement les groupes de taille \$>1\$.
  \item \verb|prochainesDates(string \$regle, DateTimeImmutable \$debut, int \$n): array|~: sans
        biblioth\`eque~: accepte \verb|chaque lundi|, \verb|1er de chaque mois|, \verb|+3 mois|,
        en utilisant \verb|modify| et \verb|DateInterval|.
  \item \'Ecris un \verb|Tableau| utilitaire avec les m\'ethodes \verb|map|, \verb|filter|,
        \verb|groupBy|, \verb|chunk|, \verb|pluck|, \verb|sortBy|, \verb|unique|, \verb|first|,
        \verb|flatten|~: c'est exactement ce que fait \verb|Illuminate\Support\Collection|,
        et le refaire une fois est la meilleure fa\c{c}on de le comprendre ensuite.
  \item Analyse de regex~: rends \verb|/(\w+)@\1/| et \verb|/(a+)+b/| inoffensives en ajoutant
        les limites (limite de backtracking via \verb|pcre.backtrack\_limit|) et explique le co\^ut
        de backtracking sur \verb|aaaaaaaaaaaaaaaaaaaaaX|.
\end{enumerate}
\end{exercice}

\begin{ressources}
\begin{itemize}
  \item Manuel --- Tableaux --- \url{https://www.php.net/manual/fr/ref.array.php}
  \item Manuel --- Fonctions de cha\^ines --- \url{https://www.php.net/manual/fr/ref.strings.php}
  \item Manuel --- Dates et heures --- \url{https://www.php.net/manual/fr/book.datetime.php}
  \item PCRE (regex) --- \url{https://www.php.net/manual/en/reference.pcre.pattern.syntax.php}
  \item \link{Collection (Laravel)}{https://laravel.com/docs/collections} --- la reference a lire apres l'exercice 5
\end{itemize}
\end{ressources}

\index{PHP!tableaux}\index{PHP!chaines}\index{regex}\index{dates}
% ============================================================================
\chapter{PHP objet~: classes, interfaces, \'enum\'erations, attributs}
\label{ch:php-oop}
% ============================================================================

\begin{objectif}
\begin{itemize}
  \item \'Ecrire des classes modernes~: constructeur promotionnel, \verb|readonly|, typage complet.
  \item Choisir entre interface, classe abstraite, trait et \'enum\'eration, et savoir pourquoi.
  \item Comprendre l'autochargement (PSR-4) et l'injection de d\'ependances sans framework.
  \item Utiliser les attributs pour d\'ecrire au lieu de configurer ailleurs.
  \item Construire un petit noyau d'application (container, routeur, repository) en 200 lignes.
\end{itemize}
\end{objectif}

\section{L'objet PHP moderne, en un exemple}

\begin{codephp}
<?php
declare(strict_types=1);

namespace App\Domain\Article;

use DateTimeImmutable;
use Stringable;

/**
 * Un article. Entite riche : les regles metier vivent ICI, pas dans les controleurs.
 */
final class Article implements Stringable
{
    // Constantes de classe : valeurs immuables partagees
    public const LONGUEUR_TITRE_MIN = 3;
    public const LONGUEUR_TITRE_MAX = 160;

    /** @var string[] */
    private array $tags;

    public function __construct(
        private readonly ?int $id,
        private string $titre,
        private string $corps,
        private Statut $statut,                       // enum : pas de chaine magique
        private readonly Auteur $auteur,
        private readonly DateTimeImmutable $creeLe,
        private ?DateTimeImmutable $publieLe = null,
        array $tags = [],
    ) {
        $this->titre = trim($titre);
        $this->validerTitre();
        $this->tags = array_values(array_unique(array_map(
            static fn (string $t): string => mb_strtolower(trim($t)),
            $tags
        )));
        if ($statut === Statut::Publie && $publieLe === null) {
            throw new InvalidArgumentException('Un article publie doit avoir une date de publication.');
        }
    }

    // --- Accesseurs : lecture seule. L'etat ne change que par des methodes qui ont un SENS. ---
    public function id(): ?int { return $this->id; }
    public function titre(): string { return $this->titre; }
    public function statut(): Statut { return $this->statut; }
    public function estPublie(): bool { return $this->statut === Statut::Publie; }

    /** @return string[] */
    public function tags(): array { return $this->tags; }

    // --- Transitions d'etat : elles retournent une NOUVELLE instance (immutabilite) ---
    public function publier(?DateTimeImmutable $quand = null): self
    {
        if ($this->statut === Statut::Archive) {
            throw new LogicException('Un article archive ne se republicie pas.');
        }
        return new self(
            $this->id, $this->titre, $this->corps, Statut::Publie, $this->auteur,
            $this->creeLe, $quand ?? new DateTimeImmutable(), $this->tags,
        );
    }

    public function renommer(string $nouveauTitre): self
    {
        $copie = clone $this;
        $copie->titre = trim($nouveauTitre);
        $copie->validerTitre();
        return $copie;
    }

    public function __toString(): string { return sprintf('#%s %s', $this->id ?? 'nouveau', $this->titre); }

    public function egal(Article $autre): bool
    {
        return $this->id !== null && $this->id === $autre->id;   // identite, pas similarite
    }

    private function validerTitre(): void
    {
        $longueur = mb_strlen($this->titre);
        if ($longueur < self::LONGUEUR_TITRE_MIN || $longueur > self::LONGUEUR_TITRE_MAX) {
            throw new DomainException(sprintf(
                'Le titre doit faire entre %d et %d caracteres (%d donne).',
                self::LONGUEUR_TITRE_MIN, self::LONGUEUR_TITRE_MAX, $longueur
            ));
        }
    }
}

// Un objet valeur : petit, immuable, comparable par valeur
final class Auteur
{
    public function __construct(public readonly string $nom, public readonly string $email)
    {
        if (!filter_var($email, FILTER_VALIDATE_EMAIL)) { throw new DomainException("email invalide : {$email}"); }
    }
    public function equals(self $autre): bool
    {
        return $this->nom === $autre->nom && $this->email === $autre->email;
    }
}

enum Statut: string
{
    case Brouillon = 'brouillon';
    case Publie = 'publie';
    case Archive = 'archive';

    public function libelle(): string
    {
        return match ($this) {
            self::Brouillon => 'Brouillon',
            self::Publie => 'Publie',
            self::Archive => 'Archive',
        };
    }
    public function peutEditer(): bool { return $this !== self::Archive; }
}
\end{codephp}

\subsection{Ce que le mot-cl\'e \texttt{final} change r\'eellement}

Marquer \verb|final| toute classe qui n'est pas faite pour \^etre \'etendue est la d\'ecision
d'architecture la rentable du PHP moderne~: tu gagnes le droit de changer l'impl\'ementation
sans casser des h\'eritiers inconnus. R\`egle simple~: \verb|final| partout, \verb|abstract| ou
interface quand tu con\c{c}ois \emph{volontairement} un point d'extension.

\begin{tabz}{llX}
  \headrow \thd{Modificateur} \& \thd{Depuis} \& \thd{Effet} \\
  \midrule
  \verb|final class| \& PHP 4/8 \& interdit d'\'etendre (et \verb|final public function| d'\'ecraser). \\
  \verb|readonly class| \& 8.2 \& toutes les propri\'et\'es sont \verb|readonly|~; interdiction d'h\'eriter. \\
  \verb|readonly public int \$x| \& 8.1 \& assign\'e une seule fois, typ\'e, non modifiable apr\`es. \\
  \verb|public private(set) int \$x| \& 8.4 \& lisible partout, \'ecriture r\'eserv\'ee \`a la classe. \\
  \verb|abstract| \& --- \& classe/m\'ethode inachev\'ee~; ne peut pas \^etre instanci\'ee. \\
  \verb|static| \& --- \& appartient \`a la classe, pas \`a l'instance. \`A utiliser avec parcimonie. \\
  \verb|clone| \& --- \& copie superficielle~; \verb|\_\_clone()| pour corriger les r\'ef\'erences partag\'ees. \\
  \verb|unset()| sur une propri\'et\'e \& --- \& la fait revenir \`a \og ~non initialis\'ee~ \fg{} (erreur si on la lit). \\
\end{tabz}

\section{H\'eritage, interfaces, traits~: quand quoi}

\begin{codephp}
<?php
interface DepotArticles                       // CONTRAT : ce que l'appelant peut faire
{
    public function trouver(int $id): ?Article;
    /** @return Article[] */
    public function recherche(?Statut $statut, int $limite = 50, int $decalage = 0): array;
    public function enregistrer(Article $a): Article;
}

interface HorizonTTL { public function expireLe(): ?DateTimeImmutable; }

abstract class Depot implements DepotArticles, HorizonTTL     // CODE PARTAGE + contrat
{
    public function __construct(protected \PDO $pdo) {}

    final public function trouver(int $id): ?Article       // non ecrasable : invariants garantis
    {
        $stmt = $this->pdo->prepare('SELECT * FROM articles WHERE id = :id LIMIT 1');
        $stmt->execute(['id' => $id]);
        $ligne = $stmt->fetch(PDO::FETCH_ASSOC);
        return $ligne === false ? null : $this->hydrater($ligne);      // methode abstraite -> la suite
    }

    abstract protected function hydrater(array $ligne): Article;
    public function expireLe(): ?DateTimeImmutable { return null; }
}

final class DepotSql extends Depot
{
    protected function hydrater(array $l): Article
    {
        return new Article(
            (int) $l['id'], $l['titre'], (string) $l['corps'], Statut::from($l['statut']),
            new Auteur($l['auteur_nom'], $l['auteur_email']),
            new DateTimeImmutable($l['cree_le']),
            $l['publie_le'] !== null ? new DateTimeImmutable($l['publie_le']) : null,
            json_decode((string) $l['tags'], true) ?: [],
        );
    }
    /* ... */
}

// Trait : REUTILISATION de code, pas un type. Pas d'etat partag\'e, pas de hierarchie.
trait InitialiseTableau
{
    private ?array $tampon = null;
    protected function tampon(): array { return $this->tampon ??= []; }
    protected function pousser(mixed $v): void { $this->tampon()[] = $v; }
}

// Conflit de traits : au choix, renommer ou exclure
class A { use Trait1, Trait2 { Trait1::auLieu insteadof Trait2; Trait2::methode as methodeB; } }
\end{codephp}

\begin{description}
  \item[Interface] \og ~je sais faire~  \fg{}. Z\'ero code, plusieurs impl\'ementations~: mock en test,
        MySQL/Postgres/InMemory, notificateur email/SMS. \`A privil\'egier.
  \item[Classe abstraite] \og ~nous sommes de la m\^eme famille, et voici le code commun~  \fg{}.
        Peut avoir de l'\'etat. Utilise-la quand tu factorises de l'impl\'ementation, pas pour imposer un type.
  \item[Trait] \og ~voici un morceau de comportement~  \fg{}. Un plan B syntaxique~;
        \'evite-le d\`es qu'il commence \`a avoir des accesseurs abstraits (signe qu'il lui manque un
        collaborateur~: injecte-le).
  \item[Enum] un ensemble \emph{ferm\'e} de valeurs nomm\'ees. Elle remplace
        \verb|const|+\verb|switch|+\verb|string| avec des m\'ethodes et une s\'ecurit\'e de type.
\end{description}

\section{\'Enum\'erations~: le plus sous-utilis\'e des mots-cl\'es}

\begin{codephp}
<?php
enum Role: string
{
    case Lecteur = 'lecteur';
    case Auteur = 'auteur';
    case Redacteur = 'redacteur';
    case Admin = 'admin';

    public const ORDRE = [self::Lecteur->value, self::Auteur->value, self::Redacteur->value, self::Admin->value];

    public function peut(?self $autre): bool
    {
        return array_search($this->value, self::ORDRE, true) >= array_search(($autre?->value ?? 'x'), self::ORDRE, true);
    }
    public static function depuisHttp(string $v): self
    {
        return self::tryFrom($v) ?? throw new InvalidArgumentException("role inconnu : {$v}");
    }
}
$role = Role::Admin;
$role->name;   // 'Admin'        $role->value : 'admin'
Role::from('redacteur');           // une ValueError si inconnu
Role::tryFrom('xx');                // null
array_map(fn (Role $r) => $r->value, Role::cases());

// Enum sans valeur = "union de types d'evenement", le motif parfait d'un bus
abstract class Evenement { }
enum TypeEvenement { case ArticlePublie; case CommandePayee; case UtilisateurInscrit; }
interface Traiteur { public function traite(TypeEvenement $type, object $payload): void; }
\end{codephp}

\begin{note}
En base de donn\'ees~: une colonne \verb|enum('brouillon','publie','archive')| (MySQL) ou
\verb|CREATE TYPE statut AS ENUM ...| (Postgres) donne une contrainte d'int\'egrit\'e gratuite.
C\^ot\'e PHP, \verb|Statut::from()| dans l'hydratation transforme une donn\'ee corrompue en exception
au lieu d'un \'etat absurde qui se propage pendant trois \'ecrans.
\end{note}

\section{Attributs~: d\'eclarer au lieu de configurer ailleurs}

\begin{codephp}
<?php
#[Attribute(Attribute::TARGET_PROPERTY | Attribute::TARGET_METHOD | Attribute::IS_REPEATABLE)]
final class Valide
{
    /** @param string[] $regles */
    public function __construct(public readonly array $regles, public readonly ?string $champ = null) {}
}

#[Attribute(Attribute::TARGET_METHOD)]
final class Route
{
    public function __construct(public readonly string $methode, public readonly string $chemin) {}
}

final class ControleurArticles
{
    public function __construct(private readonly DepotArticles $depots) {}

    #[Route('GET', '/articles/{id:\d+}')]
    #[Valide(['integer', 'min:1'], champ: 'id')]
    public function lire(Request $r, Responder $res): Response
    {
        $article = $this->depots->trouver((int) $r->params['id']);
        return $article === null ? $res->notFound() : $res->json($article->versTableau());
    }
}

// Lecture des attributs (le role d'un framework, ou ton propre micro-scanner)
$ref = new ReflectionMethod(ControleurArticles::class, 'lire');
foreach ($ref->getAttributes(Route::class) as $a) {
    $route = $a->newInstance();
    printf("%s %s -> %s::%s\n", $route->methode, $route->chemin, $ref->class, $ref->name);
}
foreach ($ref->getAttributes(Valide::class) as $a) {
    $v = $a->newInstance();
    /* appliquer $v->regles au champ $v->champ */
}
\end{codephp}

Les attributs sont exactement ce qu'utilisent Symfony (routes, validation, s\'erialisation),
Doctrine (\verb|\#[Table]|, \verb|\#[ORM\Id]|), PHPUnit (\verb|\#[Test]|, \verb|\#[DataProvider]|)
et Laravel. Les avoir \'ecrits une fois \og ~\`a la main~ \fg{} rend ces frameworks limpides.

\section{Autochargement~: PSR-4 sans myst\`ere}

Un fichier = une classe = un nom d'espace de lignes coh\'erent avec le chemin. Le reste est
l'affaire de Composer (\cref{ch:php-composer}) ou de dix lignes de \verb|spl\_autoload\_register|.

\begin{codephp}
<?php
// composer.json -> "autoload": { "psr-4": { "App\\": "src/" } }
// App\Domain\Article\Article   =>   src/Domain/Article/Article.php

// Version manuelle, pour comprendre ce que fait Composer :
spl_autoload_register(static function (string $classe): void {
    $prefixe = 'App\\';
    if (!str_starts_with($classe, $prefixe)) { return; }
    $chemin = __DIR__ . '/src/' . str_replace('\\', '/', substr($classe, strlen($prefixe))) . '.php';
    if (is_file($chemin)) { require $chemin; }
});

// Un autoloader maison est une tres bonne exercice ; en production, installe Composer,
// et gen\`ere l'optimise : composer dump-autoload -o  (une table classe -> fichier, pas de stat disque)
\end{codephp}

\section{Injection de d\'ependances~: le patron qui rend tout testable}

\begin{codephp}
<?php
// Mauvais : le controleur FABRIQUE ses dependances -> inconciliable avec un test, un mock, un swap.
final class ControleurArticlesMauvais
{
    public function __construct()
    {
        $this->depot = new DepotSql(new PDO('mysql:host=localhost', 'root', ''));  // secret + couplage
    }
}

// Bon : on recoit des ABSTRACTIONS (interfaces). Le "qui cree quoi" est ailleurs : le conteneur.
final class ControleurArticles
{
    public function __construct(
        private readonly DepotArticles $depots,
        private readonly Horloge $horloge,
        private readonly Journal $journal,
    ) {}

    public function publier(int $id): Response
    {
        $article = $this->depots->trouver($id) ?? throw new NotFound('article');
        $this->depots->enregistrer($article->publier($this->horloge->maintenant()));
        $this->journal->info('article publie', ['id' => $id]);
        return new Response(204);
    }
}

// Un conteneur minimal (30 lignes) : c'est le coeur de Symfony, Laravel, PHP-DI.
final class Conteneur
{
    /** @var array<string, callable> */ private array $definitions = [];
    /** @var array<string, object> */   private array $instances = [];

    public function set(string $id, callable $fabrique): void { $this->definitions[$id] = $fabrique; }

    public function get(string $id): object
    {
        if (isset($this->instances[$id])) { return $this->instances[$id]; }
        if (!isset($this->definitions[$id])) {
            // resolution reflexive par type : on construit les arguments des constructeurs
            return $this->instances[$id] = $this->autowire($id);
        }
        return $this->instances[$id] = ($this->definitions[$id])($this);
    }

    private function autowire(string $classe): object
    {
        $ref = new ReflectionClass($classe);
        $constructeur = $ref->getConstructor();
        $args = [];
        foreach ($constructeur?->getParameters() ?? [] as $p) {
            $type = $p->getType();
            $args[] = $type instanceof ReflectionNamedType && !$type->isBuiltin()
                ? $this->get($type->getName())
                : ($p->getDefaultValue() ?? throw new RuntimeException("impossible de resoudre {$classe}::\${$p->getName()}"));
        }
        return $ref->newInstanceArgs($args);
    }
}
$c = new Conteneur();
$c->set(PDO::class, fn () => new PDO($_ENV['DATABASE_URL'], options: [PDO::ATTR_ERRMODE => PDO::ERRMODE_EXCEPTION]));
$c->set(DepotArticles::class, fn (Conteneur $c) => new DepotSql($c->get(PDO::class)));
$controleur = $c->get(ControleurArticles::class);   // tout le graphe est construit, une fois chacun
\end{codephp}

\begin{astuce}
Trois habitudes qui font 90 \% du \og ~code propre~ \fg{} en PHP~: (1) d\'ependre d'interfaces, pas
d'impl\'ementations~; (2) ne jamais utiliser \verb|new| \`a l'int\'erieur d'une m\'ethode m\'etier pour une
collaboration~; (3) rendre les objets immuables par d\'efaut. Le reste (Design Patterns, SOLID,
DDD) est de la broderie sur ces trois points, utile \`a conna\^itre mais rarement \`a invoquer.
\end{astuce}

\subsection{La bo\^ite \`a outils d'objets utiles}

\begin{tabz}{lXl}
  \headrow \thd{\'El\'ement} \& \thd{Usage} \& \thd{Lien} \\
  \midrule
  \verb|spl\_object\_id| \& identifiant unique d'une instance, pour debug. \& --- \\
  \verb|WeakMap| / \verb|WeakReference| \& m\'etadonn\'ees attach\'ees \`a un objet, sans fuite m\'emoire. \& \url{https://www.php.net/manual/en/class.weakmap.php} \\
  \verb|SplObjectStorage| \& ensemble d'objets compar\'es par identit\'e. \& --- \\
  \verb|ArrayObject| / \verb|ArrayIterator| \& tableau avec interface objet (countable, iterable). \& --- \\
  \verb|IteratorAggregate| \& rendre ta collection it\'erable avec \verb|foreach|. \& --- \\
  \verb|Countable|, \verb|JsonSerializable|, \verb|Stringable| \& petites interfaces qui branchent le langage sur ton objet. \& --- \\
  \verb|\_\_invoke| \& rendre un objet appelable (\verb|\$obj()|)~: le mod\`ele des middlewares. \& --- \\
  \verb|Serializable| (d\'epr\'eci\'e) / \verb|\_\_serialize| + \verb|\_\_unserialize| \& contr\^ole de la s\'erialisation. \& --- \\
  \verb|Closure::bind| / \verb|Closure::fromCallable| \& ex\'ecuter du code avec un \verb|\$this| pr\^et\'e. \& --- \\
\end{tabz}

\begin{exercice}
\begin{enumerate}
  \item \'Ecris le mod\`ele complet d'un blog~: \verb|Article|, \verb|Commentaire|, \verb|Auteur|,
        \verb|Statut| (enum), \verb|DepotArticles| (interface), \verb|DepotInMemory| et
        \verb|DepotSql| (les deux impl\'ementations du m\^eme contrat), et un test qui passe
        sur les deux sans changer une ligne.
  \item Ajoute un objet valeur \verb|Email|, un \verb|Montant| (entier en centimes, op\'erations
        typ\'ees, \verb|\_\_toString|) et un \verb|Identifiant| (UUIDv4). Interdis toute cha\^ine
        \og ~brute~ \fg{} dans tes signatures.
  \item \'Ecris une classe \verb|Resultat| (succ\`es ou \'echec)~: \verb|Resultat::ok(\$v)|,
        \verb|Resultat::erreur(\$e)|, \verb|onSuccess()|, \verb|recover()|. Puis utilise-la pour
        la validation d'un formulaire, sans exceptions.
  \item Impl\'emente une \verb|Collection| immuable g\'en\'erique-ish~: \verb|map|, \verb|filter|,
        \verb|reduce|, \verb|groupBy|, \verb|first|, \verb|isEmpty|, \verb|IteratorAggregate|,
        \verb|Countable|, \verb|JsonSerializable|.
  \item \'Ecris un petit routeur avec attributs~: \verb|\#[Route('GET','/articles/{id}')]~,
        un scanner par r\'eflexion, compilation des motifs en regex, et r\'esolution des param\`etres.
  \item Ajoute un conteneur DI \`a ton projet (le mod\`ele ci-dessus suffit), et fais-en sorte que
        \verb|new| n'apparaisse plus nulle part dans \verb|src/| sauf dans le bootstrap.
        C'est le d\'eclic \og ~framework~ \fg{}~: ensuite Symfony ou Laravel te semblera \'evident.
\end{enumerate}
\end{exercice}

\begin{ressources}
\begin{itemize}
  \item Manuel --- Objets et classes --- \url{https://www.php.net/manual/fr/language.oop5.php}
  \item PHP The Right Way (sections autoloading, DI, standards) --- \url{https://phptherightway.com/}
  \item PSR-4 --- \url{https://www.php-fig.org/psr/psr-4/} --- PSR-11 (conteneur) --- \url{https://www.php-fig.org/psr/psr-11/}
  \item PHP-DI (le conteneur autonome de r\'ef\'erence) --- \url{https://php-di.org/}
  \item \link{PHPStan (analyse statique)}{https://phpstan.org/} --- \url{https://phpstan.org/blog/how-to-start-typing}
\end{itemize}
\end{ressources}

\index{PHP!objets}\index{interfaces}\index{enum}\index{attributs}\index{DI}
% ============================================================================
\chapter{HTTP en PHP~: requ\^etes, formulaires, sessions, authentification}
\label{ch:php-http}
% ============================================================================

\begin{objectif}
\begin{itemize}
  \item Lire une requ\^ete proprement~: m\'ethodes, en-t\^etes, \verb|\$\_GET|, \verb|\$\_POST|, corps JSON.
  \item Envoyer une r\'eponse correcte~: codes, en-t\^etes, redirections, streaming.
  \item Construire un formulaire robuste~: valeurs conserv\'ees, erreurs, protection CSRF.
  \item G\'erer une session et un mot de passe de bout en bout~: connexion, d\'econnexion,
        \og ~mot de passe oubli\'e~ \fg{}, r\'eg\'en\'eration d'identifiant.
  \item Uploader un fichier sans ouvrir de faille, et envoyer un email.
\end{itemize}
\end{objectif}

\section{Requ\^ete et r\'eponse~: le contrat}

\begin{codephp}
<?php
declare(strict_types=1);

// ---- LIRE LA REQUETE
$modele = $_SERVER['REQUEST_METHOD'];                       // 'GET' | 'POST' | 'PUT' ...
$uri = parse_url($_SERVER['REQUEST_URI'] ?? '/', PHP_URL_PATH);
$query = $uri;
$entete = fn (string $nom): ?string => $_SERVER['HTTP_' . strtoupper(str_replace('-', '_', $nom))] ?? null;
$contentType = $entete('content-type') ?? '';
$origine = $entete('origin');
$ip = $_SERVER['REMOTE_ADDR'] ?? 'inconnue';              // avec un proxy : X-Forwarded-For, a verifier
$estAjax = str_contains((string) $entete('x-requested-with'), 'XMLHttpRequest');

// Les superglobales, et ce qu'il faut en savoir
$_GET;              // chaine uniquement ; attention : tableau associatif brut, non filtre
$_POST;             // idem pour application/x-www-form-urlencoded et multipart/form-data
$_FILES;            // upload (voir plus bas)
$_COOKIE;           // cookies lus
$_SESSION;          // voir session
$_SERVER;           // meta : METHODE, URI, HTTP_*, DOCUMENT_ROOT, REQUEST_TIME_FLOAT
$_REQUEST;          // MELANGE des trois precedents selon variables_order : NE L'UTILISE PAS
$_ENV;              // seulement si variables_order contient "G" ; sinon : getenv() ou les constants

// Corps JSON (les API modernes ne passent plus par $_POST)
$brut = file_get_contents('php://input');
$json = $contentType !== '' && str_contains($contentType, 'json')
    ? json_decode($brut ?: 'null', true, 512, JSON_THROW_ON_ERROR)
    : null;
$entrees = $json ?? $_POST;

// Lecture securisee + validation en un point : mon petit "getter" de controleur
function champ(array $source, string $cle, ?string $defaut = null): ?string
{
    $v = $source[$cle] ?? $defaut;
    return is_string($v) ? trim($v) : ($v === null ? null : (string) $v);
}
function champInt(array $source, string $cle, ?int $defaut = null): ?int
{
    $v = $source[$cle] ?? null;
    return is_numeric($v) ? (int) $v : $defaut;
}
function champBool(array $source, string $cle): bool
{
    return in_array($source[$cle] ?? null, ['1', 1, 'on', 'true', true], true);
}
function filtre(array $source, string $cle, int $filtrePhp, array $options = []): mixed
{
    return filter_input_array($source === $_GET ? INPUT_GET : INPUT_POST, [$cle => $options ?: $filtrePhp])[$cle] ?? false;
}
$page = max(1, champInt($_GET, 'page', 1));
$email = filter_var(champ($entrees, 'email') ?? '', FILTER_VALIDATE_EMAIL);   // false si invalide
$urls = filter_var($entrees['site'] ?? '', FILTER_VALIDATE_URL);
$id = filter_var($_GET['id'] ?? null, FILTER_VALIDATE_INT, ['options' => ['min_range' => 1]]);

// ---- ENVOYER LA REPONSE
http_response_code(201);
header('content-type: application/json; charset=utf-8');
header('cache-control: private, no-store');
header('x-page: ' . $page);
header_link: header('link: </style.css>; rel=preload; as=style', false);   // false = ne pas remplacer
header('retry-after: 60');
echo json_encode(['ok' => true, 'page' => $page], JSON_UNESCAPED_UNICODE | JSON_THROW_ON_ERROR);
exit;                                            // apres une redirection : TOUJOURS exit

// Redirections : 303 apres un POST (modele POST/REDIRECT/GET), 301 si c'est definitif
function rediriger(string $url, int $code = 303): never
{
    header('location: ' . $url, true, $code);
    exit;
}
function telecharger(string $chemin, string $nom): never
{
    header('content-type: text/csv; charset=utf-8');
    header('content-disposition: attachment; filename="' . addslashes($nom) . '"');
    header('content-length: ' . (string) filesize($chemin));
    readfile($chemin);
    exit;
}
\end{codephp}

\begin{piege}
\verb|header()| \'echoue avec \og ~headers already sent~ \fg{} d\`es qu'un octet a \'et\'e \'emis~: un espace
avant \verb|<?php|, un saut de ligne apr\`es un \verb|?>|, un \verb|echo| de debug. Les
r\`egles qui \'eliminent ce bug pour toujours~: fichiers PHP \emph{sans} balise fermante finale,
\verb|ob\_start()| dans le bootstrap si tu dois bufferiser, et aucune sortie avant la logique.
\end{piege}

\section{Un formulaire qui se tient~: le motif complet}

\begin{codephp}
<?php
// src/Action/Inscription.php  ---  un seul fichier, trois cas : affichage, rejet, succes.
declare(strict_types=1);

session_start();
$csrf = $_SESSION['csrf'] ??= bin2hex(random_bytes(32));
$valeurs = ['pseudo' => '', 'email' => '', 'message' => '', 'newsletter' => '1'];
$erreurs = [];

if ($_SERVER['REQUEST_METHOD'] === 'POST') {
    // 1) protection CSRF : comparaison du jeton recu avec celui de la session
    if (!hash_equals($csrf, (string) ($_POST['_csrf'] ?? ''))) {
        http_response_code(419);
        exit('Formulaire expire, recharge la page.');
    }
    // 2) recuperer et valider
    $valeurs['pseudo'] = champ($_POST, 'pseudo') ?? '';
    $valeurs['email'] = mb_strtolower(champ($_POST, 'email') ?? '');
    $valeurs['message'] = champ($_POST, 'message') ?? '';
    $valeurs['newsletter'] = isset($_POST['newsletter']) ? '1' : '';

    if (mb_strlen($valeurs['pseudo']) < 3) { $erreurs['pseudo'] = 'Au moins 3 caracteres.'; }
    if ($valeurs['email'] === '' || filter_var($valeurs['email'], FILTER_VALIDATE_EMAIL) === false) {
        $erreurs['email'] = 'Adresse email invalide.';
    }
    if (mb_strlen($valeurs['message']) > 1000) { $erreurs['message'] = '1000 caracteres maximum.'; }

    if ($erreurs === []) {
        // 3) ecrire (requete preparee : voir chapitre base de donnees)
        $stmt = $pdo->prepare('INSERT INTO inscriptions (pseudo, email, message, newsletter, cree_le) VALUES (?, ?, ?, ?, NOW())');
        $stmt->execute([$valeurs['pseudo'], $valeurs['email'], $valeurs['message'], $valeurs['newsletter'] === '1' ? 1 : 0]);
        $_SESSION['flash']['succes'] = 'Merci, c\'est enregistre.';
        rediriger('/inscription');           // PRG : un rafraichissement ne re-soumet pas
    }
    // sinon : on retombe dans le rendu avec $valeurs et $erreurs remplis
}

function e(?string $v): string { return htmlspecialchars((string) $v, ENT_QUOTES | ENT_SUBSTITUTE, 'UTF-8'); }
?>
<?php if (isset($_SESSION['flash']['succes'])): ?>
  <p class="succes"><?= e($_SESSION['flash']['succes']) ?></p>
  <?php unset($_SESSION['flash']['succes']); ?>
<?php endif; ?>

<form method="post" action="/inscription" novalidate enctype="application/x-www-form-urlencoded">
  <input type="hidden" name="_csrf" value="<?= e($csrf) ?>">

  <label for="pseudo">Pseudo</label>
  <input id="pseudo" name="pseudo" required minlength="3" maxlength="40"
         value="<?= e($valeurs['pseudo']) ?>"
         <?php if (isset($erreurs['pseudo'])): ?>aria-invalid="true" aria-describedby="err-pseudo"<?php endif; ?>>
  <?php if (isset($erreurs['pseudo'])): ?><p id="err-pseudo" class="erreur"><?= e($erreurs['pseudo']) ?></p><?php endif; ?>

  <label for="email">Email</label>
  <input id="email" name="email" type="email" required value="<?= e($valeurs['email']) ?>">
  <?php if (isset($erreurs['email'])): ?><p class="erreur"><?= e($erreurs['email']) ?></p><?php endif; ?>

  <label for="message">Message <small>(<?= mb_strlen($valeurs['message']) ?>/1000)</small></label>
  <textarea id="message" name="message" rows="6"><?= e($valeurs['message']) ?></textarea>

  <label><input type="checkbox" name="newsletter" value="1" <?= $valeurs['newsletter'] ? 'checked' : '' ?>> Newsletter</label>

  <button type="submit">Envoyer</button>
</form>
\end{codephp}

\begin{enumerate}
  \item \textbf{Toujours} renvoyer les valeurs saisies (sinon l'utilisateur retape tout) et
        ne jamais r\'eafficher un mot de passe.
  \item \textbf{Toujours} \'echapper \`a la sortie~: \verb|htmlspecialchars| avec
        \verb|ENT\_QUOTES| et UTF-8. Une vue qui n'\'echappe pas est une faille XSS.
  \item \textbf{Toujours} faire un POST/REDIRECT/GET apr\`es un envoi r\'eussi.
  \item \textbf{Toujours} valider c\^ot\'e serveur, m\^eme si le HTML a \verb|required|~: le HTML se
        contourne en dix secondes avec DevTools ou \verb|curl|.
  \item \textbf{Jamais} faire confiance \`a \verb|\$\_FILES['type']|, \`a la longueur annonc\'ee, ni \`a
        un \verb|name|~: tout cela vient du client.
\end{enumerate}

\section{Sessions~: le stockage c\^ot\'e serveur}

\begin{codephp}
<?php
// php.ini (ou ini_set avant session_start, ou dans un pool php-fpm) : les reglages qui comptent
ini_set('session.use_strict_mode', '1');              // rejette un ID de session fourni par l'attaquant
ini_set('session.cookie_httponly', '1');              // pas accessible en JS
ini_set('session.cookie_secure', '1');                 // HTTPS uniquement
ini_set('session.cookie_samesite', 'Lax');
ini_set('session.use_only_cookies', '1');              // jamais d'ID dans l'URL
ini_set('session.gc_maxlifetime', '1800');             // 30 min d'inactivite
ini_set('session.sid_length', '48'); ini_set('session.sid_bits_per_character', '6');

session_start();                                       // cree ou reprend la session
$_SESSION['user_id'] = 7;                              // uniquement des donnees petites et serialisables
$uid = $_SESSION['user_id'] ?? null;
session_regenerate_id(true);                            // APRES toute connexion/escalade : anti-fixation
session_destroy();                                      // suppression cote serveur
setcookie('panier', '', time() - 3600, '/');             // effacer un cookie = le reemettre expire
$nom = session_name();                                   // 'PHPSESSID' par defaut ; renomme-le (minorite d'infos)

// Session en base (multiprobe : plusieurs serveurs php-fpm) : SessionHandlerInterface
final class SessionBdd implements SessionHandlerInterface
{
    public function __construct(private readonly PDO $pdo, private readonly int $duree = 1800) {}
    public function open(string $path, string $name): bool { return true; }
    public function close(): bool { return true; }
    public function read(string $id): string|false
    {
        $s = $this->pdo->prepare('SELECT data FROM sessions WHERE id = ? AND expire > NOW()');
        $s->execute([$id]);
        return (string) ($s->fetchColumn() ?: '');
    }
    public function write(string $id, string $data): int|bool
    {
        $s = $this->pdo->prepare('REPLACE INTO sessions (id, data, expire) VALUES (?, ?, NOW() + INTERVAL ? SECOND)');
        return $s->execute([$id, $data, $this->duree]);
    }
    public function destroy(string $id): int|bool { return $this->pdo->prepare('DELETE FROM sessions WHERE id = ?')->execute([$id]); }
    public function gc(int $max_lifetime): int|false { return $this->pdo->exec('DELETE FROM sessions WHERE expire < NOW()'); }
}
(new SessionHandler())->... ; // en pratique : session_set_save_handler(new SessionBdd($pdo), true);
\end{codephp}

\begin{note}
\textbf{Ce qu'on ne met jamais en session~:} un gros tableau (la session est s\'erialis\'ee et
relue \`a chaque requ\^ete), un secret non chiffr\'e, un r\^ole recalculable~; et surtout pas un jeton
d'API dans une session stock\'ee en base. Ce qu'on y met~: l'identifiant utilisateur, un
\verb|csrf|, un \'etat d'\'etapes, des flash.
\end{note}

\section{Authentification compl\`ete~: le parcours}

\begin{codephp}
<?php
declare(strict_types=1);

final class Authentification
{
    public function __construct(
        private readonly PDO $pdo,
        private readonly int $dureeVerrou = 900,
        private readonly int $essaisMax = 5,
    ) {}

    /** Inscription : verifie la force, hache, ecrit. Reponse identique que l'email existe ou non. */
    public function inscrire(string $email, string $motDePasse, string $pseudo): void
    {
        $this->verifierForce($motDePasse);
        $hash = password_hash($motDePasse, PASSWORD_ARGON2ID, [
            'memory_cost' => 1 << 15, 'time_cost' => 3, 'threads' => 2,   // ~0,3 s : assez lent pour gener, assez rapide pour toi
        ]);
        $s = $this->pdo->prepare(
            'INSERT INTO utilisateurs (email, mot_de_passe, pseudo, cree_le) VALUES (?, ?, ?, NOW())
             ON DUPLICATE KEY UPDATE id = id'     ;
        );
        $s->execute([mb_strtolower($email), $hash, $pseudo]);
    }

    public function connecter(string $email, string $motDePasse, ?string $ip = null): ?Utilisateur
    {
        $email = mb_strtolower(trim($email));
        $s = $this->pdo->prepare('SELECT id, email, mot_de_passe, pseudo, verrou_jusqua FROM utilisateurs WHERE email = ? LIMIT 1');
        $s->execute([$email]);
        $u = $s->fetch(PDO::FETCH_ASSOC);

        // 1) reponse au meme rythme et au meme message que si l'email etait inconnu
        if ($u === false) {
            password_verify($motDePasse, '$argon2id$v=19$m=65536,t=3,p=2$AAAAAAAAAAAAAAAAAAAAAA$AAAAAAAAAAAAAAAAAAAAAAAAAAAAAAAAAAAAAAAAAAA');
            throw new ErreurAuthentification('Identifiants incorrects.');
        }
        if ($u['verrou_jusqua'] !== null && strtotime((string) $u['verrou_jusqua']) > time()) {
            throw new ErreurAuthentification('Trop d\'essais. Reessaie plus tard.');
        }
        if (!password_verify($motDePasse, (string) $u['mot_de_passe'])) {
            $this->marquerEchec((int) $u['id'], $ip);
            throw new ErreurAuthentification('Identifiants incorrects.');
        }
        // 2) re-hachage si l'algorithme ou les couts ont change : la maintenance de securite offerte gracieusement
        if (password_needs_rehash((string) $u['mot_de_passe'], PASSWORD_ARGON2ID,
            ['memory_cost' => 1 << 15, 'time_cost' => 3, 'threads' => 2])) {
            $n = password_hash($motDePasse, PASSWORD_ARGON2ID, ['memory_cost' => 1 << 15, 'time_cost' => 3, 'threads' => 2]);
            $this->pdo->prepare('UPDATE utilisateurs SET mot_de_passe = ? WHERE id = ?')->execute([$n, $u['id']]);
        }
        // 3) ouverture de session securisee
        session_regenerate_id(true);
        $_SESSION = ['user_id' => (int) $u['id'], 'ip' => $ip, 'agent' => hash('sha256', (string) ($_SERVER['HTTP_USER_AGENT'] ?? '')), 'depuis' => time()];
        $this->pdo->prepare('UPDATE utilisateurs SET verrou_jusqua = NULL, tentatives = 0, derniere_connexion = NOW() WHERE id = ?')
                 ->execute([$u['id']]);
        return Utilisateur::deLigne($u);
    }

    public function deconnecter(): void
    {
        if (ini_get('session.use_cookies')) {
            $p = session_get_cookie_params();
            setcookie(session_name(), '', time() - 42000, $p['path'], $p['domain'], $p['secure'], $p['httponly']);
        }
        $_SESSION = [];
        session_destroy();
    }

    /** Reinitialisation : jeton aleatoire, HACHE en base, 15 min, usage unique, reponse neutre. */
    public function demanderReinitialisation(string $email): void
    {
        $s = $this->pdo->prepare('SELECT id FROM utilisateurs WHERE email = ?');
        $s->execute([mb_strtolower($email)]);
        $id = $s->fetchColumn();
        if ($id === false) { return; }                      // pas de fuite d'information : meme message
        $brut = bin2hex(random_bytes(32));
        $this->pdo->prepare('INSERT INTO reinitialisations (utilisateur_id, jeton_hash, expire_le) VALUES (?, ?, NOW() + INTERVAL 15 MINUTE)')
                  ->execute([$id, hash('sha256', $brut)]);
        $this->mailer->envoyer($email, 'Reinitialisation', "Lien : " . url('/mdp-oublier') . '?j=' . urlencode($brut));
    }

    public function appliquerReinitialisation(string $jetonRecu, string $nouveauMdp): void
    {
        $this->verifierForce($nouveauMdp);
        $s = $this->pdo->prepare(
            'SELECT utilisateur_id FROM reinitialisations
             WHERE jeton_hash = ? AND expire_le > NOW() AND utilise_le IS NULL LIMIT 1'
        );
        $s->execute([hash('sha256', $jetonRecu)]);
        $id = $s->fetchColumn();
        if ($id === false) { throw new ErreurAuthentification('Lien invalide ou expire.'); }
        $this->pdo->prepare('UPDATE utilisateurs SET mot_de_passe = ? WHERE id = ?')
                  ->execute([password_hash($nouveauMdp, PASSWORD_ARGON2ID), $id]);
        $this->pdo->prepare('UPDATE reinitialisations SET utilise_le = NOW() WHERE utilisateur_id = ? AND utilise_le IS NULL')->execute([$id]);
        $this->pdo->prepare('DELETE FROM sessions WHERE utilisateur_id = ?')->execute([$id]);   // deconnecter partout
    }

    private function verifierForce(string $mdp): void
    {
        if (mb_strlen($mdp) < 12) { throw new ErreurValidation('Mot de passe trop court (12 caracteres minimum).'); }
        // (bonus) comparer aux fuites connues via l'API Have I Been Pwned range, sans envoyer le mot de passe
    }
    private function marquerEchec(int $id, ?string $ip): void { /* incrementer tentatives, poser verrou */ }
}
\end{codephp}

\subsection{Le middleware d'authentification, et les autorisations}

\begin{codephp}
<?php
// Un controleur ne verifie pas qui il est : un middleware le fait, une fois, pour tout.
final class ExigerConnexion
{
    public function __construct(private readonly PDO $pdo) {}
    public function __invoke(callable $suite): Response
    {
        $id = $_SESSION['user_id'] ?? null;
        if ($id === null) { return new Response(401, ['www-authenticate' => 'session'], 'Connexion requise'); }
        // Lier la session a l'empreinte du navigateur et a l'IP : casse la session volee depuis un autre poste
        if (!hash_equals($_SESSION['agent'] ?? '', hash('sha256', (string) ($_SERVER['HTTP_USER_AGENT'] ?? '')))) {
            $this->deconnecter();
            return new Response(401, [], 'Session invalide');
        }
        $u = Utilisateur::trouver($this->pdo, (int) $id) ?? (function (): never { $this->deconnecter(); throw new RuntimeException(); })();
        $GLOBALS['utilisateur'] = $u;             // dans un vrai framework : attribut de requete
        return $suite();
    }
    private function deconnecter(): void { (new Authentification($this->pdo))->deconnecter(); }
}

// Autorisation : centralisee, testable, et qui renvoie 404 quand l'existence meme est sensible
final class Politiques
{
    public static function peutVoir(Utilisateur $u, Article $a): bool
    {
        return $a->estPublie() || $u->id() === $a->auteur()->id() || $u->role()->peut(Role::Redacteur);
    }
    public static function peutEditer(Utilisateur $u, Article $a): bool
    {
        return $u->id() === $a->auteur()->id() || $u->role()->peut(Role::Admin);
    }
}
\end{codephp}

\section{Upload de fichiers~: la version non-na\"ive}

\begin{codephp}
<?php
/**
 * @param array{name:string,type:string,tmp_name:string,error:int,size:int} $f  $_FILES['avatar']
 */
function traiterUpload(array $f, string $dossier, int $tailleMax = 5_242_880): string
{
    if ($f['error'] !== UPLOAD_ERR_OK) {
        throw new ErreurUpload(match ($f['error']) {
            UPLOAD_ERR_INI_SIZE, UPLOAD_ERR_FORM_SIZE => 'Fichier trop volumineux.',
            UPLOAD_ERR_PARTIAL => 'Envoi interrompu.',
            UPLOAD_ERR_NO_FILE => 'Aucun fichier recu.',
            default => 'Erreur d\'envoi (' . $f['error'] . ').',
        });
    }
    if ($f['size'] > $tailleMax) { throw new ErreurUpload('Taille maximale depassee.'); }

    // 1) verifier que c'est REELLEMENT une image (le MIME annonce par le client ne vaut rien)
    $info = @getimagesize($f['tmp_name']);
    if ($info === false) { throw new ErreurUpload('Ce n\'est pas une image.'); }
    $types = [IMAGETYPE_PNG => 'png', IMAGETYPE_JPEG => 'jpg', IMAGETYPE_WEBP => 'webp', IMAGETYPE_GIF => 'gif'];
    $ext = $types[$info[2]] ?? throw new ErreurUpload('Format non accepte.');

    // 2) nom genere : jamais le nom du client dans un chemin (traversee de repertoire)
    $nom = sprintf('%s-%s.%s', date('Ymd'), bin2hex(random_bytes(8)), $ext);

    // 3) hors racine web,permissions 0640, proprietaire du service ; on reencode pour neutraliser
    //    toute charge utile cachee (EXIF, script) : passer par GD ou ImageMagick
    $destination = rtrim($dossier, '/') . '/' . $nom;
    if (!is_dir($dossier)) { mkdir($dossier, 0o750, true); }
    if (!move_uploaded_file($f['tmp_name'], $destination)) {      // LA seule fonction qui prouve
        throw new ErreurUpload('Enregistrement impossible.');                   // l'origine PHP de l'upload
    }
    reencoder($destination, $ext);                                  // redimensionne + retire EXIF
    chmod($destination, 0o640);
    return '/uploads/' . $nom;
}
// En nginx/apache : le dossier d'upload ne doit JAMAIS executer de PHP :
//   php_admin_flag engine Off / ou alias en `location ^~ /uploads/ { autoindex off; }`
\end{codephp}

\subsection{En-t\^etes, cookies, emails~: les trois finitions}

\begin{codephp}
<?php
// Cookies (le meme contrat que vu en JS)
setcookie('theme', 'sombre', [
    'expires' => time() + 31536000, 'path' => '/', 'domain' => '',
    'secure' => true, 'httponly' => true, 'samesite' => 'Lax', 'partitioned' => false,
]);
setrawcookie('brut', rawurlencode($valeur), ['expires' => 0, 'path' => '/']);

// En-tetes de securite, a poser globalement (php-fpm : dans le pool, ou dans le premier include)
header_remove('x-powered-by'); header_remove('server');
header('x-content-type-options: nosniff');
header('referrer-policy: strict-origin-when-cross-origin');
header('x-frame-options: DENY');
header('permissions-policy: camera=(), microphone=(), geolocation=()');
header('content-security-policy: default-src \'self\'; img-src \'self\' data:; style-src \'self\' \'unsafe-inline\'');
header('strict-transport-security: max-age=31536000; includeSubDomains');

// Email : la fonction mail() de PHP est une source de problems (pas d'auth, pas de file,
// spam). Envoie via SMTP reel ou une API. Deux choix sains :
//   - Symfony Mailer : composer require symfony/mailer
//   - PHPMailer     : composer require phpmailer/phpmailer
use Symfony\Component\Mailer\Mailer;
use Symfony\Component\Mailer\Transport\Smtp\EsmtpTransport;
use Symfony\Component\Mime\Email;

$transport = new EsmtpTransport((string) getenv('SMTP_HOST'), (int) getenv('SMTP_PORT'), true);
$transport->setUsername((string) getenv('SMTP_USER'))->setPassword((string) getenv('SMTP_PASS'));
(new Mailer($transport))->send(
    (new Email())->from('nour@example.com')->to($email)->subject('Bienvenue')
        ->text("Salut {$pseudo},\n\nTon compte est pret.")
        ->html('<p>Salut <b>' . htmlspecialchars($pseudo) . '</b></p>')
        ->attachFromPath($cheminPdf, 'guide.pdf')
);
\end{codephp}

\begin{exercice}
\begin{enumerate}
  \item \'Ecris un \og ~formulaire g\'en\'erique~ \fg{}~: \verb|champ(type, nom, options)| qui g\'en\`ere le HTML,
        lit \verb|\$\_POST|, valide selon un tableau de r\`egles, et renvoie
        \verb|[valeurs, erreurs]|. Trois formulaires deviennent alors dix lignes chacun.
  \item Ajoute la protection \og ~honeypot~ \fg{} (un champ cach\'e que seul un robot remplit) et un
        d\'elai minimal entre affichage et envoi~; bloque les robots sans CAPTCHA.
  \item \'Ecris le module d'inscription/connexion complet ci-dessus sur ton blog~: avec verrouillage
        apr\`es cinq \'echecs, journalisation des connexions, et \og ~se souvenir de moi~ \fg{} via un jeton
        en base (jamais un cookie \og ~mot de passe hach\'e~ \fg{}).
  \item Ajoute un 2FA TOTP~: g\'en\`ere le secret, le code de v\'erification (\link{otphp}{https://github.com/Spomky-Labs/otphp}),
        et exige le code pour les r\^oles sensibles.
  \item Fais un upload multiple (jusqu'\`a 5 images)~: validit\'e, redimensionnement en trois
        tailles, nommage, suppression \`a la suppression de l'article, et tests sur des fichiers
        pi\`eges (\verb|photo.php|.jpg, image SVG, image de 40~MP, GIF anim\'e).
  \item \'Ecris une page \og ~mes sessions actives~ \fg{}~: liste des sessions de l'utilisateur (agent,
        IP, date), et bouton pour les fermer une par une~; c'est l'exercice qui fait comprendre
        le stockage de session.
\end{enumerate}
\end{exercice}

\begin{ressources}
\begin{itemize}
  \item Manuel --- Gestion des sessions --- \url{https://www.php.net/manual/en/book.session.php}
  \item Manuel --- Fonction \verb|password\_*| --- \url{https://www.php.net/manual/en/ref.password.php}
  \item PHP-FIG PSR-7 (messages HTTP) --- \url{https://www.php-fig.org/psr/psr-7/} --- PSR-15 (middlewares)
  \item \link{Symfony Mailer}{https://symfony.com/doc/current/mailer.html} --- \link{PHPMailer}{https://github.com/PHPMailer/PHPMailer}
  \item \link{OWASP~Cheat~Sheet~Authentication}{https://cheatsheetseries.owasp.org/cheatsheets/Authentication\_Cheat\_Sheet.html}
\end{itemize}
\end{ressources}

\index{PHP!formulaire}\index{PHP!session}\index{PHP!authentification}\index{CSRF}\index{upload}
% ============================================================================
\chapter{PHP et MySQL~: PDO, requ\^etes pr\'epar\'ees, transactions, migrations}
\label{ch:php-db}
% ============================================================================

\begin{objectif}
\begin{itemize}
  \item Concevoir un sch\'ema correct (types, cl\'es, index) et le versionner avec des migrations.
  \item Utiliser PDO parfaitement~: requ\^etes pr\'epar\'ees, transactions, modes de r\'ecup\'eration,
        erreurs en exception.
  \item \'Ecrire les requ\^etes d'une vraie application~: CRUD, pagination, agr\'egats, jointures.
  \item Comprendre pourquoi une requ\^ete est lente et la r\'eparer avec \verb|EXPLAIN|.
  \item \'Ecrire un d\'ep\^ot (repository) simple, testable, sans ORM, puis savoir ce qu'un ORM t'apporte.
\end{itemize}
\end{objectif}

\section{Installer, cr\'eer la base, la sch\'ematiser}

\begin{terminal}
# MariaDB locale (ou le conteneur du chapitre 1)
sudo apt install -y mariadb-server && sudo mariadb-secure-installation
mysql -u root -p
CREATE DATABASE blog CHARACTER SET utf8mb4 COLLATE utf8mb4_unicode_ci;
CREATE USER 'blog'@'localhost' IDENTIFIED BY 'un-mot-de-passe-long';
GRANT SELECT, INSERT, UPDATE, DELETE ON blog.* TO 'blog'@'localhost';
-- Surtout PAS GRANT ALL ni root pour une application.
FLUSH PRIVILEGES;
\end{terminal}

\begin{codesql}
-- schema/001_init.sql --- un schema qui tient la route
CREATE TABLE utilisateurs (
  id            BIGINT UNSIGNED NOT NULL AUTO_INCREMENT,
  email         VARCHAR(255)    NOT NULL,
  mot_de_passe  VARCHAR(255)    NOT NULL,          -- hash argon2 : ~97 caracteres, ne troncne pas a 60
  pseudo        VARCHAR(40)     NOT NULL,
  role          ENUM('lecteur','auteur','redacteur','admin') NOT NULL DEFAULT 'lecteur',
  actif         TINYINT(1)      NOT NULL DEFAULT 1,
  tentatives    TINYINT UNSIGNED NOT NULL DEFAULT 0,
  verrou_jusqua DATETIME        NULL,
  cree_le       DATETIME        NOT NULL DEFAULT CURRENT_TIMESTAMP,
  maj_le        DATETIME        NOT NULL DEFAULT CURRENT_TIMESTAMP ON UPDATE CURRENT_TIMESTAMP,
  PRIMARY KEY (id),
  UNIQUE KEY uq_utilisateur_email (email),
  KEY idx_utilisateur_role (role, actif)
) ENGINE=InnoDB;

CREATE TABLE articles (
  id          BIGINT UNSIGNED NOT NULL AUTO_INCREMENT,
  slug        VARCHAR(191)    NOT NULL,
  titre       VARCHAR(191)    NOT NULL,
  corps       MEDIUMTEXT      NOT NULL,
  statut      ENUM('brouillon','publie','archive') NOT NULL DEFAULT 'brouillon',
  auteur_id   BIGINT UNSIGNED NOT NULL,
  vues        INT UNSIGNED    NOT NULL DEFAULT 0,
  publie_le   DATETIME        NULL,
  cree_le     DATETIME        NOT NULL DEFAULT CURRENT_TIMESTAMP,
  maj_le      DATETIME        NOT NULL DEFAULT CURRENT_TIMESTAMP ON UPDATE CURRENT_TIMESTAMP,
  version     INT UNSIGNED    NOT NULL DEFAULT 1,      -- controle optimiste
  PRIMARY KEY (id),
  UNIQUE KEY uq_article_slug (slug),
  KEY idx_article_liste (statut, publie_le DESC),
  CONSTRAINT fk_article_auteur FOREIGN KEY (auteur_id) REFERENCES utilisateurs(id)
    ON DELETE RESTRICT ON UPDATE CASCADE,
  CONSTRAINT chk_article_titre CHECK (CHAR_LENGTH(titre) BETWEEN 3 AND 191)
) ENGINE=InnoDB;

CREATE TABLE tags (id SMALLINT UNSIGNED NOT NULL AUTO_INCREMENT, nom VARCHAR(40) NOT NULL,
  PRIMARY KEY (id), UNIQUE KEY uq_tag_nom (nom)) ENGINE=InnoDB;

CREATE TABLE article_tag (
  article_id BIGINT UNSIGNED NOT NULL,
  tag_id     SMALLINT UNSIGNED NOT NULL,
  PRIMARY KEY (article_id, tag_id),
  CONSTRAINT fk_at_article FOREIGN KEY (article_id) REFERENCES articles(id) ON DELETE CASCADE,
  CONSTRAINT fk_at_tag     FOREIGN KEY (tag_id)     REFERENCES tags(id)     ON DELETE CASCADE
) ENGINE=InnoDB;

CREATE INDEX idx_article_recherche ON articles ((CAST(corps AS CHAR(255)))) ;  -- astuce ; prefere FULLTEXT :
ALTER TABLE articles ADD FULLTEXT KEY ft_article (titre, corps);
-- SELECT * FROM articles WHERE MATCH(titre, corps) AGAINST ('php pdo' IN NATURAL LANGUAGE MODE);
\end{codesql}

\etape{R\`egles de sch\'ema qui t'\'epargnent des ann\'ees de douleur}
\begin{enumerate}
  \item \verb|utf8mb4| \emph{partout} (les \'emojis et l'arabe passent, contrairement \`a \verb|utf8mb3|).
  \item Une cl\'e primaire \verb|BIGINT UNSIGNED AUTO\_INCREMENT| partout, ou un UUIDv7 stock\'e en
        \verb|BINARY(16)| (meilleure localit\'e). Un id \og ~email~ \fg{} ne se met pas \`a jour~: ne le fais pas.
  \item Dates en \verb|DATETIME| \textbf{UTC}, jamais \verb|TIMESTAMP| (limit\'e \`a 2038 sur vieux syst\`emes,
        et transforme silencieusement selon le fuseau de la session).
  \item \verb|NOT NULL| par d\'efaut~; \verb|NULL| seulement quand \og ~inconnu~ \fg{} est un \'etat m\'etier
        (sinon \verb|COUNT(col)|, \verb|!=|, \verb|NOT IN| te trahissent).
  \item Un index pour chaque colonne de \verb|WHERE|/\verb|JOIN|/\verb|ORDER BY| fr\'equent, et des
        index \emph{compos\'es} dans le bon ordre (\'egalit\'e d'abord, tri ensuite).
  \item Migrations versionn\'ees, jamais d'\verb|ALTER| \`a la main en production.
\end{enumerate}

\section{PDO~: la seule API de base de donn\'ees \`a apprendre en PHP}

\begin{codephp}
<?php
declare(strict_types=1);

function creerPdo(): PDO
{
    $dsn = sprintf('mysql:host=%s;port=%d;dbname=%s;charset=utf8mb4',
        getenv('DB_HOST') ?: '127.0.0.1', (int) (getenv('DB_PORT') ?: 3306), getenv('DB_NAME') ?: 'blog');
    $pdo = new PDO($dsn, (string) getenv('DB_USER'), (string) getenv('DB_PASS'), [
        PDO::ATTR_ERRMODE => PDO::ERRMODE_EXCEPTION,               // exceptions, jamais des silences
        PDO::ATTR_DEFAULT_FETCH_MODE => PDO::FETCH_ASSOC,
        PDO::ATTR_EMULATE_PREPARES => false,                        // vraies requetes preparees (securite + types)
        PDO::ATTR_STRINGIFY_FETCHES => false,                      // int restent int
        PDO::MYSQL_ATTR_MULTI_STATEMENTS => false,                  // refuse "1; DROP TABLE"
        PDO::MYSQL_ATTR_INIT_COMMAND => "SET time_zone = '+00:00', sql_mode = 'STRICT_ALL_TABLES'",
        PDO::ATTR_PERSISTENT => false,                              // true seulement si tu sais pourquoi
    ]);
    return $pdo;
}
$pdo = creerPdo();

// ---- INSERT : prepare + execute. JAMAIS de concatenation.
$sql = 'INSERT INTO articles (slug, titre, corps, statut, auteur_id, publie_le)
        VALUES (:slug, :titre, :corps, :statut, :auteur, :publie_le)';
$stmt = $pdo->prepare($sql);
$stmt->bindValue('slug', $slug, PDO::PARAM_STR);
$stmt->bindValue('titre', $titre, PDO::PARAM_STR);
$stmt->bindValue('corps', $corps, PDO::PARAM_STR);
$stmt->bindValue('statut', 'publie', PDO::PARAM_STR);
$stmt->bindValue('auteur', $auteurId, PDO::PARAM_INT);
$stmt->bindValue('publie_le', $publieLe?->format('Y-m-d H:i:s'), $publieLe === null ? PDO::PARAM_NULL : PDO::PARAM_STR);
$stmt->execute();
$id = (int) $pdo->lastInsertId();
$modifies = $stmt->rowCount();

// Variante courte (valeurs positionnees) :
$pdo->prepare('INSERT INTO tags (nom) VALUES (?)')->execute([mb_strtolower($nom)]);

// Upsert MySQL / PostgreSQL (deux syntaxes differentes !) :
$pdo->prepare('INSERT INTO utilisateurs (email, pseudo) VALUES (?, ?)
               ON DUPLICATE KEY UPDATE pseudo = VALUES(pseudo)')->execute([$email, $pseudo]);
// Postgres : ON CONFLICT (email) DO UPDATE SET pseudo = EXCLUDED.pseudo

// ---- SELECT : les quatre facons de recuperer
$un = $pdo->query('SELECT id, titre FROM articles WHERE id = 42')->fetch();                    // array|false
$un = $pdo->prepare('SELECT id, titre FROM articles WHERE slug = ?'); $un->execute([$slug]); $article = $un->fetch();
$ligne = $un->fetch(PDO::FETCH_NUM);       // indexee par position
$objet = $un->fetch(PDO::FETCH_OBJ);      // stdClass
 typed = $un->setFetchMode(PDO::FETCH_CLASS, Article::class);
$cartographie = function (array $l): Article { return Article::deLigne($l); };                  // LA bonne habitude
foreach ($pdo->query('SELECT * FROM articles WHERE statut = "publie"') as $l) { $liste[] = $cartographie($l); }
$tous = $pdo->query('SELECT id, titre FROM articles')->fetchAll(PDO::FETCH_KEY_PAIR);           // id => titre
$colonne = $pdo->query('SELECT slug FROM articles')->fetchAll(PDO::FETCH_COLUMN);
$compte = (int) $pdo->query('SELECT COUNT(*) FROM articles')->fetchColumn();

// ---- UPDATE / DELETE : le WHERE est obligatoire, et limite
$aff = $pdo->prepare('UPDATE articles SET titre = ?, version = version + 1 WHERE id = ? AND version = ?');
$aff->execute([$titre, $id, $versionAttendue]);
if ($aff->rowCount() === 0) { throw new ConflitModification('L\'article a ete modifie entre-temps.'); }
$pdo->prepare('DELETE FROM articles WHERE id = ? AND auteur_id = ?')->execute([$id, $userId]);  // autorisation dans la requete

// ---- Parametres multiples dynamiques (IN (...)) : generer les marqueurs, pas la valeur
$ids = [1, 5, 9];
$marqueurs = implode(',', array_fill(0, count($ids), '?'));
$rows = $pdo->prepare("SELECT id, titre FROM articles WHERE id IN ($marqueurs)")->execute($ids)->fetchAll();

// ---- LIMIT / OFFSET : pas de parametre possible (ce sont des entiers : cast)
$limite = max(1, min(100, (int) $page->taille));
$decalage = max(0, (int) $page->decalage);
$rows = $pdo->prepare("SELECT * FROM articles ORDER BY publie_le DESC LIMIT {$limite} OFFSET {$decalage}")->fetchAll();
\end{codephp}

\begin{piege}
\textbf{Le pi\`ege PHP le plus co\^uteux.} Une clause \verb|LIKE '\%motif\%'| avec une valeur
contenant \verb|\%| ou \verb|\_| doit \^etre \'echapp\'ee (\verb|addcslashes(\$q, '\%\_\\')|), sinon
\og ~100\,\%~ \fg{} devient un joker. Et surtout~: un \verb|LIMIT|/\verb|ORDER BY| \emph{ne peut pas}
venir d'un param\`etre de requ\^ete pr\'epar\'ee~; ce sont les deux seuls endroits o\`u tu construis
la SQL~: donc liste blanche absolue (cf. \cref{ch:js-secu}).
\end{piege}

\section{Transactions~: la consistance, pas le luxe}

\begin{codephp}
<?php
// Exemple reel : muter un solde + ecrire une ecriture comptable, ou transf\'erer entre comptes.
function transferer(PDO $pdo, int $de, int $vers, int $centimes): void
{
    if ($centimes <= 0) { throw new InvalidArgumentException('Montant doit etre positif.'); }
    $pdo->beginTransaction();
    try {
        // verrouillage preventif : SELECT ... FOR UPDATE (InnoDB seulement)
        $sel = $pdo->prepare('SELECT id, solde FROM comptes WHERE id IN (?, ?) ORDER BY id FOR UPDATE');
        $sel->execute([$de, $vers]);
        $comptes = $sel->fetchAll(PDO::FETCH_KEY_PAIR);
        if (!isset($comptes[$de], $comptes[$vers])) { throw new RuntimeException('compte inexistant'); }
        if ($comptes[$de] < $centimes) { throw new RuntimeException('solde insuffisant'); }

        $pdo->prepare('UPDATE comptes SET solde = solde - ? WHERE id = ?')->execute([$centimes, $de]);
        $pdo->prepare('UPDATE comptes SET solde = solde + ? WHERE id = ?')->execute([$centimes, $vers]);
        $ins = $pdo->prepare('INSERT INTO ecritures (compte_id, montant, contrepartie, cree_le) VALUES (?, ?, ?, NOW())');
        $ins->execute([$de, -$centimes, $vers]);
        $ins->execute([$vers, $centimes, $de]);

        $pdo->commit();
    } catch (Throwable $e) {
        $pdo->rollBack();
        throw $e;
    }
}

// Le patron "transaction imbriquee sure" : savepoints
function transaction(PDO $pdo, callable $travail): mixed
{
    static $profondeur = 0;
    if ($profondeur++ === 0) { $pdo->beginTransaction(); } else { $pdo->exec('SAVEPOINT t' . $profondeur); }
    try {
        $r = $travail($pdo);
        $profondeur === 1 ? $pdo->commit() : $pdo->exec('RELEASE SAVEPOINT t' . $profondeur);
        return $r;
    } catch (Throwable $e) {
        $profondeur === 1 ? $pdo->rollBack() : $pdo->exec('ROLLBACK TO SAVEPOINT t' . $profondeur);
        throw $e;
    } finally {
        $profondeur--;
    }
}

// Insertions en masse : une transaction + reuses la meme requete preparee = 50 a 100x plus rapide
$t0 = microtime(true);
$pdo->beginTransaction();
$ins = $pdo->prepare('INSERT INTO logs (niveau, message, cree_le) VALUES (?, ?, NOW())');
foreach ($lignes as $l) { $ins->execute([$l['niveau'], $l['message']]); }
$pdo->commit();
printf("%d lignes en %.2f s\n", count($lignes), microtime(true) - $t0);
\end{codephp}

\section{Paginer, agr\'eger, joindre~: les requ\^etes du quotidien}

\begin{codesql}
-- 1. Pagination avec total (le "SELECT SQL_CALC_FOUND_ROWS" est mort depuis MySQL 8)
SELECT COUNT(*) AS total FROM articles WHERE statut = 'publie';
SELECT id, slug, titre, publie_le FROM articles
WHERE statut = 'publie'
ORDER BY publie_le DESC, id DESC
LIMIT 20 OFFSET 40;

-- 2. Pagination par curseur (la seule qui reste rapide a 100 000 pages)
SELECT id, slug, titre FROM articles
WHERE (publie_le, id) < ('2026-01-31 10:00:00', 4241)
ORDER BY publie_le DESC, id DESC LIMIT 20;

-- 3. Joindre les tags sans dupliquer les lignes (GROUP_CONCAT)
SELECT a.id, a.titre, GROUP_CONCAT(t.nom ORDER BY t.nom SEPARATOR ',') AS tags
FROM articles a
LEFT JOIN article_tag at ON at.article_id = a.id
LEFT JOIN tags t ON t.id = at.tag_id
WHERE a.statut = 'publie'
GROUP BY a.id, a.titre
ORDER BY a.publie_le DESC
LIMIT 20;

-- 4. Agr\'egats par auteur avec et sans articles (LEFT JOIN + COALESCE)
SELECT u.id, u.pseudo,
       COUNT(a.id) AS nb,
       COALESCE(SUM(a.vues), 0) AS vues,
       MAX(a.publie_le) AS dernier
FROM utilisateurs u
LEFT JOIN articles a ON a.auteur_id = u.id AND a.statut = 'publie'
WHERE u.actif = 1
GROUP BY u.id, u.pseudo
HAVING nb > 3
ORDER BY vues DESC
LIMIT 10;

-- 5. Top N par groupe (fenetres : MySQL 8+, MariaDB 10.2+, Postgres partout)
SELECT * FROM (
  SELECT a.*, ROW_NUMBER() OVER (PARTITION BY auteur_id ORDER BY vues DESC) AS rn
  FROM articles a WHERE statut = 'publie'
) x WHERE x.rn <= 3;

-- 6. Recherche full text + pertinence
SELECT id, titre, MATCH(titre, corps) AGAINST (:q IN BOOLEAN MODE) AS score
FROM articles WHERE MATCH(titre, corps) AGAINST (:q IN BOOLEAN MODE) > 0.1
ORDER BY score DESC LIMIT 20;

-- 7. EXPLAIN : l'outil qui repond a "pourquoi c'est lent"
EXPLAIN ANALYZE SELECT * FROM articles WHERE statut = 'publie' ORDER BY publie_le DESC LIMIT 20;
-- regarde : type (const/range/ref/ALL), key, rows, "Using filesort", "Using temporary"
SHOW INDEX FROM articles; SHOW FULL COLUMNS FROM articles;
\end{codesql}

\begin{note}
Deux index suffisent pour la page liste d'un blog~: un index compos\'e
\verb|(statut, publie\_le DESC)| pour le filtre+tri, et \verb|UNIQUE(slug)| pour la page d\'etail.
Si \verb|EXPLAIN| affiche \verb|type: ALL| (scan complet) sur une table de plus de 10 000 lignes,
ton probl\`eme n'est pas le code PHP~: il est dans l'index.
\end{note}

\section{Le d\'ep\^ot~: la couche qui prot\`ege le reste du code}

\begin{codephp}
<?php
declare(strict_types=1);

final class DepotArticles
{
    public function __construct(private readonly PDO $pdo) {}

    public function parId(int $id): ?Article
    {
        $stmt = $this->pdo->prepare(
            'SELECT a.*, u.pseudo AS auteur_pseudo, u.email AS auteur_email
             FROM articles a JOIN utilisateurs u ON u.id = a.auteur_id
             WHERE a.id = ? LIMIT 1'
        );
        $stmt->execute([$id]);
        $l = $stmt->fetch();
        return $l === false ? null : $this->hydrater($l);
    }

    /** @return array{articles: Article[], total: int, page: int, pages: int} */
    public function pageListe(?Statut $statut, string $recherche = '', int $page = 1, int $taille = 20): array
    {
        [$where, $args] = $this->contruireFiltres($statut, $recherche);
        $taille = max(1, min(100, $taille));
        $page = max(1, $page);

        $total = (int) $this->scalar("SELECT COUNT(*) FROM articles a {$where}", $args);
        $sql = "SELECT a.*, u.pseudo AS auteur_pseudo, u.email AS auteur_email
                 FROM articles a JOIN utilisateurs u ON u.id = a.auteur_id
                 {$where}
                 ORDER BY COALESCE(a.publie_le, a.cree_le) DESC, a.id DESC
                 LIMIT {$taille} OFFSET " . (($page - 1) * $taille);
        $lignes = $this->all($sql, $args);

        return [
            'articles' => array_map($this->hydrater(...), $lignes),
            'total' => $total,
            'page' => $page,
            'pages' => max(1, intdiv($total + $taille - 1, $taille)),
        ];
    }

    public function creer(Article $a): Article
    {
        $this->pdo->prepare('INSERT INTO articles (slug, titre, corps, statut, auteur_id, publie_le) VALUES (?,?,?,?,?,?)')
            ->execute([$a->slug(), $a->titre(), $a->corps(), $a->statut()->value, $a->auteur()->id(), $a->publieLe()?->format('Y-m-d H:i:s')]);
        return $a->avecId((int) $this->pdo->lastInsertId());
    }

    public function ajouterTag(int $articleId, string $tag): void
    {
        $this->pdo->prepare('INSERT INTO tags (nom) VALUES (?) ON DUPLICATE KEY UPDATE id = LAST_INSERT_ID(id)')
                  ->execute([mb_strtolower($tag)]);
        $tagId = (int) $this->pdo->lastInsertId();
        $this->pdo->prepare('INSERT IGNORE INTO article_tag (article_id, tag_id) VALUES (?, ?)')
                  ->execute([$articleId, $tagId]);
    }

    /** @return array{0:string,1:array} */
    private function construireFiltres(?Statut $statut, string $recherche): array
    {
        $conditions = []; $args = [];
        if ($statut !== null) { $conditions[] = 'a.statut = ?'; $args[] = $statut->value; }
        if ($recherche !== '') {
            $conditions[] = '(a.titre LIKE ? OR a.corps LIKE ?)';
            $q = '%' . addcslashes(mb_strtolower($recherche), '%_\\') . '%';
            array_push($args, $q, $q);
        }
        return [$conditions === [] ? '' : 'WHERE ' . implode(' AND ', $conditions), $args];
    }

    /** Raccourcis de lecture (le trait PdoPratique, plus bas). */
    private function all(string $sql, array $args = []): array
    { $s = $this->pdo->prepare($sql); $s->execute($args); return $s->fetchAll(); }
    private function scalar(string $sql, array $args = []): mixed
    { $s = $this->pdo->prepare($sql); $s->execute($args); return $s->fetchColumn(); }

    private function hydrater(array $l): Article { /* ... cf. chapitre precedent */ }
}
\end{codephp}

\begin{note}
Ce que j'ai \'ecrit ci-dessus \og ~\verb|->tap(...)|~ \fg{} et \og ~\verb|executeRet(...)|~ \fg{} n'existe pas en PDO~:
c'est une micro-extension que j'aime ajouter~:

\begin{codephp}
trait PdoPratique
{
    public function run(PDOStatement $s, array $args = []): PDOStatement { $s->execute($args); return $s; }
    public function all(PDOStatement $s, array $args = []): array { $s->execute($args); return $s->fetchAll(); }
    public function one(PDOStatement $s, array $args = []): array|false { $s->execute($args); return $s->fetch(); }
    public function scalar(PDOStatement $s, array $args = []): mixed { $s->execute($args); return $s->fetchColumn(); }
}
\end{codephp}
Ces quatre m\'ethodes suppriment la moiti\'e du bruit de PDO et rendent les d\'ep\^ots lisibles.
\end{note}

\section{Migrations~: le fichier qui fait \og ~marche~ \fg{} en production}

\begin{terminal}
composer require --dev robmorgan/phinx        # ou doctrine/migrations, ou phinx-like maison
vendor/bin/phinx init
vendor/bin/phinx create CreateArticles
vendor/bin/phinx migrate -e development
vendor/bin/phinx rollback -e development
vendor/bin/phinx status
\end{terminal}

\begin{codephp}
<?php
use Phinx\Migration\AbstractMigration;

final class CreateArticles extends AbstractMigration
{
    public function up(): void
    {
        $t = $this->table('articles', ['collation' => 'utf8mb4_unicode_ci']);
        $t->addColumn('slug', 'string', ['limit' => 191])
          ->addColumn('titre', 'string', ['limit' => 191])
          ->addColumn('corps', 'text')
          ->addColumn('statut', 'enum', ['values' => ['brouillon', 'publie', 'archive'], 'default' => 'brouillon'])
          ->addColumn('vues', 'integer', ['signed' => false, 'default' => 0])
          ->addColumn('auteur_id', 'integer', ['signed' => false])
          ->addColumn('publie_le', 'datetime', ['null' => true])
          ->addColumn('cree_le', 'datetime', ['default' => 'CURRENT_TIMESTAMP'])
          ->addIndex(['slug'], ['unique' => true])
          ->addIndex(['statut', 'publie_le'])
          ->addForeignKey('auteur_id', 'utilisateurs', 'id', ['delete' => 'RESTRICT', 'update' => 'CASCADE'])
          ->create();
    }
    public function down(): void { $this->table('articles')->drop()->save(); }
}
\end{codephp}

\begin{enumerate}
  \item Une migration = un changement d'\'etat, et elle doit \^etre \textbf{compatible avec le code
        encore en ligne} pendant le d\'eploiement~: ajoute une colonne avant de l'utiliser,
        supprime-la seulement \`a la release suivante.
  \item Ne mets jamais de donn\'ees m\'etier dans une migration de sch\'ema~: s\'epare
        \verb|schema/*| et \verb|data/*.sql| (sinon rollback impossible).
  \item Teste la \verb|down()|. Une migration sans retour est un ticket d'incident garanti.
\end{enumerate}

\section{Et les ORM~?}

\begin{tabz}{lXl}
  \headrow \thd{Outil} \& \thd{Position} \& \thd{Lien} \\
  \midrule
  PDO nu \& le contr\^ole total, le moins de magie~: mon choix pour apprendre et pour 80 \% des projets. \& --- \\
  \link{Doctrine ORM}{https://www.doctrine-project.org/projects/orm.html} \& le standard entreprise, UnitOfWork, hydratation diff\'er\'ee, migrations. Courbe d'apprentissage r\'eelle. \& \url{https://www.doctrine-project.org/} \\
  \link{Eloquent (Laravel)}{https://laravel.com/docs/eloquent} \& le plus productif, actif record, tr\`es \og ~convention~ \fg{}~; magique quand on ne le comprend pas. \& \url{https://laravel.com/docs} \\
  \link{Propel}{https://propelorm.org/} \& g\'en\`ere des classes depuis le sch\'ema, tr\`es typ\'e. \& --- \\
  \link{Doctrine DBAL}{https://www.doctrine-project.org/projects/dbal.html} \& le \og ~PDO am\'elior\'e~ \fg{} multi-plateformes, sans ORM. Bon compromis. \& --- \\
\end{tabz}

\begin{astuce}
Apprends PDO \emph{d'abord}. Un ORM est alors facile~: tu reconnais chaque \og ~lazy loading~ \fg{},
chaque \og ~N+1 requ\^etes~ \fg{}, chaque requ\^ete implicite. Sans cette base, un ORM est une source
d'erreurs silencieuses (et le \og ~N+1~ \fg{} est LA faute de performance la plus fr\'equente en
application web~: 1 requ\^ete pour la liste, puis 1 par \'el\'ement~; la correction tient en un
\verb|JOIN| ou un \verb|IN|).
\end{astuce}

\begin{exercice}
\begin{enumerate}
  \item \'Ecris le sch\'ema complet d'un gestionnaire de t\^aches~: utilisateurs, projets, t\^aches,
        \'etiquettes, affectations, commentaires~; avec les index justifi\'es par trois requ\^etes
        r\'eelles que tu \'ecris en premier.
  \item \'Ecris le d\'ep\^ot correspondant~: \verb|pageListe|, \verb|creer|, \verb|modifier| avec
        verrouillage optimiste (colonne \verb|version|), \verb|supprimerEnCascade| dans une
        transaction~; un script de \emph{seed} de 10 000 t\^aches.
  \item Rends ta pagination rapide \`a la page 5 000~: passe du \verb|OFFSET| au curseur et mesure.
  \item Trouve un \og ~N+1~ \fg{}~: affiche le nombre de requ\^etes par page (un listener SQL, ou le
        d\'ebogueur SQL de Adminer), puis \'elimine-le avec une jointure.
  \item Ajoute la recherche full-text avec pertinence, un surlignage des termes, et un temps
        de requ\^ete sous 50 ms sur 100 000 lignes.
  \item Fais une migration de donn\'ees~: ajouter une colonne \verb|slug|, la remplir, la passer
        \verb|NOT NULL| + \verb|UNIQUE|, en quatre migrations s\'epar\'ees, chacune r\'eversible.
\end{enumerate}
\end{exercice}

\begin{ressources}
\begin{itemize}
  \item Manuel PDO --- \url{https://www.php.net/manual/fr/book.pdo.php}
  \item \emph{PHP PDO: The Definitive Guide} / articles dereference --- \url{https://phpdelusions.net/pdo}
  \item MySQL 8 reference --- \url{https://dev.mysql.com/doc/refman/8.4/en/} --- MariaDB --- \url{https://mariadb.com/kb/}
  \item PostgreSQL (\`a apprendre apr\`es MySQL~: le meilleur moteur relationnel) --- \url{https://www.postgresql.org/docs/}
  \item Phinx --- \url{https://book.cakephp.org/phinx/0/en.html} --- Doctrine Migrations --- \url{https://www.doctrine-project.org/projects/migrations.html}
  \item \link{use-the-index-luke}{https://use-the-index-luke.com/} (index : le meilleur cours gratuit qui existe)
\end{itemize}
\end{ressources}

\index{PHP!PDO}\index{MySQL}\index{index}\index{migration}\index{transaction}
% ============================================================================
\chapter{S\'ecurit\'e PHP~: les failles r\'eelles et leurs rem\`edes exacts}
\label{ch:php-secu}
% ============================================================================

\begin{objectif}
\begin{itemize}
  \item Reconna\^itre les huit failles qui font l'essentiel des incidents PHP~: injection SQL,
        XSS, upload, LFI/RFI, d\'es\'erialisation, IDOR, SSRF, d\'ependance.
  \item \'Ecrire les r\'eflexes qui les bouchent toutes, ligne par ligne.
  \item Configurer PHP lui-m\^eme pour qu'il soit s\^ur par d\'efaut.
  \item Journaliser assez pour comprendre un incident, pas assez pour en cr\'eer un second.
\end{itemize}
\end{objectif}

\section{Configuration dure~: le \fichier{php.ini} de production}

\begin{codeconf}
; /etc/php/8.3/fpm/php.ini  -- les lignes qui comptent vraiment
engine = On
short_open_tag = Off                 ; sinon "<?" devient du PHP dans un template XML
expose_php = Off                     ; retire "X-Powered-By: PHP/8.3"
precision = 14

; --- Erreurs : rien au visiteur, tout au journal
display_errors = Off
display_startup_errors = Off
log_errors = On
error_log = /var/log/php/error.log
error_reporting = E_ALL
ignore_repeated_errors = On

; --- Limites : protegent contre l'epuisement
memory_limit = 256M
max_execution_time = 30
max_input_time = 60
post_max_size = 12M
upload_max_filesize = 10M
max_file_uploads = 20
max_input_vars = 1000                ; un <= tableau en entree ne doit pas pouvoir tout ecraser

; --- Securite du systeme de fichiers
open_basedir = /var/www/nour:/tmp    ; ou mieux : un pool php-fpm par site avec chroot
disable_functions = exec,passthru,shell_exec,system,proc_open,popen,curl_multi_exec,show_source,pcntl_alarm
allow_url_fopen = Off                ; sinon file_get_contents("http://...") devient un vecteur SSRF
allow_url_include = Off              ; RFI : doit etre une heresie, verifie-le

; --- Sessions durcies (cf. chapitre precedent)
session.use_strict_mode = On
session.cookie_httponly = On
session.cookie_secure = On
session.cookie_samesite = Lax
session.use_only_cookies = On
session.sid_length = 48

; --- Date
date.timezone = UTC
\end{codeconf}

\begin{note}
\verb|disable\_functions| n'est pas une s\'ecurit\'e, c'est une ceinture~: si un attaquant ex\'ecute
du code chez toi, il a d\'ej\`a gagn\'e. La vraie s\'ecurit\'e, c'est de ne jamais \'evaluer de donn\'ee
comme du code. Le point important de la liste~: \verb|allow\_url\_include=Off| et
\verb|display\_errors=Off| en production (une stack trace publi\'ee = plan du ch\^ateau).
\end{note}

\section{Injection SQL~: un seul rem\`ede}

\begin{codephp}
<?php
// NON : concatenation ou sprintf dans la SQL
$sql = "SELECT * FROM articles WHERE id = " . $_GET['id'];              // 1; DROP TABLE articles; --
$sql = sprintf('SELECT * FROM users WHERE email = "%s"', $email);       // ' OR 1=1 -- '
// OUI : marqueur nomme ou positionne, toujours
$stmt = $pdo->prepare('SELECT * FROM articles WHERE id = :id');
$stmt->execute(['id' => (int) $_GET['id']]);

// Le seul endroit ou l'on construit du texte SQL : les identifiants (table, colonne, ordre).
// Remede : liste blanche, jamais echappement.
$colonnesAutorisees = ['titre', 'vues', 'publie_le'];
$ordre = in_array($_GET['tri'] ?? '', $colonnesAutorisees, true) ? $_GET['tri'] : 'publie_le';
$sens = strtolower((string) ($_GET['ordre'] ?? 'desc')) === 'asc' ? 'ASC' : 'DESC';
$sql = "SELECT * FROM articles ORDER BY {$ordre} {$sens} LIMIT 20";     // sur, car valeurs issues d'une enum

// LIKE : l'echappement des jokers est a faire par toi
$q = '%' . addcslashes(mb_strtolower($recherche), '\\%_') . '%';
$pdo->prepare("SELECT * FROM articles WHERE LOWER(titre) LIKE ? ESCAPE '\\\\'")->execute([$q]);
\end{codephp}

\begin{piege}
Un \verb|real\_escape\_string()| + des guillemets bien plac\'es \og ~marche~ \fg{} tant que le charset de la
connexion est coh\'erent~: en \verb|latin1| ou avec \verb|emulate prepares| activ\'e, certaines
suites d'octets cassent l'\'echappement (les attaques GBK/UTF-8 historiques). Une requ\^ete
pr\'epar\'ee \textbf{r\'eelle} (\verb|PDO::ATTR\_EMULATE\_PREPARES => false|) supprime la classe enti\`ere
de probl\`emes. Il n'y a donc aucune raison de faire autrement.
\end{piege}

\section{XSS et \'echappement de sortie~: le contexte d\'ecide}

\begin{codephp}
<?php
function e(?string $v): string            // dans le HTML
{ return htmlspecialchars((string) $v, ENT_QUOTES | ENT_SUBSTITUTE | ENT_HTML5, 'UTF-8'); }

function eAttr(?string $v): string        // dans un attribut non entre guillemets : ne fait pas \c{c}a
{ return e($v); }                          // mets toujours des guillemets autour des attributs

function eUrl(?string $v): string
{
    $u = filter_var((string) $v, FILTER_VALIDATE_URL);
    if ($u === false) { return '#'; }
    $h = parse_url($u, PHP_URL_SCHEME);
    return in_array(strtolower((string) $h), ['http', 'https', 'mailto'], true) ? $u : '#';
}

function eJs(?string $v): string           // dans un <script> : JSON, pas htmlspecialchars
{ return json_encode($v, JSON_HEX_TAG | JSON_HEX_AMP | JSON_HEX_APOS | JSON_HEX_QUOT | JSON_UNESCAPED_UNICODE) ?: '""'; }

function eCss(?string $v): string { return preg_replace('/[^a-zA-Z0-9 _#.-]/', '', (string) $v) ?? ''; }
function eEmailHtml(?string $v): string    // anti-obfuscation dans les mails texte
{ return str_replace('@', '[at]', (string) $v); }
?>
<!-- dans la vue : le bon escaper pour le bon trou -->
<h1><?= e($article->titre()) ?></h1>
<a href="<?= eUrl($auteur->site()) ?>" rel="noopener nofollow ugc"><?= e($auteur->pseudo()) ?></a>
<div data-id="<?= e((string) $id) ?>" style="color: <?= eCss($couleur) ?>"></div>
<script>window.__DATA__ = <?= eJs(json_encode($donnees, JSON_THROW_ON_ERROR)) ?>;</script>
\end{codephp}

\begin{itemize}
  \item Une vue Twig/Blade \'echappe par d\'efaut~: les XSS survivent presque toujours \`a un
        \verb|{{ x|raw }}| ajout\'e \og ~pour que \c{c}a s'affiche~ \fg{}. Interdis-le en revue de code.
  \item Pour le HTML riche~: \link{HTML Purifier}{http://htmlpurifier.org/},
        \link{symfony/html-sanitizer}{https://symfony.com/doc/current/html\_sanitizer.html},
        ou \verb|strip\_tags| + une liste blanche d'\'etiquettes. Ne code jamais ton propre filtre.
  \item \verb|header('content-type: text/html; charset=utf-8')| avant toute sortie~: sans
        charset d\'eclar\'e, un navigateur peut \og ~deviner~ \fg{} un encodage qui recr\'ee du HTML depuis
        des octets valides (le vieux \verb|UTF-7|-style XSS).
\end{itemize}

\section{Inclusion de fichiers, d\'es\'erialisation, \'evaluations}

\begin{codephp}
<?php
// 1) Inclusion pilotee par l'entree utilisateur = LFI / RFI. Le pire pattern de l'histoire de PHP.
require $_GET['page'] . '.php';                      // ../../../../etc/passwd, ou une URL distante
require 'vues/' . $nom . '.php';                     // $nom vient d'un champ de formulaire ? idem

// Le remede : une table de correspondance, ou un filtre alphanumerique strict.
$vues = ['accueil' => 'vues/accueil.php', 'contact' => 'vues/contact.php'];
require $vues[(string) ($_GET['page'] ?? 'accueil')] ?? $vues['accueil'];
// ou :
$nom = preg_replace('/[^a-z0-9_-]/', '', mb_strtolower((string) $_GET['page']));
$chemin = realpath(__DIR__ . '/vues/' . $nom . '.php');
if ($chemin === false || !str_starts_with($chemin, realpath(__DIR__ . '/vues') . DIRECTORY_SEPARATOR)) {
    throw new RuntimeException('vue introuvable');       // verif de traversee de repertoire
}
require $chemin;

// 2) unserialize sur une donnee utilisateur : RCE dans de nombreux cas ("POP chain").
$etat = unserialize($_COOKIE['panier']);                 // NON
$etat = json_decode((string) $_COOKIE['panier'], true, 8, JSON_THROW_ON_ERROR);   // OUI
// Si tu dois vraiment s\'erialiser : signature + allowed_classes:false
$signe = base64_encode(serialize($objet)) . '.' . hash_hmac('sha256', $objet, $secret);
$objet = unserialize(base64_decode($partie), ['allowed_classes' => [MaClasse::class]]);

// 3) eval(), preg_replace /e (mort), create_function (mort), extract($donnees), compact()
//    dans une vue, ${$cle} dynamiques : tout ce qui transforme des donnees en code.
extract($_GET);                                          // NON : l'attaquant definit tes variables
parse_str(parse_url($url, PHP_URL_QUERY) ?: '', $out);  // $out['id'] = ... -> OK si tu lis le tableau
\end{codephp}

\section{IDOR, CSRF, SSRF, commandes~: les quatre qui restent}

\begin{codephp}
<?php
// IDOR (l'identifiant est choisi par l'attaquant) : toujours filtrer par proprietaire
$id = (int) $_GET['id'];
$facture = $pdo->prepare('SELECT * FROM factures WHERE id = ?');                  // VULNERABLE
$facture = $pdo->prepare('SELECT * FROM factures WHERE id = ? AND utilisateur_id = ?');
$facture->execute([$id, $_SESSION['user_id']]);                                     // sur
// + ne renvoie pas 403 qui confirme l'existence : 404.

// CSRF : cf. chapitre formulaires. Jeton + Origin + SameSite=Lax.
$enteteOrigine = $_SERVER['HTTP_ORIGIN'] ?? null;
if ($_SERVER['REQUEST_METHOD'] !== 'GET' && $enteteOrigine !== null
    && parse_url($enteteOrigine, PHP_URL_HOST) !== $_SERVER['HTTP_HOST']) {
    http_response_code(403); exit('origine non autorisee');
}

// SSRF : un webhook qui demande une URL au visiteur permet de scanner ton reseau interne.
$hote = parse_url($url, PHP_URL_HOST);
$ip = gethostbyname((string) $hote);
if (filter_var($ip, FILTER_VALIDATE_IP, FILTER_FLAG_NO_PRIV_RANGE | FILTER_FLAG_NO_RES_RANGE) === false) {
    throw new InvalidArgumentException('Adresse non autorisee');       // refuse 127.*, 10.*, 169.254.169.254 (metadata cloud)
}
// + desactive les redirections, limite la taille de la reponse, timeout court, https seulement.

// Commandes shell : tableaux + echappement, jamais de concatenation.
exec(sprintf('convert %s %s', escapeshellarg($entree), escapeshellarg($sortie)), $out, $code);
$proc = proc_open(['ffmpeg', '-i', $entree, $sortie], [1 => ['pipe', 'w'], 2 => ['pipe', 'w']], $pipes); // PHP 7.4+ : pas de shell
if ($code !== 0) { throw new RuntimeException('ffmpeg a echoue : ' . implode("\n", $out)); }

// Redirections ouvertes : "?next="
$next = (string) ($_GET['next'] ?? '/');
$safe = (str_starts_with($next, '/') && !str_starts_with($next, '//') && parse_url($next, PHP_URL_HOST) === null) ? $next : '/';
header('location: ' . $safe, true, 303);
\end{codephp}

\section{Journaux, gestion des erreurs et des exceptions}

\begin{codephp}
<?php
declare(strict_types=1);

// Bootstrap : rien au visiteur, tout au journal, y compris les fatal errors.
error_reporting(E_ALL);
ini_set('display_errors', getenv('APP_ENV') === 'dev' ? '1' : '0');
ini_set('log_errors', '1');
ini_set('error_log', '/var/log/php/app.log');

set_error_handler(static function (int $n, string $msg, string $f = '', int $l = 0): bool {
    if ((error_reporting() & $n) === 0) { return false; }
    throw new ErrorException($msg, 0, $n, $f, $l);            // les warnings deviennent des exceptions : + de rigueur
});
set_exception_handler(static function (Throwable $e): void {
    $id = bin2hex(random_bytes(6));                            // identifiant affiche a l'utilisateur
    error_log(sprintf('[%s] %s %s:%d\n%s', $id, $e->getMessage(), $e->getFile(), $e->getLine(), $e->getTraceAsString()));
    http_response_code(500);
    if (($_SERVER['HTTP_ACCEPT'] ?? '') . ($_SERVER['HTTP_X_REQUESTED_WITH'] ?? '') !== '' && stripos((string) ($_SERVER['HTTP_ACCEPT'] ?? ''), 'json') !== false) {
        header('content-type: application/json');
        echo json_encode(['erreur' => 'Incident ' . $id]);
    } else {
        require __DIR__ . '/vues/500.php';                      // page humaine, sans detail technique
    }
});
register_shutdown_function(static function (): void {
    $e = error_get_last();
    if ($e !== null && in_array($e['type'], [E_ERROR, E_PARSE, E_CORE_ERROR, E_COMPILE_ERROR], true)) {
        error_log('FATAL ' . json_encode($e));
        http_response_code(500);
        echo file_get_contents(__DIR__ . '/vues/500.html');
    }
});

// Logger structure maison (ou Monolog, cf. chapitre Composer)
function journal(string $niveau, string $message, array $contexte = []): void
{
    static $requestId;
    $requestId ??= substr(bin2hex(random_bytes(5)), 0, 10);
    $ligne = [
        'ts' => (new DateTimeImmutable('now'))->format(DATE_ATOM),
        'niveau' => $niveau, 'message' => $message, 'request_id' => $requestId,
        'url' => $_SERVER['REQUEST_URI'] ?? null, 'methode' => $_SERVER['REQUEST_METHOD'] ?? null,
        'user_id' => $_SESSION['user_id'] ?? null, 'ip' => $_SERVER['REMOTE_ADDR'] ?? null,
    ] + $contexte;
    error_log(json_encode($ligne, JSON_UNESCAPED_SLASHES | JSON_UNESCAPED_UNICODE), 4);   // 4 = sans horodatage doublon
}
journal('warn', 'mail lente', ['ms' => $duree, 'destinataire' => masquerEmail($email)]);

function masquer(?string $v): string
{
    if ($v === null || $v === '') { return ''; }
    return mb_substr($v, 0, 2) . str_repeat('*', max(0, mb_strlen($v) - 4)) . mb_substr($v, -2);
}
\end{codephp}

\begin{tabz}{lX}
  \headrow \thd{R\`egle de journal} \& \thd{Pourquoi} \\
  \midrule
  Aucun mot de passe, aucun jeton, aucun cookie \& le fichier de log est souvent le fichier le moins prot\'eg\'e du serveur. \\
  Email/IP tronqu\'es, pas supprim\'es \& tu en as besoin pour diagnostiquer, pas pour profiler les gens. \\
  Un identifiant de requ\^ete partout \& corr\`ele un incident avec un seul num\'ero donn\'e \`a l'utilisateur. \\
  Niveaux utiles~: \verb|error|, \verb|warn|, \verb|info| \& le \verb|debug| permanent en production est du bruit co\^uteux. \\
  Rotation et dur\'ee finie \& \verb|logrotate|, 30 \`a 90 jours, ou un agr\'egateur. \\
\end{tabz}

\section{D\'ependances et fichiers~: les deux derni\`eres portes}

\begin{terminal}
# Audit de securite des paquets PHP
composer audit                                  # natif depuis Composer 2.4
composer require --dev roave/security-advisories  # refuse d'installer une version connue vulnerable
# Outils complementaires
composer require --dev sensiolabs/security-checker   # (deprecie) -> utilise `composer audit`
# Packagist signale les versions vulnerables : https://packagist.org/packages/roave/security-advisories
# Analyses statiques orientees securite
composer require --dev vimeo/psalm --dev ; vendor/bin/psalm --taint-analysis
composer require --dev phpcs ; vendor/bin/phpcs --standard=PHPCompatibility
composer require --dev phpstan/phpstan --level=max
# Fichiers sensibles : ne JAMAIS exposer
#   .env, storage/, *.sql, .git/, backup.zip, wp-config.php~
\end{terminal}

\begin{codeconf}
# nginx : la liste noire minimale au-dessus d'une app PHP
location ~ /\..*                       { deny all; }              # .git, .env, .htaccess
location ~* \.(sql|log|bak|old|zip|tar|gz)$ { deny all; }
location ~ ^/(storage|vendor|config|database)/ { deny all; }
# La racine web doit pointer sur le dossier public/ uniquement. Jamais sur la racine du projet.
\end{codeconf}

\begin{exercice}
\begin{enumerate}
  \item \'Ecris un mini-IDS de revue~: une liste de motifs (concat\'enation dans \verb|prepare|,
        \verb|eval|, \verb|unserialize|, \verb|?>| final, \verb|\$\_GET| non typ\'e) que tu passes
        en regex sur \verb|src/|~; applique-la \`a un ancien projet.
  \item Sur ton blog PHP~: injecte-toi toi-m\^eme un XSS stock\'e dans un commentaire, puis
        corrige (\'echappement de sortie + \verb|strip\_tags| pour le gras) et rev\'erifie.
  \item Force une URL \verb|file:///etc/passwd| dans ton champ \og ~aper\c{c}u distant~ \fg{}~: \'echoue,
        puis impl\'emente la validation SSRF ci-dessus~: v\'erifie qu'elle \'echoue.
  \item \'Ecris la page \og ~500~ \fg{} de ton application avec un identifiant d'incident, et un test
        automatis\'e qui v\'erifie qu'aucune stack trace n'est visible en production.
  \item Passe ton \verb|php.ini| en revue contre la liste, puis utilise
        \link{PHP~Hardening}{https://github.com/lockfale/PHP-Security}~: liste les \'ecarts restants
        dans un fichier \verb{SECURITE.md}.
  \item Lance \verb|composer audit| et \verb|psalm --taint-analysis| sur un paquet du commerce~:
        lis un r\'esultat, explique-le, et \'ecris la correction (pas seulement le bump de version).
\end{enumerate}
\end{exercice}

\begin{ressources}
\begin{itemize}
  \item OWASP Cheat Sheet Series (la s\'erie compl\`ete, 15 min par fiche) --- \url{https://cheatsheetseries.owasp.org/}
  \item PHP Security Advisories DB --- \url{https://github.com/Roave/SecurityAdvisories}
  \item \link{PHP~:~Security~Book}{https://phptherightway.com/pages/The-Fundamentals-Of-Programming.html} (chapitre s\'ecurit\'e)
  \item \url{https://www.php.net/manual/fr/security.overview.php} --- la doc officielle, courte et juste
  \item PortSwigger Academy (pour manipuler les attaques web dans un bac \`a sable l\'egal) --- \url{https://portswigger.net/web-security}
\end{itemize}
\end{ressources}

\index{PHP!securite}\index{injection SQL}\index{LFI}\index{SSRF}\index{php.ini}
% ============================================================================
\chapter{Composer, standards PSR, qualit\'e et outils de development}
\label{ch:php-composer}
% ============================================================================

\begin{objectif}
\begin{itemize}
  \item Utiliser Composer comme un adulte~: manifeste, lock, scripts, paquets, publication.
  \item Conna\^itre les PSR qui changent r\'eellement le code (1, 3, 4, 6, 7, 11, 12, 15, 17).
  \item Automatiser le style et la correction~: PHP-CS-Fixer/PHP~Codesniffer, PHPStan, Rector.
  \item Installer l'environnement serveur~: nginx + php-fpm, OPcache, Xdebug, journaux.
  \item Tester~: PHPUnit, Mockery, DBunit-like, et les patterns pour rendre le code testable.
\end{itemize}
\end{objectif}

\section{Composer}

\begin{terminal}
# Installation officielle (Linux/macOS)
curl -sS https://getcomposer.org/installer | php -- --install-dir=/usr/local/bin --filename=composer
composer --version
# Windows : https://getcomposer.org/Composer-Setup.exe
# Alt : php -r "copy('https://getcomposer.org/installer', 'composer-setup.php');" ... (verifie la signature)

cd mon-projet
composer init
composer require guzzlehttp/guzzle monolog/monolog          # prod
composer require --dev phpunit/phpunit phpstan/phpstan rector/rector friendsofphp/php-cs-fixer
composer update                              # respecte les contraintes, puis reecrit composer.lock
composer update vendor/paquet --with-all-dependencies
composer outdated --direct                   # uniquement les dependances directes
composer show --tree
composer why symfony/polyfill-ctype          # qui tire ca, et pourquoi
composer validate --strict
composer dump-autoload -o --classmap-authoritative    # en prod : table classe->fichier, pas de stat
composer audit                               # CVE connues dans l'arbre
composer check-platform-reqs
composer install --no-dev --optimize-autoloader        # la commande de deploiement
\end{terminal}

\begin{codejson}
{
  "name": "nour/blog",
  "description": "Un blog PHP moderne, sans framework.",
  "type": "project",
  "license": "MIT",
  "require": {
    "php": ">=8.3",
    "ext-pdo": "*",
    "ext-mbstring": "*",
    "ext-json": "*",
    "monolog/monolog": "^3.5",
    "symfony/dotenv": "^7.0",
    "vlucas/phpdotenv": "^5.6"
  },
  "require-dev": {
    "phpunit/phpunit": "^11.0",
    "phpstan/phpstan": "^1.11",
    "rector/rector": "^1.2",
    "friendsofphp/php-cs-fixer": "^3.54",
    "squizlabs/php_codesniffer": "^3.9",
    "roave/security-advisories": "dev-latest"
  },
  "autoload": { "psr-4": { "App\\": "src/" } },
  "autoload-dev": { "psr-4": { "Tests\\": "tests/" } },
  "scripts": {
    "test": "phpunit --colors=always",
    "stan": "phpstan analyse --memory-limit=512M",
    "fix": ["php-cs-fixer fix", "rector process"],
    "check": ["@fix", "@stan", "@test"],
    "serve": "php -S 127.0.0.1:8000 -t public public/router.php",
    "post-install-cmd": ["@php bin/console migrate"]
  },
  "scripts-descriptions": { "check": "Formate, analyse, teste." },
  "config": {
    "sort-packages": true,
    "optimize-autoloader": true,
    "allow-plugins": false,
    "platform": { "php": "8.3.0" }
  },
  "minimum-stability": "stable",
  "prefer-stable": true
}
\end{codejson}

\etape{Semver et contraintes---les m\^emes notions qu'en npm, deux pi\`eges en plus}
\begin{itemize}
  \item \verb|^7.0| accepte \verb|7.x| (comme npm). \verb|~7.0| accepte \verb|7.0.x| -- et
        \emph{aussi} \verb|7.1| en philosophie Composer, ce qui surprend~: sois explicite.
  \item \verb|dev-main| et les \verb|@dev| ne se m\'elangent pas avec des versions stables~; \'evite
        les branches stables en production (impossible de reconstruire \`a l'identique sans le lock).
  \item \texttt{composer.lock} se versionne~: pour une \emph{biblioth\`eque} il est ignor\'e, pour une
        \emph{application} il est sacr\'e. C'est la seule r\`egle qui distingue les deux cas.
  \item \verb|platform.php| g\`ele la version \og ~cible~ \fg{}~: utile quand ton poste n'a pas la m\^eme
        version que le serveur (et qu'un paquet exige \og ~php \textgreater=~8.4~ \fg{}).
\end{itemize}

\etape{Publier une biblioth\`eque}
\begin{terminal}
# sur GitHub : composer.json correct, tag sement, README, licence
git tag v1.0.0 && git push --tags
# rendre le paquet decouvrable : https://packagist.org/submit (souvent auto via GitHub Actions)
# chez les autres : composer require nour/blog-utils "^1.0"
# paquet prive : repository "vcs" (https://repo.privie.git) + config gitlab-token,
#   ou Satispress/Private Packagist, ou un depot artefacts.
\end{terminal}

\section{Les standards PSR, en pratique}

\begin{tabz}{llX}
  \headrow \thd{PSR} \& \thd{Nom} \& \thd{Ce que \c{c}a change chez toi} \\
  \midrule
  1/12 \& Coding Style \& \verb|declare(strict\_types=1);|, un fichier = une classe, \verb|namespace App\;|, \verb|use| tri\'es, accolade ouvrante \`a la ligne suivante pour les m\'ethodes, indent 4 espaces. \\
  3 \& Log \& \verb|Psr\\Log\\LoggerInterface| : tu remplaces Monolog par n'importe quoi sans toucher le code. \\
  4 \& Autoload \& un dossier \verb|src/| mapp\'e sur un pr\'efixe d'espace de noms. \\
  6 \& Contexte d'interface \& \verb|getContainer()| etc. : rare en app, fr\'equent en biblioth\`eque. \\
  7/17 \& Messages HTTP \& objets \verb|Request|/\verb|Response| immuables, \verb|ServerRequestInterface| ; PSR-17 les fabrique. \\
  11 \& Conteneur \& \verb|get()|/\verb|has()|~: tous les conteneurs se valent. \\
  14 \& Exceptions \& \verb|ThrowableInterface|~: rare, mais les biblioth\`eques s\'erieuses le font. \\
  15 \& Middlewares \& \verb|MiddlewareInterface|~: le mod\`ele de Slim, Laminas, et des frameworks modernes. \\
  16 \& Cache \& \verb|CacheItemPoolInterface|~: code ton cache une fois, change d'impl\'ementation \`a volont\'e. \\
  18 \& HTTP Client \& \verb|HttpClientInterface|~: d\'ecouple de Guzzle. \\
  20/23 \& Environnement \& erreurs \& lecture d'\verb|.env|~; \verb|bootstrap.php| commun. \\
\end{tabz}

Les liens de lecture~: \link{PHP-FIG}{https://www.php-fig.org/psr/},
\link{PHP~PSR~visual~guide}{https://php-fig.github.io/},
\link{PHP Standard Recommendations sur packagist}{https://packagist.org/providers/psr/log-implementation}.

\section{Style et analyse automatique}

\begin{terminal}
# PHP-CS-Fixer : aligne le style sur un preset (Symfony, PSR-12)
vendor/bin/php-cs-fixer fix --diff --using-cache=no
# PHP_CodeSniffer : verifier sans corriger (id\'eal en CI)
vendor/bin/phpcs --standard=PSR12 --colors src tests
vendor/bin/phpcbf --standard=PSR12 src            # autocorrection
# PHPStan : la vraie valeur ajout\'ee --- trouve les bugs, pas les espaces
vendor/bin/phpstan analyse --level=8 src
# Psalm (analyse plus profonde, plus bavarde)
vendor/bin/psalm --level=2
# Rector : refonte automatique (monter de version PHP, moderniser une API)
vendor/bin/rector process --dry-run
\end{terminal}

\begin{codeconf}
# .php-cs-fixer.dist.php
<?php
$finder = (new PhpCsFixer\Finder())->in(['src', 'tests'])->append(['config']);
return (new PhpCsFixer\Config())
    ->setRiskyAllowed(true)
    ->setRules([
        '@Symfony' => true,
        '@PHP83Migration' => true,
        'declare_strict_types' => true,
        'global_namespace_import' => ['import_classes' => true],
        'ordered_imports' => ['sort_algorithm' => 'alpha'],
        'native_function_invocation' => ['include' => ['@compiler_optimized']],
        'phpdoc_to_comment' => false,
        'yoda_style' => false,
        'no_unused_imports' => true,
        'single_line_throw' => false,
    ])
    ->setFinder($finder);
\end{codeconf}

\begin{codeconf}
# phpstan.neon
parameters:
    level: 8
    paths: [src, tests]
    checkMissingIterableValueType: true
    treatPhpDocTypesAsCertain: false
    reportUnmatchedIgnoredErrors: false
    ignoreErrors:
        - identifier: missingType.iterableValue
    bootstrapFiles: [tests/bootstrap.php]
\end{codeconf}

\begin{note}
Ordre d'adoption qui rapporte~: (1) PHP-CS-Fixer (5 minutes, aucun d\'ebat)~; (2) PHPStan niveau
5, puis 8 progressivement~; (3) \verb|declare(strict\_types=1)| partout~; (4) Rector pour les
mont\'ees de version~; (5) Psalm si le projet est critique. Le niveau 8 de PHPStan est le vrai
d\'etecteur de bugs de nullit\'e, de types fant\^omes et d'oublis de v\'erification---l'\'equivalent du
\verb|strict| de TypeScript, en plus s\'ev\`ere.
\end{note}

\section{Le serveur~: nginx + php-fpm (ou Apache)}

\begin{codeconf}
# /etc/nginx/sites-available/nour
server {
    listen 80;
    server_name nour.dev www.nour.dev;
    root /var/www/nour/public;                 # LA RACINE WEB = public/ , jamais le projet
    index index.php;

    location / {
        try_files $uri $uri/ /index.php$is_args$args;   # le "front controller"
    }
    location ~ \.php$ {
        include fastcgi_params;
        fastcgi_pass unix:/run/php/php8.3-fpm.sock;
        fastcgi_param SCRIPT_FILENAME $realpath_root$fastcgi_script_name;
        fastcgi_param DOCUMENT_ROOT $realpath_root;
        fastcgi_read_timeout 60s;
    }
    location = /favicon.ico { log_not_found off; access_log off; }
    location ~* \.(css|js|woff2|svg|png|jpg|webp)$ {
        expires 1y; add_header Cache-Control "public, immutable"; access_log off;
    }
    location ~ /\.(?!well-known) { deny all; }
}
\end{codeconf}

\begin{codeconf}
; /etc/php/8.3/fpm/pool.d/nour.conf  -- un pool par application
[nour]
user = nour
group = nour
listen = /run/php/php8.3-fpm-nour.sock
listen.owner = www-data
listen.group = www-data
pm = dynamic
pm.max_children = 25                 ; memoire_par_processus ~60 Mo => (RAM*0.8)/60
pm.start_servers = 6
pm.min_spare_servers = 4
pm.max_spare_servers = 10
pm.max_requests = 500                 ; recycle pour limiter les fuites
clear_env = no                        ; sans cela, tes variables d'environnement ne passent PAS
php_admin_value[upload_max_filesize] = 10M
php_admin_value[memory_limit] = 256M
php_admin_value[error_log] = /var/log/php/nour-error.log
catch_workers_output = yes
decorate_workers_output = no
\end{codeconf}

\begin{terminal}
sudo systemctl reload php8.3-fpm nginx
php -i | grep -E "opcache.enable|memory"        # OPcache doit etre actif
# OPcache en production (les memes opcodes recompiles a chaque requete = gaspillage) :
#   opcache.enable=1, opcache.memory_consumption=192, opcache.max_accelerated_files=20000,
#   opcache.validate_timestamps=0  (+ `php-fpm -tt && kill -USR2 <master>` apres un deploiement)
\end{terminal}

\section{D\'ebogueur Xdebug---et comment s'en servir sans pleurer}

\begin{terminal}
# Linux (paquet) : sudo apt install php-xdebug ; ou via le gestionnaire de la distro
# macOS : pecl install xdebug  |  Windows : download.coderix.com/xdebug (choisir le .dll
#         correspondant a php -i : version, NTS/TS, xdebug)
php --ini ; php -i | grep xdebug.mode
\end{terminal}

\begin{codeconf}
; /etc/php/8.3/mods-available/xdebug.ini  (ou dans php.ini)
zend_extension=xdebug
xdebug.mode=debug
xdebug.start_with_request=trigger        ; ne s'active que si XDEBUG_TRIGGER present : + de controle
xdebug.client_host=127.0.0.1
xdebug.client_port=9003
xdebug.idekey=VSCODE
xdebug.log=/tmp/xdebug.log
xdebug.max_nesting_level=512
xdebug.profiler_output_dir=/tmp
xdebug.mode=profile,debug                 ; les deux a la fois en dev local seulement
\end{codeconf}

\begin{codejson}
{
  "name": "Ecouter Xdebug",
  "type": "php",
  "request": "launch",
  "port": 9003,
  "pathMappings": { "/var/www/nour": "${workspaceFolder}" },
  "log": true
}
\end{codejson}

D\'eclencher une session~: \link{xdebug-trigger}{https://xdebug.org/wizard}~; en pratique~: extension
\og ~Xdebug Helper~ \fg{}, ou \cmd{curl -H "XDEBUG-TRIGGER: 1" http://nour.dev/article?id=7}.
\verb|xdebug\_break()| dans le code, \verb|xdebug\_print\_function\_stack()| pour voir qui m'a appel\'e,
\verb|xdebug\_start\_code\_coverage()| pour mesurer ce qui manque.

\subsection{Profilage~: o\`u part le temps}

\begin{terminal}
# 1. Xdebug profile (cachegrind) puis QCachegrind / KCacheGrind / WebCacheGrind
curl -H "XDEBUG-PROFILE: 1" http://nour.dev/rapport -o /dev/null
kcachegrind /tmp/cachegrind.out.*

# 2. blackfire.io (versions locales gratuites) ou xhprof
composer require --dev symfony/xhprof

# 3. Le minitraceur maison, 10 lignes, suffisant 9 fois sur 10
#    $t=microtime(true); ... ; error_log(sprintf('[perf] %s: %.1f ms', $etape, 1e3*(microtime(true)-$t)));
\end{terminal}

\section{Tests PHP~: PHPUnit (et Pest)}

\begin{codephp}
<?php
declare(strict_types=1);
namespace Tests\Unit\Domain;

use App\Domain\Article\Article;
use App\Domain\Article\Statut;
use DateTimeImmutable;
use PHPUnit\Framework\Attributes\CoversClass;
use PHPUnit\Framework\Attributes\DataProvider;
use PHPUnit\Framework\TestCase;

#[CoversClass(Article::class)]
final class ArticleTest extends TestCase
{
    private function article(string $titre = 'Un titre correct', Statut $s = Statut::Brouillon): Article
    {
        return new Article(null, $titre, 'corps', $s, new Auteur('N', 'n@e.com'), new DateTimeImmutable());
    }

    public function testLaPublicationFixeLaDate(): void
    {
        $a = $this->article()->publier();
        self::assertSame(Statut::Publie, $a->statut());
        self::assertNotNull($a->publieLe());
    }

    #[DataProvider('titresInvalides')]
    public function testRefuseLesTitresInvalides(string $titre, string $attendu): void
    {
        $this->expectException(\DomainException::class);
        $this->expectExceptionMessageMatches($attendu);
        $this->article($titre);
    }
    public static function titresInvalides(): iterable
    {
        yield 'vide' => ['', '/minimum/'];
        yield 'trop court' => ['ab', '/minimum/'];
        yield 'trop long' => [str_repeat('a', 200), '/maximum/'];
    }

    public function testRenommerNeMutePasLOriginal(): void
    {
        $original = $this->article('A');
        $modifie = $original->renommer('B');
        self::assertNotSame($original, $modifie);
        self::assertSame('A', $original->titre());
    }
}
\end{codephp}

\begin{codephp}
<?php
// tests/bootstrap.php --- le minimum pour tester sans toucher a la vraie base
require __DIR__ . '/../vendor/autoload.php';
putenv('APP_ENV=test');
// SQLite en memoire pour les tests d'integration : rapide, isole, reel
$pdo = new PDO('sqlite::memory:');
$pdo->exec(file_get_contents(__DIR__ . '/../database/schema.sqlite.sql'));
\end{codephp}

\begin{codeconf}
; phpunit.xml.dist
<?xml version="1.0" encoding="UTF-8"?>
<phpunit xmlns:xsi="http://www.w3.org/2001/XMLSchema-instance"
         xsi:noNamespaceSchemaLocation="vendor/phpunit/phpunit/phpunit.xsd"
         bootstrap="tests/bootstrap.php" colors="true" cacheDirectory=".phpunit.cache"
         failOnWarning="true" failOnNotice="true" beStrictAboutOutputDuringTests="true">
  <testsuites>
    <testsuite name="unit"><directory>tests/Unit</directory></testsuite>
    <testsuite name="integration"><directory>tests/Integration</directory></testsuite>
  </testsuites>
  <source>
    <include><directory>src</directory></include>
  </source>
  <coverage>
    <report><text outputFile="php://stdout"/><html outputDirectory="coverage"/></report>
  </coverage>
  <php>
    <ini name="error_reporting" value="-1"/>
    <server name="APP_ENV" value="test" force="true"/>
  </php>
</phpunit>
\end{codeconf}

\begin{terminal}
vendor/bin/phpunit --testdox --coverage-text
vendor/bin/phpunit --filter=Article --coverage-clover coverage.xml
vendor/bin/phpunit --order-by=defects --reverse-list    # les echecs d'abord
# Pest (syntaxe plus courte, tres appreciee) : composer require --dev pestphp/pest && vendor/bin/pest
\end{terminal}

\subsection{Paquets PHP qui valent leur poids}

\begin{tabz}{lX}
  \headrow \thd{Paquet} \& \thd{Usage} \\
  \midrule
  \link{monolog/monolog}{https://github.com/Seldaek/monolog} \& journalisation PSR-3, handlers fichier/Sentry/Slack.\\
  \link{symfony/dotenv}{https://symfony.com/components/Dotenv} \& lecture de \fichier{.env} par environnements.\\
  \link{nesbot/carbon}{https://carbon.nesbot.com/} \& dates agr\'eables (et \verb|DateTimeImmutable| natif, souvent suffisant).\\
  \link{ramsey/uuid}{https://github.com/ramsey/uuid} \& UUIDv7/v4.\\
  \link{vlucas/phpdotenv}{https://github.com/vlucas/phpdotenv} \& l'alternative simple.\\
  \link{guzzlehttp/guzzle}{https://docs.guzzlephp.org/} ou \link{symfony/http-client}{https://symfony.com/components/HttpClient} \& client HTTP.\\
  \link{phpmailer/phpmailer}{https://github.com/PHPMailer/PHPMailer} ou \link{symfony/mailer}{https://symfony.com/doc/current/mailer.html} \& emails.\\
  \link{respect/validation}{https://respectvalidation.readthedocs.io/} \& validation d\'eclarative.\\
  \link{league/flysystem}{https://flysystem.thephpleague.com/} \& stockage abstrait (local, S3, FTP).\\
  \link{symfony/console}{https://symfony.com/components/Console} \& commandes CLI, sans framework.\\
  \link{robmorgan/phinx}{https://book.cakephp.org/phinx/} \& migrations.\\
  \link{mockery/mockery}{https://docs.mockery.io/} \& mocks (surtout pour les interfaces).\\
  \link{friendsofphp/php-cs-fixer}{https://github.com/PHP-CS-Fixer/PHP-CS-Fixer} \& style.\\
  \link{rector/rector}{https://getrector.com/} \& refactorisations en masse.\\
  \link{filp/whoops}{https://github.com/filp/whoops} \& page d'erreur de d\'eveloppement magnifique.\\
\end{tabz}

\begin{exercice}
\begin{enumerate}
  \item Cr\'ee un projet \og ~sans framework~ \fg{} propre~: \verb|composer init|, PSR-4, \verb|public/index.php|,
        \verb|src/|, \verb|tests/|, \verb|data/|, scripts \verb|serve|, \verb|check|, \verb|test|.
  \item Ajoute PHP-CS-Fixer et PHPStan niveau~8~: fais passer le projet au vert en partant de
        \og ~plein d'erreurs~ \fg{}. Note les bugs r\'eels trouv\'es par PHPStan (il y en a).
  \item \'Ecris 20 tests PHPUnit~: d'abord sur le domaine pur (aucune d\'ependance), puis sur le d\'ep\^ot
        avec SQLite en m\'emoire, puis un test d'int\'egration HTTP (une vraie requ\^ete sur le routeur).
  \item Configure Xdebug~: arr\^ete-toi dans une action, inspecte la pile, utilise
        \verb|xdebug\_break()| dans une fermeture. Puis profile une page lente et nomme la fonction
        qui co\^ute~; corrige, reprofile.
  \item Rends ton d\'ep\^ots compatible PSR-3 et PSR-16~: ajoute un cache m\'emoire avec TTL~; passe
        du \verb|apcu| au \verb|files| sans changer une ligne de m\'etier.
  \item Publie une petite biblioth\`eque sur Packagist~: tests, CI, style, analyse, tag, README~;
        installe-la dans un autre projet.
\end{enumerate}
\end{exercice}

\begin{ressources}
\begin{itemize}
  \item Composer documentation --- \url{https://getcomposer.org/doc/}
  \item PHP-FIG --- \url{https://www.php-fig.org/psr/}
  \item PHPStan --- \url{https://phpstan.org/} --- Psalm --- \url{https://psalm.dev/} --- Rector --- \url{https://getrector.com/}
  \item PHPUnit --- \url{https://docs.phpunit.de/} --- Pest --- \url{https://pestphp.com/}
  \item Xdebug --- \url{https://xdebug.org/docs/step\\\_debug} --- PHP-FPM tuning --- \url{https://php.watch/configuration/php-fpm}
  \item \link{PHP~The~Right~Way}{https://phptherightway.com/} --- relis la partie \og ~Dependencies~ \fg{} et \og ~Tooling~ \fg{}.
\end{itemize}
\end{ressources}

\index{Composer}\index{PSR}\index{PHPStan}\index{Xdebug}\index{PHPUnit}
% ============================================================================
\chapter{Construire une API REST en PHP (et l'authentifier avec un JWT)}
\label{ch:php-api}
% ============================================================================

\begin{objectif}
\begin{itemize}
  \item Concevoir une API coh\'erente~: ressources, verbes, codes, pagination, erreurs, versionnage.
  \item \'Ecrire en PHP pur un mini-framework~: routeur, middlewares, validation, s\'erialisation.
  \item Authentifier par jeton (JWT) et par cookie, et savoir quand employer l'un ou l'autre.
  \item Livrer une documentation exploitable (OpenAPI) et un client de test.
  \item Prot\'eger le point d'entr\'ee~: CORS, d\'ebit, tailles, quotas, idempotence.
\end{itemize}
\end{objectif}

\section{D'abord le contrat (10 minutes qui en sauvent 100)}

\begin{tabz}{lXX}
  \headrow \thd{Question} \& \thd{R\'eponse par d\'efaut} \& \thd{Et si tu fais autrement} \\
  \midrule
  Mod\`ele d'URL \& \verb|/api/v1/ressource/{id}/sous-ressource| \& \verb|/articles/lire?id=..| devient invendable. \\
  Verbes \& GET lire, POST cr\'eer, PUT remplacer, PATCH modifier partiellement, DELETE supprimer \& un POST partout cache la s\'emantique et casse le cache. \\
  Codes \& 200, 201, 204, 400, 401, 403, 404, 409, 422, 429, 500 \& 200 avec \verb|{erreur}| force le client \`a lire le corps pour savoir. \\
  Erreurs \& RFC 9457 \verb|application/problem+json| \& un format maison par endpoint~: enfer. \\
  Pagination \& \verb|?page=\&limit=| + en-t\^ete \verb|Link|, ou curseur pour les gros volumes \& sans limite maximale, une requ\^ete tue la base. \\
  Dates \& ISO~8601 UTC (\verb|2026-01-31T10:00:00Z|) \& \verb|31/01/2026|~: ambiguity garantie. \\
  Ids \& opaque (\verb|string| ou UUIDv7) ; jamais auto-incr\'ement\'e expos\'e seul \& \'enum\'erable : les voisins lisent ton catalogue. \\
  Version \& dans l'URL \verb|/v1/| ou un en-t\^ete accept ; une seule r\`egle, la m\^eme partout \& \verb|?v=1| partout : le chaos. \\
\end{tabz}

\begin{codesql}
-- Exemple de contrat d'erreur unique (RFC 9457) :
{
  "type": "https://api.nour.dev/problemes/validation",
  "title": "Donnees invalides",
  "status": 422,
  "detail": "3 champs rejettent la validation.",
  "instance": "/api/v1/articles",
  "trace": "a91f3c",
  "errors": [
    { "champ": "titre", "code": "trop_court", "message": "Au moins 3 caracteres." },
    { "champ": "tags[2]", "code": "trop_long", "message": "30 caracteres maximum." }
  ]
}
\end{codesql}

\section{Le mini-framework~: 150 lignes qui expliquent tout}

\begin{codephp}
<?php
// src/Noyau/App.php --- un noyau d'API, sans dependance, a comprendre ligne par ligne
declare(strict_types=1);
namespace App\Noyau;

use JsonException;

final class Requete
{
    /** @param array<string,mixed> $params @param array<string,mixed> $corps */
    public function __construct(
        public readonly string $methode,
        public readonly string $chemin,
        public readonly array $params = [],
        public readonly array $corps = [],
        public readonly array $query = [],
        public readonly array $entetes = [],
        public ?object $utilisateur = null,
    ) {}
    public function entete(string $nom): ?string { return $this->entetes[strtolower($nom)] ?? null; }
}

final class Reponse
{
    /** @param array<string,string> $entetes */
    public function __construct(
        public readonly int $code = 200,
        public readonly mixed $corps = null,
        public readonly array $entetes = [],
        public readonly ?string $brut = null,
    ) {}
    public static function json(mixed $donnees, int $code = 200, array $entetes = []): self
    {
        return new self($code, $donnees, ['content-type' => 'application/json; charset=utf-8'] + $entetes);
    }
    public static function probleme(int $code, string $titre, string $detail = '', array $extra = []): self
    {
        return self::json(array_merge([
            'type' => 'https://api.nour.dev/problemes/' . slug($titre),
            'title' => $titre, 'status' => $code, 'detail' => $detail,
        ], $extra), $code, ['content-type' => 'application/problem+json; charset=utf-8']);
    }
    public static function vide(int $code = 204): self { return new self($code); }
}

function slug(string $s): string
{
    return strtolower(preg_replace('/[^a-z0-9]+/i', '-', iconv('UTF-8', 'ASCII//TRANSLIT', $s) ?: $s) ?? '') ?? 'erreur';
}

final class App
{
    /** @var array<int, array{methode:string, motif:string, gestionnaire:callable, middlewares:array}> */
    private array $routes = [];
    /** @var callable[] */ private array $middlewares = [];

    public function get(string $chemin, callable $h): void { $this->ajouter('GET', $chemin, $h); }
    public function post(string $chemin, callable $h): void { $this->ajouter('POST', $chemin, $h); }
    public function patch(string $chemin, callable $h): void { $this->ajouter('PATCH', $chemin, $h); }
    public function delete(string $chemin, callable $h): void { $this->ajouter('DELETE', $chemin, $h); }

    public function middleware(callable $m): void { $this->middlewares[] = $m; }

    private function ajouter(string $methode, string $chemin, callable $h): void
    {
        $motif = preg_replace_callback('/\{(\w+)(?::([^}]+))?\}/', static function (array $m): string {
            return sprintf('(?<%s>%s)', $m[1], $m[2] ?? '[^/]+');
        }, $chemin);
        $this->routes[] = ['methode' => $methode, 'motif' => '#^' . $motif . '$#', 'gestionnaire' => $h, 'middlewares' => []];
    }

    public function traiter(): void
    {
        $methode = $_SERVER['REQUEST_METHOD'] ?? 'GET';
        $chemin = parse_url($_SERVER['REQUEST_URI'] ?? '/', PHP_URL_PATH) ?: '/';
        $entetes = [];
        foreach ($_SERVER as $k => $v) {
            if (str_starts_with((string) $k, 'HTTP_')) {
                $entetes[strtolower(str_replace('_', '-', substr((string) $k, 5)))] = (string) $v;
            }
        }
        $corps = $this->lireCorps($entetes);

        foreach ($this->routes as $r) {
            if ($r['methode'] !== $methode || !preg_match($r['motif'], $chemin, $m)) { continue; }
            $params = array_filter($m, 'is_string', ARRAY_FILTER_USE_KEY);
            $req = new Requete($methode, $chemin, $params, $corps, $_GET, $entetes);
            $chain = array_reduce(
                array_reverse([...$this->middlewares, static fn (Requete $q) => ($r['gestionnaire'])($q)]),
                static fn (callable $suivant, callable $mw): callable => static fn (Requete $q): Reponse => $mw($q, $suivant),
                static fn (Requete $q): Reponse => new Reponse(500)
            );
            $this->envoyer($chain($req));
            return;
        }
        $this->envoyer(Reponse::probleme(404, 'Ressource introuvable', $methode . ' ' . $chemin));
    }

    /** @return array<string,mixed> */
    private function lireCorps(array $entetes): array
    {
        $type = $entetes['content-type'] ?? '';
        $brut = file_get_contents('php://input') ?: '';
        if (str_contains($type, 'json')) {
            if ($brut === '') { return []; }
            try { return (array) json_decode($brut, true, 32, JSON_THROW_ON_ERROR); }
            catch (JsonException $e) { throw new ErreurRequete('JSON invalide : ' . $e->getMessage(), 400); }
        }
        if (str_contains($type, 'x-www-form-urlencoded')) { parse_str($brut, $out); return $out; }
        return $_POST;
    }

    private function envoyer(Reponse $r): void
    {
        http_response_code($r->code);
        foreach ($r->entetes as $nom => $valeur) { header($nom . ': ' . $valeur, true); }
        if ($r->corps !== null) {
            echo json_encode($r->corps, JSON_UNESCAPED_UNICODE | JSON_UNESCAPED_SLASHES | JSON_THROW_ON_ERROR);
        }
    }
}

final class ErreurRequete extends \RuntimeException { public function __construct(string $m, public readonly int $code = 400) { parent::__construct($m); } }
\end{codephp}

\begin{codephp}
<?php
// public/index.php --- branchement : middlewares, gestion d'erreur, routes
declare(strict_types=1);
require __DIR__ . '/../vendor/autoload.php';

use App\Noyau\{App, Requete, Reponse, ErreurRequete};

$app = new App();
$app->middleware(function (Requete $r, callable $suite): Reponse {
    header('x-request-id: ' . bin2hex(random_bytes(6)));
    header('content-type: application/json; charset=utf-8');
    return $suite($r);
});
$app->middleware(App\Securite::cors(...));                 // cf. chapitre securite
$app->middleware(App\Securite::limiteurDebit(120, 60));    // 120 req/min/IP (apcu)
$app->middleware(App\Auth::porte('bearer'));                // remplit $r->utilisateur si jeton valide

$articles = new App\Depot\DepotArticles(pdo());
$valider = new App\Validation\Valideur();

$app->get('/api/v1/sante', fn () => Reponse::json(['ok' => true, 'php' => PHP_VERSION, 'base' => pdo()->query('SELECT 1')->fetchColumn() ? 'ok' : 'ko']));

$app->get('/api/v1/articles', function (Requete $r) use ($articles, $valider) {
    $q = $valider->verifier($r->query, [
        'page' => 'int|min:1', 'limite' => 'int|min:1|max:100', 'recherche' => 'string|max:120',
    ]);
    return Reponse::json($articles->pageListe($q['page'] ?? 1, $q['limite'] ?? 20, $q['recherche'] ?? ''));
});

$app->post('/api/v1/articles', function (Requete $r) use ($articles, $valider) {
    $donnees = $valider->verifier($r->corps, [
        'titre' => 'requis|string|min:3|max:160',
        'corps' => 'requis|string|min:20',
        'tags' => 'array|max:12',
    ], [
        'tags.*' => 'string|min:1|max:30',
    ]);
    if ($r->utilisateur === null) { return Reponse::probleme(401, 'Connexion requise'); }
    $cree = $articles->creer($donnees, $r->utilisateur);
    return Reponse::json($cree, 201)->avecEntete('location', '/api/v1/articles/' . $cree['id']);
});

$app->patch('/api/v1/articles/{id:\d+}', function (Requete $r) use ($articles) {
    $maj = $articles->modifier((int) $r->params['id'], $r->corps, $r->utilisateur);
    return $maj === null ? Reponse::probleme(404, 'Introuvable') : Reponse::json($maj);
});

$app->delete('/api/v1/articles/{id:\d+}', fn (Requete $r) => $articles->supprimer((int) $r->params['id'], $r->utilisateur)
    ? Reponse::vide() : Reponse::probleme(404, 'Introuvable'));

try {
    $app->traiter();
} catch (ErreurRequete $e) {
    (function (Reponse $r): void { http_response_code($r->code); foreach ($r->entetes as $n => $v) { header("$n: $v"); } echo json_encode($r->corps); })
    (Reponse::probleme($e->code, 'Requete invalide', $e->getMessage()));
} catch (Throwable $e) {
    error_log((string) $e);
    (static function (): void { http_response_code(500); header('content-type: application/problem+json');
        echo json_encode(['title' => 'Erreur interne', 'status' => 500, 'trace' => substr(hash('crc32b', (string) $e), 0, 8)]); })();
}
exit;
\end{codephp}

\subsection{Un validateur minimal}

\begin{codephp}
<?php
declare(strict_types=1);
namespace App\Validation;

final class Valideur
{
    /** @param array<string,string> $regles  @return array<string,mixed> @throws ErreurValidation */
    public function verifier(array $donnees, array $regles, array $reglesNichees = []): array
    {
        $erreurs = []; $out = [];
        foreach ($regles as $champ => $chaine) {
            $r\`egles = explode('|', $chaine);
            $valeur = $donnees[$champ] ?? null;
            $obligatoire = in_array('requis', $r\`egles, true);
            if ($valeur === null || $valeur === '') {
                if ($obligatoire) { $erreurs[$champ][] = 'Champ obligatoire.'; }
                continue;
            }
            foreach ($r\`egles as $r) {
                [$nom, $arg] = array_pad(explode(':', $r, 2), 2, null);
                $message = match ($nom) {
                    'requis' => null,
                    'string' => is_string($valeur) ? null : 'Doit etre une chaine.',
                    'int' => filter_var($valeur, FILTER_VALIDATE_INT) === false ? 'Nombre entier attendu.' : null,
                    'bool' => null,
                    'array' => is_array($valeur) ? null : 'Liste attendue.',
                    'email' => filter_var($valeur, FILTER_VALIDATE_EMAIL) === false ? 'Email invalide.' : null,
                    'url' => filter_var($valeur, FILTER_VALIDATE_URL) === false ? 'URL invalide.' : null,
                    'min' => mb_strlen((string) $valeur) < (int) $arg ? sprintf('Minimum %s caract\`eres.', $arg) : null,
                    'max' => mb_strlen((string) $valeur) > (int) $arg ? sprintf('Maximum %s caract\`eres.', $arg) : null,
                    'max_count' => is_array($valeur) && count($valeur) > (int) $arg ? sprintf('Maximum %s elements.', $arg) : null,
                    'in' => !in_array((string) $valeur, explode(',', (string) $arg), true) ? 'Valeur non autoris\'ee.' : null,
                    default => throw new InvalidArgumentException("regle inconnue : {$r}"),
                };
                if ($message !== null) { $erreurs[$champ][] = $message; break; }
            }
            $out[$champ] = match (true) {
                in_array('int', $r\`egles, true) => (int) $valeur,
                in_array('bool', $r\`egles, true) => filter_var($valeur, FILTER_VALIDATE_BOOL),
                in_array('string', $r\`egles, true) => trim((string) $valeur),
                default => $valeur,
            };
        }
        if ($erreurs !== []) { throw new ErreurValidation($erreurs); }
        return $out;
    }
}

final class ErreurValidation extends \RuntimeException
{
    /** @param array<string,list<string>> $erreurs */
    public function __construct(public readonly array $erreurs)
    {
        parent::__construct('Validation echouee');
    }
    public function versProblem(): array
    {
        return ['errors' => array_map(
            static fn (array $l): array => ['message' => $l[0]],
            $this->erreurs
        )];
    }
}
\end{codephp}

\section{JWT~: le fabriquer, le v\'erifier, et savoir quand s'en passer}

\begin{codephp}
<?php
declare(strict_types=1);
namespace App\Auth;

/**
 * JWT minimal (HS256) en 40 lignes : tu comprends enfin ce que font les bibliotheques.
 * En production, prends firebase/php-jwt ou lcobucci/jwt -- et surtout la config qui va avec.
 */
final class Jeton
{
    public function __construct(private readonly string $secret, private readonly int $duree = 900) {}

    /** @param array<string,mixed> $revendications */
    public function emettre(int $utilisateurId, array $revendications = []): string
    {
        $entete = ['typ' => 'JWT', 'alg' => 'HS256'];
        $maintenant = time();
        $corps = [
            'iss' => 'https://api.nour.dev',
            'aud' => 'nour-web',
            'sub' => (string) $utilisateurId,
            'iat' => $maintenant,
            'nbf' => $maintenant,
            'exp' => $maintenant + $this->duree,
            'jti' => bin2hex(random_bytes(8)),
        ] + $revendications;
        $segments = [
            self::b64($entete),
            self::b64($corps),
        ];
        $signature = hash_hmac('sha256', implode('.', $segments), $this->secret, true);
        $segments[] = self::b64($signature, brut: true);
        return implode('.', $segments);
    }

    /** @return array<string,mixed> @throws ErreurJeton */
    public function verifier(string $jeton): array
    {
        $parties = explode('.', $jeton);
        if (count($parties) !== 3) { throw new ErreurJeton('Format invalide'); }
        [$h, $c, $s] = $parties;
        $attendu = hash_hmac('sha256', $h . '.' . $c, $this->secret, true);
        $recu = base64_decode(strtr($s, '-_', '+/'), true);
        if ($recu === false || !hash_equals(hash('sha256', $attendu), hash('sha256', $recu))) {
            throw new ErreurJeton('Signature invalide');
        }
        $entete = self::deb64($h);
        // Contre l'attaque "alg: none" : algorithme epingle
        if (($entete['alg'] ?? null) !== 'HS256') { throw new ErreurJeton('Algorithme non autorise'); }
        $corps = self::deb64($c);
        $now = time();
        foreach (['nbf' => '<', 'iat' => '<', 'exp' => '>'] as $cle => $sens) {
            if (isset($corps[$cle]) && ($sens === '<' ? $now < (int) $corps[$cle] : $now > (int) $corps[$cle])) {
                throw new ErreurJeton("Jeton non valide pour {$cle}");
            }
        }
        foreach (['iss' => 'https://api.nour.dev', 'aud' => 'nour-web'] as $cle => $valeur) {
            if (($corps[$cle] ?? null) !== $valeur) { throw new ErreurJeton("Revendication {$cle} invalide"); }
        }
        return $corps;
    }

    private static function b64(array|string $donnees, bool $brut = false): string
    {
        $chaine = $brut ? $donnees : json_encode($donnees, JSON_UNESCAPED_SLASHES | JSON_THROW_ON_ERROR);
        return rtrim(strtr(base64_encode($chaine), '+/', '-_'), '=');
    }
    private static function deb64(string $s): array
    {
        $d = base64_decode(strtr($s, '-_', '+/'), true);
        if ($d === false) { throw new ErreurJeton('Base64 invalide'); }
        return (array) json_decode($d, true, 16, JSON_THROW_ON_ERROR);
    }
}
final class ErreurJeton extends \RuntimeException {}

// Le middleware : autorisation = verifier le jeton, puis CHARGER l'utilisateur en base.
// Ne fais jamais confiance aux roles stockes dans le jeton pour une decision sensible
// (le jeton vit 15 min ; un role revoque doit etre effectif immediatement).
final class Auth
{
    public static function porte(string $type): callable
    {
        return static function (\App\Noyau\Requete $r, callable $suite) use ($type): \App\Noyau\Reponse {
            $jeton = null;
            if ($type === 'bearer' && preg_match('/^Bearer\s+(\S+)$/i', (string) $r->entete('authorization'), $m)) {
                $jeton = $m[1];
            }
            if ($jeton !== null) {
                try {
                    $corps = (new Jeton($_ENV['JWT_SECRET']))->verifier($jeton);
                    $r = $r->avecUtilisateur(Utilisateurs::trouver((int) $corps['sub']));
                } catch (ErreurJeton $e) {
                    return \App\Noyau\Reponse::probleme(401, 'Jeton invalide', $e->getMessage());
                }
            }
            return $suite($r);
        };
    }
}
\end{codephp}

\begin{tabz}{llX}
  \headrow \thd{Point} \& \thd{Cookie de session} \& \thd{JWT bearer} \\
  \midrule
  R\'evoquer \& instantan\'e (on supprime la ligne) \& impossible avant expiration~: pr\'evoir une liste de r\'evocation \\
  CSRF \& \`a prot\'eger (jeton) \& quasi inexistant (le navigateur ne l'ajoute pas tout seul) \\
  Multi-sous-domaines \& possible (\verb|Domain=.nour.dev|) \\
  Mobile / API tierces \& awkward \& naturel \\
  \'Etat serveur \& requis (Redis/BDD) \& aucun (mais cache quand m\^eme) \\
  S\'ecurit\'e du stockage \& le JS ne le voit jamais (\verb|HttpOnly|) \& dans \verb|localStorage| = vol\'e par le moindre XSS \\
\end{tabz}

\begin{astuce}
Pour une application web JS + PHP~: \textbf{cookie de session} (Simple, CSRF pr\`es, r\'evocable) pour le
navigateur, \textbf{JWT} pour le mobile et les int\'egrations tierces. Les deux pointent sur le m\^eme
contr\^oleur~: le middleware produit un objet utilisateur, le reste du code ne sait pas comment il a \'et\'e
authentifi\'e.
\end{astuce}

\subsection{CORS, quotas, idempotence~: les trois finitions}

\begin{codephp}
<?php
final class Securite
{
    public static function cors(\App\Noyau\Requete $r, callable $suite): \App\Noyau\Reponse
    {
        $origines = ['http://localhost:5173', 'https://nour.dev'];
        $o = $r->entete('origin');
        $autorisee = $o !== null && in_array($o, $origines, true);
        $enTetes = static function (array $h) use ($autorisee, $o): array {
            if (!$autorisee) { return $h; }
            return $h + [
                'access-control-allow-origin' => $o,
                'vary' => 'origin',
                'access-control-allow-credentials' => 'true',
                'access-control-max-age' => '600',
            ];
        };
        if ($r->methode === 'OPTIONS') {
            return new \App\Noyau\Reponse(204, null, $enTetes([
                'access-control-allow-methods' => 'GET,POST,PUT,PATCH,DELETE,OPTIONS',
                'access-control-allow-headers' => 'content-type,authorization,idempotency-key',
            ]));
        }
        return $suite($r);
    }

    public static function limiteurDebit(int $max, int $fenetre): callable
    {
        return static function (\App\Noyau\Requete $r, callable $suite) use ($max, $fenetre): \App\Noyau\Reponse {
            if (function_exists('apcu_add') && $r->utilisateur === null) {
                $cle = 'rl:' . ($_SERVER['REMOTE_ADDR'] ?? '?') . ':' . (int) floor(time() / $fenetre);
                if (!apcu_add($cle, 1, $fenetre)) {
                    $n = apcu_inc($cle);
                    if ($n > $max) {
                        return \App\Noyau\Reponse::probleme(429, 'Trop de requetes', "limite : {$max}/{$fenetre}s",
                            ['retry_after' => $fenetre - (time() % $fenetre)]);
                    }
                }
            }
            return $suite($r);
        };
    }
}

// Idempotence pour les POST critiques (paiement, envoi d'email) :
// une cle client -> on rejoue la meme reponse au lieu de creer un doublon.
function avecCleIdempotence(callable $creation): callable
{
    return static function (array $donnees, string $cle) use ($creation) {
        $pdo = pdo();
        $pdo->prepare('INSERT IGNORE INTO idempotence (cle, cree_le) VALUES (?, NOW())')->execute([$cle]);
        $dejaLa = $pdo->prepare('SELECT reponse FROM idempotence WHERE cle = ? AND reponse IS NOT NULL');
        $dejaLa->execute([$cle]);
        if ($deja = $dejaLa->fetchColumn()) { return json_decode((string) $deja, true, 512, JSON_THROW_ON_ERROR); }
        $reponse = $creation($donnees);
        $pdo->prepare('UPDATE idempotence SET reponse = ? WHERE cle = ?')
            ->execute([json_encode($reponse, JSON_THROW_ON_ERROR), $cle]);
        return $reponse;
    };
}
\end{codephp}

\section{Documenter et tester~: OpenAPI + un client}

\begin{codeyml}
# openapi.yaml (extrait) --- a editer ou a generer depuis les attributs + zircote/swagger-php
openapi: 3.1.0
info: { title: "API Nour", version: "1.0.0", description: "Articles, tags, auth." }
servers: [{ url: "https://api.nour.dev/v1" }]
paths:
  /articles:
    get:
      summary: Liste paginee
      parameters:
        - { name: page, in: query, schema: { type: integer, minimum: 1, default: 1 } }
        - { name: limite, in: query, schema: { type: integer, maximum: 100, default: 20 } }
      responses:
        "200":
          description: OK
          content:
            application/json:
              schema: { $ref: "#/components/schemas/PageArticles" }
        "429": { $ref: "#/components/responses/TropDeRequetes" }
components:
  schemas:
    Article:
      type: object
      required: [id, titre, statut]
      properties:
        id: { type: integer, format: int64, readOnly: true }
        titre: { type: string, minLength: 3, maxLength: 160 }
        statut: { type: string, enum: [brouillon, publie, archive] }
        publie_le: { type: string, format: date-time, nullable: true }
\end{codeyml}

\begin{terminal}
# Generer la spec depuis le code : composer require --dev zircote/swagger-php
vendor/bin/openapi src -o public/openapi.json --format yaml
# Servir la doc : https://github.com/scalar/scalar (Scalar, le meilleur lecteur moderne)
npx @scalar/api-reference --url /openapi.json
# Valider la spec + generer des exemples de contrats :
npx @redocly/cli lint openapi.yaml
# Client de test dans VS Code : api.http
\end{terminal}

\begin{codeplain}
### api.http --- extension "REST Client" (humao.rest-client)
@base = http://127.0.0.1:8000/api/v1

GET {{base}}/sante

###
# @name login
POST {{base}/auth/connexion
content-type: application/json

{ "email": "nour@example.com", "mot_de_passe": "un mot de passe costaud" }

###
@token = {{login.response.body.jeton}}

GET {{base}}/articles?page=1
authorization: Bearer {{token}}
x-request-id: 01HTQ...

###
POST {{base}}/articles
authorization: Bearer {{token}}
idempotency-key: 01HTQ6M2K9
content-type: application/json

{ "titre": "Apprendre PHP vraiment", "corps": "Un article long et honnete...", "tags": ["php"] }
\end{codeplain}

\begin{projet}[API compl\`ete]
Construis l'API de ton blog~: authentification (inscription, connexion, refresh, d\'econnexion,
r\'einitialisation), articles CRUD avec upload d'image de couverture, tags, commentaires
(mod\'erables), recherche, pagination par curseur, quota par utilisateur, OpenAPI servi sur
\verb|/docs|, 40 tests d'int\'egration (HTTP r\'eel + SQLite en m\'emoire), et un client JS g\'en\'er\'e
\`a partir de la spec (\link{openapi-typescript-codegen}{https://github.com/feross/openapi-typescript-codegen}
ou \link{oazapfts}{https://github.com/oazapfts/oazapfts}). C'est le projet qui, sur un CV, se
distingue des dizaines de \og ~to-do app~ \fg{}.
\end{projet}

\begin{ressources}
\begin{itemize}
  \item RFC 9457 (Problem Details) --- \url{https://datatracker.ietf.org/doc/html/rfc9457}
  \item OpenAPI~3.1 --- \url{https://spec.openapis.org/oas/latest.html} --- Scalar \url{https://scalar.com/}
  \item \link{Firebase~PHP~JWT}{https://github.com/firebase/php-jwt} --- \link{LCobucci~JWT}{https://github.com/lcobucci/jwt}
  \item API design guide (Microsoft) --- \url{https://github.com/microsoft/api-guidelines}
  \item \link{Zircote Swagger-PHP (attributs -> OpenAPI)}{https://github.com/zircote/swagger-php}
\end{itemize}
\end{ressources}

\index{API REST}\index{JWT}\index{OpenAPI}\index{CORS}
% ============================================================================
\chapter{Frameworks PHP~: Laravel, Symfony, et les autres}
\label{ch:php-frameworks}
% ============================================================================

\begin{objectif}
\begin{itemize}
  \item Situer les cadres PHP en 2026, et choisir sans dogme.
  \item Reconna\^itre leurs concepts communs~: routage, DI, ORM, migrations, vues, files, tests.
  \item \'Ecrire, dans chacun, le CRUD qu'on vient de faire \og ~\`a la main~ \fg{}, et mesurer le gain.
  \item Comprendre les templates et l'\'ecosyst\`eme (bundles, packages, Forge, Octane...).
  \item Savoir quand s'en passer, et comment migrer une application legacy vers l'un d'eux.
\end{itemize}
\end{objectif}

\begin{note}
Ce chapitre n'est pas un cours de Laravel ni de Symfony~: ces cadres ont chacun une
documentation vaste et un livre entier leur serait consacr\'e. Il te donne la \textbf{carte}~:
tu sauras o\`u chercher, et surtout, tout ce que tu as \'ecrit \og ~\`a la main~ \fg{} dans les chapitres
pr\'ec\'edents porte un nom dans ces cadres. C'est cet aller-retour~: comprendre le m\'ecanisme,
puis utiliser l'outil.
\end{note}

\section{Le paysage}

\begin{tabz}{lXll}
  \headrow \thd{Cadre} \& \thd{Ce qui le d\'efinit} \& \thd{Version} \& \thd{Lien} \\
  \midrule
  \link{Laravel}{https://laravel.com/} \& la productivit\'e~: Eloquent, Artisan, \'ecosyst\`eme \'enorme (Nova, Forge, Livewire, Inertia), tr\`es bonne documentation, le plus d'offres \og ~web agency~ \fg{}. \& 12 \& \url{https://laravel.com/docs} \\
  \link{Symfony}{https://symfony.com/} \& les composants r\'eutilisables + un cadre~: rigueur, DI explicite, long support LTS, tr\`es pr\'esent en France/Europe et dans le secteur public. \& 7.x (6.4 LTS) \& \url{https://symfony.com/doc/current/index.html} \\
  \link{Slim}{https://www.slimframework.com/} \& un routeur + PSR-7/15, rien d'autre~: id\'eal pour une API l\'eg\`ere. \& 4.x \& \url{https://www.slimframework.com/docs/} \\
  \link{Laminas / Mezzio}{https://mezzio.dev/} \& PSR pur, modulaire, tr\`es typ\'e entreprise. \& --- \& \url{https://docs.laminas.dev/} \\
  \link{CodeIgniter~4}{https://codeigniter.com/} \& l\'eger, simple, h\'eritage du mutualis\'e. \& 4.x \& \url{https://codeigniter.com/user\\\_guide/} \\
  \link{Yii~2}{https://www.yiiframework.com/} \& Gii (g\'en\'eration de CRUD), tr\`es rapide, populaire en Asie. \& 2.0 \& \url{https://www.yiiframework.com/doc/} \\
  \link{Phalcon}{https://phalcon.io/} \& extension C~: performances brutes \'elev\'ees. \& 5.x \& \url{https://docs.phalcon.io/} \\
  \link{CakePHP}{https://book.cakephp.org/} \& \og ~convention over configuration~ \fg{}, le doyen. \& 5.x \& \\
  \link{WordPress}{https://wordpress.org/develop/} \& pas un cadre MVC, mais 40 \% du web~: l'API Plugin/Theme vaut le d\'etour pro. \& 6.x \& \url{https://developer.wordpress.org/} \\
\end{tabz}

\section{Le dictionnaire commun~: le m\^eme objet, trois noms}

\begin{tabz}{lXXX}
  \headrow \thd{Concept} \& \thd{En PHP nu (tes chapitres)} \& \thd{Laravel} \& \thd{Symfony} \\
  \midrule
  Point d'entr\'ee unique \& \verb|public/index.php| \& \verb|public/index.php| + \verb|bootstrap/app.php| \& \verb|public/index.php| + \verb|Kernel| \\
  Routage \& ta table \verb|\$routes| \& \verb|routes/web.php|, \verb|Route::get()| \& attributs \verb|\#[Route]| sur un contr\^oleur \\
  Contr\^oleur \& ta fonction \& \verb|app/Http/Controllers| \& \verb|src/Controller| \\
  Injection de d\'ependances \& ton conteneur maison \& le conteneur/service container auto-wired \& \verb|services.yaml| + autoconfigure \\
  ORM \& ton d\'ep\^ot PDO \& Eloquent (ActiveRecord) \& Doctrine ORM (Data Mapper) \\
  Migrations \& Phinx \& \verb|php artisan make:migration| \& \verb|bin/console doctrine:migrations:migrate| \\
  Validation \& ton \verb|Valideur| \& Form Requests + \verb|validate()| \& symfony/validator + attributs \\
  Vues \& \verb|require| + \'echappement \& Blade \& Twig \\
  Journal \& \verb|error\_log| \& \verb|Log::channel()| \& Monolog (c'est un composant Symfony) \\
  Files d'attente \& Redis + cron \& \verb|dispatch()| + Horizon \& Messenger + transport \\
  T\^aches planifi\'ees \& crontab \& scheduler \& Messenger scheduled tasks \\
  Tests \& PHPUnit \& \verb|php artisan test| (Pest par d\'efaut) \& PHPUnit + \verb|WebTestCase|/Panther \\
\end{tabz}

\subsection{Le m\^eme CRUD dans les trois styles}

\begin{codephp}
<?php
// ---------- LARAVEL : routes/api.php
use App\Http\Controllers\ArticleController;
Route::apiResource('articles', ArticleController::class)->parameters(['articles' => 'slug']);

// ---------- LARAVEL : ArticleController.php
namespace App\Http\Controllers;

use App\Http\Requests\StoreArticle;
use App\Models\Article;

class ArticleController extends Controller
{
    public function index(): \Illuminate\Http\Resources\Json\AnonymousResourceCollection
    {
        return Article::with('auteur:id,nom,email')          // evite le N+1
            ->where('statut', 'publie')
            ->latest('publie_le')
            ->paginate(20)
            ->through(fn (Article $a) => ArticleResource::make($a));
    }

    public function store(StoreArticle $request): \Symfony\Component\HttpFoundation\Response
    {
        $article = Article::create($request->validated() + ['auteur_id' => $request->user()->id]);
        return response()->json(ArticleResource::make($article), 201)
            ->header('Location', route('articles.show', $article));
    }
}

// ---------- LARAVEL : StoreArticle.php (la validation devient un objet reutilisable)
class StoreArticle extends \Illuminate\Foundation\Http\FormRequest
{
    public function rules(): array
    {
        return [
            'titre' => ['required', 'string', 'min:3', 'max:160'],
            'corps' => ['required', 'string', 'min:20', 'max:100000'],
            'tags' => ['array', 'max:12'],
            'tags.*' => ['string', 'max:30', 'distinct'],
        ];
    }
}

// ---------- LARAVEL : Models/Article.php (ActiveRecord : le modele porte la persistance)
class Article extends \Illuminate\Database\Eloquent\Model
{
    protected $fillable = ['titre', 'corps', 'slug', 'statut', 'publie_le', 'auteur_id'];
    protected $casts = ['publie_le' => 'datetime', 'tags' => 'array'];

    public function auteur(): \Illuminate\Database\Eloquent\Relations\BelongsTo
    { return $this->belongsTo(Auteur::class); }
    public function commentaires(): \Illuminate\Database\Eloquent\Relations\HasMany
    { return $this->hasMany(Commentaire::class)->where('approuve', true); }

    public function scopePublies($q) { return $q->where('statut', 'publie'); }
}
\end{codephp}

\begin{codephp}
<?php
// ---------- SYMFONY : src/Controller/ArticleController.php
namespace App\Controller;

use App\Entity\Article;
use App\Repository\ArticleRepository;
use App\Form\ArticleType;
use Doctrine\ORM\EntityManagerInterface;
use Symfony\Bundle\FrameworkBundle\Controller\AbstractController;
use Symfony\Component\HttpFoundation\{Request, Response, JsonResponse};
use Symfony\Component\Routing\Attribute\{Route, Method};
use Symfony\Component\Validator\Constraints as Assert;

final class ArticleController extends AbstractController
{
    public function __construct(private readonly ArticleRepository $articles) {}

    #[Route('/api/articles', name: 'api_articles', methods: ['GET'])]
    public function index(Request $request): JsonResponse
    {
        $page = max(1, (int) $request->query->getInt('page', 1));
        $items = $this->articles->parPage($page, min(100, max(1, $request->query->getInt('limite', 20))));
        return $this->json([
            'items' => $items,
            'page' => $page,
            'total' => $this->articles->compterPublies(),
        ], headers: ['cache-control' => 'public, max-age=30']);
    }

    #[Route('/articles/new', name: 'article_new', methods: ['GET', 'POST'])]
    public function new(Request $request, EntityManagerInterface $em): Response
    {
        $article = new Article();
        $form = $this->createForm(ArticleType::class, $article);
        $form->handleRequest($request);
        if ($form->isSubmitted() && $form->isValid()) {
            $article->publier(new \DateTimeImmutable());
            $em->persist($article);
            $em->flush();                       // une seule ecriture, la transaction est implicite
            $this->addFlash('succes', 'Article cree.');
            return $this->redirectToRoute('article_voir', ['slug' => $article->getSlug()]);
        }
        return $this->render('article/new.html.twig', ['form' => $form]);
    }
}

// ---------- SYMFONY : src/Entity/Article.php (Data Mapper : l'entite est pure)
#[\Doctrine\ORM\Mapping\Entity]
#[\Doctrine\ORM\Mapping\Table(name: 'articles')]
#[\Doctrine\ORM\Mapping\HasLifecycleCallbacks]
class Article
{
    #[\Doctrine\ORM\Mapping\Id, \Doctrine\ORM\Mapping\GeneratedValue, \Doctrine\ORM\Mapping\Column]
    private ?int $id = null;

    #[Assert\Length(min: 3, max: 160)]
    #[\Doctrine\ORM\Mapping\Column(length: 191, unique: true)]
    private string $slug = '';

    /** @phpstan-ignore-next-line */
    #[\Doctrine\ORM\Mapping\ManyToOne(targetEntity: Utilisateur::class)]
    #[\Doctrine\ORM\Mapping\JoinColumn(nullable: false)]
    private ?Utilisateur $auteur = null;

    public function publier(\DateTimeImmutable $quand): void
    {
        $this->publieLe = $quand;
        $this->statut = Statut::Publie;
    }
    /* getters/setters : php bin/console make:entity les genere */
}
\end{codephp}

\begin{codephp}
<?php
// ---------- Slim 4 : le "presque rien" pour une API
use Psr\Http\Message\ServerRequestInterface as Requete;
use Psr\Http\Message\ResponseInterface as Reponse;
use Slim\Factory\AppFactory;

$app = AppFactory::create();
$app->addBodyParsingMiddleware();
$app->addRoutingMiddleware();
$app->addErrorMiddleware(true, true, true);          // affiche les details en dev

$app->get('/api/v1/articles', function (Requete $r, Reponse $res, array $args) use ($depot) {
    $q = $r->getQueryParams();
    $donnees = $depot->page((int) ($q['page'] ?? 1), min(100, (int) ($q['limite'] ?? 20)));
    $res->getBody()->write(json_encode($donnees, JSON_THROW_ON_ERROR | JSON_UNESCAPED_UNICODE));
    return $res->withHeader('content-type', 'application/json');
});
$app->run();
\end{codephp}

\section{Les templates~: Blade et Twig}

\begin{codeplain}
{{-- Blade (Laravel) --}}
@extends('layouts.app')
@section('title', $article->titre)

@section('content')
    <article>
        <h1>{{ $article->titre }}</h1>          {{-- echappe par defaut ; {!! !!} = brut (dangereux) --}}
        <p>{{ Str::limit($article->extrait, 200) }}</p>
        @can('update', $article)
            <a href="{{ route('articles.edit', $article) }}">Modifier</a>
        @endcan
        @foreach ($article->tags as $tag)
            <x-ui.badge>{{ $tag->nom }}</x-ui.badge>       {{-- un composant Blade --}}
        @endforeach
        @empty($article->commentaires)
            <p>Aucun commentaire.</p>
        @endempty
    </article>
@endsection
\end{codeplain}

\begin{codeplain}
{# Twig (Symfony) #}
{% extends 'base.html.twig' %}
{% block title %}{{ article.titre }}{% endblock %}
{% block content %}
  <article>
    <h1>{{ article.titre }}</h1>                          {# echappe par defaut #}
    <p>{{ article.extrait|u.truncate(200, '...') }}</p>
    {% is_granted('update', article) %}
      <a href="{{ path('article_edit', { slug: article.slug }) }}">Modifier</a>
    {% endis_granted %}
    {% for tag in article.tags %}
      {% include 'components/_badge.html.twig' with { label: tag.nom } %}
    {% else %}
      <p>Aucun tag.</p>
    {% endfor %}
    {{ form_start(form) }}{{ form_end(form) }}
  </article>
{% endblock %}
\end{codeplain}

\section{Les outils qui accompagnent chaque cadre}

\begin{tabz}{XX}
  \headrow \thd{Laravel} \& \thd{Symfony} \\
  \midrule
  \verb|artisan| (g\'en\'erateurs, migrations, tinker) \& \verb|bin/console| (+ debug:router, debug:container) \\
  \link{Eloquent}{https://laravel.com/docs/eloquent} \& \link{Doctrine}{https://www.doctrine-project.org/projects/orm.html} \\
  \link{Livewire}{https://livewire.laravel.com/} / \link{Inertia}{https://inertiajs.com/} (interactive sans SPA) \& \link{Stimulus + Turbo (Hotwire)}{https://symfony.com/doc/current/frontend/htmx.html} \\
  \link{Forge}{https://forge.laravel.com/} / \link{Envoyer}{https://envoyer.io/} (d\'eploiement) \& \link{Symfony~Cloud}{https://symfony.com/cloud/} (PaaS officiel) \\
  \link{Octane}{https://laravel.com/docs/octane} (FrankenPHP/Swoole, ~5x d\'ebit) \& \link{FrankenPHP}{https://frankenphp.dev/} (app worker, idem, et c'est un serveur web \`a lui seul) \\
  Nova (back-office), Horizon (files), Reverb (WebSocket) \& Blackfire, Messenger, UX, Flex, MakerBundle \\
  \link{Name~of~Zen~docs}{https://laravel-news.com/} \& \url{https://symfony.com/blog/} + la newsletter \\
\end{tabz}

\begin{astuce}
\textbf{FrankenPHP} m\'erite ton attention~: un serveur applicatif PHP (fork de Caddy) qui garde le
processus vivant~: d\'emarrage d'app quasi instantan\'e, mont\'ee en d\'ebit, mont\'ees de charge \og ~\`a la
Node~ \fg{}, et int\'egration avec Mercure/SSE. C'est la nouveaut\'e d'infrastructure PHP la plus
int\'eressante depuis longtemps, et elle marche avec Laravel (Octane), Symfony et Slim.
\end{astuce}

\section{Choisir --- et migrer}

\begin{enumerate}
  \item \textbf{Tu veux livrer vite, seul, avec peu de lignes~?} Laravel. La communaut\'e est tr\`es
        active, presque tout existe d\'ej\`a en paquet.
  \item \textbf{Tu vises une entreprise, un march\'e fran\c{c}ais/europ\'een, un projet \`a longue dur\'ee de
        vie, ou tu aimes le typage fort~?} Symfony. Ses composants sont utilisables sans le cadre.
  \item \textbf{API minuscule, ou tu veux garder la main~?} Slim, ou ton noyau du chapitre~7.
  \item \textbf{Site \'editorial, clients non techniques~?} WordPress (avec Roots Bedrock + Acorn si tu veux du moderne). C'est un m\'etier \`a part~; ne le m\'eprise pas, mais ne l'appelle pas un cadre.
  \item \textbf{Portfolio et apprentissage~?} Fais le m\^eme projet en PHP nu, puis en Laravel~: tu
        sauras expliquer \emph{pourquoi} le cadre est utile. C'est exactement ce qui distingue un
        d\'eveloppeur d'un utilisateur de tutoriel.
\end{enumerate}

\etape{Migration d'un projet legacy~: la strat\'egie r\'ealiste}
\begin{enumerate}
  \item Ne r\'e\'ecris rien d'abord~: ajoute \verb|declare(strict\_types=1)|, Composer, PSR-4,
        PHPStan, et les tests sur le comportement actuel (tests \og ~de caract\'erisation~ \fg{}).
  \item Extrais le routeur et la liaison vers les contr\^oleurs (une table, une fonction par page).
  \item Introduis un conteneur de service~; remplace les \verb|new| globaux et les
        \verb|include| par de l'injection~; enrobe l'acc\`es \`a la base dans un d\'ep\^ot.
  \item Remplace les \verb|echo| \'epars par un moteur de template (Twig~: il s'installe seul).
  \item Garde l'application \og ~enrob\'ee~ \fg{} dans un cadre (Symfony/Laravel peuvent servir un script
        legacy via une route)~: les pages migrent une par une, pendant des mois, sans gros plan.
\end{enumerate}

\begin{exercice}
\begin{enumerate}
  \item Installe Laravel (\cmd{composer create-project laravel/laravel blog}) et Symfony
        (\cmd{composer create-project symfony/skeleton:"7.2.\*" blog}), fais tourner les deux,
        et trouve la route d'accueil de chacun dans le code~: tu dois pouvoir dire o\`u le
        fichier\ est et pourquoi.
  \item Dans Laravel~: mod\`ele \verb|Article| + migration + factory + seeder, contr\^oleur resource,
        validation par FormRequest, ressource JSON, paginate, autorisation avec Policies~;
        10 tests (feature tests avec SQLite en m\'emoire).
  \item Dans Symfony~: m\^eme chose avec MakerBundle (\verb|make:entity|, \verb|make:controller|),
        Doctrine fixtures, formulaire, validation par attributs, un test fonctionnel avec
        \verb|WebTestCase|.
  \item Compare la production~: nombre de fichiers, lignes \'ecrites, temps pass\'e, courbe
        d'apprentissage~; r\'edige une d\'ecision argument\'ee dans \verb|ANALYSE.md| (c'est un livrable
        que les recruteurs lisent).
  \item Ajoute dans le projet Laravel une file d'attente~: \`a la publication, envoyer un email en
        arri\`ere-plan avec \verb|ShouldQueue|, un job retry et un \'echec visible dans Horizon.
  \item D\'eploie le projet Symfony ou Laravel sur un VPS avec Deployer ou Ploi~: z\'ero
        \verb|git pull| manuel, migrations automatiques, OPcache invalid\'e, HTTPS, et rollback.
\end{enumerate}
\end{exercice}

\begin{ressources}
\begin{itemize}
  \item Laravel --- docs, Laracasts, News --- \url{https://laravel.com/docs/12.x}, \url{https://laracasts.com/}, \url{https://laravel-news.com/}
  \item Symfony --- les guides (excellents, en fran\c{c}ais partiel) --- \url{https://symfony.com/doc/current/page\\\_blocks.html} \\
        et la formation gratuite Open Classrooms --- \url{https://openclassrooms.com/fr/search?query=symfony}
  \item \link{Symfony~Casts}{https://symfonycasts.com/} (les meilleures vid\'eos de l'\'ecosyst\`eme, beaucoup sont gratuites)
  \item \link{Laravel~Daily~Learn}{https://laraveldaily.com/learn}, \link{Relay}{https://laravelrelay.app/}
  \item Comparatifs honn\^etes et r\'ecents --- \url{https://packagist.org/stats/} (parts d'usage des paquets)
\end{itemize}
\end{ressources}

\index{Laravel}\index{Symfony}\index{Blade}\index{Twig}\index{migration legacy}

\part{Java : robuste, partout, pour de vrai}
% ============================================================================
\chapter{Java~: le langage, la JVM, les bases}
\label{ch:java-bases}
% ============================================================================

\begin{objectif}
\begin{itemize}
  \item Installer un JDK, comprendre ce que sont JVM, bytecode, GC, et le mod\`ele \og ~\'ecrire une fois,
        ex\'ecuter partout~ \fg{}.
  \item Compiler et ex\'ecuter un programme, avec et sans outil de build.
  \item Ma\^itriser le typage statique, les primitifs, les conversions et la nullit\'e.
  \item \'Ecrire du Java moderne~: \verb|var|, enregistrements, expressions \verb|switch|, blocs de texte.
  \item Savoir lire une stack trace et un message du compilateur sans paniquer.
\end{itemize}
\end{objectif}

\section{La plateforme, en cinq mots}

\begin{description}
  \item[JDK] l'outil de d\'eveloppement~: compilateur (\verb|javac|), lanceur (\verb|java|),
        \verb|jshell|, \verb|javadoc|, \verb|jar|, analyeurs (\verb|jcmd|, \verb|jstack|, \verb|jmap|).
  \item[JVM] la machine virtuelle qui ex\'ecute le \emph{bytecode} (\fichier{.class}), just-\`a-temps
        (JIT), avec ramasse-miettes, et un \'ecosyst\`eme d'outils de profilage.
  \item[Version] une nouvelle version \textbf{tous les six mois}~; des versions \textbf{LTS}
        (17, 21, \emph{25}) soutenues des ann\'ees. En entreprise~: 17 ou 21 dominent.
        En 2026, \textbf{Java~21} est le choix par d\'efaut (virtual threads, records, patterns,
        s\'equenc\'e), \textbf{Java~25} apporte les versions de preview stabilis\'ees (structured
        concurrency, primitive types en preview, etc.).
  \item[Distributions] les binaires OpenJDK~: \link{Temurin (Adoptium)}{https://adoptium.net/fr/}
        (le choix par d\'efaut), \link{Amazon~Corretto}{https://aws.amazon.com/corretto/},
        \link{Azul~Zulu}{https://www.azul.com/downloads/}, \link{BellSoft~Liberica}{https://bell-sw.com/pages/downloads/},
        \link{GraalVM}{https://www.graalvm.org/} (compilation native, polyglotte).
  \item[Sp\'ecification] \link{JLS}{https://docs.oracle.com/javase/specs/} et les
        \link{JEP}{https://openjdk.org/jeps/0}--- chaque nouveaut\'es y est \'ecrite noir sur blanc.
\end{description}

\begin{terminal}
# Linux (Debian/Ubuntu) : le depot Temurin
sudo apt install -y temurin-21-jdk
# macOS
brew install --cask temurin@21
# Windows : https://adoptium.net/temurin/releases/?os=windows (MSI : cocher JAVA_HOME + PATH)
# Ou, multi-versions (recommande) :
brew install sdkman-cli && sdk install java 21.0.5-tem && sdk list java     # Linux/macOS
# Windows : scoop bucket add java && scoop install temurin21-jdk
# Verification :
java -version && javac -version && echo $JAVA_HOME
which java && readlink -f $(which java)     # ne pas etre surpris : /usr/bin/java est un lien alternatives
\end{terminal}

\subsection{Compiler, ex\'ecuter, et le \og ~single-file source launcher~ \fg{}}

\begin{codejava}
// HelloWorld.java --- le nom du fichier doit correspondre a la classe publique
public class HelloWorld {
    public static void main(String[] args) {
        System.out.println("Bonjour, Java !");
    }
}
\end{codejava}

\begin{terminal}
# Facon classique : compiler puis lancer
javac HelloWorld.java            # produit HelloWorld.class (le bytecode)
java HelloWorld                  # la JVM execute ; pas de ".class"

# Depuis Java 11 : lancer le fichier source directement (scripts, apprentissage)
java HelloWorld.java
# Avec arguments, proprietes, et classpath :
java -Xmx256m -Dapp.env=dev -cp lib/\*:target/classes com.exemple.App --port 8080
# JAR executable :
jar cfe app.jar com.exemple.App -C target/classes . && java -jar app.jar
# Decouvrir la JVM :
java --list-modules ; java -XX:+PrintFlagsFinal -version | grep -i heapsize ; jcmd VM.info
\end{terminal}

\begin{note}
Java est compil\'e \emph{et} ex\'ecut\'e par une JVM qui compile \`a chaud (C1 puis C2 JIT)~: le
d\'emarrage d'une application Java un peu grosse est lent, puis elle devient plus rapide que du JS
ou du PHP sur un calcul long. C'est pour cela qu'on utilise des cadres avec \og ~compilation
native~ \fg{} (GraalVM) ou des \emph{app/crac} pour le serverless~; et que dans ce livre, les mesures
de performance Java se font apr\`es \'echauffement (\link{JMH}{https://github.com/openjdk/jmh}).
\end{note}

\section{Types primitifs, variables et conversions}

\begin{codejava}
// Primitives : 8, pas d'objet, pas de null, empilees dans la pile d'appel.
boolean vif = true;
byte octet = 127;                    // -128..127 (rare, utile pour les protocols)
short petit = 32_000;                // -32 768..32 767
int n = 42;                          // -2^31..2^31-1 : c'est LE type entier par defaut
long grand = 9_000_000_000L;         // 64 bits : pour les id, les tailles, les timestamps
float virgule = 1.5f;                // 32 bits : ne l'utilise presque jamais
double d = 3.14159;                  // 64 bits IEEE-754 : le par defaut pour le flottant
char c = 'A';                        // 16 bits UTF-16 ; 0..65535 (les emoji debordent)

// 0b1010 = 10 ; 0x1F = 31 ; 1_000_000 s'ecrit avec des underscores ; 0777 est OCTAL (piege).
var inference = 42;                  // Java 10+ : le type est DEDUIT (ici int), pas " dynamique"
final int CONSTANTE = 7;             // non reassignable ; pour une veritable constante :
public static final double TVA = 0.20;

// Conversions
long l = n;                          // elargissement : implicite, sur
int revenu = (int) grand;             // retrecissement : EXPLICITE obligatoire (sinon erreur)
int arrondi = (int) Math.round(d);
long precise = Math.roundExact(2.5); // version 19+ : leve au lieu de silencieusement tronquer

// Division entiere : le piege n1 de Java
double demi = 1 / 2;                 // 0.0  <-- la division se fait en int AVANT conversion
double vraiDemi = 1.0 / 2;           // 0.5
double pourcentage = 5.0 * 2 / 100;  // 0.1 (et non 0)
int reste = 7 % -2;                   // 1 ; signe = celui du dividende

double infini = 1.0 / 0;              // Infinity : pas d'exception en flottant !
double nan = 0.0 / 0;                 // NaN
System.out.println(Double.isNaN(nan));          // true
System.out.println(nan == Double.NaN);          // false : NaN != NaN, toujours
System.out.println(0.1 + 0.2 == 0.3);           // false
System.out.println(Math.abs(0.1 + 0.2 - 0.3) < 1e-9);  // true
// En monetaire : JAMAIS double. Utilise BigDecimal (exact) ou des centimes en long.
var montant = java.math.BigDecimal.valueOf(12.50).multiply(java.math.BigDecimal.valueOf(3));
var arrondiUp = montant.setScale(2, java.math.RoundingMode.HALF_UP);   // 37.50

// Strings
String s = "Bonjour";
String bloc = """
    SELECT id, titre
    FROM   articles
    WHERE  statut = 'publie'
    """;                               // Java 15+ : text block, sans concatenations moches
var concat = "a" + 1 + true;           // "a1true" : evaluation gauche a droite
var interpo = "Resultat : " + arrondiUp.stripTrailingZeros().toPlainString();
System.out.printf("%-10s | %6.2f | %05d%n", "libelle", 12.34, 7);  // n = saut de ligne
String[] mots = bloc.lines().toList().toArray(String[]::new);
\end{codejava}

\subsection{Nullit\'e, types enveloppes, et \texttt{Optional}}

\begin{codejava}
Integer boxee = null;                  // les wrappers (Integer, Long, Double...) peuvent etre null
int casse = boxee;                      // NullPointerException a l'unboxing : le classique
long somme = liste.stream().mapToLong(Moteur::getDuree).sum();  // si un getDuree est Long null -> NPE

// Comparaison : equals/toujours pour les objets, == seulement pour les primitives
Integer a = 128, b = 128;
System.out.println(a == b);             // FALSE  (cache des Integer uniquement -128..127 !)
System.out.println(Objects.equals(a, b));   // true : LE reflexe a prendre
System.out.println("x".equals(s));         // true ; "x" == s : comparaison de reference -> faux

// Optional : un conteneur de valeur absente, pas un champ de classe
public Optional<Article> trouver(long id) {
    return Optional.ofNullable(repo.sel(id));
}
Optional<Article> o = trouver(1);
o.ifPresent(a -> System.out.println(a.titre()));
String titre = o.map(Article::titre).orElse("sans titre");
Article article = o.orElseThrow(() -> new Introuvable("id 1"));
o.filter(a -> a.publie()).ifPresent(this::notifier);
// Anti-patterns : o.get() sans isPresent(), Optional comme parametre ou champ, .orElse(calculCouteux())
// (le calcul est evalue dans tous les cas : orElseGet le rend paresseux)
\end{codejava}

\section{Op\'erateurs, conditions et boucles}

\begin{codejava}
// Arithmetique et incrementation : comme en JS/C, aux differences notees
int i = 0; i++; ++i; i += 2; i <<= 3; i >>>= 1;         // >>> : decale logique (sans signe)
boolean b = !b;                                         // pas de surcharge operateur en Java
long hachage = "cle".hashCode();                        // pas de `+` sur des types arbitraires

// Comparaison a trois voies : le spaceship de Java
Integer.compare(x, y);                                  // -1/0/1 (utilise dans les comparators)
Long.compare(a, b);
"b".compareTo("a");
Comparator.comparing(Article::date).thenComparing(Article::titre).reversed();

// Conditions
if (a > 0 && b < 10) { } else if (a == 0) { } else { }
String libelle = switch (note) {                        // Java 14+ : switch EXPRESSION
    case 20 -> "parfait";
    case 16, 17, 18, 19 -> "tres bien";
    case 10, 11, 12, 13, 14, 15 -> "bien";
    default -> {
        String r = note < 0 ? "impossible" : "a revoir";
        yield r;                                        // yield = "return" dans un switch
    }
};
// switch sur des types et des motifs (Java 21) : le pattern matching
String decrire(Object o) {
    return switch (o) {
        case Integer n when n > 100 -> "grand entier " + n;
        case Integer n -> "petit entier " + n;
        case String s when s.isBlank() -> "vide";
        case String s -> "texte de " + s.length() + " caracteres";
        case int[] tab -> "tableau de " + tab.length + " entiers";
        case null -> "rien";                             // Java 21 : le cas null est autorise
        default -> "autre chose";
    };
}

// Boucles
for (int j = 0; j < 10; j++) { }
for (var mot : "un deux".split(" ")) { }
for (int k = 0, m = 10; k < m; k++, m--) { }
while (it.hasNext()) { }
do { travail(); } while (tropTot());
// Etiquettes : a connaitre pour lire du code, a eviter d'ecrire
externe: for (int x = 0; x < 10; x++) { for (int y = 0; y < 10; y++) { if (y > x) continue externe; } }
\end{codejava}

\section{Tableaux~: le type brut dont tout le monde h\'erite dans sa t\^ete}

\begin{codejava}
int[] notes = new int[5];                 // taille fixe, initialisee a 0
int[] litteral = {14, 9, 17, 12};
String[] villes = { "Dakar", "Thies", "Saint-Louis" };
notes.length;                              // une PROPRIETE, pas une methode
villes[0] = "Thi\`es";
//ville[5];                                 // ArrayIndexOutOfBoundsException (jamais de lecture floue)

var matrice = new int[3][4];               // tableaux imbriques (lignes de tailles differentes OK)
matrice[0][1] = 7;

Arrays.sort(litteral);
Arrays.copyOf(litteral, 10);
Arrays.equals(litteral, copie);
Arrays.fill(notes, -1);
Arrays.mismatch(a, b);                      // Java 11+ : index de difference, -1 si egales
Arrays.asList("a", "b");                    // vue FIXE sur le tableau -- pas un ArrayList !
List.of(1, 2, 3);                            // liste immuable (Java 9+)
new ArrayList<>(List.of(1, 2, 3));           // la version modifiable

// Impression (jamais System.out.println(tab) : tu obtiens [I@1a2b3c)
System.out.println(Arrays.toString(litteral));
System.out.println(Arrays.deepToString(matrice));
System.out.printf("notes = %s%n", Arrays.toString(notes));

// Copie rapide et allocation : les deux utilitaires bas niveau
System.arraycopy(src, 0, dest, 10, n);     // ou : Arrays.copyOfRange(src, debut, fin)
var montee = new java.util.RandomGenerator.JdkRandom(1).ints(1000, 0, 50).toArray();
// Tableaux et genericite : impossible de creer un List<String>[] typable -> on utilise des Collections.
\end{codejava}

\section{M\'ethodes, arguments, et la surcharge}

\begin{codejava}
public final class Calculs {
    private Calculs() { }                        // classe utilitaire : pas d'instance

    // Surcharge : meme nom, signatures differentes (les types, pas les noms de parametres)
    public static int max(int a, int b) { return Math.max(a, b); }
    public static double max(double a, double b) { return Math.max(a, b); }
    public static int max(int... valeurs) {       // varargs : un tableau en entree
        if (valeurs.length == 0) { throw new IllegalArgumentException("au moins une valeur"); }
        int m = valeurs[0];
        for (var v : valeurs) { if (v > m) { m = v; } }
        return m;
    }
    public static void melanger(List<String> l) { java.util.Collections.shuffle(l); }
}

// Pas de passage par reference : les objets passent par copie de la reference
static void tente(Article a) { a = new Article("autre"); }   // l'appelant ne le voit pas
static void marche(Article a) { a.marquerVu(); }              // mais muter l'objet, si

// En Java, une methode renvoie UN seul resultat : pour en renvoyer plusieurs -> un record
public record DivResultat(int quotient, int reste) { }
public static DivResultat diviser(int a, int b) { return new DivResultat(a / b, a % b); }

// Recursion : pas de TCO garantie par la JVM -> sur de tres grandes profondeurs, une boucle.
static long fib(int n, long a, long b) { return n == 0 ? a : fib(n - 1, b, a + b); }

// Memo : les methodes d'une classe ne sont jamais globales ; tout est statique ou membre.
\end{codejava}

\section{Lire les messages~: le compilateur et la pile d'appels}

\begin{codejava}
/*
  Exemples reellement frequents, avec leur traduction :

  error: cannot find symbol                -> classe non importee, faute de frappe, ou methode
                                             inexistante ; verifie la casse et le package
  error: variable x might not have been initialized -> champ local non initialise (en Java, un
                                             local n'a AUCUNE valeur par defaut)
  error: incompatible types: int cannot be converted to String -> String.valueOf(n) ou "" + n
  error: unreported exception IOException  -> les exceptions verifiees : try/catch ou throws
  error: method does not override or implement a method from a supertype
                                          -> signature legerement differente : mets @Override
  Exception in thread "main" java.lang.NullPointerException: Cannot invoke "String.length()"
       because "s" is null                 -> depuis Java 14, le message dit QUI est null. Lis-le.
  java.lang.ArrayIndexOutOfBoundsException: Index 5 out of bounds for length 5
  java.util.ConcurrentModificationException -> tu modifies une collection pendant son iteration
  java.lang.StackOverflowError              -> recursion sans cas de base (ou cycle d'objets)
*/
\end{codejava}

\begin{exercice}
\begin{enumerate}
  \item \'Ecris \verb|FizzBuzz|, \verb|estPremier|, \verb|pgcd| (Euclide, puis version binaire),
        \verb|converser(int[])| (sur place), \verb|moyenneMobile(double[], int)|~: tous en
        m\'ethodes statiques d'une classe utilitaire priv\'ee.
  \item Pr\'edis puis v\'erifie~: \verb|1/2|, \verb|1.0/2|, \verb|7\%3|, \verb|-7\%3|,
        \verb|(int) 3.99|, \verb|Math.round(-2.5)|, \verb|Double.doubleToLongBits(0.0)|,
        \verb|Integer.MAX\_VALUE + 1|, \verb|0.1 + 0.2 == 0.3|, \verb|"1" + 2|, \verb|2 + "1"|.
  \item Convertis un nombre en romains (\verb|1..3999|), puis l'inverse~; puis un nombre en
        base quelconque (2 \`a 36) et reviens, en utilisant \verb|Integer.toString(n, base)|
        et \verb|Integer.parseInt(s, base)|~: compare \`a ta version manuscrite.
  \item \'Ecris \verb|Textes|~: \verb|inverser|, \verb|estPalindrome| (en ignorant ponctuation et
        accents via \verb|java.text.Normalizer|), \verb|compterMots|, \verb|capitaliser|,
        \verb|estAnagramme|, \verb|slug|~: en utilisant \verb|StringBuilder| et les m\'ethodes
        \verb|chars()|/ \verb|codePoints()|.
  \item \'Ecris une classe \verb|Matrice| (tableau 2D) avec \verb|ajoute|, \verb|multiplier|,
        \verb|transposer|, \verb|toString| et \verb|equals|/ \verb|hashCode| corrects.
  \item Ouvre \verb|jshell|~: teste chaque snippet de ce chapitre une ligne \`a la fois.
        C'est le meilleur moyen de v\'erifier une hypoth\`ese sans cr\'eer de projet.
\end{enumerate}
\end{exercice}

\begin{ressources}
\begin{itemize}
  \item Dev.java (la porte officielle moderne) --- \url{https://dev.java/learn/}
  \item Documentations JDK (javadoc de r\'ef\'erence) --- \url{https://docs.oracle.com/en/java/javase/21/docs/api/index.html}
  \item \link{Baeldung}{https://www.baeldung.com/java-tutorial} (recettes concises, enorme index)
  \item \emph{Introduction to Java Programming} -- \url{https://java-programming.baoyj.com/} (libre)
  \item \link{JEP index}{https://hiphopjava.github.io/jep-in-Java-FR/} (JEP traduits en francais)
\end{itemize}
\end{ressources}

\index{Java!bases}\index{JDK}\index{JVM}\index{primitives}\index{Optional}
% ============================================================================
\chapter{Java orient\'e objet~: classes, encapsulation, h\'eritage, interfaces}
\label{ch:java-oop}
% ============================================================================

\begin{objectif}
\begin{itemize}
  \item \'Ecrire une classe correcte~: invariants, encapsulation, \verb|equals|, \verb|toString|.
  \item Utiliser les \verb|record| pour les donn\'ees et les classes immuables par d\'efaut.
  \item Choisir entre h\'eritage, interface et composition~: et savoir que \verb|sealed| change la donne.
  \item Mod\'eliser un domaine avec des ADT (types sommes) et un \verb|switch| exhaustif.
  \item Conna\^itre les classiques~: \verb|static|, classes imbriqu\'ees, initialisation, paquets.
\end{itemize}
\end{objectif}

\section{Anatomie d'une classe}

\begin{codejava}
package com.exemple.blog.modele;                       // = arborescence de dossiers, obligatoire en tete

import java.time.Instant;
import java.util.ArrayList;
import java.util.Collections;
import java.util.List;
import java.util.Objects;

/**
 * Un article de blog. Modele riche : les regles vivent ici.
 */
public final class Article {

    // Constantes : public static final, nom en MAJUSCULES_SNAKE
    public static final int TITRE_MIN = 3;
    public static final int TITRE_MAX = 160;

    // Champs : prives par defaut. static = appartient a la classe.
    private static final java.util.concurrent.atomic.AtomicLong SEQUENCE = new java.util.concurrent.atomic.AtomicLong();
    private final long id;
    private String titre;
    private String corps;
    private Statut statut;
    private final Instant creeLe;
    private Instant publieLe;
    private final List<String> tags;                    // liste mutable -> encapsulee

    // Constructeur : meme nom, aucun type de retour. Surcharges = mots-cles optionnels.
    public Article(String titre, String corps, Statut statut) {
        this(SEQUENCE.incrementAndGet(), titre, corps, statut, Instant.now(), null, List.of());
    }
    public Article(long id, String titre, String corps, Statut statut,
                   Instant creeLe, Instant publieLe, List<String> tags) {
        this.id = id;
        this.creeLe = Objects.requireNonNull(creeLe, "creeLe");
        this.tags = new ArrayList<>(tags);              // COPIE defensive : jamais la reference recue
        setTitre(titre);                                // passe par le setter : la regle s'applique une fois
        this.corps = Objects.requireNonNull(corps, "corps");
        this.statut = Objects.requireNonNull(statut, "statut");
        this.publieLe = publieLe;
        verifierInvariant();
    }

    // Bloc d'initialisation de classe (une fois, au chargement) et d'instance
    static { /* lecture d'un fichier de config, construction d'un regex compile, etc. */ }
    { /* normalisation commune a tous les constructeurs -- rarement necessaire */ }

    public long id() { return id; }
    public String titre() { return titre; }

    public void setTitre(String titre) {
        var t = titre == null ? "" : titre.trim();
        if (t.length() < TITRE_MIN || t.length() > TITRE_MAX) {
            throw new IllegalArgumentException(
                "titre de %d caracteres attendu entre %d et %d".formatted(t.length(), TITRE_MIN, TITRE_MAX));
        }
        this.titre = t;
    }

    /** Vue en lecture seule : la liste reelle reste privee. */
    public List<String> tags() { return Collections.unmodifiableList(tags); }

    public void ajouterTag(String tag) {
        var t = Objects.requireNonNull(tag, "tag").trim().toLowerCase(java.util.Locale.ROOT);
        if (t.isEmpty() || t.length() > 30) { throw new IllegalArgumentException("tag invalide : " + tag); }
        if (!tags.contains(t)) { tags.add(t); }
    }

    public boolean estPublie() { return statut == Statut.PUBLIE; }

    public Article publier(Instant quand) {
        if (statut == Statut.ARCHIVE) { throw new IllegalStateException("un article archive ne se republicie pas"); }
        return new Article(id, titre, corps, Statut.PUBLIE, creeLe, quand, tags);  // nouvel objet
    }

    private void verifierInvariant() {
        if (statut == Statut.PUBLIE && publieLe == null) {
            throw new IllegalStateException("article publie sans date de publication");
        }
        if (publieLe != null && publieLe.isBefore(creeLe)) {
            throw new IllegalStateException("publication anterieure a la creation");
        }
    }

    @Override public String toString() {
        return "Article #%d \"%s\" (%s, %d tags)".formatted(id, titre, statut, tags.size());
    }
    @Override public boolean equals(Object o) {
        if (this == o) { return true; }
        if (!(o instanceof Article a)) { return false; }   // pattern matching (Java 16+)
        return id == a.id;                                  // identite metier : l'id suffit ici
    }
    @Override public int hashCode() { return Long.hashCode(id); }
}

enum Statut {
    BROUILLON("Brouillon"), PUBLIE("Publie"), ARCHIVE("Archive");
    private final String libelle;
    Statut(String libelle) { this.libelle = libelle; }
    public String libelle() { return libelle; }
    public static Statut depuis(String v) {
        return valueOf(v.trim().toUpperCase(Locale.ROOT));      // IllegalArgumentException si inconnu
    }
}
\end{codejava}

\subsection{Modificateurs d'acc\`es~: la table honn\^ete}

\begin{tabz}{lXXXX}
  \headrow \thd{Modificateur} \& \thd{Classe} \& \thd{Sous-classe} \& \thd{Paquet} \& \thd{Monde} \\
  \midrule
  \verb|private| \& oui \& non \& non \& non \\
  (aucun, \og ~package-private~ \fg{}) \& oui \& non \& oui \& non \\
  \verb|protected| \& oui \& oui \& oui \& non \\
  \verb|public| \& oui \& oui \& oui \& oui \\
  \verb|module| (Java~9, \verb|exports|) \& \multicolumn{4}{X}{uniquement avec \verb|module-info.java| : contr\^ole entre modules} \\
\end{tabz}

Les trois r\'eflexes qui rendent une conception saine~: (1) \verb|final| sur les classes et les
m\'ethodes qui n'ont pas vocation \`a l'\^etre~; (2) \verb|private| sur les champs, \verb|public|
uniquement sur ce qui est une capacit\'e~; (3) aucun setter \verb|public| automatique~: un setter
g\'en\'er\'e par l'IDE est une d\'ecision d'architecture par d\'efaut, rarement la bonne.

\section{\texttt{equals}, \texttt{hashCode}, \texttt{toString}}

Le contrat, litt\'eralement~: deux objets \verb|equals| \textbf{doivent} avoir le m\^eme
\verb|hashCode| (l'inverse n'est pas requis). Sans cela, \verb|HashMap|/ \verb|HashSet|
ne retrouvent plus la cl\'e, et tu passes une journ\'ee \`a chercher pourquoi.

\begin{codejava}
// La version courte, moderne, et sans faute :
@Override public boolean equals(Object o) {
    return o instanceof Point p && Double.compare(p.x, x) == 0 && Double.compare(p.y, y) == 0;
}
@Override public int hashCode() { return Objects.hash(x, y); }      // ou Objects.hashCode(x) seul

// Les objets flottants : compare avec une tolerance, jamais ==
static boolean pres(Point a, Point b) {
    return Math.abs(a.x() - b.x()) < 1e-9 && Math.abs(a.y() - b.y()) < 1e-9;
}

// A ne JAMAIS mettre dans equals/hashCode : un champ mutable, un identifiant genere
// a l'instance (System.identityHashCode), un objet sans equals coherent (tableau -> Arrays.equals).
// Collections : Objects.equals pour les champs, Arrays.equals/deepEquals pour les tableaux.
\end{codejava}

\section{Les \texttt{record}~: la donn\'ee bien faite, en une ligne}

\begin{codejava}
public record Point(double x, double y) { }                      // immuable, avec equals/hashCode/toString
public record ArticleResume(long id, String titre, String slug, Instant publieLe) { }

// Constructeur compact : validation centrale, sans field assignations repetees
public record Email(String valeur) {
    public Email {
        valeur = valeur == null ? "" : valeur.trim().toLowerCase(Locale.ROOT);
        if (!valeur.matches("^[^@\\s]+@[^@\\s]+\\.[a-z]{2,}$")) {
            throw new IllegalArgumentException("email invalide : " + valeur);
        }
    }
    public boolean estGmail() { return valeur.endsWith("@gmail.com"); }
}

// Constructeurs canoniques et surcharges, methodes, impl\'ementation d'interfaces : OK
public record Periode(Instant debut, Instant fin) implements Comparable<Periode> {
    public Periode {
        if (fin.isBefore(debut)) { throw new IllegalArgumentException("fin < debut"); }
    }
    public Periode(Instant debut, java.time.Duration duree) { this(debut, debut.plus(duree)); }
    public long secondes() { return java.time.Duration.between(debut, fin).toSeconds(); }
    @Override public int compareTo(Periode o) { return debut.compareTo(o.debut); }
}

// Limite : pas d'h eritage possible (extends Record implicite), pas de champ non-final.
// -> un record = DONNEES. Des que tu as besoin de mutabilite ou d'heritage -> une classe.
var p = new Point(1, 2);
System.out.println(p.x() + " " + p);            // accesseur = nom du champ ; toString genere
var points = List.of(new Point(0, 0), new Point(1, 1));
var parX = points.stream().sorted(Comparator.comparingDouble(Point::x)).toList();
\end{codejava}

\begin{astuce}
Utilise un \verb|record| comme DTO, comme cl\'e de Map, comme r\'esultat multiple d'une m\'ethode,
comme \'ev\'enement, comme requ\^ete/r\'eponse d'API. En revanche, une entit\'e m\'etier avec des r\`egles et
des transitions (\verb|publier()|, \verb|payer()|) m\'erite une classe~: le record n'emp\^eche pas
les m\'ethodes, mais il encourage \`a tout rendre public et plat.
\end{astuce}

\section{H\'eritage, polymorphisme, \texttt{abstract}}

\begin{codejava}
public abstract class Forme implements Comparable<Forme> {
    private final String couleur;
    private int x, y;

    protected Forme(String couleur) { this.couleur = Objects.requireNonNull(couleur); }   // jamais appelee seule

    public abstract double aire();                       // la sous-classe doit fournir
    public abstract Forme deplace(int dx, int dy);

    public String decrire() {                          // implementation partagee
        return "%s (%s, aire %.2f)".formatted(getClass().getSimpleName(), couleur, aire());
    }
    @Override public int compareTo(Forme o) { return Double.compare(aire(), o.aire()); }
    // equals/hashCode sur une hierarchie : ne les definit que si tu peux le faire proprement.
    // Regle de bonne conduite : si tu ecris equals dans une classe heritable, documente-le.
}

public final class Cercle extends Forme {
    private final double rayon;
    public Cercle(double rayon, String couleur) { super(couleur); this.rayon = rayon; }
    @Override public double aire() { return Math.PI * rayon * rayon; }
    @Override public Cercle deplace(int dx, int dy) { return new Cercle(rayon, couleur()); }  // covariance du retour
}

// Une hierarchie n'est PAS un moyen de partager du code : c'est un contrat de substitution
// (principe de Liskov). Si une sous-classe casse le contrat, la hierarchie etait fausse.
Forme f = new Cercle(2, "rouge");
f.aire();                       // dispatch virtuel : le code de Cercle s'execute
if (f instanceof Cercle c) { System.out.println(c.rayon()); }   // jamais de cast sans test

// Pas d'heritage multiple de classes. Les interfaces ont des methodes par defaut :
interface Dessinable {
    void dessiner();
    default void dessinerEnRouge() { System.out.println("rouge"); dessiner(); }
    static Dessinable vide() { return () -> { }; }
    sealed interface FormeDessinable extends Dessinable permits Cercle, Rectangle { }
}
\end{codejava}

\section{Interfaces~: le vrai pilier de la conception en Java}

\begin{codejava}
// 1. Un contrat, pas un champ de donnees : pas de constantes "magiques" dans une interface.
public interface Depot<T, Id> extends AutoCloseable {
    T sauver(T entite);
    java.util.Optional<T> trouver(Id id);
    List<T> lister(Pageable page);
    default boolean existe(Id id) { return trouver(id).isPresent(); }
    // Une methode "default" est la pour evoluer sans casser les impls ; pas pour stocker du logique.
}

// 2. Impl\'ementation par composition : les capacites sont des objets qu'on injecte.
public final class DepotArticlesJdbc implements Depot<Article, Long> { /* ... */ }
public final class DepotArticlesCache implements Depot<Article, Long> {
    private final Depot<Article, Long> delegue;          // pattern decorator
    private final java.util.concurrent.ConcurrentHashMap<Long, Article> cache = new java.util.concurrent.ConcurrentHashMap<>();
    public DepotArticlesCache(Depot<Article, Long> delegue) { this.delegue = delegue; }
    @Override public Article sauver(Article a) { cache.put(a.id(), a); return delegue.sauver(a); }
    /* ... */
}

// 3. Les interfaces-fonctionnelles du JDK (une seule methode abstraite, @FunctionalInterface)
@FunctionalInterface
public interface Validation<T> {
    Resultat valider(T objet);
    default Validation<T> puis(Validation<T> autre) {
        return o -> { var r = valider(o); return r.valide() ? autre.valider(o) : r; };
    }
}
Runnable r = () -> System.out.println("go");
Comparator<Article> parTitre = Comparator.comparing(Article::titre, String.CASE_INSENSITIVE_ORDER);
java.util.function.Predicate<Article> publiees = a -> a.estPublie();
java.util.function.Function<Article, String> enJson = a -> a.titre();
java.util.function.Consumer<String> affiche = System.out::println;
// Les references de methode : Article::titre, System.out::println, List::size, Integer::parseInt,
// String::new (fabrique), Cercle::new -- voir chapitre Streams.
\end{codejava}

\section{\texttt{sealed}~: la somme de types, et le switch exhaustif}

C'est LE motif qui change la fa\c{c}on de mod\'eliser un domaine en Java~: un ensemble \emph{ferm\'e}
de variantes, et un compilateur qui te pr\'evient quand tu en oublies une.

\begin{codejava}
public sealed interface ErreurApi permits ErreurHttp, ErreurValidation, ErreurReseau { }

public record ErreurHttp(int statut, String message) implements ErreurApi { }
public record ErreurValidation(Map<String, List<String>> champs) implements ErreurApi {
    public ErreurValidation { champs = Map.copyOf(champs); }
}
public record ErreurReseau(String hote, Throwable cause) implements ErreurApi { }

// Traitement : pas de "default" necessaire -> le compilateur sait que c'est complet
public static String decrire(ErreurApi e) {
    return switch (e) {
        case ErreurHttp h -> "HTTP " + h.statut() + " : " + h.message();
        case ErreurValidation v -> v.champs().size() + " champ(s) en erreur";
        case ErreurReseau r -> "joins pas " + " + r.hote() + " (" + r.cause().getMessage() + ")";
        // aucune branche default : si j'ajoute une variante, TOUTES les occurrences cassent a la compilation
    };
}

// Hierarchie fermee partielle : "non-sealed" autorise les sous-types elsewhere
public sealed class Compte permits CompteCourant, CompteEpargne { }
public non-sealed class CompteCourant extends Compte { }

// ADT recursif + pattern matching : un mini-AST, 15 lignes, sans une seule instanceof
sealed interface Expr permits Const, Plus, Mult { }
record Const(double v) implements Expr { }
record Plus(Expr g, Expr d) implements Expr { }
record Mult(Expr g, Expr d) implements Expr { }
static double eval(Expr e) {
    return switch (e) {
        case Const c -> c.v();
        case Plus(Expr g, Expr d) -> eval(g) + eval(d);       // "record patterns" (Java 21)
        case Mult(Expr g, Expr d) -> eval(g) * eval(d);
    };
}
// eval(new Mult(new Plus(new Const(1), new Const(2)), new Const(4))) = 12
\end{codejava}

\section{Immutabilit\'e~: le d\'efaut qui rend tout plus simple}

\begin{codejava}
public final class Money implements Comparable<Money> {
    private final long centimes;
    private final java.util.Currency devise;

    public Money(long centimes, java.util.Currency devise) {
        this.centimes = centimes;
        this.devise = Objects.requireNonNull(devise);
    }
    public static Money centimes(long c, java.util.Currency d) { return new Money(c, d); }
    public static Money euros(double e, java.util.Currency d) { return new Money(Math.round(e * 100), d); }

    public Money plus(Money autre) {
        requireSameDevise(autre);
        return new Money(Math.addExact(centimes, autre.centimes), devise);   // addExact : OverflowException, pas de silence
    }
    public Money fois(int facteur) { return new Money(Math.multiplyExact(centimes, facteur), devise); }
    public Money partage(int parts) { /* repartir le centime restant : la vraie difficulte du monetaire */ }

    private void requireSameDevise(Money autre) {
        if (!devise.equals(autre.devise)) { throw new IllegalArgumentException("devises differentes"); }
    }
    @Override public int compareTo(Money o) { return Long.compare(centimes, o.centimes); }
    @Override public String toString() { return "%s %s".formatted(devise, centimes / 100.0); }
}
// Pourquoi c'est important : un objet immuable est sur et partageable entre threads sans
// aucune synchronisation, se met en cache sans peur, et ne peut pas etre corrompu
// par un appelant. C'est la forme par defaut d'un modele metier en Java.
\end{codejava}

\begin{note}
Le mot \textbf{final} ne rend pas l'objet profond\'ement immuable~: un champ \verb|final List<X>|
peut rester mutable. Les trois techniques de durcissement~: copie d\'efensive \`a l'entr\'ee et \`a la
sortie~; \verb|List.copyOf(...)| / \verb|Map.copyOf(...)| (vraiment immuables, Java~10+)~; ou
ne pas exposer la collection~: renvoyer un \verb|Stream|, une taille, ou un accesseur typ\'e
\verb|int tailleTags()|.
\end{note}

\begin{exercice}
\begin{enumerate}
  \item \'Ecris le mod\`ele complet d'une biblioth\`eque~: \verb|Livre| (ISBN, titre, auteurs, ann\'ee),
        \verb|Exemplaire| (code-barres, \'etat en enum), \verb|Emprunt| (dates, membre),
        \verb|Membre|~: tout immuable, invariants dans les constructeurs, pas un seul setter public.
  \item Ajoute \verb|equals|/ \verb|hashCode|/ \verb|toString| corrects, puis refais la m\^eme
        chose avec des records~: compare les deux fichiers (le record en fait la moiti\'e).
  \item \'Ecris une classe \verb|Duree| (secondes) avec \verb|plus|, \verb|moins|, \verb|compare|,
        \verb|formatage| (\og ~1 h 05~ \fg{}), et des tests sur les cas limites~; puis remplace-la par
        \verb|java.time.Duration| et note ce que tu as gagn\'e (ou pas).
  \item Hi\'erarchie \verb|Forme| : abstract \verb|aire|, \verb|perimetre|, \verb|deplacer|, et une
        interface \verb|Comparateur| avec \verb|byArea()| et \verb|byPerimeter()| par d\'efaut.
        Ajoute un test~: un tableau mixte tri\'e par aire.
  \item Refais la mod\'elisation des erreurs avec \verb|sealed| + \verb|record| + \verb|switch|
        exhaustif, puis ajoute une nouvelle variante et compte les lignes rouges que le compilateur
        produit~: c'est la valeur de ce patron, en chair et en os.
  \item Construis une petite \verb|Resultat<T>| (succ\`es/\'echec, sans exception)~:
        \verb|ok(v)|, \verb|err(e)|, \verb|map|, \verb|flatMap|, \verb|recover|, \verb|onSuccess|,
        \verb|getOrDefault|. C'est le d\'ebut de la programmation fonctionnelle en Java~; le
        chapitre suivant t'en montrera l'utilit\'e concr\`ete avec Streams.
\end{enumerate}
\end{exercice}

\begin{ressources}
\begin{itemize}
  \item Dev.java --- \emph{Object-Oriented Programming} --- \url{https://dev.java/learn/java-platform/object-oriented-programming/}
  \item Javadoc \verb|Object| (le contrat d'equals/hashCode y est \'ecrit noir sur blanc) --- \url{https://docs.oracle.com/en/java/javase/21/docs/api/java.base/java/lang/Object.html}
  \item Records --- \url{https://openjdk.org/jeps/395} --- Sealed classes --- \url{https://openjdk.org/jeps/409}
  \item Patterns for package-private classes --- \url{https://inside.java/} (le blog des \'equipes OpenJDK)
  \item \emph{Effective Java} (Joshua~Bloch) --- les items 1--15 valent le livre : c'est LA r\'ef\'erence de style objet Java.
\end{itemize}
\end{ressources}

\index{Java!classes}\index{record}\index{sealed}\index{equals}\index{immutabilite}
% ============================================================================
\chapter{G\'en\'eriques, collections et exceptions}
\label{ch:java-collections}
% ============================================================================

\begin{objectif}
\begin{itemize}
  \item Choisir la bonne structure~: \verb|ArrayList|, \verb|HashMap|, \verb|TreeMap|, \verb|Deque|,
        \verb|Set| ordonn\'e, files concurrentes.
  \item Comprendre les generics~: bornes, wildcards, PECS, effacement.
  \item \'Ecrire des exceptions qui racontent une histoire~: v\'erifi\'ees, non v\'erifi\'ees, causes,
        \verb|try-with-resources|.
  \item \'Eviter les dix erreurs de collection qui reviennent dans chaque revue de code.
\end{itemize}
\end{objectif}

\section{Collections~: l'arbre de d\'ecision}

\begin{tabz}{lllX}
  \headrow \thd{Besoin} \& \thd{Choisis} \& \thd{Complexit\'e cl\'e} \& \thd{Notes} \\
  \midrule
  liste index\'ee, append fr\'equent \& \verb|ArrayList| \& get O(1), add O(1) amorti, remove O(n) \& le d\'efaut. Cro\^it par recopie (\$\times\$1{,}5). \\
  ajouts/retraits fr\'equents aux bords \& \verb|ArrayDeque| \& O(1) \& \textbf{ne jamais utiliser \verb|LinkedList| comme file}~: plus lent en pratique (cache). \\
  liste li\'ee \& \verb|LinkedList| \& O(1) aux it\'erateurs \& rarement justifi\'e~; API \verb|Deque|+ \verb|List| \`a la fois = pi\`ege. \\
  unicit\'e \& \verb|HashSet| \& O(1) \& ordre non garanti~; d\'epend de \verb|hashCode|. \\
  unicit\'e tri\'ee \& \verb|TreeSet| \& O(\$\log n\$) \& exige \verb|Comparable| ou un \verb|Comparator|. \\
  unicit\'e avec ordre d'insertion \& \verb|LinkedHashSet| \& O(1) \& le \og ~cache LRU-lite~ \fg{}, d\'eduction des doublons. \\
  dictionnaire \& \verb|HashMap| \& O(1) \& cl\'es \verb|equals|/\verb|hashCode| coh\'erentes. \\
  dictionnaire avec insertion en fin \& \verb|LinkedHashMap| \& O(1) \& \verb|removeEldestEntry| = cache LRU en 5 lignes. \\
  dictionnaire tri\'e / plages \& \verb|TreeMap| \& O(\$\log n\$) \& \verb|subMap|, \verb|floorKey|, \verb|ceilingKey|. \\
  dictionnaire immuable \& \verb|Map.of(...)| / \verb|Map.copyOf(..)| \& O(1) \& limite de 10 paires pour \verb|of| ; \verb|ofEntries| au-del\`a. \\
  concurrent \& \verb|ConcurrentHashMap| \& O(1) sans verrou global \& \verb|computeIfAbsent|, \verb|merge|, \verb|reduce*. \\
  files bloquantes (threads) \& \verb|ArrayBlockingQueue| / \verb|LinkedBlockingQueue| \& \& born\'ee = dos-pression. Base de tout producteur/consommateur. \\
  tas de priorit\'e \& \verb|PriorityQueue| \& offre O(\$\log n\$) \& k-plus-grands, ordonnanceur, Dijkstra. \\
  cl\'es par identit\'e (pas equals) \& \verb|IdentityHashMap| \& \& \verb|WeakHashMap| pour le cache sans fuite. \\
\end{tabz}

\begin{codejava}
// Creation et initialisation : la forme moderne
var noms = new ArrayList<String>(List.of("Awa", "Nour", "Ibra"));
var notes = new HashMap<String, Integer>(Map.of("alg\`ebre", 14, "analyse", 17));
var immuable = List.copyOf(noms);
var setDeTag = new LinkedHashSet<>(List.of("php", "java", "php"));   // [php, java]

// Lecture : les methodes qui remplacent une boucle
notes.getOrDefault("geometrie", 0);
notes.computeIfAbsent("geometrie", k -> new java.util.ArrayList<>() /* valeur */);
notes.merge("analyse", 2, Integer::sum);                 // ajoute 2 a la cle (ou la cree)
notes.putIfAbsent("proto", 9);
notes.forEach((k, v) -> System.out.printf("%-10s %d%n", k, v));
for (var e : notes.entrySet()) { }                        // entrySet : la seule facon propre d'iterer un Map
var cles = notes.keySet(); var valeurs = notes.values();
notes.containsKey("proto"); notes.containsValue(9);

// List : tri, recherche, sous-listes
noms.sort(Comparator.comparingInt(String::length).thenComparing(Comparator.naturalOrder()));
Collections.sort(noms);                                    // ancien ; noms.sort(...) est preferable
List<String> sous = noms.subList(1, 3);                     // VUE : ecrire dedans ecrit dans noms !
int idx = Collections.binarySearch(trie, "Nour");            // <0 -> (-point d'insertion - 1)
noms.removeIf(s -> s.length() < 3);
noms.replaceAll(String::toUpperCase);
var copieTriee = noms.stream().sorted().toList();

// Les trois collections utilitaires a connaitre
List<String> vid = List.of();                              // ou : Collections.emptyList()
List<String> vueFixe = Collections.unmodifiableList(noms); // vue, pas une copie !
var synchro = Collections.synchronizedList(new ArrayList<String>());  // legacy : prefere les classes java.util.concurrent

// Ensembles : les operations qui ressemblent a ce qu'on veut vraiment
var a = Set.of("php", "java", "js");
var b = Set.of("java", "go");
var union = new HashSet<>(a); union.addAll(b);
var intersection = new HashSet<>(a); intersection.retainAll(b);   // ou : a.stream().filter(b::contains)
var difference = new HashSet<>(a); difference.removeAll(b);
a.containsAll(b);                                            // inclusion
\end{codejava}

\subsection{Ce qui casse les collections (les 10 classiques)}

\begin{enumerate}
  \item \verb|ConcurrentModificationException|~: modifier une collection pendant une
        \verb|for-each|. Correction~: \verb|removeIf|, un \verb|Iterator\#remove|, ou une copie.
  \item Cl\'e de \verb|HashMap| mutable~: l'objet change de \verb|hashCode| apr\`es insertion, et il est
        perdu \`a jamais. Utilise des cl\'es immuables (records, \verb|String|, \verb|Long|).
  \item \verb|equals| sans \verb|hashCode| (ou l'inverse)~: le Set accepte des doublons, le Get
        \'echoue. Les records \'evitent le probl\`eme.
  \item \verb|Arrays.asList(tab)|~: vue de taille fixe~; un \verb|add| l\`eve
        \verb|UnsupportedOperationException|. \'Ecris \verb|new ArrayList<>(Arrays.asList(...))|.
  \item \verb|List.of(...)| contient \verb|null|~: \verb|NullPointerException| imm\'ediate (c'est voulu).
  \item Confondre \verb|isEmpty()| et \verb|size() == 0| avec un \verb|Optional|~: le premier est
        une question, le second un d\'etail.
  \item Utiliser un \verb|Vector|, un \verb|Hashtable|, un \verb|Stack|~: legacies synchronis\'es,
        \`a remplacer respectivement par \verb|ArrayList| (d\'ej\`a s\^ur en lecture), \verb|ConcurrentHashMap|,
        \verb|ArrayDeque|.
  \item \verb|subList| et \verb|stream().toList()|~: le premier est une vue, le second immuable
        depuis Java~16 (avant~: pas d'ordre garanti).
  \item Comparaison de cha\^ines avec \verb|==| (voir chapitre pr\'ec\'edent) et tri non localis\'e~:
        utilise \verb|Collator.getInstance(Locale.FRENCH)| pour un tri humain correct.
  \item Modifier une collection partag\'ee sans synchronisation~: \verb|ConcurrentHashMap| et les
        classes \verb|concurrent| existent pour \c{c}a~; \verb|Collections.synchronizedX| verrouille
        tout l'objet (et n'est pas suffisant pour \verb|forEach|~: il faut verrouiller \`a la main).
\end{enumerate}

\section{G\'en\'eriques~: sans peur}

\begin{codejava}
// 1. Classe generique : un type parametre, pas un Object + casts
public record Pageable<T>(int page, int taille, List<T> contenus, long total) {
    public <R> Pageable<R> map(java.util.function.Function<? super T, ? extends R> f) {
        return new Pageable<>(page, taille, contenus.stream().map(f).toList(), total);
    }
}

// 2. Methode generique : le type est infer\'e a l'appel
public static <T extends Comparable<T>> T max(List<T> l) {
    return l.stream().max(Comparator.naturalOrder()).orElseThrow();
}
public static <T, K, V> Map<K, V> grouper(Collection<T> src, Function<T, K> cle, Function<T, V> valeur) {
    return src.stream().collect(Collectors.toMap(cle, valeur, (a, b) -> a, LinkedHashMap::new));
}

// 3. Bornes : "T extends" veut aussi bien dire "implements"
interface Identifiable { long id(); }
static <T extends Identifiable & Comparable<T>> List<T> trierParId(List<T> l) { ... }
static <T> void copier(List<? super T> dest, List<? extends T> src) {   // PECS : Producer Extends, Consumer Super
    for (T t : src) { dest.add(t); }
}
// Lecture depuis une collection -> ? extends ; ecriture vers une collection -> ? super.

// 4. Ce qui n'existe pas, et pourquoi (effacement des types)
// new T()            -> impossible : T n'existe plus a l'execution
// T[] tab = new T[n] -> impossible : prefere List<T> ; ou Array.newInstance(clazz, n)
// catch (T e)        -> impossible
// instanceof List<String> -> interdit ; autorise : instanceof List<?>
// static T champ     -> impossible dans une classe generique (le type est par instance)

// 5. Le "unchecked" honnete : circonscrire le warning a une ligne
@SuppressWarnings("unchecked")
private static <T> T[] tableauDe(List<T> l, Class<T> type) {
    @SuppressWarnings("unchecked")
    T[] tab = (T[]) java.lang.reflect.Array.newInstance(type, l.size());
    return l.toArray(tab);
}

// 6. Les collections generiques de la bibliotheque : lis-les comme des signatures
List<String> l = new ArrayList<>();          // <> : inf\'erence depuis le type d\'eclar\'e (Java 7+)
var l2 = new ArrayList<String>();             // var : inf\'erence depuis l'initialisateur
Map<String, List<Integer>> m = new HashMap<>();
Comparator<? super Article> cmp = Comparator.comparingLong(Article::id);
\end{codejava}

\section{Exceptions~: le mod\`ele}

\begin{codejava}
/*
  Throwable
   \boxv \boxh \boxh Error                    : VMStackOverflow, OutOfMemoryError -- ne jamais capturer
   \boxe \boxh \boxh Exception
        \boxv \boxh \boxh RuntimeException    : NON v\'erifi\'ees (le compilateur n'exige rien)
        \boxv \boxv \boxv   IllegalArgumentException, IllegalStateException, NullPointerException,
        \boxv \boxv \boxv   IndexOutOfBoundsException, ClassCastException, NumberFormatException,
        \boxv \boxv \boxv   ConcurrentModificationException, ArithmeticException, UnsupportedOperationException
        \boxe \boxh \boxh Exception (autres)  : VERIFIEES ("checked") -> try/catch ou throws
              IOException, SQLException, InterruptedException,
              ClassNotFoundException, FileNotFoundException, URISyntaxException, ...

  Regle de conception :
    - verifiee   = "cela peut arriver, et l'appelant PEUT en faire quelque chose" (reessayer,
                   demander un autre fichier, basculer sur un miroir).
    - non verifiee = "bug, ou etat irr\'ecuperable ici" -> remonter, logguer, faire echouer.
*/

// 1. Le bloc complet
public List<Article> charger(java.nio.file.Path chemin) throws IOException {
    try {
        return java.nio.file.Files.readAllLines(chemin).stream()
                .map(String::trim).filter(s -> !s.isBlank()).map(Article::deLigne).toList();
    } catch (java.nio.file.NoSuchFileException e) {
        throw new IOException("configuration absente : " + chemin, e);   // CONTEXTE + CAUSE
    } catch (java.nio.file.AccessDeniedException e) {
        logger.log(Level.WARNING, "acces refuse, repli sur les defauts", e);
        return List.of();
    }
    // pas de "catch (Exception e) { }" : c'est une ligne qui annule la securite du langage
}

// 2. finally : s'execute toujours (y compris sur un return du try) -- a reserver au nettoyage
java.io.InputStream flux = null;
try {
    flux = new java.io.FileInputStream(fichier);
    return lire(flux);
} finally {
    if (flux != null) { flux.close(); }
}
// 3. try-with-resources : la version correcte, ferme dans l'ordre inverse, supprime les exceptions
//    masquees (elles sont ajoutees en "suppressed")
try (var entree = new java.io.BufferedInputStream(new java.io.FileInputStream(fichier));
     var sortie = new java.io.OutputStreamWriter(java.lang.System.out, java.nio.charset.StandardCharsets.UTF_8)) {
    entree.transferTo(java.io.OutputStream.nullOutputStream());
} catch (java.io.IOException e) {
    throw new RuntimeException(e);
}
// Une ressource = AutoCloseable ; la methode close() ne doit pas lancer d'exception ->
// pour les fermetures "douces" : try { c.close(); } catch (Exception ignore) { }

// 4. Multi-catch : quand le traitement est identique, sans perdre le type
try { traitement(); }
catch (SQLException | IOException e) { throw new ErreurDonnees("acces aux donnees", e); }

// 5. Les exceptions metier : nommees, specifiques, transportant le contexte
public sealed class ErreurDomaine extends RuntimeException permits ErreurValidation, ErreurIntrouvable, ErreurConflit {
    private final String code;
    ErreurDomaine(String message, String code, Throwable cause) { super(message, cause); this.code = code; }
    public String code() { return code; }
}
public final class ErreurValidation extends ErreurDomaine {
    private final Map<String, List<String>> champs;
    public ErreurValidation(Map<String, List<String>> champs) {
        super(champs.size() + " champ(s) invalides", "validation", null);
        this.champs = Map.copyOf(champs);
    }
    public Map<String, List<String>> champs() { return champs; }
}

// 6. Interruption : le seul cas ou le "swallow" est une faute grave
try {
    Thread.sleep(1000);
} catch (InterruptedException e) {
    Thread.currentThread().interrupt();      // RESTAURE le drapeau : l'appelant doit pouvoir reagir
    throw new IllegalStateException("interrompu", e);
}

// 7. Precondition en une ligne (le code d\'efensif propre)
Objects.requireNonNull(liste, "liste");
Objects.checkFromIndexSize(debut, longueur, tab.length);
if (montant <= 0) { throw new IllegalArgumentException("montant=%d (attendu > 0)".formatted(montant)); }
\end{codejava}

\begin{piege}
Trois anti-patterns d'exception, par ordre de nocivit\'e~: (1) \verb|catch (Exception e) {}| vide ---
tu transformes une panne visible en corruption silencieuse~; (2) \verb|throw new Exception("erreur")|
--- message sans contexte, cause perdue, type inutilisable~; (3) utiliser une exception pour un
flux normal (une \og ~page suivante~ \fg{} est-elle une anomalie~? non~: renvoie \verb|Optional| ou un
bool\'een). Une exception co\^ut\'e la construction d'un stack~: ne les utilise pas comme \og ~return~ \fg{}.
\end{piege}

\begin{astuce}
Le message d'exception est une interface~: il doit contenir \emph{les valeurs} qui ont pos\'e probl\`eme
(\verb|"montant=-3 pour la commande 42"| et non \verb|"montant invalide"|). Et conserve toujours la
cause~: \verb|new ErreurX("...", e)|. Un \verb|RuntimeException| sans cause en production est une
enqu\^ete qui commence par \og ~de toute fa\c{c}on on ne sait pas~ \fg{}.
\end{astuce}

\begin{exercice}
\begin{enumerate}
  \item \'Ecris \verb|LRUCache<K,V>| sur \verb|LinkedHashMap| (\verb|removeEldestEntry|)~: capacit\'e,
        \verb|get|/ \verb|put|/ \verb|evict\'e s|, et 8 tests qui prouvent l'ordre d'\'eviction.
  \item \'Ecris un \verb|Trie| (arbre de pr\'efixes) g\'en\'erique~: \verb|ajouter|, \verb|contient|,
        \verb|commencePar|, \verb|autocompleter(prefixe, n)|~; compare la complexit\'e \`a un
        \verb|HashSet| sur 100 000 mots.
  \item Construis un \verb|TopK<T>| avec une \verb|PriorityQueue| \`a capacit\'e fix\'ee~; puis la
        m\^eme approche avec un \verb|TreeMap<Score, Set<Id>| (ordre d\'ecroissant, scores \'egaux).
  \item \'Ecris une hi\'erarchie d'exceptions m\'etier compl\`ete pour une boutique~: inventaire insuffisant,
        coupon expir\'e, paiement refus\'e, adresse non livrable~: avec codes, donn\'ees contextuelles,
        et un mapping vers les codes HTTP 400/404/409/422.
  \item Migre un morceau de code du style \og ~\verb|null| retourn\'e~ \fg{} vers le style Java moderne~:
        \verb|Optional| en retour, \verb|orElseThrow| avec exception m\'etier, \verb|Objects.requireNonNull|
        aux entr\'ees~; et ajoute \verb|@Nullable| (JSpecify ou checkerframework) pour les fronti\`eres.
  \item Fais tourner un \verb|-Xmx32m| avec une liste de 10 millions d'\'el\'ements~: observe
        \verb|OutOfMemoryError|~; explique-le \`a voix haute~; puis r\'eduis la m\'emoire r\'eelle avec un
        \verb|int[]| et compare les chiffres.
\end{enumerate}
\end{exercice}

\begin{ressources}
\begin{itemize}
  \item Trails de la doc Collections --- \url{https://docs.oracle.com/javase/tutorial/collections/}
  \item Javadoc HashMap (l'impl\'ementation, les seuils d'arbre, la charge) --- \url{https://docs.oracle.com/en/java/javase/21/docs/api/java.base/java/util/HashMap.html}
  \item Generics (tutoriel officiel + Angelika Langer's FAQ pour les cas tordus) --- \url{https://www.angelikalanger.com/GenericsFAQ/JavaGenericsFAQ.html}
  \item \url{https://dev.java/learn/exceptions/} --- et \link{EJCH}{https://www.oreilly.com/library/view/effective-java-3rd/}
  \item \link{JSpecify}{https://jspecify.dev/} pour les annotations de nullit\'e
\end{itemize}
\end{ressources}

\index{collections}\index{generiques}\index{exceptions}\index{HashMap}
% ============================================================================
\chapter{Fichiers, flux, dates et internationalisation}
\label{ch:java-io}
% ============================================================================

\begin{objectif}
\begin{itemize}
  \item Lire et \'ecrire des fichiers proprement~: \verb|Files|, \verb|Path|, flux tamponn\'es,
        copies, parcours r\'ecursifs.
  \item Ne plus jamais se tromper d'encodage (UTF-8 explicite) ni de s\'eparateur de chemin.
  \item Manipuler dates et fuseaux avec \verb|java.time| et formater nombres/monnaies par locale.
  \item G\'erer s\'erialisation JSON, CSV, propri\'et\'es, ressources embarqu\'ees.
  \item Conna\^itre les bons outils pour les gros fichiers (flux, \verb|MappedByteBuffer|, limites).
\end{itemize}
\end{objectif}

\section{Chemin et fichiers : \texttt{java.nio.file}}

\begin{codejava}
import java.io.IOException;
import java.io.UncheckedIOException;
import java.nio.charset.StandardCharsets;
import java.nio.file.*;
import java.util.Comparator;
import java.util.stream.Stream;

// Construction : JAMAIS de concatenation avec "/" (Windows casse) ; les separators sont geres
Path racine = Path.of(System.getProperty("user.dir"), "data");
Path fichier = racine.resolve("notes.txt");            // /chemin/vers/data/notes.txt
Path absolut = fichier.toAbsolutePath().normalize();     // resout ".." et symlinks
Path parent = absolut.getParent();
Path relatif = racine.relativize(absolut);                // "notes.txt"
Path maison = Path.of(System.getProperty("user.home"));
Files.exists(fichier); Files.isDirectory(fichier); Files.isRegularFile(fichier);
Files.isReadable(fichier); Files.isWritable(fichier); Files.isHidden(fichier);
fichier.getFileName(); fichier.getName(1); fichier.getNameCount();
fichier.endsWith(".txt"); fichier.startsWith(racine);
Files.size(fichier);
FileTime t = Files.getLastModifiedTime(fichier);

// Ecrire (les trois modes)
Files.writeString(fichier, "Ligne 1\nLigne 2\n", StandardCharsets.UTF_8,
        StandardOpenOption.CREATE, StandardOpenOption.APPEND);
Files.write(fichier, List.of("a", "b"), StandardCharsets.UTF_8);      // remplace
Files.createDirectories(racine);                                        // mkdir -p
Files.createFile(fichier);                                               // leverait FileAlreadyExists
Files.copy(Paths.get("/tmp/a"), fichier, StandardCopyOption.REPLACE_EXISTING);
Files.move(fichier, fichier.resolveSibling("notes.md"), StandardCopyOption.ATOMIC_MOVE);
Files.delete(fichier);                                                   // FileNotFound -> exception
Files.deleteIfExists(fichier);

// Lire
String contenu = Files.readString(fichier, StandardCharsets.UTF_8);   // tout (<= quelques Mo)
List<String> lignes = Files.readAllLines(fichier, StandardCharsets.UTF_8);
byte[] octets = Files.readAllBytes(fichier);
try (Stream<String> flux = Files.lines(fichier, StandardCharsets.UTF_8)) {   // GROS fichier
    long vides = flux.filter(s -> !s.isBlank()).count();
} catch (IOException e) { throw new UncheckedIOException(e); }              // idiomatique : wrapper

// Parcourir
try (Stream<Path> marche = Files.walk(racine, 3, FileVisitOption.FOLLOW_LINKS)) {
    marche.filter(Files::isRegularFile)
          .filter(p -> p.toString().endsWith(".java"))
          .sorted()
          .forEach(p -> System.out.printf("%8d %s%n", p.toFile().length(), racine.relativize(p)));
}
try (var entree = Files.newDirectoryStream(racine, "*.{java,kt}")) { for (Path p : entree) { } }

// Attributs POSIX et suivi de modifications
PosixFilePermissions.fromString("rwxr-xr-x");
Files.getOwner(fichier);
var obs = FileSystems.getDefault().watch(racine, StandardWatchEventKinds.ENTRY_MODIFY);
while (true) { WatchKey cle = obs.poll(2, java.util.concurrent.TimeUnit.SECONDS); /* ... cle.reset() */ }

// Copie a la volee entre flux (le "pipe" de la bibliotheque)
try (var in = Files.newInputStream(source); var out = Files.newOutputStream(dest)) {
    in.transferTo(out);
}
// Fichier temporaire propre (toujours supprime a la fermeture du JVM) :
Path tmp = Files.createTempFile("rapport", ".csv");
tmp.toFile().deleteOnExit();     // ou : new FileOutputStream(f, true) ; ou try-with-resources
\end{codejava}

\subsection{Flux : octets, caract\`eres, tampons}

\begin{codejava}
// Les deux familles : *Stream (octets) et *Reader/*Writer (caracteres, avec charset).
// Regle : ne JAMAIS lire des caracteres avec un InputStream sans charset explicite.
try (var flux = Files.newBufferedReader(fichier, StandardCharsets.UTF_8)) {
    String ligne;
    while ((ligne = flux.readLine()) != null) { traiter(ligne); }
}

try (var in = new java.io.BufferedInputStream(Files.newInputStream(fichier), 64 * 1024);
     var out = new java.io.BufferedOutputStream(Files.newOutputStream(dest), 64 * 1024)) {
    in.transferTo(out);
    out.flush();                                   // obligatoire avant de lire ailleurs (ou auto \`a close)
}

// Data : primitives en binaire (protocoles, fichiers de sauvegarde maison)
try (var d = new java.io.DataOutputStream(new java.io.FileOutputStream("sauvegarde.bin"))) {
    d.writeInt(42); d.writeUTF("Bonjour"); d.writeLong(System.currentTimeMillis()); d.writeBoolean(true);
}
try (var d = new java.io.DataInputStream(new java.io.FileInputStream("sauvegarde.bin"))) {
    int n = d.readInt(); String s = d.readUTF(); long t = d.readLong(); boolean b = d.readBoolean();
}

// PrintWriter / Scanner : parfaits pour les scripts et l'apprentissage console
var sc = new java.util.Scanner(System.in, StandardCharsets.UTF_8);
System.out.print("Entier ? "); int valeur = sc.hasNextInt() ? sc.nextInt() : -1;
System.out.printf("%,d en euros%n", 1234);         // 1 234 (espace ins\'ecable selon la locale)
// Java 21+ : java.io.StreamScanner si tu \'ecrivais un parseur \`a la main
\end{codejava}

\section{Dates et heures---\texttt{java.time}}

\begin{codejava}
import java.time.*;
import java.time.format.*;
import java.time.temporal.*;

// Instant : un point sur la ligne du temps (UTC, precis a la nanoseconde) -- pour stocker
Instant maintenant = Instant.now();
Instant epoch = Instant.EPOCH;
Instant date = Instant.parse("2026-01-31T10:15:30Z");

// LocalDateTime / LocalDate : sans fuseau -- pour un horaire de train, une date de naissance
LocalDate jour = LocalDate.of(2026, java.time.Month.JANVIER, 31);
LocalDate aujourdhui = LocalDate.now();
LocalTime heure = LocalTime.of(9, 30);
LocalDateTime horloge = LocalDateTime.of(jour, heure);
ZonedDateTime aDakar = ZonedDateTime.of(LocalDateTime.now(), ZoneId.of("Africa/Dakar"));
OffsetDateTime offset = OffsetDateTime.parse("2026-01-31T10:15:30+01:00");

// Manipulations : tout est immuable, tout renvoie un nouvel objet
LocalDate demain = jour.plusDays(1);
LocalDate finMois = jour.with(TemporalAdjusters.lastDayOfMonth());
LocalDate prochainLundi = jour.with(TemporalAdjusters.next(DayOfWeek.MONDAY));
Period periode = Period.between(jour, jour.plusMonths(3).plusDays(5));   // annees/mois/jours
Duration duree = Duration.between(heure, heure.plusMinutes(90));
long secondes = duree.toSeconds();
LocalDate memeJourLAnProchain = jour.plusYears(1).withDayOfMonth(Math.min(jour.getDayOfMonth(),
        Year.of(jour.getYear() + 1).length()));      // le 29 fevrier existe... ou pas
boolean anneeBissextile = Year.of(jour.getYear()).isLeap();
ZoneId sys = ZoneId.systemDefault();

// Conversions : les trois fronts reels
LocalDate enDate = aDakar.toLocalDate();
Instant depuisLocal = LocalDate.of(2026, 1, 31).atStartOfDay(ZoneId.of("UTC")).toInstant();
Date ancien = java.util.Date.from(instant);                    // pont vers JDBC / API heritees
Instant depuisSql = java.sql.Timestamp.valueOf(local).toInstant();

// Formatage / analyse : DateTimeFormatter est couteux -> une constante statique
static final DateTimeFormatter JOUR_FR = DateTimeFormatter.ofPattern("EEEE d MMMM yyyy", Locale.FRENCH);
static final DateTimeFormatter COMPACT = DateTimeFormatter.BASIC_ISO_DATE;        // 20260131
jour.format(JOUR_FR);                              // "samedi 31 janvier 2026"
jour.format(DateTimeFormatter.ISO_LOCAL_DATE);       // 2026-01-31
jour.format(DateTimeFormatter.ofLocalizedDate(FormatStyle.FULL));
LocalDate analyse = LocalDate.parse("31.01.2026", DateTimeFormatter.ofPattern("dd.MM.yyyy"));
var options = new DateTimeFormatterBuilder()
        .appendPattern("dd/MM/yyyy")
        .optionalStart().appendPattern(" HH:mm").optionalEnd()
        .parseDefaulting(ChronoField.HOUR_OF_DAY, 0)
        .toFormatter(Locale.FRENCH);
// Toujours verifier le parse : DateTimeParseException avec index du caractere fautif.
\end{codejava}

\begin{piege}
\verb|LocalDateTime| ne contient \textbf{pas} de fuseau~: une minute de diff\'erence de configuration
du serveur et tes factures sont fausses. R\`egle~: on \emph{stocke} en \verb|Instant| (ou
\verb|OffsetDateTime|) en UTC, on \emph{affiche} en fuseau local de l'utilisateur, et on ne fait
jamais d'arithm\'etique de \og ~jours~ \fg{} sur un \verb|Instant|~: un jour civil peut faire 23 ou 25 heures
(changement d'heure~! ).
\end{piege}

\section{Nombres, monnaie et i18n}

\begin{codejava}
import java.text.*;
import java.util.*;

Locale francais = Locale.forLanguageTag("fr-FR");
NumberFormat nombre = NumberFormat.getIntegerInstance(francais);
nombre.format(1234567);                                   // "1 234 567" (espace ins\'ecable)
NumberFormat devise = NumberFormat.getCurrencyInstance(francais);
devise.format(12.5);                                       // "12,50 EUR" / "12,50 \euro "
NumberFormat pourcent = NumberFormat.getPercentInstance();
pourcent.setMinimumFractionDigits(1); pourcent.format(0.1234);   // "12,3 %"
DecimalFormat modele = new DecimalFormat("#,##0.00 \euro ", new DecimalFormatSymbols(francais));
modele.format(1234.5);
ParsePosition pp = new ParsePosition(0);
Number lu = nombre.parse("1 234 567", pp);                    // attention aux espaces ins\'ecables
Collator collation = Collator.getInstance(francais);
collation.setStrength(Collator.PRIMARY);
List.of("\'ecole", "\'elan", "effort").stream().sorted(collation).toList();
BreakIterator mots = BreakIterator.getWordInstance(francais);  // decoupage linguistique reel
ChoiceFormat pluriel = new ChoiceFormat("0#pas d'article|1#un article|1<#%d articles");
MessageFormat msg = new MessageFormat(frMessage, francais);      // "Bonjour {0}, {1,number,integer} article(s)"
msg.format(new Object[]{"Nour", 3});

// Ressources par locale : src/main/resources/messages_fr.properties, messages_en.properties...
ResourceBundle rb = ResourceBundle.getBundle("messages", francais);
String bienvenue = rb.getString("titre");
// La pratique moderne : les fichiers .properties encodes en UTF-8, ou un fichier unique JSON/YAML
// charge a la main ; et java.text/icu4j (IBM) pour les besoins serieux.
var icu = com.ibm.icu.text.RuleBasedNumberFormat(Locale.FRENCH, com.ibm.icu.text.RuleBasedNumberFormat.SPELLOUT);
icu.format(42);                                              // "quarante-deux" (bibliotheque externe)
\end{codejava}

\section{S\'erialiser : JSON, propri\'et\'es, CSV, objets}

\begin{codejava}
// 1. Proprietes : la config simple, sans dependance
Properties props = new Properties();
try (var in = Files.newInputStream(Path.of("config.properties"))) { props.load(in); }
String url = props.getProperty("db.url", "jdbc:h2:mem:test");
props.store(Files.newOutputStream(Path.of("out.properties")), "genere par l'application");

// 2. JSON : Jackson (le standard de fait). Ajouter la dependance, puis :
//    com.fasterxml.jackson.core:jackson-databind
var mapper = new com.fasterxml.jackson.databind.ObjectMapper()
        .findAndRegisterModules()                       // java.time, etc.
        .configure(com.fasterxml.jackson.databind.DeserializationFeature.FAIL_ON_UNKNOWN_PROPERTIES, false)
        .setSerializationInclusion(com.fasterxml.jackson.annotation.JsonInclude.Include.NON_NULL);
var json = mapper.writeValueAsString(article);
Article relu = mapper.readValue(json, Article.class);
List<Article> liste = mapper.readValue(json, mapper.getTypeFactory()
        .constructCollectionType(List.class, Article.class));   // ou : new TypeReference<List<Article>>() {}

// DTO d'API : un record + annotations -- c'est propre et immuable
public record RequeteArticle(
        @jakarta.validation.constraints.NotBlank @jakarta.validation.constraints.Size(min = 3, max = 160) String titre,
        @jakarta.validation.constraints.NotBlank String corps,
        List<String> tags,
        @jakarta.annotation.Nullable Instant publieLe) { }

// 3. CSV : pas de bibliotheque necessaire pour du simple, mais-open commons-csv ou
//    objects/uni-csv des qu'il y a des guillemets ou des retours a la ligne.
String ligne = "id,nom\n1,\"O'Brien, junior\",42\n".replace("\r", "");

// 4. Serialisation d'objets (Serializable) : A EVITER pour de la persistance ou du reseau.
//    Fragile (versionnement de classe), lent, vecteur de failles (deserialization gadget chains).
//    Utilise JSON, Protobuf, Avro, ou un format texte define par toi.
//    (La serialization reste utile pour RMI/ehcache, et la bibliotheque la remplace ailleurs.)
\end{codejava}

\begin{astuce}
Pour de la configuration typ\'ee sans d\'ependance~: \verb|Map<String,String>| lue depuis un
\verb|.env| (15 lignes de parsing) + des enregistrements \verb|record Config(...) | valid\'es au
d\'emarrage. C'est plus lisible que 40 appels \`a \verb|System.getenv()| \'eparpill\'es dans le code, et
\c{c}a \'echoue imm\'ediatement si une variable manque -- au lieu de tomber en production un vendredi.
\end{astuce}

\subsection{Gros fichiers~: la bo\^ite \`a outils}

\begin{codejava}
// 1. Flux de lignes : memoire constante, integration parfaite avec Streams
try (var lignes = Files.lines(Path.of("gros.log"))) {
    lignes.map(l -> l.split(" ")[0]).distinct().parallel().forEach(System.out::println);
}
// 2. BufferedReader manuel quand il faut compter/avancer
try (var r = Files.newBufferedReader(path)) {
    char[] tampon = new char[8192];
    int lus;
    while ((lus = r.read(tampon)) != -1) { traiter(new String(tampon, 0, lus)); }
}
// 3. Ecriture en masse : un BufferedWriter + des sauts de tampon de 64 Ko
// 4. memory-mapped (lecture quasi nulle copie) : FileChannel.map -- attention a la taille
//    de la zone (limite de 2 Go par region) ; GarbageCollector ne libere pas la carte a coup sur :
//    a debloquer explicitement ou a eviter si tu ne maitrises pas.
// 5. Canaux + ByteChannel.transferFrom : le zero-copy reel
//    try (var src = FileChannel.open(a); var dst = FileChannel.open(b, WRITE, CREATE, TRUNCATE_EXISTING)) {
//        long taille = src.size(); long pos = 0; do { pos += src.transferTo(pos, taille - pos, dst); } while (pos < taille);
//    }
// 6. Compression integree : GZIP/ZIP via ZipOutputStream, Deflater/Inflater,
//    ou Apache Commons Compress (tar.gz, 7z, zstd) si besoin.
\end{codejava}

\begin{exercice}
\begin{enumerate}
  \item \'Ecris un utilitaire \verb|Fichiers|~: \verb|lireEnLignes|, \verb|ecrireAtomiquement|
        (\'ecrire dans un fichier temporaire puis \verb|Files.move(..., ATOMIC\_MOVE)|),
        \verb|trouver(racine, glob)|, \verb|tailleTotale(dossier)|, \verb|copierArborescence|
        (avec permissions conserv\'ees).
  \item Construis un journaliseur maison~: rotation quotidienne et par taille (10~Mo),
        niveaux, tampon de 32~Ko vid\'es p\'eriodiquement, format JSON sur une ligne,
        et gestion correcte des fermetures~; compare le d\'ebit \`a \verb|System.out.println|.
  \item Analyseur de logs~: agr\`ege par heure, par statut, par IP~; top 10 des URL les plus lentes
        (percentiles p50/p95/p99~! )~; affiche un rapport align\'e avec \verb|printf|~: le tout
        en flux (un fichier de 2~Go ne doit pas \^etre charg\'e en RAM).
  \item Convertis un CSV en JSON et inversement, en g\'erant les champs entre guillemets, les
        retours \`a la ligne internes, les nombres localis\'es, les dates \verb|dd/MM/yyyy|,
        et en \'ecrivant les deux formats en UTF-8 avec BOM pour l'export Excel.
  \item \'Ecris un calculateur d'anniversaires~: jours restants, \^age exact (ann\'ees/mois/jours),
        gestion du 29~f\'evrier, fuseau du pays de l'utilisateur, et formatage localis\'e.
        Teste avec les changements d'heure (Europe/Paris le dernier dimanche de mars).
  \item G\'en\`ere un petit rapport PDF~: \link{iText}{https://itextpdf.com/} ou
        \link{OpenPDF}{https://github.com/LibrePDF/OpenPDF}~: en-t\^ete localis\'e, tableau tri\'e,
        totaux, pagination. Puis compare avec l'export HTML + impression navigateur (souvent
        suffisant, et dix fois moins de code).
\end{enumerate}
\end{exercice}

\begin{ressources}
\begin{itemize}
  \item Tutorial officiel I/O --- \url{https://dev.java/learn/io/}
  \item \url{https://docs.oracle.com/en/java/javase/21/docs/api/java.base/java/nio/file/Files.html}
  \item \url{https://docs.oracle.com/en/java/javase/21/docs/api/java.base/java/time/package-summary.html}
  \item Jackson --- \url{https://github.com/FasterXML/jackson-databind} --- Gson (alternative) --- \url{https://github.com/google/gson}
  \item OpenCSV --- \url{http://opencsv.sourceforge.net/} --- Apache Commons CSV --- \url{https://commons.apache.org/proper/commons-csv/}
\end{itemize}
\end{ressources}

\index{fichiers}\index{NIO}\index{java.time}\index{i18n}\index{Jackson}
% ============================================================================
\chapter{Programmation fonctionnelle~: lambdas, Streams, collecteurs}
\label{ch:java-fonctionnel}
% ============================================================================

\begin{objectif}
\begin{itemize}
  \item \'Ecrire des lambdas lisibles, comprendre la capture et les interfaces fonctionnelles.
  \item Ma\^itriser le pipeline Streams~: sources, op\'erations interm\'ediaires, terminales.
  \item Utiliser \texttt{Collectors} sans recopier un exemple au hasard~: groupage, partition,
        agr\'egats, maps imbriqu\'ees, collecteurs personnalis\'es.
  \item Savoir quand le flux parall\`ele aide vraiment, et quand il nuit.
  \item \'Ecrire du code fonctionnel qui reste d\'ebogable.
\end{itemize}
\end{objectif}

\section{Lambda~: la syntaxe compl\`ete, puis la r\`egle}

\begin{codejava}
// Les cinq formes, de la plus verbeuse a la plus concise (toutes equivalentes)
Comparator<Integer> c1 = (Integer a, Integer b) -> { return Integer.compare(a, b); };
Comparator<Integer> c2 = (a, b) -> Integer.compare(a, b);
Comparator<Integer> c3 = Comparator.comparingInt(Integer::intValue);
Comparator<Integer> c4 = Integer::compare;
Comparator<Integer> c5 = (a, b) -> a - b;             // correct pour des petits entiers, pas generique !

// Le type cible est infer\'e par le contexte : "var" ne fonctionne pas avec un lambda nu
var r = (Runnable) () -> System.out.println("ok");    // le cast donne le type cible
// Runnable r2 = () -> { };                             // correct
// var r3 = () -> { };                                  // ERREUR de compilation

// Une expression lambda peut etre une expression ou un bloc ; un bloc doit avoir un return.
Runnable vide = () -> { };
java.util.function.Supplier<String> consigne = () -> { var s = "a"; return s + "b"; };

// Les interfaces fonctionnelles du JDK (java.util.function) -- a memoriser par famille :
//   Predicate<T>  : boolean test(T)                 | Function<T,R> : R apply(T)
//   Consumer<T>   : void accept(T)                  | Supplier<T>   : T get()
//   UnaryOperator<T> / BinaryOperator<T>            | BiFunction<T,U,R>
//   IntFunction<R>, ToIntFunction<T>, IntToDoubleFunction, ... (variantes primitives)
//   Runnable, Callable<V> (peut lancer une exception), Comparator<T>, Predicate combinables
Predicate<Integer> pair = i -> i % 2 == 0;
Predicate<Integer> positif = i -> i > 0;
var combo = pair.and(positiv).or(i -> i == 0);
Function<Integer, Integer> fois2 = i -> i * 2;
Function<Integer, Integer> chaine = fois2.andThen(String::valueOf).andThen(s -> "[" + s + "]");
// composition : andThen (apres) vs compose (avant) -- la source d'erreur numero 1

// Capture : uniquement des variables "effectivement finales"
int facteur = 3;
Function<Integer, Integer> mult = x -> x * facteur;      // OK (facteur non modifie ensuite)
// facteur = 4;                                           // -> erreur de compilation : pas de reassignation
final StringBuilder acc = new StringBuilder();            // contournement propre : un objet mutable
Consumer<String> collecter = s -> acc.append(s).append('\n');
// Ce n'est pas une limitation technique mais un choix : capturer une variable locale mutable
// dans un lambda partage entre threads serait une bombe.
\end{codejava}

\begin{note}
En Java, un lambda est une instance d'interface fonctionnelle (pas une fonction de premi\`ere
classe flottante). Le bytecode utilise \verb|invokedynamic| + \verb|LambdaMetafactory|~:
l'objet est cr\'e\'e une fois, et l'appel est aussi rapide qu'une m\'ethode virtuelle. Concr\`etement~:
n'aie pas peur d'un lambda dans une boucle chaude~; aie peur de l'\emph{allocation} d'un
\verb|ArrayList| dans la m\^eme boucle.
\end{note}

\section{Streams~: le pipeline}

Un \verb|Stream| n'est \textbf{pas} une collection~: c'est un plan de calcul sur une source,
\emph{consommable une fois}, paresseux jusqu'\`a l'op\'eration terminale.

\begin{codejava}
Stream<Article> src = articles.stream();

// Intermediaires (paresseux, renvoie un Stream) :
src.filter(a -> a.estPublie() && a.tags().contains("java"))
   .map(Article::titre)
   .map(String::trim)
   .distinct()
   .sorted(Comparator.reverseOrder())
   .peek(t -> log.info("traitement de {}", t))        // peek = debug seulement
   .limit(20)
   .skip(5)
   .flatMap(t -> Stream.of(t.split(" ")))              // un element -> plusieurs
   .mapMulti(String.class, (t, out) -> { for (var w : t.split(" ")) { out.accept(w); } });

// Terminales (\'evaluent tout, une seule fois) :
List<String> l = stream.toList();                       // Java 16+ (immuable)
var copieModifiable = stream.collect(Collectors.toCollection(ArrayList::new));
Set<String> set = stream.collect(Collectors.toSet());
Map<Long, String> map = stream.collect(Collectors.toMap(Article::id, Article::titre));
Optional<Article> premier = stream.findFirst();          // rencontre le 1er (ordre = source)
Article un = stream.findAny().orElseThrow();              // en parallele : findAny est moins restrictif
long n = stream.count();
boolean aucun = stream.noneMatch(pred);                  // anyMatch/allMatch/noneMatch = courts-circuits
OptionalInt min = stream.mapToInt(Article::vues).min();
int somme = stream.map(Article::vues).reduce(0, Integer::sum);
String joint = stream.map(Article::titre).collect(Collectors.joining(", ", "[", "]"));
stream.forEach(System.out::println);                     // pas forEachOrdered (garde l'ordre en parallele)
var tableau = stream.toArray(String[]::new);
var optionnel = Optional.ofNullable(stream.findFirst().orElse(null));
\end{codejava}

\subsection{Les collecteurs, l\`a o\`u le talent se voit}

\begin{codejava}
import java.util.stream.Collectors;
import static java.util.stream.Collectors.*;

// Agr\'egations
long nbPublies = articles.stream().filter(Article::estPublie).count();
DoubleSummaryStatistics stats = articles.stream().mapToDouble(Article::vues).summaryStatistics();
stats.getCount(); stats.getAverage(); stats.getMin(); stats.getMax(); stats.getSum();
IntSummaryStatistics notes = eleves.stream().mapToInt(Eleve::note).summaryStatistics();

// Groupages
Map<Statut, List<Article>> parStatut = articles.stream().collect(groupingBy(Article::statut));
Map<Statut, Long> compteParStatut = articles.stream().collect(groupingBy(Article::statut, counting()));
Map<String, List<String>> titresParAnnee = articles.stream().collect(
        groupingBy(a -> String.valueOf(a.publieLe().getYear()),
                   collectingAndThen(mapping(Article::titre, toList()), l -> new ArrayList<>(l))));
Map<Statut, Map<String, List<Article>>> doubleCle = articles.stream()
        .collect(groupingBy(Article::statut, groupingBy(a -> a.tags().isEmpty() ? "sans-tag" : a.tags().get(0))));
Map<Boolean, List<Article>> partition = articles.stream().collect(partitioningBy(Article::estPublie));

// Maps et transformations
Map<Long, Article> parId = articles.stream().collect(toMap(Article::id, a -> a));
Map<Long, Article> parIdOrdonne = articles.stream().collect(toMap(Article::id, a -> a, (x, y) -> x, TreeMap::new));
Map<String, Integer> totauxParTag = articles.stream()
        .flatMap(a -> a.tags().stream())
        .collect(groupingBy(t -> t, summingInt(t -> 1)));
// (plus simple : groupingBy + counting)

// Collections cibles
var parSet = articles.stream().collect(toCollection(LinkedHashSet::new));
String csv = articles.stream().map(Article::titre).collect(joining(",", "\"", "\""));
Map<String, Long> top5 = stream.collect(groupingBy(f, counting())).entrySet().stream()
        .sorted(Map.Entry.<String, Long>comparingByValue().reversed()).limit(5)
        .collect(toMap(Map.Entry::getKey, Map.Entry::getValue, (a, b) -> a, LinkedHashMap::new));

// Collecteur sur mesure (reduction mutable propre)
record Compacteur(int total, int max, List<String> noms) {
    static Compacteur creer() { return new Compacteur(0, 0, new ArrayList<>()); }
    void ajouter(String nom, int valeur) { noms.add(nom); total += valeur; max = Math.max(max, valeur); }
    Compacteur fusionner(Compacteur o) {   // utilise en parallele
        noms.addAll(o.noms); return new Compacteur(total + o.total, Math.max(max, o.max), noms);
    }
}
var resume = flux.collect(Collector.of(
        Compacteur::creer,
        (c, a) -> c.ajouter(a.titre(), a.vues()),
        Compacteur::fusionner,
        Collector.Characteristics.CONCURRENT));

// Un stream ne se relit pas : si tu as besoin du resultat deux fois, materialise-le.
List<Article> materialise = articles.stream().filter(c::pertinent).toList();
\end{codejava}

\begin{tabz}{lXl}
  \headrow \thd{Collecteur} \& \thd{\`A la place de...} \& \thd{Note} \\
  \midrule
  \verb|toList()| (Java~16+) \& \verb|collect(toList())| \& immuable~; \verb|collect(Collectors.toList())| reste mutable. \\
  \verb|toUnmodifiableList()| \& -- \& Java~10+, et \'echoue si un \'el\'ement est \verb|null|. \\
  \verb|toMap(k, v)| \& une boucle + put | \textbf{l\`eve} sur cl\'e dupliqu\'ee~: passe le mergeur \verb|(a,b)->a|. \\
  \verb|toCollection(HashSet::new)| \& -- \& type de collection explicite. \\
  \verb|groupingBy(classifier)| \& \verb|Map<K, List<V>>| rempli \`a la main \& accepte un downstream. \\
  \verb|groupingBy(classifier, downstream)| \& -- \& \verb|counting|, \verb|summingInt|, \verb|reducing|, \verb|minByBy|... \\
  \verb|partitioningBy(pred)| \& deux listes \& 2 cl\'es seulement, plus rapide qu'un groupage. \\
  \verb|joining(delim, prefix, suffix)| \& \verb|StringBuilder| \& attention : lent sur de gros volumes. \\
  \verb|mapping(f, d)|, \verb|filtering(p, d)|, \verb|collectingAndThen(...)| \& -- \& la combinatoire qui \'evite un stream interm\'ediaire. \\
  \verb|summarizingInt/Long/Double| \& trois parcours | un seul. \\
  \verb|reducing(id, f, op)| \& un accumulateur | associativit\'e obligatoire pour le parall\`ele. \\
\end{tabz}

\section{Flux primitifs~: la performance sans emballage}

\begin{codejava}
// Stream<Integer> cree des objets Integer pour chaque element. Les flux primitifs non.
int total = articles.stream().mapToInt(Article::vues).sum();
long totalLong = articles.stream().mapToLong(Article::vues).asLongStream().sum();
double moyenne = articles.stream().mapToDouble(a -> a.vues()).average().orElse(0);
DoubleStream.of(1.1, 2.2).map(Math::sqrt).filter(d -> d > 1).summaryStatistics();
IntStream.range(0, n).map(i -> i * i).filter(i -> i % 2 == 0).boxed().toList();
IntStream.rangeClosed(1, 100).sum();                      // 1..100 (range = exclut la fin)
IntStream.iterate(1, v -> v <= 1000, v -> v * 2).toArray();
IntStream.generate(() -> (int) (Math.random() * 6) + 1).limit(10).toArray();
Stream.generate(() -> "constant").limit(3);
Stream.iterate(new int[]{0, 1}, t -> new int[]{t[1], t[0] + t[1]}).limit(10).map(t -> t[0]);  // Fibonacci
// Boite / deboite : .boxed() (primitif -> objet) ; .mapToInt() (objet -> primitif)
List<Integer> entiers = IntStream.of(1, 2, 3).boxed().toList();
OptionalInt vide = IntStream.empty().min();
int avecDefaut = IntStream.of().max().orElse(-1);        // les Optional primitifs ont orElse/orElseGet
\end{codejava}

\section{Parall\'elisme de flux~: le mode d'emploi honn\^ete}

\begin{codejava}
int somme = fluxEnMemoire.parallel().mapToInt(String::length).sum();

// Le parallel() utilise le FORK/JOIN COMMUN (threads = nb de coeurs - 1). Donc :
//  - interdit d'y faire du bloquant (IO, sommeil, appel HTTP) : tu affames les autres ;
//  - inutile sur une liste de 5 000 elements : le cout de partition > le gain ;
//  - rentable sur des donnees en memoire, sans \'etat partag\'e, avec une operation courte et pure ;
//  - a mesurer : parfois 8x, parfois 0.9x (contention, allocation, ordre impos\'e).
List<String> ok = stream.parallel().map(String::trim).filter(s -> !s.isEmpty()).toList();
// Operations "sans ordre" (findAny, noneMatch, unordered) sont les amies du parallele :
Optional<String> nImporteLequel = stream.parallel().unordered().filter(c::estBon).findAny();
// Limiter le pool pour un calcul isol\'e, sans polluer le commun :
try (var pool = new ForkJoinPool(4)) {                  // Java 19+: AutoCloseable
    resultat = pool.submit(() -> stream.parallel().map(...).reduce(...)).get();
}
// Regle d'or : si un calcul est trop gros pour un thread et trop petit pour un cluster,
// un stream parallele est la bonne idee. Sinon, les virtual threads (chapitre suivant) sont
// presque toujours meilleurs.
\end{codejava}

\subsection{\'Ecrire du code fonctionnel qui reste lisible}

\begin{enumerate}
  \item Une op\'eration par ligne dans les pipelines de plus de trois \'etapes~; retour \`a la ligne
        apr\`es chaque \verb|.|, indent\'es de quatre espaces.
  \item Un lamba de plus de trois lignes~: nomme-le (m\'ethode priv\'ee)~: \verb|a -> estPublieRecent(a)|.
  \item \verb|peek| n'est pas un outil de production~: il peut \^etre ignor\'e par le moteur
        (court-circuit). Utilise un compteur ou un vrai point de log.
  \item Pas d'effet de bord dans \verb|map|/ \verb|filter|~: le compilateur ne te sauvera pas,
        mais le parall\`ele et le r\'eordonnancement, si.
  \item Un stream consomm\'e est mort~: si tu it\`eres deux fois, mat\'erialise une \verb|List|.
  \item Ferme les flux issus de \verb|Files.lines| (try-with-resources)~: ils d\'etiennent un handle.
  \item La boucle \verb|for| n'est pas un aveu d'\'echec~: parcours avec indices, mutation d'un
        tableau, algorithme avec break conditionnel~? La boucle gagne. Les streams excellent sur
        \og ~transformer une collection en une autre~ \fg{}.
\end{enumerate}

\begin{codejava}
// Refactorisation type : de "Java 5" a "Java 21", en gardant la clarte
// AVANT
Map<String, List<Article>> res = new HashMap<>();
for (Article a : articles) {
    if (a.estPublie()) {
        for (String tag : a.tags()) {
            res.computeIfAbsent(tag, k -> new ArrayList<>()).add(a);
        }
    }
}
// APRES
Map<String, List<Article>> res = articles.stream()
        .filter(Article::estPublie)
        .flatMap(a -> a.tags().stream().map(t -> Map.entry(t, a)))
        .collect(groupingBy(Map.Entry::getKey, mapping(Map.Entry::getValue, toList())));
// et un jour, tu realise que computeIfAbsent etait plus rapide a lire : tu reviens en arriere.
// La bonne question n'est pas "fonctionnel ou boucle" mais "le lecteur comprend-il en 5 secondes".
\end{codejava}

\begin{exercice}
\begin{enumerate}
  \item \`A partir d'une liste de 100 000 enregistrements \verb|Vente(id, produit, region, montant, date)|,
        produis~: le chiffre par r\'egion~; le top 10 produits~: la croissance mensuelle~; le panier
        moyen~; les r\'egions sans vente~; et un export CSV tri\'e~: le tout en streams~: puis en boucles~;
        compare la lisibilit\'e, pas seulement le temps.
  \item \'Ecris un collecteur \verb|versStatistiques()| qui renvoie un \verb|record Stats(n, somme,
        moyenne, min, max, ecartType)| en un seul passage~: puis valide-le avec des propri\'et\'es
        (valeurs al\'eatoires, comparaison avec trois parcours s\'epar\'es).
  \item Construis un mini-\verb|Collectors| : \verb|toMapMulti|, \verb|toGroupedCount|, \verb|toSortedSet|.
  \item Utilise \verb|mapMulti| pour aplatisir un arbre (\verb|record Noeud(String nom, List<Noeud> enfants)|)
        en un flux de chemins \og ~a/b/c~ \fg{}~: puis compare \`a un \verb|flatMap| r\'ecursif.
  \item Calcule les premiers 10 000 nombres premiers~: (a) filtre na\"if en stream~; (b) crible
        d'\'Eratosth\`ene avec \verb|int[]|~; (c) version parall\`ele. Donne les trois temps et explique
        l'\'ecart. (Spoiler~: le stream sera 50 fois plus lent, et c'est formateur.)
  \item Ajoute une m\'ethode par d\'efaut \`a une de tes interfaces fonctionnelles~: \verb| puis |,
        \verb|enSecurite| (capture les exceptions~: \verb|Function<T,R>| -> \verb|Resultat<R>|)~;
        utilise-la pour transformer un parsing fragile en composition propre.
\end{enumerate}
\end{exercice}

\begin{ressources}
\begin{itemize}
  \item Tutorial officiel Streams --- \url{https://dev.java/learn/streams/}
  \item Javadoc \verb|Collectors| (le tableau des collecteurs vaut le d\'etour) --- \url{https://docs.oracle.com/en/java/javase/21/docs/api/java.base/java/util/stream/Collectors.html}
  \item \emph{Java Streams} (P. Sotuyodhy, en fran\c{c}ais) --- \url{https://blog.sotuyodhy.com/java-streams/}
  \item Article de Brian Goetz sur les limites du parall\'elisme --- \url{https://www.oracle.com/technical-resources/articles/java/architect-streams.html}
  \item Vavr (fonctionnel avanc\'e en Java) --- \url{https://www.vavr.io/}
\end{itemize}
\end{ressources}

\index{lambda}\index{Streams}\index{Collectors}\index{parallel stream}
% ============================================================================
\chapter{Concurrence~: threads, synchronisation, virtual threads}
\label{ch:java-concurrency}
% ============================================================================

\begin{objectif}
\begin{itemize}
  \item Comprendre le mod\`ele m\'emoire Java~: ce qui est visible, dans quel ordre, et pourquoi.
  \item Cr\'eer des t\^aches, les ex\'ecuter dans un pool, attendre, annuler, propager les erreurs.
  \item Choisir le bon outil de synchronisation~: \texttt{synchronized}, lock, atomique,
        file bloquante, \texttt{ConcurrentHashMap}.
  \item Utiliser les threads virtuels (Java~21) et savoir quand ils remplacent un pool.
  \item \'Ecrire du code concurrent testable et d\'ebogable.
\end{itemize}
\end{objectif}

\section{Le mod\`ele, sans romantisme}

Deux threads ex\'ecutent deux instructions~: rien ne garantit l'ordre d'interleaving, ni que la
\'ecriture de l'un soit visible de l'autre, \emph{sauf} s'il existe une relation
\textbf{happens-before}. Les points qui la cr\'eent~: lancer un thread (\verb|start|) se passe
avant son ex\'ecution~; la fin d'un thread se passe avant que \verb|join| revienne~; une sortie de
bloc \verb|synchronized| se passe avant l'entr\'ee du m\^eme moniteur~; un \verb|volatile write| avant
le \verb|volatile read|~; la lib\'eration d'un \verb|Lock| avant son acquisition~; tout est transitif.

\begin{codejava}
static class CompteurCasse { int n; void incrementer() { n++; } }     // Jamais sur : lecture+ecriture non atomiques
static class CompteurSur {
    private int n;
    synchronized void incrementer() { n++; }                     // verrou sur this
    synchronized int valeur() { return n; }
}
static class CompteurAtomique {
    private final java.util.concurrent.atomic.AtomicInteger n = new java.util.concurrent.atomic.AtomicInteger();
    void incrementer() { n.incrementAndGet(); }
    int valeur() { return n.get(); }
}
static class CompteurAdder {                                      // meilleur encore sous forte concurrence
    private final java.util.concurrent.atomic.LongAdder n = new java.util.concurrent.atomic.LongAdder();
    void incrementer() { n.increment(); }
    long valeur() { return n.sum(); }
}
// Verdict : 4 threads x 10 millions d'incrementations -> CompteurCasse renvoie une valeur au hasard
// (toujours <= 40M), CompteurSur ~140 ms, CompteurAtomique ~90 ms, LongAdder ~20 ms. A mesurer chez toi.

// Le defaut de visibilit\'e le plus vicieux :
class Travailleur extends Thread {
    private boolean arrete = false;                       // PAS volatile : le compilateur peut hoister
    public void run() { while (!arrete) { traiter(); } }
    void stopper() { arrete = true; }                      // peut ne jamais etre vu !
}
// Corriger par : volatile, ou un Lock, ou une variable atomique, ou une queue bloquante.
\end{codejava}

\section{Trois fa\c{c}ons de lancer une t\^ache}

\begin{codejava}
// 1. Runnable/Thread : pour un programme d'exemple, jamais pour une application.
var t = new Thread(() -> System.out.println("travail sur " + Thread.currentThread()), "worker-1");
t.setDaemon(true); t.start(); t.join(java.time.Duration.ofSeconds(5));

// 2. ExecutorService : un pool de threads de travail (la voie historique, reste pertinente)
java.util.concurrent.ExecutorService pool = java.util.concurrent.Executors.newFixedThreadPool(
        Runtime.getRuntime().availableProcessors(),
        r -> { var th = new Thread(r, "metier-%d".formatted(CPT.incrementAndGet())); th.setDaemon(true); return th; });
Future<Resultat> futur = pool.submit(() -> calculer());
Resultat r = futur.get(30, java.util.concurrent.TimeUnit.SECONDS);   // TimeoutException, ExecutionException
List<Future<Resultat>> lots = pool.invokeAll(taches);                 // tous, puis lire les futures
Resultat premier = pool.invokeAny(taches);                            // le premier fini (annule les autres)
pool.shutdown();
if (!pool.awaitTermination(10, TimeUnit.SECONDS)) { pool.shutdownNow(); }

// 3. Virtual threads (Java 21) : un thread par tache, pas de pool a dimensionner
try (var executor = java.util.concurrent.Executors.newVirtualThreadPerTaskExecutor()) {
    var futures = urls.stream()
            .map(u -> executor.submit(() -> charger(u)))
            .toList();
    for (var f : futures) { traiter(f.get()); }           // bloquer est gratuit : le thread est virtuel
}
// Structured concurrency (preview Java 21-24, en cours de stabilisation) : le groupe meurt ensemble
try (var scope = new java.util.concurrent.StructuredTaskScope.ShutdownOnFailure()) {
    var user = scope.fork(() -> trouverUser(id));
    var ordres = scope.fork(() -> trouverCommandes(id));
    scope.joinUntil(java.time.Instant.now().plusSeconds(2)).throwOnFailure();
    return new Vue(user.get(), ordres.get());
}
// Le pool reste utile quand : (a) le nombre de coeurs borne le CPU ; (b) tu dois limiter
// l'acces a une ressource rare (pool JDBC = 10 connexions) ; (c) code base sur Java 17.
\end{codejava}

\subsection{CompletableFuture~: composer des \'etapes asynchrones}

\begin{codejava}
import java.util.concurrent.CompletableFuture;
import static java.util.concurrent.CompletableFuture.*;

// Creer
CompletableFuture<Article> f1 = supplyAsync(() -> depot.trouver(id), pool);
CompletableFuture<Article> f2 = completedFuture(article);
CompletableFuture<Void> f3 = runAsync(() -> journaliser(), pool);

// Chainer (les variantes "Async" repartent sur un executor ; sans suffixe, elles restent
// sur le thread qui a termine l'etape precedente -- souvent ce qu'on veut)
var chaine = supplyAsync(() -> trouverClient(id), pool)
        .thenApply(client -> Facture.pour(client, panier))
        .thenCompose(facture -> envoyerAsync(facture))            // etape asynchrone suivante
        .thenAcceptAsync(envoyee -> log.info("facture {}", envoyee.numero()))
        .exceptionally(ex -> { log.error("echec", ex); return null; })   // recuperation
        .whenComplete((res, err) -> mesurer(res, err));                    // toujours execute

// Combiner
var lesDeux = thenCombine(fA, fB, (a, b) -> merge(a, b));
var soitLUn = anyOf(fA, fB, fC).thenApply(o -> (Article) o);
List<Article> tous = allOf(listeDeFutures.toArray(CompletableFuture[]::new))
        .thenApply(v -> listeDeFutures.stream().map(CompletableFuture::join).toList()).join();
// Aversion : ne JAMAIS laisser un CompletableFuture sans terminaison (ni exceptionally) ;
// ne pas bloquer (.join()) dans une chaine qui tourne sur le pool commun : tu cr\'ees un deadlock.

// Timeout (Java 9+)
Article a = f1.orTimeout(2, TimeUnit.SECONDS)
              .completeOnException(null, new ArticleParDefaut())   // repli gracieux
              .join();
// Helper que j'aime pour une liste de futures -> une future de liste :
static <T> CompletableFuture<List<T>> tous(Collection<CompletableFuture<T>> fs) {
    return CompletableFuture.allOf(fs.toArray(CompletableFuture[]::new))
            .thenApply(v -> fs.stream().map(CompletableFuture::join).toList());
}
\end{codejava}

\section{Synchronisation~: les outils, par taille de probl\`eme}

\begin{codejava}
// 1. synchronized : le verrou intrinseque (simple, correct, et suffisant 8 fois sur 10)
synchronized (this) {                 // verrou sur l'objet ; ou methode synchronized (eviter : trop large)
    if (!init) { charger(); init = true; }
}
static synchronized void globale() { }      // verrou sur la Class -- attention aux conflits mondes

// 2. Lock : quand tu as besoin de tryLock, d'un delai, de deux conditions, ou d'un verrou equitable
var lock = new java.util.concurrent.locks.ReentrantReadWriteLock(true);
lock.readLock().lock(); try { return lecture(); } finally { lock.readLock().unlock(); }
lock.writeLock().lock(); try { ecrire(); } finally { lock.writeLock().unlock(); }
var tryLock = new java.util.concurrent.locks.ReentrantLock();
if (tryLock.tryLock(50, TimeUnit.MILLISECONDS)) { try { critique(); } finally { tryLock.unlock(); } }

// 3. Sections critiques atomiques : le CAS et ses wrappers
var valeur = new java.util.concurrent.atomic.AtomicReference<>(initial);
valeur.updateAndGet(courante -> courante.fusionner(nouveau));     // CAS en boucle interne
var map = new java.util.concurrent.ConcurrentHashMap<Long, Compteur>();
map.computeIfAbsent(cle, k -> new Compteur()).incrementer();      // atomique, sans verrou global
map.merge(cle, 1L, Long::sum);

// 4. Communication entre threads : une file bloquante remplace 90 % des wait/notify
var tampon = new java.util.concurrent.ArrayBlockingQueue<Task>(64);   // BORNEE = backpressure
producteur: while (actif) { tampon.put(tache()); }                  // bloque si plein
consommateur: while (true) { var t = tampon.poll(1, TimeUnit.SECONDS); if (t == null) break; traiter(t); }
tampon.offer(t, 200, TimeUnit.MILLISECONDS);                        // refuse au lieu de bloquer a l'infini
var barriere = new java.util.concurrent.Semaphore(8);               // N acces simultanes (limite d'API)
barriere.acquire(); try { appelExterne(); } finally { barriere.release(); }
var latch = new java.util.concurrent.CountDownLatch(3);            // attendre 3 evenements
var cycle = new java.util.concurrent.CyclicBarrier(4);              // demarrage synchronise
// 5. wait/notify a l'ancienne : seulement pour comprendre le code des autres
synchronized (monitor) { while (!condition) { monitor.wait(500); } }   // le while est OBLIGATOIRE
synchronized (monitor) { condition = true; monitor.notifyAll(); }
\end{codejava}

\begin{tablez}[p{0.32\textwidth}X]
  \headrow \thd{Probl\`eme} \& \thd{Solution} \\
  \midrule
  compter / accumuler entre threads \& \verb|LongAdder|, \verb|AtomicInteger| \\
  cache partag\'e \& \verb|ConcurrentHashMap| + \verb|computeIfAbsent| \\
  limitation de concurrence (N appels simultan\'es) \& \verb|Semaphore| \\
  phase de pr\'eparation commune \& \verb|CyclicBarrier|, \verb|Phaser| \\
  attente d'un r\'esultat \& \verb|CompletableFuture|, \verb|ExecutorService.submit| \\
  producteur/consommateur \& \verb|BlockingQueue| born\'ee \\
  arr\^et coop\'eratif d'une t\^ache \& \verb|Thread.interrupt()| + test de l'interruption \\
  \og ~une seule fois par JVM~ \fg{} \& \verb|enum Holder|, \verb|ClassValue|, \verb|volatile|+double test \\
  t\^ache planifi\'ee \& \verb|ScheduledExecutorService| (et \verb|Executors.newScheduledThreadPool|) \\
  blocage mutuel \& \textbf{ordre global d'acquisition} des verrous, timeout, aucune entr\'ee utilisateur sous verrou \\
\end{tablez}

\section{Threads virtuels~: ce que \c{c}a change vraiment (Java 21)}

Un thread virtuel est une t\^ache l\'eg\`ere manag\'ee par la JVM, pos\'ee sur un petit nombre de threads
porteurs. Co\^ut de cr\'eation~: l'ordre du microseconde, m\'emoire~: quelques kilo-octets. On peut en
cr\'eer un million. Quand il bloque sur une E/S (socket, fichier, \verb|sleep|), il se d\'epose et
lib\`ere le porteur~: le d\'ebit d'une application \og ~une requ\^ete = beaucoup d'E/S~ \fg{} explose sans
changer une ligne de logique.

\begin{codejava}
// 1. Lancer une tache
Thread.ofVirtual().name("tache-", 0).start(() -> traiter());
// 2. Executor dedie (le motif d'une servlet "un thread par requete")
try (var executor = Executors.newVirtualThreadPerTaskExecutor()) {
    sources.forEach(url -> executor.submit(() -> { var html = HttpClient.newHttpClient()
            .send(HttpRequest.newBuilder(URI.create(url)).build(), BodyHandlers.ofString());
        return html.body().length(); }));
}    // le try-with-resources attend la fin de toutes les taches
// 3. Serveur HTTP : un thread virtuel par connexion -- la meme ecriture que le code bloquant
try (var serveur = HttpServer.create(new InetSocketAddress(8080), 0)) {
    serveur.setExecutor(Executors.newVirtualThreadPerTaskExecutor());
    serveur.createContext("/", e -> {
        try { e.sendResponseHeaders(200, 5); e.getResponseBody().write("hello".getBytes()); }
        catch (IOException ex) { throw new UncheckedIOException(ex); }
        finally { e.close(); }
    });
    serveur.start();
}
// 4. Fixed : les virtual threads n'ont PAS de pool ; si tu dois limiter (API externe, connexions
//    BDD), utilise un Semaphore -- pas newFixedThreadPool(N) sur des virtual.

// Ce que ca ne change pas :
//  - le code CPU-bound ne devient pas plus rapide (toujours nb coeurs) ;
//  - synchronized qui bloque longtemps epingle le thread porteur : remplace par ReentrantLock
//    dans les sections critiques longues (corrige en 24 via JEP 491) ;
//  - ca ne remplace pas CompletableFuture quand tu veux composer des resultats ;
//  - un thread local par tache legere = des millions d'objets : prefere les Scoped Values (JEP 446+).
Thread.sleep(Duration.ofSeconds(1));           // desormais aussi efficace sur un instant que sur un thread
\end{codejava}

\section{D\'eboguer et mesurer une application concurrente}

\begin{terminal}
jps -l                                  # lister les JVM en cours
jstack <pid> > threads.txt              # l'etat de TOUS les threads : cherche "BLOCKED", "deadlock"
jcmd <pid> Thread.print
jcmd <pid> VM.native_memory summary      # (avec -XX:NativeMemoryTracking=summary)
jstat -gcutil <pid> 1000                 # activite GC en direct
jmap -histo:live <pid> | head -20
# Outils graphiques : JDK Mission Control (gratuit), async-profiler (flammes), VisualVM
async-profiler -d 30 -f cpu.html <pid>
\end{terminal}

\etape{Recettes anti-race}
\begin{itemize}
  \item \textbf{Immutabilit\'e par d\'efaut}~: une donn\'ee qui ne change pas ne se synchronise pas.
  \item \textbf{Confine}~: un objet mutable appartient \`a un seul thread (ou une seule
        section critique explicite)~; passe des copies entre threads, pas des r\'ef\'erences.
  \item \textbf{Une seule voie de mutation}~: une file + un worker, ou un
        \verb|ConcurrentHashMap| + op\'erations atomiques -- jamais les deux m\'elang\'es.
  \item \textbf{Timeouts partout}~: un appel r\'eseau sans timeout est un thread bloqu\'e pour
        l'\'eternit\'e, donc un pool \'epuis\'e, donc un site mort.
  \item \textbf{Backpressure explicite}~: une file born\'ee qui refuse (429~HTTP, rejet, m\'etrique).
  \item \textbf{Test}~: \verb|ExecutorService| remplac\'e par un ex\'ecuteur direct dans les tests
        (d\'eterminisme)~; \link{jcstress}{https://github.com/openjdk/jcstress} pour les
        micro-tests de m\'emoire~; \link{ThreadSanitizer}{https://github.com/google/sanitizers} (natif
        en JDK~24+, \verb+-XX:+UseTSO+) ; \verb|-agentlib:jdwp| pour les points d'arr\^et multi-threads.
\end{itemize}

\begin{tablez}[p{0.34\textwidth}X]
  \headrow \thd{Sympt\^ome} \& \thd{Cause la plus probable} \\
  \midrule
  L'application se fige, CPU \`a 0 \& deadlock (ordre d'acquisition incoh\'erent) ou verrou jamais lib\'er\'e \\
  CPU \`a 100 \% sans travail \& boucle d'attente sans \verb|wait|/\verb|sleep| (spin) \\
  R\'esultats parfois faux \& champ partag\'e non volatile sans verrou~; \verb|HashMap| mut\'ee \`a deux threads \\
  Pool \'epuis\'e, timeouts en cha\^ine \& appels bloquants sans timeout~; ou \verb|join| dans le pool commun \\
  \verb|OutOfMemoryError: unable to create native thread| \& threads cr\'e\'es sans limite~: utilise un pool ou des virtuels \\
  Lent au d\'emarrage, rapide ensuite \& JIT froid~: profile apr\`es \'echauffement (JMH, \verb|-XX:TieredStopAtLevel|) \\
\end{tablez}

\begin{exercice}
\begin{enumerate}
  \item \'Ecris un tableau \verb|Rates| sur \verb|ConcurrentHashMap| qui compte des \'ev\'enements par
        minute~: lecture du taux glissant~; puis la version avec \verb|LongAdder|~; puis compare
        avec un \verb|synchronized| (10 threads, 10\textsuperscript{6} writes chacun).
  \item Construis un producteur/consommateur avec \verb|ArrayBlockingQueue|~: capacit\'e 100,
        arr\^et propre (poison pill ou interrupt), mesure du d\'ebit, et un indicateur de file pleine.
  \item T\'el\'echarge 1 000 URLs~: (a) s\'equentiel~; (b) pool fixe de 20~; (c) virtual threads
        illimit\'es~; (d) virtual threads + \verb|Semaphore(20)|. Rapporte les dur\'ees et explique
        chaque chiffre~: celui-ci est l'exercice qui fait comprendre la diff\'erence latence/d\'ebit.
  \item \'Ecris un cache TTL thread-safe~: \verb|computeIfAbsent|, \'eviction par access
        (\verb|LinkedHashMap| + \verb|removeEldestEntry|) ou horodatage~; 20 tests~: dont deux
        tests de concurrence (barrier + compteurs atomiques v\'erifiant qu'une valeur n'est calcul\'ee
        qu'une fois).
  \item Reproduis volontairement un deadlock \`a deux verrous~: observe-le avec \verb|jstack|, puis
        corrige-le par l'ordre global, puis par un \verb|tryLock| avec timeout~; garde la capture
        dans ton carnet.
  \item \'Ecris une impl\'ementation de \verb|RateLimiter| \og ~jeton~ \fg{} avec \verb|AtomicLong| et
        horodatage (pas de \verb|synchronized|), puis un test qui v\'erifie que 100 threads
        obtenants au plus 10 jetons/seconde pendant 5 secondes.
\end{enumerate}
\end{exercice}

\begin{ressources}
\begin{itemize}
  \item Tutorial officiel Concurrency --- \url{https://dev.java/learn/concurrency/}
  \item JEP 444 (virtual threads) --- \url{https://openjdk.org/jeps/444} --- Structured concurrency --- \url{https://openjdk.org/jeps/453}
  \item \emph{Java Concurrency in Practice} (Brian Goetz et al.) --- toujours le meilleur livre du domaine
  \item JCStress --- \url{https://github.com/openjdk/jcstress} --- Async-profiler --- \url{https://github.com/async-profiler/async-profiler}
  \item JDK Mission Control --- \url{https://www.oracle.com/java/technologies/jdk-mission-control.html}
\end{itemize}
\end{ressources}

\index{threads}\index{synchronized}\index{CompletableFuture}\index{virtual threads}
% ============================================================================
\chapter{Outillage~: Maven, Gradle, JDK multi-versions, CI}
\label{ch:java-build}
% ============================================================================

\begin{objectif}
\begin{itemize}
  \item Lire et \'ecrire un \fichier{pom.xml} et un \fichier{build.gradle.kts} sans les copier \`a l'aveugle.
  \item Comprendre le cycle de vie, la r\'esolution de d\'ependances, les p\'erim\`etres et les BOM.
  \item Utiliser le wrapper~: \verb|./mvnw| / \verb|./gradlew|--- le build fonctionne chez le voisin.
  \item Produire un ex\'ecutable~: jar, uber-jar, image native GraalVM, Docker.
  \item Formater, analyser et publier~: Spotless, Checkstyle, Error Prone, Release, CI GitHub Actions.
\end{itemize}
\end{objectif}

\section{Cr\'eer le projet}

\begin{terminal}
# Maven : archetypes (le plus simple pour un projet nu)
mvn archetype:generate -DgroupId=com.exemple -DartifactId=app \
    -DarchetypeArtifactId=maven-archetype-quickstart -DarchetypeVersion=1.4 -DinteractiveMode=false
# Gradle (apr\`es installation du CLI)
gradle init --type java-application --dsl kts --test-framework junit-jupiter
# Generateurs en ligne (le plus courant en pratique) :
#   Spring Initializr https://start.spring.io  |  https://start.quarkus.io  |  https://micronaut.io/launch
# Mise a jour des wrappers (a commiter !) :
mvn wrapper:wrapper -Dmaven=3.9.9
gradle wrapper --gradle-version 8.11.1
\end{terminal}

\subsection{La structure~: un contrat, pas une mode}

\begin{codeplain}
app/
\boxv \boxh \boxh mvnw, mvnw.cmd, .mvn/wrapper/          # le wrapper (toujours)
\boxv \boxh \boxh pom.xml            ou   build.gradle.kts + settings.gradle.kts
\boxv \boxh \boxh src/main/java/com/exemple/app/         # code
\boxv \boxh \boxh src/main/resources/                    # config, i18n, schemas SQL, templates
\boxv \boxh \boxh src/test/java/...                       # tests (miroir des paquets)
\boxv \boxh \boxh src/test/resources/                      # jeux de donnees de test
\boxv \boxh \boxh .editorconfig  .gitignore  README.md
\boxe \boxh \boxh (Gradle) build/    (Maven) target/       # JAMAIS versionnes
\end{codeplain}

\section{Maven~: le \texttt{pom.xml} utile}

\begin{codeconf}
<?xml version="1.0" encoding="UTF-8"?>
<project xmlns="http://maven.apache.org/POM/4.0.0"
         xmlns:xsi="http://www.w3.org/2001/XMLSchema-instance"
         xsi:schemaLocation="http://maven.apache.org/POM/4.0.0 https://maven.apache.org/xsd/maven-4.0.0.xsd">
  <modelVersion>4.0.0</modelVersion>

  <!-- 1. Identite : groupId = espace de noms, artifactId = nom, version = semver -->
  <groupId>com.exemple</groupId>
  <artifactId>blog-app</artifactId>
  <version>1.4.0-SNAPSHOT</version>          <!-- -SNAPSHOT = en developpement -->
  <packaging>jar</packaging>

  <!-- 2. Proprietes : encoding, version du langage, versions centralisees -->
  <properties>
    <project.build.sourceEncoding>UTF-8</project.build.sourceEncoding>
    <project.reporting.outputEncoding>UTF-8</project.reporting.outputEncoding>
    <maven.compiler.release>21</maven.compiler.release>   <!-- mieux que source/target -->
    <junit.version>5.11.4</junit.version>
  </properties>

  <!-- 3. BOM : une dependance "import" qui aligne les versions de toute une famille -->
  <dependencyManagement>
    <dependencies>
      <dependency>                            <!-- ne la mets pas dans <dependencies> -->
        <groupId>org.junit</groupId>
        <artifactId>junit-bom</artifactId>
        <version>${junit.version}</version>
        <type>pom</type>
        <scope>import</scope>
      </dependency>
      <dependency>
        <groupId>com.fasterxml.jackson</groupId>
        <artifactId>jackson-bom</artifactId>
        <version>2.18.2</version>
        <type>pom</type><scope>import</scope>
      </dependency>
    </dependencies>
  </dependencyManagement>

  <!-- 4. Depandances -->
  <dependencies>
    <dependency>
      <groupId>org.slf4j</groupId><artifactId>slf4j-simple</artifactId><version>2.0.16</version>
    </dependency>
    <dependency>
      <groupId>org.postgresql</groupId><artifactId>postgresql</artifactId><version>42.7.4</version>
      <scope>runtime</scope>                  <!-- necessaire a l'execution seulement -->
    </dependency>
    <dependency>
      <groupId>org.assertj</groupId><artifactId>assertj-core</artifactId><version>3.27.2</version>
      <scope>test</scope>
    </dependency>
    <dependency>
      <groupId>org.projectlombok</groupId><artifactId>lombok</artifactId><version>1.18.36</version>
      <scope>provided</scope>                 <!-- vu a la compilation, pas embarque -->
    </dependency>
    <dependency>                              <!-- exclusion : l'artefact "le plus faux" -->
      <groupId>commons-logging</groupId><artifactId>commons-logging</artifactId>
      <version>1.3.4</version>
      <exclusions><exclusion><groupId>log4j</groupId><artifactId>log4j</artifactId></exclusion></exclusions>
    </dependency>
  </dependencies>

  <!-- 5. Plugins : le cycle de vie est fait de goals branches sur des phases -->
  <build>
    <finalName>${project.artifactId}</finalName>
    <plugins>
      <plugin>
        <groupId>org.apache.maven.plugins</groupId><artifactId>maven-surefire-plugin</artifactId>
        <version>3.5.2</version>
        <configuration><argLine>-Xshare:off -XX:+EnableDynamicAgentLoading</argLine></configuration>
      </plugin>
      <plugin>
        <groupId>org.apache.maven.plugins</groupId><artifactId>maven-failsafe-plugin</artifactId>
        <version>3.5.2</version>
        <executions><execution><goals><goal>integration-test</goal><goal>verify</goal></goals></execution></executions>
      </plugin>
      <plugin>
        <groupId>org.apache.maven.plugins</groupId><artifactId>maven-shade-plugin</artifactId>
        <version>3.6.0</version>                <!-- uber-jar executable -->
        <executions><execution><phase>package</phase><goals><goal>shade</goal></goals>
          <configuration><transformers>
            <transformer implementation="org.apache.maven.plugins.shade.resource.ManifestResourceTransformer">
              <mainClass>com.exemple.app.App</mainClass>
            </transformer>
          </transformers></configuration>
        </execution></executions>
      </plugin>
      <plugin>                                 <!-- execution : mvn exec:java -->
        <groupId>org.codehaus.mojo</groupId><artifactId>exec-maven-plugin</artifactId><version>3.5.0</version>
        <configuration><mainClass>com.exemple.app.App</mainClass></configuration>
      </plugin>
      <plugin>                                 <!-- versionner proprement -->
        <groupId>org.apache.maven.plugins</groupId><artifactId>maven-release-plugin</artifactId><version>3.1.1</version>
      </plugin>
    </plugins>
  </build>

  <!-- 6. Profils : <profile> avec <id>dev/prod</id>, active par -Pprod ou par activation -->
</project>
\end{codeconf}

\begin{terminal}
mvn clean                # supprime target/
mvn compile              # src/main/java -> target/classes
mvn test                 # tests unitaires (surefire) : src/test/java
mvn package              # jar dans target/
mvn verify               # + tests d'integration (failsafe) + analyses
mvn install              # dans ~/.m2/repository (consomme par les autres projets)
mvn site                 # javadoc, rapports de qualite
mvn -Pprod -DskipTests package
# Diagnostic (les commandes qui te sauvent) :
mvn -X package                                   # log de debug
mvn dependency:tree -Dverbose -Dincludes=com.fasterxml.jackson*   # qui tire quoi
mvn dependency:analyze                            # inutilisees / undeclared
mvn versions:display-dependency-updates ; mvn versions:use-latest-releases
mvn help:effective-pom                            # le POM reel apres heritage et BOM
mvn enforcer:enforce                              # regles (version de Java, interdits)
\end{terminal}

\begin{tablez}[p{0.14\textwidth}p{0.3\textwidth}X]
  \headrow \thd{Scope} \& \thd{Pr\'esent o\`u ?} \& \thd{Effet sur ton artefact} \\
  \midrule
  \verb|compile| (d\'efaut) \& compile + tests + ex\'ecution \& il se retrouve dans le classpath de tes consommateurs : chaque \verb|compile| est une dette publique \\
  \verb|provided| \& compile + tests seulement \& le conteneur ou le JDK le fournit (servlet-api, annotation processors) \\
  \verb|runtime| \& tests + ex\'ecution \& parfait pour un driver JDBC ou une impl\'ementation de log \\
  \verb|test| \& tests seulement \& JUnit, Mockito, Testcontainers, assertions \\
  \verb|system| \& --- \& \`a fuir (jar pos\'e sur le disque) \\
\end{tablez}

\section{Gradle~: le build en Kotlin DSL}

\begin{codeconf}
// build.gradle.kts
plugins {
    java
    application
    `maven-publish`
    id("com.diffplug.spotless") version "6.25.0"
    id("org.graalvm.buildtools.native") version "0.10.4"
    checkstyle
}

group = "com.exemple"
version = "1.4.0"

java {
    toolchain { languageVersion = JavaLanguageVersion.of(21) }   // Gradle telecharge le JDK s'il manque
    withSourcesJar()
}

repositories {
    mavenCentral()
    mavenLocal()                       // pour tes artefacts SNAPSHOT locaux
}

dependencies {
    implementation(platform("org.junit:junit-bom:5.11.4"))
    implementation(platform("com.fasterxml.jackson:jackson-bom:2.18.2"))
    implementation("org.slf4j:slf4j-simple")
    implementation("org.postgresql:postgresql:42.7.4")
    compileOnly("org.projectlombok:lombok:1.18.36")
    annotationProcessor("org.projectlombok:lombok:1.18.36")
    testImplementation("org.junit.jupiter:junit-jupiter")
    testImplementation("org.assertj:assertj-core:3.27.2")
    testRuntimeOnly("org.junit.platform:junit-platform-launcher")
    testImplementation("org.testcontainers:postgresql")
}

application { mainClass.set("com.exemple.app.App") }

testing {
    suites {
        val test by getting(JvmTestSuite::class) { useJUnitJupiter() }
        val integrationTest by registering(JvmTestSuite::class) {
            dependencies { implementation(project()) ; implementation("org.testcontainers:postgresql") }
        }
    }
}

tasks.withType<Test>().configureEach {
    useJUnitPlatform()
    maxParallelForks = (Runtime.getRuntime().availableProcessors() / 2).coerceAtLeast(1)
    testLogging { events("failed", "skipped") }
    reports.html.required = false
}

tasks.withType<JavaCompile>().configureEach { options.encoding = "UTF-8"; options.compilerArgs += "-parameters" }

spotless { java { target("src/*/java/**/*.java"); googleJavaFormat("1.24.0"); removeUnusedImports() } }
check.dependsOn("spotlessCheck")
\end{codeconf}

\begin{codeconf}
// settings.gradle.kts -- multi-module + catalogue de versions centralise
rootProject.name = "blog"
include("app", "domaine", "infra", "api")
dependencyResolutionManagement {
    versionCatalogs {
        create("libs") {
            from(files("gradle/libs.versions.toml"))
        }
    }
}
\end{codeconf}

\begin{codeconf}
# gradle/libs.versions.toml -- le "catalogue" : une seule source de verite
[versions]
junit = "5.11.4"
jackson = "2.18.2"
[libraries]
junit-bom = { module = "org.junit:junit-bom", version.ref = "junit" }
jackson-databind = { module = "com.fasterxml.jackson.core:jackson-databind", version.ref = "jackson" }
[plugins]
spotless = { id = "com.diffplug.spotless", version = "6.25.0" }
\end{codeconf}

\begin{terminal}
./gradlew build                # compile + testes + assemble + verifs (le reflexe)
./gradlew run --args="--port 9090"
./gradlew test --tests "*DepotArticles*" --rerun-tasks
./gradlew --scan               # rapport de build en ligne (Gradle Enterprise/Develocity public)
./gradlew build --offline      # sans reseau (cache deja peuple)
./gradlew build --stacktrace --info
./gradlew dependencies --configuration runtimeClasspath | head -60
./gradlew buildEnvironment ; ./gradlew properties | grep version
./gradlew clean build -x test -PskipChecks
gradle wrapper --gradle-version 8.11.1 --distribution-type bin   # et commite gradle/wrapper/*
# Config globale : ~/.gradle/gradle.properties
org.gradle.jvmargs=-Xmx2g -XX:MaxMetaspaceSize=512m
org.gradle.parallel=true ; org.gradle.caching=true ; org.gradle.configuration-cache=true
org.gradle.daemon=true ; org.gradle.console=rich ; org.gradle.workers.max=8
\end{terminal}

\subsection{Maven ou Gradle~?}

\begin{tablez}[p{0.2\textwidth}p{0.28\textwidth}X]
  \headrow Crit\`ere \& Maven \& Gradle \\
  \midrule
  Courbe d'apprentissage \& douce, convention forte, XML verbeux mais lisible \& plus exigeante (DSL, t\^aches), mais expressive \\
  Vitesse \& correcte, r\'esolutions s\'equentielles \& incr\'ementale + cache + parall\`ele~: 2--10\,\$\times\$ sur les gros builds \\
  Multi-module \& \verb|<modules>| simple \& excellent, graphes de t\^aches fins \\
  Plugins \& \'enorme d\'ep\^ot, versions fig\'ees par super-POM \& \'ecosyst\`eme riche, API de plugin en JVM/Kotlin \\
  Reprisibilit\'e \& \'elev\'ee (XML d\'eclaratif) \& \'elev\'ee \emph{si} le script est disciplin\'e \\
  IDE \& IntelliJ natif, VS~Code correct \& IntelliJ/Android natifs, import parfois capricieux \\
\end{tablez}

\begin{note}
Choisis Gradle pour~: un monorepo, Android, des builds lents, des besoins sur-mesure (g\'en\'eration de
code, Docker, graal). Choisis Maven pour~: un projet serveur classique, une \'equipe h\'et\'erog\`ene, une
entreprise qui centralise d\'ej\`a un d\'ep\^ot Nexus. Les deux se valent techniquement~; la pire option est
le \og ~on a mis les deux~ \fg{}.
\end{note}

\section{Produire un ex\'ecutable~: les quatre formats}

\begin{terminal}
# 1. jar simple (librairie) : mvn package -> target/app.jar ; java -cp app.jar:lib/* com.exemple.App
# 2. uber-jar (tout embarque) -- Maven shade/assembly, Gradle shadow ; attention aux META-INF/services
java -jar app.jar
# 3. distribution avec launcher + classpath (recommande pour un serveur) :
mvn -Pdist package         # ou Gradle : installDist -> build/install/app/bin/app
# 4. image native GraalVM (d\'emarrage ~10 ms, ~40 Mo de RSS, mais build long et r\'eflexion \`a configurer)
nativeCompile          # Gradle : ./gradlew nativeCompile
mvn -Pnative native:compile
native-image -jar app.jar -H:Name=app -march=native --initialize-at-build-time
# 5. Conteneur (la vraie cible en prod) : un Dockerfile multi-etapes
\end{terminal}

\begin{codeconf}
# Dockerfile -- Maven (adapte pour Gradle avec gradle build --no-daemon nativeCompile)
FROM maven:3.9-eclipse-temurin-21 AS build
WORKDIR /m
COPY . .
RUN mvn -B -DskipTests package
FROM eclipse-temurin:21-jre-jammy
RUN useradd -r -u 10001 app && mkdir -p /data && chown app /data
USER app
WORKDIR /app
COPY --from=build /m/target/app.jar .
EXPOSE 8080
HEALTHCHECK --interval=10s --timeout=3s CMD ["sh","-c","wget -qO- http://localhost:8080/health || exit 1"]
ENTRYPOINT ["java","-XX:MaxRAMPercentage=75","-XX:+ExitOnOutOfMemoryError","-jar","/app/app.jar"]
\end{codeconf}

\begin{codeconf}
# .mvn/jvm.config (options de build) ; .mvn/maven.config : -T 1C -Dstyle.color=always
# settings.xml (~/.m2/settings.xml) : miroir, credentials, profils
<settings>
  <mirrors><mirror>
    <id>central-miroir</id><mirrorOf>central</mirrorOf>
    <url>https://maven.aliyun.com/repository/central</url>   <!-- ou un Nexus interne -->
  </mirror></mirrors>
  <servers><server><id>repo-interne</id><username>${env.NEXUS_USER}</username>
           <password>${env.NEXUS_PASS}</password></server></servers>
</settings>
\end{codeconf}

\section{Qualit\'e du build~: la cha\^ine qui rend le code propre tout seul}

\begin{terminal}
# Formatage (le m\^eme partout, sans discussion en revue)
./gradlew spotlessApply            # ou : mvn spotless:apply
mvn com.diffplug.spotless:spotless-maven-plugin:apply
# Analyse statique
mvn -Panalysis verify              # checkstyle + spotbugs + pmd + dependency-check
./gradlew checkstyleMain spotbugsMain pmdMain
mvn org.owasp:dependency-check-maven:check -DfailBuildOnCVSS=7
# Reproduire et debugger le build
mvn -Dmaven.repo.local=/tmp/m2 -o package      # depot local isole / hors ligne
\end{terminal}

\begin{codeconf}
<!-- Checkstyle + PMD + SpotBugs + JaCoCo : les reglages qui rendent le build utile -->
<plugin>
  <groupId>org.apache.maven.plugins</groupId><artifactId>maven-checkstyle-plugin</artifactId><version>3.6.0</version>
  <configuration>
    <configLocation>google_checks.xml</configLocation>       <!-- ou sun_checks.xml, ou le tien -->
    <consoleOutput>true</consoleOutput><failOnViolation>true</failOnViolation>
    <violationSeverity>warning</violationSeverity>
  </configuration>
</plugin>
<plugin>
  <groupId>com.github.spotbugs</groupId><artifactId>spotbugs-maven-plugin</artifactId><version>4.8.3.1</version>
  <configuration><effort>Max</effort><threshold>Medium</threshold><xmlOutput>true</xmlOutput></configuration>
</plugin>
<plugin>
  <groupId>org.jacoco</groupId><artifactId>jacoco-maven-plugin</artifactId><version>0.8.12</version>
  <executions>
    <execution><goals><goal>prepare-agent</goal></goals></execution>
    <execution><id>rapport</id><phase>test</phase><goals><goal>report</goal></goals></execution>
    <execution><id>seuil</id><phase>verify</phase><goals><goal>check</goal></goals>
      <configuration><rules><rule><limits><limit><counter>LINE</counter><minimum>0.80</minimum></limit></limits></rule></configuration>
    </execution>
  </executions>
</plugin>
\end{codeconf}

\section{CI GitHub Actions~: le mod\`ele Java}

\begin{codeyml}
name: build
on:
  push: { branches: [main] }
  pull_request:
jobs:
  java:
    runs-on: ubuntu-latest
    steps:
      - uses: actions/checkout@v4
      - uses: actions/setup-java@v4
        with:
          distribution: temurin
          java-version: 21
          cache: maven            # ou: gradle  (le cache fait gagner 2 a 4 min)
      - name: Build + tests + verification
        run: mvn -B -ntp verify                # ou: ./gradlew build --no-daemon
      - name: Rapport de couverture (JaCoCo)
        run: mvn -B -ntp jacoco:report site -DskipTests
        if: always()
      - uses: actions/upload-artifact@v4
        if: always()
        with:
          name: rapports-${{ matrix.os ?? 'linux' }}
          path: |
            target/site
            target/surefire-reports
            build/reports
  matrice:
    runs-on: ${{ matrix.os }}
    strategy:
      fail-fast: false
      matrix:
        os: [ubuntu-latest, macos-latest, windows-latest]
        java: [17, 21]
    steps:
      - uses: actions/checkout@v4
      - uses: actions/setup-java@v4
        with: { distribution: temurin, java-version: "${{ matrix.java }}" }
      - run: mvn -B -ntp verify
\end{codeyml}

\etape{Publication}
\begin{itemize}
  \item En interne~: \verb|mvn deploy| / \verb|publish| vers un Nexus, Artifactory, GitHub Packages,
        \verb|packagecloud|--- ou un d\'ep\^ot \verb|file:///home/deploy/m2| si tu es seul.
  \item Sur Maven Central~: compte \link{Sonatype Central Portal}{https://central.sonatype.com/},
        groupes valid\'es, \verb|gpg|-signature, javadoc + sources obligatoires,
        \link{central-publishing-maven-plugin}{https://github.com/central-publishing/central-publishing-maven-plugin}
        ou \verb|./gradlew publishToMavenLocal| puis le portal web.
  \item Versionnement~: \verb|x.y.z-SNAPSHOT| en dev~; \verb|mvn release:prepare release:perform|
        (tag \verb|v1.4.0|, passage en \verb|1.4.1-SNAPSHOT|)~; ou \link{Revapi}{https://github.com/revapi/revapi}
        pour refuser une rupture de compatibilit\'e binaire dans une \verb|minor|.
\end{itemize}

\begin{piege}
Cinq erreurs de build classiques~: (1) le wrapper non committ\'e~: chacun a une version de Maven et
les builds diff\`erent~; (2) \verb|maven.compiler.source| sans \verb|release|~: tu compiles avec le
JDK~21 en visant~17 et tu utilises des API de~21~; (3) \verb|<scope>compile| partout~: tes
d\'ependances deviennent publiques et cassent les consommateurs~; (4) deux versions de Jackson
dans l'arbre~: \verb|NoClassDefFoundError| al\'eatoire~; (5) le cache CI sans \verb|~/.m2/repository|.
\end{piege}

\begin{exercice}
\begin{enumerate}
  \item Prends le projet du chapitre pr\'ec\'edent (ou un carnet d'adresses maison) et fais-le passer
        de \og ~\verb|javac| dans un terminal~ \fg{} \`a un projet Maven complet~: BOM, tests, javadoc,
        surefire, failsafe, uber-jar \verb|package|, README avec captures.
  \item Ajoute la m\^eme chose c\^ot\'e Gradle~: toolchain, catalogue de versions, t\^ache
        \verb|int\'egrationTest|~; garde les deux builds un moment et tranche avec des chiffres
        (temps de \verb|clean build| \`a froid et \`a chaud).
  \item Fais \'echouer un build volontairement et r\'epare-le~: d\'ependance en double avec versions
        incompatibles, \verb|NoSuchMethodError|~: puis \verb|dependency:tree|~; un \verb|pom.xml|
        sans UTF-8~: les accents du fichier de propri\'et\'es deviennent illisibles.
  \item Ajoute la cha\^ine qualit\'e~: Spotless (google-java-format), Checkstyle avec une config
        maison de 10 r\`egles, JaCoCo avec seuil bloquant \`a 80 \%, OWASP dependency-check.
        Fais en sorte que \verb|mvn verify| soit ton r\'eflexe local avant chaque push.
  \item \'Ecris le \verb|Dockerfile| multi-\'etapes~: puis une image \verb|-distroless| ou
        \verb|jre-alpine|~: compare taille et temps de d\'emarrage~; ajoute un \verb|HEALTHCHECK|
        et fais tourner le conteneur avec \verb|-XX:MaxRAMPercentage|.
  \item Compile en natif (GraalVM)~: mesure le temps de d\'emarrage et l'empreinte m\'emoire face au
        jar~; documente les r\'eflexions n\'ecessaires (\verb|reflect-config.json|) et ce que tu as
        perdu (diagnostic \`a chaud, JIT qui s'\'echauffe, \verb|jcmd|).
\end{enumerate}
\end{exercice}

\begin{ressources}
\begin{itemize}
  \item Maven --- \url{https://maven.apache.org/guides/} --- Reference Model (le POM, par c\oe{}ur en 2 h)
  \item Gradle --- \url{https://docs.gradle.org/current/userguide/userguide.html} --- Kotlin DSL --- \url{https://docs.gradle.org/current/userguide/kotlin\\\_dsl.html}
  \item \link{Maven Central Search}{https://central.sonatype.com/search} (trouver l'\verb|artifactId| exact)
  \item \link{mvnrepository}{https://mvnrepository.com/} (utile, mais la source de v\'erit\'e reste Central)
  \item GraalVM Native Image --- \url{https://www.graalvm.org/latest/reference-manual/native-image/}
  \item \link{Setup Java action}{https://github.com/actions/setup-java} --- \link{Spotless}{https://github.com/diffplug/spotless}
\end{itemize}
\end{ressources}

\index{Maven}\index{Gradle}\index{wrapper}\index{uber-jar}\index{GraalVM}\index{CI Java}
% ============================================================================
\chapter{Tests et qualit\'e en Java}
\label{ch:java-tests}
% ============================================================================

\begin{objectif}
\begin{itemize}
  \item Structurer une suite JUnit~5~: cycles de vie, assertions expressives, tests param\'etr\'es.
  \item Tester ce qui compte~: unit\'es avec Mockito, int\'egration avec Testcontainers, HTTP avec
        RestAssured/MockWebServer.
  \item Faire respecter l'architecture avec \texttt{ArchUnit}, et la couverture avec JaCoCo.
  \item Mesurer ce qui doit l'\^etre avec JMH, et la robustesse r\'eelle avec PIT (mutation testing).
  \item Rendre les tests rapides et non flaky~: c'est la condition de survie de la suite.
\end{itemize}
\end{objectif}

\section{JUnit~5~: ce qu'il faut conna\^itre}

\begin{codejava}
import org.junit.jupiter.api.*;
import org.junit.jupiter.params.ParameterizedTest;
import org.junit.jupiter.params.provider.*;
import static org.junit.jupiter.api.Assertions.*;
import static org.assertj.core.api.Assertions.*;      // la vraie API d'assertion (voir plus bas)

@DisplayNameGeneration(DisplayNameGenerator.ReplaceUnderscores.class)
@Tag("rapide")
class CalculateurDeTests {

    @BeforeAll static void uneFoisPourLaClasse() { }        // doit etre static (sinon @TestInstance(PER_CLASS))
    @BeforeEach void avantChaqueTest(TestInfo info) { }
    @AfterEach void apres() { }
    @AfterAll static void fin() { }

    @Test void publier_un_brouillon_fixe_la_date() {
        var article = new Article("titre", "corps", Statut.BROUILLON);
        var publie = article.publier(Instant.parse("2026-02-01T10:00:00Z"));
        assertThat(publie.statut()).isEqualTo(Statut.PUBLIE);
        assertThat(publie.publieLe()).isEqualTo(Instant.parse("2026-02-01T10:00:00Z"));
        assertThat(article).isSameAs(article).isNotSameAs(publie);            // immuabilite verifiee
        assertAll(
            () -> assertThat(publie.titre()).isEqualTo("titre"),
            () -> assertThat(publie.tags()).isEmpty(),
            () -> assertThatCode(() -> publie.publier(null)).doesNotThrowAnyException()
        );
    }

    @Test void publier_un_article_archive_refuse() {
        var a = new Article("t", "c", Statut.ARCHIVE);
        assertThatThrownBy(() -> a.publier(Instant.now()))
                .isInstanceOf(IllegalStateException.class)
                .hasMessageContaining("archive");
        // equivalents JUnit : assertThrows / assertTimeoutPreemptively / assertLinesMatch
    }

    @Test @Timeout(value = 200, unit = java.util.concurrent.TimeUnit.MILLISECONDS)
    void ne_pas_bloquer() { }
    @Test @Disabled("attend la correction du lotisseur #142")
    void en_attente() { }
    @Test @RepeatedTest(value = 5, name = RepeatedTest.CURRENT_REPETITION + "/" + RepeatedTest.TOTAL_REPETITIONS)
    @Execution(ExecutionMode.CONCURRENT)                    // avec @ResourceLock pour les ressources partagees
    void idempotence() { }
    @Test @EnabledOnJre({JRE_VERSION_17, JRE_VERSION_21}) @EnabledIfEnvironmentVariable(named = "CI", matches = "true")
    void conditionnel() { }

    // Les tests parametres : la moitie de la valeur de JUnit 5
    @ParameterizedTest(name = "[{index}] {0} -> {1}")
    @CsvSource({ "0, zero", "1, un", "42, quarante-deux", "-3, moins trois" })
    void convertir(int n, String attendu) { assertThat(Nombres.enToutesLettres(n)).isEqualTo(attendu); }

    @ParameterizedTest @MethodSource("casLimites")
    void refuser(String titre, Class<? extends Exception> ex) {
        assertThatThrownBy(() -> Article.de("titre", titre)).isInstanceOf(ex);
    }
    static Stream<Arguments> casLimites() {
        return Stream.of(arguments("", IllegalArgumentException.class),
                          arguments("  ", IllegalArgumentException.class),
                          arguments("x".repeat(161), IllegalArgumentException.class));
    }
    @ParameterizedTest @EnumSource(Statut.class)
    void chaque_statut_a_un_libelle(Statut s) { assertThat(s.libelle()).isNotBlank(); }
    @ParameterizedTest @ValueSource(ints = {2, 3, 5, 7, 11}) @NullSource @EmptySource
    void nombres_premiers(int n) { }

    // voir Mockito plus bas pour @Mock/@InjectMocks
    @TestFactory @DisplayName("test dynamique") DynamicNode[] usine() {
        // Des noeuds generes a l'execution (ex. un test par fichier de fixtures, par ligne d'un CSV)
        return java.nio.file.Files.list(dossier).map(p ->
                DynamicTest.dynamicTest(p.toString(), () -> verifier(p))).toArray(DynamicNode[]::new);
    }
    @Test void message_daide_clair() {
        var map = new java.util.HashMap<String, String>();
        assertNotNull(map.get("absent"), () -> "touche 'absent' manquante dans " + map.keySet());
        // Le fournisseur de message n'est evalue qu'en cas d'echec : pense a la forme lambda.
    }
}
\end{codejava}

\subsection{AssertJ~: l'API qui rend les tests lisibles}

\begin{codejava}
assertThat(liste).hasSize(3).contains("a", "b").doesNotContainNull().extracting(Article::titre)
        .containsExactly("un", "deux").isSortedAccordingTo(comparator);
assertThat(map).containsOnlyKeys("a", "b").doesNotContainValue(null)
        .allSatisfy((k, v) -> assertThat(v).isPositive());
assertThat(string).isEqualToIgnoringCase("OK").startsWith("O").matches("(?i)ok")
        .normalizer(Normalizer.normalize()).isEqualTo("ca");
assertThat(instant).isBetween(debut, fin).hasSameHourAs(attendu).isCloseTo(attendu, within(1, ChronoUnit.SECONDS));
assertThat(code).hasFieldOrProperty("id").isNotSameAs(other);
assertThat(optionnel).isPresent().get().isEqualTo(attendu);
assertThat(uneFuture).succeedsWithin(Duration.ofSeconds(2)).isEqualTo(attendu);
assertThatThrownBy(...).isInstanceOf(RuntimeException.class).hasRootCauseInstanceOf(SQLException.class)
        .hasMessageContaining("cle").hasFieldOrPropertyWithValue("code", "validation");
// Le livre des assertions completes : org.assertj.core.api / .api.soft (SoftAssertions pour
// cumuler les echecs dans un test de parcours de formulaire).
\end{codejava}

\section{Mockito~: simuler les fronti\`eres, pas le syst\`eme}

\begin{codejava}
@ExtendWith(MockitoExtension.class)
@MockitoSettings(strictness = Strictness.STRICT_STUBS)      // le mode qui evite les faux positifs
class ServiceArticlesTest {

    @Mock private DepotArticles depot;
    @Mock private Journal journal;
    @Captor private ArgumentCaptor<Article> capteur;
    @InjectMocks private ServiceArticles service;            // injection par constructeur

    @Test void publier_delegue_et_notifie() {
        when(depot.trouver(1L)).thenReturn(Optional.of(unArticle(Statut.BROUILLON)));
        when(depot.sauver(any())).thenAnswer(i -> i.getArgument(0));
        // ou, pour un mock "bete" : lenient().when(...) -- mais ca cache une stub inutile

        var resultat = service.publier(1L, Instant.parse("2026-02-01T00:00:00Z"));

        assertThat(resultat.statut()).isEqualTo(Statut.PUBLIE);
        verify(depot).sauver(capteur.capture());
        assertThat(capteur.getValue().publieLe()).isEqualTo(Instant.parse("2026-02-01T00:00:00Z"));
        verify(journal).info(eq("article publie {}"), eq(1L));
        verifyNoMoreInteractions(depot);
    }

    @Test void introuvable_remonte_le_bon_code() {
        when(depot.trouver(9L)).thenReturn(Optional.empty());
        assertThatThrownBy(() -> service.publier(9L, Instant.now()))
                .isInstanceOf(ErreurIntrouvable.class);
        verify(depot, never()).sauver(any());
    }
    // Spy = objet reel avec methodes detournees : a eviter (test fragile).
    // Fais une fausse implementation explicite quand le comportement compte (un FakeDepot en 20 lignes) ;
    // un mock ne sert qu'a verifier une interaction.
}
// Mockito + inline mock maker (depasse le final) :
// org.mockito:mockito-core + argLine "-javaagent:..." pour Java 21+ (avoid the warning).
// Reseaux : MockWebServer (OkHttp) ou WireMock -- voir le chapitre API.
// Horloge : ne mocke pas Instant.now() : passe une Horloge/Supplier<Instant> (voir chapitre precedent).
\end{codejava}

\section{Int\'egration~: Testcontainers et la vraie base}

\begin{codejava}
import org.testcontainers.containers.PostgreSQLContainer;
import org.testcontainers.junit.jupiter.*;

@Testcontainers
class DepotArticlesIT {
    @Container
    static final PostgreSQLContainer<?> PG = new PostgreSQLContainer<>("postgres:17-alpine")
            .withDatabaseName("test").withUsername("test").withPassword("test")
            .withInitScript("db/schema.sql")        // puis, si besoin : "db/seed.sql"
            .withReuse(true);                       // + testcontainers.reuse.enable=true dans ~/.testcontainers.properties
    private DepotArticles depot;
    @BeforeEach void monter() {
        var cfg = new com.zaxxer.hikari.HikariConfig();
        cfg.setJdbcUrl(PG.getJdbcUrl()); cfg.setUsername(PG.getUsername()); cfg.setPassword(PG.getPassword());
        cfg.setMaximumPoolSize(4);
        var ds = new com.zaxxer.hikari.HikariDataSource(cfg);
        depot = new DepotArticlesJdbc(ds);
    }
    @Test void inserer_et_relire_avec_tags() throws Exception {
        var a = depot.sauver(unArticle(Statut.PUBLIE).avecTag("java"));
        assertThat(depot.trouver(a.id())).orElseThrow().tags().contains("java");
        assertThat(depot.lister(new Pageable(0, 10))).hasSize(1);
    }
}
// Les conteneurs demarrent en ~2 s (image deja presente) : une suite IT de 40 tests en < 1 min.
// @Testcontainers(disabledWithoutDocker = true) pour ne pas casser le build d'un pote sans Docker.
// Autres images qui changent la vie : "redis:7-alpine", "minio/minio", "hashicorp/vault",
// "localstack/localstack", "elasticsearch:8", "maildev/maildev" (SMTP + HTTP, parfait pour les tests d'email).
\end{codejava}

\begin{codejava}
// Test HTTP de bout en bout, sans serveur : RestAssured ou java.net.http + MockWebServer
import static io.restassured.RestAssured.*;
given().header("Authorization", "Bearer " + jeton).queryParam("page", 0)
   .when().get("/api/articles")
   .then().statusCode(200).contentType("application/json")
   .body("total", equalTo(2)).body("contenus[0].titre", equalTo("Bonjour"))
   .time(lessThan(1500L));                                   // une assertion de latence : oui, ca se fait
given().body("""
    {"titre":"x"}""").contentType("application/json")
   .when().post("/api/articles")
   .then().statusCode(400).body("type", equalTo("https://exemple.fr/erreurs/validation"));
\end{codejava}

\section{Architecture, couverture, mutations~: les trois filets}

\begin{codejava}
// 1. ArchUnit : les regles que la revue de code oublie, le test ne les oublie pas
@AnalyzeClasses(packages = "com.exemple", importOptions = ImportOption.DoNotIncludeTests.class)
class ArchitectureTest {
    @ArchTest static final ArchRule couches = layeredArchitecture()
            .consideringOnlyDependenciesInLayers()
            .layer("Domaine").definedBy("..domaine..")
            .layer("Application").definedBy("..application..")
            .layer("Infra").definedBy("..infra..")
            .whereLayer("Domaine").mayNotBeAccessedByAnyLayerExcept("Application")
            .whereLayer("Infra").mayOnlyBeAccessedByLayers("Amoracage");
    @ArchTest static final ArchRule domaine_pur = noClasses().that().resideInAPackage("..domaine..")
            .should().dependOnClassesThat().resideOutsideOfPackages("..domaine..", "java..", "..record..");
    @ArchTest static final ArchRule pas_de_field_injection =
            noFields().should().beAnnotatedWith(Autowired.class);
    @ArchTest static final ArchRule services_finaux = classes().that().resideInAPackage("..application..")
            .and().haveSimpleNameEndingWith("Service").should().beFinal();
}
// 2. JaCoCo : la couverture est un signal d'alerte, pas un objectif.
//    "80 % de lignes" ne dit rien sur les chemins d'erreur. Regarde les BRANCHES, et les paquets
//    les moins couverts : ce sont eux, pas la moyenne, qui predisent les bugs.
// 3. PIT (mutation testing) : le seul indicateur qui repond a "mes tests detecteraient-ils
//    ce bug ?". Il injecte des fautes (>= remplace par >, constantes inversees) et exige
//    qu'au moins un test echoue. Un score de mutations de 70 % vaut mieux que 95 % de lignes.
\end{codejava}

\begin{terminal}
# PIT avec Maven (l'edition la plus simple)
mvn org.pitest:pitest-maven:mutationCoverage -DtargetClasses=com.exemple.domaine.* -DtimeoutConstant=4000
# PIT avec Gradle : plugin "info.solidsoft.pitest"
./gradlew pitest
# JMH (micro-benchmarks) : ne jamais mesurer avec System.nanoTime() a la main dans un main()
./gradlew jmh ; mvn -Pbenchmark generate-sources  # ou creer un module benchmarks/
\end{terminal}

\begin{codejava}
// Un JMH honnete, en 25 lignes (org.openjdk.jmh:jmh-core + generator)
@BenchmarkMode(Mode.AverageTime) @OutputTimeUnit(TimeUnit.NANOSECONDS)
@Warmup(iterations = 5, time = 1) @Measurement(iterations = 5, time = 1) @Fork(1) @State(Scope.Benchmark)
public class ParsingBench {
    @Param({"100", "10000"}) private int n;
    private List<String> entrees;
    @Setup public void monter() { entrees = genere(n); }
    @Benchmark public long avecStream() { return entrees.stream().mapToLong(String::length).sum(); }
    @Benchmark public long avecBoucle() { long s = 0; for (var e : entrees) { s += e.length(); } return s; }
    @Benchmark public void blackhole(Blackhole bh) { for (var e : entrees) { bh.consume(e); } }
}
// Sans @Setup/@Param/Blackhole, la mesure est du bruit ; sans @Fork, c'est faux.
\end{codejava}

\section{Anti-flaky~: les r\`egles de survie}

\begin{enumerate}
  \item Aucun \verb|Thread.sleep|~: attends une \emph{condition} avec \link{Awaitility}{https://github.com/awaitility/awaitility}
        (\verb|await().atMost(5s).untilAsserted(...)|) ou un \verb|CountDownLatch|.
  \item Horloge et al\'ea inject\'es~: \verb|Horloge| (\verb|Clock.fixed(...)|), \verb|Random(seed)| ou
        \verb|RandomGenerator.JdkRandom(1)|. Deux ex\'ecutions, m\^emes r\'esultats.
  \item Aucune d\'ependance d'ordre entre tests~: un test pris isol\'e doit passer
        (\verb|-Djunit.jupiter.testmethod.order.default=...Random| pour v\'erifier).
  \item R\'eseau et disques isol\'es~: \verb|@TempDir| pour les fichiers~; MockWebServer pour l'HTTP~;
        un conteneur pour la BDD. Jamais l'instance de dev d'un coll\`egue.
  \item Une donn\'ee de test par test (builder ou m\'ethode \verb|unArticle()|)~: pas de fixture
        globale mut\'ee en s\'erie.
  \item Locale et fuseau fix\'es~: \verb|@DefaultLocale("fr-FR")|, \verb|@DefaultTimeZone("UTC")|
        (JUnit 5.12+) ou \verb|TimeZone.setDefault| dans un extension.
  \item Un test qui \'echoue une fois sur 50 est un bug~: traite-le comme un P1, sinon il d\'evore
        la confiance de toute l'\'equipe en trois semaines.
\end{enumerate}

\begin{tablez}[p{0.24\textwidth}X]
  \headrow \thd{Ce que tu testes} \& \thd{Outil adapt\'e} \\
  \midrule
  une fonction pure (calcul, parsing, r\`egle m\'etier) \& JUnit + AssertJ, en table de v\'erit\'e param\'etr\'ee \\
  une collaboration~: \og ~le service appelle le d\'ep\^ot avec X~ \fg{} \& Mockito \emph{ou} un fake explicite \\
  SQL~: requ\^etes, contraintes, migrations, plans \& Testcontainers + un \verb|@Sql|/script d'init \\
  S\'erialisation~: JSON entr\'ee/sortie, compatibilit\'e \& test sur un fichier d'or (\verb|.json| en resources) \\
  HTTP~: une API, des en-t\^etes, des codes de statut \& RestAssured / MockMvc / WebTestClient \\
  Une application compl\`ete qui d\'emarre \& \verb|@SpringBootTest| (ou \verb|quarkus test dev-services|) \\
  Interface graphique/console \& le plus souvent~: ne pas tester l'UI, tester la couche en dessous (Playwright pour le web, tests de rendu pour Swing/JavaFX) \\
  Performance d'un algorithme \& JMH~; \textbf{jamais} dans la CI bloquante~: une alarme s\'epar\'ee \\
\end{tablez}

\begin{exercice}
\begin{enumerate}
  \item \'Ecris la suite de tests du \verb|CalculateurDe| du chapitre sur les collections~: table de
        v\'erit\'e \verb|@CsvSource|, un test d'erreur par invariante, un test de propri\'et\'e avec 1 000
        entr\'ees al\'eatoires (graine fixe).
  \item Couvre le \verb|LRUCache| de l'exercice du chapitre pr\'ec\'edent~: ajoute un test de
        concurrence (\verb|CyclicBarrier| + \verb|ExecutorService|) qui prouve que deux threads
        ne calculent pas deux fois la m\^eme cl\'e.
  \item Fais passer la couverture de branches d'un module de 40 \% \`a 80 \%~: montre le rapport
        JaCoCo avant/apr\`es~; puis lance PIT~: le score de mutations t'apprendra quelles assertions
        \'etaient d\'ecoratives.
  \item Ajoute un test d'int\'egration Testcontainers sur le d\'ep\^ot JDBC~: migration avec Flyway/Liquibase
        dans le conteneur, jeu de donn\'ees~; puis un test qui v\'erifie le plan d'ex\'ecution
        (\verb|EXPLAIN|) d'une requ\^ete~: une ligne sur la pr\'esence de l'index.
  \item \'Ecris un test d'architecture~: le paquet \verb|domaine| ne d\'epend ni de Spring, ni de JDBC,
        ni de Jackson~; aucune classe publique hors \verb|api|~; toutes les exceptions h\'eritent de
        \verb|ErreurDomaine|. Fais-le \'echouer volontairement, puis corrige le code~: c'est le but.
  \item Prends un benchmark JMH qui mesure deux impl\'ementations d'un m\^eme algorithme~:
        \'ecris la ligne de commande, la description, et ce que tu conclurais si l'\'ecart \'etait de
        3 \%~; puis fais la m\^eme mesure \`a froid et note la diff\'erence~: tu viens d'apprendre \`a ne
        pas mentir avec un micro-benchmark.
\end{enumerate}
\end{exercice}

\begin{ressources}
\begin{itemize}
  \item JUnit~5 User Guide --- \url{https://junit.org/junit5/docs/current/user-guide/}
  \item AssertJ --- \url{https://assertj.github.io/doc/} --- Mockito --- \url{https://javadoc.io/doc/org.mockito/mockito-core/latest/org/mockito/Mockito.html}
  \item Testcontainers --- \url{https://java.testcontainers.org/} (les modules officiels par image)
  \item ArchUnit --- \url{https://www.archunit.org/} --- PIT --- \url{https://pitest.org/} --- JMH --- \url{https://github.com/openjdk/jmh}
  \item Awaitility --- \url{https://github.com/awaitility/awaitility} --- RestAssured --- \url{https://restassured.io/}
\end{itemize}
\end{ressources}

\index{JUnit}\index{AssertJ}\index{Mockito}\index{Testcontainers}\index{ArchUnit}\index{JMH}
% ============================================================================
\chapter{Java et les bases de donn\'ees~: JDBC, pools, JPA/Hibernate, migrations}
\label{ch:java-db}
% ============================================================================

\begin{objectif}
\begin{itemize}
  \item \'Ecrire du JDBC correct~: \texttt{DataSource}, requ\^etes pr\'epar\'ees, cl\'es g\'en\'er\'ees, lots,
        transactions, isolation.
  \item Cartographier proprement ResultSet \$\rightarrow\$ objets, sans ORM et sans N+1.
  \item Comprendre ce que Hibernate fait pour toi~: unit\'e de travail, cache de premier niveau,
        \og d\'etachement~\fg{}, et le pi\`ege N+1.
  \item Versionner le sch\'ema avec Flyway/Liquibase~: la migration est du code.
  \item Savoir quand se passer d'un ORM (et quand il est indispensable).
\end{itemize}
\end{objectif}

\section{Le sch\'ema de travail (r\'eutilis\'e dans tout le livre)}

\begin{codesql}
-- PostgreSQL 17 / MySQL 8 : les deux sont cites, les differences sont signalees.
CREATE TABLE auteurs (
  id         BIGINT GENERATED ALWAYS AS IDENTITY PRIMARY KEY,   -- MySQL : BIGINT AUTO_INCREMENT
  nom        TEXT         NOT NULL,
  email      TEXT         NOT NULL UNIQUE,
  cree_le    TIMESTAMPTZ  NOT NULL DEFAULT now()                 -- MySQL : TIMESTAMP(3) NOT NULL
);
CREATE TABLE articles (
  id         BIGINT GENERATED ALWAYS AS IDENTITY PRIMARY KEY,
  auteur_id  BIGINT       NOT NULL REFERENCES auteurs(id) ON DELETE RESTRICT,
  titre      TEXT         NOT NULL CHECK (char_length(titre) BETWEEN 3 AND 160),
  slug       TEXT         NOT NULL UNIQUE,
  corps      TEXT         NOT NULL,
  statut     TEXT         NOT NULL DEFAULT 'BROUILLON'
             CHECK (statut IN ('BROUILLON','PUBLIE','ARCHIVE')),
  vues       INTEGER      NOT NULL DEFAULT 0,
  publie_le  TIMESTAMPTZ,
  cree_le    TIMESTAMPTZ  NOT NULL DEFAULT now(),
  modifie_le TIMESTAMPTZ  NOT NULL DEFAULT now()
);
CREATE INDEX idx_articles_publie_le ON articles (publie_le DESC) WHERE statut = 'PUBLIE';
CREATE INDEX idx_articles_auteur    ON articles (auteur_id);
CREATE TABLE tags (
  id   SERIAL PRIMARY KEY,
  nom  TEXT NOT NULL UNIQUE
);
CREATE TABLE articles_tags (
  article_id BIGINT NOT NULL REFERENCES articles(id) ON DELETE CASCADE,
  tag_id     INT    NOT NULL REFERENCES tags(id)     ON DELETE CASCADE,
  PRIMARY KEY (article_id, tag_id)
);
CREATE TABLE comptes (
  id            BIGINT GENERATED ALWAYS AS IDENTITY PRIMARY KEY,
  identifiant   TEXT NOT NULL UNIQUE,
  mot_de_passe  TEXT NOT NULL,                       -- empreinte Argon2id, jamais le mot de passe
  role          TEXT NOT NULL DEFAULT 'AUTEUR',
  version       INTEGER NOT NULL DEFAULT 0,          -- verrouillage optimiste (Hibernate @Version)
  dernier_acc   TIMESTAMPTZ
);
CREATE OR REPLACE VIEW v_articles_publies AS
  SELECT a.id, a.titre, a.slug, a.publie_le, a.vues, u.nom AS auteur
  FROM articles a JOIN auteurs u ON u.id = a.auteur_id
  WHERE a.statut = 'PUBLIE';
\end{codesql}

\section{JDBC nu : la connexion et l'\'egouttoir}

\begin{codejava}
// On ne cree JAMAIS une connexion a la main dans une application : on passe par un DataSource.
var cfg = new com.zaxxer.hikari.HikariConfig();
cfg.setJdbcUrl("jdbc:postgresql://localhost:5432/blog?currentSchema=public&ApplicationName=blog");
cfg.setUsername("blog"); cfg.setPassword(System.getenv("DB_PASS"));
cfg.setMaximumPoolSize(10);            // (2 x coeurs) + nombre de disques, en regle generale
cfg.setMinimumIdle(2);
cfg.setConnectionTimeout(3_000);        // echec rapide > attente infinie
cfg.setMaxLifetime(30 * 60_000);         // < au timeout du proxy/pool cote serveur
cfg.setKeepaliveTime(60_000);
cfg.setLeakDetectionThreshold(20_000);   // ALERTE si une connexion n'est pas rendue : detecte les fuites
cfg.setPoolName("blog-pool");
cfg.addDataSourceProperty("prepareThreshold", "1");     // postges : server-side prepared statements
var pool = new com.zaxxer.hikari.HikariDataSource(cfg);
// MySQL : jdbc:mysql://hote:3306/blog?useSSL=false&serverTimezone=UTC&characterEncoding=UTF-8
//        &useServerPrepStmts=true&cachePrepStmts=true&rewriteBatchedStatements=true&connectTimeout=3000
// H2 (tests) : jdbc:h2:mem:blog;MODE=PostgreSQL;DATABASE_TO_LOWER=TRUE ; SQLite : jdbc:sqlite:app.db

try (var c = pool.getConnection()) {                 // le try-with-resources REND la connexion au pool
    c.setAutoCommit(false);
    try {
        // ... requetes ...
        c.commit();
    } catch (SQLException e) {
        try { c.rollback(); } catch (SQLException sup) { e.addSuppressed(sup); }
        throw e;
    }
}
// La regle d'or : une connexion = un "unite de travail". On l'ouvre tard, on la rend tot.
// Ne JAMAIS garder une connexion dans un champ d'une singleton ; ne JAMAIS creer un pool par requete.
\end{codejava}

\subsection{Requ\^etes~: le bon code, ligne \`a ligne}

\begin{codejava}
// SELECT avec parametres (JAMAIS de concatenation : injection SQL)
record ResumeArticle(long id, String titre, String slug, Instant publieLe, int vues, String auteur) { }

public List<ResumeArticle> recents(int limite) throws SQLException {
    var sql = """
            SELECT id, titre, slug, publie_le, vues, auteur
            FROM   v_articles_publies
            ORDER  BY publie_le DESC
            LIMIT  ?
            """;
    try (var c = pool.getConnection();
         var ps = c.prepareStatement(sql)) {
        ps.setInt(1, limite);
        ps.setFetchSize(200);                       // postgres : active le mode curseur (memoire!)
        try (var rs = ps.executeQuery()) {
            var out = new ArrayList<ResumeArticle>(limite);
            while (rs.next()) {
                out.add(new ResumeArticle(rs.getLong("id"), rs.getString("titre"), rs.getString("slug"),
                        rs.getTimestamp("publie_le").toInstant(), rs.getInt("vues"), rs.getString("auteur")));
            }
            return out;
        }
    }
}

// INSERT + recuperation de la cle generee (getGeneratedKeys, pas de "SELECT lastval" maison)
public long inserer(Article a) throws SQLException {
    var sql = """
            INSERT INTO articles (auteur_id, titre, slug, corps, statut, publie_le, cree_le, modifie_le)
            VALUES (?, ?, ?, ?, ?, ?, now(), now())
            """;
    try (var c = pool.getConnection();
         var ps = c.prepareStatement(sql, new String[] { "id" })) {
        ps.setLong(1, a.auteurId());
        ps.setString(2, a.titre());
        ps.setString(3, a.slug());
        ps.setString(4, a.corps());
        ps.setString(5, a.statut().name());
        ps.setTimestamp(6, a.publieLe() == null ? null : Timestamp.from(a.publieLe()));
        // Timestamp.set... : ne passe jamais un Instant par setObject sans connaitre le type cible
        ps.executeUpdate();
        try (var rs = ps.getGeneratedKeys()) {
            if (!rs.next()) { throw new SQLException("cle generee absente"); }
            return rs.getLong(1);
        }
    }
}

// UPSERT (PostgreSQL) / ON DUPLICATE KEY UPDATE (MySQL)
var up = """
        INSERT INTO tags (nom) VALUES (?) ON CONFLICT (nom) DO UPDATE SET nom = EXCLUDED.nom RETURNING id
        """;
// Lots : l'ecart de performance est de 1 a 2 ordres de grandeur
try (var c = pool.getConnection(); var ps = c.prepareStatement("INSERT INTO articles_tags (article_id, tag_id) VALUES (?, ?)")) {
    c.setAutoCommit(false);
    for (var tag : tags) {
        ps.setLong(1, id); ps.setInt(2, tag.id());
        ps.addBatch();
        if (++compteur % 1000 == 0) { ps.executeBatch(); }   // decoupe : evite un tampon geant
    }
    ps.executeBatch();
    c.commit();
}
// Requete dynamique sur un IN : generer N placeholders, jamais de "IN (" + join + ")"
var ids = List.of(1L, 7L, 42L);
var sqlIn = "SELECT id, titre FROM articles WHERE id IN (" +
        ids.stream().map(x -> "?").collect(Collectors.joining(",")) + ")";
try (var ps = c.prepareStatement(sqlIn)) {
    for (int i = 0; i < ids.size(); i++) { ps.setLong(i + 1, ids.get(i)); }
}
// LIKE securise : echapper %, _ et \ avant de lier le parametre
var motif = "%" + likeEscaper(recherche) + "%";
ps.setString(1, motif);
ps.setString(2, "\\");                       // ESCAPE '\' (Postgres : standard_conforming_strings)
// Paginations : OFFSET est OK jusqu'a ~10 000 lignes, au-dela utilise un curseur
//   OFFSET :   ORDER BY publie_le DESC LIMIT ? OFFSET ?
//   CURSEUR :  WHERE (publie_le, id) < (?, ?) ORDER BY publie_le DESC, id DESC LIMIT ?
//   (MySQL : WHERE publie_le < ? OR (publie_le = ? AND id < ?) ...)
// Compter pour la pagination : SELECT count(*) ... ; ou : SELECT ..., count(*) OVER() AS total (1 requete)
// Nullabilite : rs.getObject(...) puis rs.wasNull() ; ou setNull(i, Types.TIMESTAMP) pour ecrire un NULL.
// Types delicats : BigDecimal (NUMERIC), java.time via setTimestamp/getTimestamp, tableaux via createArrayOf,
// JSON via setObject(..., PGobject) ou ps.setString + colonne jsonb.
\end{codejava}

\begin{tablez}[p{0.2\textwidth}X]
  \headrow \thd{Sympt\^ome} \& \thd{Diagnostic et correction} \\
  \midrule
  \verb|connection timed out| / pool \'epuis\'e \& fuite de connexion (non rendue)~: \verb|leakDetectionThreshold|~; connexion gard\'ee hors transaction \\
  \verb|This result set is closed| \& le \verb|ResultSet| meurt avec son \verb|Statement|~: mat\'erialis avant de fermer \\
  R\'esultats partiels ou lents sur de gros volumes \& \verb|setFetchSize| et transaction en lecture seule (\verb|setReadOnly(true)|) \\
  \verb|Batch update returned ...| \& un trigger/contrainte casse le lot~: \verb|executeBatch| renvoie \verb|SUCCESS\_NO\_INFO|, teste ligne par ligne une fois \\
  Chiffres/textes cass\'es \`a l'insertion \& encodage de connexion (MySQL \verb|utf8mb4|, URL \verb|characterEncoding=UTF-8|) \\
  Dates d\'ecal\'ees d'une heure \& fuseau du serveur vs JVM~: explicite (\verb|serverTimezone=UTC|, \verb|TIMESTAMPTZ|) \\
  \verb| deadlock detected| \& ordre d'acc\`es aux lignes dans les transactions~: trie les ids, raccourcis la transaction, isol READ COMMITTED \\
\end{tablez}

\section{Le mapping maison~: la couche \og ~d\'ep\^ot~ \fg{} en 60 lignes}

\begin{codejava}
public interface Mappage<T> { T de(java.sql.ResultSet rs) throws SQLException; }

public abstract class Dep\^otJdbc {
    protected final javax.sql.DataSource source;
    protected Dep\^otJdbc(javax.sql.DataSource source) { this.source = source; }

    protected <T> List<T> liste(String sql, Mappage<T> m, Consumer<PreparedStatement> lieurs) throws SQLException {
        try (var c = source.getConnection(); var ps = c.prepareStatement(sql)) {
            lieurs.accept(ps);
            try (var rs = ps.executeQuery()) {
                var out = new ArrayList<T>();
                while (rs.next()) { out.add(m.de(rs)); }
                return out;
            }
        }
    }
    protected <T> Optional<T> un(String sql, Mappage<T> m, Consumer<PreparedStatement> lieurs) throws SQLException {
        var l = liste(sql, m, lieurs);
        return l.isEmpty() ? Optional.empty() : Optional.of(l.get(0));
    }
    protected int modifier(String sql, Consumer<PreparedStatement> lieurs) throws SQLException {
        try (var c = source.getConnection(); var ps = c.prepareStatement(sql)) {
            lieurs.accept(ps);
            return ps.executeUpdate();
        }
    }
}

public final class DepotArticlesJdbc extends Dep\^otJdbc implements Depot<Article, Long> {
    private static final Mappage<Article> MAPPEUR = rs -> new Article(
            rs.getLong("id"), rs.getString("titre"), rs.getString("corps"),
            Statut.valueOf(rs.getString("statut")),
            rs.getTimestamp("cree_le").toInstant(),
            rs.getTimestamp("publie_le") == null ? null : rs.getTimestamp("publie_le").toInstant(),
            rs.getObject("auteur_id", Long.class));

    @Override public Article sauver(Article a) { /* upsert + tags en lot, dans une transaction */ }
    @Override public Optional<Article> trouver(Long id) {
        return un("""
                SELECT a.id, a.titre, a.corps, a.statut, a.cree_le, a.publie_le, a.auteur_id
                FROM articles a WHERE a.id = ?""", MAPPEUR, ps -> ps.setLong(1, id));
    }
    // Le piege N+1, cote JDBC : charger les tags article par article. Les deux corrections :
    //  (a) une seule requete avec JOIN + regroupage en memoire ; (b) une requete "WHERE article_id = ANY(?)"
    //      qui recupere les tags des 20 articles de la page en un aller-retour.
    public List<Article> pageAvecTags(int page, int taille) throws SQLException {
        var sql = """
            WITH p AS (SELECT * FROM articles WHERE statut = 'PUBLIE'
                       ORDER BY publie_le DESC, id DESC LIMIT ? OFFSET ?)
            SELECT p.id, p.titre, p.corps, p.statut, p.cree_le, p.publie_le, p.auteur_id,
                   t.id AS tag_id, t.nom AS tag_nom
            FROM p LEFT JOIN articles_tags at ON at.article_id = p.id
                   LEFT JOIN tags t ON t.id = at.tag_id
            ORDER BY p.publie_le DESC, p.id DESC, t.nom
            """;
        // puis : regrouper par id (un LinkedHashMap<Long, Article> + accumulateur de tags)
    }
}
\end{codejava}

\begin{astuce}
Tu n'as pas besoin d'un ORM. Pour une application avec 5 \`a 30 requ\^etes, un paquet de d\'ep\^ots JDBC
\'ecrits \`a la main est plus rapide \`a comprendre, plus facile \`a d\'eboguer, et la performance est celle
du SQL que tu as \'ecrit. Les biblioth\`eques qui r\'eduisent la corv\'ee sans te prendre le contr\^ole~:
\link{jOOQ}{https://www.jooq.org/} (SQL typ\'e et compil\'e, tr\`es bon avec PostgreSQL),
\link{Jdbi}{https://jdbi.org/} (mapping l\'eger sur JDBC), \link{Spring JDBC}{https://docs.spring.io/spring-framework/reference/data-access/jdbc.html}
(\verb|JdbcTemplate|/ \verb|NamedParameterJdbcTemplate|).
\end{astuce}

\section{JPA et Hibernate~: l'unit\'e de travail}

\begin{codejava}
// entite : la classe est proxyabble -> non-final, constructeur sans argument requis
@Entity
@Table(name = "articles", indexes = { @Index(name = "idx_publie", columnList = "publie_le") })
public class ArticleEntity {
    @Id @GeneratedValue(strategy = GenerationType.IDENTITY)
    private Long id;

    @Version private int version;                    // verrouillage optimiste : evite l'ecrasement

    @Column(nullable = false, length = 160) private String titre;
    @Column(unique = true, nullable = false, length = 180) private String slug;
    @Column(nullable = false, columnDefinition = "text") @Basic(fetch = FetchType.LAZY) private String corps;
    @Enumerated(EnumType.STRING) @Column(nullable = false, length = 16) private Statut statut = Statut.BROUILLON;
    @Column(name = "publie_le") private Instant publieLe;
    @Column(name = "cree_le", nullable = false, updatable = false) private Instant creeLe = Instant.now();
    @ManyToMany(fetch = FetchType.LAZY)
    @JoinTable(name = "articles_tags",
               joinColumns = @JoinColumn(name = "article_id"), inverseJoinColumns = @JoinColumn(name = "tag_id"))
    private Set<Tag> tags = new HashSet<>();
    @ManyToOne(fetch = FetchType.LAZY, optional = false)
    // la BDD porte deja les contraintes : NO_CONSTRAINT evite que Hibernate en recree en doublon
    @JoinColumn(name = "auteur_id", foreignKey = @ForeignKey(ConstraintMode.NO_CONSTRAINT))
    private AuteurEntity auteur;
    @PreUpdate void avantMaj() { modifieLe = Instant.now(); }     // idem @PrePersist sur creeLe
    // getters/setters (JPA accede aux champs si @Id est sur le champ : sois coherent)
}

// Le repository (Spring Data JPA) : interface + methodes derivees + requetes ecrites a la main
public interface ArticleRepository extends JpaRepository<ArticleEntity, Long>, JpaSpecificationExecutor<ArticleEntity> {
    Optional<ArticleEntity> findBySlug(String slug);
    @Query("select a from ArticleEntity a left join fetch a.tags where a.statut = :statut order by a.publieLe desc")
    List<ArticleEntity> parStatut(@Param("statut") Statut statut, Pageable page);   // fetch = anti-N+1
    @Modifying @Query("update ArticleEntity a set a.vues = a.vues + 1 where a.id = :id")
    int incrementerVues(@Param("id") long id);
}
// Projection : evite de charger une entite complete quand tu veux 3 colonnes
public interface Entete { long id(); String titre(); String slug(); }
List<Entete> entetes = em.createQuery("select a.id as id, a.titre as titre, a.slug as slug from ArticleEntity a", Entete.class).getResultList();

// Le cycle de vie : TRANSIENT -> MANAGED (dans l'em) -> DETACHED (fin de transaction) -> REMOVED
// "could not initialize proxy - no Session" = acces a un chargement paresseux hors transaction.
// Corrections : (1) fetch join ; (2) mapper en projection dans la requete ; (3) @EntityGraph ;
// (4) DTO ; (5) ne pas sortir une entite du service. Sortir un DTO est LA bonne pratique.
// batch : hibernate.jdbc.batch_size=50, order_inserts/order_updates, fetch=128 (hibernate.default_batch_fetch_size)
// logs SQL : logging.level.org.hibernate.SQL=DEBUG + use_sql_comments=true (le source est affiche)
// N+1 : hibernate.enable_lazy_load_no_trans = JAMAIS (cache le probleme en ouvrant des sessions)
\end{codejava}

\begin{tablez}[p{0.3\textwidth}p{0.22\textwidth}X]
  \headrow \& \thd{JDBC + mapping} \& \thd{JPA/Hibernate} \\
  \midrule
  Mod\`ele riche (r\`egles, objets immuables) \& excellent \& contraint (champs, proxy, setters) \\
  Requ\^etes analytiques, SQL fin \& \textbf{le terrain naturel} \& requ\^etes natives, projections \\
  CRUD d'un back-office, many-to-many \& verbeux \& \textbf{brillant} \\
  Migration de sch\'ema \& \`a toi (Flyway) \& \verb|ddl-auto=validate| (jamais \verb|update| en prod) \\
  Performance \& pr\'evisible, aucune surprise \& cache, batch, plans--- mais exige de la rigueur \\
\end{tablez}

\section{Migrations~: le sch\'ema est du code}

\begin{terminal}
# Flyway (le plus repandu) -- src/main/resources/db/migration
#   V1__schema_initial.sql   V2__ajout_tags.sql   R__vue_articles.sql   U1__annule_xxx.sql
flyway migrate -url=jdbc:postgresql://localhost/blog -user=blog -password=$DB_PASS
flyway info ; flyway validate ; flyway baseline -baselineVersion=1
#   undo (payant) : flyway undo ; communautaires : Flyway 10+ "community" gratuit pour postgres/mysql
# Liquibase : le XML/YAML/SQL avec changelog et rollback (les entreprises l'exigent parfois)
liquibase update --changeLogFile=db/changelog.xml
\end{terminal}

\begin{codeconf}
<!-- pom.xml : flyway-maven-plugin -->
<plugin>
  <groupId>org.flywaydb</groupId><artifactId>flyway-maven-plugin</artifactId><version>11.2.0</version>
  <configuration>
    <url>${db.url}</url><user>${db.user}</user><password>${db.password}</password>
    <locations><location>filesystem:src/main/resources/db/migration</location></locations>
    <validateMigrationNaming>true</validateMigrationNaming><cleanDisabled>true</cleanDisabled>
  </configuration>
</plugin>
\end{codeconf}

\begin{codeconf}
# application.properties (Spring Boot) : les reglages qui comptent
spring.datasource.url=jdbc:postgresql://localhost:5432/blog
spring.datasource.hikari.maximum-pool-size=10
spring.datasource.hikari.leak-detection-threshold=20000
spring.jpa.hibernate.ddl-auto=validate
spring.jpa.open-in-view=false                       # IMPORTANT : sinon la session vit toute la requete
spring.jpa.properties.hibernate.jdbc.batch_size=50
spring.jpa.properties.hibernate.order_inserts=true
spring.jpa.properties.hibernate.query.in_clause_parameter_padding=true
spring.jpa.properties.hibernate.session_factory.interceptor=...   # pour loguer le N+1
spring.flyway.enabled=true
spring.jpa.show-sql=false
logging.level.org.hibernate.orm.jdbc.bind=TRACE      # les valeurs liees (debug seulement)
\end{codeconf}

\begin{piege}
Trois blessures fr\'equentes en JPA~: (1) \verb|ddl-auto=update| en production~: Hibernate modifie le
sch\'ema, perd des contraintes, ou \'echoue \`a 3 h~; (2) \verb|open-in-view=true|~: la session reste
ouverte pendant le rendu, un appel lent c\^ot\'e client garde une connexion du pool~; (3) \verb|List| au
lieu de \verb|Set| dans un \verb|@OneToMany|~: les lots d'insertion se transforment en delete+insert
\`a chaque sauvegarde. Les trois sont document\'es, et les trois font mal en silence.
\end{piege}

\begin{exercice}
\begin{enumerate}
  \item \'Ecris le d\'ep\^ot JDBC complet de l'application du livre~: \verb|sauver|, \verb|trouver|,
        \verb|lister(pageable, filtre)|, \verb|compter|, \verb|rechercher(texte)| (avec \'echappement
        LIKE~; puis tsvector si tu utilises PostgreSQL)~: avec un seul aller-retour par page.
  \item Ajoute la transaction~: publier un article = mise \`a jour + tag + journal~;rollback en cas
        d'\'echec~; d\'emontre-le avec un test qui fait \'echouer volontairement le 3\ieme insert.
  \item Fais tourner le m\^eme d\'ep\^ot avec H2, PostgreSQL (Testcontainers) et SQLite~: observe les
        divergences de syntaxe (upsert, \verb|LIMIT|, \verb|SERIAL|) et d\'ecide o\`u tu mets la ligne.
  \item \'Ecris la migration Flyway qui ajoute une colonne \verb|lecture\_minutes| calcul\'ee~: (a) avec
        une valeur par d\'efaut pour les lignes existantes~; (b) sans verrouiller la table~: en trois
        \'etapes (add nullable~: update batch~: set not null). Puis v\'erifie la taille du WAL et le temps.
  \item Reprends la m\^eme entit\'e en Hibernate~: mappe, charge une page, affiche le nombre de
        requ\^etes \'emises (logs) avant/apr\`es un \verb|join fetch|~: le gain doit \^etre mesurable.
  \item Compare objectivement~: le CRUD en JDBC~; jOOQ~; Spring Data JPA~: sur les m\^emes crit\`eres
        (lignes de code, nombre de requ\^etes, lisibilit\'e du diff en revue, temps de r\'eponse p95).
        C'est l'exercice qui te fait choisir un outillage pour trois ans.
\end{enumerate}
\end{exercice}

\begin{ressources}
\begin{itemize}
  \item JDBC API --- \url{https://docs.oracle.com/en/java/javase/21/docs/api/java.sql/java/sql/package-summary.html}
  \item HikariCP --- \url{https://github.com/brettwooldridge/HikariCP} (le wiki \og ~pool sizing~ \fg{} est une lecture obligatoire)
  \item Hibernate ORM --- \url{https://hibernate.org/orm/documentation/} --- Vlad Mihalcea's blog --- \url{https://vladmihalcea.com/}
  \item Flyway --- \url{https://documentation.red-gate.com/flyway} --- Liquibase --- \url{https://docs.liquibase.com/}
  \item jOOQ --- \url{https://www.jooq.org/} --- Jdbi --- \url{https://jdbi.org/}
  \item PostgreSQL --- \url{https://www.postgresql.org/docs/current/} (le chapitre sur les index partiels et \verb|EXPLAIN|)
\end{itemize}
\end{ressources}

\index{JDBC}\index{HikariCP}\index{Hibernate}\index{JPA}\index{Flyway}\index{N+1}
% ============================================================================
\chapter{Le web en Java~: serveurs, API REST, s\'ecurit\'e, observabilit\'e}
\label{ch:java-web}
% ============================================================================

\begin{objectif}
\begin{itemize}
  \item Comprendre ce qu'un serveur HTTP Java fait d'une requ\^ete~: servlet, filtre, gestionnaire.
  \item \'Ecrire une API REST compl\`ete~: routes, validation, erreurs RFC 9457, pagination, CORS.
  \item Authentifi\'e~: sessions, JWT, OIDC~; et prot\'eger (en-t\^etes, limitation, secrets).
  \item Rendre l'application observable~: sant\'e, mesures, traces, journaux structur\'es.
  \item D\'eployer proprement~: conteneur, arr\^et gracieux, dimensionnement du pool.
\end{itemize}
\end{objectif}

\section{Le mod\`ele~: un thread (ou un virtual thread) par requ\^ete}

\begin{codeplain}
Client --TCP--> Serveur HTTP (Jetty/Tomcat/Undertow/Netty/Vert.x/Kestrel-like)
                    |
                    +-- filtres (CORS, auth, logs, compression, limitation)
                    +-- gestionnaire / servlet / contr\^oleur   <-- ton code
                            +-- validation, service m\'etier, d\'ep\^ot (pool JDBC)
                            +-- s\'erialisation JSON (Jackson)
                    +-- r\'eponse : statut, en-t\^etes, corps
Le "~servlet container~" historique (Jakarta Servlet 6) a \'et\'e supplant\'e dans les nouveaux
projets par des cadres embarqu\'es : le serveur est une dependance du projet, pas un
serveur d'application ou un WAR deploye a la main.
\end{codeplain}

\begin{tablez}[p{0.16\textwidth}p{0.14\textwidth}X]
  \headrow Cadre \& Poids \& Usage typique \\
  \midrule
  Spring Boot 3.x \& lourd (mais complet) \& le choix par d\'efaut en entreprise~: tout est l\`a \\
  Quarkus 3.x \& l\'eger, natif, r\'eactif \& cloud-native, serverless, Kotlin/Java, dev services \\
  Micronaut 4.x \& compilation-time DI, natif \& microservices, peu de r\'eflexion, AWS Lambda \\
  Helidon 4 \& Netty/loom, orient\'e API \& Oracle, r\'eactif et threads virtuels \\
  Javalin / JFinal \& minimal (Jetty) \& petite API, outils internes, prototypes rapides \\
  Vert.x \& \'ev\'enements non bloquants \& flux, websockets, d\'ebit extr\^eme \\
  Dropwizard \& Jetty + Jackson + JDBI \& services monolithiques sobres \\
\end{tablez}

\section{Une API compl\`ete avec Javalin (le minimum vital, sans magie)}

\begin{codejava}
// Dependances : io.javalin:javalin:6.x + ch.qos.logback:logback-classic + jackson-databind
public final class App {
    public static void main(String[] args) {
        var config = Config.lire();                          // env + application.yml
        var pool = HikariFactory.create(config);              // cf. chapitre precedent
        var depot = new DepotArticlesJdbc(pool);
        var service = new ServiceArticles(depot, Clock.systemUTC());

        var app = Javalin.create(cfg -> {
            cfg.http.defaultContentType = "application/json";
            cfg.bundledPlugins.enableCors(c -> c.addPolicy(p -> p.allowOrigins("https://blog.exemple.fr")
                    .allowMethods("GET,POST,PATCH,DELETE").allowHeaders("Authorization,Content-Type")
                    .maxAge(3600)));
            cfg.plugins.enableDevLogging();                   // en dev seulement
            cfg.staticFiles.add("/public", Location.CLASSPATH);
            cfg.showJettyServerStartLogs = false;
        });

        // Filtre : requete -> contexte (chronometrage, identifiant de trace, correlation)
        app.before(ctx -> {
            ctx.attribute("debut", System.nanoTime());
            ctx.res.setHeader("X-Request-Id", ctx.header("X-Request-Id") != null
                    ? ctx.header("X-Request-Id") : UUID.randomUUID().toString());
            ctx.res.setHeader("X-Content-Type-Options", "nosniff");
            ctx.res.setHeader("Referrer-Policy", "same-origin");
            ctx.res.setHeader("Permissions-Policy", "camera=(), microphone=(), geolocation=()");
            // CSP : adaptee a ton frontend ; "default-src 'self'" est un bon point de depart
            ctx.res.setHeader("Content-Security-Policy", "default-src 'self'; img-src 'self' data:");
        });
        app.after(ctx -> METRIQUES.enregistrer(ctx.method(), ctx.path(), ctx.status(),
                System.nanoTime() - (long) ctx.attribute("debut")));

        // Routes
        app.get("/health", App::sante);
        app.get("/ready", ctx -> ctx.result("pr\^et"));
        app.get("/api/articles", ctx -> {
            var page = PageRequest.de(ctx.queryInt("page", 0), ctx.queryInt("taille", 20));
            ctx.json(Page.of(service.recents(page), page));
        });
        app.get("/api/articles/{id}", ctx -> ctx.json(service.obtenir(ctx.pathLong("id"))));
        app.post("/api/articles", ctx -> {
            var requete = BodyValidator.lire(ctx, RequeteArticle.class);      // 400 + details si invalide
            var cree = service.creer(requete, ctx.attribute("utilisateur"));
            ctx.status(201).header("Location", "/api/articles/" + cree.id()).json(cree);
        });
        app.patch("/api/articles/{id}", ctx -> ctx.json(service.modifier(ctx.pathLong("id"),
                BodyValidator.lire(ctx, ModificationArticle.class))));
        app.delete("/api/articles/{id}", ctx -> { service.supprimer(ctx.pathLong("id")); ctx.status(204); });

        // Authentification : Bearer -> decodage -> sujet dans l'attribut de contexte
        app.before("/api/*", ctx -> {
            var entete = ctx.header("Authorization");
            if (entete == null || !entete.startsWith("Bearer ")) { throw new ErreurHttp(401, "jeton requis"); }
            ctx.attribute("utilisateur", jwt.verifier(entete.substring(7)));  // see chapitre Spring Security
        });
        app.get("/api/mes-articles", ctx -> ctx.json(service.parAuteur(ctx.attribute("utilisateur").id())));

        // Gestion d'erreurs : UN SEUL endroit, format RFC 9457 (problem+json)
        app.exception(ErreurValidation.class, (e, ctx) -> ctx.status(422).json(Probleme.de(ctx, 422,
                "Donnees invalides", "validation", e.champs())));
        app.exception(ErreurIntrouvable.class, (e, ctx) -> ctx.status(404).json(
                Probleme.de(ctx, 404, e.getMessage(), "introuvable")));
        app.exception(ErreurHttp.class, (e, ctx) -> ctx.status(e.statut()).json(Probleme.de(ctx, e.statut(), e.getMessage(), e.code())));
        app.throwable((e, ctx) -> {
            log.error("non ger\'e", e);
            ctx.status(500).json(Probleme.de(ctx, 500, "erreur interne", "interne")
                    .detail(config.isDev() ? e.toString() : null));     // jamais la stack en prod
        });

        app.start(config.port());
        Runtime.getRuntime().addShutdownHook(new Thread(() -> {
            app.stop();
            pool.close();                       // arret gracieux : finir les requetes en cours
        }, "arret"));
    }
}
\end{codejava}

\subsection{Les formes que doit avoir une API (rappel utile en Java)}

\begin{tabz}{p{0.24\textwidth}X}
  \headrow \thd{Sujet} \& \thd{Conventions \`a tenir} \\
  \midrule
  Noms \& Pluriels, substantifs, pas de verbes~: \verb|/articles/12/tags| et non \verb|/getArticle|. \\
  Verbes \& \verb|GET| (s\^ur, idempotent), \verb|POST|, \verb|PUT| (remplacement), \verb|PATCH| (partiel), \verb|DELETE| (204). \\
  Codes \& 200/201/204~; 400 (corps illisible) 401 403 404 409 (conflit) 422 (s\'emantique) 429~; 500/503. \\
  Erreurs \& \verb|application/problem+json| (RFC 9457)~: \verb|type, title, status, detail, instance| + \verb|errors|. \\
  Pagination \& \verb|page/taille| ou curseur opaque~; renvoyer \verb|total| ou un \verb|Link: rel="next"|. \\
  Filtrage \& \verb|?statut=PUBLIE\&depuis=2026-01-01|~; tri \verb|?tri=-publie\_le,titre|. \\
  Versionnage \& \verb|/api/v1| \emph{ou} en-t\^ete \verb|Accept-API-Version|~: un seul style, document\'e. \\
  Idempotence \& \verb|Idempotency-Key| sur les \verb|POST| qui \'ecrivent (paiements, publications). \\
  D\'ebit \& limitation par token bucket~; \verb|Retry-After|~; file born\'ee en interne. \\
\end{tabz}

\section{Authentification et autorisation~: la gamme compl\`ete}

\begin{codejava}
// 1. Mots de passe : Argon2id (recommande), bcrypt (acceptable), NEVER plain/md5/sha1.
//    dependance : de.mkammerer:argon2-jvm (liens natifs) ou Bouncy Castle.
var argon = Argon2Factory.create(Argon2Factory.Argon2Types.ARGON2id);
String empreinte = argon.hash(3, 64 * 1024, 1, motDePasse.toCharArray());   // t=3 m=64Mo p=1
boolean ok = argon.verify(empreinte, motDePasse.toCharArray());
if (argon.requiresRehash(empreinte)) { depot.nouvelleEmpreinte(user, argon.hash(...)); }  // cout qui baisse

// 2. JWT (cf. chapitre PHP 308 pour les pieges : les memes s'appliquent)
//    - verifie alg (pas "none"), exp/nbf/iat avec tolerance, iss/aud, kid ;
//    - HS256 : secret >= 256 bits, tourne ; RS256/ES256 : cle privee cote auteur, JWK publics exposes.
//    - ne JAMAIS mettre de donnees sensibles dans le payload (il est lisible).
var lib = com.nimbusds.jose.JOSEObjects...       // nimbus-jose-jwt : la lib serieuse

// 3. OIDC : deleguer l'authentification a un fournisseur (Keycloak, Auth0, Entra, Google, GitHub).
//    Flux : authorization code + PKCE ; cote Java : spring-boot-starter-oauth2-resource-server
//    ou Quarkus openid-connect. Tu ne geranges plus de mot de passe : c'est le but.
//    Le jeton d'acces est opaque ou JWT ; verifie la signature via la JWK Set du fournisseur,
//    mets les roles en claim, et fais l'autorisation dans TON code (pas dans le jeton seul).

// 4. Sessions : pour un site avec navigateur, c'est souvent le MEILLEUR choix (revocation possible).
//    Cookie : Secure, HttpOnly, SameSite=Lax, Path=/ ; expire ; rotation apres login ;
//    cote serveur : table de sessions (ou Redis) avec TTL, et suppression a la deconnexion.

// 5. Autorisation : centralise, declasre, testable.
enum Permission { ARTICLE_CREER, ARTICLE_PUBLIER, ARTICLE_SUPPRIMER, USER_ADMIN }
record Sujet(long id, String nom, Set<Permission> permissions) {
    boolean peut(Permission p) { return permissions.contains(p); }
}
// Middleware : ctx.attribute(sujet) puis, dans chaque route : autoriser(ctx, Permission.ARTICLE_PUBLIER)
// => l'oubli d'un controle devient une faute visible en revue, pas une surprise en production.
\end{codejava}

\section{Observabilit\'e~: ce qui te sauvera \`a 3 h}

\begin{codejava}
// Micrometer : la facade de mesures (le Prometheus client, OTLP, Datadog via registries)
var registry = PrometheusMeterRegistry.prometheusNativeRegistry ? null : new PrometheusMeterRegistry(PrometheusConfig.DEFAULT);
Timer.Builder.builder("http.server.requests", registry)
        .tag("uri", uri).tag("methode", meth).tag("statut", String.valueOf(code)).register(registry)
        .record(duree, TimeUnit.NANOSECONDS);
registry.summary("articles.taille", "tags").record(nbTags);
// Exposition : GET /metrics (texte Prometheus) + une alerte si 5xx > 1 % sur 5 min.

// Journaux structures (JSON) : logback + net.logstash.logback:logstash-logback-encoder
//    <encoder class="net.logstash.logback.encoder.LogstashEncoder"><includeMdcKeyName>requestId</includeMdcKeyName>
MDC.put("requestId", id);                       // propage a travers les threads (ou Scoped Values / OTel)
try (var scope = Tracer.current().with(span)) { traitement(); }

// OpenTelemetry : traces + metriques + logs, le standard qui remplace les proprietaires
// io.opentelemetry.javaagent (agent zero-code) : java -javaagent:otel.jar -jar app.jar
//    OTEL_SERVICE_NAME, OTEL_EXPORTER_OTLP_ENDPOINT, OTEL_TRACES_SAMPLER=parentbased_traceidratio
// Correlation : un meme requestId dans le log, la trace et le reponse HTTP (X-Request-Id).
\end{codejava}

\begin{terminal}
# JVM utile en production
java -XX:MaxRAMPercentage=75 -XX:+ExitOnOutOfMemoryError \
     -Xlog:gc*:file=/var/log/app/gc.log:time,uptime:filecount=5,filesize=20M \
     -XX:+HeapDumpOnOutOfMemoryError -XX:HeapDumpPath=/var/data/dumps \
     --enable-native-access=ALL-UNNAMED \
     -javaagent:/app/opentelemetry-javaagent.jar \
     -jar /app/app.jar
# Diagnostic a chaud
jcmd 1 GC.run ; jcmd 1 Thread.print ; jcmd 1 VM.flags ; jcmd 1 VM.system_properties
jstat -gcutil 1 1s 10            # old gen qui ne redescend pas = fuite : heap dump puis MAT/JProfiler
\end{terminal}

\etape{Le d\'eploiement qui tient la route}
\begin{enumerate}
  \item Conteneur non-root, \verb|readOnlyRootFilesystem|, \verb|drop ALL capabilities|, limites CPU/RAM~;
  \item sonde \verb|/health| (liveness~: le processus vit) et \verb|/ready| (readiness~: pool,
        migration, d\'ependances pr\^etes)~: deux endpoints diff\'erents, jamais un seul~;
  \item arr\^et gracieux~: SIGTERM \$\rightarrow\$ stop accepting, attendre 15 s, rendre les connexions BDD, sortir.
        En Kubernetes~: \verb|terminationGracePeriodSeconds: 30|, sinon tu perds des requ\^etes \`a chaque rollout.
  \item Secrets hors de l'image~: variables d'environnement, fichiers mont\'es, Vault/Sealed Secrets.
  \item Migration de sch\'ema \emph{avant} le d\'eploiement du code~: backwards-compatible toujours
        (ajouter une colonne~: OK~; renommer~: en deux versions~: add+copy+switch, drop plus tard).
  \item Build reproductible~: pas de \verb|SNAPSHOT| en prod, digest d'image \'epingl\'e, SBOM
        (\link{CycloneDX Maven plugin}{https://github.com/CycloneDX/cyclonedx-maven-plugin})
        + \link{OWASP dependency-check}{https://github.com/jeremylong/DependencyCheck} ou
        \link{Trivy}{https://github.com/aquasecurity/trivy}.
\end{enumerate}

\begin{piege}
Cinq blessures d'API Java en production~: (1) \verb|List| enti\`ere charg\'ee puis pagin\'ee en m\'emoire~;
(2) \verb|@Transactional| sur une m\'ethode publique appel\'ee depuis la m\^eme classe~: l'intercepteur
ne passe pas, la transaction n'existe pas~; (3) Jackson qui s\'erialise l'entit\'e JPA et d\'eclenche un
chargement paresseux en cascade~: 40 requ\^etes par r\'eponse~; (4) pool JDBC plus grand que
\verb|max\_connections| du serveur~: le pic de trafic tue la base~; (5) aucun timeout HTTP sortant
sur les appels d'autres services~: un thread bloqu\'e pour toujours, puis tout le pool.
\end{piege}

\begin{exercice}
\begin{enumerate}
  \item Reprise le projet du chapitre~9 et construis l'API en Javalin~: les 6 routes CRUD, les
        erreurs 9457, la pagination~; et le fichier \verb|openapi.yaml| (valid\'e avec
        \url{https://code.visualstudio.com/} ou Scalar).
  \item Ajoute l'authentification~: inscription/connexion avec Argon2id, sessions en base,
        cookies durcis, limitation \`a 5 tentatives~: puis une variante JWT (access 15 min + refresh).
  \item \'Ecris les tests d'int\'egration de l'API~: RestAssured contre une instance r\'eelle avec
        Testcontainers~; 25 cas~: dont chaque code d'erreur, chaque r\`egle d'autorisation.
  \item Instrumente~: timer HTTP, compteur d'erreurs par code, jauge de taille du pool
        (\verb|active|, \verb|idle|, \verb|pending|)~; puis \verb|/metrics| scrap\'e par Prometheus
        local~: et un panneau Grafana de trois graphiques.
  \item Fais le test de charge~: \link{k6}{https://k6.io/} ou \link{wrk}{https://github.com/wg/wrk}
        sur \verb|GET /api/articles|~: trouve le point de rupture~; augmente la taille du pool,
        puis le nombre de threads~: note laquelle des deux variables d\'eplace le d\'ebit, et explique.
  \item Reprends la m\^eme API en Spring Boot (chapitre suivant)~: compte les fichiers et les lignes
        n\'ecessaires~; c'est l'int\'er\^et (et le co\^ut) d'un cadre~: tu le verras en un apr\`es-midi.
\end{enumerate}
\end{exercice}

\begin{ressources}
\begin{itemize}
  \item Javalin --- \url{https://javalin.io/documentation} --- Jetty --- \url{https://www.eclipse.org/jetty/}
  \item Micrometer --- \url{https://micrometer.io/} --- OpenTelemetry Java --- \url{https://opentelemetry.io/docs/languages/java/}
  \item Nimbus JOSE+JWT --- \url{https://connect2id.com/products/libraries/nimbus-jose-jwt} --- Keycloak --- \url{https://www.keycloak.org/}
  \item OWASP API Security Top 10 --- \url{https://owasp.org/API-Security/editions/2023/en/0x00-header/}
  \item RFC 9457 (Problem Details) --- \url{https://www.rfc-editor.org/rfc/rfc9457.html} --- \link{12-factor App}{https://12factor.net/fr/}
\end{itemize}
\end{ressources}

\index{serveur HTTP}\index{API REST}\index{JWT}\index{Micrometer}\index{OpenTelemetry}
% ============================================================================
\chapter{Spring Boot~: l'atelier complet}
\label{ch:java-spring}
% ============================================================================

\begin{objectif}
\begin{itemize}
  \item Comprendre ce que Spring Boot automatise (et ce qu'il ne fait pas \`a ta place).
  \item Construire une application MVC compl\`ete~: contr\^oleurs, validation, gestion d'erreurs,
        s\'ecurit\'e, donn\'ees, tests.
  \item Lire une configuration Spring sans se noyer~: beans, profils, propri\'et\'es, auto-configuration.
  \item Choisir Spring~: ou d\'ecider de ne pas le choisir.
\end{itemize}
\end{objectif}

\section{Pourquoi Spring~: le contrat}

\begin{description}
  \item[Boot] des \emph{starters} (d\'ependances coh\'erentes), de l'\emph{auto-configuration} (des beans
        cr\'e\'es si les classes sont l\`a et rien ne les configure), un serveur embarqu\'e, un fichier
        de configuration, un \verb|actuator| de supervision.
  \item[Framework] IoC/DI (\verb|ApplicationContext|), MVC, Data, Security, Transactions, Messaging,
        Test --- tous coh\'erents et test\'es ensemble.
  \item[Co\^ut] abstractions, r\'eflexion, taille de la stack, courbe de lecture des messages,
        d\'emarrage lent hors image native. Une \og ~simple~ \fg{} \verb|@SpringBootApplication| monte
        une centaine de beans.
\end{description}

\begin{terminal}
# Creer le projet : https://start.spring.io (Dependencies : Web, Data JPA, Validation,
#   Security, Actuator, PostgreSQL, Flyway, Testcontainers(?) -> puis "Generate Project")
# ou en CLI :
curl "https://start.spring.io/starter.zip" -d type=maven-project -d javaVersion=21 \
  -d dependencies=web,data-jpa,validation,security,actuator,postgresql,flyway \
  -d groupId=com.exemple -d artifactId=blog -d name=blog -o blog.zip && unzip blog.zip
# L'outil de ligne de commande officiel (Boot 3.5+) :
sdk install springboot ; spring init --dependencies=web,data-jpa blog.zip
./mvnw spring-boot:run          # http://localhost:8080
./mvnw spring-boot:build-image  # image OCI via buildpacks (pas de Dockerfile a ecrire)
\end{terminal}

\begin{codeplain}
src/main/java/com/exemple/blog/
\boxv \boxh \boxh BlogApplication.java                 @SpringBootApplication
\boxv \boxh \boxh api/
\boxv \boxv \boxv \boxv \boxh \boxh ArticleController.java           @RestController : le monde HTTP
\boxv \boxv \boxv \boxv \boxh \boxh dto/                             records de requete/reponse (+ mapper)
\boxv \boxv \boxv \boxe \boxh \boxh GlobalExceptionHandler.java      @RestControllerAdvice
\boxv \boxh \boxh domain/                              entites, regles, services : AUCUN import Spring ici
\boxv \boxh \boxh infra/
\boxv \boxv \boxv \boxv \boxh \boxh ArticleRepository.java           Spring Data (interface)
\boxv \boxv \boxv \boxe \boxh \boxh config/                          @Configuration (securite, HTTP client, Jackson, openapi)
\boxe \boxh \boxh src/test/java/...                     miroir : *Test (unites), *IT (integration)
src/main/resources/
\boxv \boxh \boxh application.yml  application-dev.yml  application-prod.yml
\boxv \boxh \boxh db/migration/V1__init.sql
\boxe \boxh \boxh static/ , templates/ (si Thymeleaf)
\end{codeplain}

\section{Le conteneur de beans, en 30 lignes de t\^ete}

\begin{codejava}
@SpringBootApplication                // = @Configuration + @EnableAutoConfiguration + @ComponentScan
public class BlogApplication {
    public static void main(String[] args) {
        var ctx = SpringApplication.run(BlogApplication.class, args);
        // ctx.getBeanDefinitionCount()  -->  ~150 beans pour une app "vide" : c'est normal
        for (String n : ctx.getBeanDefinitionNames()) { if (n.contains("DataSource")) System.out.println(n); }
    }
    @Bean
    Clock horloge() { return Clock.systemDefaultZone(); }            // une methode = un bean
    @Bean
    @Profile("prod")
    Filter corsProd() { /* ... */ }
    @Bean
    @ConditionalOnProperty(name = "blog.feature.export", havingValue = "true")
    Exporteur exporteur() { return new ExporteurCsv(); }             // bean conditionnel
}

// L'injection : par constructeur, toujours (finale, testable, impossible d'oublier une dependance)
@Service
@RequiredArgsConstructor                 // Lombok ; sinon ecris le constructeur a la main
class ServiceArticles {
    private final ArticleRepository repo;
    private final Journal journal;
    private final Clock horloge;
    // @Autowired sur les champs = lisibilite en moins, immuabilite en moins, tests plus lourds : a eviter.
    // @Value / environnement : un composant @ConfigurationProperties (typ\'e) > des @Value parsemes.
}

// Scopes : singleton (defaut), prototype (nouvelle instance a chaque injection/getBean),
// request / session / websocket (actifs uniquement dans un contexte web).
// Le cycle de vie : constructeur -> @Autowired -> @PostConstruct (ou InitializingBean)
//                   -> utilisation -> @PreDestroy -> DisposableBean.
// Circulaires : Spring les refuse par defaut depuis Boot 2.6 ("circular reference") -- c'est un
// signal de conception, pas un obstacle a franchir avec @Lazy.
// Les extensions qui t'evitent du code : @ConfigurationProperties, @EventListener,
// @Async (attention au pool), @Scheduled, @Cacheable (avec Caffeine/Redis), @Validated.
\end{codejava}

\section{Configuration : le fichier, les profils, le typage}

\begin{codeconf}
# application.yml
spring:
  application.name: blog
  threads.virtual.enabled: true            # Boot 3.2+ : Tomcat sur threads virtuels
  datasource:
    url: ${DB_URL:jdbc:postgresql://localhost:5432/blog}
    username: ${DB_USER:blog}
    password: ${DB_PASS:}                  # jamais de mot de passe en dur dans le depot
    hikari: { maximum-pool-size: 10, leak-detection-threshold: 20000 }
  jpa:
    open-in-view: false
    hibernate.ddl-auto: validate
    properties.hibernate.jdbc.batch_size: 50
  flyway: { enabled: true, locations: classpath:db/migration }
  security.user.name: ${admin}             # (dev seulement ; en prod : OIDC ou base)
server:
  port: ${PORT:8080}
  shutdown: graceful
  error.include-message: never
  compression.enabled: true
  forward-headers-strategy: framework       # derriere un proxy : sinon les URLs sont fausses
management:
  endpoints.web.exposure.include: health,info,prometheus,metrics
  endpoint.health.show-details: when_authorized
blog:
  upload: { taille-max: 5MB, dossiers-autorises: [image, pdf] }
  rate-limit: { par-minute: 120, burst: 40 }
---
spring.config.activate.on-profile: dev
logging.level.com.exemple: DEBUG
spring.jpa.show-sql: true
---
spring.config.activate.on-profile: prod
logging.structured.format.console: ecs          # ou gelf/ecs/otel : des logs que Logstash avale
management.endpoint.health.probes.enabled: true # /actuator/health/readiness + /liveness
\end{codeconf}

\begin{codejava}
// La configuration typee : la seule facon propre de lire du YAML
@ConfigurationProperties(prefix = "blog")
@Validated
record ConfigBlog(@NotNull Upload upload, @NotNull RateLimit rateLimit) {
    record Upload(@DataSize.Max(10_000_000) long tailleMax, List<String> dossiersAutorises) { }
    record RateLimit(@Min(1) int parMinute, @Min(1) int burst) { }
}
// Enregistre : @EnableConfigurationProperties(ConfigBlog.class) ou @ConfigurationPropertiesScan.
// 3 sources, dans l'ordre de priorite (la plus forte ecrase) :
//   ligne de commande (--blog.rate-limit.par-minute=50) > variables d'env (BLOG_RATELIMIT_...) >
//   application-{profil}.yml > application.yml > valeurs par defaut du @Bean.
// Secret : jamais dans le YAML versionne ; env, Vault, AWS Secrets, ou un fichier monte en lecture seule.
\end{codejava}

\section{MVC~: contr\^oleurs, validation, erreurs, contenu}

\begin{codejava}
@RestController
@RequestMapping(path = "/api/articles", produces = MediaType.APPLICATION_JSON_VALUE)
public class ArticleController {
    private final ServiceArticles service;
    public ArticleController(ServiceArticles service) { this.service = service; }

    @GetMapping
    @Cacheable(value = "articlesPage")                       // si tu as mis un cache
    public Page<ArticleResume> liste(@RequestParam(defaultValue = "0") @Min(0) int page,
                                     @RequestParam(defaultValue = "20") @Min(1) @Max(100) int taille,
                                     @RequestParam(required = false) Statut statut,
                                     @PageableDefault(size = 20, sort = "publieLe", direction = Sort.Direction.DESC)
                                     Pageable pageable) {
        return service.rechercher(statut, pageable);
    }

    @GetMapping("/{id}")
    public ArticleDetail detail(@PathVariable @Positive long id) { return service.detail(id).orElseThrow(() -> new Introuvable(id)); }

    @PostMapping
    @ResponseStatus(HttpStatus.CREATED)
    public ResponseEntity<ArticleResume> creer(@Valid @RequestBody CreerArticle requete,
                                               @AuthenticationPrincipal UtilisateurCourant moi,
                                               UriComponentsBuilder uri) {
        var cree = service.creer(requete, moi.id());
        return ResponseEntity.created(uri.path("/api/articles/{id}").buildAndExpand(cree.id()).toUri()).body(cree);
    }

    @PatchMapping("/{id}")
    public ArticleResume modifier(@PathVariable long id, @Valid @RequestBody ModifierArticle requete) { return service.modifier(id, requete); }

    @DeleteMapping("/{id}")
    @ResponseStatus(HttpStatus.NO_CONTENT)
    public void supprimer(@PathVariable long id, @RequestHeader(name = "If-Match", required = false) String etag) { }

    // Fichiers : upload multipart + telechargement streaming (le Content-Disposition fait le nom)
    @PostMapping(path = "/{id}/illustration", consumes = MediaType.MULTIPART_FORM_DATA_VALUE)
    public void illustration(@PathVariable long id, @RequestPart MultipartFile fichier) throws IOException {
        service.enregistrerIllustration(id, fichier.getOriginalFilename(), fichier.getSize(), fichier.getInputStream());
    }
    @GetMapping(path = "/{id}/illustration", produces = MediaType.IMAGE_PNG_VALUE)
    public ResponseEntity<InputStreamResource> image(@PathVariable long id) {
        return ResponseEntity.ok().contentType(MediaType.IMAGE_PNG).eTag("\"" + service.etag(id) + "\"")
                .cacheControl(CacheControl.maxAge(Duration.ofDays(7)).cachePublic()).body(new InputStreamResource(service.flux(id)));
    }
    // N\'egociation de contenu : @RequestMapping(produces = "text/csv") ou un HttpMessageConverter maison.
}

// La validation : annotations + groupes + validation croisee
public record CreerArticle(
        @NotBlank @Size(min = 3, max = 160) String titre,
        @NotBlank @Size(max = 100_000) String corps,
        @Size(max = 8) List<@Pattern(regexp = "[a-z0-9\\-]{1,30}") String> tags,
        @Nullable @Future Instant publieLe) { }
@Target({ FIELD, TYPE }) @Retention(RUNTIME) @Constraint(validatedBy = SlugUnique.class)
@interface SlugUnique { String message() default "slug deja pris" ; Class<?>[] groups() default { } ; Class<? extends Payload>[] payload() default { } }
class SlugUnique implements ConstraintValidator<SlugUnique, String> {
    @Autowired transient ArticleRepository repo;              // seul endroit ou @Autowired est tolere
    public boolean isValid(String slug, ConstraintValidatorContext c) { return slug == null || !repo.existsBySlug(slug); }
}

// Erreurs : une classe, un format, partout. (Boot 3 expose deja /error en ProblemDetail.)
@RestControllerAdvice
class GlobalExceptionHandler {
    private static final Logger log = LoggerFactory.getLogger(GlobalExceptionHandler.class);

    @ExceptionHandler(MethodArgumentNotValidException.class)
    ProblemDetail surValidation(MethodArgumentNotValidException e) {
        var pd = ProblemDetail.forStatusAndDetail(HttpStatus.UNPROCESSABLE_ENTITY, "Donnees invalides");
        pd.setTitle("Validation"); pd.setProperty("erreurs", e.getFieldErrors().stream()
                .collect(Collectors.toMap(FieldError::getField, f -> defaultString(f.getDefaultMessage()), (a, b) -> a, LinkedHashMap::new)));
        return pd;
    }
    @ExceptionHandler(Introuvable.class)
    ProblemDetail surIntrouvable(Introuvable e) {
        var pd = ProblemDetail.forStatusAndDetail(HttpStatus.NOT_FOUND, e.getMessage()); pd.setTitle("Introuvable"); return pd;
    }
    @ExceptionHandler(ErreurDomaine.class)
    ResponseEntity<ProblemDetail> surDomaine(ErreurDomaine e) { /* map code -> statut */ }
    @ExceptionHandler(Exception.class)
    ProblemDetail surImprevu(Exception e, HttpServletRequest req) {
        log.error("non ger\'e uri={}", req.getRequestURI(), e);
        var pd = ProblemDetail.forStatusAndDetail(HttpStatus.INTERNAL_SERVER_ERROR, "Erreur interne");
        pd.setProperty("trace", currentTraceId());            // un id a donner au support, pas une stack
        return pd;
    }
    // Les convertisseurs d'URL/JSON (MethodArgumentTypeMismatch, HttpMessageNotReadable) merite
    // aussi un handler : sinon le client recoit un 400 vide et inexplicable.
}
\end{codejava}

\section{Spring Data~: les repos qui font le travail}

\begin{codejava}
public interface ArticleRepository extends JpaRepository<ArticleEntity, Long>, JpaSpecificationExecutor<ArticleEntity> {
    Optional<ArticleEntity> findBySlug(String slug);
    List<ArticleEntity> findByStatutOrderByPublieLeDesc(Statut statut, Pageable page);
    @Query("""
           select a from ArticleEntity a left join fetch a.tags
           where (:statut is null or a.statut = :statut)
             and (:depuis is null or a.publieLe >= :depuis)
             and lower(a.titre) like lower(concat('%', :q, '%'))
           order by a.publieLe desc
           """)
    Page<ArticleEntity> rechercher(@Param("q") String q, @Param("statut") Statut statut,
                                   @Param("depuis") Instant depuis, Pageable pageable);
    @Query(value = "select * from articles where to_tsvector('french', titre || ' ' || corps) @@ plainto_tsquery('french', :q)",
           countQuery = "select count(*) from articles where ...", nativeQuery = true)
    Page<ArticleEntity> pleineTexte(@Param("q") String q, Pageable pageable);
    @Modifying(clearAutomatically = true, flushAutomatically = true)
    @Query("update ArticleEntity a set a.vues = a.vues + 1 where a.id = :id")
    int incrementerVues(@Param("id") long id);
    boolean existsBySlug(String slug);
    @EntityGraph(attributePaths = { "tags", "auteur" }) List<ArticleEntity> findTop10ByStatutOrderByVuesDesc(Statut s);
}
// Projection fermee (Hibernate ne charge que ces colonnes) :
interface Entete { Long getId(); String getTitre(); String getAuteurNom(); }
@Query("select a.id as id, a.titre as titre, u.nom as auteurNom from ArticleEntity a join a.auteur u")
Page<Entete> entetes(Pageable pageable);
// Transactions : @Transactional au niveau service, readOnly = true sur les lectures ;
// PROPAGATION_REQUIRES_NEW rarement, NEVER pour "pas de transaction ici".
// Le piege : @Transactional appelee en interne (this.methode()) ne passe pas par le proxy ->
// solution : injecter soi-meme (self), ou separer les classes.
// Auditing : @EnableJpaAuditing + @CreatedDate/@LastModifiedBy sur des champs @CreatedBy.
// Specifications ( predicats composables ) pour les filtres dynamiques ; Criteria API si tu es seul juge.
// Spring Data JDBC (sans Hibernate, sans proxy) merite un coup d'oeil : souvent suffisant.
\end{codejava}

\section{Spring Security~: tout est un filtre}

\begin{codejava}
@Configuration
@EnableWebSecurity
@EnableMethodSecurity                                    // @PreAuthorize / @PostAuthorize
class SecurityConfig {
    @Bean
    SecurityFilterChain api(HttpSecurity http, JwtDecoder decoder, Limitateur limitateur) throws Exception {
        http.csrf(csrf -> csrf.ignoringRequestMatchers("/api/**").csrfTokenRepository(CookieCsrfTokenRepository.withHttpOnlyFalse()))
            .cors(Customizer.withDefaults())
            .securityContext(sc -> sc.requireExplicitSave(false))
            .authorizeHttpRequests(auth -> auth
                .requestMatchers(HttpMethod.GET, "/api/articles/**", "/health", "/assets/**").permitAll()
                .requestMatchers("/api/auth/**", "/v3/api-docs/**", "/swagger-ui/**").permitAll()
                .requestMatchers("/actuator/health/**", "/actuator/info").permitAll()
                .requestMatchers("/actuator/**").hasRole("OPS")
                .requestMatchers("/api/admin/**").hasRole("ADMIN")
                .anyRequest().authenticated())
            .oauth2ResourceServer(o -> o.jwt(j -> j.jwtAuthenticationConverter(rolesConverter())))
            .sessionManagement(s -> s.sessionCreationPolicy(SessionCreationPolicy.STATELESS))
            .exceptionHandling(e -> e.authenticationEntryPoint(new BearerTokenAuthenticationEntryPoint())
                                    .accessDeniedHandler(new BearerTokenAccessDeniedHandler()))
            .headers(h -> h.contentSecurityPolicy(csp -> csp.policyDirectives("default-src 'self'; frame-ancestors 'none'")))
            .requiresChannel(c -> c.requestMatchers("/api/auth/**").requiresSecure())
            .addFilterBefore(limitateur, RateLimitFilter.class)
            .logout(l -> l.logoutUrl("/api/auth/logout").logoutSuccessHandler((q, r, a) -> r.setStatus(204)));
        return http.build();
    }
    @Bean PasswordEncoder encoder() {
        return Argon2PasswordEncoder.defaultsForSpringSecurity_v5_8();   // ou BCryptPasswordEncoder(12)
    }
    @Bean AuthenticationManager authManager(List<AuthenticationProvider> providers) { return new ProviderManager(providers); }
    // Pour un resource server OIDC : spring.security.oauth2.resourceserver.jwt.issuer-uri suffit ;
    // un decoder sur mesure se declare par JwtDecoders.fromIssuerLocation(URI).
    // rolesConverter() : JWT -> GrantedAuthority en lisant un claim ("roles" ou "groups")
    @Bean
    CorsConfigurationSource cors() {
        var c = new CorsConfiguration();
        c.setAllowedOrigins(List.of("https://blog.exemple.fr"));
        c.setAllowedMethods(List.of("GET", "POST", "PATCH", "DELETE", "OPTIONS"));
        c.setAllowedHeaders(List.of("Authorization", "Content-Type", "If-Match"));
        c.setAllowCredentials(true); c.setMaxAge(3600L);
        var src = new UrlBasedCorsConfigurationSource(); src.registerCorsConfiguration("/api/**", c); return src;
    }
}
// Methode : @PreAuthorize("hasRole('ADMIN') or @services.proprietaire(#id, authentication)")
// Form login + sessions (si c'est un site, pas une API) : formLogin(...).loginProcessingUrl("/login")
//   + CsrfTokenRequestAttributeHandler ; password manager avec DelegatingPasswordEncoder ({argon2}...)
//   => le prefixe dans la base choisit l'algo : c'est ainsi qu'on fait migrer les vieux hashes.
// OAuth2 client (se connecter avec Google/GitHub) : spring.security.oauth2.client.registration.*
// Les pieges : CSRF active par defaut sur les mutations ; un permitAll oublie sur /actuator/env
//   (fuite de secrets) ; hasRole vs hasAuthority (le prefixe ROLE_) ; un AuthenticationManager
//   chiffre a la main dans le controleur au lieu de deleguer.
\end{codejava}

\section{Tester Spring~: les quatre niveaux}

\begin{codejava}
@ImportTestcontainers(PgTest.class)
@SpringBootTest(webEnvironment = WebEnvironment.RANDOM_PORT)
@ActiveProfiles("test")
@AutoConfigureMockMvc
class ArticleApiIT {
    @Autowired MockMvc mvc;
    @Autowired ArticleRepository repo;
    @LocalServerPort int port;

    @Test @WithMockUser(roles = "AUTEUR")
    void creer_un_article_repond_201_et_le_rend() throws Exception {
        var corps = """
                {"titre":"Bien demarrer","corps":"...","tags":["java"]}""";
        var res = mvc.perform(post("/api/articles").contentType(APPLICATION_JSON).content(corps))
                .andExpect(status().isCreated())
                .andExpect(jsonPath("$.titre").value("Bien demarrer"))
                .andReturn();
        var location = res.getResponse().getHeader("Location");
        mvc.perform(get(location)).andExpect(status().isOk()).andExpect(jsonPath("$.statut").value("BROUILLON"));
    }
    @Test void validation_renvoie_les_champs_en_erreur() throws Exception {
        mvc.perform(post("/api/articles").contentType(APPLICATION_JSON).content("{\"titre\":\"x\"}"))
           .andExpect(status().isUnprocessableEntity())
           .andExpect(jsonPath("$.erreurs.titre").exists());
    }
    @Test void non_authentifie_renvoie_401() throws Exception {
        mvc.perform(get("/api/mes-articles")).andExpect(status().isUnauthorized());
    }
    // Les tranches : @WebMvcTest(ArticleController.class) (service mocke) ; @DataJpaTest (repo, H2 ou
    // conteneur) ; @JsonTest (serialisation) ; @RestClientTest (appels sortants, MockRestServiceServer).
    // Reseaux sortants : @MockBean/Mockito sur le RestClient, ou WireMock ; jamais d'Internet en CI.
    // Doc genere : springdoc-openapi + "mvn test" qui verifie que chaque route est decrite.
    // Charge : Gatling (scala/java DSL) ou k6 -- un scenario de 30 s a 500 req/s, avec p95 < 250 ms.
}
\end{codejava}

\section{Actuator, natif, et les limites du cadre}

\begin{terminal}
curl localhost:8080/actuator/health      # {"status":"UP"} + db, ping, diskSpace (si autorise)
curl localhost:8080/actuator/info
curl -s localhost:8080/actuator/metrics/http.server.requests | jq .
curl localhost:8080/actuator/metrics/http.server.requests?tag=uri:/api/articles   # details
curl localhost:8080/actuator/env ; curl localhost:8080/actuator/configprops        # A NE JAMAIS exposer
curl localhost:8080/actuator/scheduledtasks ; curl localhost:8080/actuator/threaddump
curl -X POST localhost:8080/actuator/loggers/com.exemple -d '{"configuredLevel":"DEBUG"}' -H 'Content-Type: application/json'
# En prod : exposer health, info, prometheus seulement ; proteger le reste (hasRole('OPS') ci-dessus).
# Boot 3 : /actuator/health/liveness et /readiness -> Kubernetes sonde dessus.
\end{terminal}

\begin{codeconf}
<!-- Image native : Boot 3 genere les hints, mais c'est un contrat (pas de reflexion sauvage) -->
<plugin><groupId>org.graalvm.buildtools</groupId><artifactId>native-maven-plugin</artifactId></plugin>
<profile><id>native</id><build><plugins>
  <plugin><groupId>org.springframework.boot</groupId><artifactId>spring-boot-maven-plugin</artifactId>
    <configuration><image><builder>paketobuildpacks/builder-jammy-base:latest</builder></image></configuration></plugin>
</plugins></build></profile>
\end{codeconf}

\begin{terminal}
./mvnw -Pnative native:compile          # ~5 min, RAM genereuse necessaire
du -h target/blog ; ./target/blog &     # demarrage en ~0,1 s ; ~60 Mo de RSS
./mvnw -Dspring.aot.enabled=true test   # verifier le AOT avant de builder
./mvnw processAotResources              # les hints de reflexion generes
\end{terminal}

\begin{tablez}[p{0.3\textwidth}X]
  \headrow Situations o\`u Spring est le bon choix \& Pourquoi \\
  \midrule
  application d'entreprise, \'equipe de 3 \`a 30 \& tout est standardis\'e, la revue et l'onboarding sont faciles \\
  besoin de Security/Data/Transaction/AMQP/Kafka \& ces briques sont les mieux finies du march\'e JVM \\
  int\'egrations (LDAP, OAuth2, JMS, mail, ws) \& starters + docs + r\'eponses Stack Overflow \\
\end{tablez}
Et l'inverse~: un microservice \`a 4 routes et 2 requ\^etes SQL, un CLI, une fonction serverless \`a
d\'emarrage critique, une \'equipe qui veut z\'ero magie~: un Javalin/Helidon + Hikari + Flyway bien
rang\'es fait le m\^eme travail avec 5 fois moins de d\'ependances. Spring n'est pas une obligation~;
c'est un pari sur l'\'ecosyst\`eme.

\begin{exercice}
\begin{enumerate}
  \item Reprends le projet du chapitre~10 (API Javalin) et porte-le sur Spring Boot~: la m\^eme API,
        les m\^emes tests. Note ce que Boot t'a \'evit\'e, et ce que Boot t'a impos\'e.
  \item Ajoute Flyway et une entit\'e JPA~: puis active \verb|ddl-auto=validate| et v\'erifie que Boot
        refuse de d\'emarrer si tu retires une colonne \`a la main~: c'est le filet que beaucoup n'ont jamais.
  \item Mets la s\'ecurit\'e en place~: resource server OIDC contre Keycloak en conteneur (le
        \verb|docker-compose| officiel fait 20 lignes)~: puis ajoute \verb|@PreAuthorize| et un test
        par r\^ole (AUTEUR peut modifier son article, pas celui du voisin).
  \item Fais tourner les 4 tranches de test sur la m\^eme application~: \verb|@WebMvcTest|,
        \verb|@DataJpaTest|, \verb|@SpringBootTest|, et un test de contr\^ole d'architecture~: compare
        les temps d'ex\'ecution et d\'ecide quoi garder dans la CI bloquante.
  \item Configure l'observabilit\'e~: actuator + Prometheus + Grafana en local (compose), puis la
        sonde Kubernetes de lecture~; casse volontairement la BDD et observe le passage \`a
        \verb|/readiness| en \verb|OUT\_OF\_SERVICE|.
  \item Compile en natif~: la plupart du temps il faudra ajouter
        \verb|@RegisterReflectionForBinding(...)| ou hint de Jackson~: documente ce que la magie
        t'avait \'epargn\'e. Compare RSS, d\'emarrage et d\'ebit apr\`es \'echauffement avec le mode JVM.
\end{enumerate}
\end{exercice}

\begin{ressources}
\begin{itemize}
  \item Spring Boot Reference --- \url{https://docs.spring.io/spring-boot/index.html}
  \item Guides Spring (les ~90 \og ~Getting Started~ \fg{}) --- \url{https://spring.io/guides}
  \item Spring Security architecture --- \url{https://docs.spring.io/spring-security/reference/servlet/architecture.html}
  \item Spring Data JPA --- \url{https://docs.spring.io/spring-data/jpa/reference/} --- springdoc-openapi --- \url{https://springdoc.org/}
  \item Testcontainers + Boot --- \url{https://docs.spring.io/spring-boot/reference/testing/testcontainers.html}
  \item AOT/Natif --- \url{https://docs.spring.io/spring-boot/reference/aot/} --- Start --- \url{https://start.spring.io/}
\end{itemize}
\end{ressources}

\index{Spring Boot}\index{auto-configuration}\index{Spring Data}\index{Spring Security}\index{actuator}

\part{Aller plus loin}
% ============================================================================
\chapter{Cinq projets complets, du cahier des charges \`a la mise en ligne}
\label{ch:projets}
% ============================================================================

Cinq projets \`a faire \emph{en entier}~: sp\'ecification, d\'ecoupage, code, tests, s\'ecurit\'e,
d\'eploiement, README. Chacun indique les chapitres n\'ecessaires~; si tu bloques, c'est qu'un
pr\'erequis est \`a relire, pas que le projet est trop dur.

\begin{tablez}[p{0.19\textwidth}p{0.115\textwidth}p{0.12\textwidth}X]
  \headrow Projet \& Langage \& Dur\'ee \& Ce qu'il t'apprend vraiment \\
  \midrule
  \thd{P1} Gestion de t\^aches \& JS pur \& 8--14 h \& DOM, \'etat, events, \verb|localStorage|, build Vite, sans framework \\
  \thd{P2} Chat temps r\'eel \& Node + JS \& 20--30 h \& HTTP/WS, flux, \verb|WebSocket|, files, auth par session, d\'eploiement \\
  \thd{P3} Blog complet \& PHP + MySQL \& 30--45 h \& MVC maison, requ\^etes pr\'epar\'ees, auth, upload, admin, SEO, s\'ecurit\'e \\
  \thd{P4} Raccourcisseur d'URL \& Java + Spring \& 25--40 h \& API REST, JPA, rate-limit, tests d'int\'egration, Docker \\
  \thd{P5} Mini-ERP polyglotte \& les trois \& 60 h+ \& fronti\`eres propres~: API Java, back-office PHP, front JS, contrats \\
\end{tablez}

\section{P1 --- Gestionnaire de t\^aches (JavaScript, z\'ero framework)}

\etape{Cahier des charges}
Liste de t\^aches, plusieurs projets (\og ~pays~ \fg{} de la liste), \'edition en place, filtres
(toutes / actives / termin\'ees), recherche, tri (manuel par glisser, puis date), \'etiquettes,
\'ech\'eance avec coloration, archivage, undo (Ctrl+Z sur la derni\`ere action), export/import JSON,
mode hors-ligne complet, th\`eme sombre.

\etape{Contraintes techniques}
\begin{itemize}
  \item Aucun framework~: modules ES, \verb|<script type="module">|, Vite pour le build.
  \item Architecture en trois couches~: \verb|etat.js| (le mod\`ele + les r\'educteurs),
        \verb|rendu.js| (du mod\`ele au DOM), \verb|events.js| (du DOM au mod\`ele). Aucune fonction
        du rendu ne mutante~; aucun \'ev\'enement qui touche directement au \verb|localStorage|.
  \item Persistance~: \verb|localStorage| avec un sch\'ema versionn\'e (\verb|{version: 2, taches: []}|)
        et une fonction de migration. \'Ecriture amortie (debounce 400 ms) et propre
        (\verb|beforeunload| non fiable~: \'ecrire \`a chaque mutation, sans attendre).
  \item Accessibilit\'e~: navigation clavier compl\`ete, \verb|aria-live| pour le compteur, labels
        r\'eels, focus pi\'eg\'e dans le dialogue d'ajout, contraste AA.
\end{itemize}

\etape{D\'ecoupage en 8 livrables testables}
\begin{enumerate}
  \item Squelette Vite + structure des modules + rendu d'une liste cod\'ee en dur.
  \item \'Etat central \verb|store| (\verb|subscribe|/ \verb|dispatch|)~; aucun DOM dans l'\'etat.
  \item Ajout/\'edition/suppression + formulaire accessible (\verb|<dialog>| natif).
  \item Filtres et recherche (URL hash ou \verb|searchParams| pour partager un filtre~! ).
  \item Bascule, tri manuel (glisser avec \verb|draggable| + \verb|aria|)~; puis raccourcis clavier.
  \item Persistance + migrations + gestion d'un stockage plein/interdit (try/catch, mode m\'emoire).
  \item Undo/redo (pile de snapshots born\'ee \`a 50)~; export/import JSON valid\'e.
  \item Build \verb|vite build|, Lighthouse \`a 100, tests Vitest sur le r\'educteur et les migrations,
        DPWA (\verb|manifest.webmanifest|) et publication sur Netlify/GitHub Pages.
\end{enumerate}

\etape{Crit\`eres d'acceptation (\`a cocher)}
Recharger la page ne perd aucune donn\'ee~; supprimer la derni\`ere t\^ache d'un projet vide ne casse
rien~; l'application fonctionne sans r\'eseau~; un fichier JSON corrompu au d\'emarrage est ignor\'e
avec un message~; le clique sur \og ~terminer~ \fg{} est possible au clavier~; Lighthouse Performance
\$\geq\$ 95 et Accessibilit\'e \$=\$ 100~; \verb|npm test| passe~; le README explique les 8 livrables.

\etape{Extensions}
Sync multi-onglets via \verb|BroadcastChannel|~; statistiques (t\^aches/jour, heatmap) dessin\'ees en
\verb|<canvas>|~; une PWA avec Service Worker et file d'attente hors-ligne~; un mode \og ~semestrier~ \fg{}
avec dates r\'ecurentes (\`a faire \emph{avec} \verb|Intl| et jamais avec des millisecondes brutes).

\section{P2 --- Chat temps r\'eel (Node.js)}

\etape{Sp\'ecification}
Salons multiples, inscription/connexion, pseudonyme, statut en ligne, envoi de message, historique
pagin\'e de 50 messages, indicateurs \og ~en train d'\'ecrire~ \fg{}, pr\'esence, reactions, coupures r\'eseau
reconnect\'ees avec rejeu, mod\'eration (ignorer/bannir), salon g\'en\'eral en lecture seule.

\etape{Choix assum\'es}
Serveur HTTP \verb|node:http| ou \verb|express|~; \verb|ws| (pas socket.io) pour comprendre le
protocole~; \verb|better-sqlite3| ou PostgreSQL~; \verb|jose| ou \verb|iron-session| pour la session~;
\verb|zod| ou un validateur maison pour chaque trame re\c{c}ue~; \verb|pino| pour les logs structur\'es.
Frontend en JS pur (le m\^eme socle que P1) pour que le serveur reste le seul cerveau.

\etape{Les 4 difficult\'es p\'edagogiques (ne les \'evite pas)}
\begin{enumerate}
  \item \textbf{Protocole sur flux~:} un message WebSocket est un paquet~; tu dois cadrer
        \verb|{type, id, payload}|, valider, r\'epondre par \verb|ack|~; et g\'erer les trames
        trop grosses (limite 64 Ko~: refuser, journaliser, d\'econnecter proprement).
  \item \textbf{Reconnexion~:} garder la file des messages non acquitt\'es~; au rejeu, le client
        envoie son dernier \verb|seq|~; le serveur renvoie le delta~: test avec
        \verb|navigator.onLine| forc\'e et red\'emarrage du serveur.
  \item \textbf{Pr\'esence~:} un \verb|ping|/ \verb|pong| toutes les 25 s avec purge des sockets
        morts~; diffusion \og ~\'ecrit\ldots~ \fg{} throttl\'ee \`a 1 par seconde et 500 ms de silence.
  \item \textbf{Isolation~:} un utilisateur ne lit que son salon~; une requ\^ete HTTP ne doit jamais
        lister les messages d'un autre~: \'ecris le test qui prouve la fuite, avant l'impl\'ementation.
\end{enumerate}

\etape{S\'ecurit\'e (liste non n\'egociable)}
Hachage Argon2id/bcrypt, cookies \verb|\_\_Host-| \verb|Secure| \verb|HttpOnly| \verb|SameSite=Lax|,
CSRF sur les routes HTTP mutantes, limitation de connexion (5/min/IP + verrouillage progressif),
longueur maximale des messages, \'echappement HTML au rendu (\verb|textContent|, jamais
\verb|innerHTML| sur du contenu utilisateur), refus des origines de WebSocket non list\'ees
(v\'erifie \verb|Origin|~! ), gzip uniquement sur les r\'eponses > 1 Ko, \verb|helmet|-like headers.

\etape{Livrables}
\verb|README| avec diagramme des trames~; sch\'ema SQL comment\'e~; \verb|docker-compose.yml|~;
tests (unitaires sur le parsing de trames, int\'egration sur la reconnexion)~; un script
\verb|load.js| (1 000 clients, 20 msg/s) et la mesure du d\'ebit~; le d\'eploiement derri\`ere un reverse
proxy (Caddy ou Nginx) avec TLS~; une note d'une page sur ce que tu changerais avec Redis
en cas de multi-instances.

\section{P3 --- Blog complet en PHP nu + MySQL (le vrai projet d'embauche)}

\etape{Sp\'ecification}
C\^ot\'e lecteur~: liste pagin\'ee, article, recherche pleine texte, tags, flux RSS et sitemap,
commentaires avec mod\'eration, compte avec avatar. C\^ot\'e auteur~: r\'edaction avec autosave,
brouillons, publication programm\'ee, upload d'image redimensionn\'ee, statistiques de lecture.
C\^ot\'e admin~: gestion des utilisateurs et des r\^oles, mod\'eration, journal d'audit.

\etape{Architecture}
Un mini-MVC en 250 lignes~: un point d'entr\'ee unique \verb|public/index.php|, un routeur,
des contr\^oleurs fins, une couche d\'ep\^ot en PDO, un conteneur de services minuscule, un moteur de
template (\verb|extract| interdit~: une m\'ethode \verb|render(\$vue, \$donn\'ees)|), et
l'autoloader PSR-4 de Composer. Les composants maison~: \verb|Requ\^ete|, \verb|R\'eponse|,
\verb|S\'eance|, \verb|Authentification|, \verb|Validation|, \verb|Journal|, \verb|Cache| (fichiers).

\etape{Exigences de qualit\'e (le diff qui fait la diff\'erence en entretien)}
\begin{itemize}
  \item Toutes les requ\^etes pr\'epar\'ees~; aucune variable dans le SQL~; le pr\'efixe de table dans une
        constante~; \verb|SELECT| nommant les colonnes.
  \item \'Echappement par contexte~: \verb|htmlspecialchars| pour le HTML, \verb|rawurlencode| pour
        l'URL, jamais de \verb|echo| d'un contenu utilisateur non filtr\'e~; le CSP bloque l'inline.
  \item CSRF sur tout ce qui mute, rotation d'identifiant de s\'eance \`a la connexion, r\'eg\'en\'eration
        de l'empreinte d'argon si les param\`etres changent, verrouillage apr\`es 5 \'echecs,
        r\'einitialisation par jeton \`a usage unique stock\'e hach\'e.
  \item Upload~: \verb|getimagesize| + r\'e-encodage (jamais confiance \`a l'extension), taille limit\'ee,
        nom al\'eatoire, hors \verb|docroot| ou avec \verb|Content-Disposition|, et virus-scan branchable.
  \item Pagination par curseur au-del\`a de 10 000 lignes~; index justifi\'es par un \verb|EXPLAIN|.
  \item \verb|composer phpstan| niveau 8 sans erreur~; \verb|php-cs-fixer|~; \verb|PHPUnit| sur le
        domaine et les d\'ep\^ots~; \verb|exit 1| en CI si un seul test tombe.
  \item SEO propre~: URLs cananiques, \verb|<link rel=prev/next>|, balises OG, donn\'ees structur\'ees
        JSON-LD, temps de r\'eponse < 120 ms, images en \verb|srcset|.
\end{itemize}

\etape{Extensions}
\'Edition collaborative (verrou optimiste par colonne \verb|version|)~; recherche par \verb|FULLTEXT|
avec synonymes~; i18n (2 langues, table de traductions)~; API JSON consomm\'ee par un front JS~;
webhooks sortants sign\'es~; file d'attente en base pour l'envoi d'emails~; un
\verb|docker-compose| complet (php-fpm, nginx, mysql, maildev).

\section{P4 --- Raccourcisseur d'URL + analytique (Java + Spring Boot)}

\etape{Sp\'ecification}
Cr\'eer un raccourci (code court, base 62, collision g\'er\'ee par r\'eessai), redirection 302 avec
journalisation asynchrone, liens personnalis\'es, expiration, limites par compte (quota),
tableau de bord~: clics par jour, r\'ef\'erents, pays (par IP, sans stockage de l'IP compl\`ete~: hach\'ee
et tronqu\'ee), API publique avec jeton, suppression/r\'ecup\'eration, QR code optionnel.

\etape{Mod\`ele de donn\'ees et choix}
PostgreSQL~: \verb|liens(id bigserial, code text unique, cible text, compte\_id, cree\_le, expire\_le,
supprime bool)| et \verb|clics(id bigserial, lien\_id, jour date, referer\_domaine, pays char(2),
agent\_hash bytea)|~; agr\'egats dans \verb|clics\_jour| (d\'eclencheur ou traitement planifi\'e) pour que
le tableau de bord reste un \verb|SELECT| simple. Redis optionnel pour le compteur~: vidang\'e
p\'eriodiquement vers Postgres (la bonne discussion d'architecture~: exactitude contre d\'ebit).

\etape{Les points techniques qui font progresser}
\begin{enumerate}
  \item Le g\'en\'erateur de code~: base 62, entropie (un \verb|LongAdder| + hachage de \verb|nanoTime|,
        ou un compteur chiffr\'e)~; la v\'erification d'unicit\'e par \verb|INSERT ... ON CONFLICT| et non
        par un \verb|SELECT| suivi d'un \verb|INSERT| (course).
  \item Le chemin critique~: une seule requ\^ete~: \verb|SELECT id, compte\_id, actif FROM liens WHERE code=?|~;
        le log du clic part dans une \verb|BlockingQueue| consomm\'ee par un worker
        (jamais une \'ecriture par requ\^ete sur le m\^eme thread)~: puis le 302 sort.
  \item Le cache~: Caffeine devant Postgres avec TTL 60 s~; invalidation \`a la modification.
  \item La limitation~: jeton par compte et par IP (filtre Spring~; ou Bucket4j) avec
        \verb|Retry-After|~; mesure et rapport (\verb|/actuator/metrics|).
  \item Les tests~: \verb|@SpringBootTest| + Testcontainers~; un test de propri\'et\'e sur le
        codage/d\'ecodage~; un test de concurrence qui cr\'ee 500 liens en parall\`ele (aucun doublon).
  \item Le d\'eploiement~: image via buildpacks~; \verb|/readiness| branch\'e sur la BDD~: arr\^et gracieux~;
        et (bonus) l'image native avec le temps de d\'emarrage mesur\'e.
\end{enumerate}

\section{P5 --- Mini-ERP polyglotte (l'exercice d'architecture)}

Un entrep\^ot minimal~: articles, clients, commandes, stock. Une API Java (Spring Boot) porte la
coh\'erence m\'etier et le stock~; un back-office PHP (le squelette de P3) g\`ere les fiches et les
utilisateurs internes~; un front JS (Vite + TypeScript, ou un framework choisi apr\`es les
chapitres~JS) consomme l'API. Ce n'est pas un pr\'etexte \`a tout microserviciser~: c'est
l'exercice des \textbf{fronti\`eres}.

\etape{Les cinq d\'ecisions \`a \'ecrire noir sur blanc dans un \fichier{ADR-001.md}}
\begin{enumerate}
  \item \textbf{Contrat~:} \verb|openapi.yaml| est la source de v\'erit\'e~; il est versionn\'e~; le front
        g\'en\`ere ses types (ou l'inverse~: il est test\'e par un contrat~: \link{Dredd}{https://github.com/apiaryio/dredd},
        Schemathesis, ou \verb|atdgen|-like en Java).
  \item \textbf{Authentification~:} qui \'emet~? Qui v\'erifie~? (r\'eponse saine~: Spring ressource
        server + login g\'er\'e par PHP via un jeton sign\'e court~; r\'eponse honn\^ete~: Keycloak.)
  \item \textbf{Transactions~:} le stock et la commande sont dans la m\^eme base~: une transaction
        SQL suffit~: documente ce qui changerait si ce n'\'etait plus le cas (saga, idempotence).
  \item \textbf{Erreurs~:} un format unique (9457) chez les deux backend~; les codes m\'etier dans une
        table commune~; les traductions c\^ot\'e front.
  \item \textbf{D\'eploiement~:} un \verb|docker-compose| qui monte les trois + Postgres + un proxy~;
        un \verb|make up| qui fait tout~; un \verb|make check| qui lance les trois suites.
\end{enumerate}

\etape{Bar\`eme d'auto-\'evaluation (sur 20, sois cruel)}
Tests qui prouvent un comportement et non des lignes~: /4~; gestion des erreurs et des cas limites~:
/4~; s\'ecurit\'e (cha\^ines, secrets, autorisations)~: /4~; lisibilit\'e du code et de la structure~: /3~;
documentation (README, ADR, sch\'ema)~: /2~; exploitation (logs, m\'etriques, d\'eploiement reproductible)~: /3.
Un projet \`a 12/20 qui tourne vaut mieux qu'un 18/20 qui ne se lance que sur ta machine~: mais
vise 16 en ajoutant les trois derniers points (m\'etriques, journal, secret hors d\'ep\^ot).

\begin{note}
Comment mener ces projets sans abandonner~: d\'ecoupe en livrables d'une demi-journ\'ee~; committe \`a chaque
livraison~; \'ecris le test avant la correction d'un bug~; si tu blocs plus de
30 minutes, \'ecris le probl\`eme en une phrase (le \og ~canard en caoutchouc~ \fg{}~: 70 \% des fois tu
trouves en \'ecrivant)~; et termine. Une version moche mais en ligne, test\'ee par trois personnes,
vaut infiniment plus qu'une version parfaite dans un dossier \verb|projet-3-final-v2|.
\end{note}

\index{projet}\index{capstone}\index{architecture}
% ============================================================================
\chapter{Comparatif JS / PHP / Java, et comment construire ta carri\`ere avec}
\label{ch:cheminement}
% ============================================================================

Tu as trois langages dans les mains. Ce chapitre r\'epond aux trois questions qui reviennent~:
\emph{lequel pour quoi}~? \emph{comment les combiner}~? \emph{qu'est-ce qu'un employeur regarde
vraiment}~?

\section{Le comparatif, sans caricature}

\begin{tablez}[p{0.15\textwidth}p{0.28\textwidth}p{0.28\textwidth}X]
  \headrow \thd{Dimension} \& \thd{JavaScript} \& \thd{PHP} \& \thd{Java} \\
  \midrule
  Nature \& langage du navigateur \emph{et} serveur polyglotte \& langage de script serveur, orient\'e web \& langage \`a typage statique, compil\'e, JVM \\
  Ex\'ecution \& interpr\'et\'e/JIT (V8) ; un thread + boucle d'\'ev\'enements \& JIT \`a la demande, un process par requ\^ete (mod\`ele \og ~share-nothing~ \fg{}) \& bytecode compil\'e, JIT longue dur\'ee, multithread r\'eel \\
  Type \& dynamique, optionnellement TS \& dynamique, optionnellement strict/types \& statique, nominal, generics \\
  Points forts \& l'interface, le temps r\'eel, l'\'ecosyst\`eme npm, le full-stack d'un seul langage \& le d\'eploiement trivial, l'\'ecosyst\`eme CMS, le co\^ut d'entr\'ee, le TTFB \& la robustesse \`a l'\'echelle, les outils, la typing-detector, le long-running \\
  Points faibles \& outillage mouvant, fatigue \'ecosyst\'emique, single-thread \& incoh\'erences historiques, moins adapt\'e au long-running \& verbosit\'e, d\'emarrage, m\'emoire, cycles de release d'entreprise \\
  D\'eploiement \& statique (CDN) ou Node dans un conteneur \& \verb|rsync| + Nginx, ou serveur manag\'e \& conteneur/JAR, image native, JVM manag\'ee \\
  Co\^ut d'entr\'ee \& tr\`es bas \& tr\`es bas \& moyen (JDK, build, IDE, concepts) \\
  Mont\'ee en comp\'etence \& rapide puis exponentielle (tooling, frameworks) \& rapide puis plate puis raide (archi, tests, PHP moderne) \& lente puis continue (les fondations servent 15 ans) \\
\end{tablez}

\subsection{Performance---ce qu'il faut retenir des benchmarks}

\begin{itemize}
  \item \textbf{Une requ\^ete web se joue \`a 90 \% sur la base de donn\'ees et le r\'eseau}, pas sur le langage.
        Un PHP avec un bon index bat un Java avec un N+1.
  \item L\`a o\`u Java gagne~: les calculs longs en m\'emoire, les connexions maintenues (WebSocket,
        flux, caches applicatifs), la mise en cache JIT apr\`es \'echauffement~; et les threads
        virtuels qui rendent le mod\`ele bloquant concurrentiel face \`a Node.
  \item L\`a o\`u Node gagne~: le fan-out d'appels I/O (un agr\'egateur, un proxy BFF) et le co\^ut
        marginal d'une connexion tenue.
  \item L\`a o\`u PHP reste tr\`es bon~: le rendu de pages classiques~: aucun \'etat \`a maintenir, la
        latence est celle d'un script~; et son profil m\'emoire par requ\^ete est imbattable en
        mutualis\'e.
  \item Ne crois aucun tableau de chiffres sans~: version du runtime, warmup, concurrency, type de
        charge, mat\'eriel. \verb|web.dev|, \link{TechEmpower}{https://www.techempowerbenchmarks.org/}
        et \link{Bun vs Node}{https://benchmark.js.org/} sont des points de d\'epart, pas des verdicts.
\end{itemize}

\subsection{\'Ecosyst\`eme~: les briques \'equivalentes}

\begin{tablez}[p{0.16\textwidth}p{0.26\textwidth}p{0.26\textwidth}X]
  \headrow Domaine \& JS \& PHP \& Java \\
  \midrule
  Gestionnaire de paquets \& npm / pnpm \url{https://www.npmjs.com/} \& Composer \url{https://packagist.org/} \& Maven Central / Gradle \\
  Cadre web majeur \& React, Vue, Svelte, Next.js \& Laravel, Symfony, Slim \& Spring Boot, Quarkus, Micronaut \\
  ORM / donn\'ees \& Prisma, Drizzle, Kysely \& Eloquent, Doctrine, Illuminate/Database \& Hibernate, jOOQ, Jdbi \\
  Validation \& Zod, Yup \& Respect/Validation, Symfony Validator \& Jakarta Bean Validation (Hibernate Validator) \\
  Tests \& Vitest, Playwright, Testing Library \& PHPUnit, Pest, Behat \& JUnit 5, AssertJ, Testcontainers \\
  Mise en forme \& Prettier, ESLint \& PHP-CS-Fixer, PHPStan, Rector \& Spotless, Checkstyle, Error Prone \\
  Journalisation \& pino, winston \& Monolog \& SLF4J + Logback \\
  HTTP client \& fetch natif, axios \& Guzzle, curl, HttpClient de Symfony \& \verb|java.net.http|, OkHttp, RestClient de Spring \\
  Temps r\'eel \& ws, Socket.IO, SSE \& Ratchet, Swoole/FrankenPHP + Mercure \& Quarkus WebSocket, Netty \\
  File de t\^aches \& BullMQ, pg-boss \& Laravel Queues, Symfony Messenger \& Spring Batch, RabbitMQ/Kafka clients \\
  D\'eploiement \& Vercel/Netlify/Cloudflare, Docker \& PHP-FPM+Nginx, Platform.sh, Docker \& conteneurs, Kubernetes, Quarkus/Buildpacks \\
  Documentation \& TSDoc/JSDoc \& phpDocumentor \& Javadoc \\
\end{tablez}

\section{Quelle stack pour quel projet~?}

\begin{tablez}[p{0.24\textwidth}X]
  \headrow Type de produit \& Recommandation \\
  \midrule
  Site vitrine, blog, portfolio \& HTML/CSS + un SSG (Astro, Eleventy) ou un CMS PHP statufi\'e. Pas de SPA~: tu perds le SEO et tu gagnes de la complexit\'e. \\
  Site \'editorial avec r\'edacteurs \& WordPress/Drupal (m\'etier r\'eel) ou Laravel/Symfony + back-office sur-mesure si le besoin d\'eborde. \\
  SaaS CRUD multi-tenant \& PHP (Laravel/Symfony) \emph{ou} Java (Spring) pour le c\oe{}ur~; React ou Vue pour l'interface~; PostgreSQL pour tout le monde. \\
  Application tr\`es interactive (tableaux, \'editeurs, hors-ligne) \& JS/TS d'abord~; un back l\'eger (Node ou API Java) ; le navigateur tient l'\'etat. \\
  Temps r\'eel, collaboratif \& Node (ou Go/Java + WebSockets) + Redis pub/sub + une file de messages. \\
  Traitement par lots, int\'egrations d'entreprise, facturation \& Java~: fiabilit\'e, transactions, typage, outillage d'observabilit\'e. \\
  Traitement lourd de donn\'ees, ML, Android \& Python ou Java/Kotlin. \\
  MVP en 3 semaines pour valider \& Le langage que tu connais le mieux, un seul d\'ep\^ot, un seul d\'eploiement. \\
\end{tablez}

\subsection{Les combinaisons qui marchent vraiment}

\begin{enumerate}
  \item \textbf{API Java + front TS~:} le profil \og ~grand compte~ \fg{}. Tu parles JSON, OpenAPI, JWT,
        et tu es bilingue contrat/production.
  \item \textbf{Monolithe PHP + front l\'eger~:} le profil agence~/ produit~: un d\'ep\^ot, un d\'eploiement,
        un co\^ut d'h\'ebergement ridicule, un SEO parfait.
  \item \textbf{Full-stack JS (Next.js ou Node + SPA)~:} le profil startup~: une \'equipe, un langage,
        des it\'erations rapides~; \`a condition d'aimer maintenir l'outillage.
  \item \textbf{Le trio assum\'e~:} PHP pour les pages, Java pour les services critiques, JS pour tout ce
        qui vit dans le navigateur. C'est la r\'ealit\'e de beaucoup d'entreprises~: ton atout est
        d'\^etre la personne qui comprend les trois bords.
\end{enumerate}

\section{Construire une carri\`ere~: la partie peu romantique}

\etape{Les trois premiers emplois (ou les trois premi\`eres ann\'ees)}
\begin{itemize}
  \item Un portfolio de \textbf{3 projets termin\'es et en ligne} bat 15 tutos~: code lu par un
        humain~: un README avec d\'ecisions et limites~: un historique Git lisible.
  \item Apprends \`a \textbf{lire du code que tu n'as pas \'ecrit}~: c'est 60 \% du m\'etier. Prends un
        projet open source de 3 000 lignes et \'ecris-en une carte d'architecture~: \c{c}a se passe
        mieux avec un langage \`a typage statique (Java) ou typ\'e (TS) sous la main~: c'est un
        argument pour apprendre l'un des deux.
  \item \textbf{Les fondamentaux ne p\'erissent pas~:} HTTP, TLS, SQL et index, encodage,
        concurrence, m\'emoire (pile/monceau), complexit\'e algorithmique, mise en cache, files,
        idempotence, observabilit\'e. Les cadres passent~; ces dix mots restent.
  \item Un certificat n'a d'int\'er\^et que s'il valide une comp\'etence que tu peux montrer.
        En revanche~: un \verb|git log| propre~: une issue ferm\'ee avec reproduction~: une PR
        relue --- \c{c}a~: \c{c}a compte partout.
\end{itemize}

\begin{tablez}[p{0.2\textwidth}X]
  \headrow R\^ole \& Ce qu'on attend, et le chemin \\
  \midrule
  Int\'egrateur / frontend \& HTML/CSS/JS solides~: a11y~: outils de build~: un framework~; puis TS~: tests~: perf. PHP ou Node en second langage pour dialoguer avec le back. \\
  D\'eveloppeur web g\'en\'eraliste \& PHP ou Node c\^ot\'e serveur~: SQL \'ecrit \`a la main~: un framework choisi pour ses id\'ees~: d\'eploiement~: journaux. Java en troisi\`eme langue pour la rigueur. \\
  D\'eveloppeur back-end / services \& Java (ou Go)~: mod\'elisation~: transactions~: messagerie~: observabilit\'e~: s\'ecurit\'e~: chargement. TypeScript pour lire les fronts. \\
  DevOps/SRE \& les trois~; et surtout Linux~: r\'eseau~: conteneurs~: Terraform~: CI~: et le go\^ut des post-mortems. \\
  Freelance \& un p\'erim\`etre clair~: un contrat~: la capacit\'e \`a dire non~; et la sp\'ecialisation (ex. \og ~migrations PHP 5 \$\rightarrow\$ 8~ \fg{} ou \og ~perf WordPress~ \fg{}). \\
\end{tablez}

\etape{R\'emun\'eration et march\'e~: ordres de grandeur honn\^etes (2025--2026)}
Fourchettes indicatives pour un march\'e francophone~; elles varient fortement selon la ville, la
taille de la bo\^ite et la n\'egociation~: junior 30--40 k\texteuro{}~; interm\'ediaire 40--52 k\texteuro{}~;
senior 52--70 k\texteuro{}~: expert/architecte 70 k\texteuro{} et plus. Ce qui d\'ecale le curseur n'est
jamais \og ~je connais le framework \`a la mode~ \fg{}~: c'est la capacit\'e \`a faire tourner une d\'ecision
technique co\^uteuse~: \`a la d\'efendre~: et \`a la documenter. PHP paie historiquement moins bien que Java
\`a anciennet\'e \'egale~; l'\'ecart se resserre d\`es que tu fais du Laravel/Symfony en scale-up et que tu
toucher aux bases~: \`a la perf et \`a la prod.

\etape{Le rythme d'apprentissage qui tient sur dix ans}
\begin{enumerate}
  \item \textbf{Un projet, pas une s\'erie.} Tu apprends quand \c{c}a r\'esiste~; un tuto n'est pas une
        r\'esistance~; c'est du confort avec des onglets.
  \item \textbf{Un langage par saison.} 3 mois sur un~; 3 semaines sur le suivant pour transf\'erer~;
        puis un projet qui m\'elange les deux.
  \item \textbf{Relis tes corrig\'es dans 6 mois.} Tu dois trouver ton code d'il y a six mois
        \og ~presque correct mais un peu naus\'eabond~ \fg{}~; sinon tu n'as pas progress\'e.
  \item \textbf{\'Ecris.} Un billet de 500 mots par sujet dur~: c'est la meilleure consolidation
        qui existe~: et c'est un signal d'expert gratuit.
  \item \textbf{Lis les sources.} Le JEP qui introduit les threads virtuels~: la RFC 9110, le code
        de \verb|PDO| ou de \verb|V8|~: le commit qui corrige un bug que tu as eu~: tout cela se lit
        en une soir\'ee et te rend meilleur qu'une formation de deux jours.
  \item \textbf{Garde un carnet.} Les 10 commandes que tu retapes~: les 5 erreurs qui reviennent~:
        les 3 d\'ecisions d'archi et pourquoi~: c'est l'annexe~\ref{ch:antisecche} de ce livre, \'ecrit par toi.
\end{enumerate}

\begin{astuce}
Si tu ne devais garder qu'une r\`egle~: \textbf{sois la personne qui fait le lien}. Un d\'eveloppeur PHP
qui lit du TypeScript~: un Java qui comprend le navigateur~: un front qui sait expliquer un plan
d'ex\'ecution--- ces profils sont rares, et ce sont eux qu'on promeut. Ce livre t'a donn\'e les trois
bords~: exploite la jointure.
\end{astuce}

\index{comparatif}\index{carri\`ere}\index{stack}

\backmatter
% ============================================================================
\chapter*{Annexes~: antis\`eches, messages d'erreur et listes de contr\^ole}
\label{ch:antisecche}
\addcontentsline{toc}{chapter}{Annexes~: antis\`eches, messages d'erreur et listes de contr\^ole}
% ============================================================================

Ces pages sont faites pour \^etre imprim\'ees ou gard\'ees ouvertes dans un second \'ecran.
Rien ici n'est nouveau~: tout est dans le livre, juste l\`a o\`u ton \oe{}il en a besoin.

\section{JavaScript en une page}

\begin{codejs}
// Variables, types, egalite
let x = 1; const y = { a: 1 }; typeof y;               // "object" ; typeof null === "object"
Object.freeze(o); structuredClone(o);                   // lecture seule / copie profonde
x === 1; Object.is(NaN, NaN);                           // equality stricte / NaN-sur
// Fonctions, objets, classes
const add = (a, b) => a + b;
const { a: alias = 0, ...reste } = objet;
const [premier, ...queue] = tableau;
class A extends B { #prive = 0; static c = 1; get v() { return this.#prive; } #m() { } }
// Tableaux -- les 14 qui comptent
t.map(f).filter(p).reduce((acc, x) => acc + x, 0).find(p).findIndex(p).some(p).every(p)
 .sort((a, b) => a - b).toSorted(cmp).reverse().slice(d, f).at(-1).flat(Infinity)
 .join(", ").includes(v).indexOf(v).forEach(f).toReversed()   // toSorted/toReversed: immutables
// Chaines
s.includes(sub); s.padStart(8, "0"); s.split(/[,;]/).map(t => t.trim());
s.normalize("NFD").replace(/\p{Diacritic}/gu, "");      // retirer les accents
// Promises et async
await Promise.all(p); await Promise.allSettled(p); await Promise.any(p);
const ctrl = new AbortController(); fetch(u, { signal: ctrl.signal }); ctrl.abort();
// Modules, DOM, evenements
import defaut, { nom as alias } from "./f.js"; export default ...;
const el = document.querySelector("[data-x]");
el.classList.toggle("on", force); el.setAttribute("aria-expanded", "true");
el.insertAdjacentHTML("beforeend", html); el.dataset.x = "1";
el.addEventListener("click", h, { once: true, signal: ctrl.signal });
// Reseau et stockage
const r = await fetch(u, { method: "POST", headers: { "content-type": "application/json" },
  body: JSON.stringify(d), credentials: "include" }); r.ok; await r.json();
localStorage.setItem("k", JSON.stringify(v)); sessionStorage; indexedDB;
// Utilitaires
const d = new Date(); d.toISOString(); Intl.NumberFormat("fr-FR", { style: "currency", currency: "EUR" }).format(9.9);
structuredClone; queueMicrotask; setTimeout; requestAnimationFrame; performance.now();
// Debug : console.table(donnees), console.count(label), console.timeLog(), debugger;
\end{codejs}

\section{PHP en une page}

\begin{codephp}
<?php
declare(strict_types=1);
// Tablesau (array) = liste + hashmap ; les 16 qui comptent
$arr = ['a' => 1]; count($arr); isset($arr['a']); unset($arr['a']);
array_map(fn, $arr); array_filter($arr); array_values($arr); array_keys($arr);
array_merge($a, $b); + ; [...$a, ...$b]; in_array($v, $a, true);
array_key_exists('a', $a); sort($a); usort($a, fn); array_slice($a, 0, 3);
array_column($rows, 'nom', 'id'); array_combine($cles, $valeurs);
array_sum($n); array_unique($a); array_reverse($a); array_diff($a, $b);
// Chaines -- toujours mb_* avec de l'humain dedans
strlen($s); mb_strlen($s); mb_substr($s, 0, 10); mb_strtolower($s);
str_contains($s, 'x'); str_starts_with($s, 'x'); trim($s); sprintf('%.2f', $f);
number_format($n, 2, ',', ' '); str_pad($s, 8, '0', STR_PAD_LEFT); nl2br(e($s));
// Echappement par contexte
echo htmlspecialchars($txt, ENT_QUOTES | ENT_SUBSTITUTE, 'UTF-8');   // HTML  (ou e() de Twig/Blade)
echo rawurlencode($slug);                                            // URL
json_encode($d, JSON_UNESCAPED_UNICODE | JSON_THROW_ON_ERROR);       // JSON
// Fonctions, classes, typage
function f(int $a, ?string $b = null, string ...\$reste): array { }
$fn = fn(int $i): int => $i * 2;                                     // fleche : capture auto
class C implements I { use T; public function __construct(private readonly Id $id) { } }
enum Statut: string { case Publie = 'publie'; }
interface I { public function x(): void; }   abstract class A implements I { }
// Superglobales et requete
$_GET $_POST $_FILES $_COOKIE $_SESSION $_SERVER $_ENV   // php://input pour le JSON
session_start(); $_SESSION['csrf'] = bin2hex(random_bytes(32));
header('Location: ' . $url, true, 303); exit;             // apres un POST : PRG
// PDO
$pdo = new PDO($dsn, $u, $p, [PDO::ATTR_ERRMODE => PDO::ERRMODE_EXCEPTION,
  PDO::ATTR_DEFAULT_FETCH_MODE => PDO::FETCH_ASSOC, PDO::ATTR_EMULATE_PREPARES => false]);
$st = $pdo->prepare('SELECT * FROM t WHERE id = ?'); $st->execute([$id]); $ligne = $st->fetch();
$pdo->beginTransaction(); /* ... */ $pdo->commit();
// Dates
$d = new DateTimeImmutable('now', new DateTimeZone('Europe/Paris'));
$d->format('Y-m-d H:i'); $d->modify('+1 day'); $d->diff($autre)->days;
// Divers
filter_var($email, FILTER_VALIDATE_EMAIL); password_hash($m, PASSWORD_ARGON2ID);
password_verify($m, $h); random_int(1, 100); bin2hex(random_bytes(16));
is_file($f); file_get_contents($f); json_decode($j, false, 512, JSON_THROW_ON_ERROR);
// Erreurs : display_errors=On en dev, Off en prod ; log_errors=On ; error_reporting(E_ALL)
// Tests : vendor/bin/phpunit ; Qualite : php-cs-fixer, phpstan, psalm, rector, composer audit
\end{codephp}

\section{Java en une page}

\begin{codejava}
// Types et variables
int i = 0; long l = 9L; double d = .5; boolean b = true; char c = 'x'; String s = "a";
var v = 3;                                  // inference ; final f = ...; (non reassignable)
// Collections et boucles
List<String> l = new ArrayList<>(); Set<String> se = new HashSet<>(); Map<String, Integer> m = new HashMap<>();
List<String> imm = List.of("a", "b");       // immuable ; Map.of / Set.of / List.copyOf
for (String x : l) { } for (int j = 0; j < n; j++) { } while (c) { } do { } while (c);
l.sort(Comparator.comparing(String::length)); l.removeIf(x -> x.isEmpty());
// API utiles : List.of, Map.entry, String.join, String.format("a %s", b), "a".repeat(3), strip()
l.stream().filter(p).map(f).collect(Collectors.toList()); l.stream().toList();
m.forEach((k, v2) -> System.out.println(k + v2)); m.merge(k, 1, Integer::sum);
m.computeIfAbsent(k, x -> new ArrayList<>()).add(v);
// Streams : .map .filter .flatMap .sorted .limit .distinct .peek .min .max .sum .average
// .count .anyMatch .allMatch .findFirst .reduce(id, op) .collect(groupingBy | toMap | joining)
// .mapToInt(..).sum() .summaryStatistics() .boxed() .parallel() .unordered()
// Objets et null
Objects.equals(a, b); Objects.hash(x, y); Objects.requireNonNull(o, "msg");
Optional<T> o = Optional.ofNullable(x); o.map(..).filter(..).orElseGet(..).orElseThrow(..);
// Erreurs et entrees/sorties
try (var in = Files.newInputStream(p); var out = Files.newOutputStream(q)) { in.transferTo(out); }
catch (IOException e) { throw new UncheckedIOException(e); }
String t = Files.readString(p); List<String> ls = Files.readAllLines(p);
Files.writeString(p, t, StandardCharsets.UTF_8, StandardOpenOption.CREATE, StandardOpenOption.APPEND);
try (Stream<String> s2 = Files.lines(p)) { s2.filter(x -> !x.isBlank()).forEach(System.out::println); }
// Dates
Instant.now(); LocalDate.of(2026, 1, 31); ZonedDateTime.of(ldt, ZoneId.of("Europe/Paris"));
Duration.between(a, b).toSeconds(); Period.between(d1, d2); DateTimeFormatter.ofPattern("dd/MM/yyyy", Locale.FRENCH);
// Concurrence
ExecutorService pool = Executors.newFixedThreadPool(n);
try (var vx = Executors.newVirtualThreadPerTaskExecutor()) { vx.submit(() -> tache()); }
CompletableFuture.supplyAsync(() -> a(), pool).thenApply(x -> b).exceptionally(ex -> defaut).join();
var queue = new ArrayBlockingQueue<T>(cap); queue.put(t); queue.poll(1, TimeUnit.SECONDS);
// Outils
System.out.printf("%-8s %6.2f%n", s, d); Scanner sc = new Scanner(System.in, StandardCharsets.UTF_8);
// JVM usuelles : -Xmx512m -XX:MaxRAMPercentage=75 -Xlog:gc* -Dfile.encoding=UTF-8
// Build : ./mvnw -B verify | ./gradlew build test --no-daemon | java -jar app.jar
\end{codejava}

\section{Le tableau des messages d'erreur qui reviennent}

\begin{tablez}[p{0.3\textwidth}p{0.24\textwidth}X]
  \headrow \thd{Message (exact)} \& \thd{Cause racine} \& \thd{Le geste qui r\'epare} \\
  \midrule
  \verb|Cannot read properties of undefined (reading 'x')| \& objet absent \`a ce moment (typo, async, destructuration) \& remonter \`a la source~; guard-clause \verb|?.| ou valider l'\'etat~; \textbf{ne pas} mettre un \verb|try/catch| autour \\
  \verb|ReferenceError: x is not defined| \& variable hors port\'ee, script non charg\'e, \verb|type="module"| oubli\'e \& v\'erifier l'ordre des \verb|<script>| et le nom export\'e \\
  \verb|Cannot access 'x' before initialization| \& zone morte \verb|let|/ \verb|const|, ou d\'ependance circulaire d'un module \& d\'eclarer plus haut~; casser le cycle (extraire un module commun) \\
  \verb|unhandledRejection: ...| \& promesse non attendue dans une \verb|forEach|/ un handler \& \verb|for...of| + \verb|await|, ou \verb|Promise.all(map)| \\
  \verb|CORS policy: No 'Access-Control-Allow-Origin'| \& le serveur n'autorise pas l'origine~; ou \verb|withCredentials| sans \verb|Allow-Credentials| \& corriger le serveur (jamais le front)~; v\'erifier \verb|preflight| en 401 \\
  PHP \verb|Undefined array key "x"| \& cl\'e absente, ou nom de champ different du SQL \& \verb|??| / \verb|array\_key\_exists|~; var\_dump l\'eger de \verb|array\_keys| \\
  PHP \verb|Cannot use a scalar value as an array| \& \verb|\$x[0] = 1| alors que \verb|\$x| est un scalaire \& initialiser \verb|\$x = []| explicitement \\
  PHP \verb|SQLSTATE[23000]: Integrity constraint violation: 1062 Duplicate entry| \& unicit\'e~; ou deux clics sur le bouton \& \verb|INSERT ... ON DUPLICATE|/ \verb|upsert|~; idempotence c\^ot\'e requ\^ete \\
  PHP \verb|headers already sent| \& un espace/une BOM avant \verb|<?php|, ou un \verb|echo| pr\'ealable \& supprimer la sortie~; pas de \verb|?>| final dans un fichier de classe \\
  PHP \verb|Allowed memory size of ... bytes exhausted| \& boucle qui accumule, \verb|fetchall| d'un million de lignes \& \verb|LIMIT|~; curseur~;/streaming~; mesurer avant d'augmenter la limite \\
  Java \verb|cannot find symbol| \& import manquant, casse, m\'ethode inexistante \& \verb|Ctrl+Space| sur le type~; v\'erifier le package r\'eel (les 2 \verb|List| de Java) \\
  Java \verb|variable x might not have been initialized| \& variable locale non initialis\'ee (pas de d\'efaut) \& initialiser~: ou rendre le chemin d'affectation exhaustif \\
  Java \verb|unreported exception X; must be caught or declared| \& exception v\'erifi\'ee \& propag\'ee (\verb|throws|) ou convertie en exception m\'etier~: ne pas gober \\
  Java \verb|incompatible types: String cannot be converted to int| \& parsing absent \& \verb|Integer.parseInt| avec \verb|try/catch NumberFormatException| (ou \verb|Scanner.nextInt|) \\
  Java \verb|Non-resolvable import POM| / \verb|Could not resolve ...| \& d\'ependance introuvable, mauvais \verb|groupId|, r\'eseau, cache \& \verb|dependency:tree|~: chercher sur Maven Central~: \verb|-U|~: v\'erifier le miroir \\
  Java \verb|java: package javax.servlet does not exist| \& Spring 6 = Jakarta EE 9+ \& remplacer \verb|javax.*| par \verb|jakarta.*| \\
  Java \verb|LazyInitializationException: failed to lazily initialize| \& acc\`es hors transaction / hors session \& \verb|join fetch|~: projection~: ou mapper en DTO dans la requ\^ete \\
  Java \verb|OOM: Java heap space| / \verb|GC overhead limit| \& objet trop gros~: cache non born\'e~: \verb|-Xmx| ridicule \& mesurer avec un dump~: borner le cache (Caffeine)~: adapter le tas~: ne pas grossir \verb|-Xmx| d'abord \\
\end{tablez}

\section{Terminal, Git, HTTP---les 40 commandes utiles}

\begin{terminal}
# Navigation et fichiers
ls -la ; cd - ; pwd ; tree -L 2 ; cp -r a b ; mv a b ; rm -rf dossier ; mkdir -p x/y
find . -name "*.php" -not -path "./vendor/*" ; grep -rn --include="*.js" "motif" . | head
# Pipes et recherche
cmd | tee log.txt ; cmd 2>&1 | less ; xargs -n1 -P4 ; sed -i.bak 's/an 2024/an 2026/g' f
sort -u ; uniq -c | sort -rn | head ; wc -l ; du -sh * ; lsof -i :8080 ; kill -TERM <pid>
env | grep -i proxy ; curl -i -s -o /dev/null -w '%{http_code} %{time_total}s\n' URL ; dig +short domaine
tar czf a.tgz dossier/ ; tar xzf a.tgz ; unzip f.zip -d d/ ; chmod +x script
# Git : le quotidien
git status -sb ; git diff ; git add -p ; git commit -m "portee: sujet" ; git push -u origin HEAD
git log --oneline --graph --decorate -15 ; git switch -c sujet ; git switch main ; git pull --rebase
# Git : les 5 situations qui font paniquer (et leur sortie)
git restore --staged f ; git restore f ; git clean -fd            # annuler stage / fichier / non-suis
git commit --amend --no-edit                                       # corriger le dernier commit
git revert <sha> ; git reset --soft HEAD~1                         # annuler un commit publie / le defaire localement
git rebase -i origin/main ; :q pour abandonner un rebase ; git rebase --abort
git stash push -u -m "wip" ; git stash pop ; git stash list
git bisect start HEAD v1.0 ; git bisect run npm test ; git bisect reset   # trouver le commit fautif
git blame -L 40,60 f ; git log --oneline -- src/fichier ; git show --stat <sha>
git worktree add ../sujet sujet ; git worktree list                # deux branches en parallele
# Git : signature et propres historiques
git config --global core.editor "code --wait" ; git config --global pull.rebase true
git config --global init.defaultBranch main ; git config --global alias.lg "log --oneline --graph"
git filter-repo --invert-paths --path secret.env                   # purger un secret (puis changer le secret !)
# HTTP : comprendre en une minute
curl -v -X OPTIONS URL -H 'Origin: https://x' -H 'Access-Control-Request-Method: POST'
curl -H 'Authorization: Bearer '$T -s URL | jq . ; curl -I URL | grep -Ei 'cache|etag|set-cookie|security'
\end{terminal}

\section{Codes de statut~: le r\'eflexe en 12 lignes}

\begin{tablez}[p{0.1\textwidth}p{0.2\textwidth}X]
  \headrow Code \& Nom \& L'utiliser quand \\
  \midrule
  200 \& OK \& r\'eponse avec corps \\
  201 \& Created \& ressource cr\'e\'ee~: + \verb|Location| \\
  204 \& No Content \& suppression~: ou action sans corps \\
  301/302 \& Moved \& redirect permanent (SEO) / temporaire (apr\`es POST~: 303) \\
  400 \& Bad Request \& corps illisible, JSON cass\'e, param\`etre mal form\'e \\
  401 \& Unauthorized \& identit\'e absente ou expir\'ee (+ \verb|WWW-Authenticate|) \\
  403 \& Forbidden \& identit\'e connue~: permission absente \\
  404 \& Not Found \& ressource absente (ou cach\'ee volontairement pour ne pas la r\'ev\'eler) \\
  409 \& Conflict \& doublon, version obsol\`ete, \'etat incompatible \\
  413/415 \& Payload / Media \& fichier trop gros / type non accept\'e \\
  422 \& Unprocessable \& syntaxe correcte~: s\'emantique refus\'ee (validation m\'etier) \\
  429 \& Too Many Requests \& limitation d\'epass\'ee (+ \verb|Retry-After|) \\
  500/503 \& Internal / Unavailable \& bug~: ou panne d\'ependance (avec \verb|Retry-After|) \\
\end{tablez}

\section{Le dernier carr\'e~: la liste de contr\^ole avant de dire \og ~c'est fini~ \fg{}}

\begin{enumerate}
  \item[] Une requ\^ete SQL de la liste des 5 plus lentes a \'et\'e expliqu\'ee (\verb|EXPLAIN|)~; les index n\'ecessaires existent.
  \item[] Toute entr\'ee utilisateur est valid\'ee au serveur~: toute sortie est \'echapp\'ee selon son contexte.
  \item[] Aucun secret dans le d\'ep\^ot~; aucun \verb|display\_errors|~; aucun \verb|actuator/env| expos\'e~; les erreurs 500 ne fuient pas la stack.
  \item[] Les actions mutantes sont prot\'eg\'ees (CSRF ou jeton) et idempotentes au moins par erreur~: et test\'ees.
  \item[] Les en-t\^etes de base sont l\`a~: \verb|Content-Security-Policy|, \verb|X-Content-Type-Options|, cookies \verb|Secure| \verb|HttpOnly| \verb|SameSite|, \verb|Strict-Transport-Security| si TLS.
  \item[] Une requ\^ete lente, un appel sortant cass\'e, une base injoignable~: l'application d\'egrade proprement (timeout~: repli~: message clair).
  \item[] Les journaux contiennent un identifiant de corr\'elation~; on peut retrouver la panne d'un utilisateur en une commande.
  \item[] Les tests tournent en local \emph{et} en CI~: et la CI ne peut pas \^etre contourn\'ee par un \verb|--no-verify| silencieux.
  \item[] Le d\'eploiement est reproductible par une commande~; le retour arri\`ere aussi.
  \item[] Le README dit comment lancer~: tester~: configurer~: et pourquoi l'architecture est ce qu'elle est.
  \item[] Tu as supprim\'e le code comment\'e~: les variables \verb|temp2|~: et le \verb|console.log| du vendredi soir.
  \item[] Tu sais expliquer---\`a voix haute, sans l'IDE---chaque choix de la PR. Si tu ne sais pas~: ce n'est pas fini.
\end{enumerate}

\section{Outils~: la table de rappel}

\begin{tablez}[p{0.2\textwidth}p{0.24\textwidth}X]
  \headrow Besoin \& Choisir \& Lien \\
  \midrule
  \'Editeur \& VS Code + extensions (ESLint, Prettier, PHP Intelephense, Extension Pack for Java) \& \url{https://code.visualstudio.com/} \\
  IDE \& IntelliJ IDEA Community (Java) / PhpStorm (PHP, JS) \& \url{https://www.jetbrains.com/} \\
  JDK \& Temurin 21, g\'er\'e par SDKMAN ou mise \& \url{https://adoptium.net/} \\
  Node \& LTS via \verb|nvm| ou \verb|mise| \& \url{https://nodejs.org/fr} \\
  PHP \& 8.3+ via le PPA ondrej, \verb|mise| ou MAMP/Laragon \& \url{https://www.php.net/downloads.php} \\
  Base \& PostgreSQL 17 (ou MariaDB) + DBeaver \& \url{https://dbeaver.io/} \\
  Build \& Maven wrapper / Gradle wrapper \& \url{https://maven.apache.org/} \\
  HTTP client \& Bruno (open source) ou Insomnia \& \url{https://www.usebruno.com/} \\
  Conteneurs \& Docker Desktop / OrbStack (macOS) / Podman \& \url{https://www.docker.com/} \\
  Git GUI \& GitKraken free, GitHub Desktop, ou la ligne de commande \& \url{https://git-scm.com/} \\
  Proxy local \& Caddy (TLS automatique) ou Nginx \& \url{https://caddyserver.com/} \\
  D\'eploiement \& Netlify/Vercel (front), Fly.io/Railway/Coolify, VPS + Docker \& \url{https://fly.io/} \\
  Tests \& Vitest / Playwright~: PHPUnit~: JUnit + Testcontainers \& \url{https://vitest.dev/} \\
  Qualit\'e \& ESLint+Prettier~: PHPStan+CS-Fixer~: Spotless+Checkstyle+PIT \& \url{https://phpstan.org/} \\
  Observabilit\'e \& Grafana+Prometheus~: ou Sentry pour les erreurs \& \url{https://sentry.io/} \\
\end{tablez}

\begin{note}
Si tu ne devais garder qu'une page de tout le livre~: garde celle des messages d'erreur. Un bug
n'est pas une punition~: c'est un message. Les lire en entier~: mot pour mot~: est la comp\'etence
la plus rentable de ce m\'etier--- et celle qui distingue celui qui \og ~ne comprend rien au Java~ \fg{}
de celui qui a simplement appris \`a lire les diagnostics.
\end{note}

\index{antis\`eche}\index{messages d'erreur}\index{checklist}\index{Git}

\printindex
\end{document}
